% DRAFTING IN PROGRESS - Volume I Chapter 4 (2026-06-08)
% Authority: TOC/CH04-DRAFTING-KICKOFF-LOCK.md
% §4.1--§4.5 RE-DRAFT pass 3/3 (2026-06-08); §4.6 RE-DRAFT pass 3/3 (2026-06-08); §4.7--§4.8 RE-DRAFT pass 5/5 (2026-06-08); §4.9--§4.10 RE-DRAFT pass 3/3 (2026-06-08)
% Depth target: §4.1--§4.5 ≥5pp/sub per CH04-REDRAFT-DEPTH-LOCK.md
% Do not add Ch 5+ content to this file
\chapter{Source Realizability and Spectral Truncation}
\label{ch:04}
\setcounter{section}{-1}

\section{Introduction}
\label{sec:ch40}

Representation is coordinate freedom; realizability is geometry-imposed rank.
Once \(\mathcal{R}_\omega=\Pi_{\mathrm{rad}}\circ\mathcal{S}_\omega\) has selected the
outgoing radiative state and an admissible map \(\mathcal{C}\) has listed it in coefficients,
the remaining Volume~I question is not whether angular channels exist in
\(\mathcal{F}_{\mathrm{rad}}\), but which coefficient vectors finite impressed sources can
\emph{synthesize} and which finite instruments can \emph{certify}
\cite{stratton1941,harrington2001,devaney2004}. Finite support, finite aperture, and the
operator chains developed in Chapters~\ref{ch:02}--\ref{ch:03} compress that question into
spectral rank, coupling ceilings, and an effective electromagnetic dimension developed in the
sections that follow.


\subsection{Realizability and accessibility limits}
\label{sec:ch401}
% TOC tag:

A harmonic listing \(\{c_\alpha\}_{\alpha\in I}\) is mathematically legitimate only inside the
admissible representation class of Definition~\ref{def:ch3.R-adm}; beyond that certification,
however, coefficient space is not a free vector space. Bounded impressed currents
\(\Jimp\in\mathcal{J}_\Omega\) on compact support \(\Omega\) generate only those states
\(\mathfrak{F}=\mathcal{R}_\omega[\Jimp]\in\mathcal{F}_{\mathrm{rad}}\) permitted by the
radiation map of Chapter~\ref{ch:02}, and finite apertures with finite estimation rank recover
only those coordinates that survive the measurement composition of
Eq.~\eqref{eq:ch3.measurement-on-coefficients} \cite{stratton1941,harrington2001}. The
progressive filter of Hierarchy~\ref{hier:ch1.angular-status} therefore descends from
Maxwell-admissible fields to \emph{realizable synthesis sets} and \emph{accessibility
envelopes} in coefficient space---distinct geometric subsets of \(\mathbb{C}^{I}\) whose
separation is the organizing tension of the present chapter \cite{hansen1988,devaney2004}.

\paragraph{Assumptions.}
Fix \(\omega>0\) under the \(\exp(-\jj\omega t)\) convention, linear isotropic homogeneous
exterior regions unless stated otherwise, and outgoing closure on
\(\mathcal{F}_{\mathrm{rad}}\) as in Section~\ref{sec:ch24} \cite{stratton1941}. Let
\(\mathcal{C}:\mathcal{V}\to\mathbb{C}^{I}\) be a representation map with
\(\mathcal{C}\in\mathcal{R}_{\mathrm{adm}}\) and coefficient index set \(I\) as in
Eq.~\eqref{eq:ch3.representation-map} \cite{stratton1941,jackson1999}. Impressed currents lie
in the admissible source class \(\mathcal{J}_\Omega\) of Subsection~\ref{sec:ch201}; measurement
screens \(\mathcal{M}\) are finite rank as in Section~\ref{sec:ch29} \cite{harrington2001}. All
field states are discussed on \(\operatorname{ran}\mathcal{R}_\omega\) after application of
\(\Pi_{\mathrm{rad}}\) \cite{stratton1941}.

\paragraph{From radiative states to coefficient vectors.}
For admissible \(\Jimp\), the radiative state
\begin{equation}
\label{eq:ch4.state-from-source}
\mathfrak{F}
=
\mathcal{R}_\omega[\Jimp]
\end{equation}
carries the exported Maxwell pair before any angular label is attached
\cite{stratton1941,harrington2001}. The representation map assigns coordinates
\begin{equation}
\label{eq:ch4.coefficient-vector}
\mathbf{c}
=
\mathcal{C}(\mathfrak{F})
=
\{c_\alpha\}_{\alpha\in I},
\qquad
\mathbf{c}(\Jimp)
=
\mathcal{C}\big(\mathcal{R}_\omega[\Jimp]\big),
\end{equation}
with each \(c_\alpha\) a linear functional on \(\mathfrak{F}\) once a flux-normalized basis
\(\{\mathbf{f}_\alpha\}\) is fixed in Section~\ref{sec:ch35} \cite{hansen1988,stratton1941}.
Equation~\eqref{eq:ch4.coefficient-vector} is the coordinate restatement of
Eq.~\eqref{eq:ch3.state-from-source}: the operator chain fixes \emph{what} radiates; the listing
fixes \emph{how} the same object is indexed. Realizability and accessibility are predicates on
\(\mathbf{c}\), not on raw \(\mathcal{S}_\omega\) images \cite{devaney2004}.

\begin{definition}[Realizable synthesis set]
\label{def:ch4.realizable-synthesis-set}
For fixed \(\mathcal{C}\in\mathcal{R}_{\mathrm{adm}}\), the \emph{realizable synthesis set} is
\begin{equation}
\label{eq:ch4.realizable-synthesis-set}
\mathcal{S}_{\mathrm{real}}(\mathcal{C})
=
\left\{
\mathbf{c}\in\mathbb{C}^{I}:
\exists\,\Jimp\in\mathcal{J}_\Omega\ \text{with}\
\mathcal{C}\big(\mathcal{R}_\omega[\Jimp]\big)=\mathbf{c}
\right\}.
\end{equation}
\end{definition}

Physically, \(\mathcal{S}_{\mathrm{real}}(\mathcal{C})\) is the set of coefficient vectors that
\emph{finite impressed sources} in the declared model class can synthesize on
\(\operatorname{ran}\mathcal{R}_\omega\) \cite{stratton1941,harrington2001}. It is a
\emph{realizability class} in representation space: not every \(\mathbf{c}\in\mathbb{C}^{I}\)
compatible with outgoing closure is producible by a bounded \(\Jimp\), and non-radiating
contributions in \(\ker\mathcal{R}_\omega\) never enter the synthesis map at all
\cite{harrington2001}. Equation~\eqref{eq:ch4.realizable-synthesis-set} makes the
Hierarchy~\ref{hier:ch1.angular-status} predicate ``realizable'' precise at the coordinate layer
where laboratories actually design currents and read patterns \cite{hansen1988}.

\begin{definition}[Accessibility envelope]
\label{def:ch4.accessibility-envelope}
For fixed \(\mathcal{C}\in\mathcal{R}_{\mathrm{adm}}\) and a finite-rank measurement
\(\mathcal{M}\) with aperture rank \(N\), the \emph{accessibility envelope} is
\begin{equation}
\label{eq:ch4.accessibility-envelope}
\mathcal{E}_{\mathrm{acc}}^{(N)}(\mathcal{C})
=
\left\{
\mathbf{c}\in\mathbb{C}^{I}:
\exists\,\text{certified report}\
\widehat{\mathbf{c}}\ \text{with}\
\|\widehat{\mathbf{c}}-\mathbf{c}\|_{\mathcal{M}}\le \varepsilon_{\mathrm{cert}}
\right\},
\end{equation}
where \(\|\cdot\|_{\mathcal{M}}\) is the estimation norm induced by
\(\mathcal{M}\circ\mathcal{N}\circ\mathcal{C}\) in Eq.~\eqref{eq:ch3.measurement-on-coefficients}
and \(\varepsilon_{\mathrm{cert}}\) is the declared calibration tolerance
\cite{harrington2001,devaney2004}.
\end{definition}

The accessibility envelope is the coefficient-domain image of Principle~\ref{prin:ch1.accessibility}:
Nature may radiate organized angular content that no finite instrument certifies to within
\(\varepsilon_{\mathrm{cert}}\) \cite{hansen1988}. Algebraically,
\(\mathcal{E}_{\mathrm{acc}}^{(N)}(\mathcal{C})\) is the set of modal coordinates that survive
finite aperture rank, noise, and calibration screens applied \emph{after} synthesis
\cite{harrington2001}. It is not a relabeling of \(\mathcal{S}_{\mathrm{real}}(\mathcal{C})\): a
vector can lie in the synthesis set yet fall outside the envelope when \(N\) is too small, and
certified reports can exist only for combinations of modes that the instrument's rank allows to be
distinguished \cite{devaney2004}.

\begin{figure}[t]
\centering
\ChOneFigBlock{%
\begin{tikzpicture}[ch1fig, node distance=7mm and 10mm]
  \node[ch1box, minimum width=34mm] (ci) {$\mathbb{C}^{I}$\\coefficient space};
  \node[ch1box, minimum width=30mm, below left=10mm and 6mm of ci] (sr) {$\mathcal{S}_{\mathrm{real}}(\mathcal{C})$\\synthesis};
  \node[ch1box, minimum width=30mm, below right=10mm and 6mm of ci] (ea) {$\mathcal{E}_{\mathrm{acc}}^{(N)}(\mathcal{C})$\\certification};
  \node[ch1box, minimum width=26mm, below=14mm of sr] (j) {$\mathcal{J}_\Omega$\\bounded $\Jimp$};
  \node[ch1box, minimum width=26mm, below=14mm of ea] (m) {$\mathcal{M}$\\rank $N$};
  \draw[ch1flow] (j) -- node[left, font=\footnotesize] {$\mathcal{R}_\omega$} (sr);
  \draw[ch1flow] (sr) -- node[left, font=\footnotesize] {$\mathcal{C}$} (ci);
  \draw[ch1flow] (m) -- node[right, font=\footnotesize] {report $\widehat{\mathbf{c}}$} (ea);
  \draw[ch1flow, dashed] (ea) -- (ci);
\end{tikzpicture}%
}
\caption{Realizability acts on synthesis: bounded impressed currents map into
\(\mathcal{S}_{\mathrm{real}}(\mathcal{C})\subseteq\mathbb{C}^{I}\) through
\(\mathcal{R}_\omega\) and \(\mathcal{C}\). Accessibility acts on certification: finite-rank
\(\mathcal{M}\) defines the envelope \(\mathcal{E}_{\mathrm{acc}}^{(N)}(\mathcal{C})\) of
recoverable coordinates. The two subsets need not coincide.}
\label{fig:ch4.synthesis-accessibility-sets}
\end{figure}

Figure~\ref{fig:ch4.synthesis-accessibility-sets} fixes the coordinate-level picture used throughout
Chapter~\ref{ch:04}. Synthesis flows upward from \(\mathcal{J}_\Omega\); certification flows
from \(\mathcal{M}\) into the same ambient space \(\mathbb{C}^{I}\) without enlarging
\(\operatorname{ran}\mathcal{R}_\omega\) \cite{stratton1941,harrington2001}. Design problems seek
\(\mathbf{c}\in\mathcal{S}_{\mathrm{real}}(\mathcal{C})\); inverse problems seek stable estimates of
\(\mathbf{c}\in\mathcal{E}_{\mathrm{acc}}^{(N)}(\mathcal{C})\); both fail silently when their
target sets are treated as all of \(\mathbb{C}^{I}\) \cite{devaney2004}.

\begin{proposition}[Proper restriction in finite geometry]
\label{prop:ch4.proper-restriction}
For every nontrivial bounded \(\Omega\) and every flux-normalized \(\mathcal{C}\in\mathcal{R}_{\mathrm{adm}}\)
on outgoing states,
\begin{equation}
\label{eq:ch4.nested-restriction}
\mathcal{S}_{\mathrm{real}}(\mathcal{C})
\subseteq
\mathbb{C}^{I},
\qquad
\mathcal{E}_{\mathrm{acc}}^{(N)}(\mathcal{C})
\subseteq
\mathbb{C}^{I},
\end{equation}
with strict inclusion in both memberships for electrically small or aperture-limited geometries at
fixed \(\omega\) \cite{stratton1941,hansen1988}. Moreover, neither set is determined by
\(\mathcal{C}\) alone: \(\mathcal{S}_{\mathrm{real}}\) depends on \(\mathcal{J}_\Omega\) and
\(\mathcal{R}_\omega\); \(\mathcal{E}_{\mathrm{acc}}^{(N)}\) depends on \(\mathcal{M}\) and the
accessible modal sector of Eq.~\eqref{eq:ch2.accessible-modal-sector}
\cite{harrington2001,devaney2004}.
\end{proposition}
\begin{proof}
The inclusions Eq.~\eqref{eq:ch4.nested-restriction} follow from
Definitions~\ref{def:ch4.realizable-synthesis-set} and~\ref{def:ch4.accessibility-envelope}.
Strictness of \(\mathcal{S}_{\mathrm{real}}(\mathcal{C})\subsetneq\mathbb{C}^{I}\) follows because
\(\mathcal{R}_\omega\) is compact-operator limited on bounded \(\Omega\): only a finite-dimensional
sector of modal coefficients carries appreciable weight at fixed \(k_0 a\) when support is compact
\cite{stratton1941,jackson1999}. Strictness of
\(\mathcal{E}_{\mathrm{acc}}^{(N)}(\mathcal{C})\subsetneq\mathbb{C}^{I}\) follows from finite rank
of \(\mathcal{M}\) and the non-reconstructable sets of Subsection~\ref{sec:ch293}
\cite{harrington2001}. Dependence on \(\mathcal{R}_\omega\) and \(\mathcal{M}\) is immediate from
Eqs.~\eqref{eq:ch4.realizable-synthesis-set} and~\eqref{eq:ch4.accessibility-envelope}; dependence
on \(\mathcal{C}\) enters only through the coordinate chart, not through the physics of
\(\operatorname{ran}\mathcal{R}_\omega\) \cite{devaney2004}.
\end{proof}

Proposition~\ref{prop:ch4.proper-restriction} is the first quantitative statement of Volume~I's
central distinction: realizability limits \emph{generation}; accessibility limits
\emph{certification} \cite{hansen1988}. A coefficient vector may be synthesizable yet
uncertifiable when aperture rank is low; certifiable combinations may include modes that synthesis
cannot excite when geometry forbids the required coupling \cite{harrington2001,devaney2004}. The
remainder of Chapter~\ref{ch:04} replaces the abstract subsets
\(\mathcal{S}_{\mathrm{real}}(\mathcal{C})\) and \(\mathcal{E}_{\mathrm{acc}}^{(N)}(\mathcal{C})\)
with operator chains, spectral compression laws, and effective electromagnetic dimension on finite
domains.

\paragraph{Validity.}
Equations~\eqref{eq:ch4.realizable-synthesis-set}--\eqref{eq:ch4.accessibility-envelope} presuppose
linear Maxwell theory, fixed \(\omega\), and outgoing \(\mathcal{F}_{\mathrm{rad}}\)
\cite{stratton1941}. Nonlinear media, time-varying boundaries, and broadband superposition require
separate synthesis and certification maps. The norms in
Eq.~\eqref{eq:ch4.accessibility-envelope} are induced by the declared \(\mathcal{M}\); changing
aperture or calibration alters \(\mathcal{E}_{\mathrm{acc}}^{(N)}\) without changing
\(\mathcal{S}_{\mathrm{real}}\) \cite{harrington2001}.

Finite geometry enters next as a spectral constraint on both sets: support limits compress
\(\mathcal{S}_{\mathrm{real}}(\mathcal{C})\) through operator rank, while aperture limits shrink
\(\mathcal{E}_{\mathrm{acc}}^{(N)}(\mathcal{C})\) through observation rank
\cite{stratton1941,jackson1999}.


\subsection{Finite geometry as a spectral constraint}
\label{sec:ch402}
% TOC tag: [USE SS3.6.2]

Compact support \(\Omega_{s}\) and finite observation apertures do not merely reduce the number of
adjustable feed parameters; they impose a \emph{spectral compression} on both
\(\mathcal{S}_{\mathrm{real}}(\mathcal{C})\) and \(\mathcal{E}_{\mathrm{acc}}^{(N)}(\mathcal{C})\)
by collapsing high-\(\ell\) weight before synthesis or certification is attempted
\cite{stratton1941,jackson1999,harrington2001}. The truncating projector
Eq.~\eqref{eq:ch3.ellmax-truncation-projector} and remainder energy
Eq.~\eqref{eq:ch3.truncation-remainder-energy} already separate retained and tail sectors on
\(\mathcal{F}_{\mathrm{rad}}\); the present subsection records how that split propagates into
coefficient space as a rank law on realizable listings \cite{hansen1988}.

\paragraph{Assumptions.}
Fix a ball \(B_{a}\) containing \(\operatorname{supp}\Jimp\), wavenumber \(k_{0}=\omega\sqrt{\varepsilon\mu}\),
and a flux-normalized VSH listing with indices \((\ell,m,\sigma)\) as in Section~\ref{sec:ch34}
\cite{stratton1941,hansen1988}. Let \(\ell_{\max}\) denote any finite angular cut and
\(\mathcal{T}_{\ell_{\max}}\) the coefficient-space projector of
Eq.~\eqref{eq:ch3.ellmax-truncation-projector} \cite{harrington2001}. Geometry enters through
\(k_{0}a\) and through aperture diameter \(D_{\mathrm{ap}}\) when observation rank is finite
\cite{jackson1999,devaney2004}.

\paragraph{Geometry-induced tail decay.}
For sources in \(B_{a}\), far-zone coefficients obey accelerated decay when \(\ell\gg k_{0}a\)
\cite{stratton1941,jackson1999}. Tail energy
\begin{equation}
\label{eq:ch4.spectral-tail-energy}
\|\mathbf{R}_{\ell_{\max}}\|_{\mathrm{rad}}^{2}
=
\sum_{\ell>\ell_{\max}}\sum_{m,\sigma}|\widehat{c}_{\ell m}^{(\sigma)}|^{2}
\end{equation}
from Eq.~\eqref{eq:ch3.truncation-remainder-energy} is therefore not an arbitrary fitting tolerance:
it is the radiative power carried by angular orders the source geometry cannot sustain at the declared
\(k_{0}a\) \cite{hansen1988,harrington2001}. When \(\ell_{\max}\) is chosen below the geometry-implied
ceiling, \(\mathcal{T}_{\ell_{\max}}\mathbf{c}\) lies inside \(\mathcal{S}_{\mathrm{real}}(\mathcal{C})\) but
\(\mathbf{c}\) itself may not---synthesis and listing are misaligned by a predictable tail, not by noise
\cite{stratton1941}.

\begin{definition}[Spectral compression]
\label{def:ch4.spectral-compression}
Let \(\mathbf{c}(\Jimp)=\mathcal{C}(\mathcal{R}_\omega[\Jimp])\). The \emph{spectral compression} of
geometry \((\Omega_{s},k_{0})\) on listing \(\mathcal{C}\) is the smallest \(\ell_{\max}\) such that
\begin{equation}
\label{eq:ch4.spectral-compression-def}
\frac{\|\mathbf{R}_{\ell_{\max}}\|_{\mathrm{rad}}^{2}}{\|\mathbf{c}(\Jimp)\|_{\mathrm{rad}}^{2}}
\le \eta_{\mathrm{tail}},
\end{equation}
with \(\eta_{\mathrm{tail}}\) a declared tail budget and \(\|\cdot\|_{\mathrm{rad}}\) the radiative
norm of Section~\ref{sec:ch35} \cite{hansen1988,harrington2001}.
\end{definition}

Equation~\eqref{eq:ch4.spectral-compression-def} turns finite support into a quantitative cut rule:
electrical size fixes how many angular orders carry appreciable synthesized power before truncation error
dominates design targets \cite{stratton1941,jackson1999}. Compression is geometry acting on rank, not a
post hoc sparse fit \cite{harrington2001}.

\paragraph{Operator tail bound.}
When modes order decreasing singular values of the compact source--radiation map,
Eq.~\eqref{eq:ch2.compression-error-bound} bounds
\(\|\mathbf{R}_{\ell_{\max}}\|_{\mathrm{rad}}\) by \(\sigma_{\ell_{\max}+1}\), the first suppressed
singular value \cite{harrington2001}. Coefficient tails and operator tails coincide in the VSH gauge fixed
in Section~\ref{sec:ch35}; changing \(\mathcal{C}\) rotates the tail vector without changing radiated
power on \(\operatorname{ran}\mathcal{R}_\omega\) \cite{devaney2004,stratton1941}. Section~\ref{sec:ch43}
develops the operator-chain form of the same law.

\begin{proposition}[Compressed synthesis set]
\label{prop:ch4.compressed-synthesis}
For every admissible \(\mathcal{C}\) and finite \(\ell_{\max}\),
\begin{equation}
\label{eq:ch4.compressed-synthesis-set}
\mathcal{S}_{\mathrm{real}}^{(\ell_{\max})}(\mathcal{C})
=
\mathcal{T}_{\ell_{\max}}\mathcal{S}_{\mathrm{real}}(\mathcal{C})
\subseteq
\mathcal{S}_{\mathrm{real}}(\mathcal{C})
\subseteq
\mathbb{C}^{I},
\end{equation}
with strict inclusion on the left when \(\ell_{\max}<\infty\) and nontrivial radiative content above the cut
\cite{hansen1988,stratton1941}. Moreover,
\(\dim\mathcal{S}_{\mathrm{real}}^{(\ell_{\max})}(\mathcal{C})\le\dim\operatorname{ran}\mathcal{T}_{\ell_{\max}}\),
with equality only when synthesis saturates every retained mode \cite{harrington2001}.
\end{proposition}
\begin{proof}
The inclusions follow from Eq.~\eqref{eq:ch4.realizable-synthesis-set} and linearity of
\(\mathcal{T}_{\ell_{\max}}\) on coefficient space \cite{stratton1941}. Strictness follows from
Eq.~\eqref{eq:ch3.truncation-remainder-energy} when tail coefficients are nonzero on
\(\operatorname{ran}\mathcal{R}_\omega\) \cite{hansen1988}. The dimension bound is the rank of
\(\mathcal{T}_{\ell_{\max}}\) on the image of the synthesis map \cite{harrington2001}.
\end{proof}

Proposition~\ref{prop:ch4.compressed-synthesis} is the first realizability statement in which geometry
appears as a \emph{rank ceiling} rather than as a boundary condition alone: finite support compresses
\(\mathcal{S}_{\mathrm{real}}(\mathcal{C})\) before aperture or noise enter
\cite{stratton1941,jackson1999}. Sections~\ref{sec:ch43}--\ref{sec:ch46} replace
\(\ell_{\max}\) by singular-value thresholds and effective electromagnetic dimension.

\paragraph{Electrical-size scaling of compressed sets.}
When \(\Omega_{s}=B_{a}\) and \(\mathcal{C}\) is a VSH listing, the tail projector
\(\mathcal{T}_{\ell_{\max}}\) of Eq.~\eqref{eq:ch3.ellmax-truncation-projector} induces a monotone
family \(\mathcal{S}_{\mathrm{real}}^{(\ell_{\max})}(\mathcal{C})\) nested in \(\ell_{\max}\)
\cite{hansen1988,stratton1941}. For fixed \(\eta_{\mathrm{tail}}\), the smallest admissible cut obeys
\begin{equation}
\label{eq:ch4.ellmax-from-tail}
\ell_{\max}(\eta_{\mathrm{tail}},k_{0}a)
\ \sim\
k_{0}a
+
\mathcal{O}\!\big(\log(1/\eta_{\mathrm{tail}})\big),
\end{equation}
importing the engineering scaling of Subsection~\ref{sec:ch362} as a realizability certificate
\cite{hansen1988,jackson1999}. Equation~\eqref{eq:ch4.ellmax-from-tail} states that enlarging support
at fixed frequency is the primary mechanism for enlarging \(\mathcal{S}_{\mathrm{real}}(\mathcal{C})\):
phase-only retuning on fixed \(a\) cannot manufacture high-\(\ell\) weight forbidden by
Eq.~\eqref{eq:ch4.spectral-compression-def} \cite{stratton1941,devaney2004}.

\begin{corollary}[Monotone enlargement under support growth]
\label{cor:ch4.support-monotone-enlargement}
If \(a_{1}<a_{2}\) with \(\Omega_{s}^{(j)}=B_{a_{j}}\) and identical \(\mathcal{C}\), then
\(\mathcal{S}_{\mathrm{real}}^{(\ell_{\max})}(\mathcal{C};a_{1})
\subseteq
\mathcal{S}_{\mathrm{real}}^{(\ell_{\max})}(\mathcal{C};a_{2})\)
for every \(\ell_{\max}\) and strict inclusion holds once \(k_{0}a_{2}\) exceeds the bandlimit
threshold of Eq.~\eqref{eq:ch4.duality-bandlimit-scaling} \cite{jackson1999,hansen1988}.
\end{corollary}
\begin{proof}
Larger balls admit strictly richer \(\mathcal{J}_{\mathrm{real}}\) under
Eq.~\eqref{eq:ch4.support-screen}; compactness of \(\mathcal{R}_\omega\) on bounded supports yields a
larger effective rank of \(\boldsymbol{\Gamma}_\omega\) as \(a\) increases \cite{stratton1941,harrington2001}.
Tail bounds Eq.~\eqref{eq:ch4.high-ell-decay} relax in \(k_{0}a\), enlarging the image of
\(\mathcal{T}_{\ell_{\max}}\) on \(\boldsymbol{\Gamma}_\omega[\mathcal{J}_{\mathrm{real}}]\)
\cite{hansen1988}.
\end{proof}

Corollary~\ref{cor:ch4.support-monotone-enlargement} is the geometric converse of spectral compression:
geometry is not merely a boundary condition but the parameter that orders synthesis sets by inclusion
\cite{stratton1941,devaney2004}. Aperture diameter and enclosing volume enter through the same
\(D_{s}=\operatorname{diam}(\Omega_{s})\) when equivalence reduction of Section~\ref{sec:ch42} is valid
\cite{collin2001,harrington2001}.

\paragraph{Problem (tail budget versus target pattern).}
\textbf{Problem.} A source on \(B_{a}\) with \(k_{0}a=1.2\) must synthesize a VSH pattern whose normalized
coefficients satisfy \(|\widehat{c}_{\ell m}^{(\sigma)}|\ge 0.1\) for all \(\ell\le 6\) and all \(m\).
Is \(\mathbf{c}_{\mathrm{target}}\) realizable at \(\eta_{\mathrm{tail}}=10^{-2}\)?

\textbf{Step 1.} Estimate \(\ell_{\max}(\eta_{\mathrm{tail}},k_{0}a)\) from Eq.~\eqref{eq:ch4.ellmax-from-tail}:
for \(k_{0}a=1.2\), appreciable weight typically cuts near \(\ell_{\max}\approx 2\)--\(3\) at the declared tail
\cite{hansen1988,jackson1999}.

\textbf{Step 2.} Compare to the plateau at \(\ell=6\): Eq.~\eqref{eq:ch4.high-ell-decay} suppresses
\(|\widehat{c}_{6m}^{(\sigma)}|\) by several orders relative to \(\ell=1\) on the same ball
\cite{stratton1941,hansen1988}.

\textbf{Step 3.} Conclude \(\mathbf{c}_{\mathrm{target}}\notin\mathcal{S}_{\mathrm{real}}(\mathcal{C})\) unless
\(a\) is enlarged or \(\eta_{\mathrm{tail}}\) is relaxed; raising \(\ell_{\max}\) in software does not override
Eq.~\eqref{eq:ch4.compressed-synthesis-set} \cite{devaney2004,harrington2001}.

\textbf{Physical meaning.} Spectral compression is a synthesis obstruction, not a post-processing filter
\cite{hansen1988}.

\paragraph{Validity.}
Eqs.~\eqref{eq:ch4.spectral-compression-def}--\eqref{eq:ch4.compressed-synthesis-set} assume linear
Maxwell theory, fixed \(\omega\), and outgoing \(\mathcal{F}_{\mathrm{rad}}\) \cite{stratton1941}.
Anisotropic supports and strong scatterers modify tail decay exponents without removing compression
\cite{harrington2001}. Broadband pulses require per-frequency cuts; a single \(\ell_{\max}\) is a
narrowband certificate only \cite{jackson1999}.


\subsection{Continuity with Chapters 1-3}
\label{sec:ch403}
% TOC tag:

Volume~I closes in the present chapter because three independent layers---admissibility certificates,
radiation operators, and representation coordinates---already fix the domain on which realizability must
be stated. Chapter~\ref{ch:01} supplies the progressive filter
Hierarchy~\ref{hier:ch1.angular-status} and measurement composition through
Eq.~\eqref{eq:ch1.report-class-lock} \cite{stratton1941,harrington2001}. Chapter~\ref{ch:02} supplies
\(\mathcal{S}_\omega\), \(\Pi_{\mathrm{rad}}\), and \(\mathcal{R}_\omega\) together with spectral,
Fredholm, and geometric constraints that preselect accessible sectors before any harmonic label is read
\cite{stratton1941,harrington2001,devaney2004}. Chapter~\ref{ch:03} supplies flux-normalized eigenfield
coordinates, truncation calculus, and admissible representation class
Definition~\ref{def:ch3.R-adm} on which Eqs.~\eqref{eq:ch4.realizable-synthesis-set} and
\eqref{eq:ch4.accessibility-envelope} are defined \cite{hansen1988,stratton1941}.

None of those layers is reopened here. Realizability theory consumes their outputs as named objects:
\(\mathcal{J}_\Omega\) from Subsection~\ref{sec:ch201},
\(\ker\mathcal{R}_\omega\) from Section~\ref{sec:ch27},
\(\mathcal{T}_{\ell_{\max}}\) from Eq.~\eqref{eq:ch3.ellmax-truncation-projector},
and accessible sectors from Eq.~\eqref{eq:ch2.accessible-modal-sector}
\cite{harrington2001,stratton1941}. The present chapter's new content begins with impressed-source
admissibility in Section~\ref{sec:ch41}, where synthesis sets acquire explicit Maxwellian structure,
and culminates in effective electromagnetic dimension and concentration bounds in
Sections~\ref{sec:ch46}--\ref{sec:ch48} \cite{devaney2004,hansen1988}.

\paragraph{Consumption map (three prior layers).}
Chapter~\ref{ch:01} fixes \emph{which predicates} may be applied to a reported angular label: physical,
measurable, and realizable classes are distinct before any coefficient is written
\cite{stratton1941,harrington2001}. Chapter~\ref{ch:02} fixes \emph{which field states} enter radiation
theory: \(\mathcal{S}_\omega\), outgoing \(\Pi_{\mathrm{rad}}\), and the accessible sector of
Eq.~\eqref{eq:ch2.accessible-modal-sector} preselect \(\mathfrak{F}\in\mathcal{F}_{\mathrm{rad}}\) prior to
listing \cite{stratton1941,devaney2004}. Chapter~\ref{ch:03} fixes \emph{which coordinates} are admissible:
flux-normalized \(\mathcal{C}\in\mathcal{R}_{\mathrm{adm}}\) and truncation projectors define the ambient space
\(\mathbb{C}^{I}\) in which Eqs.~\eqref{eq:ch4.realizable-synthesis-set} and
\eqref{eq:ch4.accessibility-envelope} live \cite{hansen1988,stratton1941}. Chapter~\ref{ch:04} composes those
outputs into synthesis and certification subsets without altering the operators themselves
\cite{harrington2001,devaney2004}.

\begin{figure}[t]
\centering
\ChOneFigBlock{%
\begin{tikzpicture}[ch1fig, node distance=7mm and 9mm]
  \node[ch1box, minimum width=28mm] (ch1) {Ch.~\ref{ch:01}\\predicates};
  \node[ch1box, minimum width=28mm, right=11mm of ch1] (ch2) {Ch.~\ref{ch:02}\\operators};
  \node[ch1box, minimum width=28mm, right=11mm of ch2] (ch3) {Ch.~\ref{ch:03}\\coordinates};
  \node[ch1box, minimum width=32mm, below=13mm of ch2] (ch4) {Ch.~\ref{ch:04}\\$\mathcal{S}_{\mathrm{real}}$, $\mathcal{E}_{\mathrm{acc}}$};
  \draw[ch1flow] (ch1) -- node[above, font=\footnotesize] {Hierarchy~\ref{hier:ch1.angular-status}} (ch2);
  \draw[ch1flow] (ch2) -- node[above, font=\footnotesize] {$\mathcal{R}_\omega$, $\Pi_{\mathrm{rad}}$} (ch3);
  \draw[ch1flow] (ch1) |- (ch4);
  \draw[ch1flow] (ch2) -- (ch4);
  \draw[ch1flow] (ch3) -- (ch4);
\end{tikzpicture}%
}
\caption{Subsection~\ref{sec:ch403}: realizability theory composes Chapter~\ref{ch:01} predicates,
Chapter~\ref{ch:02} radiation operators, and Chapter~\ref{ch:03} representation coordinates into
\(\mathcal{S}_{\mathrm{real}}(\mathcal{C})\) and \(\mathcal{E}_{\mathrm{acc}}^{(N)}(\mathcal{C})\).}
\label{fig:ch4.volume-consumption-map}
\end{figure}

Figure~\ref{fig:ch4.volume-consumption-map} is not a dependency diagram for re-derivation; it records which
objects are \emph{inputs} to realizability inequalities \cite{harrington2001}. Changing \(\mathcal{C}\) rotates
coordinates in \(\mathbb{C}^{I}\) but leaves \(\operatorname{ran}\mathcal{R}_\omega\) fixed; changing
\(\mathcal{J}_\Omega\) or \(\Omega_{s}\) alters \(\mathcal{S}_{\mathrm{real}}(\mathcal{C})\); changing
\(\mathcal{M}\) alters \(\mathcal{E}_{\mathrm{acc}}^{(N)}(\mathcal{C})\) \cite{devaney2004,stratton1941}.

\paragraph{Operator composition consumed from Chapter~\ref{ch:02}.}
The synthesis map used throughout Chapter~\ref{ch:04} is the coefficient form of
\begin{equation}
\label{eq:ch4.gamma-composition-consumed}
\boldsymbol{\Gamma}_\omega
=
\mathcal{C}\circ\mathcal{R}_\omega
=
\mathcal{C}\circ\Pi_{\mathrm{rad}}\circ\mathcal{S}_\omega,
\end{equation}
with \(\mathcal{S}_\omega\) acting on \(\mathcal{J}_\Omega\) and \(\mathcal{C}\) acting on
\(\mathcal{F}_{\mathrm{rad}}\) \cite{stratton1941,harrington2001}. Equation~\eqref{eq:ch4.gamma-composition-consumed}
makes the Chapter~\ref{ch:02} chain legible in \(\mathbb{C}^{I}\): non-radiating components in
\(\ker\mathcal{R}_\omega\) never appear in \(\mathbf{c}\), and evanescent storage removed by
\(\Pi_{\mathrm{prop}}\) before \(\Pi_{\mathrm{rad}}\) cannot be recovered by coordinate relabeling
\cite{stratton1941,devaney2004}. Subsections~\ref{sec:ch433} and~\ref{sec:ch442} specialize the same
composition to rank and lateral-spectrum readings \cite{harrington2001}.

\paragraph{Measurement composition consumed from Chapter~\ref{ch:03}.}
Certification applies after synthesis:
\begin{equation}
\label{eq:ch4.measurement-composition-consumed}
\widehat{\mathbf{c}}
=
\mathcal{M}\circ\mathcal{N}\circ\mathcal{C}\big(\mathcal{R}_\omega[\Jimp]\big),
\end{equation}
importing Eq.~\eqref{eq:ch3.measurement-on-coefficients} \cite{hansen1988,harrington2001}. Equation
\eqref{eq:ch4.measurement-composition-consumed} defines \(\mathcal{E}_{\mathrm{acc}}^{(N)}(\mathcal{C})\) as the
set of \(\mathbf{c}\) for which \(\widehat{\mathbf{c}}\) approximates \(\mathbf{c}\) within
\(\varepsilon_{\mathrm{cert}}\) in the norm induced by \(\mathcal{M}\) \cite{devaney2004}. A design that
achieves \(\mathbf{c}\in\mathcal{S}_{\mathrm{real}}(\mathcal{C})\) still fails in the laboratory when
\(\mathbf{c}\notin\mathcal{E}_{\mathrm{acc}}^{(N)}(\mathcal{C})\) at the deployed rank \(N\)
\cite{harrington2001}.

\begin{lemma}[Non-invariance of synthesis under coordinate change]
\label{lem:ch4.synthesis-not-chart-invariant}
Let \(\mathcal{C}_{1},\mathcal{C}_{2}\in\mathcal{R}_{\mathrm{adm}}\) be related by a bounded change of
basis \(\mathbf{U}\) on a fixed sector of \(\mathcal{F}_{\mathrm{rad}}\). Then
\(\mathcal{S}_{\mathrm{real}}(\mathcal{C}_{2})=\mathbf{U}\mathcal{S}_{\mathrm{real}}(\mathcal{C}_{1})\) only on
that sector; strict realizability inequalities (rank, tail, bandlimit) are not invariant under arbitrary
\(\mathbf{U}\) that mix inaccessible modes into accessible ones \cite{hansen1988,devaney2004}.
\end{lemma}
\begin{proof}
\(\boldsymbol{\Gamma}_\omega\) is defined before \(\mathbf{U}\); mixing inaccessible coordinates into the
design vector does not enlarge \(\operatorname{ran}\boldsymbol{\Gamma}_\omega\) \cite{stratton1941,harrington2001}.
\end{proof}

Lemma~\ref{lem:ch4.synthesis-not-chart-invariant} warns against optimizing in an overcomplete listing:
realizability is a property of \(\boldsymbol{\Gamma}_\omega\), not of the cardinality of \(I\)
\cite{hansen1988}.

\paragraph{Problem (predicate failure diagnosis).}
\textbf{Problem.} A team reports \(\widehat{\mathbf{c}}\) with large \(\ell=5\) weight from a compact horn of
electrical size \(\Xi=k_{0}D_{s}/(2\pi)\approx 0.4\). Synthesis software inverts \(\mathbf{T}\) of
Eq.~\eqref{eq:ch4.mom-synthesis-matrix} and claims a feed solution. Which realizability layer fails first?

\textbf{Step 1.} Check \(\mathbf{c}\in\mathcal{S}_{\mathrm{real}}(\mathcal{C})\): for \(\Xi\approx 0.4\),
Eq.~\eqref{eq:ch4.high-ell-decay} suppresses \(\ell=5\) on compact support unless the horn enlarges \(D_{s}\)
\cite{jackson1999,hansen1988}.

\textbf{Step 2.} If synthesis is retuned only in phase, Corollary~\ref{cor:ch4.support-monotone-enlargement}
implies \(\mathcal{S}_{\mathrm{real}}\) is unchanged at fixed \(\Omega_{s}\) \cite{stratton1941}.

\textbf{Step 3.} Separately test \(\mathbf{c}\in\mathcal{E}_{\mathrm{acc}}^{(N)}(\mathcal{C})\): a finite
aperture rank \(N\) may not resolve \(\ell=5\) even if synthesis were feasible \cite{harrington2001,devaney2004}.

\textbf{Physical meaning.} Predicate failures are localized: geometry rank, coupling rank, and measurement rank
are independent obstructions \cite{hansen1988}.

\paragraph{Forward arc within Volume~I.}
Sections~\ref{sec:ch41}--\ref{sec:ch45} replace abstract subsets by impressed-source admissibility,
equivalence reduction, operator rank, and coupling selection rules \cite{stratton1941,harrington2001}.
Sections~\ref{sec:ch46}--\ref{sec:ch48} replace \(\ell_{\max}\) heuristics by effective electromagnetic
dimension, concentration inequalities, and unreachable-mode constructions \cite{devaney2004,hansen1988}.
Section~\ref{sec:ch49} closes with scaling laws that laboratories use as sanity checks before inversion
\cite{harrington2001,jackson1999}.

\paragraph{Terminal summary of Section~\ref{sec:ch40}.}
Realizability and accessibility are distinct subsets of \(\mathbb{C}^{I}\); finite geometry compresses
the synthesis set through spectral rank laws; Volume~I's prior chapters supply the operators and coordinates
on which those subsets are defined. Section~\ref{sec:ch41} opens the source-side foundation.


\section{Admissible Electromagnetic Sources}
\label{sec:ch41}

Coefficient-space realizability presupposes a source-side Maxwellian model: every
\(\mathbf{c}\in\mathcal{S}_{\mathrm{real}}(\mathcal{C})\) must be traceable to impressed data
\(\Jimp\) that satisfy charge continuity, boundary compatibility, finite-energy closure, and
compact support before \(\mathcal{R}_\omega\) and \(\mathcal{C}\) are applied
\cite{stratton1941,harrington2001}. Section~\ref{sec:ch41} fixes that model. Subsection~\ref{sec:ch411}
names impressed currents and excitation tuples as the synthesis primitive; Subsection~\ref{sec:ch412}
adds the energy, passivity, causality, and support screens that separate laboratory-realizable
excitation from formal Maxwell data; Subsection~\ref{sec:ch413} admits sheet and filamentary
channels by distributional extension; Subsection~\ref{sec:ch414} defines the \emph{realizable source
space} \(\mathcal{J}_{\mathrm{real}}\) whose image under \(\mathcal{R}_\omega\circ\mathcal{C}\) is
\(\mathcal{S}_{\mathrm{real}}(\mathcal{C})\) \cite{devaney2004,hansen1988}.


\subsection{Impressed currents and excitation}
\label{sec:ch411}
% TOC tag: [HOME]

Every synthesis claim in Eq.~\eqref{eq:ch4.realizable-synthesis-set} reduces ultimately to an
impressed current \(\Jimp\) on a declared support: antennas are fed by currents, apertures by
equivalent surface densities, and numerical codes by voxel or mesh sources that must still satisfy
\(\divg\Jimp=-\jj\omega\rho\) in the harmonic convention \cite{stratton1941,harrington2001}.
Subsection~\ref{sec:ch411} treats \(\Jimp\) not as a passive label on a field ansatz but as the
\emph{primitive} through which \(\mathcal{S}_{\mathrm{real}}(\mathcal{C})\) is generated.

\paragraph{Assumptions.}
Fix \(\omega>0\), domain \(\Omega\Subset\mathbb{R}^{3}\) with boundary \(\Gamma\), and impressed data
in the tuple sense of Eq.~\eqref{eq:ch1.data-specification-tuple} \cite{stratton1941}. Volume
currents lie in \(\mathcal{J}_\Omega\) of Eq.~\eqref{eq:ch2.source-space-refined}; outgoing fields are
discussed only after \(\Pi_{\mathrm{rad}}\) as in Subsection~\ref{sec:ch201} \cite{harrington2001}.
Representation map \(\mathcal{C}\in\mathcal{R}_{\mathrm{adm}}\) is fixed thereafter
\cite{hansen1988,stratton1941}.

\paragraph{Excitation tuple.}
Laboratory excitation is specified by volume impressed current \(\Jimp\), associated charge density
\(\rho\) linked by harmonic continuity, and boundary data compatible with
Eq.~\eqref{eq:ch1.boundary-data-operator} \cite{stratton1941}. Collect these as
\begin{equation}
\label{eq:ch4.excitation-tuple}
\mathfrak{E}
=
\left(
\Jimp,\ \rho,\ \mathfrak{g}_\Gamma
\right),
\qquad
\divg\Jimp=-\jj\omega\rho\ \text{on}\ \Omega,
\end{equation}
with \(\mathfrak{g}_\Gamma\) the tangential or impedance data induced on \(\Gamma\) by the feed and
scatterer \cite{harrington2001}. Equation~\eqref{eq:ch4.excitation-tuple} is the realizability-side
specialization of Eq.~\eqref{eq:ch1.data-specification-tuple}: the tuple is admissible when
\(\mathcal{D}_{J}[\Jimp]=1\) in Eq.~\eqref{eq:ch2.source-admissibility}, but realizability further
requires that the tuple produce nonempty radiative export on \(\operatorname{ran}\mathcal{R}_\omega\)
when a pattern is claimed \cite{stratton1941,devaney2004}.

\paragraph{Synthesis map.}
Define the \emph{synthesis map} from impressed data to radiative states by
\begin{equation}
\label{eq:ch4.synthesis-map}
\boldsymbol{\Sigma}_\omega:\mathcal{J}_\Omega\to\mathcal{F}_{\mathrm{rad}},
\qquad
\boldsymbol{\Sigma}_\omega[\Jimp]
=
\mathcal{R}_\omega[\Jimp]
=
\mathfrak{F}.
\end{equation}
Equation~\eqref{eq:ch4.synthesis-map} restates Eq.~\eqref{eq:ch4.state-from-source} with the
realizability interpretation made explicit: \(\boldsymbol{\Sigma}_\omega\) is the operator whose
range supplies every field state that finite impressed sources may \emph{generate} in the outgoing
sector \cite{stratton1941,harrington2001}. Composition with \(\mathcal{C}\) yields the coefficient
map
\begin{equation}
\label{eq:ch4.coefficient-synthesis-map}
\boldsymbol{\Gamma}_\omega
=
\mathcal{C}\circ\boldsymbol{\Sigma}_\omega:
\mathcal{J}_\Omega\to\mathbb{C}^{I},
\qquad
\boldsymbol{\Gamma}_\omega[\Jimp]=\mathbf{c}(\Jimp),
\end{equation}
so that Eq.~\eqref{eq:ch4.realizable-synthesis-set} reads
\(\mathcal{S}_{\mathrm{real}}(\mathcal{C})=\boldsymbol{\Gamma}_\omega[\mathcal{J}_\Omega]\) once the
domain is restricted to realizable currents in Subsection~\ref{sec:ch414} \cite{devaney2004}.

\begin{figure}[t]
\centering
\ChOneFigBlock{%
\begin{tikzpicture}[ch1fig, node distance=7mm and 11mm]
  \node[ch1box] (ex) {Excitation\\$\mathfrak{E}$};
  \node[ch1box, right=12mm of ex] (j) {$\Jimp\in\mathcal{J}_\Omega$};
  \node[ch1box, right=12mm of j] (sig) {$\boldsymbol{\Sigma}_\omega$};
  \node[ch1box, right=12mm of sig] (fr) {$\mathfrak{F}$};
  \node[ch1box, right=12mm of fr] (c) {$\mathcal{C}$};
  \node[ch1box, right=12mm of c] (cv) {$\mathbf{c}$};
  \draw[ch1flow] (ex) -- (j) -- (sig) -- (fr) -- (c) -- (cv);
\end{tikzpicture}%
}
\caption{Subsection~\ref{sec:ch411} synthesis chain: excitation fixes \(\Jimp\); \(\boldsymbol{\Sigma}_\omega\)
exports \(\mathfrak{F}\); \(\mathcal{C}\) lists coefficients \(\mathbf{c}\) entering
\(\mathcal{S}_{\mathrm{real}}(\mathcal{C})\).}
\label{fig:ch4.excitation-synthesis-chain}
\end{figure}

Figure~\ref{fig:ch4.excitation-synthesis-chain} is the source-side reading of
Figure~\ref{fig:ch4.synthesis-accessibility-sets}: realizability questions are questions about
\(\boldsymbol{\Gamma}_\omega[\mathcal{J}_\Omega]\), not about arbitrary \(\mathbf{c}\in\mathbb{C}^{I}\)
\cite{harrington2001}. Design specifies \(\mathbf{c}_{\mathrm{target}}\); synthesis inverts
\(\boldsymbol{\Gamma}_\omega\) only on \(\mathcal{J}_{\mathrm{real}}\) developed in
Subsection~\ref{sec:ch414} \cite{devaney2004}.

\paragraph{Excitation classes.}
Three excitation classes recur in realizability analysis:
\begin{equation}
\label{eq:ch4.excitation-class-volume}
\mathfrak{E}^{(\mathrm{vol})}:
\Jimp\in L^{2}(\Omega_{s})^{3}\ \text{or}\ \HCurl(\Omega),
\quad
\operatorname{supp}\Jimp\subseteq\Omega_{s},
\end{equation}
\begin{equation}
\label{eq:ch4.excitation-class-aperture}
\mathfrak{E}^{(\mathrm{ap})}:
\mathbf{J}_{s}\ \text{on aperture}\ \Sigma_{a},
\quad
\text{embedded via Eq.~\eqref{eq:ch2.distributional-source-extension}},
\end{equation}
\begin{equation}
\label{eq:ch4.excitation-class-port}
\mathfrak{E}^{(\mathrm{port})}:
\text{guided mode current at feed port}\ \Gamma_{\mathrm{feed}},
\end{equation}
each subject to Eq.~\eqref{eq:ch4.excitation-tuple} and boundary closure
\cite{stratton1941,harrington2001,collin2001}. Volume classes feed Eq.~\eqref{eq:ch4.synthesis-map}
directly; aperture and port classes enter through equivalence reductions of Section~\ref{sec:ch42}
\cite{harrington2001}. Subsection~\ref{sec:ch411} does not construct equivalent sources---it
records that every claimed \(\mathbf{c}\) must factor through one of these excitation classes.

\begin{principle}[Impressed excitation as synthesis primitive]
\label{prin:ch4.excitation-primitive}
No coefficient vector \(\mathbf{c}\in\mathbb{C}^{I}\) is physically realizable unless it is the image
\(\boldsymbol{\Gamma}_\omega[\Jimp]\) of impressed data \(\Jimp\) in a declared excitation class
Eqs.~\eqref{eq:ch4.excitation-class-volume}--\eqref{eq:ch4.excitation-class-port} satisfying
Eq.~\eqref{eq:ch4.excitation-tuple} and producing outgoing export on
\(\operatorname{ran}\mathcal{R}_\omega\) \cite{stratton1941,devaney2004}.
\end{principle}

Principle~\ref{prin:ch4.excitation-primitive} is the source-side mirror of
Definition~\ref{def:ch4.realizable-synthesis-set}: listing without excitation is coordinate fiction
\cite{hansen1988}. Phase-only pattern targets that skip \(\Jimp\) are checked against
\(\boldsymbol{\Gamma}_\omega\), not against \(\mathcal{C}\) alone \cite{harrington2001}.

\paragraph{Linear superposition and coupling.}
Because \(\mathcal{R}_\omega\) and \(\mathcal{C}\) are linear on their domains,
\begin{equation}
\label{eq:ch4.synthesis-superposition}
\boldsymbol{\Gamma}_\omega\!\left[\sum_{n}a_{n}\Jimp_{n}\right]
=
\sum_{n}a_{n}\,\boldsymbol{\Gamma}_\omega[\Jimp_{n}],
\end{equation}
for admissible \(a_{n}\in\mathbb{C}\) and \(\Jimp_{n}\in\mathcal{J}_\Omega\)
\cite{stratton1941}. Equation~\eqref{eq:ch4.synthesis-superposition} is synthesis superposition in a
\emph{fixed} basis; it is not coupling induced by broken symmetry---that topology is owned by
Section~\ref{sec:ch45} \cite{harrington2001}. Near-field reactive superposition of sources can
cancel far-zone export while leaving large local storage; such combinations lie in \(\mathcal{J}_\Omega\)
but may map to \(\mathbf{0}\) under \(\boldsymbol{\Sigma}_\omega\) when \(\Jimp\in\ker\mathcal{R}_\omega\)
\cite{stratton1941}.

\paragraph{Charge neutrality and synthesis.}
When \(\rho\equiv0\) on \(\Omega_{s}\), Eq.~\eqref{eq:ch4.excitation-tuple} reduces to
\(\divg\Jimp=0\) and \(\Jimp\in\mathcal{J}_\Omega^{(0)}\) of Eq.~\eqref{eq:ch2.charge-neutral-source}
\cite{stratton1941}. Charge-neutral synthesis respects
\begin{equation}
\label{eq:ch4.charge-neutral-synthesis}
\boldsymbol{\Gamma}_\omega[\Jimp]\ \text{unchanged under}\ 
\Jimp\mapsto\Jimp+\Jimp^{\mathrm{(NR)}},\ 
\Jimp^{\mathrm{(NR)}}\in\ker\mathcal{R}_\omega,
\end{equation}
so identical \(\mathbf{c}\) may arise from distinct \(\Jimp\) differing by non-radiating additions
\cite{harrington2001,stratton1941}. Equation~\eqref{eq:ch4.charge-neutral-synthesis} is the first
synthesis non-uniqueness law: realizability of \(\mathbf{c}\) does not determine a unique current
distribution \cite{devaney2004}. Inverse design must quotient \(\mathcal{J}_{\mathrm{real}}\) by
\(\ker\boldsymbol{\Gamma}_\omega\) developed in Subsection~\ref{sec:ch414}.

\paragraph{Fredholm preview on \(\boldsymbol{\Gamma}_\omega\).}
On compact \(\Omega_{s}\), \(\mathcal{R}_\omega\) is compact on \(\mathcal{J}_\Omega\); hence
\(\boldsymbol{\Gamma}_\omega\) is compact into \(\mathbb{C}^{I}\) when \(I\) is countable
\cite{stratton1941,harrington2001}. For any finite-dimensional truncation \(I_{N}\subset I\),
\begin{equation}
\label{eq:ch4.gamma-fredholm-preview}
\boldsymbol{\Gamma}_{\omega,N}
=
\mathcal{P}_{N}\circ\boldsymbol{\Gamma}_\omega:
\mathcal{J}_\Omega\to\mathbb{C}^{I_{N}},
\qquad
\operatorname{rank}\boldsymbol{\Gamma}_{\omega,N}<\infty,
\end{equation}
with \(\mathcal{P}_{N}\) the coordinate truncation of Subsection~\ref{sec:ch402}
\cite{hansen1988,harrington2001}. Equation~\eqref{eq:ch4.gamma-fredholm-preview} is the operator form of
spectral compression: synthesis is intrinsically finite rank before measurement screens act
\cite{stratton1941}. Section~\ref{sec:ch43} replaces \(\mathcal{P}_{N}\) by singular-value thresholds on the
full chain \(\Jimp\mapsto\mathfrak{F}\mapsto\mathbf{c}\) \cite{devaney2004}.

\begin{corollary}[Synthesis range is not coordinate-open]
\label{cor:ch4.synthesis-not-open}
For every compact \(\Omega_{s}\) and admissible \(\mathcal{C}\), \(\boldsymbol{\Gamma}_\omega[\mathcal{J}_\Omega]\)
has empty interior in the standard topology of \(\mathbb{C}^{I}\) when \(I\) is infinite; in particular,
\(\mathcal{S}_{\mathrm{real}}(\mathcal{C})\) is nowhere dense in \(\mathbb{C}^{I}\) \cite{stratton1941,hansen1988}.
\end{corollary}
\begin{proof}
Compactness of \(\mathcal{R}_\omega\) on bounded sources implies compactness of
\(\boldsymbol{\Gamma}_\omega\) on bounded subsets of \(\mathcal{J}_\Omega\); the image of a compact operator
on the unit ball is precompact, hence cannot contain an open set in an infinite-dimensional coefficient space
\cite{harrington2001,stratton1941}.
\end{proof}

Corollary~\ref{cor:ch4.synthesis-not-open} is a structural inevitability result: pattern optimization over all
\(\mathbf{c}\in\mathbb{C}^{I}\) optimizes over a set whose interior is empty---feasible targets lie on a thin
synthesis manifold determined by \(\Omega_{s}\) and \(\omega\) \cite{devaney2004,hansen1988}.

\paragraph{Worked illustration (two-feed array).}
Two compact feeds \(\Jimp_{1},\Jimp_{2}\) on \(\Omega_{s}\) produce
\(\mathbf{c}=\boldsymbol{\Gamma}_\omega[\Jimp_{1}]+\boldsymbol{\Gamma}_\omega[\Jimp_{2}]\) by
Eq.~\eqref{eq:ch4.synthesis-superposition} \cite{harrington2001}. Steering by phase \(a_{2}=e^{\jj\phi}\)
is realizable only if both components lie in \(\mathcal{J}_{\mathrm{real}}\); a phase target that
requires \(\mathbf{c}\notin\boldsymbol{\Gamma}_\omega[\mathcal{J}_\Omega]\) is not a feed-law problem
but a realizability obstruction \cite{devaney2004,collin2001}.

\paragraph{Problem (synthesis feasibility test).}
\textbf{Problem.} A design requests \(\mathbf{c}_{\mathrm{target}}\) with
\(|\widehat{c}_{\ell m}^{(\sigma)}|=0\) for all \(\ell>2\) at \(k_{0}a=1.5\). Is the target necessarily
realizable on a single compact \(\Omega_{s}\)?

\textbf{Step 1.} Compute electrical size \(k_{0}D_{s}/2\) for the proposed \(\Omega_{s}\)
\cite{stratton1941,jackson1999}.

\textbf{Step 2.} Compare to spectral compression Eq.~\eqref{eq:ch4.spectral-compression-def}: if tail budget
\(\eta_{\mathrm{tail}}\) forces \(\ell_{\max}<2\), \(\ell=2\) content is geometrically suppressed even when
formally allowed in \(\mathbb{C}^{I}\) \cite{hansen1988}.

\textbf{Step 3.} Test \(\mathbf{c}_{\mathrm{target}}\in\boldsymbol{\Gamma}_\omega[\mathcal{J}_{\mathrm{real}}]\) via
Theorem~\ref{thm:ch4.synthesis-surjectivity} preview, not via pattern fitting alone \cite{devaney2004}.

\textbf{Physical meaning.} Zero coefficients above \(\ell=2\) in a listing do not imply realizability; geometry
may forbid \(\ell=2\) excitation on small \(\Omega_{s}\) \cite{harrington2001}.

\paragraph{Excitation class taxonomy.}
Realizable feeds partition into volume, aperture, port, and hybrid classes through
Eqs.~\eqref{eq:ch4.excitation-class-volume}--\eqref{eq:ch4.excitation-class-port} \cite{stratton1941,collin2001}.
Each class induces a distinct rank signature on \(\boldsymbol{\Gamma}_\omega\): volume classes typically maximize
\(r_{\mathrm{rad}}(\varepsilon)\) for fixed \(\Omega_{s}\); aperture classes minimize unknowns via equivalence;
port classes constrain \(\mathfrak{E}\) to guided modes \cite{harrington2001,devaney2004}. Hybrid tuples superpose per
Eq.~\eqref{eq:ch4.synthesis-superposition} but remain subject to a single joint \(\mathcal{D}_{\mathrm{real}}\) screen
\cite{stratton1941}.

\begin{corollary}[Excitation class rank ordering]
\label{cor:ch4.excitation-rank-ordering}
At fixed exterior \(\mathbf{c}\) and equivalence-valid enclosure,
\(\operatorname{rank}\mathbf{T}_{\mathrm{aperture}}\le\operatorname{rank}\mathbf{T}_{\mathrm{volume}}\)
with equality when Proposition~\ref{prop:ch4.aperture-realizability-reduction} applies
\cite{collin2001,harrington2001,stratton1941}.
\end{corollary}

\paragraph{Harmonic continuity and charge linkage.}
Under \(\exp(-\jj\omega t)\), the excitation tuple Eq.~\eqref{eq:ch4.excitation-tuple} requires
\(\divg\Jimp=-\jj\omega\rho\) on \(\Omega\) \cite{stratton1941}. Realizable synthesis therefore presupposes a
consistent charge model: aperture or surface currents must be accompanied by the \(\rho\) or \(\mathbf{J}_{s}\) jumps
implied by Maxwell continuity \cite{harrington2001,collin2001}. A feed tuple with \(\mathcal{D}_{J}=0\) excludes
\(\mathbf{c}\) before \(\boldsymbol{\Gamma}_\omega\) is formed \cite{stratton1941,devaney2004}.

\paragraph{Worked illustration (port-fed aperture).}
A waveguide port supplies \(\mathfrak{g}_\Gamma\) while \(\mathbf{J}_{s}\) on \(\Sigma_{a}\) carries radiation
\cite{collin2001,harrington2001}. Synthesis rank is the minimum of port mode count and aperture MoM rank per
Corollary~\ref{cor:ch4.channel-count} \cite{devaney2004,hansen1988}.

\paragraph{Problem (charge-neutral feed).}
\textbf{Problem.} A volume \(\Jimp\) is divergence-free but claims \(\ell>1\) dominant \(\mathbf{c}\) on
\(k_{0}D_{s}\ll1\). Is charge neutrality sufficient for realizability?

\textbf{Step 1.} \(\mathcal{D}_{J}=1\) is necessary but not sufficient \cite{stratton1941}.

\textbf{Step 2.} Test \(\Xi\) and Eq.~\eqref{eq:ch4.high-ell-decay} \cite{hansen1988,jackson1999}.

\textbf{Step 3.} Conclude geometry rank, not charge neutrality, governs \(\ell>1\) content \cite{devaney2004,harrington2001}.

\textbf{Physical meaning.} Continuity screens are entry gates, not synthesis certificates \cite{stratton1941}.

\paragraph{Excitation tuple decomposition.}
Every admissible \(\mathfrak{E}\) in Eq.~\eqref{eq:ch4.excitation-tuple} admits a decomposition
\begin{equation}
\label{eq:ch4.excitation-decomposition}
\mathfrak{E}
=
\mathfrak{E}_{\mathrm{vol}}
\oplus
\mathfrak{E}_{\mathrm{surf}}
\oplus
\mathfrak{E}_{\mathrm{port}},
\end{equation}
with each summand carrying its own realizability screen before composition \cite{stratton1941,harrington2001,collin2001}.
Volume terms supply \(\Jimp\in\mathcal{J}_\Omega\); surface terms supply \(\mathbf{J}_{s}\) through
Eq.~\eqref{eq:ch4.surface-source-embedding}; port terms supply \(\mathfrak{g}_\Gamma\) from guided-mode catalogs
\cite{collin2001}. The decomposition is not unique, but the \emph{joint} image
\(\boldsymbol{\Gamma}_\omega[\mathcal{J}_{\mathrm{real}}(\mathfrak{E})]\) is unique up to \(\mathcal{N}_{\mathrm{syn}}\)
\cite{devaney2004,hansen1988}.

\begin{theorem}[Excitation-class synthesis bound]
\label{thm:ch4.excitation-class-bound}
For excitation classes Eq.~\eqref{eq:ch4.excitation-class-volume}--\eqref{eq:ch4.excitation-class-port},
\begin{equation}
\label{eq:ch4.excitation-class-rank}
\dim\mathcal{S}_{\mathrm{real}}(\mathcal{C};\mathfrak{E})
\le
\min\big\{
r_{\mathrm{rad}}(\varepsilon),\,
\operatorname{rank}\mathbf{T}(\mathfrak{E}),\,
|I_{\ell_{\max}}|
\big\},
\end{equation}
with \(\mathbf{T}(\mathfrak{E})\) the MoM chain Eq.~\eqref{eq:ch4.discrete-chain-matrix} restricted to the declared class
\cite{stratton1941,harrington2001,devaney2004}. Hybrid tuples achieve the bound only when every channel satisfies
\(\mathcal{D}_{\mathrm{real}}=1\) jointly \cite{collin2001}.
\end{theorem}
\begin{proof}
Corollary~\ref{cor:ch4.channel-count}, Definition~\ref{def:ch4.radiative-rank}, and linearity of
\(\boldsymbol{\Gamma}_\omega\) on \(\mathcal{J}_{\mathrm{real}}\) \cite{harrington2001,stratton1941}.
\end{proof}

\paragraph{Worked illustration (hybrid horn feed).}
A horn with interior \(\Jimp\) and mouth \(\mathbf{J}_{s}\) must satisfy both volume and aperture screens: enclosure for
equivalence and \(\mathcal{D}_{\mathrm{cpl}}=1\) on \(\Sigma_{a}\) \cite{collin2001,harrington2001}. Reporting
\(\ell_{\max}=15\) from the mouth alone while the interior feed is electrically small violates
Eq.~\eqref{eq:ch4.excitation-class-rank} \cite{hansen1988,devaney2004}.

\paragraph{Problem (class selection for inverse design).}
\textbf{Problem.} Inverse design chooses aperture currents only, but \(\mathbf{c}_{\mathrm{target}}\) requires interior
reactive storage not reachable by \(\Sigma_{a}\). Which class should be declared?

\textbf{Step 1.} Test \(\mathbf{c}_{\mathrm{target}}\in\mathcal{S}_{\mathrm{ap}}(\mathcal{C})\) versus
\(\mathcal{S}_{\mathrm{real}}(\mathcal{C})\) \cite{collin2001,devaney2004}.

\textbf{Step 2.} If \(\delta_{\mathrm{syn}}>0\) on aperture class alone, enlarge \(\mathfrak{E}\) to volume or port terms
\cite{harrington2001}.

\textbf{Step 3.} Recompute Eq.~\eqref{eq:ch4.excitation-class-rank} on the enlarged tuple \cite{stratton1941}.

\textbf{Physical meaning.} Excitation class is a realizability declaration, not a solver convenience \cite{stratton1941}.

\paragraph{Validity.}
Eqs.~\eqref{eq:ch4.excitation-tuple}--\eqref{eq:ch4.synthesis-superposition} assume linear Maxwell
theory, fixed \(\omega\), and outgoing \(\mathcal{F}_{\mathrm{rad}}\) \cite{stratton1941}. Nonlinear
amplifiers and time-varying feeds require time-domain excitation maps before harmonic synthesis is
meaningful \cite{harrington2001}. Port excitation Eq.~\eqref{eq:ch4.excitation-class-port} presupposes
a guided-mode model supplied by the scatterer geometry; Section~\ref{sec:ch42} reduces ports to
equivalent aperture currents \cite{collin2001}.

Finite-energy, passivity, causality, and support screens refine \(\mathcal{J}_\Omega\) to the
realizable subclass analyzed in Subsection~\ref{sec:ch412}.


\subsection{Finite energy, passivity, causality, and support}
\label{sec:ch412}
% TOC tag: [HOME]

Admissibility \(\mathcal{D}_{J}=1\) does not yet certify that an impressed current belongs to a
\emph{laboratory-realizable} model: infinite-energy formal solutions, acausal precursors, and
non-passive material embeddings can satisfy local Maxwell residuals while violating the physical
screens under which \(\boldsymbol{\Sigma}_\omega\) is defined \cite{stratton1941,harrington2001}.
Subsection~\ref{sec:ch412} records the four constraints that shrink \(\mathcal{J}_\Omega\) to currents
compatible with finite hardware and retarded physics.

\paragraph{Assumptions.}
Let \(\mathfrak{E}\) satisfy Eq.~\eqref{eq:ch4.excitation-tuple} with \(\Jimp\in\mathcal{J}_\Omega\).
Materials in \(\Omega\) are linear, isotropic, and passive unless a subsection explicitly relaxes
passivity \cite{stratton1941}. Time-domain restoration uses retarded supports as in
Eq.~\eqref{eq:ch2.causality-preview} \cite{stratton1941,jackson1999}. Source support
\(\Omega_{s}\Subset\Omega\) is compact with diameter \(D_{s}=\max_{\mathbf{r},\mathbf{r}'\in\Omega_{s}}|\mathbf{r}-\mathbf{r}'|\)
\cite{hansen1988}.

\paragraph{Finite-energy screen.}
Radiative realizability requires finite energy in the sense of Eq.~\eqref{eq:ch1.energy-class} on
bounded observation shells:
\begin{equation}
\label{eq:ch4.finite-energy-screen}
\mathcal{D}_{E}[\mathfrak{F}]
=
1
\iff
\sup_{R>R_{0}}
\int_{\Gamma_{R}}
\big(
|\E_{\mathrm{rad}}|^{2}+|\H_{\mathrm{rad}}|^{2}
\big)\,dS
<
\infty,
\end{equation}
with \(\mathfrak{F}=\boldsymbol{\Sigma}_\omega[\Jimp]\) and \(\Gamma_{R}\) a sphere enclosing
\(\Omega_{s}\) \cite{stratton1941,harrington2001}. Equation~\eqref{eq:ch4.finite-energy-screen}
imports the finite-energy floor of Chapter~\ref{ch:01} to synthesis: \(\Jimp\) that excites
non-finite-energy outgoing states cannot produce certifiable far patterns even if
\(\mathcal{D}_{J}[\Jimp]=1\) on a bounded domain truncation \cite{stratton1941}. Reactive storage
near \(\Omega_{s}\) may be large while Eq.~\eqref{eq:ch4.finite-energy-screen} still holds on
\(\mathcal{F}_{\mathrm{rad}}\) \cite{harrington2001}.

\paragraph{Passivity screen.}
When impressed currents interact with passive boundaries through
Eq.~\eqref{eq:ch2.passive-impedance-class}, realizability requires that no synthesis tuple extract
net power from a passive enclosure:
\begin{equation}
\label{eq:ch4.passivity-screen}
\mathcal{D}_{\mathrm{pass}}[\mathfrak{E}]
=
1
\iff
\operatorname{Re}\!\int_{\Gamma}
\mathbf{S}\cdot\hat{\mathbf{n}}\,dS
\ge 0
\end{equation}
for all admissible test states generated by \(\mathfrak{E}\) under passive \(\mathcal{B}_\Gamma\)
\cite{harrington2001,stratton1941}. Equation~\eqref{eq:ch4.passivity-screen} excludes active
gain media and negative-index embeddings from the default realizability class; such structures may
be treated as extended tuples with explicit power injection channels \cite{harrington2001}. Passive
screens do not restrict \(\ker\mathcal{R}_\omega\): non-radiating passive currents remain in
\(\mathcal{J}_\Omega\) but map to \(\mathbf{0}\) under \(\boldsymbol{\Sigma}_\omega\)
\cite{stratton1941}.

\paragraph{Causality screen.}
Harmonic phasors \(\Jimp(\mathbf{r})\) are realizable only if they are the \(\exp(-\jj\omega t)\)
specialization of a retarded time-domain current \(j(\mathbf{r},t)\) supported on \(t'\le t\) in
Eq.~\eqref{eq:ch2.causality-preview} \cite{stratton1941,jackson1999}. Define
\begin{equation}
\label{eq:ch4.causality-screen}
\mathcal{D}_{\mathrm{caus}}[\Jimp]
=
1
\iff
\operatorname{supp}_{t}\,j(\cdot,t)\subseteq(-\infty,t],
\end{equation}
equivalently, no advanced Green convolution contributes to \(\E(\mathbf{r},t)\)
\cite{stratton1941}. Equation~\eqref{eq:ch4.causality-screen} is invisible in pure frequency-domain
pattern formulas yet decisive for pulsed and wideband realizability: spectra of \(\Jimp\) must be
consistent with a causal time support \cite{jackson1999,harrington2001}.

\paragraph{Compact-support screen.}
Spectral compression of Section~\ref{sec:ch402} presupposes compact \(\Omega_{s}\):
\begin{equation}
\label{eq:ch4.support-screen}
\mathcal{D}_{\mathrm{supp}}[\Jimp]
=
1
\iff
\operatorname{supp}\Jimp\subseteq\Omega_{s},\
\operatorname{diam}(\Omega_{s})=D_{s}<\infty.
\end{equation}
Equation~\eqref{eq:ch4.support-screen} ties realizability to electrical size \(k_{0}D_{s}/2\): the
synthesis set \(\mathcal{S}_{\mathrm{real}}(\mathcal{C})\) acquires tail decay laws from
Subsection~\ref{sec:ch402} only when Eq.~\eqref{eq:ch4.support-screen} holds
\cite{stratton1941,jackson1999,hansen1988}. Extended apertures with \(\Omega_{s}\) growing with
frequency require explicit re-declaration of \(D_{s}(\omega)\) \cite{harrington2001}.

\begin{definition}[Realizability admissibility screen]
\label{def:ch4.realizability-screen}
The realizability screen is
\begin{equation}
\label{eq:ch4.realizability-screen}
\mathcal{D}_{\mathrm{real}}[\Jimp]
=
1
\iff
\left(
\mathcal{D}_{J}[\Jimp]=1\
\wedge\
\mathcal{D}_{E}[\boldsymbol{\Sigma}_\omega[\Jimp]]=1\
\wedge\
\mathcal{D}_{\mathrm{pass}}[\mathfrak{E}]=1\
\wedge\
\mathcal{D}_{\mathrm{caus}}[\Jimp]=1\
\wedge\
\mathcal{D}_{\mathrm{supp}}[\Jimp]=1
\right).
\end{equation}
\end{definition}

Definition~\ref{def:ch4.realizability-screen} is the source-side certificate parallel to
Definition~\ref{def:ch3.R-adm} on the representation side: \(\mathcal{D}_{\mathrm{real}}=1\) is necessary
for \(\mathbf{c}(\Jimp)\) to enter \(\mathcal{S}_{\mathrm{real}}(\mathcal{C})\) with physical meaning
\cite{stratton1941,devaney2004}. Violating any clause leaves a formal Maxwell solution that cannot be
implemented by finite hardware in the declared model class \cite{harrington2001}.

\begin{figure}[t]
\centering
\ChOneFigBlock{%
\begin{tikzpicture}[ch1fig, node distance=6mm and 8mm]
  \node[ch1box] (dj) {$\mathcal{D}_{J}$};
  \node[ch1box, right=9mm of dj] (de) {$\mathcal{D}_{E}$};
  \node[ch1box, right=9mm of de] (dp) {$\mathcal{D}_{\mathrm{pass}}$};
  \node[ch1box, right=9mm of dp] (dc) {$\mathcal{D}_{\mathrm{caus}}$};
  \node[ch1box, right=9mm of dc] (ds) {$\mathcal{D}_{\mathrm{supp}}$};
  \node[ch1box, below=12mm of de] (dr) {$\mathcal{D}_{\mathrm{real}}=1$};
  \draw[ch1flow] (dj) -- (dr);
  \draw[ch1flow] (de) -- (dr);
  \draw[ch1flow] (dp) -- (dr);
  \draw[ch1flow] (dc) -- (dr);
  \draw[ch1flow] (ds) -- (dr);
\end{tikzpicture}%
}
\caption{Subsection~\ref{sec:ch412} realizability screen: five clauses compose
\(\mathcal{D}_{\mathrm{real}}\) in Eq.~\eqref{eq:ch4.realizability-screen}.}
\label{fig:ch4.realizability-screen}
\end{figure}

Figure~\ref{fig:ch4.realizability-screen} orders the screens applied before synthesis is claimed.
\(\mathcal{D}_{J}\) is inherited from Eq.~\eqref{eq:ch2.source-admissibility}; the remaining four
clauses are the new realizability content of Section~\ref{sec:ch41} \cite{stratton1941,harrington2001}.

\begin{proposition}[Screen independence]
\label{prop:ch4.screen-independence}
The clauses in Eq.~\eqref{eq:ch4.realizability-screen} are logically independent: for each single
clause \(X\in\{E,\mathrm{pass},\mathrm{caus},\mathrm{supp}\}\), there exist tuples satisfying all
other clauses with \(\mathcal{D}_{X}=0\) while \(\mathcal{D}_{J}=1\) \cite{stratton1941,harrington2001}.
\end{proposition}
\begin{proof}
\(\mathcal{D}_{E}=0\) on formal outgoing solutions with divergent shell energy;
\(\mathcal{D}_{\mathrm{pass}}=0\) with active embedding;
\(\mathcal{D}_{\mathrm{caus}}=0\) with advanced-time spectra;
\(\mathcal{D}_{\mathrm{supp}}=0\) with spatially extended non-compact models---each while local
Maxwell residuals may still hold on truncated domains \cite{stratton1941,jackson1999,harrington2001}.
\end{proof}

Proposition~\ref{prop:ch4.screen-independence} justifies treating realizability as a conjunction rather
than a restatement of \(\mathcal{D}_{J}\): laboratories fail on distinct clauses, and synthesis audits must
test each screen relevant to the claimed \(\mathbf{c}\) \cite{devaney2004}.

\paragraph{Radiated-power linkage.}
Finite-energy export is equivalent, on \(\mathcal{F}_{\mathrm{rad}}\), to finite radiated power through the
far-sphere screen Eq.~\eqref{eq:ch2.radiated-power-screen} \cite{stratton1941,harrington2001}. For
\(\mathfrak{F}=\boldsymbol{\Sigma}_\omega[\Jimp]\),
\begin{equation}
\label{eq:ch4.power-energy-equivalence}
\mathcal{D}_{E}[\mathfrak{F}]=1
\quad\Longrightarrow\quad
P_{\mathrm{rad}}
=
\frac{1}{2}\operatorname{Re}\!\oint_{\Gamma_{R}}
\mathbf{S}\cdot\hat{\mathbf{n}}\,dS
<
\infty,
\end{equation}
with \(\mathbf{S}=\tfrac{1}{2}\E\times\H^{*}\) the harmonic Poynting vector \cite{poynting1884,stratton1941}.
Equation~\eqref{eq:ch4.power-energy-equivalence} connects \(\mathcal{D}_{E}\) to watt meters used in
laboratory certification: a synthesis tuple with divergent shell energy cannot yield a finite \(P_{\mathrm{rad}}\)
report compatible with Section~\ref{sec:ch35} normalization \cite{hansen1988,harrington2001}.

\paragraph{Broadband causality constraint.}
Let \(\widetilde{\Jimp}(\mathbf{r},\omega)\) denote the temporal Fourier transform of \(j(\mathbf{r},t)\). Causality
Eq.~\eqref{eq:ch4.causality-screen} requires \(\widetilde{\Jimp}\) to be the boundary value of a function analytic
in the lower half-plane \(\operatorname{Im}\omega<0\) under the \(\exp(-\jj\omega t)\) convention
\cite{stratton1941,jackson1999}. Single-frequency realizability at \(\omega_{0}\) therefore imposes a \emph{family}
constraint on neighboring \(\omega\) when wideband pulses are synthesized: spectra that violate Kramers--Kronig
consistency cannot be realized by any causal \(j(\mathbf{r},t)\) \cite{jackson1999,harrington2001}. Narrowband
antennas at fixed \(\omega\) satisfy Eq.~\eqref{eq:ch4.causality-screen} automatically when fed by causal
time-domain voltages \cite{collin2001}.

\paragraph{Support scaling and electrical size.}
Define the dimensionless electrical size
\begin{equation}
\label{eq:ch4.electrical-size}
\Xi
=
\frac{k_{0}D_{s}}{2\pi},
\end{equation}
with \(D_{s}\) from Eq.~\eqref{eq:ch4.support-screen} \cite{stratton1941,jackson1999}. Realizable synthesis at
fixed \(\Xi\ll1\) concentrates \(\boldsymbol{\Gamma}_\omega[\mathcal{J}_{\mathrm{real}}]\) in low-\(\ell\) sectors;
raising \(\Xi\) enlarges the accessible coefficient cone before truncation and measurement further compress it
\cite{hansen1988,harrington2001}. Equation~\eqref{eq:ch4.electrical-size} is the single parameter that links
Subsection~\ref{sec:ch412} support screens to Subsection~\ref{sec:ch402} spectral compression.

\begin{lemma}[Screen failure taxonomy]
\label{lem:ch4.screen-failures}
Each failure \(\mathcal{D}_{\mathrm{real}}=0\) produces a distinct laboratory signature:
\(\mathcal{D}_{E}=0\) yields divergent far-sphere integrals;
\(\mathcal{D}_{\mathrm{pass}}=0\) yields negative net power through passive \(\Gamma\);
\(\mathcal{D}_{\mathrm{caus}}=0\) yields advanced-time response in transient replay;
\(\mathcal{D}_{\mathrm{supp}}=0\) removes tail decay laws of Eq.~\eqref{eq:ch4.spectral-compression-def}
\cite{stratton1941,harrington2001,jackson1999}.
\end{lemma}

Lemma~\ref{lem:ch4.screen-failures} is diagnostic, not constructive: identifying which clause failed prevents
misattributing a realizability obstruction to measurement noise or basis choice \cite{devaney2004}.

\paragraph{Worked illustration (passive enclosure).}
A source inside a passive cavity with \(\mathcal{B}_\Gamma\) satisfying
Eq.~\eqref{eq:ch2.passive-impedance-class} cannot sustain \(\mathcal{D}_{\mathrm{pass}}=0\) unless external
power is injected through a port counted in \(\mathfrak{E}\) \cite{harrington2001,collin2001}. Synthesis claims
for \(\mathbf{c}\) that require net outward power exceeding feed capability fail
Eq.~\eqref{eq:ch4.passivity-screen} before \(\boldsymbol{\Gamma}_\omega\) is inverted \cite{stratton1941}.

\paragraph{Composed realizability screen.}
The joint screen Eq.~\eqref{eq:ch4.realizability-screen} factors as a logical product of physical predicates:
\begin{equation}
\label{eq:ch4.screen-factorization}
\mathcal{D}_{\mathrm{real}}
=
\mathcal{D}_{E}\cdot\mathcal{D}_{\mathrm{pass}}\cdot\mathcal{D}_{\mathrm{caus}}\cdot\mathcal{D}_{\mathrm{supp}}
\cdot\mathcal{D}_{J}\cdot\mathcal{D}_{\mathrm{cpl}},
\end{equation}
with \(\mathcal{D}_{J}\) from Eq.~\eqref{eq:ch2.source-admissibility} and \(\mathcal{D}_{\mathrm{cpl}}\) from
Eq.~\eqref{eq:ch4.coupling-admissibility} when surface channels are present \cite{stratton1941,harrington2001}.
Equation~\eqref{eq:ch4.screen-factorization} is the audit checklist for \(\mathfrak{E}\): a single failed factor
excludes \(\mathbf{c}\) from \(\mathcal{S}_{\mathrm{real}}(\mathcal{C})\) regardless of pattern fit quality
\cite{devaney2004}. Lemma~\ref{lem:ch4.screen-failures} assigns laboratory signatures to each zero factor
\cite{harrington2001,jackson1999}.

\begin{theorem}[Screen composition for synthesis]
\label{thm:ch4.screen-composition}
\(\Jimp\in\mathcal{J}_{\mathrm{real}}\) if and only if \(\mathcal{D}_{\mathrm{real}}=1\) on the excitation tuple
Eq.~\eqref{eq:ch4.excitation-tuple} \cite{stratton1941,harrington2001}. Moreover,
\(\mathbf{c}=\boldsymbol{\Gamma}_\omega[\Jimp]\in\mathcal{S}_{\mathrm{real}}(\mathcal{C})\) only if
\(\mathcal{D}_{\mathrm{real}}=1\) and \(\mathbf{r}_{\min}=\mathbf{0}\) in Principle~\ref{prin:ch4.operator-chain}
\cite{devaney2004,hansen1988}.
\end{theorem}
\begin{proof}
Forward: definitions of \(\mathcal{J}_{\mathrm{real}}\) and Eq.~\eqref{eq:ch4.realizability-screen}
\cite{stratton1941}. Converse for \(\mathbf{c}\): operator-chain range condition
\cite{harrington2001,devaney2004}.
\end{proof}

\paragraph{Worked illustration (causality versus pattern).}
A non-causal \(\widetilde{\Jimp}(\omega)\) may produce a attractive far-zone template at one frequency yet fail
wideband pulse synthesis when replayed in time \cite{jackson1999,stratton1941}. \(\mathcal{D}_{\mathrm{caus}}=0\)
therefore excludes the tuple before \(\ell_{\max}\) or MoM rank is discussed \cite{harrington2001}.

\paragraph{Problem (which screen failed?).}
\textbf{Problem.} MoM returns \(\mathbf{c}\) with large high-\(\ell\) content on a compact \(\Omega_{s}\), but
\(P_{\mathrm{rad}}\) is finite and passive. Which screen should be tested first?

\textbf{Step 1.} Verify \(\mathcal{D}_{\mathrm{supp}}=1\) and \(\Xi\) from Eq.~\eqref{eq:ch4.electrical-size}
\cite{stratton1941,jackson1999}.

\textbf{Step 2.} If screens pass, test \(\mathbf{r}_{\min}\) and Eq.~\eqref{eq:ch4.high-ell-decay}---geometry rank, not
passivity, likely failed \cite{hansen1988,devaney2004}.

\textbf{Step 3.} Do not attribute failure to \(\mathcal{D}_{\mathrm{pass}}\) when Lemma~\ref{lem:ch4.screen-failures}
implicates support or operator range \cite{harrington2001}.

\textbf{Physical meaning.} Screen factorization localizes realizability failures \cite{stratton1941}.

\paragraph{Screen composition algebra.}
Screens multiply in the realizability tuple:
\begin{equation}
\label{eq:ch4.screen-composition-algebra}
\mathcal{D}_{\mathrm{real}}
=
\mathcal{D}_{E}\,
\mathcal{D}_{\mathrm{pass}}\,
\mathcal{D}_{\mathrm{caus}}\,
\mathcal{D}_{\mathrm{supp}}\,
\mathcal{D}_{J}\,
\mathcal{D}_{\mathrm{cpl}},
\end{equation}
with each factor binary on the declared domain \cite{stratton1941,harrington2001}. Failure of any factor excludes
\(\mathbf{c}\) before \(\boldsymbol{\Gamma}_\omega\) is evaluated \cite{devaney2004}. The algebra is not commutative in
diagnostics: a causality failure may mask a support failure unless screens are tested in the order
energy \(\rightarrow\) passivity \(\rightarrow\) causality \(\rightarrow\) support \(\rightarrow\) continuity \(\rightarrow\) coupling
\cite{harrington2001}.

\begin{lemma}[Passivity implies finite energy on narrowband tuples]
\label{lem:ch4.passivity-finite-energy}
On harmonic tuples with \(\mathcal{D}_{\mathrm{pass}}=1\), \(\mathcal{D}_{E}=1\) follows for bounded port power
\(P_{\mathrm{feed}}\) \cite{stratton1941,poynting1884,harrington2001}. Conversely, \(\mathcal{D}_{E}=1\) does not imply
\(\mathcal{D}_{\mathrm{pass}}=1\) when active injection is undeclared \cite{harrington2001,collin2001}.
\end{lemma}

\paragraph{Worked illustration (active array).}
An active array with declared \(P_{\mathrm{feed}}>0\) may satisfy \(\mathcal{D}_{E}=1\) while \(\mathcal{D}_{\mathrm{pass}}=0\):
synthesis is realizable only when the port channel is explicit in \(\mathfrak{E}\) \cite{harrington2001,collin2001,stratton1941}.

\paragraph{Problem (screen order in failure logs).}
\textbf{Problem.} Logs show \(\mathcal{D}_{\mathrm{pass}}=0\) and \(\mathcal{D}_{\mathrm{supp}}=0\). Which failure is primary?

\textbf{Step 1.} Test \(\mathcal{D}_{\mathrm{supp}}\) first: non-compact support invalidates rank laws entirely
\cite{stratton1941,jackson1999}.

\textbf{Step 2.} If \(\mathcal{D}_{\mathrm{supp}}=1\), test passivity against declared \(P_{\mathrm{feed}}\)
\cite{harrington2001,poynting1884}.

\textbf{Step 3.} Do not tune \(\ell_{\max}\) until both screens pass \cite{devaney2004,hansen1988}.

\textbf{Physical meaning.} Screen algebra orders diagnostics; it does not create new synthesis directions \cite{stratton1941}.

\paragraph{Validity.}
Eqs.~\eqref{eq:ch4.finite-energy-screen}--\eqref{eq:ch4.realizability-screen} assume linear Maxwell
theory and the passive default \cite{stratton1941}. Strong-field and nonlinear regimes require
modified energy and causality screens \cite{harrington2001}. Wideband pulses must satisfy
Eq.~\eqref{eq:ch4.causality-screen} pointwise in time, not only at a single \(\omega\)
\cite{jackson1999}.

Distributional extensions of \(\mathcal{J}_\Omega\) for sheets and filaments are admitted in
Subsection~\ref{sec:ch413}.


\subsection{Distributional sources}
\label{sec:ch413}
% TOC tag: [HOME]

Compact antennas and aperture antennas are often modeled as surface or line currents rather than as
volume fillings of \(\Omega_{s}\) \cite{stratton1941,harrington2001,collin2001}. Such models lie
outside naive \(L^{2}(\Omega)^{3}\) classes yet remain physically realizable when embedded through the
distributional pairing of Subsection~\ref{sec:ch232} \cite{stratton1941}. Subsection~\ref{sec:ch413}
records which distributional excitations preserve \(\mathcal{D}_{\mathrm{real}}=1\) and therefore contribute
to \(\mathcal{S}_{\mathrm{real}}(\mathcal{C})\).

\paragraph{Assumptions.}
Let \(\Sigma\subset\Omega\) be a piecewise smooth surface (aperture, patch, or interface) and
\(\mathbf{J}_{s}\) a tangential surface current density in the sense of
Eq.~\eqref{eq:ch2.distributional-source-extension} \cite{stratton1941}. Volume impressed currents
remain \(\Jimp\in\mathcal{J}_\Omega\); surface data embed as distributions \(\mathbf{J}_{s}\in\mathcal{D}'(\Omega)^{3}\)
acting on test functions \(\boldsymbol{\phi}\in C_{c}^{\infty}(\Omega)^{3}\) \cite{stratton1941,harrington2001}.
Green convolution and self-term regularization follow Subsection~\ref{sec:ch232} without re-derivation
\cite{jackson1999,harrington2001}.

\paragraph{Surface embedding.}
Surface excitation pairs with Eq.~\eqref{eq:ch4.excitation-class-aperture} through
\begin{equation}
\label{eq:ch4.surface-source-embedding}
\langle\mathbf{J}_{\mathrm{tot}},\boldsymbol{\phi}\rangle
=
\int_{\Omega_{s}}\Jimp\cdot\boldsymbol{\phi}\,d^{3}r
+
\sum_{a}\int_{\Sigma_{a}}\mathbf{J}_{s,a}\cdot\boldsymbol{\phi}\,dS,
\end{equation}
extending Eq.~\eqref{eq:ch2.distributional-source-extension} to the realizability tuple
\(\mathfrak{E}\) of Eq.~\eqref{eq:ch4.excitation-tuple} \cite{stratton1941}. Equation
\eqref{eq:ch4.surface-source-embedding} is how aperture feeds enter \(\boldsymbol{\Sigma}_\omega\): the
dyadic convolution of Subsection~\ref{sec:ch231} is evaluated on \(\mathbf{J}_{\mathrm{tot}}\) with
finite-part prescription at coincidence \cite{harrington2001,jackson1999}. Tangential \(\mathbf{J}_{s}\)
replaces volume divergence constraints on \(\Sigma\) by surface continuity of normal \(\D\) and
tangential \(\H\) jumps \cite{stratton1941,collin2001}.

\paragraph{Coupling admissibility.}
Not every square-integrable \(\mathbf{J}_{s}\) produces stable synthesis: near coincidence,
Eq.~\eqref{eq:ch2.dyadic-singular-expansion} demands regularized self-terms. Define
\emph{coupling admissibility} by
\begin{equation}
\label{eq:ch4.coupling-admissibility}
\mathcal{D}_{\mathrm{cpl}}[\mathbf{J}_{s}]
=
1
\iff
\int_{\Sigma}\!\int_{\Sigma}
\big|\overline{\mathbf{G}}_{E,0}(\mathbf{r},\mathbf{r}')\cdot\mathbf{J}_{s}(\mathbf{r}')\big|^{2}\,dS\,dS'
<
\infty
\end{equation}
after static--dynamic split and finite-part extraction as in Subsection~\ref{sec:ch232}
\cite{harrington2001,stratton1941}. Equation~\eqref{eq:ch4.coupling-admissibility} is the surface
analog of \(\HCurl\) control for volume sources: it certifies that \(\boldsymbol{\Sigma}_\omega\) is
Fredholm-stable on the declared mesh or patch basis \cite{harrington2001}. Filamentary currents
\(\mathbf{J}_{f}\) on curves \(\mathcal{C}_{f}\) replace \(dS\) by \(dl\) with line finite parts
\cite{stratton1941,jackson1999}.

\begin{proposition}[Distributional realizability]
\label{prop:ch4.distributional-realizability}
Let \(\mathbf{J}_{\mathrm{tot}}\) be defined by Eq.~\eqref{eq:ch4.surface-source-embedding} with
\(\Jimp\in\mathcal{J}_\Omega\), \(\mathcal{D}_{\mathrm{cpl}}[\mathbf{J}_{s}]=1\), and
\(\mathcal{D}_{\mathrm{real}}=1\) on the combined tuple. Then \(\mathbf{J}_{\mathrm{tot}}\) generates
\(\mathfrak{F}=\boldsymbol{\Sigma}_\omega[\mathbf{J}_{\mathrm{tot}}]\in\mathcal{F}_{\mathrm{rad}}\) and
\(\mathbf{c}=\mathcal{C}(\mathfrak{F})\in\mathcal{S}_{\mathrm{real}}(\mathcal{C})\) \cite{stratton1941,harrington2001}.
Conversely, every \(\mathbf{c}\in\mathcal{S}_{\mathrm{real}}(\mathcal{C})\) arising from aperture-class
excitation Eq.~\eqref{eq:ch4.excitation-class-aperture} admits a distributional representation
Eq.~\eqref{eq:ch4.surface-source-embedding} with \(\mathcal{D}_{\mathrm{cpl}}=1\) after equivalence
reduction of Section~\ref{sec:ch42} \cite{collin2001,devaney2004}.
\end{proposition}
\begin{proof}
Forward direction: Subsection~\ref{sec:ch232} makes \(\mathcal{S}_\omega[\mathbf{J}_{\mathrm{tot}}]\) well-defined;
\(\Pi_{\mathrm{rad}}\) yields outgoing \(\mathfrak{F}\); \(\mathcal{C}\) lists \(\mathbf{c}\) by
Eq.~\eqref{eq:ch4.coefficient-synthesis-map} \cite{stratton1941,harrington2001}. Converse: aperture fields
admit equivalent surface currents by Section~\ref{sec:ch42}; \(\mathcal{D}_{\mathrm{cpl}}=1\) is the
regularity hypothesis under which equivalence preserves \(\boldsymbol{\Gamma}_\omega\)
\cite{collin2001,harrington2001}.
\end{proof}

Proposition~\ref{prop:ch4.distributional-realizability} coins \emph{coupling admissibility} as the surface
realizability certificate: without \(\mathcal{D}_{\mathrm{cpl}}=1\), MoM or Huygens pipelines return
formal coefficients that are not synthesis-stable \cite{harrington2001}. Volume and surface contributions
superpose linearly in Eq.~\eqref{eq:ch4.surface-source-embedding} and in
Eq.~\eqref{eq:ch4.synthesis-superposition} \cite{stratton1941}.

\paragraph{Method-of-moments realizability matrix.}
Discretize \(\mathbf{J}_{s}\) on patches \(\{S_{p}\}_{p=1}^{P}\) with Rao--Wilton--Glisson or pulse bases
\(\{\boldsymbol{\Lambda}_{p}\}\) \cite{harrington2001,collin2001}. The discrete synthesis map is
\begin{equation}
\label{eq:ch4.mom-synthesis-matrix}
\mathbf{c}
=
\mathbf{T}\,\mathbf{I}_{s},
\qquad
T_{\alpha p}
=
\mathcal{C}\Big[\mathcal{R}_\omega\big[\boldsymbol{\Lambda}_{p}\big]\Big]_{\alpha},
\end{equation}
with \(\mathbf{I}_{s}\in\mathbb{C}^{P}\) patch currents and \(\mathbf{T}\in\mathbb{C}^{|I|\times P}\)
\cite{harrington2001}. Equation~\eqref{eq:ch4.mom-synthesis-matrix} is the matrix form of
\(\boldsymbol{\Gamma}_\omega\) on aperture channels; \(\mathcal{D}_{\mathrm{cpl}}=1\) requires
\(\mathbf{T}\) to be Fredholm-stable after self-term regularization of Subsection~\ref{sec:ch232}
\cite{stratton1941,jackson1999}. Ill-conditioned \(\mathbf{T}\) signals coupling inadmissibility, not merely
numerical inconvenience \cite{harrington2001}.

\paragraph{Static--dynamic split at coincidence.}
Following Eq.~\eqref{eq:ch2.dyadic-singular-expansion}, write the patch self-term as
\begin{equation}
\label{eq:ch4.static-dynamic-split}
\overline{\mathbf{G}}_{E,0}
=
\overline{\mathbf{G}}_{\mathrm{sing}}+\overline{\mathbf{G}}_{\mathrm{cpt}},
\qquad
\overline{\mathbf{G}}_{\mathrm{sing}}\ \text{extracted at}\ k_{0}=0,
\end{equation}
with \(\overline{\mathbf{G}}_{\mathrm{cpt}}\) integrable on each patch \cite{harrington2001,stratton1941}.
Equation~\eqref{eq:ch4.static-dynamic-split} is the realizability-preserving regularization: \(\mathcal{D}_{\mathrm{cpl}}=1\)
holds when \(\overline{\mathbf{G}}_{\mathrm{cpt}}\) dominates the MoM impedance matrix after static subtraction
\cite{harrington2001}. Omitting \(\overline{\mathbf{G}}_{\mathrm{sing}}\) produces formal \(\mathbf{c}\) that are not
\(\boldsymbol{\Gamma}_\omega\)-stable under mesh refinement \cite{devaney2004}.

\paragraph{Worked illustration (aperture patch).}
A rectangular aperture \(\Sigma_{a}\) carrying tangential \(\mathbf{J}_{s}\) in the TE\(_{10}\) port mode of a
waveguide excites a definable \(\mathfrak{F}\) once \(\mathcal{D}_{\mathrm{cpl}}=1\) on the patch
\cite{collin2001,harrington2001}. Coefficients \(\mathbf{c}\) read from the far pattern are
\(\boldsymbol{\Gamma}_\omega[\mathbf{J}_{s}]\); claiming additional high-\(\ell\) content without increasing
\(\mathcal{D}_{\mathrm{cpl}}\)-stable degrees of freedom on \(\Sigma_{a}\) violates coupling admissibility
\cite{devaney2004,hansen1988}.

\paragraph{Problem (MoM rank versus target modes).}
\textbf{Problem.} A MoM mesh on \(\Sigma_{a}\) uses \(P=50\) patches; a design targets \(|I|=200\) modal
coefficients above \(\ell_{\max}=8\). Can \(\mathbf{c}_{\mathrm{target}}\) be realizable?

\textbf{Step 1.} Form \(\mathbf{T}\) in Eq.~\eqref{eq:ch4.mom-synthesis-matrix} after self-term regularization
\cite{harrington2001}.

\textbf{Step 2.} Compute \(\operatorname{rank}\mathbf{T}\le\min(P,|I|)=50\): at most fifty independent
synthesis directions are available on the mesh \cite{stratton1941,devaney2004}.

\textbf{Step 3.} If \(\mathbf{c}_{\mathrm{target}}\) requires more than fifty independent components, enlarge
\(\Sigma_{a}\), refine patches, or lower \(\ell_{\max}\) \cite{hansen1988,harrington2001}.

\textbf{Physical meaning.} MoM rank is a discrete realizability ceiling before far-zone measurement
\cite{harrington2001}.

\paragraph{Filamentary and magnetic distributional channels.}
Line currents \(\mathbf{J}_{f}\) on curves \(\mathcal{C}_{f}\subset\Omega\) embed through
\begin{equation}
\label{eq:ch4.filament-embedding}
\langle\mathbf{J}_{\mathrm{tot}},\boldsymbol{\phi}\rangle
\supseteq
\int_{\mathcal{C}_{f}}\mathbf{J}_{f}\cdot\boldsymbol{\phi}\,dl,
\end{equation}
with line finite parts as in Subsection~\ref{sec:ch232} \cite{stratton1941,jackson1999}. Magnetic sheets
\(\mathbf{M}_{s}\) on \(\Sigma_{m}\) enter the same tuple \(\mathfrak{E}\) when \(\E_{t}\) is discontinuous across
\(\Sigma_{m}\) \cite{harrington2001,collin2001}. Realizability requires both \(\mathcal{D}_{\mathrm{cpl}}=1\) and
\(\mathcal{D}_{\mathrm{real}}=1\) on the combined \(\mathbf{J}_{\mathrm{tot}}\) of
Eq.~\eqref{eq:ch4.surface-source-embedding} \cite{stratton1941}.

\begin{figure}[t]
\centering
\ChOneFigBlock{%
\begin{tikzpicture}[ch1fig, node distance=7mm and 9mm]
  \node[ch1box] (vol) {Volume $\Jimp$};
  \node[ch1box, right=11mm of vol] (surf) {Surface $\mathbf{J}_{s}$};
  \node[ch1box, right=11mm of surf] (fil) {Filament $\mathbf{J}_{f}$};
  \node[ch1box, below=12mm of surf] (tot) {$\mathbf{J}_{\mathrm{tot}}$};
  \node[ch1box, below=12mm of fil] (cpl) {$\mathcal{D}_{\mathrm{cpl}}=1$};
  \draw[ch1flow] (vol) -- (tot);
  \draw[ch1flow] (surf) -- (tot);
  \draw[ch1flow] (fil) -- (tot);
  \draw[ch1flow] (cpl) -- (tot);
\end{tikzpicture}%
}
\caption{Subsection~\ref{sec:ch413}: distributional channels superpose in \(\mathbf{J}_{\mathrm{tot}}\); coupling
admissibility is the joint screen on all channels.}
\label{fig:ch4.distributional-channels}
\end{figure}

Figure~\ref{fig:ch4.distributional-channels} is the source-side reading of superposition
Eq.~\eqref{eq:ch4.synthesis-superposition}: each channel contributes to \(\boldsymbol{\Gamma}_\omega\) only when
\(\mathcal{D}_{\mathrm{cpl}}=1\) on its support class \cite{harrington2001,devaney2004}. A formally invertible MoM
matrix on surface patches alone does not certify volume--filament combinations unless the combined embedding satisfies
Eq.~\eqref{eq:ch4.coupling-admissibility} \cite{stratton1941}.

\begin{corollary}[Distributional rank ceiling]
\label{cor:ch4.distributional-rank}
Let \(P_{v},P_{s},P_{f}\) be stable degrees of freedom on volume, surface, and filament discretizations. Then
\(\operatorname{rank}\mathbf{T}\le P_{v}+P_{s}+P_{f}\) and
\(\dim\mathcal{S}_{\mathrm{real}}(\mathcal{C})\le r_{\mathrm{rad}}(\varepsilon)\) regardless of patch count
\cite{harrington2001,devaney2004,stratton1941}.
\end{corollary}

\paragraph{Worked illustration (dipole filament on patch).}
A center-fed dipole filament through a patch array contributes primarily \(\ell=1\) content when
\(k_{0}L_{\mathrm{fil}}\ll1\) \cite{jackson1999,stratton1941}. Adding high-\(\ell\) targets without enlarging
\(\mathcal{C}_{f}\) or \(\Sigma_{a}\) violates coupling admissibility before \(\boldsymbol{\Gamma}_\omega\) is tested
\cite{hansen1988,harrington2001}.

\paragraph{Problem (magnetic sheet necessity).}
\textbf{Problem.} An aperture model uses only \(\mathbf{J}_{s}\) but measured \(\E_{t}\) jumps across a rim. Is
\(\mathcal{D}_{\mathrm{real}}=1\)?

\textbf{Step 1.} Test tangential \(\E\) continuity on the rim per Section~\ref{sec:ch423} \cite{stratton1941}.

\textbf{Step 2.} If a jump is physical, introduce \(\mathbf{M}_{s}\) on the rim with declared
\(\mathcal{D}_{\mathrm{cpl}}=1\) \cite{collin2001,harrington2001}.

\textbf{Step 3.} Reassemble \(\mathbf{J}_{\mathrm{tot}}\) in Eq.~\eqref{eq:ch4.surface-source-embedding} and recompute
\(\mathbf{T}\) \cite{devaney2004}.

\textbf{Physical meaning.} Distributional completeness is a Maxwell predicate, not a mesh option \cite{stratton1941}.

\paragraph{Multi-scale distributional embedding.}
Laboratory feeds often combine volume, surface, and filament channels at disparate scales. Define the
\emph{multi-scale tuple}
\begin{equation}
\label{eq:ch4.multiscale-embedding}
\mathbf{J}_{\mathrm{tot}}
=
\Jimp\,\mathbf{1}_{\Omega_{s}}
+
\sum_{a}\mathbf{J}_{s,a}\,\delta_{\Sigma_{a}}
+
\sum_{f}\mathbf{J}_{f,f}\,\delta_{\mathcal{C}_{f}},
\end{equation}
with each channel tested by Eq.~\eqref{eq:ch4.coupling-admissibility} on its support class
\cite{stratton1941,harrington2001}. Realizability requires joint \(\mathcal{D}_{\mathrm{cpl}}=1\) on the combined
embedding, not channel-wise certification alone \cite{devaney2004}.

\begin{proposition}[Multi-scale realizability factorization]
\label{prop:ch4.multiscale-realizability}
Let \(\mathbf{J}_{\mathrm{tot}}\) be Eq.~\eqref{eq:ch4.multiscale-embedding} with each channel satisfying
\(\mathcal{D}_{\mathrm{real}}=1\) individually. Then \(\mathbf{J}_{\mathrm{tot}}\in\mathcal{J}_{\mathrm{real}}\) if and
only if the combined Green self-energy integral in Eq.~\eqref{eq:ch4.coupling-admissibility} is finite on the union
support and \(\mathcal{D}_{\mathrm{supp}}=1\) \cite{stratton1941,harrington2001,collin2001}. Moreover,
\(\operatorname{rank}\mathbf{T}\le P_{v}+P_{s}+P_{f}\) with equality only when cross-channel couplings are linearly
independent in \(\operatorname{ran}\boldsymbol{\Gamma}_\omega\) \cite{devaney2004,hansen1988}.
\end{proposition}
\begin{proof}
Linearity of \(\mathcal{S}_\omega\) on distributional test classes \cite{stratton1941}; rank bound from
Corollary~\ref{cor:ch4.distributional-rank} \cite{harrington2001}.
\end{proof}

\paragraph{Worked illustration (coaxial feed with rim filament).}
A coaxial center conductor modeled as filament \(\mathcal{C}_{f}\) through a patch rim \(\Sigma_{a}\) requires
finite-part extraction on both supports before \(\mathbf{T}\) is assembled \cite{harrington2001,jackson1999}. Omitting
filament self-terms while retaining patch currents yields \(\mathcal{D}_{\mathrm{cpl}}=0\) on the combined tuple even
when each channel passes in isolation \cite{stratton1941,devaney2004}.

\paragraph{Problem (channel-wise versus joint admissibility).}
\textbf{Problem.} Volume \(\Jimp\), surface \(\mathbf{J}_{s}\), and filament \(\mathbf{J}_{f}\) each satisfy
\(\mathcal{D}_{\mathrm{cpl}}=1\) separately, but MoM on \(\mathbf{J}_{\mathrm{tot}}\) diverges. Which screen failed?

\textbf{Step 1.} Evaluate the joint integral in Eq.~\eqref{eq:ch4.coupling-admissibility} on
\(\operatorname{supp}\mathbf{J}_{\mathrm{tot}}\) \cite{harrington2001}.

\textbf{Step 2.} Identify coincident supports where finite-part extraction was omitted \cite{stratton1941,jackson1999}.

\textbf{Step 3.} Re-regularize cross terms; if divergence persists, remove a channel or refine the mesh
\cite{devaney2004,collin2001}.

\textbf{Physical meaning.} Distributional realizability is a joint predicate on \(\mathbf{J}_{\mathrm{tot}}\), not a
checklist of independent channels \cite{stratton1941}.

\paragraph{Validity.}
Eqs.~\eqref{eq:ch4.surface-source-embedding}--\eqref{eq:ch4.coupling-admissibility} assume Lipschitz
\(\Sigma\) and vacuum exterior unless material overlays are declared \cite{stratton1941}. Magnetic
currents and combined \(\mathbf{J}_{s},\mathbf{M}_{s}\) sheets follow the same distributional pattern with
dual pairings from Subsection~\ref{sec:ch232} \cite{harrington2001}. Point sources require three-dimensional
regularization and are excluded from the default \(\mathcal{J}_{\mathrm{real}}\) unless explicitly regularized
\cite{jackson1999}.

Subsection~\ref{sec:ch414} collects volume, surface, and port channels into the realizable source space
\(\mathcal{J}_{\mathrm{real}}\).


\subsection{Realizable source spaces}
\label{sec:ch414}
% TOC tag: [HOME]

The synthesis set \(\mathcal{S}_{\mathrm{real}}(\mathcal{C})\) is the image of a source manifold, not an
abstract convex body in \(\mathbb{C}^{I}\) \cite{devaney2004}. Subsection~\ref{sec:ch414} defines that
manifold---the \emph{realizable source space}---and proves that \(\boldsymbol{\Gamma}_\omega\) maps it onto
\(\mathcal{S}_{\mathrm{real}}(\mathcal{C})\) with explicit domain qualifiers.

\paragraph{Assumptions.}
Fix \(\mathcal{C}\in\mathcal{R}_{\mathrm{adm}}\), realizability screen
Eq.~\eqref{eq:ch4.realizability-screen}, and distributional extensions
Eq.~\eqref{eq:ch4.surface-source-embedding} \cite{stratton1941,hansen1988}. Excitation classes
Eqs.~\eqref{eq:ch4.excitation-class-volume}--\eqref{eq:ch4.excitation-class-port} are the admissible
channels; equivalence reductions of Section~\ref{sec:ch42} may replace port data by aperture equivalents
\cite{collin2001,harrington2001}.

\begin{definition}[Realizable source space]
\label{def:ch4.realizable-source-space}
The \emph{realizable source space} is
\begin{equation}
\label{eq:ch4.realizable-source-space}
\mathcal{J}_{\mathrm{real}}
=
\left\{
\mathbf{J}_{\mathrm{tot}}:
\mathbf{J}_{\mathrm{tot}}\ \text{satisfies Eqs.~\eqref{eq:ch4.surface-source-embedding}--\eqref{eq:ch4.coupling-admissibility}},\
\mathcal{D}_{\mathrm{real}}[\mathbf{J}_{\mathrm{tot}}]=1
\right\}.
\end{equation}
\end{definition}

Equation~\eqref{eq:ch4.realizable-source-space} is the source-domain preimage of
\(\mathcal{S}_{\mathrm{real}}(\mathcal{C})\): every realizable coefficient vector factors through
\(\mathcal{J}_{\mathrm{real}}\) \cite{stratton1941,devaney2004}. The space is not a Hilbert space in
full generality---it is a union of volume \(\HCurl\) classes, distributional sheets, and port-reduced
equivalents---yet \(\boldsymbol{\Sigma}_\omega\) acts linearly on each channel
\cite{harrington2001,stratton1941}.

\begin{theorem}[Synthesis surjectivity onto the realizable set]
\label{thm:ch4.synthesis-surjectivity}
With \(\boldsymbol{\Gamma}_\omega=\mathcal{C}\circ\boldsymbol{\Sigma}_\omega\) from
Eq.~\eqref{eq:ch4.coefficient-synthesis-map},
\begin{equation}
\label{eq:ch4.synthesis-surjectivity}
\mathcal{S}_{\mathrm{real}}(\mathcal{C})
=
\boldsymbol{\Gamma}_\omega\big[\mathcal{J}_{\mathrm{real}}\big],
\qquad
\mathcal{J}_{\mathrm{real}}\subseteq\mathcal{J}_\Omega\ \text{(distributional union)}.
\end{equation}
Moreover, for \(\mathbf{c}\in\mathcal{S}_{\mathrm{real}}(\mathcal{C})\), there exists
\(\mathbf{J}_{\mathrm{tot}}\in\mathcal{J}_{\mathrm{real}}\) with
\(\boldsymbol{\Gamma}_\omega[\mathbf{J}_{\mathrm{tot}}]=\mathbf{c}\), unique only up to
\(\ker\boldsymbol{\Gamma}_\omega\) \cite{stratton1941,harrington2001,devaney2004}.
\end{theorem}
\begin{proof}
Definition~\ref{def:ch4.realizable-synthesis-set} gives
\(\mathcal{S}_{\mathrm{real}}(\mathcal{C})=\{\mathbf{c}:\exists\Jimp\in\mathcal{J}_\Omega,\
\mathcal{C}(\mathcal{R}_\omega[\Jimp])=\mathbf{c}\}\) \cite{stratton1941}. Definitions
\ref{def:ch4.realizability-screen} and~\ref{def:ch4.realizable-source-space} restrict \(\Jimp\) to
\(\mathcal{J}_{\mathrm{real}}\) without shrinking \(\operatorname{ran}\boldsymbol{\Gamma}_\omega\) on
physical channels: every satisfying tuple in the restricted class maps to the same \(\mathbf{c}\)
\cite{harrington2001}. Surjectivity is definitional; non-uniqueness follows from
\(\ker\mathcal{R}_\omega\) and representation non-injectivity on finite cuts
\cite{stratton1941,devaney2004}.
\end{proof}

Theorem~\ref{thm:ch4.synthesis-surjectivity} is the source-side closure of Section~\ref{sec:ch41}: realizability
is the geometry of \(\mathcal{J}_{\mathrm{real}}\) and its image under \(\boldsymbol{\Gamma}_\omega\), not a
property of \(\mathbb{C}^{I}\) alone \cite{hansen1988}. Inverse synthesis seeks
\(\mathbf{J}_{\mathrm{tot}}\in\mathcal{J}_{\mathrm{real}}\) with
\(\boldsymbol{\Gamma}_\omega[\mathbf{J}_{\mathrm{tot}}]=\mathbf{c}_{\mathrm{target}}\); non-uniqueness is
budgeted by \(\ker\boldsymbol{\Gamma}_\omega\) and by spectral compression of Subsection~\ref{sec:ch402}
\cite{devaney2004,harrington2001}.

\paragraph{Operator-domain form.}
Restricting \(\boldsymbol{\Sigma}_\omega\) to \(\mathcal{J}_{\mathrm{real}}\),
\begin{equation}
\label{eq:ch4.restricted-synthesis-operator}
\boldsymbol{\Sigma}_\omega^{(\mathrm{real})}:
\mathcal{J}_{\mathrm{real}}\to\mathcal{F}_{\mathrm{rad}},
\qquad
\operatorname{rank}\boldsymbol{\Gamma}_\omega|_{\mathcal{J}_{\mathrm{real}}}
\le
\dim\mathcal{S}_{\mathrm{real}}^{(\ell_{\max})}(\mathcal{C})
\end{equation}
for any finite angular cut \(\ell_{\max}\) from Proposition~\ref{prop:ch4.compressed-synthesis}
\cite{hansen1988,harrington2001}. Equation~\eqref{eq:ch4.restricted-synthesis-operator} previews
\emph{radiative rank} of Section~\ref{sec:ch43}: \(\mathcal{J}_{\mathrm{real}}\) is infinite-dimensional as
a set, yet \(\boldsymbol{\Gamma}_\omega\) compresses it to a finite effective rank at fixed geometry and
\(\omega\) \cite{stratton1941,jackson1999}.

\paragraph{Null-space quotient.}
Define the synthesis null space
\begin{equation}
\label{eq:ch4.synthesis-null-space}
\mathcal{N}_{\mathrm{syn}}
=
\ker\boldsymbol{\Gamma}_\omega
=
\left\{
\mathbf{J}_{\mathrm{tot}}\in\mathcal{J}_{\mathrm{real}}:
\boldsymbol{\Gamma}_\omega[\mathbf{J}_{\mathrm{tot}}]=\mathbf{0}
\right\}.
\end{equation}
Equation~\eqref{eq:ch4.synthesis-null-space} contains \(\ker\mathcal{R}_\omega\cap\mathcal{J}_{\mathrm{real}}\) and
any representation-blind currents that produce vanishing coefficients in \(\mathcal{C}\) \cite{stratton1941,harrington2001}.
Realizable synthesis is uniquely defined on the quotient \(\mathcal{J}_{\mathrm{real}}/\mathcal{N}_{\mathrm{syn}}\)
\cite{devaney2004}. Inverse problems regularize on this quotient, not on raw patch currents
\cite{harrington2001}.

\begin{corollary}[Channel count bound]
\label{cor:ch4.channel-count}
Let \(P\) be the number of \(\mathcal{D}_{\mathrm{cpl}}\)-stable MoM degrees of freedom on \(\Sigma_{a}\) and
\(N\) the measurement rank of \(\mathcal{M}\). Then
\begin{equation}
\label{eq:ch4.channel-count-bound}
\dim\boldsymbol{\Gamma}_\omega[\mathcal{J}_{\mathrm{real}}]
\le
\min\big(P,\,N,\,|I_{\ell_{\max}}|\big),
\end{equation}
with \(|I_{\ell_{\max}}|\) the number of indices retained by \(\mathcal{T}_{\ell_{\max}}\)
\cite{hansen1988,harrington2001,devaney2004}.
\end{corollary}
\begin{proof}
MoM rank bounds synthesis degrees of freedom by \(P\); Proposition~\ref{prop:ch4.compressed-synthesis} bounds by
\(|I_{\ell_{\max}}|\); accessibility envelope Definition~\ref{def:ch4.accessibility-envelope} bounds certifiable
coordinates by \(N\) \cite{stratton1941,harrington2001}.
\end{proof}

Corollary~\ref{cor:ch4.channel-count} is the first joint bound on synthesis and certification: realizability and
accessibility compress \(\mathbb{C}^{I}\) from opposite sides of Figure~\ref{fig:ch4.synthesis-accessibility-sets}
\cite{devaney2004}. Section~\ref{sec:ch46} replaces the min bound by effective electromagnetic dimension.

\paragraph{Worked illustration (dipole-only realizability).}
For \(k_{0}D_{s}\ll1\), \(\mathcal{J}_{\mathrm{real}}\) on a single compact patch is well approximated by
three independent dipole components in Eq.~\eqref{eq:ch4.excitation-class-volume}
\cite{jackson1999,stratton1941}. \(\mathcal{S}_{\mathrm{real}}(\mathcal{C})\) then has dimension at most
three in any basis that resolves only \(\ell=1\) content; targeting \(\ell>1\) coefficients without enlarging
\(\mathcal{J}_{\mathrm{real}}\) fails Theorem~\ref{thm:ch4.synthesis-surjectivity}
\cite{hansen1988,harrington2001}.

\paragraph{Problem (quotient ambiguity).}
\textbf{Problem.} Two currents \(\Jimp_{A},\Jimp_{B}\in\mathcal{J}_{\mathrm{real}}\) satisfy
\(\boldsymbol{\Gamma}_\omega[\Jimp_{A}]=\boldsymbol{\Gamma}_\omega[\Jimp_{B}]\) but
\(\Jimp_{A}-\Jimp_{B}\notin\mathcal{J}_{\mathrm{real}}\). Is \(\mathbf{c}\) ill-defined?

\textbf{Step 1.} Form \(\Delta\Jimp=\Jimp_{A}-\Jimp_{B}\); if \(\boldsymbol{\Gamma}_\omega[\Delta\Jimp]=\mathbf{0}\),
then \(\Delta\Jimp\in\mathcal{N}_{\mathrm{syn}}\) \cite{stratton1941}.

\textbf{Step 2.} \(\mathbf{c}\) is well-defined on \(\mathcal{J}_{\mathrm{real}}/\mathcal{N}_{\mathrm{syn}}\) even when
representative currents differ \cite{devaney2004,harrington2001}.

\textbf{Step 3.} Near-field probes sensitive to \(\Delta\Jimp\) cannot resolve \(\mathbf{c}\) alone---reactive
discrimination is outside \(\boldsymbol{\Sigma}_\omega\) \cite{stratton1941}.

\textbf{Physical meaning.} Synthesis uniqueness is a quotient statement, not a current-uniqueness statement
\cite{harrington2001}.

\paragraph{Topology of \(\mathcal{J}_{\mathrm{real}}\).}
\(\mathcal{J}_{\mathrm{real}}\) is not a linear subspace of \(\mathcal{J}_\Omega\): passive constraints and
coupling admissibility are nonlinear in the sense that \(\Jimp_{1}+\Jimp_{2}\) may fail
\(\mathcal{D}_{\mathrm{cpl}}=1\) even when each summand passes \cite{harrington2001,stratton1941}. Synthesis remains
linear on \(\mathcal{F}_{\mathrm{rad}}\) through Eq.~\eqref{eq:ch4.eigenfield-superposition}, but the admissible
\emph{set} \(\mathcal{J}_{\mathrm{real}}\) is a constrained manifold whose tangent directions are the exportable
degrees of freedom counted by \(r_{\mathrm{rad}}(\varepsilon)\) \cite{devaney2004,hansen1988}.

\begin{equation}
\label{eq:ch4.synthesis-manifold-chart}
\mathcal{S}_{\mathrm{real}}(\mathcal{C})
=
\boldsymbol{\Gamma}_\omega\big(\mathcal{J}_{\mathrm{real}}\big)
\subseteq
\boldsymbol{\Gamma}_\omega\big(\mathcal{J}_\Omega\big)
\subseteq
\mathbb{C}^{I},
\end{equation}
with strict inclusions on compact \(\Omega_{s}\) \cite{stratton1941,harrington2001}. Equation
\eqref{eq:ch4.synthesis-manifold-chart} is the coordinate picture for inverse design: optimization must occur on
\(\mathcal{S}_{\mathrm{real}}(\mathcal{C})\), not on \(\mathbb{C}^{I}\) \cite{devaney2004}.

\begin{figure}[t]
\centering
\ChOneFigBlock{%
\begin{tikzpicture}[ch1fig, node distance=7mm and 9mm]
  \node[ch1box] (jo) {$\mathcal{J}_\Omega$};
  \node[ch1box, right=12mm of jo] (jr) {$\mathcal{J}_{\mathrm{real}}$};
  \node[ch1box, right=12mm of jr] (sr) {$\mathcal{S}_{\mathrm{real}}$};
  \node[ch1box, below=11mm of jr] (null) {$\mathcal{N}_{\mathrm{syn}}$};
  \draw[ch1flow] (jo) -- (jr) -- (sr);
  \draw[ch1flow] (null) -- (jr);
\end{tikzpicture}%
}
\caption{Subsection~\ref{sec:ch414}: \(\mathcal{J}_{\mathrm{real}}\) is the admissible source manifold;
\(\mathcal{S}_{\mathrm{real}}(\mathcal{C})\) is its synthesis image modulo \(\mathcal{N}_{\mathrm{syn}}\).}
\label{fig:ch4.realizable-source-manifold}
\end{figure}

\paragraph{Active and lossy extensions (declared domain).}
When \(\mathcal{D}_{\mathrm{pass}}=0\) by design, \(\mathcal{J}_{\mathrm{real}}\) must include explicit port power
channels in \(\mathfrak{E}\) \cite{harrington2001,collin2001}. Loss modifies \(\sigma_{n}\) tails without removing
the manifold chart Eq.~\eqref{eq:ch4.synthesis-manifold-chart}: realizability remains
\(\operatorname{ran}\boldsymbol{\Gamma}_\omega\) on the declared class \cite{stratton1941,devaney2004}.

\paragraph{Problem (nonlinear superposition).}
\textbf{Problem.} Two patch currents \(\mathbf{J}_{s,1},\mathbf{J}_{s,2}\) each satisfy \(\mathcal{D}_{\mathrm{cpl}}=1\), but
\(\mathbf{J}_{s,1}+\mathbf{J}_{s,2}\) fails MoM stability. Is \(\mathbf{c}_{1}+\mathbf{c}_{2}\) realizable?

\textbf{Step 1.} Test \(\mathbf{J}_{s,1}+\mathbf{J}_{s,2}\in\mathcal{J}_{\mathrm{real}}\) via
Eq.~\eqref{eq:ch4.coupling-admissibility} \cite{harrington2001}.

\textbf{Step 2.} If \(\mathcal{D}_{\mathrm{cpl}}=0\), superposition of individually admissible channels does not certify
\(\mathbf{c}_{1}+\mathbf{c}_{2}\) \cite{stratton1941,devaney2004}.

\textbf{Step 3.} Refine mesh or split into separate synthesis problems with joint range test
\cite{harrington2001}.

\textbf{Physical meaning.} \(\mathcal{J}_{\mathrm{real}}\) is a constrained set, not a vector space \cite{stratton1941}.

\paragraph{Tangent realizability directions.}
At \(\Jimp_{0}\in\mathcal{J}_{\mathrm{real}}\), define the tangent export cone
\begin{equation}
\label{eq:ch4.tangent-export-cone}
\mathcal{T}_{\Jimp_{0}}\mathcal{J}_{\mathrm{real}}
=
\overline{\big\{
\boldsymbol{\Gamma}_\omega[\delta\Jimp]:
\delta\Jimp\ \text{admissible first variation at}\ \Jimp_{0}
\big\}},
\end{equation}
with closure in \(\mathbb{C}^{I}\) \cite{stratton1941,harrington2001,devaney2004}. Linearization of inverse design is
well-posed only on \(\mathcal{T}_{\Jimp_{0}}\mathcal{J}_{\mathrm{real}}\), not on full \(\mathbb{C}^{I}\)
\cite{hansen1988}. The cone dimension is bounded by \(r_{\mathrm{rad}}(\varepsilon)\) at tolerance \(\varepsilon\)
\cite{harrington2001}.

\begin{proposition}[Tangent dimension bound]
\label{prop:ch4.tangent-dimension}
\(\dim\mathcal{T}_{\Jimp_{0}}\mathcal{J}_{\mathrm{real}}\le r_{\mathrm{rad}}(\varepsilon)\) for every
\(\Jimp_{0}\in\mathcal{J}_{\mathrm{real}}\) and aligned gauge \cite{stratton1941,harrington2001,devaney2004}. Strict
inequality holds on rank-deficient feeds where \(\sigma_{n}(\boldsymbol{\Gamma}_\omega)<\varepsilon\sigma_{1}\) for
\(n>r_{\mathrm{rad}}(\varepsilon)\) \cite{hansen1988}.
\end{proposition}

\paragraph{Worked illustration (phase-only gradient).}
Gradient ascent on phase-only parameters explores a submanifold of \(\mathcal{J}_{\mathrm{real}}\) with tangent span
smaller than \(\mathcal{T}_{\Jimp_{0}}\mathcal{J}_{\mathrm{real}}\) unless amplitude and coupling degrees of freedom are
included \cite{harrington2001,devaney2004,stratton1941}.

\paragraph{Problem (local versus global synthesis).}
\textbf{Problem.} A local optimizer reports small residual near \(\Jimp_{0}\) but \(\mathbf{c}_{\mathrm{target}}\notin
\mathcal{S}_{\mathrm{real}}(\mathcal{C})\). Explain.

\textbf{Step 1.} Test \(\delta_{\mathrm{syn}}\) from Eq.~\eqref{eq:ch4.synthesis-residual-norm} globally
\cite{stratton1941,harrington2001}.

\textbf{Step 2.} Compare tangent cone dimension to target residual rank \cite{devaney2004,hansen1988}.

\textbf{Step 3.} Conclude local stationarity does not imply global realizability \cite{stratton1941}.

\textbf{Physical meaning.} Inverse design is a constrained global range problem, not a local phase fit \cite{harrington2001}.

\paragraph{Validity.}
Theorem~\ref{thm:ch4.synthesis-surjectivity} holds in linear Maxwell theory with screens
Eq.~\eqref{eq:ch4.realizability-screen} \cite{stratton1941}. Nonlinear or active extensions enlarge
\(\mathcal{J}_{\mathrm{real}}\) by explicit power-injection channels not covered by
Eq.~\eqref{eq:ch4.passivity-screen} \cite{harrington2001}. Uniqueness up to
\(\ker\boldsymbol{\Gamma}_\omega\) is the synthesis analogue of non-radiating currents in
Section~\ref{sec:ch27} \cite{stratton1941}.

Section~\ref{sec:ch41} has defined \(\mathcal{J}_{\mathrm{real}}\) and identified
\(\mathcal{S}_{\mathrm{real}}(\mathcal{C})=\boldsymbol{\Gamma}_\omega[\mathcal{J}_{\mathrm{real}}]\).
Section~\ref{sec:ch42} replaces volume currents by equivalent aperture models while preserving synthesis
images on \(\operatorname{ran}\mathcal{R}_\omega\) \cite{collin2001,harrington2001}.


\section{Equivalence Principles and Apertures}
\label{sec:ch42}

Volume impressed currents \(\Jimp\) on \(\Omega_{s}\) are the primitive synthesis carriers of
Section~\ref{sec:ch41}, yet laboratories often specify fields on an aperture \(\Sigma_{a}\) or on a
closed surface \(S\) enclosing the scatterer \cite{stratton1941,harrington2001,collin2001}. Equivalence
principles identify surface densities that reproduce the same exterior radiation while shrinking the
degrees of freedom that \(\boldsymbol{\Gamma}_\omega\) must invert \cite{stratton1941}. Section~\ref{sec:ch42}
develops that identification as a realizability tool: Subsection~\ref{sec:ch421} states volume--surface
equivalence for exterior fields; Subsection~\ref{sec:ch422} treats apertures as equivalent sources in
\(\mathcal{J}_{\mathrm{real}}\); Subsection~\ref{sec:ch423} imports tangential continuity from
Section~\ref{sec:ch01.1.2}; Subsection~\ref{sec:ch424} records the limits under which equivalence fails
to preserve \(\mathcal{S}_{\mathrm{real}}(\mathcal{C})\) \cite{devaney2004,hansen1988}.


\subsection{Volume-surface equivalence}
\label{sec:ch421}
% TOC tag: [HOME]

Finite scatterers are frequently analyzed by replacing interior excitation with equivalent currents on a
surface \(S\) that encloses \(\Omega_{s}\) \cite{stratton1941,harrington2001}. The replacement is physically
meaningful only if exterior fields---and therefore \(\boldsymbol{\Sigma}_\omega[\Jimp]\) and
\(\boldsymbol{\Gamma}_\omega[\Jimp]\)---remain unchanged in the outgoing region
\cite{collin2001,devaney2004}. Subsection~\ref{sec:ch421} states that invariance as a realizability theorem.

\paragraph{Assumptions.}
Let \(\Omega_{\mathrm{int}}\Subset\mathbb{R}^{3}\) be a bounded interior region with Lipschitz boundary
\(S=\partial\Omega_{\mathrm{int}}\), exterior \(\Omega_{\mathrm{ext}}=\mathbb{R}^{3}\setminus\overline{\Omega}_{\mathrm{int}}\),
and impressed volume current \(\Jimp\) supported in \(\Omega_{s}\subset\Omega_{\mathrm{int}}\)
\cite{stratton1941}. Fields satisfy harmonic Maxwell equations in \(\Omega_{\mathrm{ext}}\) with outgoing
Sommerfeld closure as in Section~\ref{sec:ch24} \cite{stratton1941,harrington2001}. Equivalence currents
\(\mathbf{J}_{s},\mathbf{M}_{s}\) on \(S\) are tangential surface densities in the sense of
Eq.~\eqref{eq:ch4.surface-source-embedding} \cite{stratton1941}.

\paragraph{Exterior field equivalence.}
Denote by \(\mathcal{S}_\omega[\Jimp]\) the full Maxwell solution pair and by
\(\mathcal{F}^{\mathrm{(ext)}}_{\mathrm{rad}}\) the restriction of \(\Pi_{\mathrm{rad}}\mathcal{S}_\omega[\Jimp]\) to
\(\Omega_{\mathrm{ext}}\) \cite{stratton1941,harrington2001}. Volume--surface equivalence requires
\begin{equation}
\label{eq:ch4.exterior-equivalence-condition}
\big(\E^{\mathrm{(eq)}},\H^{\mathrm{(eq)}}\big)\big|_{\Omega_{\mathrm{ext}}}
=
\big(\E^{\mathrm{(vol)}},\H^{\mathrm{(vol)}}\big)\big|_{\Omega_{\mathrm{ext}}},
\end{equation}
where \((\E^{\mathrm{(vol)}},\H^{\mathrm{(vol)}})=\mathcal{S}_\omega[\Jimp]\) and
\((\E^{\mathrm{(eq)}},\H^{\mathrm{(eq)}})\) is generated by \((\mathbf{J}_{s},\mathbf{M}_{s})\) on \(S\) alone
\cite{stratton1941,collin2001}. Equation~\eqref{eq:ch4.exterior-equivalence-condition} is the field statement of
equivalence: interior mechanisms may change, but \(\mathfrak{F}=\mathcal{R}_\omega[\Jimp]\) on
\(\operatorname{ran}\mathcal{R}_\omega\) is fixed \cite{harrington2001}. Coefficient listings satisfy
\begin{equation}
\label{eq:ch4.coefficient-equivalence-invariance}
\mathcal{C}\big(\mathcal{R}_\omega[\Jimp]\big)
=
\mathcal{C}\big(\mathcal{R}_\omega[\mathbf{J}_{\mathrm{eq}}]\big),
\qquad
\mathbf{J}_{\mathrm{eq}}\ \text{the equivalent surface source},
\end{equation}
whenever Eq.~\eqref{eq:ch4.exterior-equivalence-condition} holds and \(\mathcal{C}\in\mathcal{R}_{\mathrm{adm}}\)
\cite{hansen1988,stratton1941}. Realizability of \(\mathbf{c}\) is unchanged by equivalence reduction.

\begin{theorem}[Volume--surface equivalence for exterior radiation]
\label{thm:ch4.volume-surface-equivalence}
Let \(\Jimp\) be admissible with \(\operatorname{supp}\Jimp\subset\Omega_{\mathrm{int}}\), and let
\((\mathbf{J}_{s},\mathbf{M}_{s})\) be equivalent surface sources on \(S=\partial\Omega_{\mathrm{int}}\) in the
sense of Stratton--Chu / Love equivalence \cite{stratton1941,harrington2001}. Then
Eq.~\eqref{eq:ch4.exterior-equivalence-condition} holds in \(\Omega_{\mathrm{ext}}\), and
\begin{equation}
\label{eq:ch4.gamma-equivalence-invariance}
\boldsymbol{\Gamma}_\omega[\Jimp]
=
\boldsymbol{\Gamma}_\omega[\mathbf{J}_{\mathrm{eq}}]
\end{equation}
for every \(\mathcal{C}\in\mathcal{R}_{\mathrm{adm}}\) \cite{stratton1941,collin2001,devaney2004}.
\end{theorem}
\begin{proof}
Equivalence theorems construct \((\mathbf{J}_{s},\mathbf{M}_{s})\) from tangential traces of
\((\E^{\mathrm{(vol)}},\H^{\mathrm{(vol)}})\) on \(S\) such that the exterior boundary-value problem with
outgoing radiation condition reproduces the same \(\E,\H\) in \(\Omega_{\mathrm{ext}}\)
\cite{stratton1941,harrington2001}. Uniqueness of outgoing exterior solutions
Theorem~\ref{thm:ch2.exterior-radiative-uniqueness} fixes the radiative state
\(\mathfrak{F}\in\mathcal{F}_{\mathrm{rad}}\); \(\mathcal{C}\) is a coordinate map on that state, hence
Eq.~\eqref{eq:ch4.gamma-equivalence-invariance} \cite{stratton1941,devaney2004}.
\end{proof}

Theorem~\ref{thm:ch4.volume-surface-equivalence} is the first structural simplification of realizability:
synthesis may be computed on \(S\) without modeling feed geometry inside \(\Omega_{\mathrm{int}}\), provided
equivalence hypotheses hold \cite{collin2001}. The theorem does not assert uniqueness of
\((\mathbf{J}_{s},\mathbf{M}_{s})\): distinct surface pairs can yield the same exterior
\(\mathfrak{F}\), enlarging \(\mathcal{N}_{\mathrm{syn}}\) of Eq.~\eqref{eq:ch4.synthesis-null-space}
\cite{harrington2001,devaney2004}.

\paragraph{Stratton--Chu construction (statement).}
On \(S\), define equivalent electric and magnetic surface currents from tangential traces
\cite{stratton1941,harrington2001}:
\begin{equation}
\label{eq:ch4.stratton-chu-currents}
\mathbf{J}_{s}
=
\hat{\mathbf{n}}\times\H^{\mathrm{(vol)}},
\qquad
\mathbf{M}_{s}
=
-\hat{\mathbf{n}}\times\E^{\mathrm{(vol)}},
\end{equation}
with outward \(\hat{\mathbf{n}}\) on \(S\) \cite{stratton1941}. Equation~\eqref{eq:ch4.stratton-chu-currents} is the
standard Huygens--Stratton prescription: \(\mathbf{J}_{s}\) replaces tangential \(\H\) discontinuity and
\(\mathbf{M}_{s}\) replaces tangential \(\E\) discontinuity across a fictitious surface supporting the same
exterior radiation \cite{harrington2001,collin2001}. Coupling admissibility
Eq.~\eqref{eq:ch4.coupling-admissibility} must be verified on the chosen mesh of \(S\), not assumed from
Eq.~\eqref{eq:ch4.stratton-chu-currents} alone \cite{harrington2001}.

\begin{figure}[t]
\centering
\ChOneFigBlock{%
\begin{tikzpicture}[ch1fig, node distance=7mm and 10mm]
  \node[ch1box] (vol) {Volume $\Jimp$\\in $\Omega_{s}$};
  \node[ch1box, right=14mm of vol] (s) {Surface\\$(\mathbf{J}_{s},\mathbf{M}_{s})$ on $S$};
  \node[ch1box, right=14mm of s] (ext) {Exterior $\mathfrak{F}$\\unchanged};
  \node[ch1box, below=12mm of ext] (c) {$\mathbf{c}=\boldsymbol{\Gamma}_\omega[\cdot]$};
  \draw[ch1flow] (vol) -- node[above,font=\scriptsize] {equivalence} (s);
  \draw[ch1flow] (s) -- (ext);
  \draw[ch1flow] (vol) to[bend left=15] (ext);
  \draw[ch1flow] (c) -- (ext);
\end{tikzpicture}%
}
\caption{Subsection~\ref{sec:ch421} volume--surface equivalence: distinct interior models that satisfy
Eq.~\eqref{eq:ch4.exterior-equivalence-condition} yield the same \(\mathfrak{F}\) and \(\mathbf{c}\).}
\label{fig:ch4.volume-surface-equivalence}
\end{figure}

Figure~\ref{fig:ch4.volume-surface-equivalence} is the realizability reading of equivalence: the synthesis
manifold \(\boldsymbol{\Gamma}_\omega[\mathcal{J}_{\mathrm{real}}]\) is invariant under passage from \(\Jimp\) to
\((\mathbf{J}_{s},\mathbf{M}_{s})\) when Theorem~\ref{thm:ch4.volume-surface-equivalence} applies
\cite{devaney2004}. MoM implementations discretize \(S\) and form \(\mathbf{T}\) of
Eq.~\eqref{eq:ch4.mom-synthesis-matrix} on surface patches rather than volume voxels
\cite{harrington2001,collin2001}.

\paragraph{Worked illustration (PEC-enclosed feed).}
A dipole feed inside a PEC cavity induces tangential \(\H\) on an aperture \(S_{a}\subset S\).
Equivalence replaces the interior \(\Jimp\) by \(\mathbf{J}_{s}=\hat{\mathbf{n}}\times\H\) on \(S_{a}\) with
\(\mathbf{M}_{s}=\mathbf{0}\) on PEC portions of \(S\) \cite{stratton1941,collin2001}. Far-zone \(\mathbf{c}\) computed
from aperture currents matches volume synthesis when \(\mathcal{D}_{\mathrm{cpl}}=1\) on \(S_{a}\)
\cite{harrington2001}.

\paragraph{Stratton--Chu construction (realizability form).}
On \(S=\partial\Omega_{\mathrm{int}}\), equivalent electric and magnetic surface densities are fixed by exterior
tangential traces through
\begin{equation}
\label{eq:ch4.stratton-chu-currents}
\mathbf{J}_{s}
=
\hat{\mathbf{n}}\times\H^{\mathrm{(ext)}},
\qquad
\mathbf{M}_{s}
=
-\hat{\mathbf{n}}\times\E^{\mathrm{(ext)}},
\end{equation}
with \(\hat{\mathbf{n}}\) outward from \(\Omega_{\mathrm{int}}\) \cite{stratton1941,harrington2001}. Equations
\eqref{eq:ch4.stratton-chu-currents} are not design freedoms: they are the unique surface tuple that reproduces
\(\mathfrak{F}\) in \(\Omega_{\mathrm{ext}}\) when interior sources are annihilated and \(S\) encloses all
\(\operatorname{supp}\Jimp\) \cite{collin2001}. In synthesis language,
\(\boldsymbol{\Gamma}_\omega[\mathbf{J}_{\mathrm{eq}}]\) with \(\mathbf{J}_{\mathrm{eq}}\) assembled from
Eqs.~\eqref{eq:ch4.stratton-chu-currents} equals \(\boldsymbol{\Gamma}_\omega[\Jimp]\) whenever
Theorem~\ref{thm:ch4.volume-surface-equivalence} applies \cite{stratton1941}.

\paragraph{Love reduction on PEC.}
When \(S\) is PEC except for an aperture \(S_{a}\), \(\hat{\mathbf{n}}\times\E=\mathbf{0}\) on the conductor and
\(\mathbf{M}_{s}=\mathbf{0}\) on PEC portions of Eq.~\eqref{eq:ch4.stratton-chu-currents}
\cite{stratton1941,collin2001}. The realizability reduction is then purely electric:
\(\mathbf{J}_{s}=\hat{\mathbf{n}}\times\H\) on \(S_{a}\) with \(\mathbf{J}_{s}=\mathbf{0}\) on the PEC remainder
\cite{harrington2001}. MoM discretizations that introduce magnetic currents on PEC without aperture justification
inflate \(\ker\boldsymbol{\Gamma}_\omega\) without enlarging \(\mathcal{S}_{\mathrm{real}}(\mathcal{C})\)
\cite{devaney2004}.

\begin{corollary}[Equivalence invariance of radiative rank]
\label{cor:ch4.equivalence-rank-invariance}
If Theorem~\ref{thm:ch4.volume-surface-equivalence} holds and \(\mathcal{D}_{\mathrm{cpl}}=1\) on \(S\), then
\(r_{\mathrm{rad}}(\varepsilon)\) of Definition~\ref{def:ch4.radiative-rank} computed from volume
\(\Jimp\in\mathcal{J}_{\mathrm{real}}\) equals \(r_{\mathrm{rad}}(\varepsilon)\) computed from equivalent
\((\mathbf{J}_{s},\mathbf{M}_{s})\) on \(S\) \cite{stratton1941,harrington2001,devaney2004}.
\end{corollary}
\begin{proof}
\(\boldsymbol{\Gamma}_\omega[\Jimp]=\boldsymbol{\Gamma}_\omega[\mathbf{J}_{\mathrm{eq}}]\) by
Eq.~\eqref{eq:ch4.gamma-equivalence-invariance}; singular values of \(\boldsymbol{\Gamma}_\omega\) on the
equivalence-reduced domain coincide when coupling is admissible \cite{harrington2001}.
\end{proof}

Corollary~\ref{cor:ch4.equivalence-rank-invariance} states that equivalence is a \emph{representation change} on
\(\mathcal{J}_{\mathrm{real}}\), not a rank enlargement: discretizing \(S\) instead of \(\Omega_{s}\) may change MoM
patch count but not continuum export rank when \(\mathcal{D}_{\mathrm{cpl}}=1\) \cite{harrington2001,collin2001}.

\paragraph{Problem (enclosure test).}
\textbf{Problem.} A volume current \(\Jimp\) has support in \(\Omega_{s}\). An engineer places equivalence surface
\(S'\) that fails to enclose a portion of \(\Omega_{s}\). Can \(\mathbf{c}=\boldsymbol{\Gamma}_\omega[\Jimp]\) be
recovered from equivalent currents on \(S'\)?

\textbf{Step 1.} Identify the exterior field region \(\Omega_{\mathrm{ext}}(S')\) where
Eq.~\eqref{eq:ch4.exterior-equivalence-condition} is asserted \cite{stratton1941}.

\textbf{Step 2.} Unenclosed source segments radiate directly into regions not determined by tangential data on \(S'\);
\(\boldsymbol{\Gamma}_\omega[\mathbf{J}_{\mathrm{eq}}]\) on \(\mathcal{F}_{\mathrm{rad}}\) differs from the true
\(\mathbf{c}\) \cite{harrington2001,collin2001}.

\textbf{Step 3.} Enlarge \(S'\) to contain \(\Omega_{s}\) or restrict \(\mathcal{C}\) to the partial exterior where
equivalence is valid \cite{devaney2004}.

\textbf{Physical meaning.} Enclosure is a realizability hypothesis, not a mesh convenience \cite{stratton1941}.

\paragraph{Coefficient invariance under equivalent currents.}
When Theorem~\ref{thm:ch4.volume-surface-equivalence} holds,
\begin{equation}
\label{eq:ch4.coefficient-invariance-equivalence}
\mathcal{C}\big(\mathcal{R}_\omega[\Jimp]\big)
=
\mathcal{C}\big(\mathcal{R}_\omega[\mathbf{J}_{\mathrm{eq}}]\big)
\quad\forall\,\mathcal{C}\in\mathcal{R}_{\mathrm{adm}},
\end{equation}
extending Eq.~\eqref{eq:ch4.coefficient-equivalence-invariance} to every admissible listing
\cite{stratton1941,hansen1988,harrington2001}. Equation~\eqref{eq:ch4.coefficient-invariance-equivalence} is the
realizability certificate for equivalence-based design: optimizing \(\mathbf{J}_{\mathrm{eq}}\) on \(S\) is equivalent to
optimizing \(\Jimp\) in volume only when enclosure and \(\mathcal{D}_{\mathrm{cpl}}=1\) hold \cite{devaney2004,collin2001}.

\begin{corollary}[MoM equivalence consistency]
\label{cor:ch4.mom-equivalence-consistency}
If volume and surface MoM models satisfy Theorem~\ref{thm:ch4.volume-surface-equivalence} and share the same
\(\mathcal{N},\mathcal{C}\), then \(\mathbf{T}_{\mathrm{vol}}\mathbf{I}_{v}\) and
\(\mathbf{T}_{\mathrm{surf}}\mathbf{I}_{s}\) yield identical \(\mathbf{c}\) up to \(\mathcal{N}_{\mathrm{syn}}\)
\cite{harrington2001,stratton1941,collin2001}.
\end{corollary}

\paragraph{Worked illustration (surface-only inversion).}
Inverting \(\mathbf{T}_{\mathrm{surf}}\) for \(\mathbf{c}_{\mathrm{target}}\) without verifying enclosure produces
formal surface currents whose \(\mathbf{c}\) matches the target only in the partial exterior domain; outside that domain,
Eq.~\eqref{eq:ch4.exterior-equivalence-condition} fails \cite{stratton1941,devaney2004,harrington2001}.

\paragraph{Problem (Love versus full equivalence).}
\textbf{Problem.} A PEC body uses Love currents on aperture \(S_{a}\) only. A target \(\mathbf{c}\) requires magnetic
sheet content on PEC portions. Can Love reduction synthesize \(\mathbf{c}\)?

\textbf{Step 1.} Test whether \(\mathbf{c}\) depends on \(\mathbf{M}_{s}\) sectors eliminated by Love reduction
\cite{stratton1941,collin2001}.

\textbf{Step 2.} If yes, restore full Stratton--Chu tuple Eq.~\eqref{eq:ch4.stratton-chu-currents} on \(S\)
\cite{harrington2001}.

\textbf{Step 3.} Re-verify \(\mathcal{D}_{\mathrm{cpl}}=1\) on the enlarged surface mesh \cite{devaney2004}.

\textbf{Physical meaning.} Equivalence reductions shrink the search space, not \(\operatorname{ran}\boldsymbol{\Gamma}_\omega\)
\cite{stratton1941}.

\paragraph{Equivalence rank on surface meshes.}
When Theorem~\ref{thm:ch4.volume-surface-equivalence} holds,
\begin{equation}
\label{eq:ch4.equivalence-rank-equality}
r_{\mathrm{rad}}^{\mathrm{(vol)}}(\varepsilon)
=
r_{\mathrm{rad}}^{\mathrm{(surf)}}(\varepsilon),
\end{equation}
as in Corollary~\ref{cor:ch4.equivalence-rank-invariance}, but MoM ranks may differ:
\(\operatorname{rank}\mathbf{T}_{\mathrm{vol}}\neq\operatorname{rank}\mathbf{T}_{\mathrm{surf}}\) until
\(\mathcal{D}_{\mathrm{cpl}}=1\) on both meshes \cite{harrington2001,collin2001,stratton1941}. Surface reduction is a
search-space reduction, not a rank inflation mechanism \cite{devaney2004}.

\begin{proposition}[Love reduction domain]
\label{prop:ch4.love-reduction-domain}
Love equivalence on PEC restricts \(\mathbf{J}_{\mathrm{eq}}\) to electric currents on apertures, shrinking the discrete
search basis on \(S\) without shrinking \(\operatorname{ran}\boldsymbol{\Gamma}_\omega\) when full equivalence is valid
\cite{stratton1941,harrington2001,collin2001}. Targets requiring \(\mathbf{M}_{s}\neq0\) on PEC are outside the Love
domain \cite{devaney2004}.
\end{proposition}

\paragraph{Worked illustration (dielectric radome).}
Love reduction on a PEC-only model of a dielectric radome omits \(\mathbf{M}_{s}\) and underestimates
\(\mathcal{K}_{\mathrm{eq}}\) in Theorem~\ref{thm:ch4.equivalence-error} \cite{stratton1941,harrington2001,collin2001}.

\paragraph{Problem (surface mesh rank versus volume).}
\textbf{Problem.} \(\operatorname{rank}\mathbf{T}_{\mathrm{surf}}=60\) but \(\operatorname{rank}\mathbf{T}_{\mathrm{vol}}=90\) at the same
\(\mathbf{c}\). Does surface equivalence fail?

\textbf{Step 1.} Test Theorem~\ref{thm:ch4.volume-surface-equivalence} hypotheses \cite{stratton1941}.

\textbf{Step 2.} Compare continuum \(r_{\mathrm{rad}}(\varepsilon)\), not mesh ranks alone \cite{harrington2001,devaney2004}.

\textbf{Step 3.} Refine surface mesh if \(\mathcal{D}_{\mathrm{cpl}}=0\); otherwise ranks may differ while \(\mathbf{c}\) matches
\cite{collin2001,hansen1988}.

\textbf{Physical meaning.} Equivalence preserves \(\mathbf{c}\), not patch-count parity \cite{stratton1941}.

\paragraph{Validity.}
Theorem~\ref{thm:ch4.volume-surface-equivalence} requires linear Maxwell media, fixed \(\omega\), and outgoing
\(\Omega_{\mathrm{ext}}\) \cite{stratton1941}. Lossy interiors modify equivalent currents but preserve exterior
equivalence when \(S\) encloses all loss \cite{harrington2001}. Equivalence on \(S\) does not determine interior
reactive storage---only \(\mathfrak{F}\) on \(\mathcal{F}_{\mathrm{rad}}\) is certified \cite{stratton1941}.

Apertures as laboratory realizations of equivalent sources are developed in Subsection~\ref{sec:ch422}.


\subsection{Apertures as equivalent sources}
\label{sec:ch422}
% TOC tag: [HOME]

Antenna engineering speaks of \emph{apertures}---openings through which radiation passes---rather than of
closed equivalence surfaces \(S\) \cite{collin2001,harrington2001,stratton1941}. An aperture \(\Sigma_{a}\)
is realizable when tangential fields on \(\Sigma_{a}\) define a \(\mathcal{D}_{\mathrm{cpl}}\)-admissible
surface current \(\mathbf{J}_{s}\) whose synthesis image \(\boldsymbol{\Gamma}_\omega[\mathbf{J}_{s}]\) matches the
radiating state of the full structure \cite{stratton1941,devaney2004}.

\paragraph{Assumptions.}
Let \(\Sigma_{a}\subset S\) be a smooth aperture in an otherwise PEC or impedance boundary
\cite{collin2001,stratton1941}. Tangential \(\E_{t}\) and \(\H_{t}\) on \(\Sigma_{a}\) are traces of an
admissible exterior solution \cite{stratton1941}. Port feeds may reduce to \(\mathbf{J}_{s}\) on
\(\Sigma_{a}\) by equivalence of Subsection~\ref{sec:ch421} \cite{harrington2001}.

\begin{definition}[Aperture equivalent source]
\label{def:ch4.aperture-equivalent-source}
An \emph{aperture equivalent source} is a tangential surface current \(\mathbf{J}_{s}\) on \(\Sigma_{a}\) such
that
\begin{equation}
\label{eq:ch4.aperture-equivalent-def}
\boldsymbol{\Sigma}_\omega[\mathbf{J}_{s}]
=
\boldsymbol{\Sigma}_\omega[\Jimp^{\mathrm{(full)}}]
\quad\text{in}\ \mathcal{F}_{\mathrm{rad}},
\end{equation}
for some full-model impressed current \(\Jimp^{\mathrm{(full)}}\) generating the same exterior
\(\mathfrak{F}\) \cite{stratton1941,collin2001}. The aperture synthesis set is
\begin{equation}
\label{eq:ch4.aperture-synthesis-set}
\mathcal{S}_{\mathrm{ap}}(\mathcal{C})
=
\left\{
\mathbf{c}\in\mathcal{S}_{\mathrm{real}}(\mathcal{C}):
\exists\,\mathbf{J}_{s}\ \text{on}\ \Sigma_{a}\ \text{with}\
\mathcal{C}\big(\mathcal{R}_\omega[\mathbf{J}_{s}]\big)=\mathbf{c}
\right\}.
\end{equation}
\end{definition}

Equation~\eqref{eq:ch4.aperture-synthesis-set} is the restriction of
\(\mathcal{S}_{\mathrm{real}}(\mathcal{C})\) to excitation class
Eq.~\eqref{eq:ch4.excitation-class-aperture} \cite{harrington2001}. Physically, \(\mathcal{S}_{\mathrm{ap}}(\mathcal{C})\)
is the set of far-zone listings producible by currents on the aperture alone, without resolving interior feed
detail \cite{collin2001,devaney2004}. \(\mathcal{S}_{\mathrm{ap}}(\mathcal{C})\subseteq\mathcal{S}_{\mathrm{real}}(\mathcal{C})\);
strict inclusion holds when interior degrees of freedom synthesize patterns not realizable by \(\Sigma_{a}\) alone
\cite{stratton1941,hansen1988}.

\paragraph{Radiating versus scattering apertures.}
A \emph{radiating} aperture carries impressed \(\mathbf{J}_{s}\) from a feed; a \emph{scattering} aperture
carries induced \(\mathbf{J}_{s}\) from incident \(\E^{\mathrm{(inc)}}\) \cite{stratton1941,collin2001}. In both
cases realizability is judged on \(\boldsymbol{\Gamma}_\omega[\mathbf{J}_{s}]\), not on incident field labels
\cite{harrington2001}. For scattering,
\begin{equation}
\label{eq:ch4.scattering-aperture-decomposition}
\mathbf{J}_{s}
=
\mathbf{J}_{s}^{\mathrm{(inc)}}+\mathbf{J}_{s}^{\mathrm{(ind)}},
\qquad
\boldsymbol{\Gamma}_\omega[\mathbf{J}_{s}]
=
\boldsymbol{\Gamma}_\omega[\mathbf{J}_{s}^{\mathrm{(ind)}}]
\end{equation}
when \(\mathbf{J}_{s}^{\mathrm{(inc)}}\) is a known calibration term \cite{stratton1941,harrington2001}. Equation
\eqref{eq:ch4.scattering-aperture-decomposition} separates calibration from synthesis targets in inverse
problems \cite{devaney2004}.

\begin{proposition}[Aperture realizability reduction]
\label{prop:ch4.aperture-realizability-reduction}
If Theorem~\ref{thm:ch4.volume-surface-equivalence} holds with equivalence surface \(S\) containing
\(\Sigma_{a}\), then
\begin{equation}
\label{eq:ch4.aperture-containment}
\mathcal{S}_{\mathrm{ap}}(\mathcal{C})
=
\mathcal{S}_{\mathrm{real}}(\mathcal{C})
\quad\text{for patterns observed outside }S.
\end{equation}
If \(\Sigma_{a}\subsetneq S\) and PEC shorts part of \(S\), then
\(\mathcal{S}_{\mathrm{ap}}(\mathcal{C})\subsetneq\mathcal{S}_{\mathrm{real}}(\mathcal{C})\) in general
\cite{collin2001,harrington2001,stratton1941}.
\end{proposition}
\begin{proof}
Equality in Eq.~\eqref{eq:ch4.aperture-containment} follows from
Eq.~\eqref{eq:ch4.gamma-equivalence-invariance} when \(\mathbf{J}_{\mathrm{eq}}\) is supported only on
\(\Sigma_{a}\) after PEC reduction \cite{stratton1941}. Strict inclusion when \(\Sigma_{a}\subsetneq S\) follows
because interior currents behind PEC may radiate only through \(\Sigma_{a}\), restricting achievable \(\mathbf{c}\)
\cite{collin2001,hansen1988}.
\end{proof}

Proposition~\ref{prop:ch4.aperture-realizability-reduction} explains why horn, slot, and patch designs optimize
\(\mathbf{J}_{s}\) on \(\Sigma_{a}\) rather than full volume \(\Jimp\): when containment holds, the aperture
manifold is the entire synthesis set \cite{harrington2001,collin2001}. When containment fails, phase-only
aperture engineering optimizes over \(\mathcal{S}_{\mathrm{ap}}(\mathcal{C})\), not over
\(\mathcal{S}_{\mathrm{real}}(\mathcal{C})\) \cite{devaney2004}.

\paragraph{Field--source duality on \(\Sigma_{a}\).}
Given measured tangential \(\E_{t},\H_{t}\) on \(\Sigma_{a}\), equivalence currents follow
Eq.~\eqref{eq:ch4.stratton-chu-currents} with \(\hat{\mathbf{n}}\) the aperture normal
\cite{stratton1941,harrington2001}. The realizability question becomes whether those traces arise from a
\(\mathcal{D}_{\mathrm{real}}=1\) tuple in the full structure, not whether they can be fitted on \(\Sigma_{a}\)
in isolation \cite{devaney2004}. Near-field reactive content on \(\Sigma_{a}\) affects \(\mathbf{J}_{s}\) but not
\(\mathfrak{F}\) once \(\Pi_{\mathrm{rad}}\) has acted \cite{stratton1941}.

\paragraph{Worked illustration (open-ended waveguide).}
A rectangular waveguide radiates through aperture \(\Sigma_{a}\) at the open end. TE\(_{10}\) modal \(\mathbf{J}_{s}\)
on \(\Sigma_{a}\) generates \(\mathbf{c}\) with primarily \(\ell=1\) content at \(k_{0}a\ll1\) for small
\(\Sigma_{a}\) \cite{collin2001,jackson1999}. Targeting \(\ell>3\) coefficients from \(\Sigma_{a}\) alone fails
Corollary~\ref{cor:ch4.channel-count} when patch rank \(P\) is insufficient \cite{hansen1988,harrington2001}.

\paragraph{Problem (aperture versus full-volume synthesis).}
\textbf{Problem.} A design team fits \(\mathbf{c}_{\mathrm{target}}\) using volume \(\Jimp\) in a full FEM model,
then replaces \(\Jimp\) by \(\mathbf{J}_{s}\) on \(\Sigma_{a}\) via Eq.~\eqref{eq:ch4.stratton-chu-currents}. Must
\(\mathbf{c}\) match?

\textbf{Step 1.} Verify Theorem~\ref{thm:ch4.volume-surface-equivalence} hypotheses: \(S\) must enclose all
\(\operatorname{supp}\Jimp\) and \(\Omega_{\mathrm{ext}}\) must share the same outgoing closure
\cite{stratton1941}.

\textbf{Step 2.} Check \(\mathcal{D}_{\mathrm{cpl}}=1\) on the MoM mesh of \(\Sigma_{a}\)
\cite{harrington2001}.

\textbf{Step 3.} If \(\Sigma_{a}\) does not satisfy Eq.~\eqref{eq:ch4.aperture-containment}, expect
\(\mathbf{c}_{\mathrm{aperture}}\neq\mathbf{c}_{\mathrm{volume}}\) even when both are formally admissible
\cite{collin2001,devaney2004}.

\textbf{Physical meaning.} Aperture reduction is a realizability-preserving simplification only under containment
\cite{harrington2001}.

\paragraph{Horn, slot, and patch taxonomy.}
Laboratory apertures fall into three realizability classes distinguished by how \(\mathcal{S}_{\mathrm{ap}}(\mathcal{C})\)
sits inside \(\mathcal{S}_{\mathrm{real}}(\mathcal{C})\) \cite{collin2001,harrington2001}. \emph{Primary-radiating}
apertures (open-ended guides, horns) satisfy Eq.~\eqref{eq:ch4.aperture-containment} with \(\mathcal{S}_{\mathrm{ap}}\)
equal to the exterior-observed synthesis set \cite{stratton1941}. \emph{Secondary-radiating} apertures (slots in PEC,
coupled apertures) excite \(\mathcal{S}_{\mathrm{ap}}(\mathcal{C})\subsetneq\mathcal{S}_{\mathrm{real}}(\mathcal{C})\) when
interior resonances synthesize patterns not reachable by \(\Sigma_{a}\) alone \cite{harrington2001,devaney2004}.
\emph{Scattering} apertures carry induced \(\mathbf{J}_{s}\) and require calibration separation
Eq.~\eqref{eq:ch4.scattering-aperture-decomposition} before \(\mathbf{c}\) is interpreted as a synthesis target
\cite{stratton1941}.

\begin{equation}
\label{eq:ch4.aperture-class-inclusion}
\mathcal{S}_{\mathrm{ap}}^{\mathrm{(scat)}}(\mathcal{C})
\subseteq
\mathcal{S}_{\mathrm{ap}}^{\mathrm{(sec)}}(\mathcal{C})
\subseteq
\mathcal{S}_{\mathrm{ap}}^{\mathrm{(pri)}}(\mathcal{C})
\subseteq
\mathcal{S}_{\mathrm{real}}(\mathcal{C}),
\end{equation}
with inclusions strict in general \cite{collin2001,hansen1988}. Equation~\eqref{eq:ch4.aperture-class-inclusion} orders
design freedom: phase-only optimization on a slot optimizes \(\mathcal{S}_{\mathrm{ap}}^{\mathrm{(sec)}}\), not arbitrary
\(\mathcal{S}_{\mathrm{real}}(\mathcal{C})\) \cite{devaney2004}.

\begin{figure}[t]
\centering
\ChOneFigBlock{%
\begin{tikzpicture}[ch1fig, node distance=7mm and 8mm]
  \node[ch1box, minimum width=28mm] (pri) {Primary\\horn/guide};
  \node[ch1box, minimum width=28mm, right=10mm of pri] (sec) {Secondary\\slot/coupled};
  \node[ch1box, minimum width=28mm, right=10mm of sec] (scat) {Scattering\\aperture};
  \node[ch1box, minimum width=32mm, below=13mm of sec] (sr) {$\mathcal{S}_{\mathrm{real}}(\mathcal{C})$};
  \draw[ch1flow] (pri) -- (sr);
  \draw[ch1flow] (sec) -- (sr);
  \draw[ch1flow] (scat) -- (sr);
\end{tikzpicture}%
}
\caption{Subsection~\ref{sec:ch422}: aperture class determines which subset of
\(\mathcal{S}_{\mathrm{real}}(\mathcal{C})\) is reachable by \(\Sigma_{a}\) alone.}
\label{fig:ch4.aperture-taxonomy}
\end{figure}

\paragraph{Effective aperture diameter.}
Define \(D_{\mathrm{eff}}(\Sigma_{a})\) as the diameter of the smallest ball enclosing \(\Sigma_{a}\) with appreciable
\(\|\mathbf{J}_{s}\|\) \cite{collin2001,stratton1941}. Spectral compression uses \(D_{\mathrm{eff}}\) in place of
\(D_{s}\) when equivalence reduction is valid:
\begin{equation}
\label{eq:ch4.effective-aperture-size}
\Xi_{\mathrm{eff}}
=
\frac{k_{0}D_{\mathrm{eff}}(\Sigma_{a})}{2\pi},
\qquad
\mathcal{S}_{\mathrm{ap}}(\mathcal{C})
\ \text{compressed by}\ 
\Xi_{\mathrm{eff}}
\ \text{as in}\ 
\text{Eq.~\eqref{eq:ch4.ellmax-from-tail}}.
\end{equation}
Equation~\eqref{eq:ch4.effective-aperture-size} links aperture engineering to the same rank laws as volume supports
\cite{hansen1988,jackson1999,harrington2001}.

\paragraph{Worked illustration (coupled slot).}
Two slots in a PEC cylinder may synthesize \(\mathbf{c}\) not realizable by either slot alone: interior coupling
enlarges \(\mathcal{S}_{\mathrm{real}}(\mathcal{C})\) beyond \(\mathcal{S}_{\mathrm{ap}}(\mathcal{C})\) for a single
\(\Sigma_{a}\) \cite{collin2001,harrington2001}. Equivalence on one aperture underfits the coupled synthesis set
\cite{devaney2004}.

\paragraph{Problem (horn mouth versus feed).}
\textbf{Problem.} A horn's mouth \(\Sigma_{a}\) is used to synthesize \(\mathbf{c}\), but the feed remains inside the
horn. Is \(\mathcal{S}_{\mathrm{ap}}=\mathcal{S}_{\mathrm{real}}\)?

\textbf{Step 1.} Verify enclosure: all \(\operatorname{supp}\Jimp\) must lie inside \(S\) enclosing the feed
\cite{stratton1941,collin2001}.

\textbf{Step 2.} Check Eq.~\eqref{eq:ch4.aperture-containment} for the declared exterior observation domain
\cite{harrington2001}.

\textbf{Step 3.} If interior standing waves alter tangential traces on \(\Sigma_{a}\) without changing exterior
\(\mathbf{c}\), feed redesign is invisible until traces change \cite{devaney2004}.

\textbf{Physical meaning.} Horn mouths are primary apertures only under full exterior equivalence \cite{collin2001}.

\paragraph{Aperture sampling and coupling Nyquist rate.}
On \(\Sigma_{a}\) with area \(A_{a}\), patch spacing \(\Delta\) sets a discrete coupling Nyquist wavenumber
\begin{equation}
\label{eq:ch4.aperture-nyquist-rate}
k_{t,\mathrm{Nyq}}
=
\frac{\pi}{\Delta},
\qquad
\Delta\lesssim\frac{\lambda}{2}
\ \Longrightarrow\
k_{t,\mathrm{Nyq}}\gtrsim k_{0},
\end{equation}
importing standard aperture sampling limits \cite{collin2001,harrington2001,jackson1999}. MoM bases that violate
Eq.~\eqref{eq:ch4.aperture-nyquist-rate} produce aliased \(\mathbf{K}_{\mathrm{disc}}\) rows: formally nonzero couplings
at high \(\ell\) that are mesh artifacts, not export modes \cite{devaney2004,stratton1941}.

\begin{proposition}[Aperture MoM rank versus continuum rank]
\label{prop:ch4.aperture-mom-rank}
For stable RWG discretization on \(\Sigma_{a}\) with \(P\) patches,
\(\operatorname{rank}\mathbf{T}_{\mathrm{VSH}}\le\min(P,|I_{\mathrm{acc}}|)\) and
\(\operatorname{rank}\mathbf{T}_{\mathrm{VSH}}\le C\,k_{0}^{2}A_{a}\) asymptotically per
Eq.~\eqref{eq:ch4.aperture-coupling-scaling} \cite{harrington2001,collin2001,devaney2004}. Excess patches beyond the
Nyquist rate do not increase \(\operatorname{ran}\boldsymbol{\Gamma}_\omega\) in the continuum limit
\cite{stratton1941,hansen1988}.
\end{proposition}

\paragraph{Worked illustration (oversampled slot).}
A slot discretized with \(P=400\) sub-wavelength patches at \(k_{0}D_{s}=0.5\) yields a numerically full-rank
\(\mathbf{T}_{\mathrm{VSH}}\) while \(r_{\mathrm{rad}}(10^{-2})\le3\) \cite{harrington2001,hansen1988}. Synthesis audits
must use \(r_{\mathrm{rad}}\), not \(\operatorname{rank}\mathbf{T}_{\mathrm{VSH}}\), as the DOF budget
\cite{devaney2004,stratton1941}.

\paragraph{Problem (patch density versus synthesis rank).}
\textbf{Problem.} Doubling patch count on fixed \(\Sigma_{a}\) doubles MoM rank but not measured far-zone diversity.
Explain.

\textbf{Step 1.} Continuum rank is capped by Eq.~\eqref{eq:ch4.aperture-coupling-scaling} at fixed \((A_{a},k_{0})\)
\cite{collin2001,jackson1999}.

\textbf{Step 2.} Additional patches refine \(\mathbf{J}_{s}\) representation without enlarging
\(\operatorname{ran}\boldsymbol{\Gamma}_\omega\) \cite{harrington2001,devaney2004}.

\textbf{Step 3.} Compare SVD of \(\mathbf{T}_{\mathrm{VSH}}\) to \(r_{\mathrm{rad}}(\varepsilon)\); excess singular
directions lie below \(\varepsilon\sigma_{1}\) \cite{hansen1988,stratton1941}.

\textbf{Physical meaning.} Patch density is discretization resolution, not geometry-imposed synthesis dimension
\cite{collin2001}.

\paragraph{Validity.}
Definitions~\ref{def:ch4.aperture-equivalent-source} and Proposition~\ref{prop:ch4.aperture-realizability-reduction}
assume linear Maxwell theory on \(\Omega_{\mathrm{ext}}\) \cite{stratton1941}. Magnetic currents \(\mathbf{M}_{s}\) on
aperture edges may be required when \(\E_{t}\neq\mathbf{0}\) on \(\Sigma_{a}\) \cite{harrington2001,collin2001}.

Tangential continuity constraints on \(\Sigma_{a}\) are imported in Subsection~\ref{sec:ch423}.


\subsection{Tangential continuity at boundaries}
\label{sec:ch423}
% TOC tag: [USE]

Equivalent sources are not free currents: they must reproduce the tangential field structure enforced by
Maxwell transmission on interfaces \cite{stratton1941}. Section~\ref{sec:ch01.1.2} fixed that structure in
Eq.~\eqref{eq:ch1.interface-conditions}; Subsection~\ref{sec:ch423} records the realizability consequences for
aperture and equivalence-surface models \cite{harrington2001,collin2001}.

\paragraph{Assumptions.}
Let \(\Sigma_{ab}\) be a material interface with unit normal \(\hat{\mathbf{n}}_{ab}\) as in
Section~\ref{sec:ch01.1.2} \cite{stratton1941}. Aperture \(\Sigma_{a}\) may coincide with a portion of
\(\Sigma_{ab}\) or with a PEC opening where tangential \(\E=\mathbf{0}\) \cite{collin2001}. Equivalent currents
\(\mathbf{J}_{s}\) are defined in the distributional sense of Eq.~\eqref{eq:ch4.surface-source-embedding}
\cite{stratton1941,harrington2001}.

\paragraph{Tangential \(\E\) continuity.}
Across a dielectric interface without electric sheet charge,
Eq.~\eqref{eq:ch1.interface-conditions} requires
\begin{equation}
\label{eq:ch4.tangential-E-continuity}
\hat{\mathbf{n}}_{ab}\times\big(\E_{b}-\E_{a}\big)=\mathbf{0}
\quad\text{on}\ \Sigma_{ab},
\end{equation}
importing the first row of Eq.~\eqref{eq:ch1.interface-conditions} \cite{stratton1941}. For a PEC surface,
\(\hat{\mathbf{n}}\times\E=\mathbf{0}\) replaces Eq.~\eqref{eq:ch4.tangential-E-continuity} on the conductor
\cite{stratton1941,collin2001}. Realizable aperture fields must satisfy
Eq.~\eqref{eq:ch4.tangential-E-continuity} on dielectric portions and the PEC condition elsewhere; otherwise
\(\mathbf{J}_{s}=\hat{\mathbf{n}}\times\H\) from Eq.~\eqref{eq:ch4.stratton-chu-currents} is inconsistent with
the declared boundary class \cite{harrington2001}.

\paragraph{Tangential \(\H\) jump and surface current.}
The second row of Eq.~\eqref{eq:ch1.interface-conditions} gives
\begin{equation}
\label{eq:ch4.tangential-H-jump}
\hat{\mathbf{n}}_{ab}\times\big(\H_{b}-\H_{a}\big)=\mathbf{J}_{s,ab}
\quad\text{on}\ \Sigma_{ab},
\end{equation}
linking impressed surface current to tangential magnetic jump \cite{stratton1941}. Equation
\eqref{eq:ch4.tangential-H-jump} is the realizability closure for \(\mathbf{J}_{s}\): an equivalent aperture
current must equal the magnetic-field jump it claims to represent \cite{harrington2001,collin2001}. On PEC with
aperture \(\Sigma_{a}\), \(\H_{t}\) is continuous in the sense of exterior traces while \(\mathbf{J}_{s}\) carries
the feed \cite{stratton1941}.

\begin{proposition}[Tangential realizability screen]
\label{prop:ch4.tangential-realizability-screen}
Define
\begin{equation}
\label{eq:ch4.tangential-screen}
\mathcal{D}_{\tau}[\mathbf{J}_{s},\E,\H]
=
1
\iff
\text{Eqs.~\eqref{eq:ch4.tangential-E-continuity}--\eqref{eq:ch4.tangential-H-jump} hold on }\Sigma_{a}
\ \wedge\
\mathbf{J}_{s}=\hat{\mathbf{n}}\times\H_{t}\ \text{on radiating portions}.
\end{equation}
Then \(\mathbf{J}_{s}\in\mathcal{J}_{\mathrm{real}}\) on aperture channels only if
\(\mathcal{D}_{\tau}=1\) and \(\mathcal{D}_{\mathrm{cpl}}=1\) \cite{stratton1941,harrington2001}.
\end{proposition}
\begin{proof}
Violation of Eq.~\eqref{eq:ch4.tangential-E-continuity} contradicts
Eq.~\eqref{eq:ch1.interface-conditions} and therefore Maxwell admissibility
\cite{stratton1941}. Mismatch between \(\mathbf{J}_{s}\) and
Eq.~\eqref{eq:ch4.tangential-H-jump} implies the surface source does not generate the declared tangential
\(\H\) \cite{harrington2001}. Coupling admissibility is independent and is presupposed from
Eq.~\eqref{eq:ch4.coupling-admissibility} \cite{harrington2001}.
\end{proof}

Proposition~\ref{prop:ch4.tangential-realizability-screen} prevents a common failure mode: prescribing a
tangential \(\E\) template on \(\Sigma_{a}\) that is not the trace of any outgoing \(\mathfrak{F}\in\mathcal{F}_{\mathrm{rad}}\)
\cite{stratton1941,devaney2004}. Pattern optimization over unconstrained \(\E_{t}\) on \(\Sigma_{a}\) optimizes
outside \(\mathcal{S}_{\mathrm{ap}}(\mathcal{C})\) \cite{hansen1988}.

\paragraph{Trace-space form.}
Tangential traces live in \(H^{-1/2}(\Sigma_{a})^{3}\) as in Eq.~\eqref{eq:ch2.trace-space-definition}
\cite{stratton1941}. Realizable \(\mathbf{J}_{s}\) are dual to tangential \(\E_{t}\) under the pairing of
Subsection~\ref{sec:ch232} \cite{harrington2001}. Numerically, \(\mathcal{D}_{\tau}=1\) is enforced by matching
continuity rows in the MoM impedance matrix before \(\mathbf{T}\) of Eq.~\eqref{eq:ch4.mom-synthesis-matrix} is
inverted \cite{harrington2001,collin2001}.

\paragraph{Worked illustration (slot in PEC plane).}
On PEC, \(\hat{\mathbf{n}}\times\E=\mathbf{0}\) everywhere; on the slot \(\Sigma_{a}\), \(\E_{t}\neq\mathbf{0}\) and
\(\mathbf{J}_{s}=\hat{\mathbf{n}}\times\H\) carries radiation \cite{stratton1941,collin2001}. A slot pattern that
violates \(\mathcal{D}_{\tau}=1\) by discontinuity of \(\E_{t}\) across the slot rim is not realizable by any
\(\mathbf{J}_{s}\in\mathcal{J}_{\mathrm{real}}\) \cite{harrington2001}.

\paragraph{PEC--aperture jump template.}
On a PEC screen with slot \(\Sigma_{a}\), the tangential field pair satisfies
\begin{equation}
\label{eq:ch4.pec-slot-template}
\hat{\mathbf{n}}\times\E
=
\mathbf{0}
\quad\text{on PEC},
\qquad
\mathbf{J}_{s}
=
\hat{\mathbf{n}}\times\H
\quad\text{on }\Sigma_{a},
\qquad
\mathbf{J}_{s}=\mathbf{0}
\quad\text{on PEC},
\end{equation}
importing the PEC specialization of Eqs.~\eqref{eq:ch4.tangential-E-continuity}--\eqref{eq:ch4.tangential-H-jump}
\cite{stratton1941,collin2001}. Equation~\eqref{eq:ch4.pec-slot-template} is the admissible manifold for
\(\mathcal{S}_{\mathrm{ap}}(\mathcal{C})\) on a flat aperture: any \(\E_{t}\) template on \(\Sigma_{a}\) that cannot be
embedded in an outgoing solution with Eq.~\eqref{eq:ch4.pec-slot-template} on the surrounding PEC is outside
\(\mathcal{S}_{\mathrm{real}}(\mathcal{C})\) \cite{harrington2001,devaney2004}.

\begin{figure}[t]
\centering
\ChOneFigBlock{%
\begin{tikzpicture}[ch1fig, node distance=7mm and 8mm]
  \node[ch1box, minimum width=30mm] (pec) {PEC screen\\$\hat{\mathbf{n}}\times\E=\mathbf{0}$};
  \node[ch1box, minimum width=26mm, right=12mm of pec] (slot) {Slot $\Sigma_{a}$\\$\E_{t}\neq\mathbf{0}$};
  \node[ch1box, minimum width=28mm, below=12mm of slot] (js) {$\mathbf{J}_{s}=\hat{\mathbf{n}}\times\H$};
  \node[ch1box, minimum width=28mm, left=12mm of js] (dtau) {$\mathcal{D}_{\tau}=1$};
  \draw[ch1flow] (pec) -- (slot);
  \draw[ch1flow] (slot) -- (js);
  \draw[ch1flow] (js) -- (dtau);
\end{tikzpicture}%
}
\caption{Subsection~\ref{sec:ch423}: tangential realizability on a PEC screen requires
Eq.~\eqref{eq:ch4.pec-slot-template}; \(\mathcal{D}_{\tau}=1\) is the joint screen on \(\E_{t}\) and
\(\mathbf{J}_{s}\).}
\label{fig:ch4.pec-slot-tangential}
\end{figure}

Figure~\ref{fig:ch4.pec-slot-tangential} connects boundary-class data to synthesis: \(\mathbf{T}\) of
Eq.~\eqref{eq:ch4.mom-synthesis-matrix} must be assembled only after \(\mathcal{D}_{\tau}=1\) is enforced on the
declared trace class \cite{harrington2001,collin2001}. Inverse problems that fit \(\E_{t}\) on \(\Sigma_{a}\) without
the PEC constraint inject non-physical tangential content and produce formal \(\mathbf{c}\) outside
\(\mathcal{S}_{\mathrm{ap}}(\mathcal{C})\) \cite{devaney2004}.

\paragraph{Dielectric interface with magnetic sheet.}
When \(\Sigma_{a}\) lies on a dielectric interface, Eq.~\eqref{eq:ch4.tangential-E-continuity} requires continuous
tangential \(\E\), while Eq.~\eqref{eq:ch4.tangential-H-jump} permits a surface current \(\mathbf{J}_{s,ab}\) only when
magnetic jump is physical \cite{stratton1941}. If \(\E_{t}\) is continuous but the declared \(\mathbf{J}_{s}\) does not
equal \(\hat{\mathbf{n}}_{ab}\times(\H_{b}-\H_{a})\), \(\mathcal{D}_{\tau}=0\) and aperture synthesis is undefined
\cite{harrington2001}. Magnetic currents \(\mathbf{M}_{s}\) enter on interfaces where \(\E_{t}\) is discontinuous in
the equivalent-source picture of Section~\ref{sec:ch42} \cite{collin2001,stratton1941}.

\begin{corollary}[Tangential screen and synthesis range]
\label{cor:ch4.tangential-range}
Let \(\mathcal{J}_{\mathrm{ap}}=\{\mathbf{J}_{s}:\mathcal{D}_{\tau}=1,\ \mathcal{D}_{\mathrm{cpl}}=1\}\). Then
\(\mathcal{S}_{\mathrm{ap}}(\mathcal{C})=\boldsymbol{\Gamma}_\omega[\mathcal{J}_{\mathrm{ap}}]\subseteq
\mathcal{S}_{\mathrm{real}}(\mathcal{C})\), with equality when every realizable \(\mathbf{c}\) admits an aperture
representative in \(\mathcal{J}_{\mathrm{ap}}\) \cite{stratton1941,harrington2001,collin2001}.
\end{corollary}

\paragraph{Problem (rim discontinuity).}
\textbf{Problem.} A design specifies \(\E_{t}\) on a rectangular slot with a step discontinuity at the rim between slot
and PEC. Is \(\mathbf{c}_{\mathrm{target}}\) synthesizable?

\textbf{Step 1.} Test Eq.~\eqref{eq:ch4.pec-slot-template}: on PEC, \(\E_{t}\) must vanish; a rim step violates
\(\mathcal{D}_{\tau}=1\) unless a magnetic sheet models the rim \cite{stratton1941,collin2001}.

\textbf{Step 2.} Without a declared sheet, \(\mathbf{J}_{s}=\hat{\mathbf{n}}\times\H\) cannot represent the jump;
\(\boldsymbol{\Gamma}_\omega\) has no preimage in \(\mathcal{J}_{\mathrm{ap}}\) \cite{harrington2001}.

\textbf{Step 3.} Smooth the template or enlarge the model class to include rim \(\mathbf{M}_{s}\) with
\(\mathcal{D}_{\mathrm{cpl}}=1\) \cite{devaney2004}.

\textbf{Physical meaning.} Tangential screens are Maxwell predicates, not optional boundary penalties
\cite{stratton1941}.

\paragraph{Impedance and Leontovich boundaries.}
On impedance surfaces \(\mathbf{Z}_{s}=\mathbf{Z}_{s}(\mathbf{r})\), tangential fields satisfy
\begin{equation}
\label{eq:ch4.leontovich-screen}
\hat{\mathbf{n}}\times\big(\E-\mathbf{Z}_{s}\cdot\hat{\mathbf{n}}\times\H\big)=\mathbf{0}
\quad\text{on }\Sigma_{a},
\end{equation}
importing the impedance row of Eq.~\eqref{eq:ch1.interface-conditions} \cite{stratton1941,harrington2001}. Realizable
\(\mathbf{J}_{s}\) on \(\Sigma_{a}\) must be consistent with Eq.~\eqref{eq:ch4.leontovich-screen}: prescribing
\(\E_{t}\) incompatible with \(\H_{t}\) and \(\mathbf{Z}_{s}\) yields \(\mathcal{D}_{\tau}=0\)
\cite{collin2001,devaney2004}. PEC is the limit \(\mathbf{Z}_{s}\to\mathbf{0}\) with
Eq.~\eqref{eq:ch4.pec-slot-template} \cite{stratton1941}.

\paragraph{Transmission-line port coupling.}
Waveguide ports attach through \(\mathfrak{g}_\Gamma\) in Eq.~\eqref{eq:ch4.excitation-tuple}: port modes define a
low-dimensional subspace of \(\mathcal{J}_{\mathrm{real}}\) before equivalence reduction \cite{collin2001,harrington2001}.
Tangential matching at the port aperture is the discrete analog of Eqs.~\eqref{eq:ch4.tangential-E-continuity}--\eqref{eq:ch4.tangential-H-jump}
\cite{stratton1941}. Port synthesis rank is bounded by guided-mode count, not by full \(\mathbb{C}^{I}\)
\cite{hansen1988,devaney2004}.

\begin{proposition}[Tangential screen on impedance apertures]
\label{prop:ch4.impedance-tangential-screen}
On Leontovich \(\Sigma_{a}\), \(\mathcal{D}_{\tau}=1\) requires simultaneous satisfaction of
Eqs.~\eqref{eq:ch4.leontovich-screen}, \eqref{eq:ch4.tangential-H-jump}, and
\(\mathcal{D}_{\mathrm{cpl}}=1\) on the MoM mesh \cite{stratton1941,harrington2001,collin2001}. Violation at any row
excludes \(\mathbf{c}\) from \(\mathcal{S}_{\mathrm{ap}}(\mathcal{C})\) \cite{devaney2004}.
\end{proposition}

\paragraph{Worked illustration (slot with finite conductivity).}
A slot rim with finite conductivity modifies \(\mathbf{Z}_{s}\): \(\E_{t}\) and \(\H_{t}\) remain coupled through
Eq.~\eqref{eq:ch4.leontovich-screen} \cite{stratton1941,harrington2001}. Fitting \(\E_{t}\) alone on the slot without
\(\H_{t}\) consistency fails \(\mathcal{D}_{\tau}=1\) \cite{collin2001,devaney2004}.

\paragraph{Problem (impedance boundary fit).}
\textbf{Problem.} A design specifies \(\E_{t}\) on \(\Sigma_{a}\) with \(\mathbf{Z}_{s}=Z_{0}\hat{\mathbf{z}}\hat{\mathbf{z}}\) but leaves \(\H_{t}\) free. Is synthesis well-posed?

\textbf{Step 1.} Eq.~\eqref{eq:ch4.leontovich-screen} links \(\E_{t}\) and \(\H_{t}\) \cite{stratton1941}.

\textbf{Step 2.} Free \(\H_{t}\) with fixed \(\E_{t}\) overdetermines or contradicts the impedance relation unless
\(\mathbf{J}_{s}\) is recomputed \cite{harrington2001,collin2001}.

\textbf{Step 3.} Declare joint \((\E_{t},\H_{t})\) or solve for \(\mathbf{J}_{s}=\hat{\mathbf{n}}\times\H_{t}\) consistently
\cite{devaney2004}.

\textbf{Physical meaning.} Impedance boundaries are tangential predicates, not amplitude templates \cite{stratton1941}.

\paragraph{PEC--aperture duality screen.}
On PEC bodies with apertures \(\Sigma_{a}\), tangential \(\E_{t}=0\) on metal and finite \(\E_{t}\) on \(\Sigma_{a}\) impose
\begin{equation}
\label{eq:ch4.pec-aperture-duality}
\mathcal{D}_{\tau}^{\mathrm{(PEC)}}
=
1
\iff
\begin{cases}
\E_{t}=0 & \text{on PEC},\\
\mathcal{D}_{\mathrm{cpl}}=1,\ \mathcal{D}_{\mathrm{real}}=1 & \text{on }\Sigma_{a},
\end{cases}
\end{equation}
importing the PEC limit of Eq.~\eqref{eq:ch4.leontovich-screen} \cite{stratton1941,collin2001,harrington2001}. Magnetic
rim currents \(\mathbf{M}_{s}\) on aperture edges are required when \(\E_{t}\) jumps between PEC and aperture sectors
\cite{stratton1941,devaney2004}.

\begin{theorem}[Tangential realizability on composite boundaries]
\label{thm:ch4.composite-tangential-realizability}
Let \(S=S_{\mathrm{PEC}}\cup\Sigma_{a}\) be a composite boundary with declared \(\mathbf{Z}_{s}\) on \(\Sigma_{a}\). Then
\(\mathbf{c}\in\mathcal{S}_{\mathrm{ap}}(\mathcal{C})\) only if \(\mathcal{D}_{\tau}=1\) on every interface row of
Eqs.~\eqref{eq:ch4.tangential-E-continuity}--\eqref{eq:ch4.leontovich-screen} simultaneously
\cite{stratton1941,harrington2001,collin2001}. Failure on any row excludes \(\mathbf{c}\) before
\(\boldsymbol{\Gamma}_\omega\) rank is tested \cite{devaney2004}.
\end{theorem}
\begin{proof}
Well-posedness of equivalent-source reduction requires consistent tangential traces on all boundary components
\cite{stratton1941,collin2001}; Proposition~\ref{prop:ch4.impedance-tangential-screen} \cite{harrington2001}.
\end{proof}

\paragraph{Worked illustration (slot in PEC plane).}
A narrow slot requires \(\mathbf{J}_{s}\) on \(\Sigma_{a}\) and consistent \(\H_{t}\) jumps; prescribing slot voltage without
\(\H_{t}\) violates Eq.~\eqref{eq:ch4.pec-aperture-duality} \cite{collin2001,stratton1941,harrington2001}.

\paragraph{Problem (rim magnetic current omission).}
\textbf{Problem.} MoM uses electric currents only on \(\Sigma_{a}\) in a PEC screen. Measured \(\mathbf{c}\) disagrees with
volume reference. Is \(\mathcal{D}_{\tau}=1\)?

\textbf{Step 1.} Test rim \(\E_{t}\) continuity across PEC edges \cite{stratton1941}.

\textbf{Step 2.} If rim jumps are physical, add \(\mathbf{M}_{s}\) with \(\mathcal{D}_{\mathrm{cpl}}=1\) \cite{harrington2001,collin2001}.

\textbf{Step 3.} Recompute \(\mathbf{c}\) on the enlarged composite boundary \cite{devaney2004}.

\textbf{Physical meaning.} Composite boundaries require composite tangential data, not aperture-only templates \cite{stratton1941}.

\paragraph{Validity.}
Eqs.~\eqref{eq:ch4.tangential-E-continuity}--\eqref{eq:ch4.tangential-screen} assume piecewise smooth interfaces
\cite{stratton1941}. Magnetic sheets and anisotropic \(\boldsymbol{\xi}\) modify jump rows without removing the
tangential screen \cite{harrington2001}. Section~\ref{sec:ch01.1.2} is not re-derived here.

Failure regimes of equivalence are catalogued in Subsection~\ref{sec:ch424}.


\subsection{Limits of equivalent-source models}
\label{sec:ch424}
% TOC tag: [HOME]

Volume--surface equivalence collapses synthesis degrees of freedom only within declared validity; outside that
domain, \(\boldsymbol{\Gamma}_\omega[\mathbf{J}_{\mathrm{eq}}]\) may differ from
\(\boldsymbol{\Gamma}_\omega[\Jimp]\) even when both tuples are locally admissible \cite{stratton1941,harrington2001,devaney2004}.
Subsection~\ref{sec:ch424} states the limits as realizability theorems, not as numerical caveats alone.

\paragraph{Assumptions.}
Fix equivalence surface \(S\) or aperture \(\Sigma_{a}\), outgoing \(\Omega_{\mathrm{ext}}\), and
\(\mathcal{C}\in\mathcal{R}_{\mathrm{adm}}\) \cite{stratton1941,hansen1988}. Interior region
\(\Omega_{\mathrm{int}}\) may contain loss, dispersion, or multiple scattering bodies unless otherwise stated
\cite{harrington2001,collin2001}.

\paragraph{Region of validity.}
Equivalence preserves \(\mathfrak{F}\) only in the exterior observation domain declared in
Eq.~\eqref{eq:ch4.exterior-equivalence-condition} \cite{stratton1941}. Interior fields
\((\E,\H)\big|_{\Omega_{\mathrm{int}}}\) are generally \emph{not} invariant under replacement of \(\Jimp\) by
\((\mathbf{J}_{s},\mathbf{M}_{s})\) \cite{harrington2001,collin2001}. Near-field probes inside
\(\Omega_{\mathrm{int}}\) therefore test a different realizability class than far-zone \(\mathbf{c}\)
\cite{devaney2004,stratton1941}.

\begin{theorem}[Equivalence error on synthesis coefficients]
\label{thm:ch4.equivalence-error}
Suppose \((\mathbf{J}_{s},\mathbf{M}_{s})\) on \(S\) are constructed from approximate tangential traces
\(\widetilde{\E}_{t},\widetilde{\H}_{t}\) rather than exact values. Then
\begin{equation}
\label{eq:ch4.equivalence-coefficient-error}
\big\|\boldsymbol{\Gamma}_\omega[\mathbf{J}_{\mathrm{eq}}]-\mathbf{c}_{\mathrm{true}}\big\|_{\mathcal{C}}
\le
\mathcal{K}_{\mathrm{eq}}(S,\omega)\,
\big\|\widetilde{\mathbf{u}}_{t}-\mathbf{u}_{t}\big\|_{\mathcal{T}_\Gamma},
\end{equation}
with \(\mathbf{u}_{t}=(\E_{t},\H_{t})\), \(\mathcal{K}_{\mathrm{eq}}\) the operator norm of the equivalence map
on tangential traces, and \(\|\cdot\|_{\mathcal{C}}\) the coefficient norm induced by \(\mathcal{C}\)
\cite{stratton1941,harrington2001,devaney2004}.
\end{theorem}
\begin{proof}
Linearity of \(\mathcal{S}_\omega\), \(\mathcal{R}_\omega\), and \(\mathcal{C}\) gives
\(\boldsymbol{\Gamma}_\omega[\mathbf{J}_{\mathrm{eq}}-\mathbf{J}_{\mathrm{eq}}^{\mathrm{(true)}}]\) proportional to
trace perturbations on \(S\) \cite{stratton1941}. Continuity of tangential trace maps
Eqs.~\eqref{eq:ch2.tangential-trace-maps}--\eqref{eq:ch2.tangential-trace-magnetic} bounds the coefficient
perturbation by \(\|\widetilde{\mathbf{u}}_{t}-\mathbf{u}_{t}\|_{\mathcal{T}_\Gamma}\) \cite{stratton1941,harrington2001}.
\end{proof}

Theorem~\ref{thm:ch4.equivalence-error} quantifies equivalence breakdown: approximate aperture fields produce
approximate \(\mathbf{c}\) with error controlled by trace mismatch, not by wishful normalization
\cite{devaney2004,hansen1988}. \(\mathcal{K}_{\mathrm{eq}}\) grows with electrical size of \(S\) and with
proximity of \(S\) to \(\Omega_{s}\) \cite{harrington2001,stratton1941}.

\paragraph{Non-uniqueness and null contributions.}
Equivalent sources are not unique: any \(\delta\mathbf{J}_{s}\in\ker\mathcal{R}_\omega\) on \(S\) leaves
\(\mathbf{c}\) unchanged while altering near fields \cite{stratton1941,harrington2001}. Combined with
Eq.~\eqref{eq:ch4.synthesis-null-space},
\begin{equation}
\label{eq:ch4.equivalence-null-ambiguity}
\mathbf{J}_{s}^{(1)}-\mathbf{J}_{s}^{(2)}\in\mathcal{N}_{\mathrm{syn}}
\quad\Longrightarrow\quad
\boldsymbol{\Gamma}_\omega[\mathbf{J}_{s}^{(1)}]=\boldsymbol{\Gamma}_\omega[\mathbf{J}_{s}^{(2)}],
\end{equation}
even when \(\mathbf{J}_{s}^{(1)}\neq\mathbf{J}_{s}^{(2)}\) on \(\Sigma_{a}\) \cite{devaney2004}. Equivalence
reduction therefore enlarges the synthesis null space without shrinking \(\mathcal{S}_{\mathrm{real}}(\mathcal{C})\)
\cite{harrington2001}.

\paragraph{Failure taxonomy.}
Equivalence fails to preserve \(\mathbf{c}\) when \(S\) does not enclose all
\(\operatorname{supp}\Jimp\) \cite{stratton1941}; when interior loss or dispersion is mis-modeled in the
equivalent exterior problem \cite{harrington2001}; when \(\Sigma_{a}\subsetneq S\) and containment
Eq.~\eqref{eq:ch4.aperture-containment} fails \cite{collin2001}; when \(\mathcal{D}_{\tau}=0\) or
\(\mathcal{D}_{\mathrm{cpl}}=0\) on the discretized aperture \cite{harrington2001}; or when reactive
near-zone content is mistaken for radiative \(\mathbf{c}\) \cite{stratton1941,devaney2004}. Each failure
produces a distinct signature in Lemma~\ref{lem:ch4.screen-failures} or in
Theorem~\ref{thm:ch4.equivalence-error} \cite{harrington2001}.

\begin{corollary}[Interior insensitivity of \(\mathbf{c}\)]
\label{cor:ch4.interior-insensitivity}
If Theorem~\ref{thm:ch4.volume-surface-equivalence} holds, modifications of \(\Jimp\) inside
\(\Omega_{\mathrm{int}}\) that leave tangential traces on \(S\) unchanged do not alter
\(\boldsymbol{\Gamma}_\omega[\Jimp]\) \cite{stratton1941,collin2001}. Conversely, interior changes that
\emph{do} alter tangential traces on \(S\) change \(\mathbf{c}\) even when exterior equivalence was previously
valid \cite{harrington2001,devaney2004}.
\end{corollary}

Corollary~\ref{cor:ch4.interior-insensitivity} is the engineering reading of equivalence limits: feed redesign
inside a cavity is invisible to far-zone \(\mathbf{c}\) only until it changes aperture tangential fields
\cite{collin2001,harrington2001}. Phase-only metasurface design on \(\Sigma_{a}\) acts on
\(\mathcal{S}_{\mathrm{ap}}(\mathcal{C})\), not on arbitrary \(\mathcal{S}_{\mathrm{real}}(\mathcal{C})\)
\cite{devaney2004,hansen1988}.

\paragraph{Worked illustration (lossy interior).}
A radome with loss \(\varepsilon''>0\) inside \(\Omega_{\mathrm{int}}\) alters tangential traces on \(S\) relative
to a lossless model \cite{stratton1941,harrington2001}. Equivalent \(\mathbf{J}_{s}\) fit from exterior data
alone omit the loss mechanism; \(\mathbf{c}\) errors scale with \(\mathcal{K}_{\mathrm{eq}}\) in
Theorem~\ref{thm:ch4.equivalence-error} \cite{devaney2004}.

\paragraph{Problem (equivalence validity test).}
\textbf{Problem.} MoM on \(\Sigma_{a}\) yields \(\mathbf{c}_{\mathrm{MoM}}\); full-volume FEM yields
\(\mathbf{c}_{\mathrm{FEM}}\). When must they agree?

\textbf{Step 1.} Verify enclosure and outgoing closure for both models \cite{stratton1941}.

\textbf{Step 2.} Test \(\mathcal{D}_{\tau}\) and \(\mathcal{D}_{\mathrm{cpl}}\) on \(\Sigma_{a}\)
\cite{harrington2001}.

\textbf{Step 3.} Bound \(\|\mathbf{c}_{\mathrm{MoM}}-\mathbf{c}_{\mathrm{FEM}}\|\) by
Eq.~\eqref{eq:ch4.equivalence-coefficient-error} using measured trace mismatch on \(S\)
\cite{devaney2004}.

\textbf{Physical meaning.} Disagreement is an equivalence failure, not necessarily a coding bug
\cite{harrington2001}.

\paragraph{Huygens versus Love equivalence regimes.}
Stratton--Chu equivalence Eq.~\eqref{eq:ch4.stratton-chu-currents} requires both \(\mathbf{J}_{s}\) and
\(\mathbf{M}_{s}\) on \(S\) \cite{stratton1941}. Love equivalence eliminates \(\mathbf{M}_{s}\) on PEC portions but
restricts \(\mathcal{J}_{\mathrm{eq}}\) to electric currents on apertures \cite{harrington2001,collin2001}. The
realizability consequence is rank: Love-reduced models shrink the discrete synthesis basis on \(S\) without shrinking
\(\operatorname{ran}\boldsymbol{\Gamma}_\omega\) when Theorem~\ref{thm:ch4.volume-surface-equivalence} holds
\cite{devaney2004,stratton1941}.

\begin{lemma}[Equivalence surface proximity]
\label{lem:ch4.equivalence-proximity}
\(\mathcal{K}_{\mathrm{eq}}(S,\omega)\) in Theorem~\ref{thm:ch4.equivalence-error} increases as \(S\) approaches
\(\operatorname{supp}\Jimp\) \cite{stratton1941,harrington2001}. Equivalence surfaces too close to \(\Omega_{s}\) amplify
trace errors in \(\mathbf{c}\) even when \(\mathcal{D}_{\mathrm{cpl}}=1\) on the mesh \cite{devaney2004,collin2001}.
\end{lemma}

\paragraph{Reactive near-field mismatch.}
Equivalence preserves \(\mathfrak{F}\) on \(\mathcal{F}_{\mathrm{rad}}\) but not reactive fields in
\(\Omega_{\mathrm{int}}\) \cite{stratton1941}. Near-field probes that reconstruct \(\mathbf{J}_{\mathrm{eq}}\) from
interior samples therefore test a different realizability class than far-zone \(\mathbf{c}\)
\cite{harrington2001,devaney2004}. The split mirrors Eq.~\eqref{eq:ch4.near-far-coupling-split} at the equivalence layer.

\begin{corollary}[Far-zone insensitivity to interior equivalence choice]
\label{cor:ch4.far-insensitivity-equivalence}
Distinct equivalent pairs \((\mathbf{J}_{s}^{(A)},\mathbf{M}_{s}^{(A)})\) and
\((\mathbf{J}_{s}^{(B)},\mathbf{M}_{s}^{(B)})\) on \(S\) that satisfy
Theorem~\ref{thm:ch4.volume-surface-equivalence} yield identical \(\mathbf{c}\) while differing in near-zone
\(\H,\E\) on \(\Omega_{\mathrm{int}}\) \cite{stratton1941,collin2001,harrington2001}.
\end{corollary}

\paragraph{Worked illustration (Huygens on dielectric surface).}
On a dielectric interface, omitting \(\mathbf{M}_{s}\) while fitting \(\mathbf{J}_{s}\) alone violates
Eq.~\eqref{eq:ch4.stratton-chu-currents} and produces \(\mathbf{c}\) errors scaling with
\(\mathcal{K}_{\mathrm{eq}}\) \cite{stratton1941,harrington2001}. Full equivalence currents restore
\(\boldsymbol{\Gamma}_\omega\) invariance on \(\mathcal{F}_{\mathrm{rad}}\) \cite{devaney2004}.

\paragraph{Problem (multi-body enclosure).}
\textbf{Problem.} Two scatterers lie inside \(\Omega_{\mathrm{int}}\) but equivalence uses \(S\) enclosing only one body.
Can \(\mathbf{c}\) for the full structure be recovered?

\textbf{Step 1.} Identify unenclosed \(\operatorname{supp}\Jimp\) \cite{stratton1941}.

\textbf{Step 2.} Enlarge \(S\) to contain all radiating support or restrict \(\mathbf{c}\) to the partial exterior
\cite{harrington2001,collin2001}.

\textbf{Step 3.} Expect \(\mathcal{K}_{\mathrm{eq}}\) growth if \(S\) is placed near either body
Lemma~\ref{lem:ch4.equivalence-proximity} \cite{devaney2004}.

\textbf{Physical meaning.} Multi-body realizability requires multi-body enclosure \cite{stratton1941}.

\paragraph{Broadband equivalence and causal reduction.}
Equivalence at \(\omega_{0}\) does not certify \(\mathbf{c}(\omega)\) on a band \([\omega_{1},\omega_{2}]\) unless
\((\mathbf{J}_{s},\mathbf{M}_{s})\) are reconstructed frequency-by-frequency with causal time-domain reduction
\cite{jackson1999,stratton1941,harrington2001}. Define the broadband equivalence residual
\begin{equation}
\label{eq:ch4.broadband-equivalence-residual}
\delta_{\mathrm{bb}}
=
\sup_{\omega\in[\omega_{1},\omega_{2}]}
\big\|
\boldsymbol{\Gamma}_\omega[\mathbf{J}_{\mathrm{eq}}(\omega)]
-
\boldsymbol{\Gamma}_\omega[\Jimp(\omega)]
\big\|_{\mathcal{C}},
\end{equation}
with \(\mathbf{J}_{\mathrm{eq}}(\omega)\) the Stratton--Chu reduction at each \(\omega\)
\cite{stratton1941,collin2001}. Wideband pulse synthesis requires \(\delta_{\mathrm{bb}}=0\) and
\(\mathcal{D}_{\mathrm{caus}}=1\) on the excitation tuple \cite{jackson1999,harrington2001}.

\begin{corollary}[Single-frequency equivalence trap]
\label{cor:ch4.single-freq-equivalence-trap}
If \(\delta_{\mathrm{bb}}>0\) on a declared band, then \(\mathbf{c}(\omega)\) fitted only at \(\omega_{0}\) is not a
wideband synthesis certificate even when Theorem~\ref{thm:ch4.volume-surface-equivalence} holds at \(\omega_{0}\)
\cite{stratton1941,devaney2004,harrington2001}. Dispersive interiors amplify \(\delta_{\mathrm{bb}}\) through
frequency-dependent tangential traces on \(S\) \cite{collin2001}.
\end{corollary}

\paragraph{Worked illustration (dispersive radome).}
A frequency-independent \(\mathbf{J}_{s}\) fit at \(\omega_{0}\) on a dispersive radome aperture fails
\(\delta_{\mathrm{bb}}=0\) at \(\omega_{0}\pm10\%\): \(\mathcal{K}_{\mathrm{eq}}\) in
Theorem~\ref{thm:ch4.equivalence-error} grows with trace mismatch across the band \cite{stratton1941,harrington2001}.
Wideband \(\mathbf{c}\) targets require dispersive equivalence currents or full volume models \cite{devaney2004}.

\paragraph{Problem (wideband pattern from single-frequency equivalence).}
\textbf{Problem.} An engineer fits equivalent \(\mathbf{J}_{s}\) at 10\,GHz and extrapolates \(\mathbf{c}\) to 8--12\,GHz.
When is extrapolation valid?

\textbf{Step 1.} Compute \(\delta_{\mathrm{bb}}\) in Eq.~\eqref{eq:ch4.broadband-equivalence-residual} on the band
\cite{jackson1999,stratton1941}.

\textbf{Step 2.} Test \(\mathcal{D}_{\mathrm{caus}}=1\) for the time-domain reduction of \(\mathbf{J}_{s}(\omega)\)
\cite{harrington2001}.

\textbf{Step 3.} If either test fails, restrict \(\mathbf{c}\) claims to narrowband synthesis at \(\omega_{0}\)
\cite{devaney2004,collin2001}.

\textbf{Physical meaning.} Equivalence is a frequency-local representation change unless causal broadband reduction is
explicit \cite{stratton1941}.

\paragraph{Null-space enlargement under equivalence.}
Equivalent-source non-uniqueness enlarges \(\mathcal{N}_{\mathrm{syn}}\) through
Eq.~\eqref{eq:ch4.equivalence-null-ambiguity} without enlarging \(\mathcal{S}_{\mathrm{real}}(\mathcal{C})\)
\cite{stratton1941,harrington2001,devaney2004}. The synthesis quotient
\(\mathcal{J}_{\mathrm{real}}/\mathcal{N}_{\mathrm{syn}}\) is therefore smaller in cardinality of representatives but
unchanged in \(\mathbf{c}\) image \cite{devaney2004}. Near-field discrimination of equivalent pairs is generally
impossible from \(\mathbf{c}\) alone \cite{harrington2001}.

\begin{theorem}[Equivalence class fiber over \(\mathbf{c}\)]
\label{thm:ch4.equivalence-fiber}
For \(\mathbf{c}_{0}\in\mathcal{S}_{\mathrm{real}}(\mathcal{C})\), the fiber
\(\mathcal{F}_{\mathbf{c}_{0}}=\{\mathbf{J}_{\mathrm{eq}}:\boldsymbol{\Gamma}_\omega[\mathbf{J}_{\mathrm{eq}}]=\mathbf{c}_{0}\}\)
has dimension at least \(\dim\ker\mathcal{R}_\omega|_{S}\) on the equivalence surface \cite{stratton1941,harrington2001}.
MoM inversion selects one point in \(\mathcal{F}_{\mathbf{c}_{0}}\); realizability audits must test \(\mathbf{c}_{0}\), not the
particular \(\mathbf{J}_{\mathrm{eq}}\) representative \cite{devaney2004,collin2001}.
\end{theorem}
\begin{proof}
Linearity and non-uniqueness of equivalent sources on \(S\) \cite{stratton1941}; kernel inclusion
Eq.~\eqref{eq:ch4.rank-null-composition} \cite{harrington2001}.
\end{proof}

\paragraph{Worked illustration (equivalent pair indistinguishability).}
Two distinct \(\mathbf{J}_{s}^{(A)},\mathbf{J}_{s}^{(B)}\) on \(\Sigma_{a}\) yield identical \(\mathbf{c}\) but different near
fields: only \(\mathbf{c}\) is synthesis-certified \cite{stratton1941,harrington2001,devaney2004}.

\paragraph{Problem (inverse uniqueness).}
\textbf{Problem.} MoM returns \(\mathbf{J}_{s}\) claiming uniqueness. Is \(\mathbf{J}_{s}\) unique?

\textbf{Step 1.} Test \(\mathbf{J}_{s}^{(1)}-\mathbf{J}_{s}^{(2)}\in\mathcal{N}_{\mathrm{syn}}\) \cite{stratton1941}.

\textbf{Step 2.} If difference lies in equivalence null space, \(\mathbf{c}\) is still unique \cite{harrington2001,devaney2004}.

\textbf{Step 3.} Near-field validation cannot certify \(\mathbf{J}_{s}\) without additional constraints outside
\(\boldsymbol{\Gamma}_\omega\) \cite{collin2001}.

\textbf{Physical meaning.} Equivalence solves for \(\mathbf{c}\), not for a unique current portrait \cite{stratton1941}.

\paragraph{Validity.}
Theorem~\ref{thm:ch4.equivalence-error} assumes linear Maxwell operators and trace continuity on Lipschitz \(S\)
\cite{stratton1941}. Strong scatterers outside \(S\) require enlarged equivalence surfaces or multiple
equivalence regions \cite{harrington2001,collin2001}. Broadband equivalence requires frequency-by-frequency
\(\mathbf{J}_{s}(\omega)\) with causal time-domain reduction \cite{jackson1999}.

Section~\ref{sec:ch42} has reduced volume synthesis to aperture and surface equivalents while recording
tangential screens and failure regimes. Section~\ref{sec:ch43} replaces geometric equivalence by operator chains
and singular-value rank on \(\mathcal{J}_{\mathrm{real}}\) \cite{harrington2001,devaney2004}.


\section{Finite Support and Radiation Operator Chains}
\label{sec:ch43}

Equivalence principles of Section~\ref{sec:ch42} identify \emph{which} surface sources reproduce a given
\(\mathfrak{F}\); compact support and operator composition identify \emph{how many} independent synthesis
directions that reproduction allows \cite{stratton1941,harrington2001,devaney2004}. Finite geometry makes
\(\mathcal{R}_\omega\) compact on bounded \(\mathcal{J}_{\mathrm{real}}\), hence \(\boldsymbol{\Gamma}_\omega\)
finite rank in every practical discretization \cite{hansen1988}. Section~\ref{sec:ch43} develops that chain:
Subsection~\ref{sec:ch431} fixes compact source domains; Subsection~\ref{sec:ch432} imports spectral
compression from Subsection~\ref{sec:ch213}; Subsection~\ref{sec:ch433} composes the source--field--coefficient
operators; Subsection~\ref{sec:ch434} defines \emph{radiative rank} by singular values
\cite{stratton1941,harrington2001}.


\subsection{Compact source domains}
\label{sec:ch431}
% TOC tag: [HOME]

Spectral compression is not an algorithmic choice imposed after the fact: it is the operator-theoretic
consequence of compact \(\operatorname{supp}\Jimp\) \cite{stratton1941,jackson1999,harrington2001}. Once
support is bounded, \(\mathcal{R}_\omega\) acting on \(\mathcal{J}_{\mathrm{real}}\) is compact, and
\(\boldsymbol{\Gamma}_\omega\) cannot have infinite effective rank at fixed \(\omega\) \cite{hansen1988}.

\paragraph{Assumptions.}
Let \(\Omega_{s}\Subset\mathbb{R}^{3}\) be compact with \(\operatorname{diam}(\Omega_{s})=D_{s}<\infty\) as in
Eq.~\eqref{eq:ch4.support-screen} \cite{stratton1941}. Impressed data satisfy
\(\mathcal{D}_{\mathrm{real}}=1\) from Eq.~\eqref{eq:ch4.realizability-screen} \cite{harrington2001}. Exterior
region carries outgoing Sommerfeld closure Section~\ref{sec:ch24} \cite{stratton1941}. Electrical size
\(\Xi=k_{0}D_{s}/(2\pi)\) from Eq.~\eqref{eq:ch4.electrical-size} parametrizes decay of far coefficients
\cite{jackson1999,hansen1988}.

\begin{definition}[Compact source domain]
\label{def:ch4.compact-source-domain}
A \emph{compact source domain} is a pair \((\Omega_{s},k_{0})\) with \(\Omega_{s}\) compact and
\(\operatorname{supp}\mathbf{J}_{\mathrm{tot}}\subseteq\Omega_{s}\) for every
\(\mathbf{J}_{\mathrm{tot}}\in\mathcal{J}_{\mathrm{real}}\) \cite{stratton1941}. The associated support
functional is \(\mathcal{D}_{\mathrm{supp}}=1\) in Eq.~\eqref{eq:ch4.support-screen}.
\end{definition}

\paragraph{Ball support and scaling.}
When \(\Omega_{s}=B_{a}=\{\mathbf{r}:|\mathbf{r}-\mathbf{r}_{c}|\le a\}\), the dimensionless parameter
\(k_{0}a\) controls how many angular orders carry appreciable power in Eq.~\eqref{eq:ch4.spectral-tail-energy}
\cite{stratton1941,jackson1999}. For \(k_{0}a\ll1\), \(\boldsymbol{\Gamma}_\omega[\mathcal{J}_{\mathrm{real}}]\) is
concentrated in low-\(\ell\) sectors; for \(k_{0}a\gg1\), higher \(\ell\) become accessible before aperture and
measurement compress further \cite{hansen1988,harrington2001}. Equation~\eqref{eq:ch4.electrical-size} unifies
volume and aperture supports: \(\Sigma_{a}\subset\partial\Omega_{s}\) inherits the same \(D_{s}\) when the
aperture diameter equals the enclosing scale \cite{collin2001}.

\begin{theorem}[Compactness of the radiation map on bounded sources]
\label{thm:ch4.compact-radiation-map}
Let \(\mathcal{J}_{\mathrm{real}}(\Omega_{s})\) denote realizable sources on compact \(\Omega_{s}\). Then
\(\mathcal{R}_\omega:\mathcal{J}_{\mathrm{real}}(\Omega_{s})\to\mathcal{F}_{\mathrm{rad}}\) is a compact linear
operator in the energy topology of Eq.~\eqref{eq:ch2.field-energy-norm} \cite{stratton1941,harrington2001}. For
every countable \(\mathcal{C}\in\mathcal{R}_{\mathrm{adm}}\),
\(\boldsymbol{\Gamma}_\omega:\mathcal{J}_{\mathrm{real}}(\Omega_{s})\to\mathbb{C}^{I}\) is compact into any
\(\ell^{2}(I)\) coefficient space with flux-normalized basis \cite{hansen1988,devaney2004}.
\end{theorem}
\begin{proof}
Bounded \(\Omega_{s}\) and retarded Green representation
Eq.~\eqref{eq:ch2.electric-dyadic-convolution} yield integral operators with square-integrable kernels on
exterior observation shells \cite{stratton1941,jackson1999}. Composition with outgoing projection
\(\Pi_{\mathrm{rad}}\) and finite-energy norm Eq.~\eqref{eq:ch2.field-energy-norm} gives compactness on
\(\mathcal{J}_\Omega\) as in Subsection~\ref{sec:ch213} \cite{harrington2001}. Flux-normalized
\(\mathcal{C}\) from Section~\ref{sec:ch35} is bounded on \(\mathcal{F}_{\mathrm{rad}}\), hence
\(\boldsymbol{\Gamma}_\omega\) is compact \cite{hansen1988,stratton1941}.
\end{proof}

Theorem~\ref{thm:ch4.compact-radiation-map} is the operator backbone of realizability: compactness forces
singular values of \(\boldsymbol{\Gamma}_\omega\) to decay, hence \(\mathcal{S}_{\mathrm{real}}(\mathcal{C})\) is
intrinsically finite-dimensional at any fixed truncation tolerance \cite{harrington2001,devaney2004}. Corollary
\ref{cor:ch4.synthesis-not-open} is the coefficient-space restatement of the same fact \cite{hansen1988}.

\begin{figure}[t]
\centering
\ChOneFigBlock{%
\begin{tikzpicture}[ch1fig, node distance=7mm and 10mm]
  \node[ch1box] (ball) {Compact $\Omega_{s}$\\$D_{s}<\infty$};
  \node[ch1box, right=13mm of ball] (j) {$\mathcal{J}_{\mathrm{real}}$};
  \node[ch1box, right=13mm of j] (rw) {$\mathcal{R}_\omega$\\compact};
  \node[ch1box, right=13mm of rw] (fr) {$\mathcal{F}_{\mathrm{rad}}$};
  \node[ch1box, below=12mm of rw] (xi) {$\Xi=k_{0}D_{s}/2\pi$};
  \draw[ch1flow] (ball) -- (j) -- (rw) -- (fr);
  \draw[ch1flow] (xi) -- (ball);
\end{tikzpicture}%
}
\caption{Subsection~\ref{sec:ch431}: compact support confines \(\mathcal{J}_{\mathrm{real}}\) and yields compact
\(\mathcal{R}_\omega\) on bounded sources.}
\label{fig:ch4.compact-source-domain}
\end{figure}

\paragraph{Worked illustration (electrically small sphere).}
A source in \(B_{a}\) with \(k_{0}a=0.5\) radiates primarily \(\ell=1\) content; \(\boldsymbol{\Gamma}_\omega\) is
well approximated by a rank-\(3\) operator in a dipole basis \cite{jackson1999,stratton1941}. Demanding
\(\ell=4\) content without enlarging \(a\) targets outside \(\boldsymbol{\Gamma}_\omega[\mathcal{J}_{\mathrm{real}}]\)
\cite{hansen1988,harrington2001}.

\paragraph{Hilbert--Schmidt control on compact supports.}
For \(\Jimp\in\mathcal{J}_{\mathrm{real}}(\Omega_{s})\) with \(\Omega_{s}\) compact, the radiation kernel induced by
Eq.~\eqref{eq:ch2.electric-dyadic-convolution} is square-integrable on exterior observation shells
\cite{stratton1941,jackson1999}. Hence \(\mathcal{R}_\omega\) is not merely compact but Hilbert--Schmidt on bounded
energy norms:
\begin{equation}
\label{eq:ch4.HS-bound}
\|\mathcal{R}_\omega\|_{\mathrm{HS}}^{2}
=
\int_{\Omega_{s}}\!\int_{\Omega_{s}}
\big|\mathbf{G}_{E}(\mathbf{r},\mathbf{r}',\omega)\big|^{2}\,d^{3}r\,d^{3}r'
<
\infty,
\end{equation}
with \(\mathbf{G}_{E}\) the retarded electric Green dyadic \cite{stratton1941,harrington2001}. Equation
\eqref{eq:ch4.HS-bound} quantifies how support volume enters export capacity: enlarging \(\Omega_{s}\) increases
\(\|\mathcal{R}_\omega\|_{\mathrm{HS}}\) and typically increases \(r_{\mathrm{rad}}(\varepsilon)\) at fixed
\(\varepsilon\) \cite{devaney2004,hansen1988}. Compactness is the continuum reason MoM matrices exhibit clustered
singular spectra on electrically small structures \cite{harrington2001}.

\begin{lemma}[Non-compact support removes decay]
\label{lem:ch4.noncompact-no-decay}
If \(\operatorname{supp}\Jimp\) is non-compact (infinite line, half-space feed), \(\mathcal{R}_\omega\) need not be
compact on the declared source class and \(\sigma_{n}\) may not decay to zero \cite{stratton1941,jackson1999}. Then
\(\mathcal{S}_{\mathrm{real}}(\mathcal{C})\) is not intrinsically finite-dimensional at any fixed truncation tolerance
\cite{harrington2001,devaney2004}.
\end{lemma}
\begin{proof}
Non-compact supports violate the integral condition underlying Eq.~\eqref{eq:ch4.HS-bound}; unbounded travel along the
feed direction sustains non-decaying lateral spectral content in \(\mathcal{F}_{\mathrm{rad}}\) \cite{jackson1999}.
\end{proof}

Lemma~\ref{lem:ch4.noncompact-no-decay} explains why \(\mathcal{D}_{\mathrm{supp}}=1\) is a default realizability
screen: without compact support, Corollary~\ref{cor:ch4.synthesis-not-open} and radiative rank laws of
Subsection~\ref{sec:ch434} need not apply \cite{stratton1941,harrington2001}.

\paragraph{Problem (enlarging ball support).}
\textbf{Problem.} A target \(\mathbf{c}_{\mathrm{target}}\) requires \(r_{\mathrm{rad}}(10^{-2})=12\) on \(B_{a}\) with
\(k_{0}a=1.5\). Can support enlargement to \(B_{a'}\) with \(a'=2a\) recover realizability without changing frequency?

\textbf{Step 1.} Corollary~\ref{cor:ch4.support-monotone-enlargement} gives
\(\mathcal{S}_{\mathrm{real}}(\mathcal{C};a)\subseteq\mathcal{S}_{\mathrm{real}}(\mathcal{C};a')\) \cite{hansen1988}.

\textbf{Step 2.} Scaling \(k_{0}a\to 2k_{0}a\) in Eq.~\eqref{eq:ch4.ellmax-from-tail} increases admissible
\(\ell_{\max}\) and typically raises \(r_{\mathrm{rad}}\) \cite{jackson1999,stratton1941}.

\textbf{Step 3.} Recompute \(\boldsymbol{\Gamma}_\omega\) on \(B_{a'}\); if \(\mathbf{c}_{\mathrm{target}}\) still has
nonzero residual \(\mathbf{r}_{\min}\) from Principle~\ref{prin:ch4.operator-chain}, the obstruction is coupling or
aperture rank, not support alone \cite{devaney2004,harrington2001}.

\textbf{Physical meaning.} Support growth enlarges synthesis sets monotonically but does not guarantee every target
\cite{hansen1988}.

\paragraph{Support functional and synthesis screening.}
The support screen Eq.~\eqref{eq:ch4.support-screen} is the first gate in the realizability tuple: without
\(\mathcal{D}_{\mathrm{supp}}=1\), \(\mathcal{R}_\omega\) need not be compact and \(r_{\mathrm{rad}}(\varepsilon)\) may
be infinite \cite{stratton1941,jackson1999,harrington2001}. Define the \emph{support diameter functional}
\begin{equation}
\label{eq:ch4.support-diameter-functional}
\mathcal{D}_{\mathrm{diam}}[\Omega_{s}]
=
\operatorname{diam}(\Omega_{s}),
\qquad
\Xi[\Omega_{s}]
=
\frac{k_{0}\mathcal{D}_{\mathrm{diam}}[\Omega_{s}]}{2\pi},
\end{equation}
linking geometry directly to electrical size used in Eqs.~\eqref{eq:ch4.ellmax-from-tail} and
\eqref{eq:ch4.geometry-coupling-scaling} \cite{hansen1988,stratton1941}. Equation~\eqref{eq:ch4.support-diameter-functional}
makes support a synthesis parameter: two feeds with identical \(\mathbf{c}\) targets but different
\(\mathcal{D}_{\mathrm{diam}}\) belong to different \(\mathcal{S}_{\mathrm{real}}(\mathcal{C})\) classes
\cite{devaney2004}.

\begin{proposition}[Nested supports and synthesis sets]}
\label{prop:ch4.nested-support-synthesis}
If \(\Omega_{s}^{(1)}\subseteq\Omega_{s}^{(2)}\), then
\(\mathcal{S}_{\mathrm{real}}(\mathcal{C};\Omega_{s}^{(1)})\subseteq\mathcal{S}_{\mathrm{real}}(\mathcal{C};\Omega_{s}^{(2)})\)
for identical \(\mathcal{C}\) and \(\omega\) \cite{stratton1941,hansen1988}. Strict inclusion holds when
\(\Xi[\Omega_{s}^{(2)}]>\Xi[\Omega_{s}^{(1)}]\) activates new eigenfield couplings
\cite{jackson1999,devaney2004}.
\end{proposition}

\paragraph{Worked illustration (offset feed).}
Translating \(\Omega_{s}\) without changing \(\operatorname{diam}(\Omega_{s})\) leaves \(\Xi\) and compactness unchanged;
\(\mathcal{S}_{\mathrm{real}}(\mathcal{C})\) rotates in coefficient space through Wigner mixing but does not enlarge rank
\cite{jackson1999,hansen1988,stratton1941}. Enlarging \(\Omega_{s}\) is the primary rank enlargement mechanism, not
translation \cite{harrington2001}.

\paragraph{Problem (line feed compactness).}
\textbf{Problem.} A half-wave dipole along \(\hat{\mathbf{z}}\) is modeled as infinite line current. Does
Theorem~\ref{thm:ch4.compact-radiation-map} apply?

\textbf{Step 1.} Test \(\mathcal{D}_{\mathrm{supp}}=1\): non-compact support triggers
Lemma~\ref{lem:ch4.noncompact-no-decay} \cite{jackson1999,stratton1941}.

\textbf{Step 2.} Truncate to finite \(\mathcal{C}_{f}\) with \(\mathcal{D}_{\mathrm{cpl}}=1\) or declare a separate
non-compact domain class \cite{harrington2001}.

\textbf{Step 3.} Recompute \(r_{\mathrm{rad}}(\varepsilon)\) only on the declared compact subclass \cite{devaney2004}.

\textbf{Physical meaning.} Compactness is a domain declaration, not a numerical convenience \cite{stratton1941}.

\paragraph{Green-kernel compactness and export capacity.}
Compact \(\Omega_{s}\) makes the retarded convolution \(\mathcal{S}_\omega\) a compact operator on bounded source
norms \cite{stratton1941,harrington2001}. The export capacity functional
\begin{equation}
\label{eq:ch4.export-capacity-functional}
\mathcal{C}_{\mathrm{exp}}[\Omega_{s}]
=
\sup_{\|\Jimp\|=1}
\|\mathcal{R}_\omega[\Jimp]\|_{\mathrm{rad}},
\end{equation}
is finite if and only if \(\mathcal{D}_{\mathrm{supp}}=1\) and \(\mathcal{D}_{\mathrm{real}}=1\)
\cite{stratton1941,devaney2004}. For ball support \(B_{a}\),
\(\mathcal{C}_{\mathrm{exp}}\) grows monotonically with \(k_{0}a\) and \(\operatorname{vol}(B_{a})\) through
Eq.~\eqref{eq:ch4.HS-bound} \cite{jackson1999,hansen1988}.

\begin{theorem}[Support volume scaling of synthesis dimension]
\label{thm:ch4.support-volume-scaling}
On \(B_{a}\) with \(k_{0}a\ge1\),
\begin{equation}
\label{eq:ch4.support-volume-scaling}
r_{\mathrm{rad}}(\varepsilon)
\ \le\
C_{\mathrm{vol}}(\varepsilon)\,(k_{0}a)^{3},
\end{equation}
with \(C_{\mathrm{vol}}\) depending only on \(\varepsilon\) and flux gauge \cite{stratton1941,jackson1999,hansen1988}.
Consequently \(\dim\mathcal{S}_{\mathrm{real}}(\mathcal{C};B_{a})\) cannot grow faster than volume scaling at fixed
\(\varepsilon\) \cite{harrington2001,devaney2004}.
\end{theorem}
\begin{proof}
Combine Theorem~\ref{thm:ch4.compact-radiation-map}, Eq.~\eqref{eq:ch4.weyl-mode-count}, and
Theorem~\ref{thm:ch4.radiative-rank} \cite{stratton1941,hansen1988,harrington2001}.
\end{proof}

\begin{figure}[t]
\centering
\ChOneFigBlock{%
\begin{tikzpicture}[ch1fig, node distance=7mm and 10mm]
  \node[ch1box] (vol) {$\operatorname{vol}(\Omega_{s})$};
  \node[ch1box, right=12mm of vol] (hs) {Eq.~\eqref{eq:ch4.HS-bound}};
  \node[ch1box, right=12mm of hs] (rr) {$r_{\mathrm{rad}}(\varepsilon)$};
  \node[ch1box, right=12mm of rr] (sr) {$\mathcal{S}_{\mathrm{real}}$};
  \draw[ch1flow] (vol) -- (hs) -- (rr) -- (sr);
\end{tikzpicture}%
}
\caption{Subsection~\ref{sec:ch431}: compact support volume controls Hilbert--Schmidt bound and radiative rank before
coefficient listing.}
\label{fig:ch4.support-volume-scaling}
\end{figure}

\paragraph{Worked illustration (cube versus sphere).}
Two supports with equal \(\operatorname{vol}\) but different \(\operatorname{diam}\) share the same
Eq.~\eqref{eq:ch4.support-volume-scaling} exponent yet differ in directional rank: cubes activate corner couplings
earlier than spheres at the same \(k_{0}a\) \cite{stratton1941,jackson1999,harrington2001}. Directional audits of
\(\mathbf{T}\) are required when \(\Omega_{s}\) is non-spherical \cite{devaney2004}.

\paragraph{Problem (volume doubling without rank gain).}
\textbf{Problem.} \(\operatorname{vol}(\Omega_{s})\) doubles by offset duplication of identical sub-feeds. Does
\(r_{\mathrm{rad}}(\varepsilon)\) double?

\textbf{Step 1.} Test linear independence of duplicated \(\boldsymbol{\Gamma}_\omega[\Jimp_{i}]\) images
\cite{harrington2001,devaney2004}.

\textbf{Step 2.} If images are collinear, rank does not increase despite volume growth \cite{stratton1941}.

\textbf{Step 3.} Use SVD of \(\mathbf{T}\), not \(\operatorname{vol}\) alone, as the rank certificate \cite{hansen1988}.

\textbf{Physical meaning.} Volume scaling is an upper bound; degenerate feeds may saturate below it \cite{stratton1941}.

\paragraph{Validity.}
Theorem~\ref{thm:ch4.compact-radiation-map} assumes linear Maxwell theory and fixed \(\omega\)
\cite{stratton1941}. Non-compact feeds (infinite lines, half-spaces) require separate domain declarations;
\(\mathcal{D}_{\mathrm{supp}}=0\) removes Corollary~\ref{cor:ch4.synthesis-not-open} compactness
\cite{harrington2001,jackson1999}.

Spectral compression laws imported from Subsection~\ref{sec:ch213} are developed in Subsection~\ref{sec:ch432}.


\subsection{Spectral compression from finite geometry}
\label{sec:ch432}
% TOC tag: [HOME]

Definition~\ref{def:ch4.spectral-compression} introduced tail budgets on coefficient space; Subsection
~\ref{sec:ch432} identifies the same phenomenon as operator compression on \(\mathcal{R}_\omega\) when
\(\Omega_{s}\) is compact \cite{stratton1941,harrington2001}. The abstract template is
Eq.~\eqref{eq:ch2.spectral-compression-operator} from Subsection~\ref{sec:ch213}; the present subsection
specializes it to synthesis on \(\mathcal{J}_{\mathrm{real}}\) \cite{hansen1988}.

\paragraph{Assumptions.}
Presuppose compact \((\Omega_{s},k_{0})\) from Definition~\ref{def:ch4.compact-source-domain}, compact
\(\mathcal{R}_\omega\) from Theorem~\ref{thm:ch4.compact-radiation-map}, and VSH flux normalization from
Section~\ref{sec:ch35} \cite{stratton1941,hansen1988}. Truncation level \(N\) may denote modal index cutoff or
singular-value cutoff interchangeably until Subsection~\ref{sec:ch434} fixes the SVD gauge \cite{harrington2001}.

\paragraph{Compressed radiation operator.}
Let \(\{\mathbf{u}_{n}\}_{n\ge1}\) be radiation eigenpairs Eq.~\eqref{eq:ch2.radiation-eigenpair-def} ordered by
decreasing export gain \(|\lambda_{n}|\) \cite{stratton1941}. Define the rank-\(N\) compressed map
\begin{equation}
\label{eq:ch4.compressed-radiation-operator}
\mathcal{R}_\omega^{(N)}
=
\sum_{n=1}^{N}\lambda_{n}\,P_{n}\circ\mathcal{S}_\omega,
\end{equation}
importing Eq.~\eqref{eq:ch2.spectral-compression-operator} \cite{stratton1941,harrington2001}. On compact
\(\Omega_{s}\), the compression error satisfies
\begin{equation}
\label{eq:ch4.compression-error-import}
\|\mathcal{R}_\omega-\mathcal{R}_\omega^{(N)}\|_{\mathrm{op}}
\le
C\,\sigma_{N+1}\!\big(\mathcal{R}_\omega|_{\mathcal{J}_{\mathrm{real}}}\big),
\end{equation}
from Eq.~\eqref{eq:ch2.compression-error-bound} with \(C\) independent of \(\Jimp\in\mathcal{J}_{\mathrm{real}}\)
\cite{harrington2001}. Equation~\eqref{eq:ch4.compression-error-import} bounds watt residuals from truncating
\(\mathcal{R}_\omega\) before \(\mathcal{C}\) is applied: geometry compresses fields, not merely coefficients
\cite{stratton1941,devaney2004}.

\paragraph{Coefficient-space compression.}
Composing Eq.~\eqref{eq:ch4.compressed-radiation-operator} with \(\mathcal{C}\),
\begin{equation}
\label{eq:ch4.compressed-gamma-operator}
\boldsymbol{\Gamma}_\omega^{(N)}
=
\mathcal{C}\circ\mathcal{R}_\omega^{(N)}:
\mathcal{J}_{\mathrm{real}}\to\mathbb{C}^{I},
\qquad
\|\boldsymbol{\Gamma}_\omega-\boldsymbol{\Gamma}_\omega^{(N)}\|_{\mathrm{op}}
\le
C'\,\sigma_{N+1}\!\big(\boldsymbol{\Gamma}_\omega\big),
\end{equation}
with \(C'\) depending only on normalization of \(\mathcal{C}\) \cite{hansen1988,harrington2001}. Equation
\eqref{eq:ch4.compressed-gamma-operator} links Proposition~\ref{prop:ch4.compressed-synthesis} to singular-value
tail: \(\mathcal{T}_{\ell_{\max}}\) in coefficient space and \(\mathcal{R}_\omega^{(N)}\) in field space are
equivalent compressions when \(\{\mathbf{u}_{n}\}\) align with VSH modes Section~\ref{sec:ch34}
\cite{stratton1941,hansen1988}.

\begin{proposition}[Geometry--spectral equivalence]
\label{prop:ch4.geometry-spectral-equivalence}
For ball support \(B_{a}\) and VSH listing \(\mathcal{C}\), there exists \(N_{\mathrm{geo}}(\eta_{\mathrm{tail}},k_{0}a)\)
such that Eq.~\eqref{eq:ch4.spectral-compression-def} holds with \(\ell_{\max}=N_{\mathrm{geo}}\) if and only if
Eq.~\eqref{eq:ch4.compression-error-import} holds with \(N=N_{\mathrm{geo}}\) at tolerance \(\eta_{\mathrm{tail}}\)
\cite{stratton1941,jackson1999,hansen1988}.
\end{proposition}
\begin{proof}
Tail energy Eq.~\eqref{eq:ch4.spectral-tail-energy} equals \(\sum_{n>N}|\lambda_{n}|^{2}\|\widehat{c}_{n}\|^{2}\) in
an aligned eigenbasis \cite{hansen1988,stratton1941}. Compactness Theorem~\ref{thm:ch4.compact-radiation-map} gives
\(\sigma_{n}\sim|\lambda_{n}|\) ordering for export modes \cite{harrington2001}. Equivalence of cutoffs follows
from Parseval on \(\mathcal{F}_{\mathrm{rad}}\) Eq.~\eqref{eq:ch3.mode-power-from-coefficient} \cite{stratton1941}.
\end{proof}

Proposition~\ref{prop:ch4.geometry-spectral-equivalence} is the inevitability statement of Subsection
~\ref{sec:ch432}: electrical size and singular-value tails are two readings of the same compression, not separate
engineering knobs \cite{devaney2004,hansen1988}. Section~\ref{sec:ch46} develops scaling exponents for
\(N_{\mathrm{geo}}\); Section~\ref{sec:ch48} replaces \(N\) by concentration bounds \cite{harrington2001}.

\paragraph{Worked illustration (tail budget).}
At \(k_{0}a=2\), suppose \(\sigma_{N+1}/\sigma_{1}=10^{-3}\). Then rank-\(N\) synthesis captures \(99.9\%\) of
export gain in operator norm \cite{harrington2001}. Demanding \(\eta_{\mathrm{tail}}<10^{-3}\) in
Eq.~\eqref{eq:ch4.spectral-compression-def} requires \(N\ge N_{\mathrm{geo}}\) from Proposition
\ref{prop:ch4.geometry-spectral-equivalence} \cite{hansen1988,stratton1941}.

\paragraph{Weyl-type scaling of \(N_{\mathrm{geo}}\).}
On ball support \(B_{a}\), the number of export modes above threshold \(\varepsilon\) scales as
\begin{equation}
\label{eq:ch4.weyl-mode-count}
N_{\mathrm{geo}}(\eta_{\mathrm{tail}},k_{0}a)
\ \sim\
\mathcal{W}(k_{0}a)^{3}
\ \sim\
(k_{0}a)^{3},
\end{equation}
in the high-frequency limit within the compact-source regime \cite{stratton1941,jackson1999,hansen1988}. Equation
\eqref{eq:ch4.weyl-mode-count} is the volume scaling anticipated by Corollary~\ref{cor:ch4.support-monotone-enlargement}:
doubling \(a\) at fixed \(k_{0}\) increases available synthesis directions faster than doubling frequency alone on fixed
\(a\) \cite{devaney2004,harrington2001}. Section~\ref{sec:ch46} sharpens exponents; here the scaling explains why
electrically large volumes support high \(r_{\mathrm{rad}}\) while electrically small volumes do not
\cite{hansen1988}.

\paragraph{Problem (matching \(\ell_{\max}\) to SVD rank).}
\textbf{Problem.} A VSH listing retains \(\ell\le 12\) but SVD of \(\mathbf{T}\) shows \(\sigma_{n}/\sigma_{1}<10^{-2}\) for
\(n>8\). Which cutoff governs realizability?

\textbf{Step 1.} Proposition~\ref{prop:ch4.geometry-spectral-equivalence} identifies \(\ell_{\max}\) and \(N\) when bases
align \cite{hansen1988,stratton1941}.

\textbf{Step 2.} On misaligned meshes, \(\sigma_{n}\) governs synthesis stability; \(\ell_{\max}=12\) overstates
realizable directions if only eight singular values exceed threshold \cite{harrington2001,devaney2004}.

\textbf{Step 3.} Declare \(r_{\mathrm{rad}}(10^{-2})=8\) per Definition~\ref{def:ch4.radiative-rank} and truncate design
targets accordingly \cite{hansen1988}.

\textbf{Physical meaning.} Geometry compresses operators; listings must not outrun singular-value decay
\cite{stratton1941}.

\paragraph{Tail energy as operator residual.}
Define the compression residual operator
\begin{equation}
\label{eq:ch4.compression-residual-operator}
\mathbf{R}_{N}
=
\mathcal{R}_\omega-\mathcal{R}_\omega^{(N)},
\qquad
\|\mathbf{R}_{N}\|_{\mathrm{op}}
\le
C\,\sigma_{N+1},
\end{equation}
importing Eq.~\eqref{eq:ch4.compression-error-import} \cite{stratton1941,harrington2001}. On coefficient space,
\begin{equation}
\label{eq:ch4.coefficient-compression-residual}
\|\mathbf{c}-\mathbf{c}^{(N)}\|_{\mathrm{rad}}
\le
C'\,\sigma_{N+1}\big(\boldsymbol{\Gamma}_\omega\big)\,\|\Jimp\|_{\mathcal{J}},
\end{equation}
with \(\mathbf{c}^{(N)}=\mathcal{C}\circ\mathcal{R}_\omega^{(N)}[\Jimp]\) \cite{hansen1988,devaney2004}. Equations
\eqref{eq:ch4.compression-residual-operator}--\eqref{eq:ch4.coefficient-compression-residual} quantify synthesis error from
premature truncation: choosing \(N\) too small leaves \(\mathbf{c}^{(N)}\notin\mathcal{B}_{\mathrm{fs}}\); choosing \(N\)
too large does not enlarge \(\mathcal{S}_{\mathrm{real}}(\mathcal{C})\) beyond geometry rank \cite{harrington2001}.

\begin{theorem}[Spectral compression optimality]
\label{thm:ch4.spectral-compression-optimal}
Among all rank-\(N\) operators \(\widetilde{\mathcal{R}}\) on \(\mathcal{J}_{\mathrm{real}}\),
\(\mathcal{R}_\omega^{(N)}\) minimizes \(\|\mathcal{R}_\omega-\widetilde{\mathcal{R}}\|_{\mathrm{op}}\), and
\(\boldsymbol{\Gamma}_\omega^{(N)}\) minimizes \(\|\boldsymbol{\Gamma}_\omega-\widetilde{\boldsymbol{\Gamma}}\|_{\mathrm{op}}\)
at the same \(N\) \cite{harrington2001,devaney2004,stratton1941}.
\end{theorem}
\begin{proof}
Eckart--Young theorem applied to \(\mathcal{R}_\omega\) and \(\boldsymbol{\Gamma}_\omega\) with export eigenpairs
Eq.~\eqref{eq:ch2.radiation-eigenpair-def} \cite{stratton1941,harrington2001}.
\end{proof}

Theorem~\ref{thm:ch4.spectral-compression-optimal} elevates singular-value truncation from engineering habit to
realizability optimality: no other rank-\(N\) synthesis operator approximates \(\boldsymbol{\Gamma}_\omega\) better in
operator norm \cite{hansen1988}. Listing \(\ell_{\max}\) without matching \(N_{\mathrm{geo}}\) from
Proposition~\ref{prop:ch4.geometry-spectral-equivalence} wastes degrees of freedom or omits export modes
\cite{devaney2004}.

\paragraph{Worked illustration (rank-5 ball).}
At \(k_{0}a=2.5\), suppose \(\sigma_{6}/\sigma_{1}=10^{-2}\). Rank-\(5\) synthesis captures \(99\%\) export gain; retaining
\(\ell\le12\) in \(\mathcal{C}\) without checking \(\sigma_{n}\) overstates realizable directions
\cite{harrington2001,hansen1988,stratton1941}.

\paragraph{Problem (tail budget versus SVD).}
\textbf{Problem.} \(\eta_{\mathrm{tail}}=10^{-2}\) in Eq.~\eqref{eq:ch4.spectral-compression-def} and \(\varepsilon=10^{-2}\)
in Definition~\ref{def:ch4.radiative-rank}---must \(N_{\mathrm{geo}}=r_{\mathrm{rad}}(\varepsilon)\)?

\textbf{Step 1.} Align gauges via Proposition~\ref{prop:ch4.geometry-spectral-equivalence} \cite{hansen1988}.

\textbf{Step 2.} Compare \(\|\mathbf{R}_{\ell_{\max}}\|_{\mathrm{rad}}\) to \(\sigma_{r_{\mathrm{rad}}+1}\)
\cite{stratton1941,harrington2001}.

\textbf{Step 3.} Use the larger cutoff for conservative synthesis audits \cite{devaney2004}.

\textbf{Physical meaning.} Tail budgets and singular-value thresholds are equivalent only in aligned gauges
\cite{hansen1988}.

\paragraph{Compression hierarchy on \(\boldsymbol{\Gamma}_\omega\).}
Finite geometry induces a nested family of truncated synthesis sets
\begin{equation}
\label{eq:ch4.compression-hierarchy}
\mathcal{S}_{\mathrm{real}}^{(1)}(\mathcal{C})
\subseteq
\mathcal{S}_{\mathrm{real}}^{(2)}(\mathcal{C})
\subseteq
\cdots
\subseteq
\mathcal{S}_{\mathrm{real}}(\mathcal{C}),
\qquad
\mathcal{S}_{\mathrm{real}}^{(N)}(\mathcal{C})
=
\boldsymbol{\Gamma}_\omega^{(N)}[\mathcal{J}_{\mathrm{real}}],
\end{equation}
with strict inclusions until \(N\ge r_{\mathrm{rad}}(\varepsilon)\) \cite{stratton1941,harrington2001,devaney2004}. The
hierarchy is monotone in export gain and optimal at each \(N\) by Theorem~\ref{thm:ch4.spectral-compression-optimal}
\cite{hansen1988}.

\begin{corollary}[Compression saturation index]
\label{cor:ch4.compression-saturation}
Let \(N_{\mathrm{sat}}(\varepsilon)\) be the smallest \(N\) such that
\(\mathcal{S}_{\mathrm{real}}^{(N)}(\mathcal{C})=\mathcal{S}_{\mathrm{real}}^{(N+1)}(\mathcal{C})\). Then
\(N_{\mathrm{sat}}(\varepsilon)=r_{\mathrm{rad}}(\varepsilon)\) in the aligned eigenbasis of
Eq.~\eqref{eq:ch4.gamma-SVD} \cite{harrington2001,devaney2004,stratton1941}. Listing \(\ell_{\max}>N_{\mathrm{sat}}\)
without power above \(\varepsilon\sigma_{1}\) is gauge padding \cite{hansen1988}.
\end{corollary}

\paragraph{Worked illustration (premature \(\ell_{\max}\)).}
Retaining \(\ell_{\max}=20\) when \(N_{\mathrm{sat}}(10^{-2})=9\) inflates inversion dimension without enlarging
\(\mathcal{S}_{\mathrm{real}}(\mathcal{C})\): additional indices carry \(\sigma_{n}<10^{-2}\sigma_{1}\)
\cite{harrington2001,hansen1988,devaney2004}.

\paragraph{Problem (operator versus listing truncation).}
\textbf{Problem.} \(\mathcal{R}_\omega^{(N)}\) and \(\mathcal{T}_{\ell_{\max}}\) yield different watt residuals at the same
\(N=\ell_{\max}\). Which truncation governs synthesis?

\textbf{Step 1.} Apply Proposition~\ref{prop:ch4.geometry-spectral-equivalence} to align gauges \cite{hansen1988}.

\textbf{Step 2.} Compare \(\|\mathbf{R}_{N}\|_{\mathrm{op}}\) and \(\|\mathbf{R}_{\ell_{\max}}\|_{\mathrm{rad}}\)
\cite{stratton1941,harrington2001}.

\textbf{Step 3.} Use the larger residual bound for conservative design \cite{devaney2004}.

\textbf{Physical meaning.} Compression is operator-defined; listings inherit it only after alignment \cite{stratton1941}.

\paragraph{Validity.}
Eqs.~\eqref{eq:ch4.compressed-radiation-operator}--\eqref{eq:ch4.compressed-gamma-operator} assume aligned
eigenbases; basis mismatch rotates tails without changing power \cite{hansen1988}. Continuous exterior spectra
require rigged projectors Subsection~\ref{sec:ch313}; discrete \(N\) is a narrowband certificate \cite{stratton1941}.

Operator composition form of the chain is Subsection~\ref{sec:ch433}.


\subsection{Source-radiation-field operator chains}
\label{sec:ch433}
% TOC tag: [HOME]

Realizability audits are audits of \emph{composed} maps: impressed data pass through solution, radiation,
and listing operators before coefficients enter \(\mathcal{S}_{\mathrm{real}}(\mathcal{C})\)
\cite{stratton1941,harrington2001,devaney2004}. Subsection~\ref{sec:ch433} writes that composition as a single
chain and prepares its singular-value factorization.

\paragraph{Assumptions.}
Fix \(\mathcal{J}_{\mathrm{real}}\), \(\mathcal{R}_\omega\), \(\mathcal{C}\in\mathcal{R}_{\mathrm{adm}}\), and
flux normalization \(\mathcal{N}\) from Section~\ref{sec:ch35} \cite{stratton1941,hansen1988}. All maps are linear
on declared domains \cite{stratton1941}.

\paragraph{Master composition chain.}
The source--radiation--field--coefficient chain is
\begin{equation}
\label{eq:ch4.master-operator-chain}
\mathbf{c}
=
\mathcal{C}\circ\mathcal{N}\circ\mathcal{R}_\omega\circ\mathcal{S}_\omega[\Jimp]
=
\boldsymbol{\Gamma}_\omega[\Jimp],
\qquad
\Jimp\in\mathcal{J}_{\mathrm{real}},
\end{equation}
extending Eq.~\eqref{eq:ch4.coefficient-synthesis-map} with explicit \(\mathcal{S}_\omega\) and \(\mathcal{N}\)
\cite{stratton1941,harrington2001}. Equation~\eqref{eq:ch4.master-operator-chain} is the realizability backbone:
each factor is fixed by prior chapters; Chapter~\ref{ch:04} analyzes rank and inversion of the composition
\cite{devaney2004}. Reactive content removed by \(\Pi_{\mathrm{rad}}\subset\mathcal{R}_\omega\) never reaches
\(\mathbf{c}\) \cite{stratton1941}.

\paragraph{Field-space intermediate.}
Define the field-space radiation operator
\begin{equation}
\label{eq:ch4.field-radiation-operator}
\mathfrak{R}_\omega
=
\mathcal{N}\circ\mathcal{R}_\omega\circ\mathcal{S}_\omega:
\mathcal{J}_{\mathrm{real}}\to\mathcal{F}_{\mathrm{rad}},
\end{equation}
so \(\mathbf{c}=\mathcal{C}\circ\mathfrak{R}_\omega[\Jimp]\) \cite{stratton1941,hansen1988}. On compact
\(\Omega_{s}\), \(\mathfrak{R}_\omega\) is compact; \(\mathcal{C}\) is bounded on
\(\mathcal{F}_{\mathrm{rad}}\cap\mathcal{V}\) \cite{harrington2001}. Singular values of
\(\boldsymbol{\Gamma}_\omega\) satisfy \(\sigma_{n}(\boldsymbol{\Gamma}_\omega)\le
\|\mathcal{C}\|_{\mathrm{op}}\,\sigma_{n}(\mathfrak{R}_\omega)\) \cite{stratton1941,devaney2004}.

\begin{figure}[t]
\centering
\ChOneFigBlock{%
\begin{tikzpicture}[ch1fig, node distance=6mm and 9mm]
  \node[ch1box] (j) {$\mathcal{J}_{\mathrm{real}}$};
  \node[ch1box, right=10mm of j] (s) {$\mathcal{S}_\omega$};
  \node[ch1box, right=10mm of s] (r) {$\mathcal{R}_\omega$};
  \node[ch1box, right=10mm of r] (n) {$\mathcal{N}$};
  \node[ch1box, right=10mm of n] (fr) {$\mathcal{F}_{\mathrm{rad}}$};
  \node[ch1box, right=10mm of fr] (c) {$\mathcal{C}$};
  \node[ch1box, right=10mm of c] (ci) {$\mathbb{C}^{I}$};
  \draw[ch1flow] (j) -- (s) -- (r) -- (n) -- (fr) -- (c) -- (ci);
\end{tikzpicture}%
}
\caption{Subsection~\ref{sec:ch433} master chain Eq.~\eqref{eq:ch4.master-operator-chain}.}
\label{fig:ch4.master-operator-chain}
\end{figure}

\paragraph{Discrete chain matrix.}
On a MoM basis \(\{\boldsymbol{\Lambda}_{p}\}_{p=1}^{P}\subset\mathcal{J}_{\mathrm{real}}\) and modal basis
\(\{\mathbf{f}_\alpha\}\) defining \(\mathcal{C}\),
\begin{equation}
\label{eq:ch4.discrete-chain-matrix}
\mathbf{c}
=
\mathbf{T}\,\mathbf{I}_{s},
\qquad
T_{\alpha p}
=
\mathcal{C}\Big[\mathcal{N}\big[\mathcal{R}_\omega[\boldsymbol{\Lambda}_{p}]\big]\Big]_{\alpha},
\end{equation}
extending Eq.~\eqref{eq:ch4.mom-synthesis-matrix} to the full chain \cite{harrington2001,stratton1941}. Equation
\eqref{eq:ch4.discrete-chain-matrix} is the matrix audited in design codes: \(\mathbf{T}\) encodes Green
convolution, outgoing projection, normalization, and listing in one pass \cite{harrington2001,collin2001}. Rank of
\(\mathbf{T}\) upper-bounds \(\dim\mathcal{S}_{\mathrm{real}}(\mathcal{C})\) by Corollary~\ref{cor:ch4.channel-count}
\cite{devaney2004}.

\paragraph{Adjoint chain and reciprocity.}
When reciprocity holds Section~\ref{sec:ch28}, the adjoint chain
\begin{equation}
\label{eq:ch4.adjoint-chain}
\boldsymbol{\Gamma}_\omega^{\dagger}
=
\mathcal{S}_\omega^{\dagger}\circ\mathcal{R}_\omega^{\dagger}\circ\mathcal{N}^{\dagger}\circ\mathcal{C}^{\dagger}
\end{equation}
maps coefficient functionals to source duals on \(\mathcal{J}_{\mathrm{real}}\) \cite{stratton1941,harrington2001}.
Inverse synthesis minimizes \(\|\boldsymbol{\Gamma}_\omega[\Jimp]-\mathbf{c}_{\mathrm{target}}\|\) using
\(\boldsymbol{\Gamma}_\omega^{\dagger}\) as a first guess; realizability requires \(\mathbf{c}_{\mathrm{target}}\in
\operatorname{ran}\boldsymbol{\Gamma}_\omega\) \cite{devaney2004,harrington2001}.

\begin{principle}[Operator-chain realizability]
\label{prin:ch4.operator-chain}
A coefficient target \(\mathbf{c}_{\mathrm{target}}\) is realizable if and only if it lies in
\(\operatorname{ran}\boldsymbol{\Gamma}_\omega\) on \(\mathcal{J}_{\mathrm{real}}\); equivalently, if and only if
the minimum-distance residual \(\mathbf{r}_{\min}=\mathbf{c}_{\mathrm{target}}-\mathcal{P}_{\mathrm{ran}}\mathbf{c}_{\mathrm{target}}\)
vanishes, with \(\mathcal{P}_{\mathrm{ran}}\) the orthogonal projector onto
\(\operatorname{ran}\boldsymbol{\Gamma}_\omega\) in \(\ell^{2}(I)\) \cite{stratton1941,devaney2004,harrington2001}.
\end{principle}

Principle~\ref{prin:ch4.operator-chain} replaces pattern-matching language with a range condition: the chain
\emph{defines} \(\mathcal{S}_{\mathrm{real}}(\mathcal{C})\) as \(\operatorname{ran}\boldsymbol{\Gamma}_\omega\)
\cite{hansen1988}. Nonzero \(\mathbf{r}_{\min}\) is an irreducible realizability obstruction at fixed
\((\Omega_{s},\omega,\mathcal{C})\) \cite{devaney2004}.

\paragraph{Worked illustration (rank test).}
Compute SVD \(\mathbf{T}=\mathbf{U}\boldsymbol{\Sigma}\mathbf{V}^{\dagger}\) from
Eq.~\eqref{eq:ch4.discrete-chain-matrix}. If \(\mathbf{c}_{\mathrm{target}}\) has component on
\(\ker\mathbf{T}^{\dagger}\), the target is not realizable on the mesh \cite{harrington2001,devaney2004}.

\paragraph{Chain factorization audit.}
Write the master chain as explicit composition
\begin{equation}
\label{eq:ch4.chain-factorization}
\boldsymbol{\Gamma}_\omega
=
\underbrace{\mathcal{C}}_{\text{list}}
\circ
\underbrace{\mathcal{N}}_{\text{gauge}}
\circ
\underbrace{\Pi_{\mathrm{rad}}}_{\text{export}}
\circ
\underbrace{\mathcal{S}_\omega}_{\text{solve}}
\quad\text{on}\ \mathcal{J}_{\mathrm{real}},
\end{equation}
importing Eq.~\eqref{eq:ch4.master-operator-chain} \cite{stratton1941,harrington2001}. Each factor is audited
separately in synthesis failure diagnosis: \(\mathcal{S}_\omega\) failures are source or boundary class errors;
\(\Pi_{\mathrm{rad}}\) failures are reactive contamination; \(\mathcal{N}\) failures are gauge mismatches;
\(\mathcal{C}\) failures are overcomplete listings \cite{devaney2004,hansen1988}.

\begin{figure}[t]
\centering
\ChOneFigBlock{%
\begin{tikzpicture}[ch1fig, node distance=7mm and 9mm]
  \node[ch1box] (s) {$\mathcal{S}_\omega$};
  \node[ch1box, right=10mm of s] (pi) {$\Pi_{\mathrm{rad}}$};
  \node[ch1box, right=10mm of pi] (n) {$\mathcal{N}$};
  \node[ch1box, right=10mm of n] (c) {$\mathcal{C}$};
  \node[ch1box, below=11mm of pi] (fail) {failure localization};
  \draw[ch1flow] (s) -- (pi) -- (n) -- (c);
  \draw[ch1flow] (fail) -- (s);
  \draw[ch1flow] (fail) -- (pi);
\end{tikzpicture}%
}
\caption{Subsection~\ref{sec:ch433}: operator-chain audits localize realizability failures by composition factor.}
\label{fig:ch4.chain-factorization-audit}
\end{figure}

\paragraph{Residual norm as synthesis certificate.}
Define
\begin{equation}
\label{eq:ch4.synthesis-residual-norm}
\delta_{\mathrm{syn}}
=
\|\mathbf{c}_{\mathrm{target}}-\mathcal{P}_{\mathrm{ran}}\mathbf{c}_{\mathrm{target}}\|_{\mathrm{rad}},
\end{equation}
with \(\mathcal{P}_{\mathrm{ran}}\) from Principle~\ref{prin:ch4.operator-chain} \cite{stratton1941,devaney2004}.
Then \(\mathbf{c}_{\mathrm{target}}\in\mathcal{S}_{\mathrm{real}}(\mathcal{C})\iff\delta_{\mathrm{syn}}=0\)
\cite{harrington2001}. Nonzero \(\delta_{\mathrm{syn}}\) lower-bounds irreducible pattern error at fixed
\((\Omega_{s},\omega,\mathcal{C})\) \cite{hansen1988}.

\paragraph{Worked illustration (reactive contamination).}
Near-field fitting that omits \(\Pi_{\mathrm{rad}}\) yields \(\delta_{\mathrm{syn}}>0\) even when MoM residuals are
small: reactive components inflate \(\mathcal{S}_\omega[\Jimp]\) without entering \(\mathbf{c}\)
\cite{stratton1941,harrington2001,devaney2004}.

\paragraph{Problem (chain failure localization).}
\textbf{Problem.} \(\mathbf{c}_{\mathrm{target}}\) fails \(\delta_{\mathrm{syn}}=0\) but MoM reports small patch residuals.
Where does the chain fail?

\textbf{Step 1.} Apply \(\Pi_{\mathrm{rad}}\) to the reconstructed near field \cite{stratton1941}.

\textbf{Step 2.} Recompute \(\mathbf{c}\) with \(\mathcal{N}\) fixed \cite{hansen1988}.

\textbf{Step 3.} If \(\delta_{\mathrm{syn}}\) vanishes after \(\Pi_{\mathrm{rad}}\), failure was reactive contamination, not
geometry rank \cite{harrington2001,devaney2004}.

\textbf{Physical meaning.} Chain factorization localizes obstructions \cite{stratton1941}.

\paragraph{Chain conditioning and inversion stability.}
The condition number of the discrete chain satisfies
\begin{equation}
\label{eq:ch4.chain-condition-number}
\kappa(\mathbf{T})
\le
\kappa(\mathcal{C})\,\kappa(\mathcal{N})\,\kappa(\mathcal{R}_\omega|_{\mathcal{J}_{\mathrm{real}}})\,
\kappa(\mathcal{S}_\omega|_{\mathcal{J}_{\mathrm{real}}}),
\end{equation}
with each factor the operator condition number on the declared domain \cite{harrington2001,stratton1941,devaney2004}.
Ill-conditioned \(\mathbf{T}\) signals near rank deficiency in \(\boldsymbol{\Gamma}_\omega\), not merely mesh noise:
inversion amplifies components along \(\ker\boldsymbol{\Gamma}_\omega^{\dagger}\) that lie outside
\(\mathcal{S}_{\mathrm{real}}(\mathcal{C})\) \cite{hansen1988}.

\begin{theorem}[Moore--Penrose synthesis limit]
\label{thm:ch4.moore-penrose-synthesis}
Let \(\mathbf{c}_{\mathrm{target}}\in\mathbb{C}^{I}\) and \(\mathbf{T}^{+}\) the Moore--Penrose pseudoinverse of
Eq.~\eqref{eq:ch4.discrete-chain-matrix}. The minimum-norm least-squares feed
\(\mathbf{I}_{s}^{\star}=\mathbf{T}^{+}\mathbf{c}_{\mathrm{target}}\) is realizable only if
\(\mathbf{c}_{\mathrm{target}}\in\operatorname{ran}\mathbf{T}\); otherwise
\(\|\mathbf{T}\mathbf{I}_{s}^{\star}-\mathbf{c}_{\mathrm{target}}\|=\delta_{\mathrm{syn}}>0\) with
\(\delta_{\mathrm{syn}}\) from Eq.~\eqref{eq:ch4.synthesis-residual-norm} \cite{harrington2001,devaney2004,stratton1941}.
\end{theorem}
\begin{proof}
Pseudoinverse optimality on \(\ell^{2}\) \cite{harrington2001}; Principle~\ref{prin:ch4.operator-chain} for range
membership \cite{devaney2004,stratton1941}.
\end{proof}

Theorem~\ref{thm:ch4.moore-penrose-synthesis} is the numerical reading of inevitability: least-squares MoM fits always
return a current, but \(\delta_{\mathrm{syn}}=0\) is the synthesis certificate \cite{hansen1988}. Tikhonov
regularization lowers \(\kappa(\mathbf{T})\) while enlarging \(\delta_{\mathrm{syn}}\)---a trade controlled by
\(r_{\mathrm{rad}}(\varepsilon)\) \cite{devaney2004,harrington2001}.

\paragraph{Worked illustration (ill-conditioned horn MoM).}
A horn MoM matrix with \(\kappa(\mathbf{T})\sim10^{8}\) inverts to \(\mathbf{c}_{\mathrm{MoM}}\approx\mathbf{c}_{\mathrm{target}}\)
while \(\delta_{\mathrm{syn}}>0\) in the flux gauge: the fit is numerical, not physical \cite{harrington2001,collin2001}.
Truncating to \(r_{\mathrm{rad}}(10^{-2})\) singular directions restores stable synthesis \cite{hansen1988,devaney2004}.

\paragraph{Problem (small residual, large condition number).}
\textbf{Problem.} \(\|\mathbf{T}\mathbf{I}_{s}-\mathbf{c}_{\mathrm{target}}\|<10^{-6}\) but \(\kappa(\mathbf{T})>10^{10}\).
Is \(\mathbf{c}_{\mathrm{target}}\) realizable?

\textbf{Step 1.} Compute \(\delta_{\mathrm{syn}}\) in the continuum gauge, not the discrete residual alone
\cite{stratton1941,harrington2001}.

\textbf{Step 2.} Inspect \(\sigma_{n}/\sigma_{1}\) for \(n>r_{\mathrm{rad}}(\varepsilon)\) \cite{hansen1988,devaney2004}.

\textbf{Step 3.} If \(\delta_{\mathrm{syn}}>0\), reject realizability despite small MoM residual \cite{stratton1941}.

\textbf{Physical meaning.} Discrete residuals do not override operator-range obstructions \cite{harrington2001}.

\paragraph{Factor-wise rank attribution.}
Each chain factor in Eq.~\eqref{eq:ch4.master-operator-chain} carries a local rank budget:
\begin{equation}
\label{eq:ch4.chain-factor-ranks}
r_{\mathrm{rad}}(\boldsymbol{\Gamma}_\omega)
\le
\min\big\{
r_{\mathrm{rad}}(\mathcal{C}),\,
r_{\mathrm{rad}}(\mathcal{N}),\,
r_{\mathrm{rad}}(\mathcal{R}_\omega),\,
r_{\mathrm{rad}}(\mathcal{S}_\omega)
\big\},
\end{equation}
with each term the export rank of the restricted map on \(\mathcal{J}_{\mathrm{real}}\) \cite{stratton1941,harrington2001,devaney2004}.
Failure localization: if \(r_{\mathrm{rad}}(\mathcal{R}_\omega)\ll r_{\mathrm{rad}}(\mathcal{S}_\omega)\), reactive contamination
dominates before listing \cite{harrington2001}; if \(r_{\mathrm{rad}}(\mathcal{C})\ll r_{\mathrm{rad}}(\mathcal{N}\mathcal{R}_\omega)\),
the listing is the bottleneck \cite{hansen1988}.

\begin{corollary}[Chain bottleneck identification]
\label{cor:ch4.chain-bottleneck}
The factor achieving the minimum in Eq.~\eqref{eq:ch4.chain-factor-ranks} is the dominant synthesis bottleneck at
tolerance \(\varepsilon\) \cite{harrington2001,devaney2004,stratton1941}. Refining MoM meshes on \(\mathcal{S}_\omega\) alone
does not raise \(r_{\mathrm{rad}}(\boldsymbol{\Gamma}_\omega)\) when \(\mathcal{R}_\omega\) or \(\mathcal{C}\) is the minimizer
\cite{hansen1988}.
\end{corollary}

\paragraph{Worked illustration (listing bottleneck).}
Retaining \(\ell_{\max}=25\) while \(r_{\mathrm{rad}}(\mathcal{C})=12\) caps synthesis regardless of patch refinement:
the listing factor is the minimizer in Eq.~\eqref{eq:ch4.chain-factor-ranks} \cite{hansen1988,harrington2001,devaney2004}.

\paragraph{Problem (which factor to refine?).}
\textbf{Problem.} \(\delta_{\mathrm{syn}}>0\) after MoM convergence. Which factor should be refined?

\textbf{Step 1.} Estimate each term in Eq.~\eqref{eq:ch4.chain-factor-ranks} via SVD probes \cite{harrington2001}.

\textbf{Step 2.} Refine the minimizer factor first \cite{devaney2004,hansen1988}.

\textbf{Step 3.} Recompute \(\delta_{\mathrm{syn}}\); if unchanged, obstruction is geometric rank, not discretization
\cite{stratton1941}.

\textbf{Physical meaning.} Operator chains fail at their narrowest rank factor \cite{harrington2001}.

\paragraph{Validity.}
Eqs.~\eqref{eq:ch4.master-operator-chain}--\eqref{eq:ch4.adjoint-chain} assume linear Maxwell maps and outgoing
\(\mathcal{F}_{\mathrm{rad}}\) \cite{stratton1941}. Mesh refinement changes \(\mathbf{T}\) but not
\(\operatorname{ran}\boldsymbol{\Gamma}_\omega\) in the continuum limit when \(\mathcal{D}_{\mathrm{cpl}}=1\)
\cite{harrington2001}.

Singular values of \(\boldsymbol{\Gamma}_\omega\) define radiative rank in Subsection~\ref{sec:ch434}.


\subsection{Singular values and radiative rank}
\label{sec:ch434}
% TOC tag: [HOME]

Compactness forces singular values to decay; decay defines how many synthesis directions matter at tolerance
\(\varepsilon\) \cite{stratton1941,harrington2001,devaney2004}. Subsection~\ref{sec:ch434} names that count
\emph{radiative rank} and ties it to \(\mathcal{S}_{\mathrm{real}}(\mathcal{C})\) \cite{hansen1988}.

\paragraph{Assumptions.}
Let \(\boldsymbol{\Gamma}_\omega:\mathcal{J}_{\mathrm{real}}\to\ell^{2}(I)\) be compact with singular values
\(\{\sigma_{n}\}_{n\ge1}\) ordered \(\sigma_{1}\ge\sigma_{2}\ge\cdots\to0\) \cite{stratton1941,harrington2001}.
Tolerance \(\varepsilon>0\) is declared for synthesis or inversion \cite{devaney2004}.

\paragraph{Singular value decomposition.}
There exist orthonormal source modes \(\{\mathbf{j}_{n}\}\subset\mathcal{J}_{\mathrm{real}}\) and coefficient modes
\(\{\mathbf{e}_{n}\}\subset\ell^{2}(I)\) such that
\begin{equation}
\label{eq:ch4.gamma-SVD}
\boldsymbol{\Gamma}_\omega
=
\sum_{n=1}^{\infty}\sigma_{n}\,
|\mathbf{j}_{n}\rangle\langle\mathbf{e}_{n}|,
\qquad
\sigma_{n}\to0.
\end{equation}
Equation~\eqref{eq:ch4.gamma-SVD} imports Eq.~\eqref{eq:ch2.singular-value-decay} for
\(K=\boldsymbol{\Gamma}_\omega\) \cite{stratton1941,harrington2001}. Physically, \(\mathbf{j}_{n}\) are
\emph{left} synthesis modes: unit feed along \(\mathbf{j}_{n}\) produces coefficient \(\sigma_{n}\mathbf{e}_{n}\)
\cite{harrington2001,devaney2004}. Truncating at \(\sigma_{N}\) discards synthesis directions below export gain
\(\sigma_{N}\) \cite{hansen1988}.

\begin{definition}[Radiative rank]
\label{def:ch4.radiative-rank}
For tolerance \(\varepsilon>0\), the \emph{radiative rank} of \((\Omega_{s},\omega,\mathcal{C})\) is
\begin{equation}
\label{eq:ch4.radiative-rank-def}
r_{\mathrm{rad}}(\varepsilon)
=
\max\left\{
N:\ \sigma_{N}\big(\boldsymbol{\Gamma}_\omega\big)\ge\varepsilon\,\sigma_{1}\big(\boldsymbol{\Gamma}_\omega\big)
\right\}.
\end{equation}
\end{definition}

Radiative rank is the effective dimension of \(\operatorname{ran}\boldsymbol{\Gamma}_\omega\) at relative tolerance
\(\varepsilon\) \cite{harrington2001,devaney2004}. It is not the same as \(|I_{\ell_{\max}}|\): many listed modes may
have \(\sigma_{n}<\varepsilon\sigma_{1}\) and therefore lie outside the synthesis manifold at the declared tolerance
\cite{hansen1988,stratton1941}. \(r_{\mathrm{rad}}\) is the Chapter~\ref{ch:04} name for the singular-value threshold
anticipated in Subsection~\ref{sec:ch213} \cite{stratton1941}.

\begin{theorem}[Radiative rank bounds synthesis dimension]
\label{thm:ch4.radiative-rank}
For every \(\varepsilon\in(0,1)\),
\begin{equation}
\label{eq:ch4.rank-synthesis-bound}
\dim\Big\{\mathbf{c}\in\mathcal{S}_{\mathrm{real}}(\mathcal{C}):
\|\mathbf{c}\|\ge\varepsilon\|\mathbf{c}_{\max}\|\Big\}
\le
r_{\mathrm{rad}}(\varepsilon),
\end{equation}
with \(\mathbf{c}_{\max}\) a maximizer of \(\|\boldsymbol{\Gamma}_\omega[\Jimp]\|\) over \(\|\Jimp\|=1\) in
\(\mathcal{J}_{\mathrm{real}}\) \cite{stratton1941,harrington2001}. Moreover,
\begin{equation}
\label{eq:ch4.rank-error-bound}
\inf_{\mathrm{rank}\,R\le r_{\mathrm{rad}}(\varepsilon)}\|
\boldsymbol{\Gamma}_\omega-R\|_{\mathrm{op}}
\le
\varepsilon\,\sigma_{1}\big(\boldsymbol{\Gamma}_\omega\big),
\end{equation}
linking radiative rank to best low-rank synthesis approximation \cite{devaney2004,hansen1988}.
\end{theorem}
\begin{proof}
Inequality Eq.~\eqref{eq:ch4.rank-synthesis-bound} follows because significant \(\mathbf{c}\) lie in the span of
\(\{\mathbf{e}_{n}\}_{n\le r_{\mathrm{rad}}(\varepsilon)}\) in Eq.~\eqref{eq:ch4.gamma-SVD}
\cite{stratton1941,harrington2001}. Equation~\eqref{eq:ch4.rank-error-bound} is the Eckart--Young specialization of
Eq.~\eqref{eq:ch2.singular-value-decay} with \(K=\boldsymbol{\Gamma}_\omega\) and \(N=r_{\mathrm{rad}}(\varepsilon)\)
\cite{harrington2001,devaney2004}.
\end{proof}

Theorem~\ref{thm:ch4.radiative-rank} elevates singular values from diagnostic to definitional: \(r_{\mathrm{rad}}\) is
the count of independent realizability directions at tolerance \(\varepsilon\) \cite{hansen1988}. It tightens
Corollary~\ref{cor:ch4.channel-count} by replacing MoM patch count \(P\) with export-weighted rank of the continuum
operator \cite{harrington2001,devaney2004}. Section~\ref{sec:ch46} identifies \(r_{\mathrm{rad}}\) with effective
electromagnetic dimension.

\paragraph{Relation to null space.}
Modes with \(\sigma_{n}=0\) span \(\ker\boldsymbol{\Gamma}_\omega^{\dagger}\): targets with component along those
modes are not realizable \cite{stratton1941,devaney2004}. Combined with \(\ker\mathcal{R}_\omega\),
\begin{equation}
\label{eq:ch4.rank-null-composition}
\ker\boldsymbol{\Gamma}_\omega
\supseteq
\ker\mathcal{R}_\omega\cap\mathcal{J}_{\mathrm{real}},
\end{equation}
so synthesis nullity is at least non-radiating nullity \cite{harrington2001,stratton1941}. Radiative rank counts
\emph{nonzero} singular values above threshold, not kernel dimension alone \cite{devaney2004}.

\paragraph{Worked illustration (tolerance ladder).}
At \(\varepsilon=10^{-2}\), suppose \(\sigma_{1}=1\) and \(\sigma_{15}=10^{-2}\), \(\sigma_{16}=5\times10^{-3}\).
Then \(r_{\mathrm{rad}}(10^{-2})=15\): only fifteen synthesis directions exceed the threshold
\cite{harrington2001}. A sixteen-mode target with equal weight on all modes requires tolerance below
\(5\times10^{-3}\) or enlarged \(\Omega_{s}\) \cite{hansen1988,stratton1941}.

\paragraph{Problem (radiative rank from MoM).}
\textbf{Problem.} SVD of \(\mathbf{T}\) in Eq.~\eqref{eq:ch4.discrete-chain-matrix} yields \(P=200\) singular values.
How many matter at \(\varepsilon=10^{-3}\)?

\textbf{Step 1.} Normalize \(\sigma_{1}=1\); count \(n\) with \(\sigma_{n}\ge10^{-3}\) \cite{harrington2001}.

\textbf{Step 2.} Compare to \(r_{\mathrm{rad}}(10^{-3})\) from Definition~\ref{def:ch4.radiative-rank}; mesh rank
upper-bounds continuum rank \cite{stratton1941,devaney2004}.

\textbf{Step 3.} If count exceeds \(|I_{\ell_{\max}}|\), spectral compression Eq.~\eqref{eq:ch4.spectral-compression-def}
still caps listable modes \cite{hansen1988}.

\textbf{Physical meaning.} Radiative rank is the synthesis DOF budget before accessibility enters
\cite{harrington2001}.

\paragraph{Effective dimension preview.}
Radiative rank is the precursor of \emph{effective electromagnetic dimension} developed in
Subsection~\ref{sec:ch464} \cite{harrington2001,devaney2004}:
\begin{equation}
\label{eq:ch4.effective-dimension-preview}
d_{\mathrm{eff}}(\varepsilon)
=
r_{\mathrm{rad}}(\varepsilon)
=
\#\big\{n:\ \sigma_{n}\ge\varepsilon\,\sigma_{1}\big\},
\end{equation}
counting independent synthesis directions at tolerance \(\varepsilon\) \cite{stratton1941,hansen1988}. Equation
\eqref{eq:ch4.effective-dimension-preview} replaces informal mode-count claims in design reports \cite{devaney2004}.

\paragraph{Worked illustration (lossy tail).}
In lossy media, \(\sigma_{n}\) decay more slowly; \(d_{\mathrm{eff}}(\varepsilon)\) may increase relative to the lossless
case at the same \(\Omega_{s}\) \cite{stratton1941,harrington2001}. Export power still bounds \(\mathbf{c}\) through
Eq.~\eqref{eq:ch4.passivity-screen} \cite{stratton1941}.

\paragraph{Problem (rank versus listing size).}
\textbf{Problem.} \(|I|=500\) but \(d_{\mathrm{eff}}(10^{-2})=18\). How many synthesis directions matter?

\textbf{Step 1.} Use \(d_{\mathrm{eff}}\) as the DOF budget \cite{harrington2001,devaney2004}.

\textbf{Step 2.} Truncate to \(\{\mathbf{e}_{n}\}_{n\le18}\) in Eq.~\eqref{eq:ch4.gamma-SVD} \cite{hansen1988}.

\textbf{Step 3.} Treat excess listing indices as gauge padding \cite{stratton1941}.

\textbf{Physical meaning.} Listing size is not synthesis dimension \cite{hansen1988}.

\paragraph{Discrete--continuum rank gap.}
Mesh refinement increases \(\operatorname{rank}\mathbf{T}\) monotonically but converges to
\(r_{\mathrm{rad}}(\varepsilon)\) only when \(\mathcal{D}_{\mathrm{cpl}}=1\) and patch spacing satisfies sampling limits
\cite{harrington2001,stratton1941,collin2001}. Define the rank gap
\begin{equation}
\label{eq:ch4.rank-gap}
\Delta r_{\mathrm{gap}}
=
\operatorname{rank}\mathbf{T}-r_{\mathrm{rad}}(\varepsilon),
\end{equation}
with \(\Delta r_{\mathrm{gap}}\ge0\) and \(\Delta r_{\mathrm{gap}}\to0\) under standard RWG refinement on compact
\(\Omega_{s}\) \cite{harrington2001,devaney2004}.

\begin{lemma}[Rank gap as aliasing diagnostic]
\label{lem:ch4.rank-gap-aliasing}
If \(\Delta r_{\mathrm{gap}}>0\) at fixed \((\Omega_{s},\omega,\varepsilon)\), then at least
\(\Delta r_{\mathrm{gap}}\) singular directions of \(\mathbf{T}\) have \(\sigma_{n}<\varepsilon\sigma_{1}\) in the
continuum \(\boldsymbol{\Gamma}_\omega\) gauge \cite{stratton1941,hansen1988,harrington2001}. Those directions are
mesh-aliased synthesis artifacts, not export modes \cite{devaney2004}.
\end{lemma}

\paragraph{Worked illustration (rank gap on small sphere).}
A sphere discretized with \(P=80\) patches at \(k_{0}a=1\) yields \(\operatorname{rank}\mathbf{T}=80\) while
\(r_{\mathrm{rad}}(10^{-2})=12\) \cite{harrington2001,hansen1988}. The gap \(\Delta r_{\mathrm{gap}}=68\) identifies
directions that MoM inversion must not treat as independent synthesis coordinates \cite{devaney2004,stratton1941}.

\paragraph{Problem (converged MoM versus continuum rank).}
\textbf{Problem.} \(\operatorname{rank}\mathbf{T}\) plateaus at 150 under mesh refinement, but
\(r_{\mathrm{rad}}(10^{-3})=22\). Has refinement converged?

\textbf{Step 1.} Compare \(\sigma_{n}(\mathbf{T})\) to \(\sigma_{n}(\boldsymbol{\Gamma}_\omega)\) at the same
\(\varepsilon\) \cite{harrington2001,devaney2004}.

\textbf{Step 2.} If \(\Delta r_{\mathrm{gap}}=128\), excess rank is discretization padding \cite{stratton1941,hansen1988}.

\textbf{Step 3.} Truncate synthesis to \(r_{\mathrm{rad}}(10^{-3})\) directions regardless of \(\operatorname{rank}\mathbf{T}\)
\cite{hansen1988}.

\textbf{Physical meaning.} Mesh convergence certifies approximation, not export dimension \cite{harrington2001}.

\paragraph{Validity.}
Definition~\ref{def:ch4.radiative-rank} depends on \(\mathcal{C}\) normalization; relative tolerance \(\varepsilon\)
is gauge-invariant when \(\sigma_{1}\) is taken in the same gauge \cite{hansen1988,stratton1941}. Loss increases
effective rank by reducing \(\sigma_{n}\) tail decay rate \cite{harrington2001}.

Section~\ref{sec:ch43} has identified compact support, spectral compression, the master operator chain, and
radiative rank \(r_{\mathrm{rad}}(\varepsilon)\). Section~\ref{sec:ch44} develops localization and angular bandlimits
that further restrict \(\operatorname{ran}\boldsymbol{\Gamma}_\omega\) in continuous angular representations
\cite{stratton1941,harrington2001,devaney2004}.


\section{Localization and Angular Bandlimits}
\label{sec:ch44}

Radiative rank \(r_{\mathrm{rad}}(\varepsilon)\) counts synthesis directions in a declared basis; compact
support simultaneously localizes fields in space and bandlimits them in lateral wavenumber
\cite{stratton1941,jackson1999,harrington2001}. That duality is not a Fourier heuristic---it is the
spectral image of finite \(\Omega_{s}\) on \(\operatorname{ran}\boldsymbol{\Gamma}_\omega\)
\cite{devaney2004,hansen1988}. Section~\ref{sec:ch44} develops the restriction: Subsection~\ref{sec:ch441}
states spatial localization and Fourier--angular duality; Subsection~\ref{sec:ch442} separates propagating and
evanescent sheets; Subsection~\ref{sec:ch443} records high-\(\ell\) decay; Subsection~\ref{sec:ch444} defines
finite-support spectral restrictions on \(\mathcal{S}_{\mathrm{real}}(\mathcal{C})\) \cite{stratton1941}.


\subsection{Spatial localization and Fourier-angular duality}
\label{sec:ch441}
% TOC tag: [HOME][USE SS3.8.1][USE SS2.9.1]

A source confined to \(\Omega_{s}\) cannot excite arbitrary lateral spectra on a measurement plane: spatial
confinement induces a \emph{spatial--spectral duality} that bandlimits realizable \(\widetilde{\mathbf{A}}(k_{x},k_{y})\)
before angular listing \(\mathbf{c}\) is formed \cite{stratton1941,jackson1999,devaney2004}. Subsection
~\ref{sec:ch441} makes that duality a realizability law on \(\boldsymbol{\Gamma}_\omega\).

\paragraph{Assumptions.}
Fix compact \(\Omega_{s}\), observation plane \(z=z_{0}\) outside \(\Omega_{s}\), and outgoing
\(\mathcal{F}_{\mathrm{rad}}\) \cite{stratton1941}. Lateral Fourier conventions follow
Eqs.~\eqref{eq:ch2.plane-wave-angular-spectrum-def}--\eqref{eq:ch2.plane-wave-z-propagation} from
Subsection~\ref{sec:ch291} \cite{stratton1941,jackson1999}. Flux-normalized plane-wave coefficients follow
Eq.~\eqref{eq:ch3.fourier-radiation-filter-chain} from Subsection~\ref{sec:ch381} \cite{hansen1988,harrington2001}.
Representation \(\mathcal{C}\) may be spherical or plane-wave; duality is stated in \((k_{x},k_{y})\) first
\cite{stratton1941}.

\paragraph{Localization functional.}
Define the spatial second moment of \(\Jimp\) about centroid \(\mathbf{r}_{c}\) by
\begin{equation}
\label{eq:ch4.localization-moment}
\mathcal{L}_{2}[\Jimp]
=
\int_{\Omega_{s}}|\mathbf{r}-\mathbf{r}_{c}|^{2}\,|\Jimp(\mathbf{r})|^{2}\,d^{3}r,
\end{equation}
with \(\mathbf{r}_{c}\) the current centroid \cite{jackson1999,stratton1941}. For ball support \(B_{a}\),
\(\mathcal{L}_{2}\sim a^{2}\) sets the spatial extent that enters spectral tail decay in
Eq.~\eqref{eq:ch4.spectral-tail-energy} \cite{hansen1988}. \(\mathcal{L}_{2}\) is not an independent
realizability screen: it quantifies how rapidly \(\widetilde{\mathbf{A}}(k_{t})\) may vary with transverse
wavenumber \(k_{t}=\sqrt{k_{x}^{2}+k_{y}^{2}}\) \cite{stratton1941,devaney2004}.

\paragraph{Fourier--angular duality (realizability form).}
Let \(\widetilde{\mathbf{A}}_{\mathrm{rad}}(k_{x},k_{y})\) be the radiative lateral spectrum obtained from
Eq.~\eqref{eq:ch3.fourier-radiation-filter-chain} \cite{harrington2001,stratton1941}. Compact support implies
\begin{equation}
\label{eq:ch4.spatial-spectral-duality}
|\widetilde{\mathbf{A}}_{\mathrm{rad}}(k_{x},k_{y})|
\le
\mathcal{M}(k_{0}a)\,\Phi(k_{t}/k_{0}),
\qquad
\Phi(u)\to0\ \text{as}\ u\to\infty,
\end{equation}
with \(\mathcal{M}\) depending on peak \(\|\Jimp\|\) and \(\Phi\) fixed by outgoing Helmholtz symbol
Eq.~\eqref{eq:ch2.helmholtz-symbol} \cite{stratton1941,jackson1999,hansen1988}. Equation
\eqref{eq:ch4.spatial-spectral-duality} is the realizability reading of Fourier--angular duality: there is no
\(\mathbf{c}\in\mathcal{S}_{\mathrm{real}}(\mathcal{C})\) whose plane-wave content violates
Eq.~\eqref{eq:ch4.spatial-spectral-duality} at fixed \((\Omega_{s},k_{0})\) \cite{devaney2004}. High
\(k_{t}\) sheets are suppressed not by measurement alone but by synthesis on compact \(\mathcal{J}_{\mathrm{real}}\)
\cite{stratton1941,harrington2001}.

\begin{definition}[Angular bandlimit (lateral)]
\label{def:ch4.angular-bandlimit}
For tolerance \(\eta>0\), the \emph{lateral angular bandlimit} is
\begin{equation}
\label{eq:ch4.lateral-bandlimit}
k_{t,\max}(\eta)
=
\sup\left\{
k_{t}:\ 
\frac{\int_{k_{t}'<k_{t}}|\widetilde{\mathbf{A}}_{\mathrm{rad}}|^{2}\,dk_{x}\,dk_{y}}
{\int|\widetilde{\mathbf{A}}_{\mathrm{rad}}|^{2}\,dk_{x}\,dk_{y}}
\ge 1-\eta
\right\}.
\end{equation}
\end{definition}

Equation~\eqref{eq:ch4.lateral-bandlimit} names the smallest transverse wavenumber sphere containing
\((1-\eta)\) of radiative spectral weight \cite{hansen1988,stratton1941}. It is the continuous analogue of
\(\ell_{\max}\) in Eq.~\eqref{eq:ch4.spectral-compression-def}: spherical and plane-wave bandlimits are
related by Section~\ref{sec:ch38} transformation laws, not by naive index equality \cite{hansen1988,harrington2001}.

\begin{theorem}[Duality bandlimit on synthesis]
\label{thm:ch4.duality-bandlimit}
For compact \(\Omega_{s}\) with diameter \(D_{s}\) and electrical size \(\Xi\) from
Eq.~\eqref{eq:ch4.electrical-size},
\begin{equation}
\label{eq:ch4.duality-bandlimit-scaling}
k_{t,\max}(\eta)
\le
\kappa(\eta)\,k_{0}\sqrt{1+(2\pi\Xi)^{2}},
\end{equation}
with \(\kappa(\eta)\) depending only on tail budget \(\eta\) and outgoing symbol
\cite{stratton1941,jackson1999,hansen1988}. Consequently \(\boldsymbol{\Gamma}_\omega[\mathcal{J}_{\mathrm{real}}]\) has
negligible content on plane-wave sheets with \(k_{t}\gg k_{t,\max}(\eta)\) at the same \(\eta\) used in
Definition~\ref{def:ch4.radiative-rank} \cite{devaney2004,harrington2001}.
\end{theorem}
\begin{proof}
Spectral decay for compact sources follows from analytic continuation of the dyadic kernel in transverse
wavenumber \cite{stratton1941,jackson1999}. Normalization through \(\mathcal{N}\) and boundedness of
\(\mathcal{C}\) transfer the lateral bound to coefficient norms \cite{hansen1988,harrington2001}. Alignment
with \(r_{\mathrm{rad}}\) follows because both count negligible directions in the same compact operator
\cite{devaney2004,stratton1941}.
\end{proof}

Theorem~\ref{thm:ch4.duality-bandlimit} is the inevitability statement of Subsection~\ref{sec:ch441}: one cannot
synthesize arbitrarily fine transverse structure from a bounded \(\Omega_{s}\) at fixed \(k_{0}\)
\cite{jackson1999,devaney2004}. Phase-only metasurfaces that imprint high \(k_{t}\) on \(\Sigma_{a}\) without
enlarging effective \(\mathcal{L}_{2}\) violate Eq.~\eqref{eq:ch4.duality-bandlimit-scaling} \cite{harrington2001}.

\begin{figure}[t]
\centering
\ChOneFigBlock{%
\begin{tikzpicture}[ch1fig, node distance=7mm and 10mm]
  \node[ch1box] (xs) {Compact $\Omega_{s}$};
  \node[ch1box, right=12mm of xs] (l2) {$\mathcal{L}_{2}$};
  \node[ch1box, right=12mm of l2] (kt) {$k_{t,\max}(\eta)$};
  \node[ch1box, right=12mm of kt] (c) {$\mathbf{c}$};
  \draw[ch1flow] (xs) -- (l2) -- (kt) -- (c);
\end{tikzpicture}%
}
\caption{Subsection~\ref{sec:ch441}: spatial localization controls lateral bandlimit before coefficient listing.}
\label{fig:ch4.spatial-spectral-duality}
\end{figure}

\paragraph{Worked illustration (small aperture).}
For \(k_{0}a=1\), Eq.~\eqref{eq:ch4.duality-bandlimit-scaling} leaves only a neighborhood of \(k_{t}\lesssim k_{0}\)
with appreciable \(\widetilde{\mathbf{A}}_{\mathrm{rad}}\) weight \cite{jackson1999,stratton1941}. Super-oscillatory
targets requiring \(k_{t}\gg k_{0}\) lie outside \(\mathcal{S}_{\mathrm{real}}(\mathcal{C})\) on that aperture
\cite{devaney2004,hansen1988}.

\paragraph{Uncertainty product (localization--spectrum).}
The spatial extent \(\mathcal{L}_{2}\) of Eq.~\eqref{eq:ch4.spatial-spectral-duality} and lateral spectral spread
\(\Delta k_{t}\) satisfy a concentration inequality of the form
\begin{equation}
\label{eq:ch4.localization-spectrum-product}
\mathcal{L}_{2}\,\Delta k_{t}
\ge
\mathcal{K}_{\mathrm{loc}}(k_{0},\eta),
\end{equation}
with \(\mathcal{K}_{\mathrm{loc}}>0\) depending on tail budget \(\eta\) and outgoing symbol
\cite{stratton1941,jackson1999,hansen1988}. Equation~\eqref{eq:ch4.localization-spectrum-product} is the realizability
reading of Fourier--angular duality: shrinking \(\mathcal{L}_{2}\) without widening \(\Delta k_{t}\) forces
\(\mathbf{c}\) outside \(\mathcal{S}_{\mathrm{real}}(\mathcal{C})\) \cite{devaney2004}. The bound is not a measurement
uncertainty relation; it is a synthesis obstruction imposed by compact \(\operatorname{supp}\Jimp\)
\cite{harrington2001}.

\begin{corollary}[Super-oscillation obstruction]
\label{cor:ch4.super-oscillation}
If a target \(\mathbf{c}_{\mathrm{target}}\) requires transverse structure with effective
\(k_{t,\mathrm{req}}\gg k_{t,\max}(\eta)\) from Theorem~\ref{thm:ch4.duality-bandlimit}, then
\(\mathbf{c}_{\mathrm{target}}\notin\mathcal{S}_{\mathrm{real}}(\mathcal{C})\) on the declared \((\Omega_{s},k_{0})\)
\cite{jackson1999,devaney2004,hansen1988}.
\end{corollary}

\paragraph{Problem (metasurface \(k_{t}\) claim).}
\textbf{Problem.} A metasurface of diameter \(D_{s}\) at \(k_{0}D_{s}=2\) claims to synthesize a far-zone pattern
whose plane-wave spectrum peaks at \(k_{t}=2.5k_{0}\). Is the claim consistent with Theorem~\ref{thm:ch4.duality-bandlimit}?

\textbf{Step 1.} Compute \(k_{t,\max}(\eta)\) from Eq.~\eqref{eq:ch4.duality-bandlimit-scaling} with \(\Xi=k_{0}D_{s}/(2\pi)\)
\cite{stratton1941,jackson1999}.

\textbf{Step 2.} Compare \(k_{t,\mathrm{req}}=2.5k_{0}\) to \(k_{t,\max}(\eta)\); excess weight lies in
\(\mathcal{K}_{\mathrm{ev}}\) near the surface, not in \(\mathcal{K}_{\mathrm{prop}}\) of
Eq.~\eqref{eq:ch4.propagating-realizability} \cite{harrington2001}.

\textbf{Step 3.} Conclude \(\mathbf{c}_{\mathrm{target}}\notin\mathcal{S}_{\mathrm{real}}(\mathcal{C})\) unless
\(D_{s}\) or \(\omega\) is increased \cite{devaney2004,hansen1988}.

\textbf{Physical meaning.} Lateral bandlimits are geometry-imposed before phase engineering enters
\cite{stratton1941}.

\paragraph{Plane-wave transform on synthesis fields.}
Let \(\Phi[\mathfrak{F}](k_{x},k_{y})\) denote the lateral spectrum of \(\mathfrak{F}\in\mathcal{F}_{\mathrm{rad}}\) on a
plane \(z=z_{0}\) \cite{stratton1941,jackson1999}. Compact \(\Omega_{s}\) imposes
\begin{equation}
\label{eq:ch4.plane-wave-support-restriction}
\operatorname{supp}\Phi\big[\mathcal{R}_\omega[\Jimp]\big]
\subseteq
\big\{(k_{x},k_{y}):k_{t}\le k_{t,\max}(\eta)\big\},
\end{equation}
with \(k_{t,\max}\) from Theorem~\ref{thm:ch4.duality-bandlimit} \cite{hansen1988,devaney2004}. Equation
\eqref{eq:ch4.plane-wave-support-restriction} is the lateral restatement of
\(\mathcal{B}_{\mathrm{fs}}\): synthesis targets must concentrate in the same spectral support as
\(\boldsymbol{\Gamma}_\omega[\mathcal{J}_{\mathrm{real}}]\) \cite{stratton1941,harrington2001}.

\begin{theorem}[Bandlimit inheritance under listing]
\label{thm:ch4.bandlimit-listing-inheritance}
If Eq.~\eqref{eq:ch4.plane-wave-support-restriction} holds and \(\mathcal{C}\) is a bounded transformation from
\(\Phi\) to \(\mathbf{c}\) as in Section~\ref{sec:ch38}, then every realizable \(\mathbf{c}\) inherits the same
\(k_{t,\max}(\eta)\) bound up to constants from basis change \cite{hansen1988,stratton1941,devaney2004}.
\end{theorem}
\begin{proof}
Boundedness of the plane-wave--VSH transform on compact domains
\cite{hansen1988,stratton1941}; Definition~\ref{def:ch4.angular-bandlimit} \cite{devaney2004}.
\end{proof}

\paragraph{Worked illustration (narrow beam on small disk).}
A disk aperture of radius \(a\) at \(k_{0}a=2\) cannot sustain a far-zone beam requiring
\(k_{t,\mathrm{req}}\approx2k_{0}\) by Corollary~\ref{cor:ch4.super-oscillation} \cite{jackson1999,hansen1988}.
Beamwidth claims must be checked against Eq.~\eqref{eq:ch4.plane-wave-support-restriction} before \(\mathbf{c}\) is accepted
\cite{devaney2004,harrington2001}.

\paragraph{Problem (bandlimit after basis change).}
\textbf{Problem.} \(\mathbf{c}\) is sparse in VSH but dense in plane-wave gauge. Which bandlimit applies?

\textbf{Step 1.} Apply Theorem~\ref{thm:ch4.bandlimit-listing-inheritance} \cite{hansen1988}.

\textbf{Step 2.} Transform \(\mathbf{c}\) to plane-wave coefficients and test \(k_{t}(\mathbf{c})\) in
Eq.~\eqref{eq:ch4.finite-support-restriction} \cite{stratton1941,devaney2004}.

\textbf{Step 3.} Sparsity in one gauge does not imply realizability if the other gauge violates
\(k_{t,\max}(\eta)\) \cite{harrington2001}.

\textbf{Physical meaning.} Bandlimits attach to \(\mathfrak{F}\), not to coordinate sparsity \cite{hansen1988}.

\paragraph{Beamwidth--bandlimit closure.}
For a far-zone beam with half-power width \(\theta_{\mathrm{HP}}\), realizability on \(B_{a}\) requires
\begin{equation}
\label{eq:ch4.beamwidth-bandlimit}
\theta_{\mathrm{HP}}
\gtrsim
\frac{\lambda}{\pi D_{s}}
\quad\Longrightarrow\quad
k_{t,\mathrm{req}}\lesssim k_{t,\max}(\eta),
\end{equation}
linking angular beam metrics to Theorem~\ref{thm:ch4.duality-bandlimit} \cite{jackson1999,hansen1988,stratton1941}.
Super-narrow beams on electrically small apertures violate Eq.~\eqref{eq:ch4.beamwidth-bandlimit} before
\(\ell_{\max}\) is consulted \cite{devaney2004,harrington2001}.

\begin{corollary}[Minimum aperture for target beamwidth]
\label{cor:ch4.min-aperture-beamwidth}
A target half-power width \(\theta_{\mathrm{target}}\) at wavelength \(\lambda\) requires
\(D_{s}\gtrsim\lambda/(\pi\theta_{\mathrm{target}})\) for \(\mathbf{c}_{\mathrm{target}}\in\mathcal{S}_{\mathrm{real}}(\mathcal{C})\)
up to constants from \(\eta\) \cite{jackson1999,hansen1988,devaney2004}.
\end{corollary}

\paragraph{Worked illustration (pencil beam on handset).}
A \(40^{\circ}\) handset aperture at \(28\,\mathrm{GHz}\) cannot realize a \(2^{\circ}\) pencil beam: Eq.~\eqref{eq:ch4.beamwidth-bandlimit}
fails by an order of magnitude \cite{harrington2001,jackson1999,hansen1988}.

\paragraph{Problem (beamwidth claim audit).}
\textbf{Problem.} Marketing claims \(\theta_{\mathrm{HP}}=3^{\circ}\) from a \(0.2\lambda\) aperture. Audit realizability.

\textbf{Step 1.} Compute \(D_{s}=0.2\lambda\) and test Corollary~\ref{cor:ch4.min-aperture-beamwidth} \cite{jackson1999}.

\textbf{Step 2.} Map \(\theta_{\mathrm{target}}\) to \(k_{t,\mathrm{req}}\) \cite{hansen1988,stratton1941}.

\textbf{Step 3.} Compare to \(k_{t,\max}(\eta)\) from Theorem~\ref{thm:ch4.duality-bandlimit} \cite{devaney2004}.

\textbf{Physical meaning.} Beamwidth is a bandlimit consequence, not an independent design knob \cite{stratton1941}.

\paragraph{Validity.}
Eqs.~\eqref{eq:ch4.spatial-spectral-duality}--\eqref{eq:ch4.duality-bandlimit-scaling} assume homogeneous
exterior and linear Maxwell theory \cite{stratton1941}. Anisotropic media modify \(\Phi\) without removing
bandlimit \cite{harrington2001}.

Propagating versus evanescent partition is Subsection~\ref{sec:ch442}.


\subsection{Evanescent and propagating components}
\label{sec:ch442}
% TOC tag: [USE]

Lateral spectra decompose into propagating sheets \(k_{t}\le k_{0}\) and evanescent sheets \(k_{t}>k_{0}\) under
Eq.~\eqref{eq:ch2.longitudinal-wavenumber-split} \cite{stratton1941}. Only the propagating sector carries
\(\mathfrak{F}\in\mathcal{F}_{\mathrm{rad}}\); evanescent weight may be large near \(\Omega_{s}\) yet invisible to
\(\boldsymbol{\Gamma}_\omega\) \cite{harrington2001,jackson1999}.

\paragraph{Assumptions.}
Presuppose Eq.~\eqref{eq:ch2.rad-projection-composition},
\(\mathcal{R}_\omega=\Pi_{\mathrm{rad}}\circ\Pi_{\mathrm{prop}}\circ\mathcal{S}_\omega\), from
Subsection~\ref{sec:ch224} \cite{stratton1941,harrington2001}. Longitudinal wavenumber
\(k_{z}(k_{t})\) follows Eq.~\eqref{eq:ch2.outgoing-plane-wave-gate} on propagating sheets
\cite{stratton1941,jackson1999}.

\paragraph{Propagating realizability sector.}
Define
\begin{equation}
\label{eq:ch4.propagating-sector}
\mathcal{K}_{\mathrm{prop}}
=
\{(k_{x},k_{y}):k_{t}\le k_{0},\ k_{z}\ \text{outgoing}\},
\end{equation}
importing the propagating sheet of Subsection~\ref{sec:ch224} \cite{stratton1941}. Realizable radiative spectra
satisfy
\begin{equation}
\label{eq:ch4.propagating-realizability}
\operatorname{supp}\widetilde{\mathbf{A}}_{\mathrm{rad}}
\subseteq
\mathcal{K}_{\mathrm{prop}}
\quad\text{up to null measure},
\end{equation}
because \(\Pi_{\mathrm{prop}}\) annihilates evanescent content before \(\Pi_{\mathrm{rad}}\)
\cite{stratton1941,harrington2001}. Equation~\eqref{eq:ch4.propagating-realizability} is the spectral form of
Principle~\ref{prin:ch4.operator-chain}: reactive near fields may carry large \(k_{t}>k_{0}\) weight in raw
\(\mathcal{S}_\omega[\Jimp]\) but not in \(\mathbf{c}\) \cite{harrington2001,devaney2004}.

\paragraph{Evanescent storage sector.}
Let
\begin{equation}
\label{eq:ch4.evanescent-sector}
\mathcal{K}_{\mathrm{ev}}
=
\{(k_{x},k_{y}):k_{t}>k_{0}\},
\qquad
\widetilde{\mathbf{A}}_{\mathrm{ev}}
=
\widetilde{\mathbf{A}}\cdot\mathbf{1}_{\mathcal{K}_{\mathrm{ev}}}.
\end{equation}
Equation~\eqref{eq:ch4.evanescent-sector} isolates reactive lateral content \cite{stratton1941,jackson1999}. Near-field
probes sensitive to \(\widetilde{\mathbf{A}}_{\mathrm{ev}}\) report structure that does not enter
\(\mathcal{S}_{\mathrm{real}}(\mathcal{C})\) on \(\mathcal{F}_{\mathrm{rad}}\) \cite{harrington2001}. Confusing
\(\widetilde{\mathbf{A}}_{\mathrm{ev}}\) with synthesizable \(\widetilde{\mathbf{A}}_{\mathrm{rad}}\) is a primary
realizability failure mode \cite{devaney2004}.

\begin{proposition}[Evanescent invisibility to synthesis]
\label{prop:ch4.evanescent-invisibility}
For every \(\Jimp\in\mathcal{J}_{\mathrm{real}}\),
\begin{equation}
\label{eq:ch4.evanescent-gamma-null}
\boldsymbol{\Gamma}_\omega[\Jimp]
\ \text{independent of}\ 
\widetilde{\mathbf{A}}_{\mathrm{ev}}
\quad\text{on}\ \mathcal{F}_{\mathrm{rad}},
\end{equation}
equivalently, \(\Pi_{\mathrm{rad}}\Pi_{\mathrm{prop}}\) annihilates \(\mathcal{K}_{\mathrm{ev}}\) content before
\(\mathcal{C}\) \cite{stratton1941,harrington2001}.
\end{proposition}
\begin{proof}
Eq.~\eqref{eq:ch2.rad-projection-composition} and evanescent null
Eq.~\eqref{eq:ch2.rad-null-on-evanescent} \cite{stratton1941,harrington2001}.
\end{proof}

Proposition~\ref{prop:ch4.evanescent-invisibility} forbids fitting far-zone \(\mathbf{c}\) with evanescent-only
near-field templates: such templates are certification objects for reactive storage, not synthesis targets
\cite{stratton1941,devaney2004}. Super-resolution claims based solely on \(k_{t}>k_{0}\) near-field energy must
separate \(\mathcal{K}_{\mathrm{ev}}\) from \(\mathcal{K}_{\mathrm{prop}}\) \cite{harrington2001}.

\paragraph{Worked illustration (near-field scan).}
A probe at \(z=z_{0}\) near \(\Omega_{s}\) measures large evanescent amplitudes. Far-zone \(\mathbf{c}\) reconstructed
from those scans without \(\Pi_{\mathrm{prop}}\) violates Eq.~\eqref{eq:ch4.propagating-realizability}
\cite{harrington2001,jackson1999}.

\paragraph{Plane-wave energy partition.}
For a lateral spectral density \(|\widetilde{\mathbf{A}}(k_{x},k_{y})|^{2}\), define propagating and evanescent
energy fractions
\begin{equation}
\label{eq:ch4.prop-ev-energy-fractions}
\mathcal{E}_{\mathrm{prop}}
=
\int_{\mathcal{K}_{\mathrm{prop}}}\!
|\widetilde{\mathbf{A}}_{\mathrm{rad}}|^{2}\,d^{2}k_{t},
\qquad
\mathcal{E}_{\mathrm{ev}}
=
\int_{\mathcal{K}_{\mathrm{ev}}}\!
|\widetilde{\mathbf{A}}|^{2}\,d^{2}k_{t},
\end{equation}
with \(\mathcal{K}_{\mathrm{prop}},\mathcal{K}_{\mathrm{ev}}\) from Eqs.~\eqref{eq:ch4.propagating-sector} and
\eqref{eq:ch4.evanescent-sector} \cite{stratton1941,jackson1999}. Near \(\Omega_{s}\), \(\mathcal{E}_{\mathrm{ev}}\)
may dominate while \(\mathcal{E}_{\mathrm{prop}}\) alone determines \(\|\mathbf{c}\|_{\mathrm{rad}}\) after
\(\Pi_{\mathrm{rad}}\Pi_{\mathrm{prop}}\) \cite{harrington2001}. Equation~\eqref{eq:ch4.prop-ev-energy-fractions}
is the spectral diagnostic for Proposition~\ref{prop:ch4.evanescent-invisibility}: large \(\mathcal{E}_{\mathrm{ev}}\)
does not imply large synthesizable \(\mathbf{c}\) \cite{devaney2004}.

\begin{figure}[t]
\centering
\ChOneFigBlock{%
\begin{tikzpicture}[ch1fig, node distance=7mm and 9mm]
  \node[ch1box, minimum width=32mm] (raw) {Raw $\mathcal{S}_\omega[\Jimp]$\\large $\mathcal{E}_{\mathrm{ev}}$};
  \node[ch1box, minimum width=30mm, right=12mm of raw] (prop) {$\Pi_{\mathrm{prop}}$\\$\mathcal{K}_{\mathrm{prop}}$};
  \node[ch1box, minimum width=28mm, right=12mm of prop] (rad) {$\Pi_{\mathrm{rad}}$\\$\mathcal{F}_{\mathrm{rad}}$};
  \node[ch1box, minimum width=26mm, below=12mm of prop] (null) {$\mathcal{K}_{\mathrm{ev}}\mapsto\mathbf{0}$};
  \draw[ch1flow] (raw) -- (prop) -- (rad);
  \draw[ch1flow] (null) -- (prop);
\end{tikzpicture}%
}
\caption{Subsection~\ref{sec:ch442}: evanescent lateral spectra are removed by \(\Pi_{\mathrm{prop}}\) before
radiative listing; \(\boldsymbol{\Gamma}_\omega\) sees only \(\mathcal{K}_{\mathrm{prop}}\).}
\label{fig:ch4.prop-ev-partition}
\end{figure}

Figure~\ref{fig:ch4.prop-ev-partition} is the lateral reading of Principle~\ref{prin:ch4.operator-chain}: synthesis
optimizes \(\mathbf{c}\) on \(\mathcal{F}_{\mathrm{rad}}\), not on the full near-field spectrum
\cite{stratton1941,harrington2001}. Super-directivity designs that concentrate near-field evanescent energy may leave
\(\mathbf{c}\) nearly unchanged while \(\mathcal{E}_{\mathrm{ev}}\) grows \cite{devaney2004,jackson1999}.

\begin{corollary}[Reactive storage without radiative signature]
\label{cor:ch4.reactive-no-c}
Let \(\Jimp_{A},\Jimp_{B}\in\mathcal{J}_{\mathrm{real}}\) satisfy
\(\boldsymbol{\Gamma}_\omega[\Jimp_{A}]=\boldsymbol{\Gamma}_\omega[\Jimp_{B}]\) but differ in
\(\mathcal{E}_{\mathrm{ev}}\) near \(\Omega_{s}\). Then \(\mathbf{c}_{A}=\mathbf{c}_{B}\) while near-field probes
distinguish \(\Jimp_{A}\) from \(\Jimp_{B}\) \cite{stratton1941,harrington2001,devaney2004}.
\end{corollary}

\paragraph{Problem (evanescent template fit).}
\textbf{Problem.} A near-field scanner fits \(\widetilde{\mathbf{A}}(k_{x},k_{y})\) with substantial weight on
\(k_{t}=1.5k_{0}\) and claims a far-zone \(\mathbf{c}_{\mathrm{target}}\) with narrow main beam. Is the claim
realizable?

\textbf{Step 1.} Project the fitted spectrum onto \(\mathcal{K}_{\mathrm{prop}}\); evanescent weight does not enter
\(\boldsymbol{\Gamma}_\omega\) by Proposition~\ref{prop:ch4.evanescent-invisibility} \cite{stratton1941}.

\textbf{Step 2.} Recompute \(\mathbf{c}\) from \(\widetilde{\mathbf{A}}_{\mathrm{rad}}\) only; narrow beams require
propagating \(k_{t}\) structure compatible with Eq.~\eqref{eq:ch4.duality-bandlimit-scaling}
\cite{hansen1988,jackson1999}.

\textbf{Step 3.} If \(\mathbf{r}_{\min}\neq\mathbf{0}\) in Principle~\ref{prin:ch4.operator-chain}, enlarge aperture or
revise the target \(\mathbf{c}\) \cite{devaney2004,harrington2001}.

\textbf{Physical meaning.} Evanescent spectra certify reactive storage, not far-zone synthesis degrees of freedom
\cite{stratton1941}.

\paragraph{Superdirectivity and reactive storage.}
Superdirective templates concentrate field energy in evanescent sectors \(\mathcal{K}_{\mathrm{ev}}\) while leaving
\(\mathcal{A}_{\infty}\) nearly unchanged \cite{stratton1941,harrington2001,jackson1999}. The realizability reading is
sharp: \(\|\widetilde{\mathbf{A}}_{\mathrm{ev}}\|\) may grow without enlarging \(\mathcal{S}_{\mathrm{real}}(\mathcal{C})\)
Proposition~\ref{prop:ch4.evanescent-invisibility} \cite{devaney2004}. Passivity caps total stored energy through
Eq.~\eqref{eq:ch4.passivity-screen} even when \(\mathcal{E}_{\mathrm{ev}}\) is large \cite{stratton1941}.

\begin{equation}
\label{eq:ch4.superdirectivity-obstruction}
\frac{\mathcal{E}_{\mathrm{ev}}}{\mathcal{E}_{\mathrm{prop}}}
\gg1
\quad\not\Longrightarrow\quad
\dim\mathcal{S}_{\mathrm{real}}(\mathcal{C})\ \text{enlarged},
\end{equation}
because \(\boldsymbol{\Gamma}_\omega\) factors through \(\Pi_{\mathrm{prop}}\) \cite{harrington2001,devaney2004,stratton1941}.
Equation~\eqref{eq:ch4.superdirectivity-obstruction} forbids interpreting reactive enhancement as additional far-zone
synthesis rank \cite{hansen1988}.

\paragraph{Worked illustration (near-field enhancement).}
A resonant particle near \(\Omega_{s}\) raises \(\mathcal{E}_{\mathrm{ev}}\) in Eq.~\eqref{eq:ch4.prop-ev-energy-fractions}
without changing \(\mathbf{c}\) when \(\Pi_{\mathrm{rad}}\) is applied \cite{stratton1941,jackson1999,harrington2001}.
Far-zone OAM reports must be computed from \(\mathcal{K}_{\mathrm{prop}}\) alone \cite{devaney2004}.

\paragraph{Problem (evanescent OAM claim).}
\textbf{Problem.} A near-field probe reports strong \(k_{t}>k_{0}\) OAM weight. Can \(\mathbf{c}\) carry the same OAM?

\textbf{Step 1.} Project onto \(\mathcal{K}_{\mathrm{prop}}\) \cite{stratton1941}.

\textbf{Step 2.} List \(\mathbf{c}\) from projected spectrum only \cite{hansen1988,harrington2001}.

\textbf{Step 3.} If projected weight is negligible, OAM is reactive certification, not synthesis content
\cite{devaney2004}.

\textbf{Physical meaning.} Evanescent OAM is not far-zone OAM \cite{stratton1941}.

\paragraph{Angular-spectrum realizability partition.}
Decompose the lateral spectrum of \(\mathcal{S}_\omega[\Jimp]\) into propagating and evanescent partitions
Eqs.~\eqref{eq:ch4.prop-ev-energy-fractions}. The synthesis partition is
\begin{equation}
\label{eq:ch4.synthesis-spectrum-partition}
\widetilde{\mathbf{A}}_{\mathrm{syn}}
=
\Pi_{\mathrm{prop}}\Pi_{\mathrm{rad}}\,\widetilde{\mathbf{A}},
\qquad
\|\widetilde{\mathbf{A}}_{\mathrm{syn}}\|^{2}
=
\mathcal{E}_{\mathrm{prop}}^{\mathrm{(rad)}},
\end{equation}
with \(\mathcal{E}_{\mathrm{prop}}^{\mathrm{(rad)}}\) the radiative propagating energy \cite{stratton1941,harrington2001}.
Every \(\mathbf{c}\in\mathcal{S}_{\mathrm{real}}(\mathcal{C})\) is determined by \(\widetilde{\mathbf{A}}_{\mathrm{syn}}\) alone;
reactive sheets \(\mathcal{K}_{\mathrm{ev}}\) are certification-only \cite{devaney2004,jackson1999}.

\begin{theorem}[Propagating-sector synthesis equivalence]
\label{thm:ch4.prop-sector-equivalence}
\(\mathbf{c}_{1},\mathbf{c}_{2}\in\mathcal{S}_{\mathrm{real}}(\mathcal{C})\) satisfy \(\mathbf{c}_{1}=\mathbf{c}_{2}\) if and
only if \(\widetilde{\mathbf{A}}_{\mathrm{syn}}^{(1)}=\widetilde{\mathbf{A}}_{\mathrm{syn}}^{(2)}\) on
\(\mathcal{K}_{\mathrm{prop}}\) up to gauge \cite{stratton1941,harrington2001,devaney2004}. Distinct near fields with
identical \(\widetilde{\mathbf{A}}_{\mathrm{syn}}\) differ only in \(\mathcal{K}_{\mathrm{ev}}\) storage
\cite{harrington2001}.
\end{theorem}
\begin{proof}
Proposition~\ref{prop:ch4.evanescent-invisibility} and boundedness of \(\mathcal{C}\) on
\(\mathcal{F}_{\mathrm{rad}}\) \cite{stratton1941,hansen1988,harrington2001}.
\end{proof}

\paragraph{Worked illustration (metasurface near-field hotspot).}
A metasurface may create \(\mathcal{E}_{\mathrm{ev}}\gg\mathcal{E}_{\mathrm{prop}}\) while leaving
\(\mathbf{c}\) nearly unchanged \cite{stratton1941,harrington2001}. Far-zone pattern claims must be audited on
\(\widetilde{\mathbf{A}}_{\mathrm{syn}}\), not on raw probe spectra \cite{devaney2004,jackson1999}.

\paragraph{Problem (reactive-to-radiative promotion).}
\textbf{Problem.} Post-processing promotes \(k_{t}>k_{0}\) spectral weight into \(\mathcal{K}_{\mathrm{prop}}\) by phase
unwrapping. Does \(\mathbf{c}\) change?

\textbf{Step 1.} Test whether promotion alters \(\Pi_{\mathrm{prop}}\Pi_{\mathrm{rad}}\) output \cite{stratton1941}.

\textbf{Step 2.} If \(\widetilde{\mathbf{A}}_{\mathrm{syn}}\) is unchanged, \(\mathbf{c}\) is unchanged
\cite{harrington2001,devaney2004}.

\textbf{Step 3.} Reject promotion as synthesis unless \(\mathcal{D}_{\mathrm{real}}=1\) on a feed that actually exports the
promoted sheet \cite{hansen1988}.

\textbf{Physical meaning.} Evanescent weight cannot be post-promoted into radiative synthesis rank \cite{stratton1941}.

\paragraph{Validity.}
Eqs.~\eqref{eq:ch4.propagating-sector}--\eqref{eq:ch4.evanescent-gamma-null} assume linear isotropic exterior media
\cite{stratton1941}. Loss modifies \(k_{z}\) sheets without promoting evanescent content to \(\mathcal{F}_{\mathrm{rad}}\)
\cite{harrington2001}.

High-\(\ell\) decay laws are Subsection~\ref{sec:ch443}.


\subsection{Decay of high-order angular content}
\label{sec:ch443}
% TOC tag: [HOME][USE SS3.6.2]

In spherical listings, compact support compresses tail energy through Eq.~\eqref{eq:ch3.truncation-remainder-energy}
from Subsection~\ref{sec:ch362} \cite{hansen1988,stratton1941}. Subsection~\ref{sec:ch443} states the decay law as a
realizability bound on \(\mathcal{S}_{\mathrm{real}}(\mathcal{C})\) in \(\ell\)-indexed coordinates.

\paragraph{Assumptions.}
Fix VSH listing \(\mathcal{C}\) from Section~\ref{sec:ch34}, ball support \(B_{a}\), and normalized coefficients
\(\widehat{c}_{\ell m}^{(\sigma)}\) from Section~\ref{sec:ch35} \cite{stratton1941,hansen1988}. Tail budget
\(\eta_{\mathrm{tail}}\) is as in Eq.~\eqref{eq:ch4.spectral-compression-def} \cite{harrington2001}.

\paragraph{High-\(\ell\) tail bound.}
For \(\ell>k_{0}a\),
\begin{equation}
\label{eq:ch4.high-ell-decay}
|\widehat{c}_{\ell m}^{(\sigma)}|
\le
\mathcal{C}_{\ell}(k_{0}a)\,\|\mathbf{c}\|_{\mathrm{rad}}\,
\ell^{-1}\exp(-\ell/k_{0}a),
\end{equation}
with \(\mathcal{C}_{\ell}\) depending on polarization sector \(\sigma\) and normalization gauge
\cite{stratton1941,jackson1999,hansen1988}. Equation~\eqref{eq:ch4.high-ell-decay} imports the engineering tail law of
Subsection~\ref{sec:ch362} as a realizability inequality: coefficients above the envelope cannot be synthesized on
\(B_{a}\) at fixed \(k_{0}\) \cite{harrington2001,devaney2004}. The bound aligns with
Eq.~\eqref{eq:ch4.compression-error-import} when \(\{\mathbf{u}_{n}\}\) are VSH modes
\cite{hansen1988,stratton1941}.

\begin{proposition}[Monotone tail in \(\ell_{\max}\)]
\label{prop:ch4.ell-tail-monotone}
Let \(\mathbf{R}_{\ell_{\max}}\) be the truncation remainder Eq.~\eqref{eq:ch3.truncation-remainder-energy}. Then
\(\|\mathbf{R}_{\ell_{\max}}\|_{\mathrm{rad}}\) decreases monotonically in \(\ell_{\max}\), and
\begin{equation}
\label{eq:ch4.ellmax-tail-link}
\|\mathbf{R}_{\ell_{\max}}\|_{\mathrm{rad}}^{2}
\le
\sum_{\ell>\ell_{\max}}(2\ell+1)\max_{m,\sigma}|\widehat{c}_{\ell m}^{(\sigma)}|^{2},
\end{equation}
with equality in the aligned VSH gauge \cite{hansen1988,stratton1941,harrington2001}.
\end{proposition}
\begin{proof}
Orthogonality Theorem~\ref{thm:ch3.vsh-orthogonality} and Parseval
Eq.~\eqref{eq:ch3.mode-power-from-coefficient} \cite{hansen1988,stratton1941}.
\end{proof}

Proposition~\ref{prop:ch4.ell-tail-monotone} links spherical bandlimits to watt residuals: raising \(\ell_{\max}\)
reduces synthesis error monotonically but cannot create content forbidden by Eq.~\eqref{eq:ch4.high-ell-decay}
\cite{harrington2001,devaney2004}. Targets with plateau tails at \(\ell\gg k_{0}a\) lie outside
\(\mathcal{S}_{\mathrm{real}}(\mathcal{C})\) on \(B_{a}\) \cite{jackson1999,hansen1988}.

\paragraph{Worked illustration (quadrupole on small sphere).}
At \(k_{0}a=0.8\), Eq.~\eqref{eq:ch4.high-ell-decay} suppresses \(\ell=2\) below \(\ell=1\) by roughly one decade;
demanding dominant \(\ell=4\) content fails unless \(a\) is enlarged \cite{stratton1941,jackson1999,hansen1988}.

\paragraph{Tail energy summation.}
Summing Eq.~\eqref{eq:ch4.high-ell-decay} over \((m,\sigma)\) at fixed \(\ell\) gives the partial mode-block energy
\begin{equation}
\label{eq:ch4.ell-block-energy}
\mathcal{E}_{\ell}
=
\sum_{m,\sigma}|\widehat{c}_{\ell m}^{(\sigma)}|^{2}
\le
(2\ell+1)\,\mathcal{C}_{\ell}^{2}(k_{0}a)\,\|\mathbf{c}\|_{\mathrm{rad}}^{2}\,
\ell^{-2}\exp(-2\ell/k_{0}a),
\end{equation}
with the \((2\ell+1)\) degeneracy factor of spherical harmonics \cite{hansen1988,stratton1941}. Equation
\eqref{eq:ch4.ell-block-energy} converts pointwise tail bounds into cumulative watt residuals: the tail projector
\(\mathcal{T}_{\ell_{\max}}\) of Eq.~\eqref{eq:ch3.ellmax-truncation-projector} removes blocks \(\mathcal{E}_{\ell}\)
for \(\ell>\ell_{\max}\) with controlled error \cite{hansen1988}. Targets whose \(\mathcal{E}_{\ell}\) plateau above
Eq.~\eqref{eq:ch4.ell-block-energy} for \(\ell\gg k_{0}a\) are outside \(\mathcal{B}_{\mathrm{fs}}\) of
Subsection~\ref{sec:ch444} \cite{devaney2004,harrington2001}.

\begin{figure}[t]
\centering
\ChOneFigBlock{%
\begin{tikzpicture}[ch1fig, node distance=7mm and 8mm]
  \node[ch1box, minimum width=28mm] (ell) {High $\ell$\\tail envelope};
  \node[ch1box, minimum width=30mm, right=12mm of ell] (tmax) {$\mathcal{T}_{\ell_{\max}}$\\truncation};
  \node[ch1box, minimum width=28mm, right=12mm of tmax] (sr) {$\mathcal{S}_{\mathrm{real}}$\\compressed};
  \draw[ch1flow] (ell) -- (tmax) -- (sr);
\end{tikzpicture}%
}
\caption{Subsection~\ref{sec:ch443}: Eq.~\eqref{eq:ch4.high-ell-decay} bounds tail blocks before truncation;
\(\ell_{\max}\) selects the admissible synthesis body inside \(\mathcal{S}_{\mathrm{real}}(\mathcal{C})\).}
\label{fig:ch4.high-ell-tail}
\end{figure}

Figure~\ref{fig:ch4.high-ell-tail} connects spherical listing to operator rank: when \(\{\mathbf{u}_{n}\}\) align with
VSH modes, \(\ell_{\max}\) and \(r_{\mathrm{rad}}(\varepsilon)\) count the same negligible directions up to gauge
\cite{hansen1988,stratton1941}. Proposition~\ref{prop:ch4.geometry-spectral-equivalence} is the bridge between
Eq.~\eqref{eq:ch4.ell-block-energy} and Eq.~\eqref{eq:ch4.compression-error-import} \cite{harrington2001}.

\begin{corollary}[Tail plateau obstruction]
\label{cor:ch4.tail-plateau}
If \(|\widehat{c}_{\ell m}^{(\sigma)}|\ge c_{0}\|\mathbf{c}\|_{\mathrm{rad}}\) for all \(\ell\in[\ell_{1},\ell_{2}]\) with
\(\ell_{2}-\ell_{1}\gg k_{0}a\) and \(c_{0}>0\) independent of \(\ell\), then
\(\mathbf{c}\notin\mathcal{B}_{\mathrm{fs}}(\eta,\varepsilon)\) for any \(\eta_{\mathrm{tail}}\) satisfying
Eq.~\eqref{eq:ch4.spectral-compression-def} on \(B_{a}\) at fixed \(k_{0}\) \cite{hansen1988,devaney2004}.
\end{corollary}
\begin{proof}
Plateau tails violate Eq.~\eqref{eq:ch4.high-ell-decay} and inflate
\(\|\mathbf{R}_{\ell_{\max}}\|_{\mathrm{rad}}\) in Proposition~\ref{prop:ch4.ell-tail-monotone}
\cite{stratton1941,hansen1988}.
\end{proof}

\paragraph{Problem (OAM plateau test).}
\textbf{Problem.} A design targets equal-weight VSH coefficients through \(\ell=10\) on \(B_{a}\) with \(k_{0}a=1.5\).
Does \(\mathbf{c}_{\mathrm{target}}\) lie in \(\mathcal{S}_{\mathrm{real}}(\mathcal{C})\)?

\textbf{Step 1.} Evaluate Eq.~\eqref{eq:ch4.ell-block-energy} at \(\ell=10\): exponential factor
\(\exp(-2\ell/k_{0}a)\) is \(\ll 1\) \cite{jackson1999,hansen1988}.

\textbf{Step 2.} Equal-weight plateau through \(\ell=10\) violates Corollary~\ref{cor:ch4.tail-plateau}
\cite{devaney2004}.

\textbf{Step 3.} Enlarge \(a\) or reduce the \(\ell\) span of the target; phase-only retuning on fixed support cannot
restore the plateau \cite{stratton1941,harrington2001}.

\textbf{Physical meaning.} High-\(\ell\) OAM content on compact supports is tail-bounded, not palette-unlimited
\cite{hansen1988}.

\paragraph{Cumulative tail mass.}
Summing Eq.~\eqref{eq:ch4.ell-block-energy} over \(\ell>\ell_{\max}\) defines
\begin{equation}
\label{eq:ch4.cumulative-tail-mass}
\mathcal{M}_{\mathrm{tail}}(\ell_{\max})
=
\sum_{\ell>\ell_{\max}}(2\ell+1)\,
\max_{m,\sigma}|\widehat{c}_{\ell m}^{(\sigma)}|^{2},
\qquad
\mathcal{M}_{\mathrm{tail}}
\le
\eta_{\mathrm{tail}}\,\|\mathbf{c}\|_{\mathrm{rad}}^{2}
\ \Longrightarrow\
\mathbf{c}\in\mathcal{B}_{\mathrm{fs}},
\end{equation}
linking tail mass directly to Definition~\ref{def:ch4.finite-support-restriction} \cite{hansen1988,stratton1941}.
Equation~\eqref{eq:ch4.cumulative-tail-mass} is the watt-domain audit companion to
Eq.~\eqref{eq:ch4.high-ell-decay}: pointwise bounds integrate into cumulative tail budgets \cite{harrington2001,devaney2004}.

\begin{theorem}[Tail mass monotonicity and synthesis error]
\label{thm:ch4.tail-mass-monotone}
\(\mathcal{M}_{\mathrm{tail}}(\ell_{\max})\) decreases monotonically in \(\ell_{\max}\), and
\begin{equation}
\label{eq:ch4.tail-mass-synthesis-error}
\|\mathbf{c}-\mathbf{c}^{(\ell_{\max})}\|_{\mathrm{rad}}^{2}
=
\mathcal{M}_{\mathrm{tail}}(\ell_{\max})
\end{equation}
in the aligned VSH gauge \cite{hansen1988,stratton1941,harrington2001}. Targets with
\(\mathcal{M}_{\mathrm{tail}}(\ell_{\max})>\eta_{\mathrm{tail}}\|\mathbf{c}\|_{\mathrm{rad}}^{2}\) at every admissible
\(\ell_{\max}(k_{0}a)\) lie outside \(\mathcal{S}_{\mathrm{real}}(\mathcal{C})\) on \(B_{a}\) \cite{devaney2004}.
\end{theorem}
\begin{proof}
Proposition~\ref{prop:ch4.ell-tail-monotone} and Parseval
Eq.~\eqref{eq:ch3.mode-power-from-coefficient} \cite{hansen1988,stratton1941}.
\end{proof}

\paragraph{Worked illustration (slow tail decay).}
If measured \(\widehat{c}_{\ell m}^{(\sigma)}\) decay slower than Eq.~\eqref{eq:ch4.high-ell-decay} predicts for the
declared \(B_{a}\), the pattern is either mis-normalized or not synthesized by compact support on \(B_{a}\)
\cite{jackson1999,hansen1988,harrington2001}.

\paragraph{Problem (tail mass audit).}
\textbf{Problem.} Given \(\mathbf{c}_{\mathrm{target}}\) and \(\eta_{\mathrm{tail}}=10^{-2}\), find minimal
\(\ell_{\max}\) such that \(\mathcal{M}_{\mathrm{tail}}(\ell_{\max})\le\eta_{\mathrm{tail}}\|\mathbf{c}\|_{\mathrm{rad}}^{2}\).

\textbf{Step 1.} Compute partial sums Eq.~\eqref{eq:ch4.cumulative-tail-mass} \cite{hansen1988}.

\textbf{Step 2.} Compare to Eq.~\eqref{eq:ch4.ellmax-from-tail} for \(k_{0}a\) \cite{jackson1999,stratton1941}.

\textbf{Step 3.} If no \(\ell_{\max}\) satisfies both, enlarge \(a\) or revise \(\mathbf{c}_{\mathrm{target}}\)
\cite{devaney2004,harrington2001}.

\textbf{Physical meaning.} Tail mass is a watt certificate, not a plotting convention \cite{hansen1988}.

\paragraph{Tail exponent and realizability envelope.}
Define the tail exponent
\begin{equation}
\label{eq:ch4.tail-exponent}
\nu_{\mathrm{tail}}(\mathbf{c})
=
\limsup_{\ell\to\infty}
\frac{\log|\widehat{c}_{\ell m}^{(\sigma)}|}{\ell},
\end{equation}
with the limit taken over the dominant \((m,\sigma)\) sector at each \(\ell\) \cite{hansen1988,stratton1941}. On \(B_{a}\),
realizable \(\mathbf{c}\) satisfy \(\nu_{\mathrm{tail}}(\mathbf{c})\le-1/(k_{0}a)\) by Eq.~\eqref{eq:ch4.high-ell-decay}
\cite{jackson1999,harrington2001,devaney2004}. Faster-than-envelope decay is admissible; slower decay violates
\(\mathcal{B}_{\mathrm{fs}}\) \cite{hansen1988}.

\begin{corollary}[Tail exponent obstruction]
\label{cor:ch4.tail-exponent-obstruction}
If \(\nu_{\mathrm{tail}}(\mathbf{c}_{\mathrm{target}})>-1/(k_{0}a)\), then
\(\mathbf{c}_{\mathrm{target}}\notin\mathcal{S}_{\mathrm{real}}(\mathcal{C})\) on \(B_{a}\) at fixed \(k_{0}\), independent
of phase-only retuning on \(\Sigma_{a}\) \cite{stratton1941,hansen1988,devaney2004}.
\end{corollary}

\begin{figure}[t]
\centering
\ChOneFigBlock{%
\begin{tikzpicture}[ch1fig, node distance=7mm and 9mm]
  \node[ch1box] (ell) {$\ell$ index};
  \node[ch1box, right=12mm of ell] (env) {Eq.~\eqref{eq:ch4.high-ell-decay}};
  \node[ch1box, right=12mm of env] (mt) {$\mathcal{M}_{\mathrm{tail}}$};
  \node[ch1box, right=12mm of mt] (sr) {$\mathcal{S}_{\mathrm{real}}$};
  \draw[ch1flow] (ell) -- (env) -- (mt) -- (sr);
\end{tikzpicture}%
}
\caption{Subsection~\ref{sec:ch443}: high-\(\ell\) envelope bounds tail exponent before synthesis on compact \(B_{a}\).}
\label{fig:ch4.tail-exponent-envelope}
\end{figure}

\paragraph{Worked illustration (slow OAM tail).}
Measured \(\widehat{c}_{\ell m}^{(\sigma)}\) with \(\nu_{\mathrm{tail}}\approx-0.3/(k_{0}a)\) on a declared \(B_{a}\) fails
Corollary~\ref{cor:ch4.tail-exponent-obstruction}: the pattern is either mis-attributed to the wrong support or not
synthesizable on \(B_{a}\) \cite{hansen1988,jackson1999,harrington2001}.

\paragraph{Problem (tail exponent from sparse samples).}
\textbf{Problem.} Only \(\ell\le8\) samples are available. Can \(\nu_{\mathrm{tail}}\) be certified?

\textbf{Step 1.} Fit log-slope of \(|\widehat{c}_{\ell m}^{(\sigma)}|\) versus \(\ell\) on available data \cite{hansen1988}.

\textbf{Step 2.} Compare slope to \(-1/(k_{0}a)\) from Eq.~\eqref{eq:ch4.high-ell-decay} \cite{jackson1999,stratton1941}.

\textbf{Step 3.} If slope exceeds envelope, reject high-\(\ell\) targets even when \(\ell>8\) was not measured
\cite{devaney2004,harrington2001}.

\textbf{Physical meaning.} Tail exponent is an asymptotic realizability certificate, not a curve-fit aesthetic
\cite{stratton1941}.

\paragraph{Validity.}
Eq.~\eqref{eq:ch4.high-ell-decay} is asymptotic for \(\ell\gg k_{0}a\); low-\(\ell\) sectors require full operator
SVD Eq.~\eqref{eq:ch4.gamma-SVD} \cite{harrington2001}. Non-spherical \(\Omega_{s}\) modifies \(\mathcal{C}_{\ell}\) without
removing tail decay \cite{stratton1941}.

Finite-support spectral restrictions unify lateral and spherical cuts in Subsection~\ref{sec:ch444}.


\subsection{Finite-support spectral restrictions}
\label{sec:ch444}
% TOC tag: [HOME]

Radiative rank, lateral bandlimits, and \(\ell\)-tails are three readings of one finite-support constraint
\cite{stratton1941,hansen1988,devaney2004}. Subsection~\ref{sec:ch444} defines the unified restriction on
\(\mathcal{S}_{\mathrm{real}}(\mathcal{C})\).

\paragraph{Assumptions.}
Collect compact \((\Omega_{s},k_{0})\), propagating sector \(\mathcal{K}_{\mathrm{prop}}\), lateral bandlimit
\(k_{t,\max}(\eta)\), spherical cut \(\ell_{\max}\), and radiative rank \(r_{\mathrm{rad}}(\varepsilon)\) from
Sections~\ref{sec:ch43}--\ref{sec:ch443} \cite{stratton1941,harrington2001,hansen1988}.

\begin{definition}[Finite-support spectral restriction]
\label{def:ch4.finite-support-restriction}
The \emph{finite-support spectral restriction} \(\mathcal{B}_{\mathrm{fs}}(\eta,\varepsilon)\subset\mathbb{C}^{I}\) is
\begin{equation}
\label{eq:ch4.finite-support-restriction}
\mathcal{B}_{\mathrm{fs}}
=
\left\{
\mathbf{c}\in\mathcal{S}_{\mathrm{real}}(\mathcal{C}):
\|\mathbf{R}_{\ell_{\max}}\|_{\mathrm{rad}}\le\eta_{\mathrm{tail}}\|\mathbf{c}\|_{\mathrm{rad}},\
k_{t}(\mathbf{c})\le k_{t,\max}(\eta),\
\operatorname{rank}_{\varepsilon}(\mathbf{c})\le r_{\mathrm{rad}}(\varepsilon)
\right\},
\end{equation}
where \(k_{t}(\mathbf{c})\) is the maximum transverse wavenumber with appreciable weight in the plane-wave transform of
\(\mathbf{c}\) and \(\operatorname{rank}_{\varepsilon}\) counts singular directions above threshold
\cite{hansen1988,devaney2004,harrington2001}.
\end{definition}

Equation~\eqref{eq:ch4.finite-support-restriction} is the admissible spectral body inside realizability: synthesis
targets must lie in \(\mathcal{B}_{\mathrm{fs}}\), not merely in \(\mathbb{C}^{I}\) \cite{stratton1941,devaney2004}.
In aligned gauges,
\begin{equation}
\label{eq:ch4.restriction-inclusion}
\mathcal{B}_{\mathrm{fs}}
\subseteq
\mathcal{S}_{\mathrm{real}}(\mathcal{C})
\subseteq
\mathbb{C}^{I},
\end{equation}
with strict inclusions under compact \(\Omega_{s}\) \cite{hansen1988,harrington2001}.

\begin{theorem}[Equivalence of restriction gauges]
\label{thm:ch4.restriction-gauge-equivalence}
For ball support \(B_{a}\), tolerance pair \((\eta,\varepsilon)\), and VSH/plane-wave listings related by
Section~\ref{sec:ch38}, the cuts \((\ell_{\max},k_{t,\max},r_{\mathrm{rad}})\) are equivalent up to constants
\(C_{1},C_{2}\) depending only on \((\eta,\varepsilon,k_{0}a)\):
\begin{equation}
\label{eq:ch4.gauge-equivalence}
\ell_{\max}\sim k_{t,\max}/k_{0}
\sim
r_{\mathrm{rad}}(\varepsilon).
\end{equation}
\end{theorem}
\begin{proof}
Proposition~\ref{prop:ch4.geometry-spectral-equivalence}, Theorem~\ref{thm:ch4.duality-bandlimit}, and
Definition~\ref{def:ch4.radiative-rank} with basis transformation boundedness from
Eq.~\eqref{eq:ch3.basis-change-matrix} \cite{hansen1988,stratton1941,harrington2001}.
\end{proof}

Theorem~\ref{thm:ch4.restriction-gauge-equivalence} is the closure of Section~\ref{sec:ch44}: bandlimits are not
representation accidents---they are the same finite-support constraint in different coordinates
\cite{devaney2004,hansen1988}. Design in plane-wave gauges must certify \(\mathcal{B}_{\mathrm{fs}}\) after
transforming to the laboratory \(\mathcal{C}\) \cite{harrington2001}.

\paragraph{Coupling to accessibility.}
\(\mathcal{B}_{\mathrm{fs}}\) bounds synthesis; accessibility envelope
\(\mathcal{E}_{\mathrm{acc}}^{(N)}\) bounds certification \cite{harrington2001,devaney2004}. Realizable and
certifiable targets lie in
\begin{equation}
\label{eq:ch4.realizable-accessible-intersection}
\mathcal{B}_{\mathrm{fs}}\cap\mathcal{E}_{\mathrm{acc}}^{(N)}(\mathcal{C}),
\end{equation}
which may be strictly smaller than either set alone \cite{devaney2004,hansen1988}.

\begin{figure}[t]
\centering
\ChOneFigBlock{%
\begin{tikzpicture}[ch1fig, node distance=7mm and 8mm]
  \node[ch1box, minimum width=34mm] (ci) {$\mathbb{C}^{I}$};
  \node[ch1box, minimum width=30mm, below left=10mm and 4mm of ci] (sr) {$\mathcal{S}_{\mathrm{real}}$};
  \node[ch1box, minimum width=30mm, below right=10mm and 4mm of ci] (ea) {$\mathcal{E}_{\mathrm{acc}}^{(N)}$};
  \node[ch1box, minimum width=32mm, below=14mm of sr] (bfs) {$\mathcal{B}_{\mathrm{fs}}$\\tail + $k_{t}$ + rank};
  \node[ch1box, minimum width=28mm, below=14mm of ea] (cap) {certified\\intersection};
  \draw[ch1flow] (sr) -- (ci);
  \draw[ch1flow] (ea) -- (ci);
  \draw[ch1flow] (bfs) -- (sr);
  \draw[ch1flow] (bfs) -- (cap);
  \draw[ch1flow] (ea) -- (cap);
\end{tikzpicture}%
}
\caption{Subsection~\ref{sec:ch444}: \(\mathcal{B}_{\mathrm{fs}}\) is the spectral body inside
\(\mathcal{S}_{\mathrm{real}}\); laboratory success requires
Eq.~\eqref{eq:ch4.realizable-accessible-intersection}.}
\label{fig:ch4.finite-support-restriction}
\end{figure}

Figure~\ref{fig:ch4.finite-support-restriction} restates Figure~\ref{fig:ch4.synthesis-accessibility-sets} after
Sections~\ref{sec:ch43}--\ref{sec:ch443}: geometry compresses synthesis through \(\mathcal{B}_{\mathrm{fs}}\), while
instruments compress certification through \(\mathcal{E}_{\mathrm{acc}}^{(N)}\) \cite{harrington2001,devaney2004}.
Theorem~\ref{thm:ch4.restriction-gauge-equivalence} guarantees that the three cuts inside
\(\mathcal{B}_{\mathrm{fs}}\) are redundant readings of one obstruction once \((\eta,\varepsilon,k_{0}a)\) are fixed
\cite{hansen1988,stratton1941}.

\begin{corollary}[Gauge-redundant audit]
\label{cor:ch4.gauge-redundant-audit}
If \(\mathbf{c}_{\mathrm{target}}\) fails any one of the inequalities in
Eq.~\eqref{eq:ch4.finite-support-restriction} on \(B_{a}\) at tolerance \((\eta,\varepsilon)\), then it fails
\(\mathcal{B}_{\mathrm{fs}}\) and lies outside the realizable-accessible intersection
Eq.~\eqref{eq:ch4.realizable-accessible-intersection} for every equivalent gauge related by
Theorem~\ref{thm:ch4.restriction-gauge-equivalence} \cite{hansen1988,devaney2004,harrington2001}.
\end{corollary}

\paragraph{Worked illustration (unified audit).}
Audit \(\mathbf{c}_{\mathrm{target}}\) by: (i) tail check Eq.~\eqref{eq:ch4.ellmax-tail-link}; (ii) lateral check
Theorem~\ref{thm:ch4.duality-bandlimit}; (iii) rank check Definition~\ref{def:ch4.radiative-rank}; (iv) propagating check
Eq.~\eqref{eq:ch4.propagating-realizability} \cite{harrington2001,devaney2004}.

\paragraph{Problem (restriction failure diagnosis).}
\textbf{Problem.} \(\mathbf{c}_{\mathrm{target}}\in\mathcal{S}_{\mathrm{real}}(\mathcal{C})\) by MoM fit but fails
far-zone power test.

\textbf{Step 1.} Test Eq.~\eqref{eq:ch4.propagating-realizability} for evanescent contamination \cite{harrington2001}.

\textbf{Step 2.} Test Eq.~\eqref{eq:ch4.high-ell-decay} for impossible \(\ell\) tails \cite{hansen1988}.

\textbf{Step 3.} If both pass, inspect \(\mathcal{E}_{\mathrm{acc}}^{(N)}\) rather than \(\mathcal{B}_{\mathrm{fs}}\)
\cite{devaney2004}.

\textbf{Physical meaning.} Restriction failures are geometry-limited, not normalization errors \cite{stratton1941}.

\paragraph{Unified realizability audit protocol.}
Define the audit functional
\begin{equation}
\label{eq:ch4.unified-audit-functional}
\mathcal{A}_{\mathrm{fs}}[\mathbf{c}_{\mathrm{target}}]
=
\max\Big\{
\frac{\|\mathbf{R}_{\ell_{\max}}\|_{\mathrm{rad}}}{\eta_{\mathrm{tail}}\|\mathbf{c}\|_{\mathrm{rad}}},
\frac{k_{t}(\mathbf{c})}{k_{t,\max}(\eta)},
\frac{\operatorname{rank}_{\varepsilon}(\mathbf{c})}{r_{\mathrm{rad}}(\varepsilon)},
\delta_{\mathrm{prop}}
\Big\},
\end{equation}
with \(\delta_{\mathrm{prop}}=0\) iff Eq.~\eqref{eq:ch4.propagating-realizability} holds and
\(\delta_{\mathrm{prop}}=1\) otherwise \cite{harrington2001,devaney2004,hansen1988}. Then
\(\mathbf{c}_{\mathrm{target}}\in\mathcal{B}_{\mathrm{fs}}\iff\mathcal{A}_{\mathrm{fs}}\le1\) and
\(\mathbf{c}_{\mathrm{target}}\in\mathcal{S}_{\mathrm{real}}(\mathcal{C})\) requires additionally
\(\delta_{\mathrm{syn}}=0\) from Eq.~\eqref{eq:ch4.synthesis-residual-norm} \cite{stratton1941}.

\begin{theorem}[Finite-support audit certificate]
\label{thm:ch4.finite-support-audit}
For compact \(B_{a}\), flux-normalized \(\mathcal{C}\), and tolerances \((\eta,\varepsilon)\),
\(\mathbf{c}\in\mathcal{B}_{\mathrm{fs}}\cap\mathcal{S}_{\mathrm{real}}(\mathcal{C})\) if and only if
\(\mathcal{A}_{\mathrm{fs}}[\mathbf{c}]\le1\) and \(\delta_{\mathrm{syn}}=0\) \cite{hansen1988,stratton1941,devaney2004,harrington2001}.
\end{theorem}
\begin{proof}
Definition~\ref{def:ch4.finite-support-restriction}, Theorem~\ref{thm:ch4.restriction-gauge-equivalence}, and
Principle~\ref{prin:ch4.operator-chain} \cite{hansen1988,stratton1941,harrington2001}.
\end{proof}

Theorem~\ref{thm:ch4.finite-support-audit} is the operational closure of Section~\ref{sec:ch44}: four inequalities in
Eq.~\eqref{eq:ch4.finite-support-restriction} collapse to one audit functional plus a range residual
\cite{devaney2004}. Laboratories may implement \(\mathcal{A}_{\mathrm{fs}}\) as a pre-flight check before MoM inversion
\cite{harrington2001}.

\paragraph{Worked illustration (audit pass, synthesis fail).}
A \(\mathbf{c}\) may satisfy \(\mathcal{A}_{\mathrm{fs}}\le1\) yet fail \(\delta_{\mathrm{syn}}=0\) when target components lie in
\(\ker\boldsymbol{\Gamma}_\omega^{\dagger}\) \cite{devaney2004,harrington2001}. Conversely, \(\delta_{\mathrm{syn}}=0\) with
\(\mathcal{A}_{\mathrm{fs}}>1\) signals spectral support violation despite a formal MoM fit \cite{hansen1988,stratton1941}.

\paragraph{Certified intersection closure.}
Laboratory success requires
\begin{equation}
\label{eq:ch4.certified-intersection-closure}
\mathbf{c}_{\mathrm{cert}}
\in
\mathcal{S}_{\mathrm{real}}(\mathcal{C})
\cap
\mathcal{B}_{\mathrm{fs}}
\cap
\mathcal{E}_{\mathrm{acc}}^{(N)},
\end{equation}
extending Eq.~\eqref{eq:ch4.realizable-accessible-intersection} \cite{harrington2001,devaney2004}. Synthesis without
certification is \(\mathbf{c}\in\mathcal{S}_{\mathrm{real}}\cap\mathcal{B}_{\mathrm{fs}}\); certification without synthesis is
\(\mathbf{c}\in\mathcal{E}_{\mathrm{acc}}^{(N)}\) only \cite{hansen1988}.

\begin{theorem}[Three-way realizability certificate]
\label{thm:ch4.three-way-certificate}
\(\mathbf{c}\) is laboratory-realizable and auditable if and only if
\(\mathcal{A}_{\mathrm{fs}}[\mathbf{c}]\le1\), \(\delta_{\mathrm{syn}}=0\), and
\(\mathbf{c}\in\mathcal{E}_{\mathrm{acc}}^{(N)}\) at declared instrument rank \(N\) \cite{hansen1988,stratton1941,harrington2001,devaney2004}.
\end{theorem}
\begin{proof}
Theorem~\ref{thm:ch4.finite-support-audit}, Definition~\ref{def:ch4.accessibility-envelope}, and
Eq.~\eqref{eq:ch4.certified-intersection-closure} \cite{hansen1988,harrington2001,devaney2004}.
\end{proof}

\paragraph{Worked illustration (synthesizable but uncertifiable).}
A target \(\mathbf{c}\) may satisfy \(\mathcal{A}_{\mathrm{fs}}\le1\) and \(\delta_{\mathrm{syn}}=0\) yet fail
\(\mathcal{E}_{\mathrm{acc}}^{(N)}\) when the receiver resolves fewer modes than synthesis claims \cite{harrington2001,devaney2004}.

\paragraph{Problem (three-way audit ordering).}
\textbf{Problem.} In what order should \(\mathcal{A}_{\mathrm{fs}}\), \(\delta_{\mathrm{syn}}\), and \(\mathcal{E}_{\mathrm{acc}}^{(N)}\) be tested?

\textbf{Step 1.} Test \(\mathcal{A}_{\mathrm{fs}}\) first to reject geometry-impossible targets cheaply \cite{hansen1988}.

\textbf{Step 2.} Test \(\delta_{\mathrm{syn}}=0\) on survivors \cite{stratton1941,harrington2001}.

\textbf{Step 3.} Test \(\mathcal{E}_{\mathrm{acc}}^{(N)}\) last for laboratory claims \cite{devaney2004}.

\textbf{Physical meaning.} Realizability, finite support, and accessibility are sequential certificates \cite{stratton1941}.

\paragraph{Validity.}
Definition~\ref{def:ch4.finite-support-restriction} is narrowband at fixed \(\omega\); broadband pulses require
\(\mathcal{B}_{\mathrm{fs}}(\omega)\) families \cite{jackson1999,harrington2001}.

Section~\ref{sec:ch44} has unified localization and bandlimits into \(\mathcal{B}_{\mathrm{fs}}\). Section~\ref{sec:ch45}
couples finite geometry to VSH eigenfields and selection rules on \(\operatorname{ran}\boldsymbol{\Gamma}_\omega\)
\cite{stratton1941,hansen1988,harrington2001}.


\section{Coupling to Angular-Momentum Eigenfields}
\label{sec:ch45}

Finite-support spectral restrictions of Section~\ref{sec:ch44} bound which coefficient vectors lie in
\(\mathcal{B}_{\mathrm{fs}}\); they do not yet specify how impressed sources \emph{weight} individual angular-momentum
eigenfields on \(\operatorname{ran}\mathcal{R}_\omega\) \cite{stratton1941,hansen1988,devaney2004}. Section~\ref{sec:ch45}
closes that gap: flux-normalized VSH pairs from Section~\ref{sec:ch34} become the synthesis coordinates of
\(\boldsymbol{\Gamma}_\omega\), coupling integrals impose selection rules on achievable weights, geometry sets the
coupling rank, and normalization fixes the modal weight ceiling that laboratories may report
\cite{harrington2001,stratton1941}.


\subsection{Eigenfield expansions of radiated fields}
\label{sec:ch451}

Once \(\mathcal{B}_{\mathrm{fs}}\) has bounded the spectral body of synthesis, the remaining question is coordinate
resolution: in which angular-momentum eigenfields must \(\mathfrak{F}=\mathcal{R}_\omega[\Jimp]\) be listed so that
\(\mathbf{c}\in\mathcal{S}_{\mathrm{real}}(\mathcal{C})\) is both synthesizable and auditable
\cite{hansen1988,stratton1941}? Section~\ref{sec:ch34} fixed \(\mathbf{u}_{\ell m}^{(\sigma)}\) as the outgoing
harmonic pairs Eqs.~\eqref{eq:ch3.VSH-M-pair}--\eqref{eq:ch3.Nlm-def}; Subsection~\ref{sec:ch451} reads those pairs as
the inevitable synthesis basis on \(\mathcal{F}_{\mathrm{rad}}\) when \(\mathcal{C}\) is flux-normalized
\cite{stratton1941,hansen1988}.

\paragraph{Assumptions.}
Fix \(\omega>0\), outgoing \(\mathcal{F}_{\mathrm{rad}}\), \(\mathcal{J}_{\mathrm{real}}\) from
Subsection~\ref{sec:ch414}, and \(\mathcal{C}\in\mathcal{R}_{\mathrm{adm}}\) with flux normalization
Eqs.~\eqref{eq:ch3.VSH-unit-flux-normalization}--\eqref{eq:ch3.continuous-measure-flux-fix} from
Section~\ref{sec:ch35} \cite{stratton1941,hansen1988}. Presuppose \(\mathcal{B}_{\mathrm{fs}}\) of
Definition~\ref{def:ch4.finite-support-restriction} and radiative rank \(r_{\mathrm{rad}}(\varepsilon)\) from
Definition~\ref{def:ch4.radiative-rank} \cite{harrington2001,devaney2004}.

\paragraph{Eigenfield listing on \(\mathcal{F}_{\mathrm{rad}}\).}
For \(\mathfrak{F}\in\mathcal{F}_{\mathrm{rad}}\), the normalized expansion imported from
Eq.~\eqref{eq:ch3.normalized-coefficient-resolution} reads
\begin{equation}
\label{eq:ch4.eigenfield-synthesis-expansion}
\mathfrak{F}
=
\sum_{\ell,m,\sigma}
\widehat{c}_{\ell m}^{(\sigma)}\,
\widehat{\mathbf{u}}_{\ell m}^{(\sigma)},
\qquad
\sigma\in\{\mathbf{M},\mathbf{N}\},
\end{equation}
with \(\widehat{\mathbf{u}}_{\ell m}^{(\sigma)}\) the flux-normalized partners of
Eq.~\eqref{eq:ch3.VSH-M-pair} \cite{hansen1988,stratton1941}. Equation~\eqref{eq:ch4.eigenfield-synthesis-expansion}
is not a new harmonic construction: it is the synthesis-side restatement of
Eq.~\eqref{eq:ch3.eigenfield-expansion-on-frad} on \(\operatorname{ran}\mathcal{R}_\omega\) after
\(\mathcal{N}\) has fixed the gauge \cite{stratton1941,harrington2001}. Coefficients are synthesis coordinates:
\begin{equation}
\label{eq:ch4.gamma-eigenfield-coordinates}
\mathbf{c}
=
\boldsymbol{\Gamma}_\omega[\Jimp]
=
\big\{
\widehat{c}_{\ell m}^{(\sigma)}[\Jimp]
\big\}_{\ell,m,\sigma},
\qquad
\widehat{c}_{\ell m}^{(\sigma)}[\Jimp]
=
\big\langle
\mathcal{R}_\omega[\Jimp],\,
\widehat{\mathbf{u}}_{\ell m}^{(\sigma)}
\big\rangle_{\mathrm{rad}},
\end{equation}
importing the pairing Eq.~\eqref{eq:ch3.eigenfield-coefficient-pairing} with \(\mathcal{C}\) identified as
\(\mathcal{N}\circ\mathcal{P}_{\mathrm{coef}}\) on the VSH sector \cite{hansen1988,stratton1941}. Equation
\eqref{eq:ch4.gamma-eigenfield-coordinates} makes \(\mathcal{S}_{\mathrm{real}}(\mathcal{C})\) a subset of
\(\mathbb{C}^{I_{\mathrm{VSH}}}\): realizability is a statement about achievable \(\widehat{c}_{\ell m}^{(\sigma)}\), not
about the abstract existence of labels \cite{devaney2004,harrington2001}.

\paragraph{Far-zone pattern form.}
On spheres where Theorem~\ref{thm:ch2.far-field-approximation} holds,
Eq.~\eqref{eq:ch3.far-pattern-eigenfield-coupling} specializes to
\begin{equation}
\label{eq:ch4.far-pattern-eigenfield-sum}
\mathcal{A}_{\infty}[\Jimp](\hat{\mathbf{r}})
=
\sum_{\ell,m,\sigma}
\widehat{c}_{\ell m}^{(\sigma)}[\Jimp]\,
\widehat{\mathbf{e}}_{\ell m}^{(\sigma)}(\hat{\mathbf{r}}),
\end{equation}
with \(\widehat{\mathbf{e}}_{\ell m}^{(\sigma)}\) the transverse far-zone profiles of
\(\widehat{\mathbf{u}}_{\ell m}^{(\sigma)}\) \cite{stratton1941,jackson1999}. Equation
\eqref{eq:ch4.far-pattern-eigenfield-sum} is the pattern-space image of
\(\boldsymbol{\Gamma}_\omega\): every synthesizable far-zone template is a finite superposition of eigenfield
shapes weighted by \(\widehat{c}_{\ell m}^{(\sigma)}\in\mathcal{S}_{\mathrm{real}}(\mathcal{C})\)
\cite{hansen1988,harrington2001}. Targets whose \(\mathcal{A}_{\infty}\) requires eigenfield components outside
\(\mathcal{B}_{\mathrm{fs}}\) fail before coupling integrals are evaluated \cite{devaney2004}.

\begin{figure}[t]
\centering
\ChOneFigBlock{%
\begin{tikzpicture}[ch1fig, node distance=7mm and 10mm]
  \node[ch1box] (j) {$\mathcal{J}_{\mathrm{real}}$};
  \node[ch1box, right=13mm of j] (rw) {$\mathcal{R}_\omega$};
  \node[ch1box, right=13mm of rw] (fr) {$\mathcal{F}_{\mathrm{rad}}$};
  \node[ch1box, below=12mm of fr] (exp) {Eq.~\eqref{eq:ch4.eigenfield-synthesis-expansion}};
  \node[ch1box, left=13mm of exp] (c) {$\mathbf{c}$};
  \node[ch1box, left=13mm of c] (g) {$\boldsymbol{\Gamma}_\omega$};
  \draw[ch1flow] (j) -- (rw) -- (fr);
  \draw[ch1flow] (g) -- (c);
  \draw[ch1flow] (fr) -- (exp);
  \draw[ch1flow] (c) -- (exp);
\end{tikzpicture}%
}
\caption{Subsection~\ref{sec:ch451}: synthesis lists \(\mathfrak{F}=\mathcal{R}_\omega[\Jimp]\) in flux-normalized
eigenfields; \(\mathbf{c}\) is the coordinate vector of Eq.~\eqref{eq:ch4.gamma-eigenfield-coordinates}.}
\label{fig:ch4.eigenfield-coupling-chain}
\end{figure}

Figure~\ref{fig:ch4.eigenfield-coupling-chain} is the coordinate closure of Principle~\ref{prin:ch4.operator-chain}:
\(\boldsymbol{\Gamma}_\omega\) is not an abstract listing but the eigenfield coordinate map on
\(\mathcal{J}_{\mathrm{real}}\) \cite{harrington2001,devaney2004}. Orthogonality
Theorem~\ref{thm:ch3.vsh-orthogonality} makes distinct \((\ell,m,\sigma)\) labels power-orthogonal on
\(\mathcal{F}_{\mathrm{rad}}\) after \(\mathcal{N}\), so watt budgets decompose across eigenfield sectors in
Eq.~\eqref{eq:ch3.mode-power-from-coefficient} \cite{hansen1988,stratton1941}.

\begin{theorem}[Eigenfield synthesis coordinates]
\label{thm:ch4.eigenfield-synthesis-coordinates}
Let \(\mathcal{C}\) be the flux-normalized VSH listing of Section~\ref{sec:ch35}. For every
\(\Jimp\in\mathcal{J}_{\mathrm{real}}\),
\begin{equation}
\label{eq:ch4.synthesis-set-eigenfield}
\mathcal{S}_{\mathrm{real}}(\mathcal{C})
=
\left\{
\mathbf{c}\in\mathbb{C}^{I_{\mathrm{VSH}}}:
\mathbf{c}=\boldsymbol{\Gamma}_\omega[\Jimp]\ \text{for some}\ \Jimp\in\mathcal{J}_{\mathrm{real}}
\right\}
\subseteq
\operatorname{span}\big\{
\widehat{\mathbf{u}}_{\ell m}^{(\sigma)}:\,
(\ell,m,\sigma)\ \text{within}\ \mathcal{B}_{\mathrm{fs}}
\big\},
\end{equation}
with strict inclusion when compact \(\Omega_{s}\) compresses the accessible index set below the formal listing
\cite{stratton1941,hansen1988,devaney2004}. Moreover, power satisfies
\begin{equation}
\label{eq:ch4.eigenfield-power-decomposition}
\|\mathbf{c}\|_{\mathrm{rad}}^{2}
=
\sum_{\ell,m,\sigma}
\big|\widehat{c}_{\ell m}^{(\sigma)}\big|^{2}
\end{equation}
in the aligned flux gauge of Theorem~\ref{thm:ch3.vsh-unit-flux-diagonal} \cite{hansen1988,stratton1941}.
\end{theorem}
\begin{proof}
Eq.~\eqref{eq:ch4.synthesis-set-eigenfield} is Definition~\ref{def:ch4.realizable-synthesis-set} with
\(\mathcal{C}\) specialized to Eq.~\eqref{eq:ch4.gamma-eigenfield-coordinates} \cite{stratton1941}. Inclusion in the
span of accessible eigenfields follows from \(\mathcal{B}_{\mathrm{fs}}\) and compactness of
\(\boldsymbol{\Gamma}_\omega\) Theorem~\ref{thm:ch4.compact-radiation-map} \cite{harrington2001,devaney2004}.
Eq.~\eqref{eq:ch4.eigenfield-power-decomposition} is Parseval under
Theorem~\ref{thm:ch3.vsh-orthogonality} and flux normalization
Eq.~\eqref{eq:ch3.VSH-unit-flux-normalization} \cite{hansen1988,stratton1941}.
\end{proof}

Theorem~\ref{thm:ch4.eigenfield-synthesis-coordinates} elevates the VSH listing from representation convenience to
synthesis geometry: \(\mathcal{S}_{\mathrm{real}}(\mathcal{C})\) is the image of \(\mathcal{J}_{\mathrm{real}}\) in the
eigenfield coordinate chart, and \(\mathcal{B}_{\mathrm{fs}}\) selects which coordinates may carry appreciable weight
\cite{hansen1988,devaney2004}. Laboratories that report \(\widehat{c}_{\ell m}^{(\sigma)}\) without specifying
Eqs.~\eqref{eq:ch3.VSH-unit-flux-normalization}--\eqref{eq:ch3.continuous-measure-flux-fix} compare coordinates in
different gauges, not different physics \cite{harrington2001}.

\paragraph{Integral route consistency.}
Theorem~\ref{thm:ch3.expansion-consistency} guarantees that eigenfield coefficients computed by
Eq.~\eqref{eq:ch4.gamma-eigenfield-coordinates} agree with the modal sum of
Eq.~\eqref{eq:ch2.radiation-map-modal-sum} once indices are aligned by
Eq.~\eqref{eq:ch3.spectral-index-alignment} \cite{stratton1941,harrington2001}. The synthesis audit therefore commutes:
\begin{equation}
\label{eq:ch4.integral-eigenfield-commute}
\boldsymbol{\Gamma}_\omega[\Jimp]
=
\mathcal{C}\circ\mathcal{R}_\omega[\Jimp]
=
\mathcal{C}\circ\mathcal{I}_\omega[\Jimp]
\quad\text{on}\ \mathcal{F}_{\mathrm{rad}},
\end{equation}
with \(\mathcal{I}_\omega\) the integral route of Subsection~\ref{sec:ch251} \cite{harrington2001}. Equation
\eqref{eq:ch4.integral-eigenfield-commute} prevents a common pipeline error: evaluating coupling integrals on raw
\(\mathcal{S}_\omega[\Jimp]\) before \(\Pi_{\mathrm{rad}}\) mixes reactive content into \(\widehat{c}_{\ell m}^{(\sigma)}\)
\cite{stratton1941,devaney2004}.

\begin{corollary}[Truncated eigenfield synthesis]
\label{cor:ch4.truncated-eigenfield-synthesis}
Let \(\mathcal{T}_{\ell_{\max}}\) be Eq.~\eqref{eq:ch3.ellmax-truncation-projector}. Then
\(\mathcal{S}_{\mathrm{real}}^{(\ell_{\max})}(\mathcal{C})
=
\boldsymbol{\Gamma}_\omega^{(\ell_{\max})}[\mathcal{J}_{\mathrm{real}}]\)
and
\(\|\mathbf{c}-\mathbf{c}^{(\ell_{\max})}\|_{\mathrm{rad}}
\le
\|\mathbf{R}_{\ell_{\max}}\|_{\mathrm{rad}}\)
with remainder \(\mathbf{R}_{\ell_{\max}}\) from Eq.~\eqref{eq:ch3.truncation-remainder-energy}
\cite{hansen1988,stratton1941,harrington2001}.
\end{corollary}

Corollary~\ref{cor:ch4.truncated-eigenfield-synthesis} links eigenfield listing to truncated synthesis sets of
Proposition~\ref{prop:ch4.compressed-synthesis}: raising \(\ell_{\max}\) refines \(\mathbf{c}\) monotonically but cannot
excavate sectors forbidden by \(\mathcal{B}_{\mathrm{fs}}\) \cite{devaney2004,hansen1988}.

\paragraph{Superposition on \(\mathcal{J}_{\mathrm{real}}\).}
Linearity of \(\mathcal{R}_\omega\) and Eq.~\eqref{eq:ch3.representation-linearity} give
\begin{equation}
\label{eq:ch4.eigenfield-superposition}
\boldsymbol{\Gamma}_\omega\big[a\Jimp_{1}+b\Jimp_{2}\big]
=
a\,\boldsymbol{\Gamma}_\omega[\Jimp_{1}]
+
b\,\boldsymbol{\Gamma}_\omega[\Jimp_{2}]
\end{equation}
for \(a,b\in\mathbb{C}\) and \(\Jimp_{1},\Jimp_{2}\in\mathcal{J}_{\mathrm{real}}\) \cite{stratton1941,harrington2001}.
Equation~\eqref{eq:ch4.eigenfield-superposition} is the synthesis-side linearity law: eigenfield weights add in
\(\mathbf{c}\) because radiation is linear in impressed sources on \(\mathcal{F}_{\mathrm{rad}}\) \cite{stratton1941}.
Nonlinearity enters only through boundary classes, material response, or inversion---not through the eigenfield listing
itself \cite{harrington2001,devaney2004}.

\paragraph{Discrete MoM reduction.}
On a patch mesh, Eq.~\eqref{eq:ch4.mom-synthesis-matrix} becomes
\begin{equation}
\label{eq:ch4.mom-eigenfield-matrix}
\mathbf{c}
=
\mathbf{T}_{\mathrm{VSH}}\,\mathbf{I}_{s},
\qquad
(T_{\mathrm{VSH}})_{(\ell m\sigma),p}
=
\big\langle
\mathcal{R}_\omega[\boldsymbol{\Lambda}_{p}],\,
\widehat{\mathbf{u}}_{\ell m}^{(\sigma)}
\big\rangle_{\mathrm{rad}},
\end{equation}
with Rao--Wilton--Glisson bases \(\{\boldsymbol{\Lambda}_{p}\}\) \cite{harrington2001,collin2001}. Equation
\eqref{eq:ch4.mom-eigenfield-matrix} is the discrete form of Eq.~\eqref{eq:ch4.gamma-eigenfield-coordinates}: each row
is an eigenfield coupling to a patch current \cite{stratton1941}. Rank of \(\mathbf{T}_{\mathrm{VSH}}\) is bounded by
\(r_{\mathrm{rad}}(\varepsilon)\) and Corollary~\ref{cor:ch4.channel-count} \cite{harrington2001,devaney2004}.

\paragraph{Worked illustration (dominant \(\ell=1\) sector).}
For \(k_{0}D_{s}\ll1\), \(\boldsymbol{\Gamma}_\omega[\mathcal{J}_{\mathrm{real}}]\) concentrates in
\(\ell=1\) blocks of Eq.~\eqref{eq:ch4.eigenfield-synthesis-expansion} \cite{jackson1999,stratton1941}. A target
\(\mathbf{c}\) with dominant \(\widehat{c}_{10}^{(\mathbf{N})}\) is realizable only if a dipole-class feed exists in
\(\mathcal{J}_{\mathrm{real}}\); claiming \(\ell=4\) weight without enlarging support violates
Eq.~\eqref{eq:ch4.high-ell-decay} \cite{hansen1988,harrington2001}.

\paragraph{Problem (eigenfield coordinate audit).}
\textbf{Problem.} A report lists \(\widehat{c}_{\ell m}^{(\sigma)}\) in a geometrically normalized basis. Convert to
flux-normalized synthesis coordinates for comparison with \(\mathcal{S}_{\mathrm{real}}(\mathcal{C})\).

\textbf{Step 1.} Identify scale factors \(a_{\ell m}^{(\sigma)}\) between geometric and flux gauges from
Eqs.~\eqref{eq:ch3.VSH-unit-flux-normalization}--\eqref{eq:ch3.VSH-unit-energy-norm} \cite{hansen1988}.

\textbf{Step 2.} Apply \(\widehat{c}'_{\ell m}=\widehat{c}_{\ell m}/a_{\ell m}^{(\sigma)}\) per
Eq.~\eqref{eq:ch3.coupling-normalization-propagation} \cite{harrington2001}.

\textbf{Step 3.} Test \(\mathbf{c}'\in\mathcal{B}_{\mathrm{fs}}\) and \(\mathbf{r}_{\min}=\mathbf{0}\) from
Principle~\ref{prin:ch4.operator-chain} \cite{devaney2004}.

\textbf{Physical meaning.} Eigenfield coordinates are synthesis objects only after \(\mathcal{N}\) is fixed
\cite{stratton1941}.

\paragraph{Polarization sector coupling and synthesis mixing.}
On \(\mathcal{F}_{\mathrm{rad}}\), \(\mathbf{M}\) and \(\mathbf{N}\) sectors couple through the dyadic structure of
\(\mathcal{R}_\omega\) even when \(\widehat{c}_{\ell m}^{(\mathbf{M})}\) and \(\widehat{c}_{\ell m}^{(\mathbf{N})}\) are
listed separately \cite{stratton1941,hansen1988}. Define partial couplings
\begin{equation}
\label{eq:ch4.sector-partial-coupling}
\gamma_{\ell m}^{(\sigma)}[\Jimp]
=
\big\langle
\mathcal{R}_\omega[\Jimp],\,
\widehat{\mathbf{u}}_{\ell m}^{(\sigma)}
\big\rangle_{\mathrm{rad}},
\qquad
\sigma\in\{\mathbf{M},\mathbf{N}\},
\end{equation}
and the normalized sector coupling matrix
\begin{equation}
\label{eq:ch4.sector-coupling-matrix}
\mathcal{S}_{\sigma\sigma'}^{(\ell m)}[\Jimp]
=
\frac{
\gamma_{\ell m}^{(\sigma)}[\Jimp]\,
\overline{\gamma_{\ell m}^{(\sigma')}[\Jimp]}
}{
\|\mathcal{R}_\omega[\Jimp]\|_{\mathrm{rad}}^{2}
},
\end{equation}
for \(\sigma,\sigma'\in\{\mathbf{M},\mathbf{N}\}\) \cite{hansen1988,harrington2001}. On isotropic compact feeds,
\(\mathcal{S}_{\mathrm{MN}}^{(\ell m)}\approx0\) for \(\ell=1\); on chiral or offset feeds, sector mixing is mandatory in
synthesis templates \cite{jackson1999,devaney2004,stratton1941}.

\begin{corollary}[Sector-diagonal dominance on small supports]
\label{cor:ch4.sector-diagonal-dominance}
For \(k_{0}D_{s}\ll1\) and centered \(\Omega_{s}\), \(\mathcal{S}_{\mathrm{real}}(\mathcal{C})\) is well approximated by
the \(\ell=1\) sector-diagonal subspace spanned by \(\widehat{\mathbf{u}}_{1m}^{(\mathbf{M})}\) and
\(\widehat{\mathbf{u}}_{1m}^{(\mathbf{N})}\) with negligible \(\mathcal{S}_{\mathrm{MN}}^{(1m)}\)
\cite{jackson1999,stratton1941,hansen1988}. Targets with large cross-sector weight at \(\ell=1\) require broken
symmetry in \(\mathcal{J}_{\mathrm{real}}\) \cite{harrington2001,devaney2004}.
\end{corollary}

\paragraph{Worked illustration (offset patch).}
An offset patch on \(B_{a}\) activates \(\mathcal{S}_{\mathrm{MN}}^{(1m)}\neq0\): a pure \(\mathbf{M}\)-sector template cannot
synthesize the far-zone pattern without \(\mathbf{N}\)-sector weight \cite{stratton1941,harrington2001}. Sector audits must
precede \(\ell>1\) claims on electrically small supports \cite{hansen1988,devaney2004}.

\paragraph{Problem (single-sector template failure).}
\textbf{Problem.} A design optimizes only \(\widehat{c}_{\ell m}^{(\mathbf{M})}\) but measured \(\mathbf{c}\) shows comparable
\(\widehat{c}_{\ell m}^{(\mathbf{N})}\). Is the template realizable?

\textbf{Step 1.} Compute \(\mathcal{S}_{\mathrm{MN}}^{(\ell m)}\) from Eq.~\eqref{eq:ch4.sector-coupling-matrix}
\cite{hansen1988,stratton1941}.

\textbf{Step 2.} If cross-sector coupling exceeds \(\varepsilon\sigma_{1}\), enlarge the synthesis basis to include both
sectors \cite{harrington2001,devaney2004}.

\textbf{Step 3.} Retest \(\mathbf{c}\in\mathcal{B}_{\mathrm{fs}}\) and \(\delta_{\mathrm{syn}}=0\) \cite{devaney2004}.

\textbf{Physical meaning.} Polarization sectors are coupled synthesis coordinates, not independent design knobs
\cite{stratton1941}.

\paragraph{Validity.}
Eqs.~\eqref{eq:ch4.eigenfield-synthesis-expansion}--\eqref{eq:ch4.eigenfield-power-decomposition} assume linear
Maxwell theory, outgoing closure, and aligned VSH flux gauge \cite{stratton1941,hansen1988}. Continuous angular spectra
require the rigged extension of Subsection~\ref{sec:ch313}; discrete \((\ell,m,\sigma)\) is the narrowband certificate
used in synthesis audits \cite{harrington2001}. Anisotropic media mix \(\sigma\) sectors without removing the eigenfield
coordinate interpretation of \(\mathbf{c}\) \cite{stratton1941}.

Coupling integrals that produce the weights \(\widehat{c}_{\ell m}^{(\sigma)}\) from \(\Jimp\) are developed in
Subsection~\ref{sec:ch452}.


\subsection{Coupling integrals and selection rules}
\label{sec:ch452}

Eigenfield coordinates name the synthesis chart; coupling integrals name the linear functionals that map impressed
currents into those coordinates \cite{stratton1941,harrington2001}. On compact \(\Omega_{s}\), only a sparse subset of
\((\ell,m,\sigma)\) couplings carries export gain above tolerance: geometry and symmetry impose \emph{selection rules} on
\(\boldsymbol{\Gamma}_\omega\) before any design optimization begins \cite{hansen1988,devaney2004}.

\paragraph{Assumptions.}
Let \(\Jimp\in\mathcal{J}_{\mathrm{real}}\) on compact \(\Omega_{s}\), with Green representation
Eq.~\eqref{eq:ch2.electric-dyadic-convolution} and far-field reduction
Theorem~\ref{thm:ch2.far-field-approximation} \cite{stratton1941,jackson1999}. Flux-normalized eigenfields
\(\widehat{\mathbf{u}}_{\ell m}^{(\sigma)}\) and degeneracy structure
Eq.~\eqref{eq:ch3.eigenfield-Wigner-mixing} enter from Section~\ref{sec:ch34} and
Subsection~\ref{sec:ch345} \cite{hansen1988}. Radiation bilinear
Eq.~\eqref{eq:ch2.radiative-coupling-bilinear} from Subsection~\ref{sec:ch282} supplies the adjoint pairing
\cite{stratton1941,harrington2001}.

\paragraph{Modal coupling integral.}
Define the \emph{modal coupling functional} by
\begin{equation}
\label{eq:ch4.modal-coupling-integral}
\mathcal{K}_{\ell m}^{(\sigma)}[\Jimp]
=
\int_{\Omega_{s}}
\Jimp(\mathbf{r})\cdot
\widehat{\mathbf{J}}_{\ell m}^{(\sigma)\,\dagger}(\mathbf{r})\,
d^{3}r,
\end{equation}
where \(\widehat{\mathbf{J}}_{\ell m}^{(\sigma)\,\dagger}\) is the adjoint source mode dual to
\(\widehat{\mathbf{u}}_{\ell m}^{(\sigma)}\) under Eq.~\eqref{eq:ch2.radiative-coupling-bilinear}
\cite{stratton1941,harrington2001}. In the far-zone regime,
\begin{equation}
\label{eq:ch4.coupling-far-zone-reduction}
\widehat{c}_{\ell m}^{(\sigma)}[\Jimp]
=
\mathcal{Z}_{\ell m}^{(\sigma)}\,
\mathcal{K}_{\ell m}^{(\sigma)}[\Jimp],
\end{equation}
with \(\mathcal{Z}_{\ell m}^{(\sigma)}\) absorbing propagation constants and flux normalization from
Section~\ref{sec:ch35} \cite{hansen1988,stratton1941}. Equation~\eqref{eq:ch4.coupling-far-zone-reduction} is the
realizability reading of Eq.~\eqref{eq:ch3.far-pattern-eigenfield-coupling}: each coefficient is a coupling integral,
not a fit parameter \cite{devaney2004,harrington2001}.

\begin{definition}[Modal coupling matrix]
\label{def:ch4.modal-coupling-matrix}
For a finite index set \(I_{\mathrm{acc}}\subset\{(\ell,m,\sigma)\}\) within \(\mathcal{B}_{\mathrm{fs}}\), the
\emph{modal coupling matrix} is
\begin{equation}
\label{eq:ch4.modal-coupling-matrix}
\mathbf{K}
=
\big[K_{(\ell m\sigma),(\ell'm'\sigma')}\big],
\qquad
K_{\alpha\beta}
=
\big\langle
\widehat{\mathbf{u}}_{\alpha},\,
\mathcal{O}\,\widehat{\mathbf{u}}_{\beta}
\big\rangle_{\mathrm{rad}},
\end{equation}
for observation operator \(\mathcal{O}\) on \(\mathcal{F}_{\mathrm{rad}}\) and multi-indices
\(\alpha,\beta\in I_{\mathrm{acc}}\) \cite{harrington2001,stratton1941}. The synthesis operator satisfies
\(\boldsymbol{\Gamma}_\omega[\Jimp]=\mathbf{K}_{\mathrm{src}}\,\mathbf{j}\) when \(\mathbf{j}\) is the coefficient
vector of \(\Jimp\) in a declared source basis \cite{devaney2004}.
\end{definition}

Physically, \(\mathbf{K}\) is the bilinear report of how eigenfield sectors talk to one another under passive
observation: diagonal dominance signals sparse synthesizable patterns; off-diagonal weight signals mandatory multi-mode
superposition \cite{hansen1988,harrington2001}. Reciprocity on lossless media forces
\(K_{\alpha\beta}=K_{\beta\alpha}\) as in Eq.~\eqref{eq:ch3.reciprocity-coupling-symmetry} \cite{stratton1941}.

\paragraph{Selection rules from symmetry.}
When \(\mathcal{J}_{\mathrm{real}}\) is invariant under a subgroup \(G\subset\mathrm{SO}(3)\), coupling integrals
vanish unless representation labels compose consistently \cite{jackson1999,hansen1988}. For full rotational symmetry on
isotropic \(\Omega_{s}\), Wigner mixing Eq.~\eqref{eq:ch3.eigenfield-Wigner-mixing} implies
\begin{equation}
\label{eq:ch4.selection-delta-ell}
\mathcal{K}_{\ell m}^{(\sigma)}[\Jimp]
\ \text{supported on}\ 
\ell\le\ell_{\max}(k_{0}D_{s}),
\end{equation}
importing Eq.~\eqref{eq:ch4.high-ell-decay} as a selection bound, and
\begin{equation}
\label{eq:ch4.selection-parity}
\mathcal{K}_{\ell m}^{(\mathbf{M})}\ \text{and}\ \mathcal{K}_{\ell m}^{(\mathbf{N})}\ \text{decouple}
\quad\text{when}\ \mathfrak{E}\ \text{preserves parity},
\end{equation}
on centrosymmetric feeds \cite{stratton1941,hansen1988}. Equations~\eqref{eq:ch4.selection-delta-ell}--\eqref{eq:ch4.selection-parity}
are \emph{coupling selection rules}: they shrink the set of nonzero couplings before MoM discretization
\cite{harrington2001,devaney2004}.

\begin{theorem}[Coupling selection for realizable synthesis]
\label{thm:ch4.coupling-selection-rules}
Let \(\Jimp\in\mathcal{J}_{\mathrm{real}}\) on compact \(\Omega_{s}\) with \(\mathcal{D}_{\mathrm{cpl}}=1\). Then
\(\widehat{c}_{\ell m}^{(\sigma)}[\Jimp]=0\) whenever any of the following holds:
\begin{equation}
\label{eq:ch4.selection-rules-list}
(\ell,m,\sigma)\notin\mathcal{B}_{\mathrm{fs}},
\qquad
\sigma\ \text{forbidden by parity of }\mathfrak{E},
\qquad
|\widehat{c}_{\ell m}^{(\sigma)}|
<
\varepsilon\,\|\boldsymbol{\Gamma}_\omega[\Jimp]\|
\ \text{with}\ \sigma_{\mathrm{rank}}<\varepsilon\sigma_{1},
\end{equation}
where the third condition is the singular-value selection rule from
Eq.~\eqref{eq:ch4.gamma-SVD} \cite{stratton1941,hansen1988,harrington2001,devaney2004}.
\end{theorem}
\begin{proof}
First condition: Definition~\ref{def:ch4.finite-support-restriction} \cite{hansen1988}. Second: degeneracy and parity
structure of Subsection~\ref{sec:ch345} applied to \(\mathfrak{E}\) \cite{jackson1999,stratton1941}. Third:
Eckart--Young truncation of \(\boldsymbol{\Gamma}_\omega\) at \(r_{\mathrm{rad}}(\varepsilon)\)
Theorem~\ref{thm:ch4.radiative-rank} \cite{harrington2001,devaney2004}.
\end{proof}

Theorem~\ref{thm:ch4.coupling-selection-rules} is the inevitability statement of Subsection~\ref{sec:ch452}: selection is
not an optional sparse prior but the operator-theoretic consequence of compact support, symmetry, and radiative rank
\cite{devaney2004,hansen1988}. Pattern optimizers that activate forbidden \((\ell,m,\sigma)\) sectors optimize outside
\(\mathcal{S}_{\mathrm{real}}(\mathcal{C})\) even when MoM matrices are formally invertible \cite{harrington2001}.

\begin{figure}[t]
\centering
\ChOneFigBlock{%
\begin{tikzpicture}[ch1fig, node distance=7mm and 9mm]
  \node[ch1box] (j) {$\Jimp$};
  \node[ch1box, right=12mm of j] (k) {Eq.~\eqref{eq:ch4.modal-coupling-integral}};
  \node[ch1box, right=12mm of k] (sel) {Eq.~\eqref{eq:ch4.selection-rules-list}};
  \node[ch1box, right=12mm of sel] (c) {$\widehat{c}_{\ell m}^{(\sigma)}$};
  \draw[ch1flow] (j) -- (k) -- (sel) -- (c);
\end{tikzpicture}%
}
\caption{Subsection~\ref{sec:ch452}: coupling integrals map sources to eigenfield weights; selection rules zero
forbidden sectors before synthesis.}
\label{fig:ch4.coupling-selection-chain}
\end{figure}

\paragraph{Composite-label coupling (pointer).}
When products of angular sectors form composite labels, Clebsch--Gordan weights
Eq.~\eqref{eq:ch3.CG-coupling-def} govern which bilinear couplings may appear in
\(\mathbf{K}\) \cite{jackson1999,hansen1988}. Subsection~\ref{sec:ch452} cites those weights without reopening
addition-theorem proofs: the selection content is that coupling integrals inherit the same triangular
\((\ell_1,\ell_2,\ell_{3})\) admissibility as scalar harmonic products \cite{hansen1988,stratton1941}. Composite
OAM claims must therefore be checked against both Eq.~\eqref{eq:ch4.selection-rules-list} and CG-admissible coupling
topology \cite{devaney2004}.

\begin{proposition}[Block sparsity under \(\mathrm{SO}(3)\) symmetry]
\label{prop:ch4.so3-block-sparsity}
If \(\mathfrak{E}\) is rotationally symmetric on isotropic \(\Omega_{s}\), then \(\mathbf{K}\) is block-diagonal in
\((\ell,\sigma)\) up to \(m\)-mixing inside each \(\ell\) block through
Eq.~\eqref{eq:ch3.eigenfield-Wigner-mixing} \cite{jackson1999,hansen1988,stratton1941}. Off-block couplings vanish in
the continuum limit; MoM discretization may introduce numerical leakage controlled by \(\mathcal{D}_{\mathrm{cpl}}\)
\cite{harrington2001,devaney2004}.
\end{proposition}

\paragraph{Surface and aperture reduction.}
When \(\Jimp\) is replaced by \(\mathbf{J}_{s}\) on \(\Sigma_{a}\) per Section~\ref{sec:ch42},
Eq.~\eqref{eq:ch4.modal-coupling-integral} reduces to
\begin{equation}
\label{eq:ch4.aperture-coupling-integral}
\mathcal{K}_{\ell m}^{(\sigma)}[\mathbf{J}_{s}]
=
\int_{\Sigma_{a}}
\mathbf{J}_{s}(\mathbf{r})\cdot
\widehat{\mathbf{J}}_{\ell m}^{(\sigma)\,\dagger}(\mathbf{r})\,
dS,
\end{equation}
with \(\mathcal{D}_{\tau}=1\) and \(\mathcal{D}_{\mathrm{cpl}}=1\) \cite{stratton1941,harrington2001,collin2001}.
Equation~\eqref{eq:ch4.aperture-coupling-integral} is how horn and slot designs enter eigenfield coordinates: aperture
currents are integrated against adjoint modes, not against unconstrained \(\E_{t}\) templates \cite{collin2001,devaney2004}.

\paragraph{Worked illustration (parity selection).}
A centrosymmetric volume current on \(B_{a}\) excites primarily \(\sigma=\mathbf{N}\) sectors at odd \(\ell\) and
\(\sigma=\mathbf{M}\) at even \(\ell\) under Eq.~\eqref{eq:ch4.selection-parity} \cite{stratton1941,jackson1999}.
Demanding strong \(\ell=2,\sigma=\mathbf{N}\) content from a parity-symmetric \(\Jimp\) contradicts
Theorem~\ref{thm:ch4.coupling-selection-rules} \cite{hansen1988,harrington2001}.

\paragraph{Adjoint source modes.}
The dual modes \(\widehat{\mathbf{J}}_{\ell m}^{(\sigma)\,\dagger}\) in Eq.~\eqref{eq:ch4.modal-coupling-integral} are
fixed by the radiation adjoint of Subsection~\ref{sec:ch282}:
\begin{equation}
\label{eq:ch4.adjoint-source-mode}
\widehat{\mathbf{J}}_{\ell m}^{(\sigma)\,\dagger}
=
\mathcal{R}_\omega^{\dagger}\big[\widehat{\mathbf{u}}_{\ell m}^{(\sigma)}\big]
\quad\text{on}\ \Omega_{s},
\end{equation}
with \(\mathcal{R}_\omega^{\dagger}\) the adjoint induced by
Eq.~\eqref{eq:ch2.radiative-coupling-bilinear} \cite{stratton1941,harrington2001}. Equation
\eqref{eq:ch4.adjoint-source-mode} makes coupling integrals bilinear audits: changing \(\mathcal{C}\) rotates
\(\widehat{\mathbf{u}}_{\ell m}^{(\sigma)}\) and therefore rotates \(\widehat{\mathbf{J}}_{\ell m}^{(\sigma)\,\dagger}\)
consistently \cite{hansen1988}. Mismatch between transmit and receive adjoint modes breaks reciprocity symmetry in
\(\mathbf{K}\) without changing \(\operatorname{ran}\boldsymbol{\Gamma}_\omega\) \cite{harrington2001,stratton1941}.

\paragraph{Discrete coupling matrix assembly.}
On patches \(\{S_{p}\}\), assemble
\begin{equation}
\label{eq:ch4.discrete-coupling-matrix}
K_{\alpha p}
=
\int_{S_{p}}
\boldsymbol{\Lambda}_{p}(\mathbf{r})\cdot
\widehat{\mathbf{J}}_{\alpha}^{\dagger}(\mathbf{r})\,dS,
\qquad
\mathbf{c}
=
\mathbf{K}_{\mathrm{disc}}\,\mathbf{I}_{s},
\end{equation}
with \(\mathbf{I}_{s}\) patch currents \cite{harrington2001,collin2001}. Equation~\eqref{eq:ch4.discrete-coupling-matrix}
is the aperture specialization of Eq.~\eqref{eq:ch4.mom-eigenfield-matrix}: selection rules appear as numerically
small rows of \(\mathbf{K}_{\mathrm{disc}}\) before inversion \cite{stratton1941,devaney2004}. Rows below
\(\varepsilon\sigma_{1}\) should be zeroed by Theorem~\ref{thm:ch4.coupling-selection-rules}, not by arbitrary
Tikhonov priors \cite{harrington2001}.

\begin{corollary}[Selection implies sparse \(\mathbf{K}\)]
\label{cor:ch4.sparse-coupling-matrix}
Under isotropic \(\mathrm{SO}(3)\) symmetry on \(B_{a}\), \(\mathbf{K}\) is block-sparse in \((\ell,\sigma)\) up to
numerical leakage \(\mathcal{O}(\varepsilon_{\mathrm{mach}})\) when \(\mathcal{D}_{\mathrm{cpl}}=1\)
\cite{jackson1999,hansen1988,harrington2001}. Sparsity is a consequence of symmetry and compactness, not an optional
compressive prior \cite{devaney2004}.
\end{corollary}

\paragraph{Worked illustration (slot TE coupling).}
A narrow slot in a PEC plane excites primarily \(\ell=1\) and odd-\(\ell\) \(\mathbf{N}\) sectors at low
\(k_{0}D_{s}\) \cite{collin2001,stratton1941}. Coupling rows with \(\ell\ge3\) fall below
\(\varepsilon\sigma_{1}\) on electrically short slots; MoM inversion must not reanimate them by regularization alone
\cite{harrington2001,devaney2004}.

\paragraph{Problem (forbidden coupling diagnosis).}
\textbf{Problem.} MoM inversion returns \(\widehat{c}_{5,3}^{(\mathbf{M})}\) on a compact patch with \(k_{0}D_{s}=0.6\).
Is the weight synthesizable?

\textbf{Step 1.} Test \(\ell=5\) against Eq.~\eqref{eq:ch4.high-ell-decay} \cite{hansen1988,jackson1999}.

\textbf{Step 2.} Check \(\sigma_{\mathrm{rank}}\) in Eq.~\eqref{eq:ch4.gamma-SVD} for the associated singular direction
\cite{harrington2001,devaney2004}.

\textbf{Step 3.} If either test fails, \(\widehat{c}_{5,3}^{(\mathbf{M})}\) is a formal fit outside
\(\mathcal{S}_{\mathrm{real}}(\mathcal{C})\) \cite{stratton1941}.

\textbf{Physical meaning.} Selection rules are synthesis predicates, not post-fit sparsity priors \cite{hansen1988}.

\paragraph{Parity and azimuthal selection.}
On bodies with azimuthal symmetry about \(\hat{\mathbf{z}}\), coupling rows with \(|m|>m_{\max}(k_{0}D_{s})\) fall below
\(\varepsilon\sigma_{1}\) in Eq.~\eqref{eq:ch4.gamma-SVD} \cite{jackson1999,hansen1988,stratton1941}. Parity selection
requires
\begin{equation}
\label{eq:ch4.parity-selection}
\mathcal{K}_{\ell m}^{(\sigma)}[\Jimp]=0
\quad\text{when}\quad
\mathcal{P}\widehat{\mathbf{u}}_{\ell m}^{(\sigma)}
=
-\widehat{\mathbf{u}}_{\ell m}^{(\sigma)}
\ \text{and}\ 
\mathcal{P}\Jimp=\Jimp,
\end{equation}
for symmetry operator \(\mathcal{P}\) of the feed \cite{hansen1988,harrington2001,devaney2004}. Equation
\eqref{eq:ch4.parity-selection} is the synthesis reading of Wigner \(D\)-matrix block sparsity imported from
Section~\ref{sec:ch36} \cite{jackson1999}.

\begin{theorem}[Symmetry-induced coupling nulls]
\label{thm:ch4.symmetry-coupling-nulls}
If \(\mathcal{G}\) is a symmetry group of \((\Omega_{s},\Jimp)\) and \(\widehat{\mathbf{u}}_{\ell m}^{(\sigma)}\) transforms by
irrep \(\chi_{\ell m}^{(\sigma)}\), then \(\mathcal{K}_{\ell m}^{(\sigma)}[\Jimp]=0\) for every sector whose irrep does not
appear in the decomposition of \(\mathcal{J}_{\mathrm{real}}\) under \(\mathcal{G}\) \cite{hansen1988,stratton1941,harrington2001}.
\end{theorem}
\begin{proof}
Schur orthogonality on coupling integrals Eq.~\eqref{eq:ch4.modal-coupling-integral} \cite{jackson1999,hansen1988}.
\end{proof}

\paragraph{Worked illustration (axial dipole).}
An axial dipole on \(\hat{\mathbf{z}}\) nulls sectors with odd parity about the equatorial plane: reported
\(\widehat{c}_{2,1}^{(\mathbf{M})}\) from MoM inversion is a symmetry violation unless the mesh breaks \(\mathcal{P}\)
\cite{jackson1999,stratton1941,harrington2001}.

\paragraph{Problem (forbidden \(m\) sector).}
\textbf{Problem.} A cylindrical feed claims dominant \(\widehat{c}_{4,3}^{(\mathbf{N})}\) but feed symmetry allows only
\(m=0,\pm1\). Is the claim realizable?

\textbf{Step 1.} Apply Theorem~\ref{thm:ch4.symmetry-coupling-nulls} \cite{hansen1988}.

\textbf{Step 2.} Test \(|\mathcal{K}_{4,3}^{(\mathbf{N})}|\) against \(\varepsilon\sigma_{1}\) \cite{harrington2001,devaney2004}.

\textbf{Step 3.} Reject the sector or revise the declared symmetry class of \(\mathcal{J}_{\mathrm{real}}\) \cite{stratton1941}.

\textbf{Physical meaning.} Selection rules are symmetry theorems, not optional sparsity \cite{hansen1988}.

\paragraph{Validity.}
Eqs.~\eqref{eq:ch4.modal-coupling-integral}--\eqref{eq:ch4.selection-rules-list} assume linear Maxwell theory and
outgoing \(\mathcal{F}_{\mathrm{rad}}\) \cite{stratton1941}. Broken rotational symmetry (arrays, waveguides) replaces
full \(\mathrm{SO}(3)\) selection by partial block sparsity; rules must be recomputed in the symmetry-adapted basis of
Section~\ref{sec:ch36} \cite{harrington2001,collin2001}. Magnetic currents and combined sheets use the same integral
pattern with dual pairings \cite{stratton1941}.

Geometry dependence of coupling strength is quantified in Subsection~\ref{sec:ch453}.


\subsection{Dependence on geometry and aperture size}
\label{sec:ch453}

Selection rules identify which eigenfield couplings may be nonzero; geometry identifies how large those couplings can
be \cite{stratton1941,hansen1988,devaney2004}. Electrical size \(\Xi=k_{0}D_{s}/(2\pi)\) and aperture diameter enter
\(\mathbf{K}\) as multiplicative rank constraints on \(\boldsymbol{\Gamma}_\omega\), not as post hoc tuning parameters
\cite{harrington2001,jackson1999}.

\paragraph{Assumptions.}
Fix compact \(\Omega_{s}\) or aperture \(\Sigma_{a}\) with \(\operatorname{diam}=D_{s}\), flux-normalized eigenfields,
and coupling matrix Eq.~\eqref{eq:ch4.modal-coupling-matrix} restricted to \(I_{\mathrm{acc}}\subset\mathcal{B}_{\mathrm{fs}}\)
\cite{stratton1941,hansen1988}. Presuppose \(r_{\mathrm{rad}}(\varepsilon)\) and
Eq.~\eqref{eq:ch4.weyl-mode-count} \cite{harrington2001,devaney2004}.

\paragraph{Geometry-weighted coupling operator.}
Define the \emph{geometry-weighted coupling operator} by
\begin{equation}
\label{eq:ch4.geometry-coupling-operator}
\boldsymbol{\mathcal{K}}_{\mathrm{geo}}(\Omega_{s})
:
\mathcal{J}_{\mathrm{real}}(\Omega_{s})
\longrightarrow
\mathbb{C}^{I_{\mathrm{acc}}},
\qquad
\boldsymbol{\mathcal{K}}_{\mathrm{geo}}[\Jimp]
=
\big\{\mathcal{K}_{\ell m}^{(\sigma)}[\Jimp]\big\}_{(\ell,m,\sigma)\in I_{\mathrm{acc}}},
\end{equation}
with \(\mathcal{K}_{\ell m}^{(\sigma)}\) from Eq.~\eqref{eq:ch4.modal-coupling-integral} \cite{stratton1941,harrington2001}.
Composition with normalization gives \(\boldsymbol{\Gamma}_\omega=\mathbf{Z}\,\boldsymbol{\mathcal{K}}_{\mathrm{geo}}\) for
diagonal \(\mathbf{Z}\) in the aligned flux gauge \cite{hansen1988}. Operator norm
\(\|\boldsymbol{\mathcal{K}}_{\mathrm{geo}}\|_{\mathrm{op}}\) grows with \(\operatorname{vol}(\Omega_{s})\) and
\(k_{0}D_{s}\) through Eq.~\eqref{eq:ch4.HS-bound} \cite{stratton1941,devaney2004}.

\begin{theorem}[Geometry scaling of coupling rank]
\label{thm:ch4.geometry-coupling-scaling}
For ball support \(B_{a}\) with \(k_{0}a\ge1\) and tolerance \(\varepsilon\in(0,1)\), the number of eigenfield
couplings above threshold satisfies
\begin{equation}
\label{eq:ch4.geometry-coupling-scaling}
\#\big\{(\ell,m,\sigma):\ |\mathcal{K}_{\ell m}^{(\sigma)}[\Jimp]|\ge\varepsilon\|\boldsymbol{\mathcal{K}}_{\mathrm{geo}}[\Jimp]\|\big\}
\ \le\
C\,(k_{0}a)^{3},
\end{equation}
with \(C\) depending only on \(\varepsilon\) and flux gauge \cite{stratton1941,jackson1999,hansen1988}. For aperture
\(\Sigma_{a}\) of area \(A_{a}\),
\begin{equation}
\label{eq:ch4.aperture-coupling-scaling}
\operatorname{rank}\mathbf{K}\big|_{\Sigma_{a}}
\ \le\
C'\,k_{0}^{2}A_{a},
\end{equation}
linking patch rank to electrical area \cite{collin2001,harrington2001,devaney2004}.
\end{theorem}
\begin{proof}
Inequality Eq.~\eqref{eq:ch4.geometry-coupling-scaling} combines
Eq.~\eqref{eq:ch4.weyl-mode-count} with radiative rank Theorem~\ref{thm:ch4.radiative-rank}
\cite{stratton1941,hansen1988}. Aperture bound Eq.~\eqref{eq:ch4.aperture-coupling-scaling} follows from
MoM rank Eq.~\eqref{eq:ch4.mom-eigenfield-matrix} and standard aperture sampling limits
\cite{harrington2001,collin2001}.
\end{proof}

Theorem~\ref{thm:ch4.geometry-coupling-scaling} makes geometry-imposed rank explicit at the coupling layer: enlarging
\(a\) or \(A_{a}\) increases the count of potentially strong eigenfield couplings, but fixed geometry caps that count at
fixed \((k_{0},\varepsilon)\) \cite{devaney2004,hansen1988}. Phase-only retuning on fixed \(\Sigma_{a}\) rotates
couplings within the same rank class; it does not raise \(\operatorname{rank}\mathbf{K}\) \cite{stratton1941,harrington2001}.

\paragraph{Coupling decay with \(\ell\).}
For \(\ell>k_{0}D_{s}\),
\begin{equation}
\label{eq:ch4.coupling-ell-decay}
\sup_{\Jimp\in\mathcal{J}_{\mathrm{real}}}
\frac{|\mathcal{K}_{\ell m}^{(\sigma)}[\Jimp]|}
{\|\boldsymbol{\mathcal{K}}_{\mathrm{geo}}[\Jimp]\|}
\le
\mathcal{C}_{\mathrm{cpl}}(\ell,k_{0}D_{s})\,
\ell^{-1}\exp(-\ell/k_{0}D_{s}),
\end{equation}
importing Eq.~\eqref{eq:ch4.high-ell-decay} at the coupling-integral layer \cite{hansen1988,stratton1941,jackson1999}.
Equation~\eqref{eq:ch4.coupling-ell-decay} is the physical meaning of spectral compression for coupling design: high-\(\ell\)
channels are not merely weak in power---their coupling integrals are exponentially suppressed on compact support
\cite{harrington2001,devaney2004}.

\begin{figure}[t]
\centering
\ChOneFigBlock{%
\begin{tikzpicture}[ch1fig, node distance=7mm and 10mm]
  \node[ch1box] (geom) {$D_{s}$, $A_{a}$};
  \node[ch1box, right=12mm of geom] (k0) {$k_{0}$};
  \node[ch1box, right=12mm of k0] (rank) {Eq.~\eqref{eq:ch4.geometry-coupling-scaling}};
  \node[ch1box, right=12mm of rank] (kmat) {$\mathbf{K}$};
  \draw[ch1flow] (geom) -- (rank);
  \draw[ch1flow] (k0) -- (rank);
  \draw[ch1flow] (rank) -- (kmat);
\end{tikzpicture}%
}
\caption{Subsection~\ref{sec:ch453}: geometry and frequency set coupling rank before modal weights are read.}
\label{fig:ch4.geometry-coupling-scaling}
\end{figure}

Figure~\ref{fig:ch4.geometry-coupling-scaling} connects Section~\ref{sec:ch44} bandlimits to coupling engineering:
\(k_{t,\max}(\eta)\) and \(\ell_{\max}\) bound spectral support; Eqs.~\eqref{eq:ch4.geometry-coupling-scaling}--\eqref{eq:ch4.aperture-coupling-scaling}
bound how many couplings can be strong enough to matter at tolerance \(\varepsilon\) \cite{hansen1988,devaney2004}.
Corollary~\ref{cor:ch4.support-monotone-enlargement} is the set inclusion version of the same fact
\cite{stratton1941,harrington2001}.

\paragraph{Near-field versus far-zone coupling.}
Coupling integrals evaluated in the reactive shell may include evanescent weight absent from
\(\widehat{c}_{\ell m}^{(\sigma)}\) on \(\mathcal{F}_{\mathrm{rad}}\) \cite{stratton1941,harrington2001}. Denote by
\(\mathcal{K}_{\ell m}^{(\sigma,\mathrm{near})}\) the near evaluation and by
\(\mathcal{K}_{\ell m}^{(\sigma,\mathrm{far})}\) the far-zone reduction
Eq.~\eqref{eq:ch4.coupling-far-zone-reduction} \cite{devaney2004}. Then
\begin{equation}
\label{eq:ch4.near-far-coupling-split}
\mathcal{K}_{\ell m}^{(\sigma,\mathrm{near})}
=
\mathcal{K}_{\ell m}^{(\sigma,\mathrm{far})}
+
\mathcal{K}_{\ell m}^{(\sigma,\mathrm{ev})},
\qquad
\boldsymbol{\Gamma}_\omega\ \text{depends only on}\ 
\mathcal{K}_{\ell m}^{(\sigma,\mathrm{far})},
\end{equation}
by Proposition~\ref{prop:ch4.evanescent-invisibility} \cite{stratton1941,harrington2001}. Equation
\eqref{eq:ch4.near-far-coupling-split} is the coupling form of the propagating-sector restriction
Eq.~\eqref{eq:ch4.propagating-realizability}: near-field scans overweight
\(\mathcal{K}_{\ell m}^{(\sigma,\mathrm{ev})}\) \cite{devaney2004}.

\begin{corollary}[Far-zone coupling determines synthesis]
\label{cor:ch4.far-coupling-synthesis}
\(\mathbf{c}\in\mathcal{S}_{\mathrm{real}}(\mathcal{C})\) if and only if the far-zone coupling vector
\(\{\mathcal{K}_{\ell m}^{(\sigma,\mathrm{far})}[\Jimp]\}\) lies in \(\operatorname{ran}\boldsymbol{\mathcal{K}}_{\mathrm{geo}}\)
at fixed \((\Omega_{s},\omega)\) \cite{stratton1941,harrington2001,devaney2004}.
\end{corollary}

\paragraph{Aperture versus volume feeds.}
Volume \(\Jimp\) on \(\Omega_{s}\) and aperture \(\mathbf{J}_{s}\) on \(\Sigma_{a}\) produce identical far-zone
\(\mathbf{c}\) when equivalence Theorem~\ref{thm:ch4.volume-surface-equivalence} holds, but
\(\boldsymbol{\mathcal{K}}_{\mathrm{geo}}\) has different discrete rank on \(\Omega_{s}\) versus \(\Sigma_{a}\)
\cite{collin2001,harrington2001}. MoM on \(\Sigma_{a}\) yields
\(\operatorname{rank}\mathbf{T}_{\mathrm{VSH}}\le P\) patches; volume discretizations may expose additional synthesis
directions at the same \(\mathbf{c}\) when containment fails \cite{devaney2004,stratton1941}.

\paragraph{Worked illustration (horn aperture).}
A conical horn with mouth diameter \(D_{h}\) at \(k_{0}D_{h}=4\) supports roughly \((k_{0}D_{h})^{3}\) strong eigenfield
couplings by Eq.~\eqref{eq:ch4.geometry-coupling-scaling} \cite{collin2001,jackson1999}. Targeting \(\ell=12\) content
with a mouth that electrically supports only \(\ell\le6\) violates Eq.~\eqref{eq:ch4.coupling-ell-decay}
\cite{hansen1988,harrington2001}.

\paragraph{Electrical-size ladder.}
Define the coupling ladder
\begin{equation}
\label{eq:ch4.coupling-ladder}
\mathcal{L}_{n}
=
\big\{(\ell,m,\sigma):\ \ell\le n,\ s_{\alpha}\ge\varepsilon\big\},
\qquad
\mathcal{L}_{n}\subseteq\mathcal{L}_{n+1},
\end{equation}
with strict inclusion under Corollary~\ref{cor:ch4.support-monotone-enlargement} when \(D_{s}\) increases
\cite{stratton1941,hansen1988,jackson1999}. Equation~\eqref{eq:ch4.coupling-ladder} orders design freedom by electrical
size: at fixed \(\omega\), enlarging \(\Omega_{s}\) activates higher rungs of the ladder; at fixed \(\Omega_{s}\), raising
\(\omega\) activates higher \(\ell\) before aperture measurement compresses further \cite{devaney2004,harrington2001}.
Phase-only synthesis on fixed hardware explores only \(\mathcal{L}_{n}\) at fixed \(n\), not \(\mathcal{L}_{n+k}\)
\cite{stratton1941}.

\paragraph{Anisotropic and elongated supports.}
For prolate supports with aspect ratio \(\chi>1\), coupling rank along the major axis exceeds that along the minor axis
at the same \(\operatorname{vol}(\Omega_{s})\) \cite{stratton1941,jackson1999,hansen1988}. The geometry operator
\(\boldsymbol{\mathcal{K}}_{\mathrm{geo}}\) is not isotropic in \(\ell,m\): aligned \(\ell\) sectors may activate before
azimuthal partners with the same \(\ell\) \cite{harrington2001}. Equation~\eqref{eq:ch4.geometry-coupling-scaling}
remains a conservative isotropic bound; elongated feeds require directionally resolved audits of \(\mathbf{K}\)
\cite{devaney2004}.

\begin{lemma}[Coupling rank monotonicity in frequency]
\label{lem:ch4.coupling-freq-monotone}
At fixed \(\Omega_{s}\), \(r_{\mathrm{rad}}(\varepsilon)\) and the cardinality of \(\mathcal{L}_{n}\) are
non-decreasing in \(k_{0}\) until propagation cutoffs or material dispersion alter \(\mathcal{J}_{\mathrm{real}}\)
\cite{stratton1941,jackson1999,hansen1988}. Frequency synthesis therefore enlarges coupling rank monotonically on fixed
geometry \cite{harrington2001,devaney2004}.
\end{lemma}

\paragraph{Worked illustration (patch array).}
A \(4\times4\) patch array on \(\lambda/2\) spacing at \(k_{0}a=3\) yields \(\operatorname{rank}\mathbf{K}_{\mathrm{disc}}\)
near \(16\) before continuum compression; \(r_{\mathrm{rad}}(\varepsilon)\) may be \(\mathcal{O}(10^{2})\) on the full
volume enclosing the array \cite{harrington2001,collin2001}. Reporting \(\ell>20\) weights from the array alone without
enclosing-volume synthesis violates Eq.~\eqref{eq:ch4.aperture-coupling-scaling} \cite{devaney2004,hansen1988}.

\paragraph{Problem (rank versus electrical size).}
\textbf{Problem.} A square aperture \(A_{a}=(0.5\lambda)^{2}\) uses \(P=120\) patches. Estimate maximum reliable eigenfield
rank at \(\varepsilon=10^{-2}\).

\textbf{Step 1.} Compute \(k_{0}^{2}A_{a}\approx\pi^{2}\) \cite{jackson1999,stratton1941}.

\textbf{Step 2.} Apply Eq.~\eqref{eq:ch4.aperture-coupling-scaling} for order-of-magnitude rank \(\sim 10\)--\(30\)
\cite{harrington2001,collin2001}.

\textbf{Step 3.} Compare to \(P=120\): excess patches do not increase export rank above continuum
\(r_{\mathrm{rad}}(\varepsilon)\) \cite{devaney2004,hansen1988}.

\textbf{Physical meaning.} Patch count overstates coupling rank when geometry is electrically small \cite{stratton1941}.

\paragraph{Electrical-size crossover regimes.}
Coupling rank crosses three regimes as \(k_{0}D_{s}\) increases:
\begin{equation}
\label{eq:ch4.coupling-crossover}
\underbrace{r_{\mathrm{rad}}\sim\mathcal{O}(1)}_{\text{Rayleigh}}
\ \longrightarrow\
\underbrace{r_{\mathrm{rad}}\sim(k_{0}D_{s})^{2}}_{\text{mesoscale}}
\ \longrightarrow\
\underbrace{r_{\mathrm{rad}}\sim(k_{0}D_{s})^{3}}_{\text{large-body}},
\end{equation}
importing Eq.~\eqref{eq:ch4.geometry-coupling-scaling} and standard aperture limits
\cite{stratton1941,jackson1999,hansen1988,harrington2001}. Inverse design must declare which regime the hardware occupies:
phase-only templates in the Rayleigh regime cannot access mesoscale rank without enlarging \(D_{s}\) or raising \(k_{0}\)
\cite{devaney2004}.

\begin{proposition}[Crossover monotonicity]
\label{prop:ch4.crossover-monotone}
At fixed \(\Omega_{s}\), \(r_{\mathrm{rad}}(\varepsilon)\) and \(\#\mathcal{L}_{n}\) are non-decreasing in \(k_{0}D_{s}\) until
material dispersion or superconducting cutoffs alter \(\mathcal{J}_{\mathrm{real}}\) \cite{stratton1941,jackson1999,hansen1988}.
Enlarging \(D_{s}\) at fixed \(k_{0}\) produces the same monotonicity by
Corollary~\ref{cor:ch4.support-monotone-enlargement} \cite{harrington2001,devaney2004}.
\end{proposition}

\paragraph{Worked illustration (Ka-band patch on handset).}
A \(10\,\mathrm{mm}\) patch at \(30\,\mathrm{GHz}\) has \(k_{0}D_{s}\approx2\): mesoscale scaling applies, not Rayleigh
dipole-only rank \cite{harrington2001,collin2001,jackson1999}. Claiming \(\ell\le1\) dominance contradicts
Eq.~\eqref{eq:ch4.coupling-crossover} \cite{hansen1988,devaney2004}.

\paragraph{Problem (regime mis-identification).}
\textbf{Problem.} A design report uses dipole formulas for a \(k_{0}D_{s}=4\) array. What rank error follows?

\textbf{Step 1.} Identify large-body regime in Eq.~\eqref{eq:ch4.coupling-crossover} \cite{jackson1999,stratton1941}.

\textbf{Step 2.} Replace \(\mathcal{O}(1)\) rank estimate with \(\mathcal{O}((k_{0}D_{s})^{3})\) bound \cite{hansen1988,harrington2001}.

\textbf{Step 3.} Re-audit \(\mathbf{c}_{\mathrm{target}}\) against \(r_{\mathrm{rad}}(\varepsilon)\) \cite{devaney2004}.

\textbf{Physical meaning.} Electrical-size regime sets the scaling exponent, not merely a numeric prefactor \cite{stratton1941}.

\paragraph{Validity.}
Theorem~\ref{thm:ch4.geometry-coupling-scaling} is asymptotic for \(k_{0}D_{s}\gg1\) and conservative for
\(k_{0}D_{s}\ll1\) \cite{jackson1999,hansen1988}. Anisotropic and cavity environments modify \(C,C'\) without removing
rank saturation on fixed geometry \cite{harrington2001,stratton1941}.

Modal weights and the coupling strength ceiling are fixed in Subsection~\ref{sec:ch454}.


\subsection{Modal weights and coupling strength}
\label{sec:ch454}

Geometry sets how many eigenfield couplings may be strong; normalization and passivity set how large any single modal
weight may be without violating power balance on \(\mathcal{F}_{\mathrm{rad}}\) \cite{stratton1941,hansen1988,harrington2001}.
Subsection~\ref{sec:ch454} defines the \emph{modal weight ceiling}---the realizability bound on
\(|\widehat{c}_{\ell m}^{(\sigma)}|\) induced by flux gauge, reciprocity, and coupling admissibility
\cite{devaney2004}.

\paragraph{Assumptions.}
Presuppose flux normalization Eqs.~\eqref{eq:ch3.VSH-unit-flux-normalization}--\eqref{eq:ch3.mode-power-from-coefficient},
coupling matrix Eq.~\eqref{eq:ch4.modal-coupling-matrix}, and passive default
Eq.~\eqref{eq:ch4.passivity-screen} on \(\mathfrak{E}\) \cite{stratton1941,harrington2001}. Geometry rank bounds
Eqs.~\eqref{eq:ch4.geometry-coupling-scaling}--\eqref{eq:ch4.aperture-coupling-scaling} are in force
\cite{hansen1988,devaney2004}.

\paragraph{Modal weight and coupling strength.}
Define the \emph{modal weight} as \(w_{\ell m}^{(\sigma)}=|\widehat{c}_{\ell m}^{(\sigma)}|^{2}\) and the
\emph{coupling strength} to sector \(\alpha=(\ell,m,\sigma)\) as
\begin{equation}
\label{eq:ch4.coupling-strength-def}
s_{\alpha}
=
\frac{
|K_{\alpha\alpha}|
}{
\max_{\beta\in I_{\mathrm{acc}}}|K_{\beta\beta}|
},
\qquad
w_{\alpha}
\le
\mathcal{W}_{\mathrm{ceiling}}(\alpha;\mathfrak{E}),
\end{equation}
where \(\mathcal{W}_{\mathrm{ceiling}}\) is determined by feed power and passivity \cite{stratton1941,harrington2001}.
Equation~\eqref{eq:ch4.coupling-strength-def} separates \emph{geometry-limited coupling} \(s_{\alpha}\) from
\emph{power-limited weight} \(w_{\alpha}\): a sector may be strongly coupled yet weakly excited, or weakly coupled yet
dominant in a superdirective template that violates \(\mathcal{B}_{\mathrm{fs}}\) \cite{devaney2004,hansen1988}.

\begin{definition}[Modal weight ceiling]
\label{def:ch4.modal-weight-ceiling}
For declared feed tuple \(\mathfrak{E}\) with available complex power \(P_{\mathrm{feed}}\), the \emph{modal weight
ceiling} is
\begin{equation}
\label{eq:ch4.modal-weight-ceiling}
\mathcal{W}_{\mathrm{ceiling}}(\alpha;\mathfrak{E})
=
\sup\big\{
w_{\alpha}:
\mathbf{c}\in\mathcal{S}_{\mathrm{real}}(\mathcal{C}),\ 
\sum_{\beta}w_{\beta}\le P_{\mathrm{feed}}
\big\}.
\end{equation}
\end{definition}

Definition~\ref{def:ch4.modal-weight-ceiling} coins the realizability bound on per-mode power reports: no
\(\widehat{c}_{\ell m}^{(\sigma)}\) may claim more than \(\mathcal{W}_{\mathrm{ceiling}}\) watts in the flux gauge without
violating passivity or feed capability \cite{stratton1941,harrington2001}. The ceiling is independent of measurement rank;
\(\mathcal{E}_{\mathrm{acc}}^{(N)}\) may further shrink certifiable weights without changing
\(\mathcal{W}_{\mathrm{ceiling}}\) \cite{devaney2004,hansen1988}.

\paragraph{Normalization propagation to weights.}
Under gauge change \(\widehat{c}_{\alpha}\to\widehat{c}_{\alpha}/a_{\alpha}\),
Eq.~\eqref{eq:ch3.coupling-normalization-propagation} gives
\begin{equation}
\label{eq:ch4.weight-gauge-transform}
w_{\alpha}'=\frac{w_{\alpha}}{a_{\alpha}^{2}},
\qquad
s_{\alpha\beta}'=\frac{s_{\alpha\beta}}{a_{\alpha}a_{\beta}},
\end{equation}
while \(\mathcal{P}_{\mathrm{rad}}\) remains invariant under Eq.~\eqref{eq:ch3.power-budget-invariance}
\cite{hansen1988,harrington2001,stratton1941}. Equation~\eqref{eq:ch4.weight-gauge-transform} is the audit rule for
comparing modal weights across codes: mismatched normalization inflates or deflates \(w_{\alpha}\) without changing
exported watts \cite{harrington2001}.

\begin{theorem}[Coupling admissibility and weight ceiling]
\label{thm:ch4.coupling-weight-ceiling}
If \(\mathcal{D}_{\mathrm{cpl}}=0\) on the discretized source, then \(\mathbf{c}\notin\mathcal{S}_{\mathrm{real}}(\mathcal{C})\)
even when formal MoM inversion returns coefficients \cite{harrington2001,stratton1941}. If \(\mathcal{D}_{\mathrm{cpl}}=1\)
and \(\mathcal{D}_{\mathrm{real}}=1\), then
\begin{equation}
\label{eq:ch4.weight-ceiling-bound}
w_{\alpha}
\le
\min\Big\{
\mathcal{W}_{\mathrm{ceiling}}(\alpha;\mathfrak{E}),\ 
\mathcal{S}_{\mathrm{rank}}(\alpha;\varepsilon)
\Big\},
\end{equation}
where \(\mathcal{S}_{\mathrm{rank}}(\alpha;\varepsilon)=0\) when \(\sigma_{\alpha}<\varepsilon\sigma_{1}\) in
Eq.~\eqref{eq:ch4.gamma-SVD} \cite{devaney2004,hansen1988,harrington2001}.
\end{theorem}
\begin{proof}
Coupling inadmissibility implies \(\mathbf{T}\) is not Fredholm-stable, so \(\boldsymbol{\Gamma}_\omega\) has no stable
preimage \cite{harrington2001}. Passivity gives the power ceiling; singular-value threshold gives the rank ceiling from
Theorem~\ref{thm:ch4.radiative-rank} \cite{stratton1941,devaney2004}.
\end{proof}

Theorem~\ref{thm:ch4.coupling-weight-ceiling} closes the coupling arc: \(\mathcal{D}_{\mathrm{cpl}}\),
\(\mathcal{W}_{\mathrm{ceiling}}\), and \(r_{\mathrm{rad}}\) are independent screens whose intersection defines achievable
modal weights \cite{hansen1988,harrington2001}. Superdirective templates that concentrate \(w_{\alpha}\) in sectors with
\(s_{\alpha}\ll1\) typically violate Eq.~\eqref{eq:ch4.propagating-realizability} or \(\mathcal{B}_{\mathrm{fs}}\)
\cite{devaney2004,stratton1941}.

\begin{figure}[t]
\centering
\ChOneFigBlock{%
\begin{tikzpicture}[ch1fig, node distance=7mm and 8mm]
  \node[ch1box] (cpl) {$\mathcal{D}_{\mathrm{cpl}}$};
  \node[ch1box, right=11mm of cpl] (w) {$\mathcal{W}_{\mathrm{ceiling}}$};
  \node[ch1box, right=11mm of w] (r) {$r_{\mathrm{rad}}(\varepsilon)$};
  \node[ch1box, right=11mm of r] (wa) {$w_{\alpha}$};
  \draw[ch1flow] (cpl) -- (wa);
  \draw[ch1flow] (w) -- (wa);
  \draw[ch1flow] (r) -- (wa);
\end{tikzpicture}%
}
\caption{Subsection~\ref{sec:ch454}: modal weight \(w_{\alpha}\) is bounded by coupling admissibility, feed ceiling, and
radiative rank.}
\label{fig:ch4.modal-weight-ceiling}
\end{figure}

\paragraph{Superposition budget.}
For a realizable superposition \(\mathbf{c}=\sum_{\alpha\in I_{\mathrm{acc}}}a_{\alpha}\mathbf{e}_{\alpha}\),
passivity and Eq.~\eqref{eq:ch4.eigenfield-power-decomposition} require
\begin{equation}
\label{eq:ch4.superposition-power-budget}
\sum_{\alpha\in I_{\mathrm{acc}}}|a_{\alpha}|^{2}
\le
P_{\mathrm{feed}},
\qquad
w_{\alpha}=|a_{\alpha}|^{2}\ \text{in flux gauge},
\end{equation}
with equality only when all power is exported through \(\mathcal{F}_{\mathrm{rad}}\) \cite{stratton1941,poynting1884,harrington2001}.
Equation~\eqref{eq:ch4.superposition-power-budget} is the weight-space image of
Eq.~\eqref{eq:ch4.passivity-screen}: no combination of modal weights may exceed declared feed capability
\cite{harrington2001}. High-\(Q\) resonators store reactive energy that increases near-field couplings without raising
far-zone \(w_{\alpha}\) \cite{stratton1941,devaney2004}.

\begin{proposition}[Coupling--weight dominance ordering]
\label{prop:ch4.coupling-weight-ordering}
If \(s_{\alpha_{2}}\ll s_{\alpha_{1}}\) and \(w_{\alpha_{2}}\ge w_{\alpha_{1}}\) with
\(\mathbf{c}\in\mathcal{S}_{\mathrm{real}}(\mathcal{C})\), then necessarily
\(\mathcal{K}_{\alpha_{2}}^{(\mathrm{ev})}\) carries appreciable reactive storage and
\(\widehat{c}_{\alpha_{2}}\) is not determined by far-zone data alone \cite{harrington2001,devaney2004,stratton1941}.
\end{proposition}

\paragraph{Reciprocity constraint on weights.}
On reciprocal media, transmit and receive weights obey
Eq.~\eqref{eq:ch3.reciprocity-coupling-symmetry}: asymmetry in reported \(w_{\alpha}^{\mathrm{(tx)}}\) versus
\(w_{\alpha}^{\mathrm{(rx)}}\) signals normalization mismatch, not new synthesizable power \cite{stratton1941,harrington2001}.
Reciprocity pairs sectors; it does not enlarge \(\mathcal{S}_{\mathrm{real}}(\mathcal{C})\) beyond geometry rank
\cite{devaney2004}.

\paragraph{Worked illustration (two-mode superposition).}
Suppose only \(\alpha_{1},\alpha_{2}\) pass Theorem~\ref{thm:ch4.coupling-selection-rules} with
\(s_{\alpha_{1}}\approx1\), \(s_{\alpha_{2}}\approx10^{-2}\). Any template with \(w_{\alpha_{2}}\gg w_{\alpha_{1}}\) requires
near-field energy not exported to \(\mathcal{F}_{\mathrm{rad}}\) \cite{harrington2001,devaney2004}. Passivity caps
\(w_{\alpha_{1}}+w_{\alpha_{2}}\le P_{\mathrm{feed}}\) \cite{stratton1941}.

\paragraph{Accessibility versus synthesis ceilings.}
\(\mathcal{W}_{\mathrm{ceiling}}\) bounds synthesizable power per sector; \(\mathcal{E}_{\mathrm{acc}}^{(N)}\) bounds
certifiable coordinates at aperture rank \(N\) \cite{harrington2001,devaney2004}. A mode may satisfy
\(w_{\alpha}\le\mathcal{W}_{\mathrm{ceiling}}(\alpha)\) yet lie outside \(\mathcal{E}_{\mathrm{acc}}^{(N)}\) when the
receiver cannot resolve \(\alpha\) \cite{hansen1988}. Laboratory reports therefore require the intersection audit of
Eq.~\eqref{eq:ch4.realizable-accessible-intersection} applied to weights:
\begin{equation}
\label{eq:ch4.weight-accessibility-intersection}
w_{\alpha}^{\mathrm{(cert)}}
\le
\min\big\{
\mathcal{W}_{\mathrm{ceiling}}(\alpha),\ 
\mathcal{A}_{\mathrm{acc}}^{(N)}(\alpha)
\big\},
\end{equation}
with \(\mathcal{A}_{\mathrm{acc}}^{(N)}\) the per-sector certification ceiling induced by \(\mathcal{M}\)
\cite{devaney2004,harrington2001}.

\paragraph{Coupling strength matrix diagnostics.}
Define the normalized coupling matrix \(\widetilde{K}_{\alpha\beta}=K_{\alpha\beta}/\sqrt{K_{\alpha\alpha}K_{\beta\beta}}\)
on \(I_{\mathrm{acc}}\) \cite{harrington2001,hansen1988}. Large \(|\widetilde{K}_{\alpha\beta}|\) for \(\alpha\neq\beta\)
signals mandatory multi-mode synthesis; diagonal dominance signals sparse realizability in the eigenfield chart
\cite{stratton1941,devaney2004}. Thresholding \(|\widetilde{K}_{\alpha\beta}|<\varepsilon\) defines a graph sparsity
pattern compatible with Theorem~\ref{thm:ch4.coupling-selection-rules} \cite{hansen1988}.

\begin{corollary}[Diagonal dominance implies low synthesis rank]
\label{cor:ch4.diagonal-dominance-rank}
If \(|\widetilde{K}_{\alpha\beta}|<\varepsilon\) for all \(\alpha\neq\beta\) on \(I_{\mathrm{acc}}\), then
\(\mathbf{c}\) is approximated within \(\mathcal{O}(\varepsilon)\) by a single-sector template at the maximizing
\(s_{\alpha}\) \cite{harrington2001,devaney2004,stratton1941}. Multi-sector templates with comparable weights require
off-diagonal coupling above threshold \cite{hansen1988}.
\end{corollary}

\paragraph{Worked illustration (mismatched libraries).}
Transmit and receive mode libraries normalized with different \(a_{\alpha}\) from
Eq.~\eqref{eq:ch3.coupling-normalization-propagation} report distorted \(s_{\alpha\beta}\) and false weight ceilings
\cite{harrington2001,hansen1988}. Re-scaling both libraries jointly restores
Eq.~\eqref{eq:ch3.reciprocity-coupling-symmetry} without altering \(\mathcal{P}_{\mathrm{rad}}\)
\cite{stratton1941,poynting1884}.

\paragraph{Problem (weight ceiling audit).}
\textbf{Problem.} A feed delivers \(P_{\mathrm{feed}}=1\,\mathrm{W}\). A design claims \(w_{3,0}^{(\mathbf{N})}=0.9\,\mathrm{W}\) and
\(w_{5,2}^{(\mathbf{M})}=0.8\,\mathrm{W}\) on \(k_{0}D_{s}=1\). Is the claim realizable?

\textbf{Step 1.} Sum weights: \(0.9+0.8>1\) violates passivity unless reactive storage is declared outside
\(\mathcal{F}_{\mathrm{rad}}\) \cite{stratton1941,harrington2001}.

\textbf{Step 2.} Test \(\ell=5\) against Eq.~\eqref{eq:ch4.coupling-ell-decay} \cite{hansen1988,jackson1999}.

\textbf{Step 3.} Conclude \(\mathbf{c}\notin\mathcal{S}_{\mathrm{real}}(\mathcal{C})\) \cite{devaney2004}.

\textbf{Physical meaning.} Modal weights are watts in the flux gauge, not dimensionless pattern samples \cite{hansen1988}.

\paragraph{Multi-mode power budget closure.}
For sectors \(\alpha\in I_{\mathrm{acc}}\), passivity requires
\begin{equation}
\label{eq:ch4.multimode-power-closure}
\sum_{\alpha\in I_{\mathrm{acc}}}
w_{\alpha}
\le
P_{\mathrm{feed}},
\qquad
w_{\alpha}
=
|a_{\alpha}|^{2}
\ \text{in flux gauge},
\end{equation}
extending Eq.~\eqref{eq:ch4.superposition-power-budget} \cite{stratton1941,harrington2001,poynting1884}. Each
\(w_{\alpha}\) is additionally capped by \(\mathcal{W}_{\mathrm{ceiling}}(\alpha)\) and
Eq.~\eqref{eq:ch4.coupling-ell-decay} \cite{hansen1988,devaney2004}.

\begin{theorem}[Feasible weight polytope]
\label{thm:ch4.feasible-weight-polytope}
The set of realizable weight vectors \(\mathbf{w}=(w_{\alpha})_{\alpha\in I_{\mathrm{acc}}}\) is a convex polytope
\(\mathcal{P}_{\mathrm{w}}\subset\mathbb{R}^{|I_{\mathrm{acc}}|}\) cut by hyperplanes
\(w_{\alpha}\le\mathcal{W}_{\mathrm{ceiling}}(\alpha)\), Eq.~\eqref{eq:ch4.multimode-power-closure}, and
\(w_{\alpha}=0\) for sectors nullified by Theorem~\ref{thm:ch4.coupling-selection-rules}
\cite{stratton1941,harrington2001,devaney2004,hansen1988}. A target \(\mathbf{w}_{\mathrm{target}}\notin\mathcal{P}_{\mathrm{w}}\)
implies \(\mathbf{c}_{\mathrm{target}}\notin\mathcal{S}_{\mathrm{real}}(\mathcal{C})\) regardless of phase optimization
\cite{hansen1988}.
\end{theorem}
\begin{proof}
Linearity of \(\boldsymbol{\Gamma}_\omega\) on \(\mathcal{F}_{\mathrm{rad}}\) and convexity of passivity constraints
Eq.~\eqref{eq:ch4.passivity-screen} \cite{stratton1941,harrington2001,poynting1884}.
\end{proof}

\paragraph{Worked illustration (dual-sector budget).}
Claiming \(w_{1,0}^{(\mathbf{N})}=0.6\,\mathrm{W}\) and \(w_{3,2}^{(\mathbf{M})}=0.6\,\mathrm{W}\) on \(P_{\mathrm{feed}}=1\,\mathrm{W}\)
with \(k_{0}D_{s}=1\) violates \(\mathcal{P}_{\mathrm{w}}\): either passivity or
Eq.~\eqref{eq:ch4.coupling-ell-decay} fails \cite{stratton1941,hansen1988,harrington2001}.

\paragraph{Problem (weight vector feasibility).}
\textbf{Problem.} Given \(\mathbf{w}_{\mathrm{target}}\), test realizability without full MoM inversion.

\textbf{Step 1.} Zero weights forbidden by Theorem~\ref{thm:ch4.coupling-selection-rules} \cite{hansen1988}.

\textbf{Step 2.} Test each \(w_{\alpha}\le\mathcal{W}_{\mathrm{ceiling}}(\alpha)\) \cite{harrington2001,devaney2004}.

\textbf{Step 3.} Test \(\sum w_{\alpha}\le P_{\mathrm{feed}}\) \cite{stratton1941,poynting1884}.

\textbf{Physical meaning.} Modal weights live in a geometry-cut polytope, not in unconstrained amplitude space \cite{hansen1988}.

\paragraph{Validity.}
Eqs.~\eqref{eq:ch4.coupling-strength-def}--\eqref{eq:ch4.weight-ceiling-bound} assume narrowband linear theory and
passive feeds \cite{stratton1941}. Active amplification enlarges \(\mathcal{W}_{\mathrm{ceiling}}\) only when
\(P_{\mathrm{feed}}\) includes declared injected power \cite{harrington2001}. Lossy media modify reciprocity symmetry without
removing the rank ceiling \(\mathcal{S}_{\mathrm{rank}}\) \cite{stratton1941,devaney2004}.

Section~\ref{sec:ch45} has coupled finite geometry to angular-momentum eigenfields: synthesis coordinates are
\(\widehat{c}_{\ell m}^{(\sigma)}\), coupling integrals impose selection rules, geometry caps coupling rank, and
\(\mathcal{W}_{\mathrm{ceiling}}\) caps modal weights \cite{hansen1988,harrington2001,devaney2004}. Section~\ref{sec:ch46}
replaces heuristic cuts by truncation exponents and effective electromagnetic dimension on the same operator chain
\cite{stratton1941}.


\section{Truncation and Effective Dimension}
\label{sec:ch46}

Sections~\ref{sec:ch43}--\ref{sec:ch45} identified radiative rank \(r_{\mathrm{rad}}(\varepsilon)\), spectral compression on
\(\boldsymbol{\Gamma}_\omega\), and geometry-capped eigenfield coupling \cite{stratton1941,harrington2001,devaney2004}. Those results answer
how export information is organized; Section~\ref{sec:ch46} answers how much of that organization a finite source can actually use at
declared tolerance \cite{hansen1988,jackson1999}. Informal cutoffs at \(\ell_{\max}\) or patch count are replaced by four auditable objects:
\(\ell_{\mathrm{supp}}\), scaling exponents, coupling-tail laws, and the effective electromagnetic dimension \(d_{\mathrm{eff}}\)
\cite{stratton1941,harrington2001}. Representation may list any number of angular indices; realizability exports only those sectors whose
singular values exceed \(\varepsilon\sigma_{1}\) on \(\mathcal{J}_{\mathrm{real}}\) \cite{devaney2004,hansen1988}. Finite support fixes the
scaling exponents; admissible gauge changes move prefactors but not the inevitability of truncation \cite{jackson1999,stratton1941}.


\subsection{Maximum supported angular order}
\label{sec:ch461}
% TOC tag: [HOME][USE SS3.7.2]

Boundary quantization in Subsection~\ref{sec:ch372} replaces continuous direction spectra by countable mode sets under
\(\Gamma_{b}\) \cite{stratton1941,harrington2001}. On open regions with compact sources, the analogous realizability question is not
how many angular indices a listing can print, but which \(\ell\) sectors carry export gain above tolerance on
\(\mathcal{J}_{\mathrm{real}}\) \cite{hansen1988,devaney2004}. Representation is coordinate freedom: any \(\ell_{\max}\) may be
declared in \(\mathcal{C}\). Realizability is geometry-imposed rank: only sectors coupled through \(\boldsymbol{\Gamma}_\omega\) above
\(\varepsilon\sigma_{1}\) participate in \(\mathcal{S}_{\mathrm{real}}(\mathcal{C})\) \cite{stratton1941,harrington2001}.

\paragraph{Assumptions.}
Fix ball support \(B_{a}\) with radius \(a\), flux-normalized VSH listing \(\mathcal{C}\), and tolerance \(\varepsilon\in(0,1)\)
\cite{stratton1941,hansen1988}. Presuppose Eq.~\eqref{eq:ch4.high-ell-decay}, Definition~\ref{def:ch4.radiative-rank}, and
Eq.~\eqref{eq:ch4.gamma-SVD} \cite{harrington2001,devaney2004}. Narrowband phasors \(\exp(-j\omega t)\) and SI units are in force
\cite{jackson1999}. Synthesis audits use \(\boldsymbol{\Gamma}_\omega=\mathbf{Z}\boldsymbol{\mathcal{K}}_{\mathrm{geo}}\) in the aligned
flux gauge of Section~\ref{sec:ch34} \cite{hansen1988}.

\paragraph{Maximum supported order.}
Define the \emph{maximum supported angular order}
\begin{equation}
\label{eq:ch4.ell-max-supported}
\ell_{\mathrm{supp}}(\varepsilon,k_{0}a)
=
\max\left\{
\ell:\ \exists\,(\ell,m,\sigma)\ \text{with}\ 
\sigma_{(\ell,m,\sigma)}\ge\varepsilon\,\sigma_{1}
\ \text{in Eq.~\eqref{eq:ch4.gamma-SVD}}
\right\},
\end{equation}
where \(\sigma_{(\ell,m,\sigma)}\) is the singular value of \(\boldsymbol{\Gamma}_\omega\) aligned to the unit export vector
\(\widehat{\mathbf{u}}_{\ell m}^{(\sigma)}\) in the VSH chart \cite{harrington2001,hansen1988}. Equation~\eqref{eq:ch4.ell-max-supported}
is the synthesis-side angular cutoff: modes with \(\ell>\ell_{\mathrm{supp}}\) lie below export threshold at \(\varepsilon\) regardless of
listing labels or plotting resolution \cite{stratton1941,devaney2004}. Mathematically, \(\ell_{\mathrm{supp}}\) is the largest block index in
the singular-value ordering whose leading value still exceeds \(\varepsilon\sigma_{1}\). Physically, it is the highest azimuthal order that
can receive appreciable watts from any admissible \(\Jimp\in\mathcal{J}_{\mathrm{real}}(B_{a})\) at the declared frequency
\cite{hansen1988,jackson1999}.

\paragraph{Sector-resolved supported order.}
Electric and magnetic sectors need not share the same ceiling on anisotropic feeds \cite{harrington2001}. Define
\begin{equation}
\label{eq:ch4.ell-sector-supported}
\ell_{\mathrm{supp}}^{(\sigma)}(\varepsilon)
=
\max\big\{\ell:\ \exists m\ \text{with}\ \sigma_{(\ell,m,\sigma)}\ge\varepsilon\sigma_{1}\big\},
\qquad
\sigma\in\{\mathbf{M},\mathbf{N}\},
\end{equation}
and aggregate \(\ell_{\mathrm{supp}}(\varepsilon)=\max_{\sigma}\ell_{\mathrm{supp}}^{(\sigma)}(\varepsilon)\) \cite{hansen1988,stratton1941}.
Sector splitting precedes combined OAM claims: a dominant \(\mathbf{N}\)-sector at \(\ell=5\) may coexist with a suppressed
\(\mathbf{M}\)-sector at the same \(\ell\) on offset feeds \cite{harrington2001,devaney2004}.

\begin{theorem}[Supported order versus electrical size]
\label{thm:ch4.ell-max-scaling}
For compact \(B_{a}\) and aligned VSH gauge,
\begin{equation}
\label{eq:ch4.ell-max-scaling}
\ell_{\mathrm{supp}}(\varepsilon,k_{0}a)
\ \le\
\ell_{\max}(\eta_{\mathrm{tail}},k_{0}a)
\ \sim\
k_{0}a+\mathcal{O}\!\big(\log(1/\varepsilon)\big),
\end{equation}
importing Eq.~\eqref{eq:ch4.ellmax-from-tail} and Eq.~\eqref{eq:ch4.high-ell-decay} \cite{hansen1988,jackson1999,stratton1941}.
Equality holds in the isotropic ball case up to gauge constants \cite{harrington2001,devaney2004}.
\end{theorem}
\begin{proof}
Tail envelope Eq.~\eqref{eq:ch4.high-ell-decay} bounds nonzero \(\widehat{c}_{\ell m}^{(\sigma)}\) by
\(\exp(-\ell/k_{0}a)\) on \(B_{a}\) \cite{hansen1988,jackson1999}. Singular-value ordering of \(\boldsymbol{\Gamma}_\omega\) on VSH
sectors identifies the largest \(\ell\) whose aligned \(\sigma_{(\ell,m,\sigma)}\) remains above \(\varepsilon\sigma_{1}\)
\cite{stratton1941,harrington2001}. Combining the tail bound with Definition~\ref{def:ch4.radiative-rank} yields
Eq.~\eqref{eq:ch4.ell-max-scaling} \cite{devaney2004}.
\end{proof}

Theorem~\ref{thm:ch4.ell-max-scaling} is the open-region counterpart of Subsection~\ref{sec:ch372} quantization: geometry imposes a
countable ceiling on \(\ell\), not merely on mesh indices \cite{stratton1941}. Phase-only metasurfaces on fixed \(a\) cannot raise
\(\ell_{\mathrm{supp}}\) without raising \(k_{0}a\) or enlarging \(\mathcal{L}_{2}\) support \cite{devaney2004,hansen1988}. The theorem
also explains why high-\(\ell\) OAM marketing on handset-scale apertures fails before optimization: the tail inequality is already violated
at the coefficient level \cite{harrington2001,jackson1999}.

\begin{corollary}[Listing cutoff versus supported order]
\label{cor:ch4.listing-vs-supported}
If \(\ell_{\max}>\ell_{\mathrm{supp}}(\varepsilon,k_{0}a)\), indices \(\ell_{\mathrm{supp}}<\ell\le\ell_{\max}\) are export-padding at
tolerance \(\varepsilon\) \cite{hansen1988,harrington2001,devaney2004}. MoM inversion must not treat those rows as independent synthesis
coordinates \cite{harrington2001}.
\end{corollary}

\begin{figure}[t]
\centering
\ChOneFigBlock{%
\begin{tikzpicture}[ch1fig, node distance=7mm and 10mm]
  \node[ch1box] (ka) {$k_{0}a$};
  \node[ch1box, right=12mm of ka] (ell) {$\ell_{\mathrm{supp}}$};
  \node[ch1box, right=12mm of ell] (svd) {Eq.~\eqref{eq:ch4.gamma-SVD}};
  \node[ch1box, right=12mm of svd] (sr) {$\mathcal{S}_{\mathrm{real}}$};
  \node[ch1box, below=11mm of ell] (tail) {Eq.~\eqref{eq:ch4.high-ell-decay}};
  \draw[ch1flow] (ka) -- (ell) -- (svd) -- (sr);
  \draw[ch1flow] (tail) -- (ell);
\end{tikzpicture}%
}
\caption{Subsection~\ref{sec:ch461}: electrical size sets \(\ell_{\mathrm{supp}}\) through tail bounds before synthesis listing; only
sectors above \(\varepsilon\sigma_{1}\) enter \(\mathcal{S}_{\mathrm{real}}\).}
\label{fig:ch4.ell-max-supported}
\end{figure}

Figure~\ref{fig:ch4.ell-max-supported} compresses the audit chain: \(k_{0}a\) controls tail decay, tail decay controls
\(\ell_{\mathrm{supp}}\), and \(\ell_{\mathrm{supp}}\) controls which rows of \(\mathcal{C}\) carry realizability content
\cite{hansen1988,stratton1941}. A listing that stops at \(\ell_{\max}\) without reporting \(\ell_{\mathrm{supp}}\) confuses coordinate
cardinality with export rank \cite{harrington2001,devaney2004}.

\paragraph{Tail-to-threshold translation.}
Eq.~\eqref{eq:ch4.high-ell-decay} implies that sector power at \(\ell\) is suppressed by \(\exp(-2\ell/k_{0}a)\) relative to low-\(\ell\)
sectors after flux normalization \cite{hansen1988,jackson1999}. Requiring \(\sigma_{(\ell,m,\sigma)}\ge\varepsilon\sigma_{1}\) therefore
forces \(\ell\lesssim k_{0}a+\log(1/\varepsilon)\) \cite{stratton1941}. At \(\varepsilon=10^{-2}\) and \(k_{0}a=2\), the logarithmic
correction adds only \(\mathcal{O}(1)\) to the linear term: \(\ell_{\mathrm{supp}}\) is dominated by electrical size, not by solver tolerance
\cite{harrington2001,devaney2004}.
\begin{equation}
\label{eq:ch4.ell-supp-monotone}
\frac{\partial\ell_{\mathrm{supp}}(\varepsilon,k_{0}a)}{\partial(k_{0}a)}\ge0,\qquad
\frac{\partial\ell_{\mathrm{supp}}(\varepsilon,k_{0}a)}{\partial\varepsilon}\le0.
\end{equation}
Equation~\eqref{eq:ch4.ell-supp-monotone} is the monotone chain used in laboratory audits: enlarging the ball or raising frequency opens
higher \(\ell\) sectors; tightening \(\varepsilon\) lowers \(\ell_{\mathrm{supp}}\) \cite{hansen1988,jackson1999}.

\paragraph{Worked illustration (handset aperture).}
At \(k_{0}a=2\), \(\ell_{\mathrm{supp}}(10^{-2})\approx4\)--\(6\) by Eq.~\eqref{eq:ch4.ell-max-scaling} \cite{hansen1988,jackson1999,harrington2001}.
Marketing claims of dominant \(\ell=12\) OAM on the same aperture fail Theorem~\ref{thm:ch4.ell-max-scaling} because tail mass at
\(\ell=12\) is \(\exp(-6)\approx0.002\) relative to \(\ell\le2\) sectors after normalization \cite{devaney2004}. The failure is
coefficient-level, not a matter of feed tuning alone \cite{stratton1941}.

\paragraph{Worked illustration (sector split on offset feed).}
A patch displaced from the ball center may satisfy \(\ell_{\mathrm{supp}}^{(\mathbf{N})}(10^{-2})=5\) while
\(\ell_{\mathrm{supp}}^{(\mathbf{M})}(10^{-2})=4\) at the same \(k_{0}a\) \cite{harrington2001,jackson1999,hansen1988}. Combined OAM
reports must decompose \(\mathbf{M}\) and \(\mathbf{N}\) sectors before asserting a single dominant \(\ell\) \cite{stratton1941,devaney2004}.

\paragraph{Worked numerics (\(k_{0}a\) sweep).}
At \(\varepsilon=10^{-2}\), representative centered-feed values are \(\ell_{\mathrm{supp}}\approx3,5,7,9\) for
\(k_{0}a=1,1.5,2,2.5\) \cite{hansen1988,harrington2001}. Any design target with \(\ell>\ell_{\mathrm{supp}}(k_{0}a)\) is outside
\(\mathcal{S}_{\mathrm{real}}(\mathcal{C})\) on \(B_{a}\) before MoM or inverse design begins \cite{stratton1941,devaney2004}.

\paragraph{Problem (supported order audit).}
\textbf{Problem.} A VSH listing retains \(\ell\le15\) but \(k_{0}a=1.8\). How many \(\ell\) sectors matter at \(\varepsilon=10^{-2}\)?

\textbf{Step 1.} Compute \(\ell_{\mathrm{supp}}(10^{-2},1.8)\) from Eq.~\eqref{eq:ch4.ell-max-supported} via SVD of \(\mathbf{T}\) on the
aligned export basis \cite{harrington2001}.

\textbf{Step 2.} Compare to Eq.~\eqref{eq:ch4.ell-max-scaling} for conservative analytic bound \(\ell_{\mathrm{supp}}\lesssim1.8+\log(100)\approx6\)
\cite{hansen1988,jackson1999}.

\textbf{Step 3.} Truncate design targets and MoM templates to \(\ell\le\ell_{\mathrm{supp}}\); rows \(\ell_{\mathrm{supp}}<\ell\le15\) are
padding per Corollary~\ref{cor:ch4.listing-vs-supported} \cite{devaney2004,stratton1941}.

\textbf{Physical meaning.} \(\ell_{\mathrm{supp}}\) is a geometry-imposed ceiling on exportable azimuthal order, not a plotting convenience
\cite{hansen1988}.

\paragraph{Problem (supported order versus \(\ell_{\max}\)).}
\textbf{Problem.} \(\ell_{\max}=12\) is used in \(\mathcal{C}\) but \(\ell_{\mathrm{supp}}(10^{-2})=7\). How many rows matter?

\textbf{Step 1.} Count sectors with \(\sigma_{(\ell,m,\sigma)}\ge10^{-2}\sigma_{1}\) in Eq.~\eqref{eq:ch4.gamma-SVD} \cite{harrington2001}.

\textbf{Step 2.} Apply Corollary~\ref{cor:ch4.listing-vs-supported}: indices \(8\le\ell\le12\) are export-padding \cite{hansen1988,stratton1941}.

\textbf{Step 3.} Truncate synthesis templates to \(\ell\le7\) and recompute \(\delta_{\mathrm{syn}}\) on the truncated target
\cite{devaney2004}.

\textbf{Physical meaning.} Supported order is an export threshold; listing width is coordinate freedom \cite{hansen1988,harrington2001}.

\paragraph{Problem (OAM claim on fixed disk).}
\textbf{Problem.} A \(0.2\,\mathrm{m}\) disk at \(30\,\mathrm{GHz}\) (\(k_{0}a\approx12.6\)) claims \(\ell=20\) dominance. Audit.

\textbf{Step 1.} Compute \(\ell_{\mathrm{supp}}(10^{-2},12.6)\lesssim12.6+\log(100)\approx17\) from Eq.~\eqref{eq:ch4.ell-max-scaling}
\cite{jackson1999,hansen1988}.

\textbf{Step 2.} Test tail ratio \(\exp(-2\cdot20/12.6)\approx0.04\) for \(\ell=20\) against \(\varepsilon=10^{-2}\) \cite{hansen1988}.

\textbf{Step 3.} If SVD shows \(\sigma_{(20,m,\sigma)}<10^{-2}\sigma_{1}\), reject dominance claim \cite{harrington2001,devaney2004}.

\textbf{Physical meaning.} Even electrically large apertures saturate \(\ell\) through tail decay, not through listing preference
\cite{stratton1941}.

\paragraph{Gauge invariance of \(\ell_{\mathrm{supp}}\).}
Admissible changes of \(\mathcal{C}\) that preserve flux normalization and \(\boldsymbol{\Gamma}_\omega\) leave \(\ell_{\mathrm{supp}}\)
invariant \cite{hansen1988,stratton1941}. Changing \(\ell_{\max}\) without changing \(\boldsymbol{\Gamma}_\omega\) adds null rows to
\(\mathbf{T}\) but does not enlarge \(\mathcal{S}_{\mathrm{real}}(\mathcal{C})\) \cite{harrington2001,devaney2004}. Laboratories must
report \(\ell_{\mathrm{supp}}\) alongside \(\ell_{\max}\) to prevent coordinate inflation in synthesis reports \cite{hansen1988}.

\paragraph{Relation to \(d_{\mathrm{eff}}\).}
Each supported \(\ell\) block contributes up to \((2\ell+1)\) sector labels, but \(d_{\mathrm{eff}}(\varepsilon)\) counts only singular
directions above \(\varepsilon\sigma_{1}\) \cite{harrington2001,devaney2004}. Thus \(\ell_{\mathrm{supp}}\) is necessary but not sufficient
for rank saturation: within \(\ell\le\ell_{\mathrm{supp}}\), many \((m,\sigma)\) pairs may still lie below threshold
\cite{hansen1988,stratton1941}.
\begin{equation}
\label{eq:ch4.ell-deff-relation}
d_{\mathrm{eff}}(\varepsilon)\le\sum_{\ell=0}^{\ell_{\mathrm{supp}}}(2\ell+1)\times\#\{\sigma:\sigma\text{ active}\},
\end{equation}
with equality only when every sector in the sum exceeds \(\varepsilon\sigma_{1}\) \cite{harrington2001}.

\paragraph{MoM row audit protocol.}
Let \(\mathbf{T}\) be the discretized export map in the aligned gauge \cite{harrington2001}. Row \( (\ell,m,\sigma)\) is \emph{synthesis-active}
if \(\sigma_{(\ell,m,\sigma)}\ge\varepsilon\sigma_{1}\); otherwise it is \emph{padding} per Corollary~\ref{cor:ch4.listing-vs-supported}
\cite{hansen1988,devaney2004}. The audit protocol is: (i) form \(\mathbf{T}\) on \(\mathcal{J}_{\mathrm{real}}(B_{a})\); (ii) extract
\(\{\sigma_{(\ell,m,\sigma)}\}\); (iii) compute \(\ell_{\mathrm{supp}}\) via Eq.~\eqref{eq:ch4.ell-max-supported}; (iv) truncate templates to
active rows only \cite{stratton1941,harrington2001}. Skipping step (iii) is the dominant source of false OAM realizability claims in numerical
papers \cite{hansen1988,jackson1999}.

\begin{theorem}[Padding rows are synthesis-invisible]
\label{thm:ch4.padding-invisible}
Rows with \(\sigma_{(\ell,m,\sigma)}<\varepsilon\sigma_{1}\) carry no component of \(\mathcal{S}_{\mathrm{real}}(\mathcal{C})\) above
threshold: for any \(\mathbf{c}\in\mathcal{S}_{\mathrm{real}}(\mathcal{C})\), the coefficient on padding row \((\ell,m,\sigma)\) is
\(\mathcal{O}(\varepsilon)\|\mathbf{c}\|\) in the flux gauge \cite{harrington2001,devaney2004,stratton1941}.
\end{theorem}
\begin{proof}
Eckart--Young truncation at rank \(r_{\mathrm{rad}}(\varepsilon)\) discards singular directions below \(\varepsilon\sigma_{1}\)
\cite{harrington2001}. Padding rows span those directions \cite{hansen1988,stratton1941}.
\end{proof}

\paragraph{Worked illustration (MoM padding).}
\(\ell_{\max}=15\), \(k_{0}a=1.5\), \(\ell_{\mathrm{supp}}(10^{-2})=5\): rows with \(\ell>5\) contribute less than \(1\%\) of export power
yet inflate \(\dim\mathbf{T}\) by a factor \(\approx3\) \cite{hansen1988,harrington2001,devaney2004}. Truncating \(\mathbf{T}\) to
\(\ell\le5\) before inversion stabilizes \(\delta_{\mathrm{syn}}\) without changing realizable \(\mathbf{c}\) \cite{stratton1941}.

\paragraph{Problem (padding strip).}
\textbf{Problem.} After SVD, \(40\%\) of rows lie below \(\varepsilon\sigma_{1}\). Strip them?

\textbf{Step 1.} Identify padding via Theorem~\ref{thm:ch4.padding-invisible} \cite{harrington2001}.

\textbf{Step 2.} Rebuild \(\mathbf{T}\) on active rows only \cite{devaney2004,hansen1988}.

\textbf{Step 3.} Re-certify \(\delta_{\mathrm{syn}}=0\) on truncated target \cite{stratton1941}.

\textbf{Physical meaning.} Padding rows are coordinate artifacts, not synthesis degrees of freedom \cite{hansen1988}.

\paragraph{Philosophical closure (representation versus export).}
The VSH chart is complete on \(S^{2}\): any square-integrable far pattern admits an expansion in \(\widehat{c}_{\ell m}^{(\sigma)}\)
\cite{hansen1988,stratton1941}. Completeness is a representation theorem; \(\ell_{\mathrm{supp}}\) is a realizability theorem on
\(\mathcal{J}_{\mathrm{real}}(B_{a})\) \cite{harrington2001,devaney2004}. Confusing the two is the root error in OAM device claims that
display high-\(\ell\) coefficients produced by numerical noise or gauge padding \cite{jackson1999,hansen1988}. Export audits must always ask:
does \(\boldsymbol{\Gamma}_\omega[\mathcal{J}_{\mathrm{real}}]\) carry weight above \(\varepsilon\sigma_{1}\) in that sector?
\cite{stratton1941,harrington2001}.

\paragraph{Power-participation ratio by \(\ell\) block.}
Define block participation
\begin{equation}
\label{eq:ch4.ell-block-participation}
\mathrm{PPR}(\ell)=\frac{\sum_{m,\sigma}\sigma_{(\ell,m,\sigma)}^{2}}{\sum_{\ell',m',\sigma'}\sigma_{(\ell',m',\sigma')}^{2}},
\end{equation}
with \(\mathrm{PPR}(\ell)<\varepsilon^{2}\) for all \(\ell>\ell_{\mathrm{supp}}\) in the aligned gauge \cite{harrington2001,hansen1988}.
Equation~\eqref{eq:ch4.ell-block-participation} converts singular-value audits into per-order power language laboratories already use
\cite{devaney2004,stratton1941}.

\paragraph{Worked illustration (PPR audit).}
At \(k_{0}a=2\), \(\mathrm{PPR}(6)\approx0.008\) while \(\mathrm{PPR}(3)\approx0.22\): order \(\ell=6\) is below production tolerance
even if included in \(\ell_{\max}=12\) listing \cite{hansen1988,harrington2001,jackson1999}.

\paragraph{Problem (PPR versus marketing plot).}
\textbf{Problem.} Polar plot shows visible \(\ell=8\) sidelobe structure. Is \(\ell=8\) supported?

\textbf{Step 1.} Compute \(\mathrm{PPR}(8)\) from \(\mathbf{T}\) SVD, not from plot amplitude \cite{harrington2001,hansen1988}.

\textbf{Step 2.} Compare \(\mathrm{PPR}(8)\) to \(\varepsilon^{2}\) \cite{devaney2004}.

\textbf{Step 3.} Attribute sidelobes to padding or noise if \(\mathrm{PPR}(8)<\varepsilon^{2}\) \cite{stratton1941}.

\textbf{Physical meaning.} Visible plot features below export threshold are not synthesis resources \cite{hansen1988}.

\paragraph{Certificate table (supported-order audit).}
\begin{equation}
\label{eq:ch4.ell-supp-certificate}
\underbrace{\ell_{\mathrm{supp}}(\varepsilon,k_{0}a)}_{\text{angular ceiling}}
\ \le\
\underbrace{\ell_{\max}(\eta_{\mathrm{tail}},k_{0}a)}_{\text{tail bound}}
\ \le\
\underbrace{k_{0}a+\mathcal{O}(\log(1/\varepsilon))}_{\text{Eq.~\eqref{eq:ch4.ell-max-scaling}}}.
\end{equation}
Equation~\eqref{eq:ch4.ell-supp-certificate} is the audit chain for synthesis datasheets \cite{hansen1988,jackson1999,harrington2001,stratton1941}.

\begin{figure}[t]
\centering
\ChOneFigBlock{%
\begin{tikzpicture}[ch1fig, node distance=7mm and 8mm]
  \node[ch1box] (list) {$\ell_{\max}$};
  \node[ch1box, right=10mm of list] (supp) {$\ell_{\mathrm{supp}}$};
  \node[ch1box, right=10mm of supp] (tail) {Eq.~\eqref{eq:ch4.high-ell-decay}};
  \node[ch1box, right=10mm of tail] (ka) {$k_{0}a$};
  \draw[ch1flow] (list) -- (supp) -- (tail) -- (ka);
\end{tikzpicture}%
}
\caption{Subsection~\ref{sec:ch461}: realizability content increases rightward along the certificate chain.}
\label{fig:ch4.ell-supp-certificate-chain}
\end{figure}

\paragraph{Worked illustration (certificate failure).}
\(\ell_{\max}=15\), \(k_{0}a=1.4\), no \(\ell_{\mathrm{supp}}\) reported: audit gives \(\ell_{\mathrm{supp}}\lesssim5\) at \(\varepsilon=10^{-2}\)
\cite{hansen1988,harrington2001,devaney2004}.

\paragraph{Problem (certificate completion).}
\textbf{Problem.} Complete Eq.~\eqref{eq:ch4.ell-supp-certificate} for \(k_{0}a=2.2\), \(\varepsilon=10^{-2}\).
\textbf{Step 1.} Read \(\ell_{\mathrm{supp}}\) from SVD \cite{harrington2001}.
\textbf{Step 2.} Bound tail term via Eq.~\eqref{eq:ch4.ellmax-from-tail} \cite{hansen1988,jackson1999}.
\textbf{Step 3.} Verify chain and attach to report \cite{stratton1941,devaney2004}.
\textbf{Physical meaning.} Supported order is the certified ceiling, not the listing index \cite{hansen1988}.

\paragraph{Validity.}
Eqs.~\eqref{eq:ch4.ell-max-supported}--\eqref{eq:ch4.ell-max-scaling} assume open-region ball supports, aligned VSH gauge, and linear
narrowband Maxwell maps \cite{stratton1941,hansen1988}. Cavities require Subsection~\ref{sec:ch372} indices \((\ell,m,n_{r})\) instead of
pure spherical orders \cite{harrington2001}. Loss lowers \(Q\) and may redistribute sector weights without promoting \(\ell>\ell_{\mathrm{supp}}\)
on passive supports \cite{stratton1941,devaney2004}. Anisotropic media modify \(\ell_{\mathrm{supp}}^{(\sigma)}\) splits but not the
inevitability of a ceiling at fixed \((k_{0}a,\varepsilon)\) \cite{hansen1988,harrington2001}.


\subsection{Scaling laws for spectral truncation}
\label{sec:ch462}
% TOC tag: [HOME]

Truncation is not a numerical afterthought applied once a MoM matrix is assembled. On compact supports it is a scaling law: the count of
export directions above \(\varepsilon\sigma_{1}\) grows with electrical size and frequency according to exponents fixed by the support class
\cite{stratton1941,jackson1999,hansen1988}. Subsection~\ref{sec:ch462} states those exponents as realizability inequalities on
\(r_{\mathrm{rad}}(\varepsilon)\), \(\ell_{\mathrm{supp}}\), and \(\operatorname{rank}\mathbf{K}\) \cite{harrington2001,devaney2004}. The
philosophical point is inevitability: finite geometry cannot sustain an unbounded synthesis chart at fixed tolerance.

\paragraph{Assumptions.}
Presuppose Eqs.~\eqref{eq:ch4.weyl-mode-count}, \eqref{eq:ch4.geometry-coupling-scaling}, \eqref{eq:ch4.aperture-coupling-scaling}, and
Theorem~\ref{thm:ch4.support-volume-scaling} \cite{stratton1941,hansen1988,harrington2001}. Asymptotic regime \(k_{0}a\ge1\) unless the
Rayleigh branch is explicitly invoked \cite{jackson1999}. Tolerance \(\varepsilon\) is held fixed across comparisons \cite{devaney2004}.

\paragraph{Volume and area exponents.}
On enclosing ball \(B_{a}\),
\begin{equation}
\label{eq:ch4.truncation-volume-scaling}
r_{\mathrm{rad}}(\varepsilon)
\ \sim\
C_{\mathrm{vol}}(\varepsilon)\,(k_{0}a)^{3},
\qquad
\ell_{\mathrm{supp}}(\varepsilon,k_{0}a)
\ \sim\
C_{\ell}(\varepsilon)\,k_{0}a,
\end{equation}
importing Eq.~\eqref{eq:ch4.weyl-mode-count} and Theorem~\ref{thm:ch4.ell-max-scaling} \cite{stratton1941,jackson1999,hansen1988}. The cubic
exponent on \(r_{\mathrm{rad}}\) is the export-side image of volume mode density: each additional wavelength inside \(B_{a}\) opens new
orthogonal radiation channels \cite{harrington2001}. The linear exponent on \(\ell_{\mathrm{supp}}\) is the angular tail counterpart: high
orders are permitted only in proportion to \(k_{0}a\) \cite{hansen1988,devaney2004}. On aperture \(\Sigma_{a}\) of area \(A_{a}\),
\begin{equation}
\label{eq:ch4.truncation-area-scaling}
\operatorname{rank}\mathbf{K}\big|_{\Sigma_{a}}
\ \sim\
C_{A}(\varepsilon)\,k_{0}^{2}A_{a},
\end{equation}
from Eq.~\eqref{eq:ch4.aperture-coupling-scaling} \cite{collin2001,harrington2001,devaney2004}. Volume scaling dominates when synthesis is
declared on the enclosing ball; area scaling governs mouth-only and patch-array models \cite{stratton1941}.

\begin{proposition}[Scaling law consistency]
\label{prop:ch4.truncation-scaling-consistency}
If equivalence Theorem~\ref{thm:ch4.volume-surface-equivalence} holds and \(\mathcal{D}_{\mathrm{cpl}}=1\),
Eqs.~\eqref{eq:ch4.truncation-volume-scaling} and \eqref{eq:ch4.truncation-area-scaling} bound the same \(\mathbf{c}\) image up to
constants depending on \(\varepsilon\) and mesh alignment \cite{stratton1941,harrington2001,collin2001}. Neither exponent is raised by
phase-only retuning on fixed geometry \cite{devaney2004,hansen1988}.
\end{proposition}
\begin{proof}
\(\boldsymbol{\Gamma}_\omega\) depends on \(\mathcal{J}_{\mathrm{real}}\) support class; phase rotations preserve singular values on fixed
support \cite{harrington2001,stratton1941}. Volume--surface equivalence identifies images up to equivalence constants \cite{collin2001}.
\end{proof}

\paragraph{Joint scaling audit.}
Laboratories should report a single conservative bound
\begin{equation}
\label{eq:ch4.joint-scaling-law}
d_{\mathrm{eff}}(\varepsilon)\le\min\big\{C_{V}(\varepsilon)(k_{0}a)^{3},\,C_{A}(\varepsilon)k_{0}^{2}A_{a}\big\},
\end{equation}
with the active term determined by declared synthesis support \cite{stratton1941,hansen1988,harrington2001}. Equation~\eqref{eq:ch4.joint-scaling-law}
prevents aperture-only optimism when currents occupy a larger enclosing volume \cite{devaney2004,collin2001}.

\begin{theorem}[Scaling exponent gauge invariance]
\label{thm:ch4.scaling-gauge-invariant}
Exponents \(3\) and \(2\) in Eq.~\eqref{eq:ch4.joint-scaling-law} are invariant under admissible changes of \(\mathcal{C}\); only
prefactors \(C_{V},C_{A}\) depend on flux normalization \cite{hansen1988,stratton1941,harrington2001}.
\end{theorem}
\begin{proof}
\(\boldsymbol{\Gamma}_\omega\) is unitarily equivalent under flux gauge changes; singular-value thresholds scale homogeneously, leaving rank
growth exponents unchanged \cite{hansen1988,harrington2001}.
\end{proof}

\begin{figure}[t]
\centering
\ChOneFigBlock{%
\begin{tikzpicture}[ch1fig, node distance=7mm and 9mm]
  \node[ch1box] (vol) {$(k_{0}a)^{3}$};
  \node[ch1box, right=11mm of vol] (area) {$k_{0}^{2}A_{a}$};
  \node[ch1box, right=11mm of area] (min) {$\min$};
  \node[ch1box, right=11mm of min] (deff) {$d_{\mathrm{eff}}$};
  \draw[ch1flow] (vol) -- (min);
  \draw[ch1flow] (area) -- (min);
  \draw[ch1flow] (min) -- (deff);
\end{tikzpicture}%
}
\caption{Subsection~\ref{sec:ch462}: truncation rank is bounded by the smaller of volume-growth and area-growth laws at declared support.}
\label{fig:ch4.truncation-scaling-audit}
\end{figure}

Figure~\ref{fig:ch4.truncation-scaling-audit} is the design-room audit: doubling frequency at fixed outline moves along the appropriate
exponent; changing outline without updating \(k_{0}D_{s}\) leaves rank claims inconsistent \cite{jackson1999,harrington2001}.

\paragraph{Prefactor sensitivity.}
Constants \(C_{V}(\varepsilon)\) and \(C_{A}(\varepsilon)\) depend on \(\varepsilon\) and gauge but not on \(\ell_{\max}\)
\cite{hansen1988,stratton1941}. On a circular aperture,
\begin{equation}
\label{eq:ch4.scaling-ratio}
\frac{C_{V}(\varepsilon)(k_{0}a)^{3}}{C_{A}(\varepsilon)k_{0}^{2}A_{a}}
\sim
\frac{C_{V}}{C_{A}}\cdot\frac{k_{0}a}{\pi}.
\end{equation}
At \(k_{0}a=3\) with \(C_{V}/C_{A}\approx1\), volume and area bounds are comparable: both must be audited \cite{harrington2001,devaney2004}.

\paragraph{Worked illustration (doubling frequency).}
At fixed \(a\), doubling \(k_{0}\) doubles \(k_{0}a\): \(r_{\mathrm{rad}}\) grows by \(2^{3}=8\) in the large-body regime while
\(\ell_{\mathrm{supp}}\) grows by \(2\) \cite{jackson1999,harrington2001,hansen1988}. A handset outline unchanged from sub-6\,GHz to mmWave
therefore crosses distinct realizability classes \cite{stratton1941,collin2001}.

\paragraph{Worked illustration (horn versus ball).}
A horn feed modeled on mouth area \(A_{a}\) may satisfy \(k_{0}^{2}A_{a}\approx40\) while the enclosing \(B_{a}\) yields
\((k_{0}a)^{3}\approx30\) \cite{collin2001,harrington2001}. Eq.~\eqref{eq:ch4.joint-scaling-law} selects the area bound for aperture-only
synthesis and the volume bound when \(\mathcal{J}_{\mathrm{real}}\) fills the enclosure \cite{stratton1941,devaney2004}.

\paragraph{Worked numerics (5G panel).}
Panel with \(k_{0}W=3\) gives \(k_{0}^{2}A_{a}\sim9\pi\approx28\) and \(d_{\mathrm{eff}}(10^{-2})\lesssim30\) before element-count claims
\cite{hansen1988,collin2001,harrington2001}.

\paragraph{Problem (volume versus area rank).}
\textbf{Problem.} Patch array on \(\Sigma_{a}\) reports \(k_{0}^{2}A_{a}\approx50\); enclosing ball gives \((k_{0}a)^{3}\approx30\). Which
bound governs?

\textbf{Step 1.} Declare whether synthesis is on \(\Sigma_{a}\) or \(B_{a}\) \cite{collin2001,stratton1941}.

\textbf{Step 2.} Apply matching term in Eq.~\eqref{eq:ch4.joint-scaling-law} \cite{harrington2001,devaney2004}.

\textbf{Step 3.} Use \(\min(\cdot)\) for conservative \(d_{\mathrm{eff}}\) report \cite{hansen1988}.

\textbf{Physical meaning.} Scaling exponents attach to support class, not CAD convenience \cite{stratton1941}.

\paragraph{Problem (frequency extrapolation).}
\textbf{Problem.} Rank measured at \(3\,\mathrm{GHz}\) is extrapolated to \(30\,\mathrm{GHz}\) on fixed hardware. Valid?

\textbf{Step 1.} Compute \(k_{0}D_{s}\) at each frequency \cite{jackson1999}.

\textbf{Step 2.} Apply Eq.~\eqref{eq:ch4.truncation-volume-scaling} with factor \((10)^{3}=1000\) only in large-body regime \cite{hansen1988,harrington2001}.

\textbf{Step 3.} Recompute \(d_{\mathrm{eff}}(\varepsilon)\) by SVD at each band \cite{devaney2004}.

\textbf{Physical meaning.} Rank scales with electrical size, not with outline drawings alone \cite{stratton1941}.

\paragraph{Problem (phase-only rank claim).}
\textbf{Problem.} Metasurface retunes phases; designer claims doubled \(d_{\mathrm{eff}}\). Audit.

\textbf{Step 1.} Verify \(\operatorname{diam}(\Omega_{s})\) and \(k_{0}\) unchanged \cite{stratton1941,harrington2001}.

\textbf{Step 2.} Invoke Proposition~\ref{prop:ch4.truncation-scaling-consistency} \cite{hansen1988,devaney2004}.

\textbf{Step 3.} Reject rank doubling; phase rotates within fixed export subspace \cite{harrington2001}.

\textbf{Physical meaning.} Exponents are geometry-imposed; phase is coordinate freedom on fixed rank \cite{stratton1941}.

\paragraph{Rayleigh crossover reminder.}
When \(k_{0}a\ll1\), Eq.~\eqref{eq:ch4.truncation-volume-scaling} is superseded by dipole-class rank \(\le3\) \cite{jackson1999,hansen1988}.
Scaling laws must name the active regime before extrapolation \cite{stratton1941,harrington2001}.

\paragraph{Design frequency ladder.}
Hold \(\operatorname{diam}(\Omega_{s})\) fixed and sample \(k_{0}\) at bands \(f_{1}<f_{2}<f_{3}\). Then
\begin{equation}
\label{eq:ch4.freq-ladder}
\frac{d_{\mathrm{eff}}(f_{3})}{d_{\mathrm{eff}}(f_{1})}
\approx
\left(\frac{f_{3}}{f_{1}}\right)^{3}
\quad\text{on }B_{a},\qquad
\frac{\ell_{\mathrm{supp}}(f_{3})}{\ell_{\mathrm{supp}}(f_{1})}
\approx
\frac{f_{3}}{f_{1}},
\end{equation}
in the large-body regime \cite{jackson1999,hansen1988,harrington2001}. Handset roadmaps that quote a single ``mode count'' across decades of
frequency without recomputing Eq.~\eqref{eq:ch4.freq-ladder} are not auditable \cite{stratton1941,devaney2004,collin2001}.

\begin{corollary}[Outline invariance does not imply rank invariance]
\label{cor:ch4.outline-rank-noninvariant}
Mechanical outline invariance with frequency change leaves \(\operatorname{diam}(\Omega_{s})\) fixed but increases \(k_{0}D_{s}\), hence
increases \(d_{\mathrm{eff}}(\varepsilon)\) per Eq.~\eqref{eq:ch4.truncation-volume-scaling} \cite{jackson1999,harrington2001,stratton1941}.
\end{corollary}

\paragraph{Worked illustration (band ladder).}
\(150\,\mathrm{mm}\) device: \(d_{\mathrm{eff}}(10^{-2})\le3\) near \(700\,\mathrm{MHz}\); same outline near \(28\,\mathrm{GHz}\) gives
\(d_{\mathrm{eff}}\sim\mathcal{O}(10^{2})\) before element-count marketing \cite{harrington2001,collin2001,hansen1988,jackson1999}.

\paragraph{Problem (single mode-count spec).}
\textbf{Problem.} Vendor quotes ``32 modes'' for a phone outline without frequency. Audit.

\textbf{Step 1.} Request \((f,k_{0}D_{s},\varepsilon)\) triple \cite{jackson1999,hansen1988}.

\textbf{Step 2.} Apply Eq.~\eqref{eq:ch4.joint-scaling-law} at that triple \cite{harrington2001}.

\textbf{Step 3.} Reject spec without frequency; rank is not outline-invariant \cite{stratton1941,devaney2004}.

\textbf{Physical meaning.} Mode count is a function of electrical size, not mechanical drawing alone \cite{hansen1988}.

\paragraph{Derivation sketch (volume exponent).}
Eq.~\eqref{eq:ch4.truncation-volume-scaling} inherits the Weyl counting exponent from Eq.~\eqref{eq:ch4.weyl-mode-count} on \(B_{a}\)
\cite{stratton1941,jackson1999}. Each wavelength cubed inside the ball adds \(\mathcal{O}(1)\) orthogonal radiation channels at fixed
\(\varepsilon\) because \(\boldsymbol{\Gamma}_\omega\) is a compact operator whose essential rank tracks volume mode density
\cite{hansen1988,harrington2001,devaney2004}. The linear \(\ell_{\mathrm{supp}}\) exponent in the same equation pair is the angular shadow of
that volume law: high orders cannot participate until \(k_{0}a\) is large enough to defeat tail decay \cite{jackson1999,hansen1988}.

\paragraph{Area exponent for aperture synthesis.}
When \(\mathcal{J}_{\mathrm{real}}\) is surface-supported on \(\Sigma_{a}\), rank growth is quadratic in wavenumber times area because
tangential currents on a two-dimensional manifold carry a two-dimensional spatial mode budget \cite{collin2001,stratton1941,harrington2001}.
Patch arrays sample that budget; they do not create new budget beyond \(k_{0}^{2}A_{a}\) \cite{hansen1988,devaney2004}.

\begin{theorem}[Exponent selection rule]
\label{thm:ch4.exponent-selection}
Let \(d_{\mathrm{eff}}^{\mathrm{(vol)}}\sim(k_{0}a)^{3}\) and \(d_{\mathrm{eff}}^{\mathrm{(ap)}}\sim k_{0}^{2}A_{a}\). The operative
truncation exponent is determined by \(\mathcal{J}_{\mathrm{real}}\) support dimension: volume-fed synthesis uses exponent \(3\);
aperture-fed synthesis uses exponent \(2\) \cite{stratton1941,collin2001,hansen1988,harrington2001}. Mixed models require explicit
domain declaration \cite{devaney2004}.
\end{theorem}
\begin{proof}
\(\boldsymbol{\Gamma}_\omega\) is defined on the declared \(\mathcal{J}_{\mathrm{real}}\) class; rank growth follows support Hausdorff
dimension through Weyl and sampling laws cited above \cite{stratton1941,jackson1999,hansen1988}.
\end{proof}

\paragraph{Worked illustration (mixed model).}
Horn interior currents fill \(B_{a}\) but mouth synthesis lists only \(\Sigma_{a}\): volume exponent governs physics, area exponent governs
reported port modes \cite{collin2001,harrington2001,stratton1941}. Audits must separate the two domains \cite{devaney2004,hansen1988}.

\paragraph{Problem (exponent selection).}
\textbf{Problem.} Full-wave model uses volume mesh; synthesis spec quotes aperture modes only. Which scaling law certifies the spec?

\textbf{Step 1.} Identify \(\mathcal{J}_{\mathrm{real}}\) support in the synthesis claim \cite{stratton1941,collin2001}.

\textbf{Step 2.} Apply Theorem~\ref{thm:ch4.exponent-selection} \cite{hansen1988,harrington2001}.

\textbf{Step 3.} Reject specs that cite area exponent while optimizing volume currents \cite{devaney2004}.

\textbf{Physical meaning.} Truncation exponents follow support dimension, not mesh convenience \cite{stratton1941}.

\paragraph{Weyl-to-realizability bridge (pass 3).}
Eq.~\eqref{eq:ch4.weyl-mode-count} counts modes in a phase-space ball; Eq.~\eqref{eq:ch4.truncation-volume-scaling} counts export directions
above \(\varepsilon\sigma_{1}\) \cite{stratton1941,jackson1999,hansen1988}. The bridge inequality
\begin{equation}
\label{eq:ch4.weyl-to-realizability}
r_{\mathrm{rad}}(\varepsilon)\le C_{\mathrm{Weyl}}(\varepsilon)\,N_{\mathrm{mode}}^{\mathrm{(vol)}}
\le
C_{\mathrm{Weyl}}(\varepsilon)\,C_{V}(\varepsilon)(k_{0}a)^{3}
\end{equation}
states that realizability rank cannot exceed Weyl mode budget at the same tolerance \cite{harrington2001,devaney2004}. Philosophically,
Weyl counting is a necessary condition for export; sufficient strength requires \(\boldsymbol{\Gamma}_\omega\) singular values
\cite{hansen1988,stratton1941}.

\paragraph{Worked illustration (Weyl bridge).}
\(k_{0}a=2.5\) gives \(N_{\mathrm{mode}}^{\mathrm{(vol)}}\sim\mathcal{O}(15)\) while \(r_{\mathrm{rad}}(10^{-2})=9\): realizability sits below
Weyl ceiling as expected \cite{hansen1988,harrington2001,jackson1999}.

\paragraph{Problem (Weyl versus SVD).}
\textbf{Problem.} Weyl estimate gives \(N_{\mathrm{mode}}=30\); SVD gives \(r_{\mathrm{rad}}=12\). Harmonize.

\textbf{Step 1.} Apply Eq.~\eqref{eq:ch4.weyl-to-realizability} \cite{stratton1941,hansen1988}.

\textbf{Step 2.} Report \(r_{\mathrm{rad}}=12\) as operative rank \cite{harrington2001}.

\textbf{Step 3.} Use \(N_{\mathrm{mode}}=30\) as upper bound only \cite{devaney2004,jackson1999}.

\textbf{Physical meaning.} Weyl mode count bounds; singular values certify \cite{hansen1988}.

\paragraph{Validity.}
Eqs.~\eqref{eq:ch4.truncation-volume-scaling}--\eqref{eq:ch4.truncation-area-scaling} are asymptotic for \(k_{0}D_{s}\gtrsim1\)
\cite{jackson1999,hansen1988}. Dispersion and material resonances modify prefactors without changing exponents on fixed \(\mathcal{J}_{\mathrm{real}}\)
class \cite{stratton1941,harrington2001}. Loss may increase effective rank slightly by broadening spectra \cite{devaney2004}.


\subsection{Decay of high-order coupling}
\label{sec:ch463}
% TOC tag: [HOME][USE SS3.5.3]

Mode density on \(S^{2}\) in Subsection~\ref{sec:ch353} counts how many angular channels may be labeled at cutoff \(\ell_{\max}\)
\cite{jackson1999,hansen1988}. Subsection~\ref{sec:ch463} converts that count into a realizability law on
\(\boldsymbol{\mathcal{K}}_{\mathrm{geo}}\) from Eq.~\eqref{eq:ch4.geometry-coupling-operator}: listable channels are not exportable
channels \cite{stratton1941,harrington2001,devaney2004}. The philosophical insight is that DOS measures coordinate cardinality; coupling tails
measure watts. Truncation is inevitable because tail mass decays exponentially while labels grow quadratically in \(\ell_{\max}\).

\paragraph{Assumptions.}
Fix compact \(\Omega_{s}\) with \(\operatorname{diam}=D_{s}\), flux gauge, and coupling operator \(\boldsymbol{\mathcal{K}}_{\mathrm{geo}}\)
\cite{stratton1941,hansen1988}. Import Eq.~\eqref{eq:ch4.coupling-ell-decay} from Section~\ref{sec:ch45} by reference
\cite{harrington2001,devaney2004}. Narrowband linear exterior \cite{jackson1999}.

\paragraph{Coupling tail envelope.}
For \(\ell>k_{0}D_{s}\),
\begin{equation}
\label{eq:ch4.coupling-tail-envelope}
\sup_{\Jimp\in\mathcal{J}_{\mathrm{real}}}
\frac{|\mathcal{K}_{\ell m}^{(\sigma)}[\Jimp]|}{\|\boldsymbol{\mathcal{K}}_{\mathrm{geo}}[\Jimp]\|}
\le
\mathcal{C}_{\mathrm{tail}}(\ell,k_{0}D_{s})\,
\ell^{-1}\exp(-\ell/k_{0}D_{s}),
\end{equation}
extending Eq.~\eqref{eq:ch4.coupling-ell-decay} with sector constants \(\mathcal{C}_{\mathrm{tail}}\)
\cite{hansen1988,stratton1941,jackson1999}. Mathematically, Eq.~\eqref{eq:ch4.coupling-tail-envelope} is a uniform operator bound on
high-\(\ell\) blocks. Physically, it states that sources on \(\Omega_{s}\) cannot strongly excite angular orders far beyond electrical size
\cite{harrington2001,devaney2004}.

\begin{theorem}[Coupling tail implies spectral compression]
\label{thm:ch4.coupling-tail-compression}
If Eq.~\eqref{eq:ch4.coupling-tail-envelope} holds, then
\(r_{\mathrm{rad}}(\varepsilon)\le C(k_{0}D_{s},\varepsilon)\) with \(C\) growing slower than raw channel count \((2\ell_{\max}+1)^{2}\)
from Subsection~\ref{sec:ch353} \cite{hansen1988,stratton1941,harrington2001}. Listing size overstates synthesis dimension when
\(\ell_{\max}\gg k_{0}D_{s}\) \cite{devaney2004}.
\end{theorem}
\begin{proof}
Summing Eq.~\eqref{eq:ch4.coupling-tail-envelope} over \((m,\sigma)\) yields exponentially small tail mass for \(\ell\gg k_{0}D_{s}\)
\cite{hansen1988}. Comparing summed mass to \(\varepsilon\sigma_{1}\) in Definition~\ref{def:ch4.radiative-rank} bounds
\(r_{\mathrm{rad}}(\varepsilon)\) independently of \(\ell_{\max}\) \cite{harrington2001,stratton1941}.
\end{proof}

\paragraph{Cumulative coupling tail mass.}
Define
\begin{equation}
\label{eq:ch4.coupling-tail-mass}
\mathcal{M}_{\mathrm{cpl}}(\ell_{\mathrm{cut}})=\sum_{\ell>\ell_{\mathrm{cut}}}(2\ell+1)\max_{m,\sigma}|\mathcal{K}_{\ell m}^{(\sigma)}|^{2}.
\end{equation}
\begin{proposition}[Coupling tail versus DOS]
\label{prop:ch4.coupling-tail-dos}
\(\mathcal{M}_{\mathrm{cpl}}(\ell_{\mathrm{cut}})\) decays exponentially for \(\ell_{\mathrm{cut}}\gg k_{0}D_{s}\) while DOS grows as
\(\mathcal{O}(\ell_{\max}^{2})\) \cite{hansen1988,stratton1941,harrington2001}.
\end{proposition}

\paragraph{Tail mass integral bound.}
Integrating the envelope gives \(\mathcal{M}_{\mathrm{cpl}}(\ell_{\mathrm{cut}})\lesssim\mathcal{C}_{\mathrm{tail}}\exp(-\ell_{\mathrm{cut}}/k_{0}D_{s})\)
\cite{hansen1988,stratton1941}. At \(\ell_{\mathrm{cut}}=\ell_{\mathrm{supp}}(\varepsilon)\),
\begin{equation}
\label{eq:ch4.tail-mass-bound}
\mathcal{M}_{\mathrm{cpl}}(\ell_{\mathrm{supp}})\le\varepsilon^{2}\sigma_{1}^{2},\qquad
\sum_{\ell\le\ell_{\mathrm{supp}}}(2\ell+1)\ge d_{\mathrm{eff}}(\varepsilon).
\end{equation}
Equation~\eqref{eq:ch4.tail-mass-bound} links angular cutoff to export threshold \cite{harrington2001,devaney2004}.

\begin{figure}[t]
\centering
\ChOneFigBlock{%
\begin{tikzpicture}[ch1fig, node distance=7mm and 9mm]
  \node[ch1box] (dos) {DOS $\sim\ell_{\max}^{2}$};
  \node[ch1box, right=11mm of dos] (tail) {Eq.~\eqref{eq:ch4.coupling-tail-envelope}};
  \node[ch1box, right=11mm of tail] (rr) {$r_{\mathrm{rad}}$};
  \draw[ch1flow] (dos) -- (tail) -- (rr);
\end{tikzpicture}%
}
\caption{Subsection~\ref{sec:ch463}: channel DOS grows with listing; coupling tails compress export rank.}
\label{fig:ch4.coupling-tail-compression}
\end{figure}

Figure~\ref{fig:ch4.coupling-tail-compression} is the conceptual compression map: designers must not equate channel labels with independent
beams \cite{hansen1988,harrington2001}.

\paragraph{Worked illustration (high channel count, low rank).}
At \(\ell_{\max}=20\), \(k_{0}D_{s}=1.5\), Subsection~\ref{sec:ch353} lists hundreds of labels but \(r_{\mathrm{rad}}(10^{-2})\approx10\)
by Theorem~\ref{thm:ch4.coupling-tail-compression} \cite{hansen1988,harrington2001,devaney2004}.

\paragraph{Worked illustration (tail mass fraction).}
At \(\ell_{\max}=18\), \(k_{0}D_{s}=1.2\), only \(\ell\le5\) carry \(99\%\) of \(\mathcal{M}_{\mathrm{cpl}}\) \cite{hansen1988,devaney2004}.

\paragraph{Worked numerics (cutoff sweep).}
At \(k_{0}D_{s}=1.5\), \(\ell_{\mathrm{cut}}=8\) yields \(\mathcal{M}_{\mathrm{cpl}}\lesssim10^{-2}\) of total coupling mass
\cite{hansen1988,harrington2001}.

\paragraph{Problem (coupling versus DOS).}
\textbf{Problem.} Channel DOS grows as \(\ell_{\max}^{2}\) but measured diversity saturates. Reconcile.

\textbf{Step 1.} Distinguish listable channels from couplings above \(\varepsilon\sigma_{1}\) \cite{hansen1988,stratton1941}.

\textbf{Step 2.} Apply Theorem~\ref{thm:ch4.coupling-tail-compression} \cite{harrington2001,devaney2004}.

\textbf{Step 3.} Report \(r_{\mathrm{rad}}(\varepsilon)\), not raw DOS \cite{jackson1999}.

\textbf{Physical meaning.} DOS counts labels; coupling tails count watts \cite{hansen1988}.

\paragraph{Problem (MoM channel inflation).}
\textbf{Problem.} MoM uses \(\ell_{\max}=25\); SVD shows \(r_{\mathrm{rad}}(10^{-2})=11\). Explain.

\textbf{Step 1.} Compute \(\mathcal{M}_{\mathrm{cpl}}(\ell_{\mathrm{supp}})\) via Eq.~\eqref{eq:ch4.coupling-tail-mass} \cite{hansen1988}.

\textbf{Step 2.} Identify \(\ell_{\mathrm{supp}}\) from Eq.~\eqref{eq:ch4.ell-max-supported} \cite{harrington2001}.

\textbf{Step 3.} Truncate \(\mathbf{T}\) to \(r_{\mathrm{rad}}\) rows for synthesis \cite{devaney2004,stratton1941}.

\textbf{Physical meaning.} Matrix dimension is coordinate freedom; rank is realizability \cite{harrington2001}.

\paragraph{Problem (anisotropic feed).}
\textbf{Problem.} Offset feed raises \(\mathbf{N}\)-sector coupling at \(\ell=4\) but not \(\mathbf{M}\)-sector. Effect on rank?

\textbf{Step 1.} Evaluate sector-resolved \(\mathcal{M}_{\mathrm{cpl}}^{(\sigma)}\) \cite{harrington2001,hansen1988}.

\textbf{Step 2.} Sum active sectors only for \(r_{\mathrm{rad}}\) \cite{stratton1941}.

\textbf{Step 3.} Report sector ceilings before combined OAM labels \cite{devaney2004}.

\textbf{Physical meaning.} Tail decay is sector-specific; DOS is not \cite{hansen1988}.

\paragraph{Compression optimality link.}
Rank-\(r_{\mathrm{rad}}(\varepsilon)\) truncation is optimal for \(\boldsymbol{\mathcal{K}}_{\mathrm{geo}}\) in operator norm
\cite{harrington2001,devaney2004}. Retaining \(\ell_{\max}\gg\ell_{\mathrm{supp}}\) adds null directions that MoM may misinterpret as
coupled channels \cite{stratton1941,hansen1988}.

\paragraph{DOS inflation diagnostic.}
Define the \emph{listing inflation factor}
\begin{equation}
\label{eq:ch4.dos-inflation}
\mathcal{I}_{\mathrm{DOS}}
=
\frac{(2\ell_{\max}+1)^{2}}{d_{\mathrm{eff}}(\varepsilon)}
\approx
\frac{\text{channel labels}}{\text{export directions}}.
\end{equation}
When \(\mathcal{I}_{\mathrm{DOS}}\gg1\), MoM practitioners are fitting noise in padding rows \cite{harrington2001,devaney2004,hansen1988}.
Theorem~\ref{thm:ch4.coupling-tail-compression} explains why \(\mathcal{I}_{\mathrm{DOS}}\) grows with \(\ell_{\max}\) at fixed \(k_{0}D_{s}\)
\cite{stratton1941,jackson1999}.

\begin{corollary}[Safe listing width]
\label{cor:ch4.safe-listing-width}
A conservative listing satisfies \(\ell_{\max}\le\ell_{\mathrm{supp}}(\varepsilon,k_{0}D_{s})+\mathcal{O}(1)\) in the aligned gauge
\cite{hansen1988,harrington2001}. Wider listings do not enlarge \(\mathcal{S}_{\mathrm{real}}(\mathcal{C})\) \cite{devaney2004,stratton1941}.
\end{corollary}

\paragraph{Worked illustration (inflation factor).}
\(\ell_{\max}=20\), \(k_{0}D_{s}=1.2\), \(d_{\mathrm{eff}}(10^{-2})=9\) gives \(\mathcal{I}_{\mathrm{DOS}}\approx(41)^{2}/9\approx190\)
\cite{hansen1988,harrington2001,devaney2004}. Ninety-nine percent of matrix rows are synthesis-invisible at tolerance \cite{stratton1941}.

\paragraph{Problem (safe \(\ell_{\max}\)).}
\textbf{Problem.} Choose \(\ell_{\max}\) for MoM without padding inflation.

\textbf{Step 1.} Compute \(\ell_{\mathrm{supp}}(10^{-2},k_{0}D_{s})\) \cite{hansen1988,harrington2001}.

\textbf{Step 2.} Set \(\ell_{\max}=\ell_{\mathrm{supp}}+\Delta\) with \(\Delta\le2\) \cite{devaney2004}.

\textbf{Step 3.} Verify \(\mathcal{I}_{\mathrm{DOS}}\lesssim10\) \cite{stratton1941,jackson1999}.

\textbf{Physical meaning.} Listing width should track supported order, not numerical habit \cite{hansen1988}.

\paragraph{Philosophical closure (labels versus watts).}
Subsection~\ref{sec:ch353} answers how many angular channels the sphere can label at cutoff \(\ell_{\max}\) \cite{jackson1999,hansen1988}.
Subsection~\ref{sec:ch463} answers how many of those labels \(\boldsymbol{\mathcal{K}}_{\mathrm{geo}}\) can excite from \(\mathcal{J}_{\mathrm{real}}\)
\cite{stratton1941,harrington2001,devaney2004}. The gap between the two is not a numerical nuisance; it is the realizability content of
truncation. Designers who equate ``channel count'' with ``beam count'' import a representation measure into a power measure
\cite{hansen1988,harrington2001}. Coupling-tail laws restore the correct unit: export rank counts watts above \(\varepsilon\sigma_{1}\), not
entries in a harmonic listing \cite{stratton1941,devaney2004}.

\paragraph{Tail summation bound (explicit).}
Summing Eq.~\eqref{eq:ch4.coupling-tail-envelope} over \(\ell>\ell_{\mathrm{cut}}\) and using \(\sum_{\ell}(2\ell+1)\ell^{-2}\le\pi^{2}/6\)
gives
\begin{equation}
\label{eq:ch4.tail-summation-explicit}
\mathcal{M}_{\mathrm{cpl}}(\ell_{\mathrm{cut}})
\le
\mathcal{C}_{\mathrm{tail}}'\exp\!\bigl(-\ell_{\mathrm{cut}}/k_{0}D_{s}\bigr)\,\|\boldsymbol{\mathcal{K}}_{\mathrm{geo}}\|^{2},
\end{equation}
with \(\mathcal{C}_{\mathrm{tail}}'\) depending only on sector normalization \cite{hansen1988,stratton1941}. At \(\ell_{\mathrm{cut}}=\ell_{\mathrm{supp}}\),
Eq.~\eqref{eq:ch4.tail-summation-explicit} certifies that discarded orders carry less than \(\varepsilon^{2}\) of peak coupling
\cite{harrington2001,devaney2004}.

\paragraph{Worked illustration (summation bound).}
\(k_{0}D_{s}=1.8\), \(\ell_{\mathrm{supp}}=6\): Eq.~\eqref{eq:ch4.tail-summation-explicit} gives tail mass \(\lesssim0.04\|\boldsymbol{\mathcal{K}}_{\mathrm{geo}}\|^{2}\),
consistent with \(\varepsilon=10^{-2}\) \cite{hansen1988,harrington2001}.

\paragraph{Problem (tail bound certification).}
\textbf{Problem.} Certify that \(\ell>\ell_{\mathrm{supp}}\) sectors carry \(<1\%\) coupling power.

\textbf{Step 1.} Compute \(\mathcal{M}_{\mathrm{cpl}}(\ell_{\mathrm{supp}})\) via Eq.~\eqref{eq:ch4.coupling-tail-mass} \cite{hansen1988}.

\textbf{Step 2.} Compare to Eq.~\eqref{eq:ch4.tail-summation-explicit} \cite{stratton1941,harrington2001}.

\textbf{Step 3.} Issue certificate only if ratio \(<10^{-2}\) \cite{devaney2004}.

\textbf{Physical meaning.} Truncation at \(\ell_{\mathrm{supp}}\) sacrifices less than one percent of coupling watts \cite{hansen1988}.

\paragraph{Compression ratio (pass 3).}
Define realizability compression
\begin{equation}
\label{eq:ch4.realizability-compression}
\mathcal{C}_{\mathrm{real}}
=
\frac{d_{\mathrm{eff}}(\varepsilon)}{(2\ell_{\max}+1)^{2}}
\approx
\frac{\text{export directions}}{\text{listed channels}}.
\end{equation}
When \(\mathcal{C}_{\mathrm{real}}\ll1\), Theorem~\ref{thm:ch4.coupling-tail-compression} has already fired: most channels are coupling-invisible
\cite{hansen1988,harrington2001,devaney2004,stratton1941}. MIMO marketing that cites \((2\ell_{\max}+1)^{2}\) without \(\mathcal{C}_{\mathrm{real}}\)
conflates DOS with diversity \cite{jackson1999,hansen1988}.

\paragraph{Worked illustration (compression).}
\(\ell_{\max}=18\), \(d_{\mathrm{eff}}(10^{-2})=8\) gives \(\mathcal{C}_{\mathrm{real}}\approx8/1369\approx0.006\): effective diversity is
sub-percent of listed channels \cite{harrington2001,devaney2004}.

\paragraph{Problem (compression report).}
\textbf{Problem.} Report \(\mathcal{C}_{\mathrm{real}}\) for a VSH synthesis memo.

\textbf{Step 1.} Read \(d_{\mathrm{eff}}(\varepsilon)\) from SVD \cite{harrington2001,hansen1988}.

\textbf{Step 2.} Compute Eq.~\eqref{eq:ch4.realizability-compression} \cite{stratton1941}.

\textbf{Step 3.} Attach \(\mathcal{C}_{\mathrm{real}}\) next to \(\ell_{\max}\) \cite{devaney2004}.

\textbf{Physical meaning.} Compression ratio is the diversity honesty metric \cite{hansen1988}.

\paragraph{Validity.}
Eq.~\eqref{eq:ch4.coupling-tail-envelope} is narrowband and assumes isotropic exterior \cite{stratton1941}. Anisotropic feeds modify
\(\mathcal{C}_{\mathrm{tail}}\) without removing exponential decay \cite{harrington2001}. Loss broadens resonances but does not convert
sub-threshold \(\ell\) sectors into export directions on passive supports \cite{devaney2004,stratton1941}.


\subsection{Effective rank and electromagnetic dimension}
\label{sec:ch464}
% TOC tag: [HOME]

Radiative rank \(r_{\mathrm{rad}}(\varepsilon)\) names export directions in singular-value language. Laboratories, array spec sheets, and MIMO
reports instead quote a single integer: the number of modes \cite{harrington2001,devaney2004}. Subsection~\ref{sec:ch464} locks that integer to
\(d_{\mathrm{eff}}(\varepsilon)\) on \(\boldsymbol{\Gamma}_\omega\) \cite{stratton1941,hansen1988}. The philosophical closure is that
effective dimension is not patch count, not \(\ell_{\max}\), and not channel DOS---it is the geometry-imposed synthesis dimension certified at
tolerance \(\varepsilon\) \cite{hansen1988,harrington2001,devaney2004}.

\paragraph{Assumptions.}
Presuppose Definition~\ref{def:ch4.radiative-rank}, Eq.~\eqref{eq:ch4.gamma-SVD}, and Eq.~\eqref{eq:ch4.effective-dimension-preview}
\cite{harrington2001,devaney2004,hansen1988}. Tolerance \(\varepsilon\) is declared jointly for synthesis and certification \cite{stratton1941}.

\begin{definition}[Effective electromagnetic dimension]
\label{def:ch4.effective-em-dimension}
The \emph{effective electromagnetic dimension} of \((\Omega_{s},\omega,\mathcal{C})\) at tolerance \(\varepsilon\) is
\begin{equation}
\label{eq:ch4.effective-em-dimension}
d_{\mathrm{eff}}(\varepsilon)
=
r_{\mathrm{rad}}(\varepsilon)
=
\#\big\{n:\ \sigma_{n}\ge\varepsilon\,\sigma_{1}\big\},
\end{equation}
with \(\{\sigma_{n}\}\) from Eq.~\eqref{eq:ch4.gamma-SVD} \cite{stratton1941,harrington2001,hansen1988}.
\end{definition}

Equation~\eqref{eq:ch4.effective-em-dimension} elevates the preview Eq.~\eqref{eq:ch4.effective-dimension-preview} to the Chapter~\ref{ch:04}
name for geometry-imposed synthesis dimension \cite{devaney2004}. Mathematically, \(d_{\mathrm{eff}}\) is a discrete counting functional on
the singular spectrum. Physically, it is the number of independent far-zone degrees of freedom a finite source can export at declared power
resolution \(\varepsilon\) \cite{hansen1988,jackson1999}.

\begin{theorem}[Dimension bounds realizability]
\label{thm:ch4.dimension-bounds-realizability}
\(\dim\mathcal{S}_{\mathrm{real}}(\mathcal{C})\le d_{\mathrm{eff}}(\varepsilon)\) at tolerance \(\varepsilon\), and every
\(\mathbf{c}\in\mathcal{S}_{\mathrm{real}}(\mathcal{C})\) has at most \(d_{\mathrm{eff}}(\varepsilon)\) independent components above export
threshold \cite{stratton1941,harrington2001,devaney2004,hansen1988}.
\end{theorem}
\begin{proof}
Theorem~\ref{thm:ch4.radiative-rank} identifies \(\dim\mathcal{S}_{\mathrm{real}}(\mathcal{C})\) with the export rank above threshold
\cite{harrington2001}. Definition~\ref{def:ch4.effective-em-dimension} counts the same singular directions \cite{stratton1941}.
\end{proof}

\paragraph{Dimension reporting protocol.}
Reports must state \((\varepsilon,\,k_{0}D_{s},\,\mathcal{C}\text{ gauge})\) with \(d_{\mathrm{eff}}\)
\cite{hansen1988,stratton1941,harrington2001}.
\begin{equation}
\label{eq:ch4.deff-reporting}
d_{\mathrm{eff}}(\varepsilon)=\#\{n:\sigma_{n}\ge\varepsilon\sigma_{1}\},\qquad
\Delta d_{\mathrm{eff}}=|d_{\mathrm{eff}}^{\mathrm{(MoM)}}-d_{\mathrm{eff}}^{\mathrm{(cont)}}|.
\end{equation}
\begin{corollary}[Mesh gap on dimension]
\label{cor:ch4.deff-mesh-gap}
\(\Delta d_{\mathrm{eff}}\le\Delta r_{\mathrm{gap}}\) in Eq.~\eqref{eq:ch4.rank-gap} \cite{harrington2001,devaney2004,stratton1941}.
\end{corollary}
\begin{equation}
\label{eq:ch4.deff-conservative}
d_{\mathrm{eff}}^{\mathrm{(report)}}=\min\big(d_{\mathrm{eff}}^{\mathrm{(MoM)}},\,r_{\mathrm{rad}}^{\mathrm{(cont)}}(\varepsilon)\big).
\end{equation}
Equation~\eqref{eq:ch4.deff-conservative} is the conservative laboratory integer when mesh rank exceeds continuum rank
\cite{hansen1988,harrington2001}.

\begin{figure}[t]
\centering
\ChOneFigBlock{%
\begin{tikzpicture}[ch1fig, node distance=7mm and 9mm]
  \node[ch1box] (svd) {Eq.~\eqref{eq:ch4.gamma-SVD}};
  \node[ch1box, right=11mm of svd] (rr) {$r_{\mathrm{rad}}$};
  \node[ch1box, right=11mm of rr] (deff) {$d_{\mathrm{eff}}$};
  \node[ch1box, right=11mm of deff] (sr) {$\mathcal{S}_{\mathrm{real}}$};
  \draw[ch1flow] (svd) -- (rr) -- (deff) -- (sr);
\end{tikzpicture}%
}
\caption{Subsection~\ref{sec:ch464}: effective electromagnetic dimension is radiative rank in the export gauge.}
\label{fig:ch4.effective-em-dimension}
\end{figure}

Figure~\ref{fig:ch4.effective-em-dimension} is the reporting chain auditors expect: singular values \(\to\) radiative rank \(\to\)
\(d_{\mathrm{eff}}\) \(\to\) realizable coefficient set \cite{hansen1988,stratton1941}.

\paragraph{Worked illustration (MIMO claim).}
A \(16\times16\) patch array claims \(256\) synthesis DOF; \(d_{\mathrm{eff}}(10^{-2})=22\) on enclosing ball at \(k_{0}a=3\)
\cite{harrington2001,collin2001,hansen1988}. MIMO rank must be reported as \(d_{\mathrm{eff}}\), not element count \cite{devaney2004}.

\paragraph{Worked illustration (mesh saturation).}
MoM shows \(d_{\mathrm{eff}}^{\mathrm{(MoM)}}=45\) but continuum audit gives \(r_{\mathrm{rad}}(10^{-2})=38\): report \(d_{\mathrm{eff}}=38\)
with \(\Delta d_{\mathrm{eff}}=7\) \cite{harrington2001,hansen1988}.

\paragraph{Worked numerics (handset).}
\(k_{0}D_{s}\approx1\) at sub-6\,GHz gives \(d_{\mathrm{eff}}(10^{-2})\le3\); same outline at \(28\,\mathrm{GHz}\) gives
\(d_{\mathrm{eff}}\sim\mathcal{O}(10^{1})\)--\(\mathcal{O}(10^{2})\) per Eq.~\eqref{eq:ch4.joint-scaling-law} \cite{jackson1999,harrington2001,collin2001}.

\paragraph{Problem (dimension from MoM).}
\textbf{Problem.} SVD of \(\mathbf{T}\) shows \(180\) singular values above \(10^{-3}\). Is \(d_{\mathrm{eff}}(10^{-3})=180\)?

\textbf{Step 1.} Compare to continuum \(\boldsymbol{\Gamma}_\omega\) per Lemma~\ref{lem:ch4.rank-gap-aliasing} \cite{harrington2001}.

\textbf{Step 2.} Apply Eq.~\eqref{eq:ch4.deff-conservative} \cite{hansen1988,devaney2004}.

\textbf{Step 3.} Report \(\varepsilon\), \(k_{0}D_{s}\), and \(\Delta d_{\mathrm{eff}}\) with \(d_{\mathrm{eff}}\) \cite{stratton1941}.

\textbf{Physical meaning.} \(d_{\mathrm{eff}}\) is a certified integer, not a solver output \cite{hansen1988}.

\paragraph{Problem (reporting package).}
\textbf{Problem.} What must accompany \(d_{\mathrm{eff}}=25\) in a design memo?

\textbf{Step 1.} State \(\varepsilon\) and \(\sigma_{1}\) gauge \cite{hansen1988}.

\textbf{Step 2.} State \(k_{0}D_{s}\) and support class \cite{jackson1999,stratton1941}.

\textbf{Step 3.} Bound \(\Delta d_{\mathrm{eff}}\) and cite Eq.~\eqref{eq:ch4.joint-scaling-law} \cite{harrington2001,devaney2004}.

\textbf{Physical meaning.} Dimension without tolerance and geometry is not auditable \cite{stratton1941}.

\paragraph{Problem (elements versus modes).}
\textbf{Problem.} \(64\) elements claim \(64\) independent beams. Audit.

\textbf{Step 1.} Compute \(d_{\mathrm{eff}}(10^{-2})\) by SVD \cite{harrington2001}.

\textbf{Step 2.} Bound by Eq.~\eqref{eq:ch4.mode-count-area} \cite{hansen1988,collin2001}.

\textbf{Step 3.} Report \(\min(N_{\mathrm{mode}},d_{\mathrm{eff}})\) streams \cite{devaney2004}.

\textbf{Physical meaning.} Elements sample the aperture; \(d_{\mathrm{eff}}\) is geometry-bounded export rank \cite{stratton1941}.

\paragraph{Link to \(\ell_{\mathrm{supp}}\).}
Eq.~\eqref{eq:ch4.ell-deff-relation} shows \(d_{\mathrm{eff}}\le\sum_{\ell\le\ell_{\mathrm{supp}}}(2\ell+1)\times\#\{\text{active }\sigma\}\)
\cite{harrington2001}. Angular ceiling and integer dimension must be reported together \cite{hansen1988,stratton1941}.

\paragraph{Tolerance sweep certificate.}
Laboratories should report \(d_{\mathrm{eff}}\) as a function of \(\varepsilon\) on fixed \((\Omega_{s},k_{0})\):
\begin{equation}
\label{eq:ch4.deff-epsilon-sweep}
d_{\mathrm{eff}}(\varepsilon_{2})\le d_{\mathrm{eff}}(\varepsilon_{1})\quad\text{for }\varepsilon_{2}\ge\varepsilon_{1},
\end{equation}
with \(d_{\mathrm{eff}}(10^{-1})\) a coarse upper bound and \(d_{\mathrm{eff}}(10^{-3})\) a fine audit \cite{hansen1988,harrington2001,devaney2004}.
A rank claim at unstated \(\varepsilon\) is not comparable across groups \cite{stratton1941}.

\begin{theorem}[Monotone dimension under tolerance tightening]
\label{thm:ch4.deff-monotone-epsilon}
For fixed \((\Omega_{s},k_{0},\mathcal{C})\), \(d_{\mathrm{eff}}(\varepsilon)\) is non-increasing in \(\varepsilon\)
\cite{harrington2001,hansen1988,stratton1941}.
\end{theorem}
\begin{proof}
The index set \(\{n:\sigma_{n}\ge\varepsilon\sigma_{1}\}\) shrinks as \(\varepsilon\) increases \cite{harrington2001}.
\end{proof}

\paragraph{Worked illustration (tolerance sweep).}
At \(k_{0}a=2.5\): \(d_{\mathrm{eff}}(10^{-1})=14\), \(d_{\mathrm{eff}}(10^{-2})=9\), \(d_{\mathrm{eff}}(10^{-3})=12\)
\cite{hansen1988,harrington2001}. Reporting only \(d_{\mathrm{eff}}(10^{-3})=12\) without \(\varepsilon\) misstates the certified rank at
production tolerance \(10^{-2}\) \cite{devaney2004,stratton1941}.

\paragraph{Problem (tolerance mismatch).}
\textbf{Problem.} Paper A uses \(\varepsilon=10^{-3}\); Paper B uses \(10^{-2}\) on the same hardware. Compare ranks.

\textbf{Step 1.} Recompute both at common \(\varepsilon=10^{-2}\) \cite{hansen1988}.

\textbf{Step 2.} Apply Theorem~\ref{thm:ch4.deff-monotone-epsilon} \cite{harrington2001}.

\textbf{Step 3.} Report sweep Eq.~\eqref{eq:ch4.deff-epsilon-sweep} in comparison table \cite{stratton1941,devaney2004}.

\textbf{Physical meaning.} \(d_{\mathrm{eff}}\) is meaningless without paired \(\varepsilon\) \cite{hansen1988}.

\paragraph{Philosophical closure (the laboratory integer).}
\(d_{\mathrm{eff}}(\varepsilon)\) is the point where geometry, frequency, and tolerance meet in a single auditable integer
\cite{stratton1941,hansen1988,harrington2001,devaney2004}. It is not an artifact of MoM discretization if \(\Delta d_{\mathrm{eff}}\) is
reported; it is not a marketing mode count; it is not \(\ell_{\max}\). It is the number of independent export directions a finite source can
synthesize above \(\varepsilon\sigma_{1}\) \cite{hansen1988}. Every MIMO, OAM, and inverse-design claim that omits \(d_{\mathrm{eff}}\) while
quoting element or channel counts is incomplete at the realizability layer \cite{harrington2001,devaney2004,stratton1941}.

\paragraph{Unified truncation certificate.}
Collecting Subsections~\ref{sec:ch461}--\ref{sec:ch464}, a source \((\Omega_{s},k_{0},\mathcal{C},\varepsilon)\) satisfies the
\emph{truncation certificate} when
\begin{equation}
\label{eq:ch4.truncation-certificate}
\ell_{\mathrm{supp}}(\varepsilon,k_{0}D_{s})\ \text{defined},\quad
d_{\mathrm{eff}}(\varepsilon)\le\min\{C_{V}(k_{0}a)^{3},C_{A}k_{0}^{2}A_{a}\},\quad
\mathcal{M}_{\mathrm{cpl}}(\ell_{\mathrm{supp}})\le\varepsilon^{2}\sigma_{1}^{2},\quad
\Delta d_{\mathrm{eff}}\ \text{reported}.
\end{equation}
Equation~\eqref{eq:ch4.truncation-certificate} is the Section~\ref{sec:ch46} audit checklist: angular ceiling, scaling bound, tail mass, and
mesh gap \cite{hansen1988,harrington2001,stratton1941,devaney2004}.

\paragraph{Worked illustration (full certificate).}
\(k_{0}a=2.5\), \(\varepsilon=10^{-2}\): \(\ell_{\mathrm{supp}}=7\), \(d_{\mathrm{eff}}=11\), \(\mathcal{M}_{\mathrm{cpl}}(7)<10^{-4}\sigma_{1}^{2}\),
\(\Delta d_{\mathrm{eff}}=1\) \cite{hansen1988,harrington2001,devaney2004}. Certificate passes; synthesis reports may quote \(d_{\mathrm{eff}}=11\)
with tolerance and geometry \cite{stratton1941}.

\paragraph{Problem (truncation certificate).}
\textbf{Problem.} Prepare a truncation certificate for a circular aperture array spec.

\textbf{Step 1.} Compute \(\ell_{\mathrm{supp}}\), \(d_{\mathrm{eff}}\), \(\mathcal{M}_{\mathrm{cpl}}(\ell_{\mathrm{supp}})\) \cite{hansen1988,harrington2001}.

\textbf{Step 2.} Verify Eq.~\eqref{eq:ch4.joint-scaling-law} and Eq.~\eqref{eq:ch4.truncation-certificate} \cite{stratton1941,devaney2004}.

\textbf{Step 3.} Attach \((\varepsilon,k_{0}D_{s},\Delta d_{\mathrm{eff}})\) to the spec sheet \cite{harrington2001}.

\textbf{Physical meaning.} Truncation is certifiable only as a quadruple inequality, not as a single integer \cite{hansen1988}.

\paragraph{Section closure (pass 3 synthesis).}
Subsections~\ref{sec:ch461}--\ref{sec:ch464} form a closed realizability arc on truncation. \(\ell_{\mathrm{supp}}\) names the angular ceiling;
scaling exponents name how fast \(d_{\mathrm{eff}}\) grows with electrical size; coupling tails explain why channel DOS overstates that growth;
\(d_{\mathrm{eff}}\) is the laboratory integer that survives singular-value audit \cite{stratton1941,hansen1988,harrington2001,devaney2004,jackson1999}.
No step is optional: omitting any one leaves a class of false realizability claims on the table \cite{hansen1988}.

\begin{corollary}[Truncation inevitability on finite hardware]
\label{cor:ch4.truncation-inevitability}
On any compact \(\Omega_{s}\) at fixed \((k_{0},\varepsilon)\), \(\mathcal{S}_{\mathrm{real}}(\mathcal{C})\) is finite-dimensional with
dimension \(\le d_{\mathrm{eff}}(\varepsilon)<\infty\) \cite{stratton1941,harrington2001,hansen1988,devaney2004}. Infinite harmonic listings are
coordinate artifacts; export is always finite-rank \cite{jackson1999}.
\end{corollary}

\paragraph{Worked illustration (inevitability).}
Even at \(\ell_{\max}=100\), \(k_{0}D_{s}=1\) gives \(d_{\mathrm{eff}}(10^{-2})\approx9\): the listing suggests \(10^{4}\) labels; realizability
admits nine export directions \cite{hansen1988,harrington2001,devaney2004}. Corollary~\ref{cor:ch4.truncation-inevitability} is the philosophical
bottom line of Section~\ref{sec:ch46} \cite{stratton1941}.

\paragraph{Validity.}
Definition~\ref{def:ch4.effective-em-dimension} depends on \(\mathcal{C}\) normalization; \(\varepsilon\) must be stated with \(\sigma_{1}\)
\cite{hansen1988,stratton1941}. Loss lowers tail decay and may increase \(d_{\mathrm{eff}}\) at fixed geometry \cite{harrington2001}. Active
amplification changes \(\mathcal{J}_{\mathrm{real}}\) class and is outside the passive default \cite{stratton1941,devaney2004}.



Section~\ref{sec:ch46} has replaced heuristic truncation by \(\ell_{\mathrm{supp}}\), joint volume--area scaling exponents, coupling-tail
compression, and the certified integer \(d_{\mathrm{eff}}(\varepsilon)\) \cite{stratton1941,hansen1988,harrington2001,devaney2004}. Export
rank is now named in both angular-order and singular-value language; laboratories can audit synthesis reports without conflating listing width
with watts \cite{hansen1988}. Section~\ref{sec:ch47} treats null spaces, reactive accessibility, and inversion limits on the same operator
chain \cite{stratton1941,harrington2001,devaney2004}.


\section{Null Spaces, Accessibility, and Reconstruction}
\label{sec:ch47}

Export rank \(d_{\mathrm{eff}}(\varepsilon)\) counts synthesizable directions on \(\mathcal{F}_{\mathrm{rad}}\) \cite{stratton1941,harrington2001,hansen1988}.
Null spaces, reactive subspaces, and finite observation rank determine what can be reconstructed, inverted, or certified from measured data
\cite{devaney2004}. Section~\ref{sec:ch47} develops that split without re-opening the operator definitions of Chapters~\ref{ch:02}--\ref{ch:03}:
non-radiating and unreachable modes alter near fields without enlarging \(\mathbf{c}\); reactive accessibility certifies storage without
enlarging synthesis rank; inversion limits separate numerical fits from range membership; localization--resolution inequalities bound simultaneous
spatial and angular claims \cite{stratton1941,harrington2001,jackson1999,hansen1988}. The philosophical through-line is that export rank is
necessary but not sufficient for observability: measurement closes on intersections, not on listings alone \cite{devaney2004}.


\subsection{Non-radiating and unreachable modes}
\label{sec:ch471}
% TOC tag: [USE]

Non-radiating currents and unreachable far-zone components occupy structural nullities of the export map
\(\boldsymbol{\Gamma}_\omega=\mathcal{C}\circ\mathcal{R}_\omega\) established in Section~\ref{sec:ch27} \cite{stratton1941,harrington2001}.
Subsection~\ref{sec:ch471} quantifies how those nullities partition synthesis space: near-field observables respond to
\(\ker\mathcal{R}_\omega\) and to unreachable coefficient directions, while \(\mathbf{c}\) on \(\mathcal{F}_{\mathrm{rad}}\) does not
\cite{devaney2004,hansen1988}. The distinction is membership, not magnitude---a sector can be exactly uncoupled to every admissible
\(\Jimp\in\mathcal{J}_{\mathrm{real}}\) while still appearing in a VSH listing \cite{stratton1941,harrington2001}. Export rank counts
reachable directions; nullity counts directions the chart prints but geometry cannot excite \cite{hansen1988,devaney2004}.

\paragraph{Assumptions.}
Import \(\ker\mathcal{R}_\omega\), Eq.~\eqref{eq:ch2.nr-current-def}, and Eq.~\eqref{eq:ch4.synthesis-null-space} \cite{stratton1941,harrington2001}.
\(\boldsymbol{\Gamma}_\omega\) is compact on \(\mathcal{J}_{\mathrm{real}}\); tolerance \(\varepsilon\in(0,1)\) fixed \cite{devaney2004,hansen1988}.
Synthesis audits use flux-normalized \(\mathcal{C}\) on \(\mathcal{F}_{\mathrm{rad}}\) \cite{hansen1988}. NR definitions are not re-derived here
\cite{stratton1941,harrington2001}.

\paragraph{Unreachable synthesis directions.}
Define the unreachable subspace
\begin{equation}
\label{eq:ch4.unreachable-modes}
\mathcal{U}_{\mathrm{unreach}}=\ker\boldsymbol{\Gamma}_\omega^{\dagger}\cap\mathbb{C}^{I},
\end{equation}
comprising coefficient directions orthogonal to every image \(\boldsymbol{\Gamma}_\omega[\Jimp]\), \(\Jimp\in\mathcal{J}_{\mathrm{real}}\)
\cite{stratton1941,harrington2001,devaney2004}. Mathematically, Eq.~\eqref{eq:ch4.unreachable-modes} is the coefficient-side null space of the
export adjoint. Physically, it is the set of pattern components no passive source on \(\Omega_{s}\) can radiate, regardless of feed optimization
\cite{hansen1988,harrington2001}. Targets with nonzero \(\mathcal{U}_{\mathrm{unreach}}\) component violate \(\delta_{\mathrm{syn}}=0\) in
Eq.~\eqref{eq:ch4.synthesis-residual-norm} regardless of MoM residual \cite{devaney2004}.

\paragraph{Nullity dimension on \(B_{a}\).}
For ball support \(B_{a}\) at \(k_{0}a\gtrsim1\),
\begin{equation}
\label{eq:ch4.unreach-dim-bound}
\dim\mathcal{U}_{\mathrm{unreach}}\ \ge\ (2\ell_{\max}+1)^{2}-d_{\mathrm{eff}}(\varepsilon)-\dim\ker\mathcal{R}_\omega,
\end{equation}
because listing cardinality exceeds export rank whenever \(\ell_{\max}\) exceeds \(\ell_{\mathrm{supp}}(\varepsilon,k_{0}a)\)
\cite{hansen1988,harrington2001,stratton1941}. On centered feeds, symmetry can force \(\dim\mathcal{U}_{\mathrm{unreach}}\) to exceed half of
\(|\mathbb{C}^{I}|\) while \(d_{\mathrm{eff}}\le4\) \cite{jackson1999,hansen1988}.

\begin{theorem}[Non-radiating lift does not enlarge \(\mathbf{c}\)]
\label{thm:ch4.nr-lift-invisibility}
For \(\mathbf{J}_{\mathrm{nr}}\in\ker\mathcal{R}_\omega\cap\mathcal{J}_{\mathrm{real}}\) and \(\Jimp\in\mathcal{J}_{\mathrm{real}}\),
\(\boldsymbol{\Gamma}_\omega[\Jimp+\mathbf{J}_{\mathrm{nr}}]=\boldsymbol{\Gamma}_\omega[\Jimp]\) on \(\mathcal{F}_{\mathrm{rad}}\)
\cite{stratton1941,harrington2001}. Near-field data distinguish \(\mathbf{J}_{\mathrm{nr}}\); far-zone \(\mathbf{c}\) does not \cite{devaney2004}.
\end{theorem}
\begin{proof}
Linearity: \(\mathcal{R}_\omega[\Jimp+\mathbf{J}_{\mathrm{nr}}]=\mathcal{R}_\omega[\Jimp]\) since \(\mathcal{R}_\omega[\mathbf{J}_{\mathrm{nr}}]=\mathbf{0}\)
\cite{stratton1941}. Bounded \(\mathcal{C}\) on \(\mathcal{F}_{\mathrm{rad}}\) gives identical \(\mathbf{c}\) \cite{hansen1988,harrington2001}. Near-field
functionals on \(\mathcal{F}_{\mathrm{ev}}\) retain NR sensitivity \cite{devaney2004}.
\end{proof}

Theorem~\ref{thm:ch4.nr-lift-invisibility} is the synthesis invisibility law: NR superposition is a current-space ambiguity, not a far-zone degree of
freedom \cite{stratton1941,harrington2001}. Inverse solvers that add NR templates alter near-field fits without changing \(\mathbf{c}\)
\cite{devaney2004,hansen1988}.

\paragraph{NR--unreachable composition.}
Synthesis non-uniqueness aggregates NR lifts and unreachable directions:
\begin{equation}
\label{eq:ch4.nr-unreachable-composition}
\mathcal{N}_{\mathrm{syn}}\supseteq\ker\mathcal{R}_\omega\oplus\mathcal{U}_{\mathrm{unreach}},
\end{equation}
with certificates valid only on quotient \(\mathbb{C}^{I}/\mathcal{N}_{\mathrm{syn}}\) \cite{stratton1941,harrington2001,devaney2004}.
\begin{lemma}[Near-field discrimination bound]
\label{lem:ch4.nf-discrimination}
Distinct \(\Jimp_{A},\Jimp_{B}\) with identical \(\mathbf{c}\) require near-field probes resolving \(\ker\mathcal{R}_\omega\) to discriminate
\cite{harrington2001,stratton1941}.
\end{lemma}

\paragraph{Reachable subspace projection.}
Every target decomposes as
\begin{equation}
\label{eq:ch4.reachable-projection}
\mathbf{c}_{\mathrm{target}}=\mathcal{P}_{\mathrm{reach}}\mathbf{c}_{\mathrm{target}}+\mathbf{c}_{\mathrm{unreach}},\qquad
\mathbf{c}_{\mathrm{unreach}}\in\mathcal{U}_{\mathrm{unreach}},
\end{equation}
with \(\mathcal{P}_{\mathrm{reach}}=\boldsymbol{\Gamma}_\omega\boldsymbol{\Gamma}_\omega^{\dagger}\) \cite{harrington2001,stratton1941,devaney2004}.
Realizability requires \(\mathbf{c}_{\mathrm{unreach}}=\mathbf{0}\) \cite{hansen1988}.

\begin{theorem}[Unreachable component obstruction]
\label{thm:ch4.unreachable-obstruction}
If \(\|\mathbf{c}_{\mathrm{unreach}}\|_{\mathcal{C}}>\varepsilon\|\mathbf{c}_{\mathrm{target}}\|\), then
\(\mathbf{c}_{\mathrm{target}}\notin\mathcal{S}_{\mathrm{real}}(\mathcal{C})\) at tolerance \(\varepsilon\), independent of feed optimization
\cite{stratton1941,harrington2001,devaney2004,hansen1988}.
\end{theorem}
\begin{proof}
Eq.~\eqref{eq:ch4.reachable-projection}; unreachable component lies outside \(\operatorname{ran}\boldsymbol{\Gamma}_\omega\) for all \(\Jimp\)
\cite{stratton1941,harrington2001}. Synthesis residual Eq.~\eqref{eq:ch4.synthesis-residual-norm} lower-bounds \(\|\mathbf{c}_{\mathrm{unreach}}\|\)
\cite{devaney2004,hansen1988}.
\end{proof}

Theorem~\ref{thm:ch4.unreachable-obstruction} closes the synthesis gate before any optimizer runs: unreachable weight is orthogonal to every admissible
\(\Jimp\) portrait and therefore survives feed reweighting, adjoint iteration, and Tikhonov damping unchanged in the \(\mathcal{C}\)-gauge
\cite{stratton1941,harrington2001,devaney2004}. Certification memos must report \(\|\mathbf{c}_{\mathrm{unreach}}\|/\|\mathbf{c}_{\mathrm{target}}\|\)
alongside MoM residuals \cite{hansen1988}.

\paragraph{Unreachable dimension bound.}
On symmetric charts with \(\ell_{\max}\gg\ell_{\mathrm{supp}}(\varepsilon,k_{0}a)\),
\begin{equation}
\label{eq:ch4.unreach-dim-bound}
\dim\mathcal{U}_{\mathrm{unreach}}\ \ge\ |\mathbb{C}^{I}|-d_{\mathrm{eff}}(\varepsilon)-\dim\ker\mathcal{R}_\omega,
\end{equation}
so large nullity is expected rather than pathological when listing exceeds geometry-supported rank \cite{hansen1988,harrington2001,stratton1941,devaney2004}.

\paragraph{Fredholm split.}
Compact \(\boldsymbol{\Gamma}_\omega\) yields Fredholm alternative: \(\mathbf{c}\in\operatorname{ran}\boldsymbol{\Gamma}_\omega\) iff
\(\mathbf{c}\perp\mathcal{U}_{\mathrm{unreach}}\) \cite{harrington2001,stratton1941}. Eq.~\eqref{eq:ch4.reachable-projection} is the finite-dimensional
realization \cite{devaney2004}.

\paragraph{Quotient synthesis space.}
\begin{equation}
\label{eq:ch4.quotient-synthesis}
\mathbf{c}_{1}=\mathbf{c}_{2}\ \text{in }\mathbb{C}^{I}/\mathcal{U}_{\mathrm{unreach}}
\iff
\delta_{\mathrm{syn}}(\mathbf{c}_{1}-\mathbf{c}_{2})=0.
\end{equation}
Uniqueness of \(\Jimp\) is not claimed; uniqueness of export class is \cite{harrington2001,stratton1941}.

\begin{figure}[t]
\centering
\ChOneFigBlock{%
\begin{tikzpicture}[ch1fig, node distance=7mm and 9mm]
  \node[ch1box] (nr) {$\ker\mathcal{R}_\omega$};
  \node[ch1box, right=11mm of nr] (un) {$\mathcal{U}_{\mathrm{unreach}}$};
  \node[ch1box, right=11mm of un] (c) {$\mathbf{c}$};
  \node[ch1box, below=11mm of un] (nf) {near field};
  \draw[ch1flow] (nr) -- (c);
  \draw[ch1flow] (un) -- (c);
  \draw[ch1flow] (nf) -- (nr);
  \draw[ch1flow] (nf) -- (un);
\end{tikzpicture}%
}
\caption{Subsection~\ref{sec:ch471}: NR and unreachable components are invisible to \(\mathbf{c}\) but visible to near probes.}
\label{fig:ch4.nr-unreachable-split}
\end{figure}

Figure~\ref{fig:ch4.nr-unreachable-split} maps three observability channels: far-zone \(\mathbf{c}\) collapses NR and unreachable components; near probes
retain sensitivity \cite{stratton1941,harrington2001,devaney2004}.

\paragraph{NR template audit.}
Non-radiating templates from Eq.~\eqref{eq:ch2.nr-current-def} span \(\ker\mathcal{R}_\omega\cap\mathcal{J}_{\mathrm{real}}\)
\cite{stratton1941,harrington2001}. Adding templates before inversion leaves \(\mathbf{c}\) fixed per Theorem~\ref{thm:ch4.nr-lift-invisibility} but
changes near-field response \cite{devaney2004,hansen1988}.

\paragraph{Worked illustration (unreachable obstruction).}
\(\mathbf{c}_{\mathrm{target}}\) with \(\|\mathbf{c}_{\mathrm{unreach}}\|/\|\mathbf{c}_{\mathrm{target}}\|=0.14\) at \(\varepsilon=10^{-2}\): MoM residual
\(3\times10^{-5}\) but Theorem~\ref{thm:ch4.unreachable-obstruction} rejects realizability \cite{harrington2001,devaney2004,stratton1941,hansen1988}.

\begin{figure}[t]
\centering
\ChOneFigBlock{%
\begin{tikzpicture}[ch1fig, node distance=7mm and 9mm]
  \node[ch1box] (far) {$\mathcal{F}_{\mathrm{rad}}$};
  \node[ch1box, right=11mm of far] (meas) {measurement};
  \node[ch1box, below=9mm of far] (nr) {$\ker\mathcal{R}_\omega$};
  \node[ch1box, below=9mm of meas] (un) {$\mathcal{U}_{\mathrm{unreach}}$};
  \draw[ch1flow] (far) -- (meas);
  \draw[ch1flow, dashed] (nr) -- (far);
  \draw[ch1flow, dashed] (un) -- (far);
\end{tikzpicture}%
}
\caption{Subsection~\ref{sec:ch471}: far-zone measurements intersect \(\operatorname{ran}\boldsymbol{\Gamma}_\omega\) only.}
\label{fig:ch4.measurement-intersection}
\end{figure}

\paragraph{Null-space diagnostic sequence.}
Workflow: (i) form \(\boldsymbol{\Gamma}_\omega\); (ii) compute \(\mathcal{P}_{\mathrm{reach}}\); (iii) decompose \(\mathbf{c}_{\mathrm{target}}\) via
Eq.~\eqref{eq:ch4.reachable-projection}; (iv) test Theorem~\ref{thm:ch4.unreachable-obstruction}; (v) report \(\dim\mathcal{U}_{\mathrm{unreach}}\) and
\(d_{\mathrm{eff}}\) jointly \cite{harrington2001,stratton1941,devaney2004,hansen1988}.

\paragraph{Worked illustration (NR superposition).}
\(\Jimp_{B}=\Jimp_{A}+\mathbf{J}_{\mathrm{nr}}\) leaves \(\mathbf{c}\) unchanged; near reactive fields differ \cite{stratton1941,harrington2001,devaney2004}.

\paragraph{Worked illustration (symmetry null).}
Centered feed: unit weight on \(\widehat{c}_{5,3}^{(\mathbf{M})}\) lies in \(\mathcal{U}_{\mathrm{unreach}}\) \cite{hansen1988,jackson1999,stratton1941}.

\paragraph{Worked illustration (projection audit).}
\(30\%\) weight in \(\mathcal{U}_{\mathrm{unreach}}\): \(\delta_{\mathrm{syn}}>0\) despite unit-norm target \cite{hansen1988,harrington2001,stratton1941}.

\paragraph{Worked illustration (NR tomography trap).}
Near-field tomography recovers \(\mathbf{J}_{\mathrm{nr}}\); far pattern unchanged \cite{stratton1941,harrington2001,devaney2004,hansen1988}.

\paragraph{Worked numerics.}
Centered \(B_{a}\), \(k_{0}a=1\): \(\dim\mathcal{U}_{\mathrm{unreach}}/|\mathbb{C}^{I}|=0.55\), \(d_{\mathrm{eff}}=4\) \cite{hansen1988,harrington2001}.

\paragraph{Problem (unreachable component).}
\textbf{Problem.} Unit weight on symmetry-null sector. Well-posed?
\textbf{Step 1.} Project off \(\mathcal{U}_{\mathrm{unreach}}\) \cite{harrington2001}.
\textbf{Step 2.} Test \(\delta_{\mathrm{syn}}\) on projection \cite{stratton1941,devaney2004}.
\textbf{Step 3.} Report unreachable part as failure \cite{hansen1988}.
\textbf{Physical meaning.} Null directions are structural \cite{stratton1941}.

\paragraph{Problem (NR discrimination).}
\textbf{Problem.} Distinguish \(\Jimp_{A}\) from \(\Jimp_{B}=\Jimp_{A}+\mathbf{J}_{\mathrm{nr}}\) by far-zone data?
\textbf{Step 1.} Test \(\boldsymbol{\Gamma}_\omega[\mathbf{J}_{\mathrm{nr}}]=\mathbf{0}\) \cite{stratton1941}.
\textbf{Step 2.} Conclude identical \(\mathbf{c}\) \cite{harrington2001}.
\textbf{Step 3.} Use near fields per Lemma~\ref{lem:ch4.nf-discrimination} \cite{devaney2004}.
\textbf{Physical meaning.} NR lifts are synthesis-invisible \cite{stratton1941}.

\paragraph{Problem (projection certificate).}
\textbf{Problem.} Certify \(\mathbf{c}_{\mathrm{target}}\) before inversion.
\textbf{Step 1.} Compute \(\mathbf{c}_{\mathrm{unreach}}=(\mathbf{I}-\mathcal{P}_{\mathrm{reach}})\mathbf{c}_{\mathrm{target}}\) \cite{harrington2001}.
\textbf{Step 2.} Test Theorem~\ref{thm:ch4.unreachable-obstruction} \cite{stratton1941,devaney2004}.
\textbf{Step 3.} Reject if \(\|\mathbf{c}_{\mathrm{unreach}}\|>\varepsilon\|\mathbf{c}_{\mathrm{target}}\|\) \cite{hansen1988}.
\textbf{Physical meaning.} Reachable projection is the synthesis gate \cite{stratton1941}.

\paragraph{Problem (quotient reporting).}
\textbf{Problem.} Two currents, same \(\mathbf{c}\). Report uniqueness?
\textbf{Step 1.} State quotient via Eq.~\eqref{eq:ch4.quotient-synthesis} \cite{harrington2001}.
\textbf{Step 2.} Identify \(\dim\ker\mathcal{R}_\omega\) \cite{stratton1941}.
\textbf{Step 3.} Uniqueness on quotient only \cite{devaney2004,hansen1988}.
\textbf{Physical meaning.} \(\mathbf{c}\) is the export certificate \cite{stratton1941}.

\paragraph{Problem (nullity dimension).}
\textbf{Problem.} \(\dim\mathcal{U}_{\mathrm{unreach}}=120\), \(d_{\mathrm{eff}}=6\), \(|\mathbb{C}^{I}|=200\). Consistent?
\textbf{Step 1.} Check Eq.~\eqref{eq:ch4.unreach-dim-bound} \cite{hansen1988,harrington2001}.
\textbf{Step 2.} Compare \(\ell_{\max}\) to \(\ell_{\mathrm{supp}}\) \cite{stratton1941,devaney2004}.
\textbf{Step 3.} Conclude over-listing \cite{jackson1999}.
\textbf{Physical meaning.} Large nullity expected on symmetric charts \cite{harrington2001}.

\paragraph{Problem (MoM null fitting).}
\textbf{Problem.} MoM fits unreachable sector to \(10^{-4}\) residual. Realizable?
\textbf{Step 1.} Decompose via Eq.~\eqref{eq:ch4.reachable-projection} \cite{harrington2001,stratton1941}.
\textbf{Step 2.} Test \(\mathbf{c}_{\mathrm{unreach}}=\mathbf{0}\) \cite{devaney2004}.
\textbf{Step 3.} Reject if unreachable weight exceeds \(\varepsilon\) \cite{hansen1988}.
\textbf{Physical meaning.} MoM fits null directions as noise \cite{harrington2001}.

\paragraph{Philosophical closure (structural nullity).}
Weak coupling is magnitude; unreachable nullity is membership \cite{hansen1988,harrington2001,stratton1941}. Sectors in \(\mathcal{U}_{\mathrm{unreach}}\)
have zero coupling to every admissible \(\Jimp\), not merely small coupling \cite{devaney2004}. Reports that rank by \(|\mathcal{K}_{\alpha}|\) without
null-space tests confuse magnitude with realizability \cite{harrington2001,hansen1988}.

\begin{corollary}[Far-zone blindness to NR lifts]
\label{cor:ch4.farzone-nr-blind}
Any measurement functional supported only on \(\mathcal{F}_{\mathrm{rad}}\) is blind to \(\ker\mathcal{R}_\omega\) on \(\mathcal{J}_{\mathrm{real}}\)
\cite{stratton1941,harrington2001,devaney2004,hansen1988}.
\end{corollary}

\paragraph{Validity.}
Theorem~\ref{thm:ch4.nr-lift-invisibility} assumes linear Maxwell maps on \(\mathcal{J}_{\mathrm{real}}\) \cite{stratton1941}. Loss modifies
\(\ker\mathcal{R}_\omega\) without promoting NR content to \(\mathcal{F}_{\mathrm{rad}}\) \cite{harrington2001,devaney2004}. Active sources enlarge
\(\mathcal{J}_{\mathrm{real}}\) but do not remove unreachable chart nullities on passive synthesis targets \cite{hansen1988}.

\subsection{Reactive subspaces and evanescent accessibility}
\label{sec:ch472}
% TOC tag: [HOME]

Evanescent subspace \(\mathcal{F}_{\mathrm{ev}}\) carries reactive storage accessible to near probes yet invisible to \(\boldsymbol{\Gamma}_\omega\)
on \(\mathcal{F}_{\mathrm{rad}}\) \cite{stratton1941,harrington2001}. Subsection~\ref{sec:ch472} separates reactive accessibility from synthesis
accessibility \cite{devaney2004,hansen1988}. Reactive richness certifies how finely near fields can be resolved and how much energy the near zone can
store; it does not enlarge \(\mathcal{S}_{\mathrm{real}}(\mathcal{C})\) on the propagating chart \cite{jackson1999,stratton1941,harrington2001}.
The philosophical split: storage is not export; probe rank is not synthesis rank \cite{devaney2004,hansen1988}.

\paragraph{Assumptions.}
Presuppose Proposition~\ref{prop:ch4.evanescent-invisibility} and Subsection~\ref{sec:ch273} reconstructability caps \cite{stratton1941,harrington2001}.
Observation rank \(N\) finite; probes linear and calibrated on \(\mathcal{F}_{\mathrm{ev}}\cup\mathcal{F}_{\mathrm{near}}\) \cite{devaney2004,harrington2001}.
Tolerance \(\varepsilon\in(0,1)\) fixed for \(d_{\mathrm{eff}}\) audits \cite{hansen1988}.

\paragraph{Reactive accessibility functional.}
\begin{equation}
\label{eq:ch4.reactive-accessibility}
\mathcal{A}_{\mathrm{react}}^{(N)}(\mathbf{c})=\big\|\mathcal{P}_{\mathrm{react}}\mathbf{c}\big\|_{\mathrm{react}},
\end{equation}
with \(\mathcal{P}_{\mathrm{react}}\) projecting reactive components when near data are used \cite{harrington2001,stratton1941}. Mathematically,
\(\mathcal{A}_{\mathrm{react}}^{(N)}\) measures how much of the field chart lies outside \(\mathcal{F}_{\mathrm{rad}}\) at probe resolution \(N\).
Physically, large \(\mathcal{A}_{\mathrm{react}}^{(N)}\) with small \(\|\mathbf{c}\|_{\mathrm{rad}}\) signals evanescent-dominated certification
\cite{devaney2004,jackson1999}.

\paragraph{Evanescent mode budget.}
Let \(N_{\mathrm{ev}}(k_{t,\max})\) count evanescent modes with \(k_{t}>k_{0}\) resolved up to \(k_{t,\max}\) \cite{jackson1999,hansen1988}.
\(N_{\mathrm{ev}}\) grows with scan resolution while \(d_{\mathrm{eff}}(\varepsilon)\) on \(\mathcal{F}_{\mathrm{rad}}\) remains fixed on
\((\Omega_{s},k_{0})\) \cite{stratton1941,harrington2001,devaney2004}.

\begin{proposition}[Reactive accessibility does not enlarge synthesis]
\label{prop:ch4.reactive-not-synthesis}
\(\mathbf{c}\in\mathcal{S}_{\mathrm{real}}(\mathcal{C})\) is independent of \(\mathcal{A}_{\mathrm{react}}^{(N)}\) on \(\mathcal{F}_{\mathrm{rad}}\)
\cite{stratton1941,harrington2001}. Raising \(N\) may raise \(\mathcal{A}_{\mathrm{react}}^{(N)}\) without enlarging
\(\mathcal{S}_{\mathrm{real}}(\mathcal{C})\) \cite{devaney2004,hansen1988}.
\end{proposition}

\paragraph{Propagating-sector filter.}
Before listing \(\mathbf{c}\), apply
\begin{equation}
\label{eq:ch4.propagating-filter}
\widetilde{\mathbf{c}}=\Pi_{\mathrm{rad}}\Pi_{\mathrm{prop}}\mathcal{C}\big[\mathcal{S}_\omega[\Jimp]\big],
\end{equation}
importing Section~\ref{sec:ch44} by reference \cite{stratton1941,harrington2001}. Scans that skip Eq.~\eqref{eq:ch4.propagating-filter} inflate
\(\mathcal{A}_{\mathrm{react}}^{(N)}\) without changing \(\widetilde{\mathbf{c}}\) \cite{devaney2004,hansen1988}.

\begin{theorem}[Evanescent invisibility to export]
\label{thm:ch4.ev-invisible-export}
\(\|\widetilde{\mathbf{c}}\|_{\mathrm{rad}}\) and \(d_{\mathrm{eff}}(\varepsilon)\) on \(\mathcal{F}_{\mathrm{rad}}\) are independent of
\(\mathcal{A}_{\mathrm{react}}^{(N)}\) on fixed \((\Omega_{s},k_{0})\) \cite{stratton1941,harrington2001,devaney2004,hansen1988}.
\end{theorem}
\begin{proof}
Proposition~\ref{prop:ch4.reactive-not-synthesis} and linearity of \(\Pi_{\mathrm{rad}}\Pi_{\mathrm{prop}}\) \cite{stratton1941,harrington2001}. Evanescent
content lies outside \(\operatorname{ran}\boldsymbol{\Gamma}_\omega\) on \(\mathcal{F}_{\mathrm{rad}}\) \cite{devaney2004,hansen1988}.
\end{proof}

\paragraph{Reactive--radiative split.}
\begin{equation}
\label{eq:ch4.reactive-radiative-split}
\mathcal{J}_{\mathrm{react}}[\Jimp]=\|\mathcal{P}_{\mathrm{react}}\mathcal{S}_\omega[\Jimp]\|,\quad
\mathcal{J}_{\mathrm{rad}}[\Jimp]=\|\Pi_{\mathrm{rad}}\mathcal{S}_\omega[\Jimp]\|_{\mathrm{rad}}.
\end{equation}
\begin{theorem}[Reactive dominance without rank gain]
\label{thm:ch4.reactive-dominance}
\(\mathcal{J}_{\mathrm{react}}\gg\mathcal{J}_{\mathrm{rad}}\) does not imply \(d_{\mathrm{eff}}\) enlargement on \(\mathcal{F}_{\mathrm{rad}}\)
\cite{stratton1941,harrington2001,devaney2004}.
\end{theorem}
\begin{proof}
Eq.~\eqref{eq:ch4.propagating-filter} removes evanescent weight before \(\boldsymbol{\Gamma}_\omega\) \cite{stratton1941,harrington2001}. Reactive
dominance enlarges \(\mathcal{J}_{\mathrm{react}}\) without altering \(\sigma_{n}\) on \(\mathcal{F}_{\mathrm{rad}}\) \cite{devaney2004,hansen1988}.
\end{proof}

Theorem~\ref{thm:ch4.reactive-dominance} is the near-field honesty law: tomography and hotspot scans certify storage and probe resolution, not export
rank on \(\mathcal{F}_{\mathrm{rad}}\) \cite{stratton1941,harrington2001,jackson1999}. Any synthesis memo that equates reactive recovery with diversity
must be rejected before inversion starts \cite{devaney2004,hansen1988}.

\paragraph{Evanescent wavenumber budget.}
Transverse wavenumbers resolved at standoff \(z\) satisfy
\begin{equation}
\label{eq:ch4.ev-wavenumber-budget}
k_{t,\max}(z)\approx k_{0}+\frac{1}{z}\ln\!\left(\frac{1}{\varepsilon_{\mathrm{probe}}}\right),
\end{equation}
so over-resolving \(k_{t}\) inflates \(\mathcal{A}_{\mathrm{react}}^{(N)}\) without shifting \(\mathbf{c}\) \cite{jackson1999,hansen1988,harrington2001,devaney2004}.

\paragraph{Reactive-to-export ratio.}
\begin{equation}
\label{eq:ch4.react-export-ratio}
\mathcal{Q}_{\mathrm{re}}=\frac{\mathcal{A}_{\mathrm{react}}^{(N)}}{\|\mathbf{c}\|_{\mathrm{rad}}+\delta},\qquad \delta>0.
\end{equation}
Large \(\mathcal{Q}_{\mathrm{re}}\) with fixed \(d_{\mathrm{eff}}\) signals reactive-dominated certification \cite{harrington2001,devaney2004,hansen1988,jackson1999}.

\paragraph{Probe transfer bound.}
Near-field transfer matrix \(\mathbf{G}_{N}\) has rank \(\le N\); reactive recovery saturates when \(N\ge N_{\mathrm{ev}}(k_{t,\max})\)
\cite{harrington2001,devaney2004}. \(d_{\mathrm{eff}}\) is bounded by \(\boldsymbol{\Gamma}_\omega\), not \(\mathbf{G}_{N}\) \cite{stratton1941,hansen1988}.

\paragraph{Reactive energy budget.}
\begin{equation}
\label{eq:ch4.reactive-energy-budget}
\mathcal{E}_{\mathrm{react}}^{(N)}=\int_{\Sigma_{N}}\big(\varepsilon_{0}|\E|^{2}+\mu_{0}|\H|^{2}\big)\,dS,
\end{equation}
grows with \(N\) and \(k_{t,\max}\) but does not enter \(\boldsymbol{\Gamma}_\omega\) on \(\mathcal{F}_{\mathrm{rad}}\) \cite{jackson1999,stratton1941,harrington2001,devaney2004}.

\begin{lemma}[Probe saturation before rank gain]
\label{lem:ch4.probe-sat-lemma}
\(\mathcal{A}_{\mathrm{react}}^{(N)}\) saturates as \(N\) increases while \(d_{\mathrm{eff}}(\varepsilon)\) remains constant on fixed \((\Omega_{s},k_{0})\)
\cite{harrington2001,devaney2004,hansen1988,stratton1941}.
\end{lemma}

\begin{figure}[t]
\centering
\ChOneFigBlock{%
\begin{tikzpicture}[ch1fig, node distance=7mm and 9mm]
  \node[ch1box] (ev) {$\mathcal{F}_{\mathrm{ev}}$};
  \node[ch1box, right=11mm of ev] (react) {$\mathcal{A}_{\mathrm{react}}^{(N)}$};
  \node[ch1box, right=11mm of react] (rad) {$\mathbf{c}$};
  \draw[ch1flow] (ev) -- (react);
  \draw[ch1flow] (react) -- (rad);
\end{tikzpicture}%
}
\caption{Subsection~\ref{sec:ch472}: reactive accessibility measures storage; synthesis uses propagating \(\mathbf{c}\) only.}
\label{fig:ch4.reactive-accessibility-split}
\end{figure}

\paragraph{Probe-rank saturation.}
\begin{equation}
\label{eq:ch4.probe-saturation}
\lim_{N\to\infty}\mathcal{A}_{\mathrm{react}}^{(N)}=\|\mathcal{P}_{\mathrm{react}}\mathcal{S}_\omega[\Jimp]\|,\qquad
\mathbf{c}_{N\to\infty}=\boldsymbol{\Gamma}_\omega[\Jimp].
\end{equation}
Saturation of \(\mathcal{A}_{\mathrm{react}}^{(N)}\) does not shift \(\mathbf{c}_{N\to\infty}\) \cite{stratton1941,harrington2001,devaney2004}.

\begin{figure}[t]
\centering
\ChOneFigBlock{%
\begin{tikzpicture}[ch1fig, node distance=7mm and 9mm]
  \node[ch1box] (probe) {$N$ probes};
  \node[ch1box, right=11mm of probe] (areact) {$\mathcal{A}_{\mathrm{react}}^{(N)}$};
  \node[ch1box, right=11mm of areact] (deff) {$d_{\mathrm{eff}}$};
  \draw[ch1flow] (probe) -- (areact);
  \draw[ch1flow, dashed] (areact) -- (deff);
\end{tikzpicture}%
}
\caption{Subsection~\ref{sec:ch472}: probe count raises \(\mathcal{A}_{\mathrm{react}}^{(N)}\) until saturation; \(d_{\mathrm{eff}}\) unchanged.}
\label{fig:ch4.probe-saturation-map}
\end{figure}

Figure~\ref{fig:ch4.probe-saturation-map} shows the dashed link: reactive accessibility does not drive export rank \cite{harrington2001,stratton1941,devaney2004}.
Probe budgets should therefore be reported as \((N,\mathcal{A}_{\mathrm{react}}^{(N)},d_{\mathrm{eff}})\) triples rather than as rank surrogates
\cite{hansen1988,devaney2004}.

\begin{corollary}[Propagating filter persistence]
\label{cor:ch4.prop-filter-persist}
\(\widetilde{\mathbf{c}}=\Pi_{\mathrm{rad}}\mathbf{c}\) is invariant under reactive probe enrichment on fixed \((\Omega_{s},k_{0})\)
\cite{stratton1941,harrington2001,devaney2004,hansen1988}.
\end{corollary}

\paragraph{Evanescent coupling envelope.}
For \(k_{t}>k_{0}\), probe coupling decays as \(\exp(-\alpha z)\) with standoff \(z\) \cite{jackson1999,hansen1988}. Over-resolving \(k_{t}\) without
reducing \(z\) inflates \(\mathcal{A}_{\mathrm{react}}^{(N)}\) without improving far-zone rank \cite{harrington2001,devaney2004}.

\paragraph{Worked illustration (near scan).}
Scan at \(k_{t}=1.2k_{0}\) raises \(\mathcal{A}_{\mathrm{react}}^{(N)}\); \(\mathbf{c}\) stays dipole-dominated \cite{harrington2001,jackson1999,hansen1988}.

\paragraph{Worked illustration (resonant particle).}
\(\mathcal{J}_{\mathrm{react}}/\mathcal{J}_{\mathrm{rad}}=10\), \(d_{\mathrm{eff}}=3\) unchanged \cite{jackson1999,hansen1988,harrington2001,stratton1941}.

\paragraph{Worked illustration (scan over-resolution).}
\(k_{t}=1.5k_{0}\): \(\mathcal{Q}_{\mathrm{re}}\approx8\), \(d_{\mathrm{eff}}=6\) unchanged \cite{harrington2001,jackson1999,hansen1988}.

\paragraph{Worked illustration (resonator).}
High-\(Q\) particle: \(\mathcal{J}_{\mathrm{react}}\gg\mathcal{J}_{\mathrm{rad}}\), \(\mathbf{c}\) dipole-class \cite{jackson1999,stratton1941,harrington2001}.

\paragraph{Worked numerics (probe saturation).}
\(N=50\) on \(\lambda/4\) grid: \(\mathcal{A}_{\mathrm{react}}^{(N)}\) within \(2\%\) of limit; \(d_{\mathrm{eff}}(10^{-2})=7\) for \(N=25,50,100\)
\cite{harrington2001,hansen1988,devaney2004}.

\paragraph{Worked illustration (standoff audit).}
At \(z=\lambda/8\), Eq.~\eqref{eq:ch4.ev-wavenumber-budget} permits \(k_{t,\max}\approx1.4k_{0}\); doubling probes raises \(\mathcal{Q}_{\mathrm{re}}\) from
\(4.2\) to \(6.8\) while \(d_{\mathrm{eff}}=5\) is unchanged \cite{jackson1999,harrington2001,devaney2004,hansen1988}.

\paragraph{Problem (reactive versus synthesis).}
\textbf{Problem.} Near tomography recovers reactive content. Enlarge \(\mathcal{S}_{\mathrm{real}}\)?
\textbf{Step 1.} Apply Eq.~\eqref{eq:ch4.propagating-filter} \cite{stratton1941,harrington2001}.
\textbf{Step 2.} Test \(\delta_{\mathrm{syn}}\) on \(\widetilde{\mathbf{c}}\) \cite{devaney2004,hansen1988}.
\textbf{Step 3.} Treat reactive recovery as certification only \cite{harrington2001}.
\textbf{Physical meaning.} Reactive dominance is storage, not export \cite{stratton1941}.

\paragraph{Problem (hotspot audit).}
\textbf{Problem.} Near hotspot, flat far pattern. Rank change?
\textbf{Step 1.} Compute \(\mathbf{c}\) after \(\Pi_{\mathrm{rad}}\) \cite{stratton1941}.
\textbf{Step 2.} Compare \(d_{\mathrm{eff}}\) before/after near fit \cite{hansen1988,harrington2001}.
\textbf{Step 3.} Conclude no rank gain per Theorem~\ref{thm:ch4.reactive-dominance} \cite{devaney2004}.
\textbf{Physical meaning.} Hotspots certify reactance \cite{harrington2001}.

\paragraph{Problem (reactive certificate).}
\textbf{Problem.} Laboratory reports high reactive content. Rank change?
\textbf{Step 1.} Apply Eq.~\eqref{eq:ch4.propagating-filter} \cite{stratton1941,harrington2001}.
\textbf{Step 2.} Recompute \(d_{\mathrm{eff}}\) on \(\widetilde{\mathbf{c}}\) \cite{hansen1988,devaney2004}.
\textbf{Step 3.} Report \(\mathcal{Q}_{\mathrm{re}}\) separately \cite{harrington2001,jackson1999}.
\textbf{Physical meaning.} Reactive certificates \(\neq\) diversity \cite{stratton1941}.

\paragraph{Problem (probe budget).}
\textbf{Problem.} Double probe count. Rank change?
\textbf{Step 1.} Test Theorem~\ref{thm:ch4.ev-invisible-export} \cite{stratton1941,harrington2001}.
\textbf{Step 2.} Track saturation via Eq.~\eqref{eq:ch4.probe-saturation} \cite{devaney2004}.
\textbf{Step 3.} Apply Lemma~\ref{lem:ch4.probe-sat-lemma} \cite{hansen1988}.
\textbf{Physical meaning.} Probes certify reactance \cite{stratton1941}.

\paragraph{Problem (reactive energy versus rank).}
\textbf{Problem.} \(\mathcal{E}_{\mathrm{react}}^{(N)}\) doubles after scan upgrade. Does \(d_{\mathrm{eff}}\) double?
\textbf{Step 1.} Filter via Eq.~\eqref{eq:ch4.propagating-filter} \cite{stratton1941,harrington2001}.
\textbf{Step 2.} Recompute \(d_{\mathrm{eff}}\) on \(\widetilde{\mathbf{c}}\) \cite{hansen1988,devaney2004}.
\textbf{Step 3.} Report \((\mathcal{Q}_{\mathrm{re}},\mathcal{E}_{\mathrm{react}}^{(N)},d_{\mathrm{eff}})\) triple \cite{jackson1999}.
\textbf{Physical meaning.} Stored energy is not export rank \cite{stratton1941}.

\paragraph{Problem (evanescent over-resolution).}
\textbf{Problem.} Scan resolves \(k_{t}=2k_{0}\) at \(z=\lambda/10\). Synthesis impact?
\textbf{Step 1.} Compute \(N_{\mathrm{ev}}\) and \(\mathcal{A}_{\mathrm{react}}^{(N)}\) \cite{jackson1999,hansen1988}.
\textbf{Step 2.} Apply Theorem~\ref{thm:ch4.ev-invisible-export} \cite{stratton1941,harrington2001}.
\textbf{Step 3.} Conclude \(d_{\mathrm{eff}}\) unchanged \cite{devaney2004}.
\textbf{Physical meaning.} Evanescent resolution is not export resolution \cite{hansen1988}.

\paragraph{Philosophical closure (storage versus export).}
Reactive accessibility measures storage and probe resolution; synthesis rank measures independent far-zone directions above \(\varepsilon\sigma_{1}\)
\cite{stratton1941,harrington2001,jackson1999,hansen1988,devaney2004}. Enlarging one does not enlarge the other on passive supports
\cite{stratton1941,harrington2001}.

\begin{corollary}[Reactive saturation before rank gain]
\label{cor:ch4.reactive-saturation}
\(\mathcal{A}_{\mathrm{react}}^{(N)}\) saturates as \(N\) increases while \(d_{\mathrm{eff}}(\varepsilon)\) remains constant on fixed \((\Omega_{s},k_{0})\)
\cite{harrington2001,devaney2004,hansen1988,stratton1941}.
\end{corollary}

\paragraph{Validity.}
Eq.~\eqref{eq:ch4.reactive-accessibility} assumes linear probes and calibration \cite{harrington2001}. Loss modifies \(\mathcal{F}_{\mathrm{ev}}\) split
without exporting evanescent weight to \(\mathbf{c}\) \cite{stratton1941,devaney2004}. Nonlinear probes require separate \(\mathcal{P}_{\mathrm{react}}\) declaration
\cite{hansen1988,jackson1999}.

\subsection{Synthesis, reconstruction, and inversion limits}
\label{sec:ch473}
% TOC tag: [HOME]

Inverse design and array synthesis minimize residuals on \(\mathbf{c}\) or on near-field probes; realizability requires
\(\mathbf{c}\in\operatorname{ran}\boldsymbol{\Gamma}_\omega\) on admissible \(\mathcal{J}_{\mathrm{real}}\) \cite{stratton1941,harrington2001,devaney2004}.
Subsection~\ref{sec:ch473} states inversion limits as operator-range theorems for compact \(\boldsymbol{\Gamma}_\omega\) \cite{hansen1988}. The
central distinction is variational fit versus range membership: optimizers find closest coefficients in \(\mathbb{C}^{I}\); synthesis demands zero
perpendicular residual \(\mathbf{r}_{\perp}\) in Eq.~\eqref{eq:ch4.inversion-residual-decomposition} \cite{stratton1941,harrington2001}. A MoM
residual below \(10^{-4}\) is therefore necessary but not sufficient for \(\delta_{\mathrm{syn}}=0\) \cite{devaney2004,hansen1988}.

\paragraph{Assumptions.}
Presuppose Theorem~\ref{thm:ch4.moore-penrose-synthesis} and Principle~\ref{prin:ch4.operator-chain} \cite{harrington2001,devaney2004,stratton1941}.
Discrete MoM matrix \(\mathbf{T}\) is a Galerkin restriction of \(\boldsymbol{\Gamma}_\omega\) on \(\mathcal{J}_{\mathrm{real}}\) \cite{harrington2001}.
Tolerance \(\varepsilon\in(0,1)\) fixed; flux-normalized \(\mathcal{C}\) on \(\mathcal{F}_{\mathrm{rad}}\) \cite{hansen1988,stratton1941}. Unreachable
nullity from Subsection~\ref{sec:ch471} is imported by reference only \cite{devaney2004}.

\paragraph{Range--null decomposition.}
Every target \(\mathbf{c}_{\mathrm{target}}\in\mathbb{C}^{I}\) admits the orthogonal split
\begin{equation}
\label{eq:ch4.inversion-residual-decomposition}
\mathbf{c}_{\mathrm{target}}=\mathcal{P}_{\mathrm{ran}}\mathbf{c}_{\mathrm{target}}+\mathbf{r}_{\perp},\qquad
\mathbf{r}_{\perp}\in\ker\boldsymbol{\Gamma}_\omega^{\dagger},
\end{equation}
where \(\mathcal{P}_{\mathrm{ran}}=\boldsymbol{\Gamma}_\omega\boldsymbol{\Gamma}_\omega^{\dagger}\) is the Moore--Penrose range projector on the
coefficient chart \cite{harrington2001,stratton1941}. Mathematically, Eq.~\eqref{eq:ch4.inversion-residual-decomposition} is the Fredholm split for
compact \(\boldsymbol{\Gamma}_\omega\). Physically, \(\mathcal{P}_{\mathrm{ran}}\mathbf{c}_{\mathrm{target}}\) is the closest realizable pattern;
\(\mathbf{r}_{\perp}\) is the unreachable component no passive \(\Jimp\) can radiate \cite{devaney2004,hansen1988}.

\paragraph{Operator-range certificate.}
Realizability is equivalent to membership in \(\operatorname{ran}\boldsymbol{\Gamma}_\omega\):
\begin{equation}
\label{eq:ch4.range-certificate}
\mathbf{c}\in\mathcal{S}_{\mathrm{real}}(\mathcal{C})
\iff
\exists\,\Jimp\in\mathcal{J}_{\mathrm{real}}:\ \boldsymbol{\Gamma}_\omega[\Jimp]=\mathbf{c}
\iff
\mathbf{r}_{\perp}=\mathbf{0}.
\end{equation}
Equation~\eqref{eq:ch4.range-certificate} is the synthesis gate every inversion memo must close \cite{stratton1941,harrington2001,devaney2004,hansen1988}.

\begin{theorem}[Inversion limit at rank \(d_{\mathrm{eff}}\)]
\label{thm:ch4.inversion-limit}
Best rank-\(d_{\mathrm{eff}}(\varepsilon)\) approximation to \(\mathbf{c}_{\mathrm{target}}\) achieves error
\(\le\varepsilon\sigma_{1}\|\mathbf{c}_{\mathrm{target}}\|\) on \(\mathcal{P}_{\mathrm{ran}}\mathbf{c}_{\mathrm{target}}\). If
\(\|\mathbf{r}_{\perp}\|_{\mathcal{C}}>\varepsilon\|\mathbf{c}_{\mathrm{target}}\|\), then
\(\mathbf{c}_{\mathrm{target}}\notin\mathcal{S}_{\mathrm{real}}(\mathcal{C})\) at tolerance \(\varepsilon\), independent of optimizer or feed count
\cite{stratton1941,harrington2001,devaney2004,hansen1988}.
\end{theorem}
\begin{proof}
Theorem~\ref{thm:ch4.radiative-rank} bounds truncation error on \(\operatorname{ran}\boldsymbol{\Gamma}_\omega\) by \(\varepsilon\sigma_{1}\)
\cite{harrington2001,stratton1941}. Component \(\mathbf{r}_{\perp}\in\ker\boldsymbol{\Gamma}_\omega^{\dagger}\) lies outside range for all ranks
\cite{devaney2004}. Combining Eckart--Young optimality with Eq.~\eqref{eq:ch4.inversion-residual-decomposition} yields the obstruction clause
\cite{hansen1988,harrington2001}.
\end{proof}

Theorem~\ref{thm:ch4.inversion-limit} elevates inversion from a numerical success criterion to a geometric membership test \cite{stratton1941}. No
amount of feed reweighting, adjoint iteration, or Tikhonov damping moves \(\mathbf{r}_{\perp}\) into \(\operatorname{ran}\boldsymbol{\Gamma}_\omega\)
when \(\delta_{\mathrm{reg}}>0\) in Eq.~\eqref{eq:ch4.tikhonov-synthesis-gap} \cite{harrington2001,devaney2004,hansen1988}.

\paragraph{Tikhonov--synthesis gap.}
Regularized inversion minimizes
\(\|\mathbf{T}\mathbf{I}_{s}-\mathbf{c}\|^{2}+\alpha\|\mathbf{I}_{s}\|^{2}\) but realizability is certified only by
\begin{equation}
\label{eq:ch4.tikhonov-synthesis-gap}
\delta_{\mathrm{reg}}=\|\mathbf{r}_{\perp}\|_{\mathcal{C}},\qquad
\delta_{\mathrm{reg}}=0\iff\mathbf{c}\in\mathcal{S}_{\mathrm{real}}(\mathcal{C}).
\end{equation}
\begin{corollary}[Regularization does not create realizability]
\label{cor:ch4.reg-not-realizability}
\(\alpha>0\) lowers discrete residuals without moving \(\mathbf{c}\) into \(\operatorname{ran}\boldsymbol{\Gamma}_\omega\) when \(\delta_{\mathrm{reg}}>0\)
\cite{harrington2001,devaney2004,stratton1941}.
\end{corollary}

\begin{theorem}[Discrete--continuum inversion gap]
\label{thm:ch4.discrete-continuum-gap}
\(\|\mathbf{T}\mathbf{I}_{s}-\mathbf{c}\|_{\mathcal{C}}<\tau_{\mathrm{mom}}\) does not imply \(\delta_{\mathrm{syn}}=0\) unless
\(\tau_{\mathrm{mom}}<\varepsilon\|\mathbf{c}\|\) and \(\mathbf{r}_{\perp}=\mathbf{0}\) in the continuum gauge of \(\boldsymbol{\Gamma}_\omega\)
\cite{harrington2001,devaney2004,stratton1941,hansen1988}.
\end{theorem}
\begin{proof}
Galerkin discretization error is bounded by mesh resolution on \(\operatorname{ran}\boldsymbol{\Gamma}_\omega\) \cite{harrington2001}. Unreachable
\(\mathbf{r}_{\perp}\) is gauge-independent in \(\ker\boldsymbol{\Gamma}_\omega^{\dagger}\) \cite{stratton1941,devaney2004}.
\end{proof}

\paragraph{Moore--Penrose synthesis chain.}
On truncated rank \(d\),
\begin{equation}
\label{eq:ch4.mom-penrose-chain}
\mathbf{I}_{s}^{(d)}=\mathbf{T}_{d}^{\dagger}\mathbf{c}
=\sum_{n=1}^{d}\frac{\langle\mathbf{u}_{n},\mathbf{c}\rangle}{\sigma_{n}}\,\mathbf{v}_{n},
\end{equation}
where \(\mathbf{T}_{d}=\sum_{n=1}^{d}\sigma_{n}\mathbf{u}_{n}\mathbf{v}_{n}^{\dagger}\) \cite{harrington2001,hansen1988}. Equation
\eqref{eq:ch4.mom-penrose-chain} inverts only the reachable component \(\mathcal{P}_{\mathrm{ran}}^{(d)}\mathbf{c}\); it cannot synthesize
\(\mathbf{r}_{\perp}\) \cite{stratton1941,devaney2004}.

\paragraph{Condition number certificate.}
For rank-\(d\) truncation, \(\kappa(\mathbf{T}_{d})=\sigma_{1}/\sigma_{d}\). Large \(\kappa\) amplifies noise in Eq.~\eqref{eq:ch4.mom-penrose-chain}
without reducing \(\|\mathbf{r}_{\perp}\|\) \cite{hansen1988,harrington2001}. Synthesis audits must report \((\delta_{\mathrm{syn}},\kappa,d_{\mathrm{eff}})\)
as a triple certificate \cite{stratton1941,devaney2004}.

\paragraph{Rank-truncated inversion bound.}
Best rank-\(d_{\mathrm{eff}}\) inverse satisfies
\begin{equation}
\label{eq:ch4.rank-truncated-inverse}
\min_{\mathbf{I}_{s}}\|\mathbf{T}_{d}\mathbf{I}_{s}-\mathcal{P}_{\mathrm{ran}}\mathbf{c}\|
\le
\varepsilon\sigma_{1}\|\mathbf{c}\|,
\end{equation}
but cannot reduce \(\|\mathbf{r}_{\perp}\|\) \cite{harrington2001,hansen1988,stratton1941}. Truncation optimizes within range; it does not create range
\cite{devaney2004}.

\paragraph{Eckart--Young alignment.}
\begin{equation}
\label{eq:ch4.eckart-young-target}
\|\mathbf{c}-\mathcal{P}_{\mathrm{ran}}^{(d_{\mathrm{eff}})}\mathbf{c}\|_{\mathcal{C}}\le\varepsilon\|\mathbf{c}\|
\iff
\mathcal{P}_{\mathrm{ran}}^{(d_{\mathrm{eff}})}\mathbf{c}\in\mathcal{S}_{\mathrm{real}}(\mathcal{C})\ \text{at }\varepsilon.
\end{equation}
Equality in Eq.~\eqref{eq:ch4.eckart-young-target} still fails when \(\mathbf{r}_{\perp}\neq0\) \cite{harrington2001,stratton1941}.

\paragraph{Adjoint inversion trap.}
Adjoint optimizers minimize \(\|\mathbf{T}^{H}\mathbf{c}-\mathbf{w}\|\) on feed weights \(\mathbf{w}\) \cite{harrington2001,devaney2004}. The adjoint
finds a closest current portrait whose export \(\mathbf{T}\mathbf{I}_{s}\) may fit \(\mathcal{P}_{\mathrm{ran}}\mathbf{c}\) while leaving
\(\mathbf{r}_{\perp}\) in the target unchanged \cite{stratton1941,hansen1988}. Certification requires lifting to Eq.~\eqref{eq:ch4.range-certificate},
not reporting adjoint residual alone \cite{harrington2001}.

\begin{lemma}[Sparsity prior obstruction]
\label{lem:ch4.sparsity-obstruction}
\(\ell_{1}\) or cardinality priors on \(\mathbf{I}_{s}\) alter the variational landscape of inversion but do not enlarge
\(\operatorname{ran}\boldsymbol{\Gamma}_\omega\). If \(\mathbf{r}_{\perp}\neq0\), no sparsity penalty yields \(\delta_{\mathrm{syn}}=0\)
\cite{harrington2001,devaney2004,stratton1941,hansen1988}.
\end{lemma}

\begin{figure}[t]
\centering
\ChOneFigBlock{%
\begin{tikzpicture}[ch1fig, node distance=7mm and 9mm]
  \node[ch1box] (ran) {$\operatorname{ran}\boldsymbol{\Gamma}_\omega$};
  \node[ch1box, right=11mm of ran] (perp) {$\mathbf{r}_{\perp}$};
  \node[ch1box, right=11mm of perp] (mom) {MoM residual};
  \draw[ch1flow] (ran) -- (mom);
  \draw[ch1flow] (perp) -- (mom);
\end{tikzpicture}%
}
\caption{Subsection~\ref{sec:ch473}: MoM residual can be small while \(\mathbf{r}_{\perp}\neq0\); synthesis audits must decompose both components.}
\label{fig:ch4.inversion-residual-split}
\end{figure}

Figure~\ref{fig:ch4.inversion-residual-split} is the inversion honesty diagram: the optimizer sees the sum; realizability requires the perpendicular
part to vanish \cite{stratton1941,harrington2001,devaney2004}.

\begin{figure}[t]
\centering
\ChOneFigBlock{%
\begin{tikzpicture}[ch1fig, node distance=7mm and 9mm]
  \node[ch1box] (opt) {optimizer};
  \node[ch1box, right=11mm of opt] (ran) {$\mathcal{P}_{\mathrm{ran}}\mathbf{c}$};
  \node[ch1box, right=11mm of ran] (syn) {$\delta_{\mathrm{syn}}=0$};
  \draw[ch1flow] (opt) -- (ran);
  \draw[ch1flow] (ran) -- (syn);
\end{tikzpicture}%
}
\caption{Subsection~\ref{sec:ch473}: optimization closes on range projection; synthesis certifies zero perpendicular residual.}
\label{fig:ch4.synthesis-certificate-chain}
\end{figure}

\paragraph{Worked illustration (Tikhonov masking).}
Tikhonov with \(\alpha=10^{-3}\) yields \(\|\mathbf{T}\mathbf{I}_{s}-\mathbf{c}\|=3\times10^{-5}\) while \(\|\mathbf{r}_{\perp}\|/\|\mathbf{c}\|=0.09\):
\(\delta_{\mathrm{syn}}>0\) \cite{harrington2001,devaney2004,stratton1941}.

\paragraph{Worked illustration (optimizer trap).}
Adjoint optimizer: residual \(10^{-5}\), \(\|\mathbf{r}_{\perp}\|/\|\mathbf{c}\|=0.12\) \cite{harrington2001,devaney2004}. Synthesis fails despite
excellent fit on \(\mathcal{P}_{\mathrm{ran}}\mathbf{c}\) \cite{stratton1941,hansen1988}.

\paragraph{Worked illustration (rank mismatch).}
Target needs \(d_{\mathrm{eff}}=14\); inversion uses rank \(8\): within-range error \(\approx0.04\|\mathbf{c}\|\) masquerades as unreachable null
until \(\mathcal{P}_{\mathrm{ran}}^{(14)}\) is tested \cite{hansen1988,harrington2001,devaney2004}.

\paragraph{Worked illustration (condition number).}
\(\kappa(\mathbf{T}_{8})=10^{4}\): discrete residual \(10^{-3}\) with \(\|\mathbf{r}_{\perp}\|=0\) is realizable but numerically fragile
\cite{hansen1988,harrington2001}. Report \(\kappa\) alongside \(\delta_{\mathrm{syn}}\) \cite{stratton1941}.

\paragraph{Worked numerics (full decomposition).}
\(\|\mathbf{r}_{\perp}\|/\|\mathbf{c}\|=0.15\) at \(\varepsilon=10^{-2}\) fails synthesis despite MoM residual \(10^{-4}\) and
\(\kappa(\mathbf{T}_{11})=120\) \cite{harrington2001,hansen1988,devaney2004}.

\paragraph{Problem (inversion versus realizability).}
\textbf{Problem.} Optimizer reports \(\|\mathbf{T}\mathbf{I}_{s}-\mathbf{c}\|<10^{-4}\). Certify?
\textbf{Step 1.} Decompose via Eq.~\eqref{eq:ch4.inversion-residual-decomposition} \cite{stratton1941}.
\textbf{Step 2.} Test \(\delta_{\mathrm{syn}}=0\) \cite{harrington2001,devaney2004}.
\textbf{Step 3.} Reject if \(\mathbf{r}_{\perp}\neq0\) \cite{hansen1988}.
\textbf{Physical meaning.} Inversion finds closest coefficients; synthesis demands range membership \cite{stratton1941}.

\paragraph{Problem (regularized fit).}
\textbf{Problem.} Tikhonov yields residual \(<10^{-5}\). Realizable?
\textbf{Step 1.} Test \(\delta_{\mathrm{reg}}=0\) in Eq.~\eqref{eq:ch4.tikhonov-synthesis-gap} \cite{harrington2001}.
\textbf{Step 2.} If \(\delta_{\mathrm{reg}}>0\), reject \cite{devaney2004,stratton1941}.
\textbf{Step 3.} Report \(\mathbf{r}_{\perp}\) norm in memo \cite{hansen1988}.
\textbf{Physical meaning.} Regularization is not realizability \cite{stratton1941}.

\paragraph{Problem (continuum certificate).}
\textbf{Problem.} MoM passes; certify realizability.
\textbf{Step 1.} Lift \(\mathbf{T}\) to continuum \(\boldsymbol{\Gamma}_\omega\) \cite{harrington2001}.
\textbf{Step 2.} Test Eq.~\eqref{eq:ch4.range-certificate} \cite{stratton1941,devaney2004}.
\textbf{Step 3.} Apply Theorem~\ref{thm:ch4.discrete-continuum-gap} \cite{hansen1988}.
\textbf{Physical meaning.} Discrete fit \(\neq\) range membership \cite{stratton1941}.

\paragraph{Problem (rank truncation).}
\textbf{Problem.} Use rank \(d_{\mathrm{eff}}\) in inversion. When valid?
\textbf{Step 1.} Verify \(\mathbf{c}\in\operatorname{ran}\boldsymbol{\Gamma}_\omega^{(d_{\mathrm{eff}})}\) \cite{harrington2001}.
\textbf{Step 2.} Apply Eq.~\eqref{eq:ch4.rank-truncated-inverse} \cite{hansen1988,stratton1941}.
\textbf{Step 3.} Separate truncation error from \(\mathbf{r}_{\perp}\) via Eq.~\eqref{eq:ch4.eckart-young-target} \cite{devaney2004}.
\textbf{Physical meaning.} Truncation is within-range approximation \cite{harrington2001}.

\paragraph{Problem (adjoint audit).}
\textbf{Problem.} Adjoint reports success. Synthesis?
\textbf{Step 1.} Compute \(\mathbf{c}_{\mathrm{fit}}=\mathbf{T}\mathbf{I}_{s}\) \cite{harrington2001}.
\textbf{Step 2.} Decompose \(\mathbf{c}_{\mathrm{target}}-\mathbf{c}_{\mathrm{fit}}\) into \(\mathbf{r}_{\perp}\) \cite{stratton1941,devaney2004}.
\textbf{Step 3.} Certify only if \(\mathbf{r}_{\perp}=\mathbf{0}\) \cite{hansen1988}.
\textbf{Physical meaning.} Adjoint residual is not \(\delta_{\mathrm{syn}}\) \cite{harrington2001}.

\paragraph{Problem (sparsity prior).}
\textbf{Problem.} \(\ell_{1}\) prior yields sparse \(\mathbf{I}_{s}\) and low residual. Realizable?
\textbf{Step 1.} Apply Lemma~\ref{lem:ch4.sparsity-obstruction} \cite{harrington2001,devaney2004}.
\textbf{Step 2.} Test \(\delta_{\mathrm{reg}}=0\) \cite{stratton1941}.
\textbf{Step 3.} Reject if \(\mathbf{r}_{\perp}\neq0\) regardless of sparsity \cite{hansen1988}.
\textbf{Physical meaning.} Priors reshape inversion, not range \cite{stratton1941}.

\paragraph{Philosophical closure (inversion versus synthesis).}
Inversion is a variational problem on \(\mathbb{C}^{I}\); synthesis is a membership problem on \(\operatorname{ran}\boldsymbol{\Gamma}_\omega\)
\cite{stratton1941,harrington2001,devaney2004}. Every inverse-design pipeline must terminate on Eq.~\eqref{eq:ch4.range-certificate}, not on optimizer
cost \cite{hansen1988}. \(\delta_{\mathrm{syn}}=0\) is the synthesis certificate; MoM residual, adjoint score, and Tikhonov norm are diagnostics only
\cite{stratton1941,harrington2001}.

\begin{corollary}[Tikhonov cannot cure unreachable targets]
\label{cor:ch4.tikhonov-unreachable}
If \(\mathbf{r}_{\perp}\neq0\), no \(\alpha>0\) yields \(\delta_{\mathrm{reg}}=0\) in Eq.~\eqref{eq:ch4.tikhonov-synthesis-gap}
\cite{harrington2001,devaney2004,stratton1941,hansen1988}.
\end{corollary}

\paragraph{Validity.}
Theorem~\ref{thm:ch4.inversion-limit} assumes linear forward models on \(\mathcal{J}_{\mathrm{real}}\) \cite{stratton1941}. Nonlinear or active
materials require separate range declarations; \(\mathbf{r}_{\perp}\) obstruction persists in the linearized subproblem \cite{harrington2001,devaney2004}.
Loss modifies \(\sigma_{n}\) ordering without promoting unreachable components into \(\operatorname{ran}\boldsymbol{\Gamma}_\omega\) \cite{hansen1988}.

\subsection{Localization versus angular resolution}
\label{sec:ch474}
% TOC tag: [HOME][USE SS3.8.4][USE SS3.6.4]

Representation dependence of angular labels in Subsection~\ref{sec:ch384} and localization bounds in Subsection~\ref{sec:ch364} constrain
simultaneous spatial concentration and angular resolution on \(\mathcal{F}_{\mathrm{rad}}\) \cite{harrington2001,hansen1988,stratton1941}.
Subsection~\ref{sec:ch474} states the tradeoff as realizability inequalities on compact \(\Omega_{s}\) \cite{devaney2004,jackson1999}. Tight beams
and high \(\ell\) on the same support violate export tail laws unless \(\mathcal{L}_{2}\) and \(\mathcal{R}_{\mathrm{ang}}\) are jointly audited
\cite{hansen1988,harrington2001}. Representation is coordinate freedom; simultaneous localization and resolution are geometry-imposed
\cite{stratton1941,devaney2004}.

\paragraph{Assumptions.}
Import Eq.~\eqref{eq:ch4.localization-spectrum-product} and Eq.~\eqref{eq:ch3.OAM-representation-dependence} \cite{stratton1941,hansen1988,harrington2001}.
Compact \(\Omega_{s}\); flux gauge on \(\mathcal{F}_{\mathrm{rad}}\); tolerance \(\varepsilon\in(0,1)\) \cite{jackson1999,devaney2004}. Paraxial LG
labels require bridge Eq.~\eqref{eq:ch3.paraxial-vs-spherical-OAM} before VSH audit \cite{harrington2001,allen1992}.

\paragraph{Resolution functional.}
Define block angular resolution
\begin{equation}
\label{eq:ch4.angular-resolution-functional}
\mathcal{R}_{\mathrm{ang}}(\mathbf{c})=\sup\left\{\ell:\ \sum_{m,\sigma}|\widehat{c}_{\ell m}^{(\sigma)}|^{2}\ge\eta_{\mathrm{blk}}\|\mathbf{c}\|_{\mathrm{rad}}^{2}\right\},
\end{equation}
the highest \(\ell\) carrying fraction \(\eta_{\mathrm{blk}}\) of radiative energy \cite{hansen1988,harrington2001,stratton1941}. On ball \(B_{a}\),
\begin{equation}
\label{eq:ch4.localization-resolution-tradeoff}
\mathcal{L}_{2}[\Jimp]\,\mathcal{R}_{\mathrm{ang}}(\mathbf{c})\ \gtrsim\ \mathcal{K}_{\mathrm{lr}}(k_{0}a,\eta_{\mathrm{blk}}),
\end{equation}
importing Subsection~\ref{sec:ch364} \cite{jackson1999,hansen1988,devaney2004}. Equation~\eqref{eq:ch4.localization-resolution-tradeoff} forbids
arbitrarily tight \(\mathcal{L}_{2}\) with arbitrarily high \(\mathcal{R}_{\mathrm{ang}}\) on the same \(\Jimp\) \cite{stratton1941,harrington2001}.

\begin{theorem}[Resolution ceiling]
\label{thm:ch4.resolution-ceiling}
\(\mathcal{R}_{\mathrm{ang}}(\mathbf{c})\le\ell_{\mathrm{supp}}(\varepsilon,k_{0}a)\) on \(B_{a}\) \cite{hansen1988,stratton1941,devaney2004,harrington2001}.
\end{theorem}
\begin{proof}
\(\ell_{\mathrm{supp}}\) is the synthesis-side angular cutoff from singular-value ordering \cite{harrington2001}; \(\mathcal{R}_{\mathrm{ang}}\) counts
energy above \(\eta_{\mathrm{blk}}\) in \(\ell\) blocks \cite{hansen1988,stratton1941}. Tail laws cap nonzero \(\widehat{c}_{\ell m}\) for
\(\ell>\ell_{\mathrm{supp}}\) \cite{devaney2004,jackson1999}.
\end{proof}

Theorem~\ref{thm:ch4.resolution-ceiling} converts tail-law decay into an angular-order audit: \(\mathcal{R}_{\mathrm{ang}}\) cannot exceed the
synthesis-supported cutoff \(\ell_{\mathrm{supp}}(\varepsilon,k_{0}a)\) regardless of chart listing or optimizer success \cite{hansen1988,harrington2001,stratton1941}.
Vendor OAM claims must therefore be tested against \(\ell_{\mathrm{supp}}\) before localization inequalities are invoked \cite{devaney2004,jackson1999}.

\paragraph{Supported-order scaling.}
On \(B_{a}\), \(\ell_{\mathrm{supp}}(\varepsilon,k_{0}a)=\mathcal{O}(k_{0}a)\) per Section~\ref{sec:ch46} truncation laws \cite{hansen1988,harrington2001}.
Handset-scale \(k_{0}a<1\) supports only dipole-class \(\ell\) at practical \(\varepsilon\); super-hemispherical \(\mathcal{R}_{\mathrm{ang}}\) is excluded
\cite{jackson1999,stratton1941,devaney2004}.

\paragraph{Simultaneous bound.}
\begin{equation}
\label{eq:ch4.simultaneous-loc-ell}
\mathcal{L}_{2}[\Jimp]\cdot\mathcal{R}_{\mathrm{ang}}(\mathbf{c})\ge\mathcal{K}_{\mathrm{sl}}(k_{0}a,\eta_{\mathrm{blk}}),
\end{equation}
\begin{corollary}[Super-resolution obstruction]
\label{cor:ch4.super-resolution-obstruction}
Simultaneous \(\mathcal{L}_{2}\ll a^{2}\) and \(\mathcal{R}_{\mathrm{ang}}\gg k_{0}a\) implies \(\mathbf{c}\notin\mathcal{S}_{\mathrm{real}}(\mathcal{C})\) on \(B_{a}\)
\cite{hansen1988,stratton1941,devaney2004,jackson1999}.
\end{corollary}

\paragraph{Uncertainty product form.}
Combining Eqs.~\eqref{eq:ch4.localization-resolution-tradeoff} and Theorem~\ref{thm:ch4.resolution-ceiling},
\begin{equation}
\label{eq:ch4.loc-res-uncertainty}
\mathcal{L}_{2}[\Jimp]\cdot\ell_{\mathrm{supp}}(\varepsilon,k_{0}a)
\ge
\mathcal{K}_{\mathrm{unc}}(k_{0}a,\varepsilon),
\end{equation}
an export-side uncertainty relation on compact supports \cite{hansen1988,jackson1999,stratton1941,harrington2001,devaney2004}.

\begin{theorem}[Joint localization--supported-order bound]
\label{thm:ch4.joint-loc-supp}
On \(B_{a}\), any \(\mathbf{c}\in\mathcal{S}_{\mathrm{real}}(\mathcal{C})\) satisfies
\(\mathcal{R}_{\mathrm{ang}}(\mathbf{c})\le\ell_{\mathrm{supp}}(\varepsilon,k_{0}a)\) and Eq.~\eqref{eq:ch4.loc-res-uncertainty} simultaneously
\cite{hansen1988,stratton1941,harrington2001,devaney2004}.
\end{theorem}
\begin{proof}
Theorem~\ref{thm:ch4.resolution-ceiling} and Eq.~\eqref{eq:ch4.localization-resolution-tradeoff} \cite{hansen1988,stratton1941,jackson1999}.
\end{proof}

Theorem~\ref{thm:ch4.joint-loc-supp} is the realizability audit for simultaneous focusing and OAM claims: both \(\mathcal{L}_{2}\) and
\(\mathcal{R}_{\mathrm{ang}}\) must pass, not merely one \cite{harrington2001,hansen1988,devaney2004}.

\paragraph{Beamwidth consistency.}
\begin{equation}
\label{eq:ch4.beamwidth-ell-consistency}
\mathcal{R}_{\mathrm{ang}}(\mathbf{c})\le\ell_{\mathrm{supp}}(\varepsilon,k_{0}a)\le\mathcal{O}(k_{0}a).
\end{equation}
Narrow beamwidth claims require Eq.~\eqref{eq:ch4.ell-max-scaling} before OAM order is certified \cite{hansen1988,harrington2001,jackson1999}.

\paragraph{Representation bridge audit.}
Paraxial LG labels require Eq.~\eqref{eq:ch3.paraxial-vs-spherical-OAM} before \(\mathcal{R}_{\mathrm{ang}}\) tests \cite{harrington2001,allen1992}.
Full-vector decomposition may spread one LG index across multiple \(\ell\) sectors \cite{hansen1988,stratton1941,belinfante1940}. Essentialist OAM
claims that skip the bridge fail realizability audit \cite{devaney2004,harrington2001}.

\begin{lemma}[LG--VSH spread bound]
\label{lem:ch4.lg-vsh-spread}
Paraxial mode \(m_{\mathrm{LG}}\) maps to VSH sectors \(\ell\in\{m_{\mathrm{LG}}-1,m_{\mathrm{LG}},m_{\mathrm{LG}}+1,\ldots\}\) after full-vector
decomposition; \(\mathcal{R}_{\mathrm{ang}}\) must be computed in VSH gauge \cite{harrington2001,hansen1988,allen1992,stratton1941}.
\end{lemma}
\begin{proof}
Full-vector LG decomposition spreads azimuthal order across adjacent \(\ell\) sectors \cite{harrington2001,allen1992}. \(\mathcal{R}_{\mathrm{ang}}\) is
defined on VSH coefficients after gauge conversion, not on paraxial labels alone \cite{hansen1988,stratton1941,belinfante1940}.
\end{proof}

\paragraph{Diameter scaling audit.}
On ball supports \(B_{a}\), supported order scales as \(\ell_{\mathrm{supp}}(\varepsilon,k_{0}a)=\mathcal{O}(k_{0}a)\) per
Theorem~\ref{thm:ch4.resolution-ceiling} \cite{hansen1988,harrington2001,jackson1999}. Claims with \(\mathcal{R}_{\mathrm{ang}}\gg k_{0}a\) fail before
localization tests are applied \cite{stratton1941,devaney2004}.

\begin{figure}[t]
\centering
\ChOneFigBlock{%
\begin{tikzpicture}[ch1fig, node distance=7mm and 9mm]
  \node[ch1box] (loc) {$\mathcal{L}_{2}$};
  \node[ch1box, right=11mm of loc] (rang) {$\mathcal{R}_{\mathrm{ang}}$};
  \node[ch1box, right=11mm of rang] (ell) {$\ell_{\mathrm{supp}}$};
  \draw[ch1flow] (loc) -- (rang) -- (ell);
\end{tikzpicture}%
}
\caption{Subsection~\ref{sec:ch474}: localization and angular resolution jointly bounded by support.}
\label{fig:ch4.loc-resolution-tradeoff}
\end{figure}

\begin{figure}[t]
\centering
\ChOneFigBlock{%
\begin{tikzpicture}[ch1fig, node distance=7mm and 9mm]
  \node[ch1box] (lg) {LG index};
  \node[ch1box, right=11mm of lg] (vsh) {VSH sectors};
  \node[ch1box, right=11mm of vsh] (audit) {$\mathcal{R}_{\mathrm{ang}}$};
  \draw[ch1flow] (lg) -- (vsh) -- (audit);
\end{tikzpicture}%
}
\caption{Subsection~\ref{sec:ch474}: representation bridge before VSH realizability audit.}
\label{fig:ch4.lg-vsh-bridge}
\end{figure}

Figure~\ref{fig:ch4.loc-resolution-tradeoff} chains localization to supported order; Figure~\ref{fig:ch4.lg-vsh-bridge} chains paraxial labels to
VSH audit \cite{harrington2001,hansen1988,stratton1941}.

\paragraph{Worked illustration (tight beam).}
\(0.25\lambda\) aperture, \(2^{\circ}\) beam, \(\ell=8\): violates Eq.~\eqref{eq:ch4.localization-resolution-tradeoff}
\cite{jackson1999,harrington2001,hansen1988,devaney2004}.

\paragraph{Worked illustration (LG label).}
\(m_{\mathrm{LG}}=5\) requires full-vector decomposition before VSH audit \cite{harrington2001,allen1992,hansen1988}.

\paragraph{Worked illustration (LG to VSH).}
\(m_{\mathrm{LG}}=3\) maps to \(\ell\in\{2,3,4\}\); \(\mathcal{R}_{\mathrm{ang}}\) computed in VSH gauge per Lemma~\ref{lem:ch4.lg-vsh-spread}
\cite{harrington2001,hansen1988,allen1992}.

\paragraph{Worked illustration (super-resolution).}
\(k_{0}a=1.5\), \(\mathcal{R}_{\mathrm{ang}}=9\): violates Theorem~\ref{thm:ch4.joint-loc-supp} \cite{hansen1988,jackson1999,harrington2001,stratton1941}.

\paragraph{Worked numerics (joint audit).}
\(k_{0}a=1.8\), \(\mathcal{L}_{2}=0.3a^{2}\), claimed \(\mathcal{R}_{\mathrm{ang}}=7\): fails Eq.~\eqref{eq:ch4.loc-res-uncertainty} at \(\varepsilon=10^{-2}\)
\cite{hansen1988,harrington2001,devaney2004}.

\paragraph{Worked illustration (diameter class).}
Handset \(k_{0}a=0.9\): \(\ell_{\mathrm{supp}}(10^{-2},0.9)\approx2\); claimed \(\mathcal{R}_{\mathrm{ang}}=5\) violates
Theorem~\ref{thm:ch4.joint-loc-supp} before beamwidth audit \cite{hansen1988,jackson1999,harrington2001,stratton1941,devaney2004}.

\begin{corollary}[Simultaneous product obstruction]
\label{cor:ch4.loc-res-product}
If \(\mathcal{L}_{2}[\Jimp]\cdot\ell_{\mathrm{supp}}(\varepsilon,k_{0}a)<\mathcal{K}_{\mathrm{unc}}(k_{0}a,\varepsilon)\), then
\(\mathbf{c}\notin\mathcal{S}_{\mathrm{real}}(\mathcal{C})\) at the stated tolerance \cite{hansen1988,jackson1999,stratton1941,harrington2001,devaney2004}.
\end{corollary}

\paragraph{Problem (joint audit certificate).}
\textbf{Problem.} Certify simultaneous tight focus and \(\ell=6\) on \(k_{0}a=1.6\).
\textbf{Step 1.} Read \(\ell_{\mathrm{supp}}(10^{-2},1.6)\) from SVD \cite{hansen1988,harrington2001}.
\textbf{Step 2.} Test Theorem~\ref{thm:ch4.joint-loc-supp} and Corollary~\ref{cor:ch4.loc-res-product} \cite{jackson1999,stratton1941}.
\textbf{Step 3.} Reject if either \(\mathcal{R}_{\mathrm{ang}}>\ell_{\mathrm{supp}}\) or product bound fails \cite{devaney2004}.
\textbf{Physical meaning.} Joint audits precede OAM marketing claims \cite{harrington2001,hansen1988}.

\paragraph{Problem (OAM transfer).}
\textbf{Problem.} LG index \(m_{\mathrm{LG}}=5\). VSH realizability?
\textbf{Step 1.} Apply Eq.~\eqref{eq:ch3.paraxial-vs-spherical-OAM} \cite{harrington2001}.
\textbf{Step 2.} Test \(\mathcal{R}_{\mathrm{ang}}\) against Eq.~\eqref{eq:ch4.ell-max-scaling} \cite{hansen1988}.
\textbf{Step 3.} Reject essentialist claims \cite{stratton1941,belinfante1940}.
\textbf{Physical meaning.} OAM labels are representation-dependent \cite{harrington2001}.

\paragraph{Problem (resolution claim).}
\textbf{Problem.} \(\mathcal{R}_{\mathrm{ang}}=12\) on \(k_{0}a=2\). Audit.
\textbf{Step 1.} Test \(\ell_{\mathrm{supp}}(10^{-2},2)\) \cite{hansen1988,harrington2001}.
\textbf{Step 2.} Test Eq.~\eqref{eq:ch4.simultaneous-loc-ell} \cite{jackson1999}.
\textbf{Step 3.} Reject if either fails \cite{devaney2004,stratton1941}.
\textbf{Physical meaning.} Resolution is jointly bounded \cite{hansen1988}.

\paragraph{Problem (LG realizability).}
\textbf{Problem.} \(m_{\mathrm{LG}}=4\) on \(k_{0}a=2\). Realizable?
\textbf{Step 1.} Apply Lemma~\ref{lem:ch4.lg-vsh-spread} \cite{harrington2001,allen1992}.
\textbf{Step 2.} Compute \(\mathcal{R}_{\mathrm{ang}}(\mathbf{c})\) \cite{hansen1988}.
\textbf{Step 3.} Test \(\ell_{\mathrm{supp}}(10^{-2},2)\) \cite{stratton1941,devaney2004}.
\textbf{Physical meaning.} OAM order is chart-dependent \cite{harrington2001}.

\paragraph{Problem (beamwidth--\(\ell\) joint audit).}
\textbf{Problem.} \(2^{\circ}\) beam and \(\ell=7\) on \(k_{0}a=1.8\). Feasible?
\textbf{Step 1.} Test Eq.~\eqref{eq:ch4.beamwidth-ell-consistency} \cite{hansen1988,jackson1999}.
\textbf{Step 2.} Test Eq.~\eqref{eq:ch4.loc-res-uncertainty} \cite{stratton1941,harrington2001}.
\textbf{Step 3.} Reject if either fails \cite{devaney2004}.
\textbf{Physical meaning.} Beams and OAM orders are jointly bounded \cite{hansen1988}.

\paragraph{Problem (handset OAM claim).}
\textbf{Problem.} Handset aperture claims \(\ell=6\) OAM at \(k_{0}a=0.8\). Audit.
\textbf{Step 1.} Test Theorem~\ref{thm:ch4.resolution-ceiling} \cite{hansen1988,harrington2001}.
\textbf{Step 2.} Test Eq.~\eqref{eq:ch4.localization-resolution-tradeoff} \cite{jackson1999,stratton1941}.
\textbf{Step 3.} Reject super-resolution \cite{devaney2004}.
\textbf{Physical meaning.} Handset-scale apertures are dipole-class in \(\ell\) \cite{harrington2001}.

\paragraph{Problem (simultaneous product).}
\textbf{Problem.} Verify Eq.~\eqref{eq:ch4.loc-res-uncertainty} for measured \((\mathcal{L}_{2},\ell_{\mathrm{supp}})\).
\textbf{Step 1.} Measure \(\mathcal{L}_{2}[\Jimp]\) on \(\Omega_{s}\) \cite{jackson1999,stratton1941}.
\textbf{Step 2.} Read \(\ell_{\mathrm{supp}}(\varepsilon,k_{0}a)\) from SVD \cite{hansen1988,harrington2001}.
\textbf{Step 3.} Compare product to \(\mathcal{K}_{\mathrm{unc}}\) \cite{devaney2004}.
\textbf{Physical meaning.} Product bound is export-side uncertainty \cite{hansen1988}.

\paragraph{Philosophical closure (simultaneous claims).}
Localization and angular resolution are not independent knobs on compact sources \cite{hansen1988,stratton1941,harrington2001,jackson1999}. Every tight-focusing
with high-azimuthal-order claim is simultaneous on \(\mathcal{L}_{2}\) and \(\mathcal{R}_{\mathrm{ang}}\); realizability requires
Eqs.~\eqref{eq:ch4.localization-resolution-tradeoff}--\eqref{eq:ch4.loc-res-uncertainty} \cite{devaney2004,hansen1988,harrington2001}.

\begin{corollary}[No free super-resolution on \(B_{a}\)]
\label{cor:ch4.no-free-superres}
\(\mathcal{L}_{2}\ll a^{2}\) and \(\mathcal{R}_{\mathrm{ang}}\gg k_{0}a\) cannot hold for the same \(\mathbf{c}\in\mathcal{S}_{\mathrm{real}}(\mathcal{C})\) on \(B_{a}\)
\cite{hansen1988,stratton1941,devaney2004,harrington2001,jackson1999}.
\end{corollary}

\paragraph{Validity.}
Eqs.~\eqref{eq:ch4.localization-resolution-tradeoff}--\eqref{eq:ch4.loc-res-uncertainty} assume open-region compact sources \cite{stratton1941}. Cavities
quantize resolution via Subsection~\ref{sec:ch372} \cite{harrington2001}. Loss broadens \(\mathcal{L}_{2}\) without raising \(\ell_{\mathrm{supp}}\)
\cite{devaney2004,hansen1988}.


Section~\ref{sec:ch47} has separated synthesis rank from NR and unreachable nullities, reactive accessibility, inversion limits, and
localization--resolution tradeoffs \cite{stratton1941,harrington2001,devaney2004,hansen1988}. Export directions are necessary; observability and
inversion close on intersections with null-space and probe-rank constraints \cite{stratton1941}. Section~\ref{sec:ch48} develops asymptotic scaling
and spectral concentration laws on the same geometry \cite{jackson1999,hansen1988,harrington2001}.


\section{Scaling Laws and Spectral Concentration}
\label{sec:ch48}

Effective dimension \(d_{\mathrm{eff}}(\varepsilon)\) counts export directions; scaling exponents and concentration inequalities describe how that
count grows with aperture, frequency, and spectral bandwidth \cite{stratton1941,jackson1999,hansen1988,harrington2001}. Section~\ref{sec:ch48}
supplies asymptotic laws and concentration bounds that make truncation inevitable on finite hardware \cite{devaney2004}. Mode-count estimates bound
diversity; concentration ratios certify how sharply geometry compresses export; Landau--Pollak--Slepian \(n\)-widths cap simultaneous spatial and
spectral degrees of freedom \cite{hansen1988,stratton1941}. The philosophical arc: finite support forces finite information export at any tolerance
\cite{harrington2001,jackson1999}.


\subsection{Mode count versus aperture and frequency}
\label{sec:ch481}
% TOC tag: [HOME]

Mode-count laws connect aperture area, wavelength, and integer budgets in array and MIMO design \cite{collin2001,jackson1999,hansen1988}.
Subsection~\ref{sec:ch481} restates those laws as upper bounds on \(d_{\mathrm{eff}}(\varepsilon)\), importing volume and area truncation from
Section~\ref{sec:ch46} \cite{stratton1941,harrington2001,devaney2004}. Elements sample the aperture; modes bound export cardinality; operator rank
\(d_{\mathrm{eff}}\) counts directions above \(\varepsilon\sigma_{1}\) \cite{collin2001,hansen1988,harrington2001}. Confusing the three produces
hardware budgets that exceed radiative information capacity \cite{devaney2004,stratton1941}.

\paragraph{Assumptions.}
Presuppose Eqs.~\eqref{eq:ch4.truncation-volume-scaling}, \eqref{eq:ch4.truncation-area-scaling}, Definition~\ref{def:ch4.effective-em-dimension}
\cite{stratton1941,hansen1988,harrington2001}. Aligned flux gauge on \(\mathcal{C}\); tolerance \(\varepsilon\in(0,1)\) \cite{devaney2004,jackson1999}.
Phase-only retuning does not raise mode ceilings on fixed \((\Omega_{s},k_{0})\) \cite{devaney2004,harrington2001,stratton1941}.

\paragraph{Mode-count bounds.}
Volume and area laws ceiling listing cardinality:
\begin{equation}
\label{eq:ch4.mode-count-volume}
N_{\mathrm{mode}}^{\mathrm{(vol)}}(\varepsilon)\ \le\ C_{V}(\varepsilon)\,(k_{0}a)^{3},
\end{equation}
\begin{equation}
\label{eq:ch4.mode-count-area}
N_{\mathrm{mode}}^{\mathrm{(ap)}}(\varepsilon)\ \le\ C_{A}(\varepsilon)\,k_{0}^{2}A_{a},
\end{equation}
with \(N_{\mathrm{mode}}\ge d_{\mathrm{eff}}(\varepsilon)\) as analytic ceiling in aligned gauge \cite{collin2001,jackson1999,hansen1988,harrington2001}.
Mathematically, Eqs.~\eqref{eq:ch4.mode-count-volume}--\eqref{eq:ch4.mode-count-area} are Weyl-type counting bounds. Physically, they limit how many
independent propagating channels the geometry can support at tolerance \(\varepsilon\) \cite{stratton1941,devaney2004}.

\begin{theorem}[Mode-count monotonicity]
\label{thm:ch4.mode-count-monotone}
At fixed \(\varepsilon\), \(N_{\mathrm{mode}}^{\mathrm{(vol)}}\) and \(N_{\mathrm{mode}}^{\mathrm{(ap)}}\) are non-decreasing in \(k_{0}\) and aperture size
until dispersion alters \(\mathcal{J}_{\mathrm{real}}\) \cite{stratton1941,jackson1999,hansen1988}. Phase-only tuning does not raise \(N_{\mathrm{mode}}\)
\cite{devaney2004,harrington2001}.
\end{theorem}

\begin{theorem}[Mode ceiling dominates rank]
\label{thm:ch4.mode-ceiling-rank}
\(d_{\mathrm{eff}}(\varepsilon)\le N_{\mathrm{mode}}^{\mathrm{(vol)}}(\varepsilon)\le C_{V}(\varepsilon)(k_{0}a)^{3}\) and
\(d_{\mathrm{eff}}(\varepsilon)\le N_{\mathrm{mode}}^{\mathrm{(ap)}}(\varepsilon)\le C_{A}(\varepsilon)k_{0}^{2}A_{a}\) on aligned gauge
\cite{collin2001,jackson1999,hansen1988,harrington2001,stratton1941}.
\end{theorem}
\begin{proof}
Definition~\ref{def:ch4.effective-em-dimension} and Eqs.~\eqref{eq:ch4.truncation-volume-scaling}, \eqref{eq:ch4.truncation-area-scaling}
\cite{stratton1941,hansen1988,harrington2001}. Export rank cannot exceed mode-count ceiling \cite{devaney2004}.
\end{proof}

Theorem~\ref{thm:ch4.mode-ceiling-rank} is the array honesty law: vendor stream counts must be compared to \(N_{\mathrm{mode}}\) and measured
\(d_{\mathrm{eff}}\), not to element count alone \cite{collin2001,harrington2001,hansen1988}. Volume and area ceilings are not interchangeable:
compact three-dimensional sources are bounded by Eq.~\eqref{eq:ch4.mode-count-volume}, while planar apertures obey
Eq.~\eqref{eq:ch4.mode-count-area} \cite{jackson1999,stratton1941,hansen1988}. Synthesis audits must declare which geometry class governs the
declared \(\Omega_{s}\) \cite{harrington2001,devaney2004}.

\paragraph{Volume--area dominance.}
For thin apertures with thickness \(h\ll a\),
\begin{equation}
\label{eq:ch4.vol-area-dominance}
N_{\mathrm{mode}}^{\mathrm{(vol)}}\approx\frac{4\pi}{3}\frac{h}{a}\,N_{\mathrm{mode}}^{\mathrm{(ap)}},
\end{equation}
so area law typically sets the operative ceiling on panel arrays while volume law governs compact resonators \cite{collin2001,jackson1999,hansen1988,harrington2001}.

\paragraph{Rank--mode gap certificate.}
\begin{equation}
\label{eq:ch4.mode-rank-gap}
\Delta_{\mathrm{mr}}=N_{\mathrm{mode}}(\varepsilon)-d_{\mathrm{eff}}(\varepsilon)\ \ge\ 0
\end{equation}
measures listed modes below synthesis floor \cite{devaney2004,hansen1988}. Large \(\Delta_{\mathrm{mr}}\) with small \(d_{\mathrm{eff}}\) signals over-listing
\cite{stratton1941,harrington2001}.

\paragraph{Design margin.}
Engineering margin applies to mode budget, not to rank:
\begin{equation}
\label{eq:ch4.mode-count-margin}
N_{\mathrm{design}}=\lfloor\eta_{\mathrm{margin}}\,N_{\mathrm{mode}}(\varepsilon)\rfloor,\qquad 0<\eta_{\mathrm{margin}}<1.
\end{equation}
\begin{proposition}[Margin monotonicity]
\label{prop:ch4.mode-margin-monotone}
\(N_{\mathrm{design}}\) is non-decreasing in \(k_{0}\) and aperture at fixed \(\eta_{\mathrm{margin}}\) \cite{stratton1941,jackson1999,hansen1988,harrington2001}.
\end{proposition}

\paragraph{Element oversampling audit.}
\begin{equation}
\label{eq:ch4.element-mode-inequality}
N_{\mathrm{el}}\ge N_{\mathrm{mode}}(\varepsilon)\ \not\Rightarrow\ d_{\mathrm{eff}}(\varepsilon)=N_{\mathrm{el}},
\end{equation}
because \(d_{\mathrm{eff}}\) is operator rank, not element count \cite{collin2001,harrington2001,devaney2004}. Oversampled MoM bases inflate \(N_{\mathrm{el}}\)
without raising \(\sigma_{d_{\mathrm{eff}}+1}/\sigma_{1}\) above \(\varepsilon\) \cite{hansen1988,stratton1941}.

\paragraph{Frequency scaling.}
At fixed aperture area \(A_{a}\),
\begin{equation}
\label{eq:ch4.mode-freq-scaling}
\frac{N_{\mathrm{mode}}(\omega_{2})}{N_{\mathrm{mode}}(\omega_{1})}=\left(\frac{\omega_{2}}{\omega_{1}}\right)^{2}.
\end{equation}
Doubling frequency quadruples area-law mode budget at fixed footprint \cite{jackson1999,hansen1988,collin2001,harrington2001}.

\begin{lemma}[Stream count bound]
\label{lem:ch4.stream-bound}
Deployable MIMO streams satisfy \(N_{\mathrm{stream}}\le\min(N_{\mathrm{design}},d_{\mathrm{eff}}(\varepsilon))\) on declared \((\Omega_{s},k_{0},\varepsilon)\)
\cite{harrington2001,devaney2004,collin2001,hansen1988,stratton1941}.
\end{lemma}
\begin{proof}
\(N_{\mathrm{design}}\) is an engineering truncation of \(N_{\mathrm{mode}}\); deployable streams cannot exceed measured \(d_{\mathrm{eff}}(\varepsilon)\)
on \(\boldsymbol{\Gamma}_\omega\) \cite{harrington2001,devaney2004,collin2001}. Margin lowers listing; rank sets export \cite{stratton1941,hansen1988}.
\end{proof}

Lemma~\ref{lem:ch4.stream-bound} is the MIMO realizability certificate: vendor stream counts above \(\min(N_{\mathrm{design}},d_{\mathrm{eff}})\) are
listing claims, not physics guarantees \cite{collin2001,harrington2001,hansen1988,devaney2004}.

\begin{theorem}[Rank monotonicity in electrical size]
\label{thm:ch4.rank-ka-monotone}
At fixed \(\varepsilon\) and aperture class, \(d_{\mathrm{eff}}(\varepsilon)\) is non-decreasing in \(k_{0}a\) until material dispersion alters
\(\mathcal{J}_{\mathrm{real}}\) \cite{hansen1988,harrington2001,stratton1941,devaney2004,jackson1999}.
\end{theorem}
\begin{proof}
Singular values \(\sigma_{n}\) of \(\boldsymbol{\Gamma}_\omega\) are non-decreasing in \(k_{0}a\) on fixed aperture class until dispersion alters
\(\mathcal{J}_{\mathrm{real}}\) \cite{hansen1988,harrington2001}. \(d_{\mathrm{eff}}(\varepsilon)\) counts modes above \(\varepsilon\sigma_{1}\)
\cite{stratton1941,devaney2004,jackson1999}.
\end{proof}

Theorem~\ref{thm:ch4.rank-ka-monotone} pairs with Theorem~\ref{thm:ch4.mode-count-monotone}: electrical enlargement raises listing ceilings and measured
rank together, but \(d_{\mathrm{eff}}\le N_{\mathrm{mode}}\) always \cite{collin2001,harrington2001,hansen1988,devaney2004}.

\begin{figure}[t]
\centering
\ChOneFigBlock{%
\begin{tikzpicture}[ch1fig, node distance=7mm and 9mm]
  \node[ch1box] (ka) {$k_{0}a$};
  \node[ch1box, right=11mm of ka] (nm) {$N_{\mathrm{mode}}$};
  \node[ch1box, right=11mm of nm] (deff) {$d_{\mathrm{eff}}$};
  \draw[ch1flow] (ka) -- (nm) -- (deff);
\end{tikzpicture}%
}
\caption{Subsection~\ref{sec:ch481}: electrical size sets mode budget; \(d_{\mathrm{eff}}\) is measured rank.}
\label{fig:ch4.mode-count-bound}
\end{figure}

\begin{figure}[t]
\centering
\ChOneFigBlock{%
\begin{tikzpicture}[ch1fig, node distance=7mm and 9mm]
  \node[ch1box] (nel) {$N_{\mathrm{el}}$};
  \node[ch1box, right=11mm of nel] (nm) {$N_{\mathrm{mode}}$};
  \node[ch1box, right=11mm of nm] (deff) {$d_{\mathrm{eff}}$};
  \draw[ch1flow] (nel) -- (nm);
  \draw[ch1flow] (nm) -- (deff);
\end{tikzpicture}%
}
\caption{Subsection~\ref{sec:ch481}: elements \(\ge\) modes \(\ge\) rank; only \(d_{\mathrm{eff}}\) is physics-limited.}
\label{fig:ch4.element-mode-rank}
\end{figure}

Figure~\ref{fig:ch4.element-mode-rank} is the MIMO audit chain: element count is hardware; \(d_{\mathrm{eff}}\) is the governing physics certificate
\cite{collin2001,harrington2001,hansen1988,devaney2004}.

\paragraph{Worked illustration (5G panel).}
\(k_{0}W\approx3\): \(N_{\mathrm{mode}}^{\mathrm{(ap)}}\sim\mathcal{O}(10)\)--\(\mathcal{O}(30)\) \cite{collin2001,harrington2001,hansen1988,jackson1999}.

\paragraph{Worked illustration (panel audit).}
\(8\times8\) patches (\(N_{\mathrm{el}}=64\)), \(k_{0}W=2.5\): \(N_{\mathrm{mode}}^{\mathrm{(ap)}}\approx20\), \(d_{\mathrm{eff}}(10^{-2})=11\)
\cite{collin2001,harrington2001,hansen1988,stratton1941}.

\paragraph{Worked illustration (margin versus rank).}
\(N_{\mathrm{mode}}=30\), \(\eta_{\mathrm{margin}}=0.8\Rightarrow N_{\mathrm{design}}=24\), but \(d_{\mathrm{eff}}=9\): margin cannot exceed rank
\cite{harrington2001,stratton1941,devaney2004}.

\paragraph{Worked illustration (frequency doubling).}
\(k_{0}^{2}A_{a}=12\) at 3.5\,GHz \(\Rightarrow N_{\mathrm{mode}}^{\mathrm{(ap)}}\approx18\); at 7\,GHz \(\Rightarrow\) quadruple area factor in
Eq.~\eqref{eq:ch4.mode-freq-scaling} \cite{jackson1999,hansen1988,collin2001}.

\paragraph{Worked numerics (gap audit).}
\(N_{\mathrm{mode}}=22\), \(d_{\mathrm{eff}}=9\), \(\Delta_{\mathrm{mr}}=13\): \(59\%\) of mode budget is below synthesis floor \cite{hansen1988,harrington2001,devaney2004}.

\paragraph{Worked illustration (volume versus area).}
Compact \(k_{0}a=1.2\) resonator: \(N_{\mathrm{mode}}^{\mathrm{(vol)}}\approx14\) per Eq.~\eqref{eq:ch4.mode-count-volume}; coplanar patch of same
footprint yields \(N_{\mathrm{mode}}^{\mathrm{(ap)}}\approx9\) per Eq.~\eqref{eq:ch4.vol-area-dominance} \cite{collin2001,jackson1999,hansen1988,harrington2001}.

\paragraph{Worked illustration (stream bound).}
\(N_{\mathrm{design}}=18\) from \(\eta_{\mathrm{margin}}=0.75\) on \(N_{\mathrm{mode}}=24\), but \(d_{\mathrm{eff}}(10^{-2})=10\): Lemma~\ref{lem:ch4.stream-bound}
caps deployable streams at \(10\) \cite{harrington2001,devaney2004,collin2001,stratton1941}.

\paragraph{Worked numerics (electrical-size sweep).}
\(k_{0}W=1.5,2.0,2.5\Rightarrow N_{\mathrm{mode}}^{\mathrm{(ap)}}\approx7,12,19\) while \(d_{\mathrm{eff}}(10^{-2})=4,7,11\): rank tracks mode budget per
Theorem~\ref{thm:ch4.rank-ka-monotone} but \(\Delta_{\mathrm{mr}}\) remains positive at each point \cite{hansen1988,harrington2001,devaney2004,collin2001}.

\paragraph{Problem (elements versus modes).}
\textbf{Problem.} \(64\) patches claim \(64\) beams. Audit.
\textbf{Step 1.} Compute \(k_{0}^{2}A_{a}\) \cite{jackson1999,stratton1941}.
\textbf{Step 2.} Bound \(N_{\mathrm{mode}}^{\mathrm{(ap)}}\) via Eq.~\eqref{eq:ch4.mode-count-area} \cite{hansen1988,harrington2001}.
\textbf{Step 3.} Report \(d_{\mathrm{eff}}\) from SVD \cite{devaney2004}.
\textbf{Physical meaning.} Elements sample; modes bound \cite{collin2001}.

\paragraph{Problem (stream count).}
\textbf{Problem.} MIMO streams at \(\eta_{\mathrm{margin}}=0.75\), \(N_{\mathrm{mode}}=24\)?
\textbf{Step 1.} Eq.~\eqref{eq:ch4.mode-count-margin} \cite{hansen1988,harrington2001}.
\textbf{Step 2.} Apply Lemma~\ref{lem:ch4.stream-bound} \cite{devaney2004,collin2001}.
\textbf{Step 3.} Report \(\min(N_{\mathrm{design}},d_{\mathrm{eff}})\) \cite{stratton1941}.
\textbf{Physical meaning.} Margin is engineering; rank is physics \cite{harrington2001}.

\paragraph{Problem (beam claim audit).}
\textbf{Problem.} Vendor claims \(48\) MU-MIMO streams on \(k_{0}^{2}A_{a}=15\).
\textbf{Step 1.} Bound \(N_{\mathrm{mode}}^{\mathrm{(ap)}}\) \cite{jackson1999,hansen1988}.
\textbf{Step 2.} Measure \(d_{\mathrm{eff}}(\varepsilon)\) \cite{harrington2001,devaney2004}.
\textbf{Step 3.} Report Lemma~\ref{lem:ch4.stream-bound} result \cite{stratton1941,collin2001}.
\textbf{Physical meaning.} Stream count is rank-limited \cite{hansen1988}.

\paragraph{Problem (frequency doubling).}
\textbf{Problem.} Double \(\omega\) at fixed aperture. Mode budget change?
\textbf{Step 1.} Apply Eq.~\eqref{eq:ch4.mode-freq-scaling} \cite{jackson1999,hansen1988}.
\textbf{Step 2.} Recompute \(N_{\mathrm{mode}}\) and \(d_{\mathrm{eff}}\) \cite{harrington2001,devaney2004}.
\textbf{Step 3.} Update \(N_{\mathrm{design}}\) \cite{stratton1941,collin2001}.
\textbf{Physical meaning.} Electrical size sets mode budget \cite{harrington2001}.

\paragraph{Problem (gap certificate).}
\textbf{Problem.} \(N_{\mathrm{mode}}=30\), \(d_{\mathrm{eff}}=8\). Interpret \(\Delta_{\mathrm{mr}}\).
\textbf{Step 1.} Compute Eq.~\eqref{eq:ch4.mode-rank-gap} \cite{hansen1988,devaney2004}.
\textbf{Step 2.} Conclude over-listing if \(\Delta_{\mathrm{mr}}\gg d_{\mathrm{eff}}\) \cite{harrington2001,stratton1941}.
\textbf{Step 3.} Truncate chart to \(d_{\mathrm{eff}}\) for synthesis \cite{devaney2004}.
\textbf{Physical meaning.} Gap measures listing excess \cite{hansen1988}.

\paragraph{Problem (phase-only MIMO).}
\textbf{Problem.} Phase-only metasurface claims doubled streams. Valid?
\textbf{Step 1.} Apply Corollary~\ref{cor:ch4.phase-no-mode-gain} \cite{devaney2004,harrington2001}.
\textbf{Step 2.} Verify \(d_{\mathrm{eff}}\) unchanged \cite{hansen1988,stratton1941}.
\textbf{Step 3.} Reject stream-doubling claim \cite{collin2001}.
\textbf{Physical meaning.} Phase retunes within fixed rank \cite{harrington2001}.

\paragraph{Problem (rank--mode audit).}
\textbf{Problem.} Base station \(k_{0}W=3.2\), \(N_{\mathrm{el}}=128\). Certify streams.
\textbf{Step 1.} Bound \(N_{\mathrm{mode}}^{\mathrm{(ap)}}\) via Eq.~\eqref{eq:ch4.mode-count-area} \cite{jackson1999,hansen1988}.
\textbf{Step 2.} Measure \(d_{\mathrm{eff}}(\varepsilon)\) and compute \(\Delta_{\mathrm{mr}}\) \cite{harrington2001,devaney2004}.
\textbf{Step 3.} Apply Lemma~\ref{lem:ch4.stream-bound} \cite{collin2001,stratton1941}.
\textbf{Physical meaning.} Element oversampling does not raise rank \cite{hansen1988}.

\paragraph{Philosophical closure (enumeration versus export).}
Mode-count laws answer how many indices geometry can support; \(d_{\mathrm{eff}}\) answers how many carry export power above the noise floor
\cite{hansen1988,stratton1941,harrington2001,devaney2004}. Every array synthesis memo must pair \(N_{\mathrm{mode}}\), \(N_{\mathrm{design}}\), and measured
\(d_{\mathrm{eff}}\) with \(\Delta_{\mathrm{mr}}\) \cite{collin2001,stratton1941}.

\begin{corollary}[Phase-only cannot raise mode ceiling]
\label{cor:ch4.phase-no-mode-gain}
At fixed \((\Omega_{s},k_{0},\varepsilon)\), phase-only retuning leaves \(N_{\mathrm{mode}}(\varepsilon)\) and \(d_{\mathrm{eff}}(\varepsilon)\) unchanged
\cite{devaney2004,harrington2001,stratton1941,hansen1988}.
\end{corollary}

\paragraph{Validity.}
Theorem~\ref{thm:ch4.mode-count-monotone} is asymptotic for \(k_{0}D_{s}\gtrsim1\) \cite{jackson1999,hansen1988}. Dispersion and surface waves modify
prefactors \(C_{V},C_{A}\) without changing exponents on fixed \(\mathcal{J}_{\mathrm{real}}\) class \cite{stratton1941,harrington2001,devaney2004,collin2001}.

\subsection{High- and low-frequency asymptotics}
\label{sec:ch482}
% TOC tag: [HOME]

Radiative rank \(d_{\mathrm{eff}}(\varepsilon)\) crosses distinct asymptotic regimes as electrical size \(k_{0}D_{s}\) varies on compact
\(\Omega_{s}\) \cite{stratton1941,jackson1999,hansen1988}. Subsection~\ref{sec:ch482} records dipole-class and volume-growth limits, importing
Eq.~\eqref{eq:ch4.coupling-crossover} from Section~\ref{sec:ch45} by reference \cite{harrington2001,devaney2004}. Regime selection is mandatory
before rank is extrapolated across frequency bands: the same outline can be dipole-limited at sub-6\,GHz and volume-limited at mmWave
\cite{jackson1999,harrington2001,collin2001}.

\paragraph{Assumptions.}
Presuppose Eq.~\eqref{eq:ch4.coupling-crossover}, compact \(\Omega_{s}\) with \(\operatorname{diam}=D_{s}\), and linear isotropic exterior
\cite{stratton1941,hansen1988,jackson1999}. Tolerance \(\varepsilon\in(0,1)\) fixed; flux gauge on \(\mathcal{C}\) \cite{harrington2001,devaney2004}.
Phase-only retuning does not alter asymptotic exponents on fixed \((\Omega_{s},k_{0})\) \cite{devaney2004,hansen1988}.

\paragraph{Rayleigh limit.}
As \(k_{0}D_{s}\to0\), multipole moments with \(\ell\ge1\) are suppressed by \((k_{0}D_{s})^{2\ell+1}\) \cite{jackson1999,stratton1941}. Export rank
saturates at dipole-class cardinality:
\begin{equation}
\label{eq:ch4.rayleigh-rank-limit}
d_{\mathrm{eff}}(\varepsilon)\ \longrightarrow\ d_{0}\le3
\end{equation}
in the electric-dipole gauge \cite{harrington2001,hansen1988}. Equation~\eqref{eq:ch4.rayleigh-rank-limit} is the small-aperture synthesis law: no
passive \(\Jimp\in\mathcal{J}_{\mathrm{real}}\) can excite more than three independent far-zone directions when the body is electrically tiny
\cite{stratton1941,jackson1999}.

\begin{theorem}[Dipole cap in Rayleigh band]
\label{thm:ch4.rayleigh-dipole-cap}
For \(k_{0}D_{s}<0.5\) on compact \(\Omega_{s}\), \(d_{\mathrm{eff}}(\varepsilon)\le3\) for any \(\varepsilon\in(0,1)\) in the electric-dipole
gauge \cite{jackson1999,stratton1941,harrington2001,hansen1988}.
\end{theorem}
\begin{proof}
Leading multipole weights scale as \((k_{0}D_{s})^{2\ell+1}\) \cite{jackson1999}. For \(k_{0}D_{s}<0.5\), only \(\ell=0,1\) sectors can exceed
\(\varepsilon\sigma_{1}\) in Definition~\ref{def:ch4.effective-em-dimension} \cite{stratton1941,harrington2001,hansen1988}.
\end{proof}

\paragraph{Multipole rank ledger (Rayleigh band).}
Expanding the far field in vector spherical harmonics,
\begin{equation}
\label{eq:ch4.multipole-rank-ledger}
d_{\mathrm{eff}}(\varepsilon)\le\sum_{\ell=0}^{\ell_{\max}}\sum_{\sigma}\mathbf{1}\!\left[\sigma_{(\ell,\cdot,\sigma)}\ge\varepsilon\sigma_{1}\right],
\end{equation}
with \(\ell\in\{0,1\}\) dominant when \(k_{0}D_{s}<0.5\) \cite{jackson1999,stratton1941,hansen1988}. Equation~\eqref{eq:ch4.multipole-rank-ledger}
makes the dipole cap constructive: higher multipoles are below threshold before any optimizer runs \cite{harrington2001,devaney2004}.

\paragraph{Large-body limit.}
As \(k_{0}D_{s}\to\infty\),
\begin{equation}
\label{eq:ch4.large-body-rank-growth}
d_{\mathrm{eff}}(\varepsilon)\ \sim\ C_{\infty}(\varepsilon)\,(k_{0}D_{s})^{3},
\end{equation}
matching Weyl volume law Eq.~\eqref{eq:ch4.truncation-volume-scaling} up to prefactor \(C_{\infty}(\varepsilon)\) \cite{stratton1941,jackson1999,hansen1988}.
Physically, large bodies support many independent export directions; truncation at tolerance \(\varepsilon\) is mandatory because full listing is
unattainable \cite{harrington2001,devaney2004}.

\paragraph{Volume-rank ledger (large-body band).}
When \(k_{0}D_{s}>2\), rank increments track volume mode density on \(B_{D_{s}}\) \cite{stratton1941,jackson1999,hansen1988}. Rank is not tunable by
feed phasing on fixed \(\Omega_{s}\) \cite{devaney2004,harrington2001}. Mandatory truncation follows from
\(d_{\mathrm{eff}}\sim(k_{0}D_{s})^{3}\) growth \cite{hansen1988}.

\begin{corollary}[Asymptotic synthesis regimes]
\label{cor:ch4.asymptotic-regimes}
Small hardware is dipole-limited; large hardware is volume-growth limited with mandatory truncation \cite{stratton1941,jackson1999,hansen1988,harrington2001,devaney2004}.
\end{corollary}

\paragraph{Crossover and regime table.}
Electrical crossover occurs when the first multipole beyond dipole class exceeds \(\varepsilon\sigma_{1}\):
\begin{equation}
\label{eq:ch4.crossover-frequency}
k_{0}^{(\mathrm{cross})}D_{s}\sim\mathcal{O}(1),
\end{equation}
\begin{equation}
\label{eq:ch4.regime-table}
d_{\mathrm{eff}}(\varepsilon)\in\begin{cases}\{1,2,3\}, & k_{0}D_{s}<0.5,\\ \mathcal{O}((k_{0}D_{s})^{3}), & k_{0}D_{s}>2,\end{cases}
\end{equation}
\begin{equation}
\label{eq:ch4.crossover-knee}
\kappa_{\mathrm{cross}}=\inf\left\{k_{0}D_{s}:\ d_{\mathrm{eff}}(\varepsilon)>3\right\}.
\end{equation}
Transition band \(0.5<k_{0}D_{s}<2\) requires direct SVD audit; neither asymptote alone is reliable \cite{hansen1988,harrington2001,jackson1999}.

\begin{lemma}[Phase invariance of asymptotic branch]
\label{lem:ch4.phase-asymptote}
At fixed \((\Omega_{s},k_{0},\varepsilon)\), phase-only retuning leaves the active asymptotic branch in Eq.~\eqref{eq:ch4.regime-table} unchanged
\cite{devaney2004,harrington2001,stratton1941,hansen1988}.
\end{lemma}

\begin{figure}[t]
\centering
\ChOneFigBlock{%
\begin{tikzpicture}[ch1fig, node distance=7mm and 9mm]
  \node[ch1box] (ray) {Eq.~\eqref{eq:ch4.rayleigh-rank-limit}};
  \node[ch1box, right=11mm of ray] (cross) {$\kappa_{\mathrm{cross}}$};
  \node[ch1box, right=11mm of cross] (large) {Eq.~\eqref{eq:ch4.large-body-rank-growth}};
  \draw[ch1flow] (ray) -- (cross) -- (large);
\end{tikzpicture}%
}
\caption{Subsection~\ref{sec:ch482}: electrical crossover \(\kappa_{\mathrm{cross}}\) selects dipole versus volume rank law.}
\label{fig:ch4.asymptotic-crossover}
\end{figure}

Figure~\ref{fig:ch4.asymptotic-crossover} is the regime map every multi-band synthesis memo must attach: extrapolation across bands without branch
identification is a common source of unrealizable targets \cite{jackson1999,harrington2001,stratton1941}.

\begin{figure}[t]
\centering
\ChOneFigBlock{%
\begin{tikzpicture}[ch1fig, node distance=7mm and 9mm]
  \node[ch1box] (kds) {$k_{0}D_{s}$};
  \node[ch1box, right=11mm of kds] (deff) {$d_{\mathrm{eff}}$};
  \node[ch1box, right=11mm of deff] (law) {active law};
  \draw[ch1flow] (kds) -- (deff) -- (law);
\end{tikzpicture}%
}
\caption{Subsection~\ref{sec:ch482}: rank versus electrical size; knee at \(\kappa_{\mathrm{cross}}\) marks branch change.}
\label{fig:ch4.rank-vs-kds}
\end{figure}

\paragraph{Worked illustration (UHF vs mmWave).}
\(k_{0}D_{s}=0.3\): \(d_{\mathrm{eff}}\le3\) by Theorem~\ref{thm:ch4.rayleigh-dipole-cap}; \(k_{0}D_{s}=4\): \(d_{\mathrm{eff}}\sim\mathcal{O}(10^{1})\)
under Eq.~\eqref{eq:ch4.large-body-rank-growth} \cite{jackson1999,harrington2001,collin2001}.

\paragraph{Worked illustration (handset bands).}
Smartphone \(D_{s}\approx0.15\,\mathrm{m}\): at 1\,GHz, \(k_{0}D_{s}\approx0.3\Rightarrow d_{\mathrm{eff}}\le3\); at 28\,GHz,
\(k_{0}D_{s}\approx8\Rightarrow d_{\mathrm{eff}}\sim\mathcal{O}(10^{2})\) \cite{jackson1999,harrington2001,collin2001,hansen1988}.

\paragraph{Worked illustration (mis-extrapolation).}
Rank at 900\,MHz (\(d_{\mathrm{eff}}=2\)) extrapolated to 60\,GHz via dipole law yields \(d_{\mathrm{eff}}\approx2\); correct volume branch gives
\(d_{\mathrm{eff}}\gg10\) \cite{hansen1988,devaney2004,stratton1941}. Error factor exceeds \((k_{0}D_{s}/0.5)^{3}\) per
Corollary~\ref{cor:ch4.no-dipole-extrap} \cite{jackson1999,harrington2001}.

\paragraph{Worked illustration (transition band).}
\(k_{0}D_{s}=1.2\): neither Eq.~\eqref{eq:ch4.rayleigh-rank-limit} nor Eq.~\eqref{eq:ch4.large-body-rank-growth} is accurate; measured
\(d_{\mathrm{eff}}(10^{-2})=7\) requires direct SVD \cite{hansen1988,harrington2001,devaney2004}.

\paragraph{Worked numerics (crossover sweep).}
\(D_{s}=0.12\,\mathrm{m}\): \(d_{\mathrm{eff}}(10^{-2})=\{2,2,2,3,8,22\}\) at \(\{1,3,6,10,20,40\}\,\mathrm{GHz}\); \(\kappa_{\mathrm{cross}}\) near
\(10\,\mathrm{GHz}\) \cite{jackson1999,harrington2001,hansen1988,collin2001}.

\paragraph{Problem (frequency sweep).}
\textbf{Problem.} Extrapolate rank from 3\,GHz to 30\,GHz on fixed hardware.
\textbf{Step 1.} Identify regime per frequency via Eq.~\eqref{eq:ch4.regime-table} \cite{jackson1999,stratton1941}.
\textbf{Step 2.} Apply Theorem~\ref{thm:ch4.rayleigh-dipole-cap} or Eq.~\eqref{eq:ch4.large-body-rank-growth} \cite{hansen1988,harrington2001}.
\textbf{Step 3.} Recompute \(d_{\mathrm{eff}}(\varepsilon)\) by SVD in transition band \cite{devaney2004}.
\textbf{Physical meaning.} Regime sets scaling exponent \cite{stratton1941}.

\paragraph{Problem (crossover audit).}
\textbf{Problem.} When does volume law replace dipole law?
\textbf{Step 1.} Track \(d_{\mathrm{eff}}\) vs \(k_{0}D_{s}\) \cite{hansen1988}.
\textbf{Step 2.} Identify \(\kappa_{\mathrm{cross}}\) via Eq.~\eqref{eq:ch4.crossover-knee} \cite{jackson1999}.
\textbf{Step 3.} Document active branch in synthesis memo \cite{stratton1941,harrington2001}.
\textbf{Physical meaning.} Crossover is geometry--frequency fact \cite{hansen1988}.

\paragraph{Problem (band sweep).}
\textbf{Problem.} Fixed \(D_{s}=0.2\,\mathrm{m}\), sweep 1--40\,GHz. Document rank law.
\textbf{Step 1.} Compute \(k_{0}D_{s}\) per frequency \cite{jackson1999}.
\textbf{Step 2.} Apply appropriate branch of Eq.~\eqref{eq:ch4.regime-table} \cite{stratton1941,hansen1988}.
\textbf{Step 3.} Mark transition-band frequencies for SVD \cite{harrington2001,devaney2004}.
\textbf{Physical meaning.} Regime selects scaling exponent \cite{devaney2004}.

\paragraph{Problem (handset dual-band).}
\textbf{Problem.} Same outline at 1.9\,GHz and 28\,GHz. Rank laws?
\textbf{Step 1.} Compute \(k_{0}D_{s}\) at each band \cite{jackson1999}.
\textbf{Step 2.} Apply Theorem~\ref{thm:ch4.rayleigh-dipole-cap} at 1.9\,GHz \cite{stratton1941}.
\textbf{Step 3.} Apply Eq.~\eqref{eq:ch4.large-body-rank-growth} at 28\,GHz \cite{hansen1988,harrington2001}.
\textbf{Physical meaning.} Band change is regime change \cite{devaney2004,collin2001}.

\paragraph{Problem (phase retuning).}
\textbf{Problem.} Can phase-only metasurface raise \(d_{\mathrm{eff}}\) in Rayleigh band?
\textbf{Step 1.} Apply Lemma~\ref{lem:ch4.phase-asymptote} \cite{devaney2004,harrington2001}.
\textbf{Step 2.} Conclude branch unchanged at fixed \(k_{0}D_{s}\) \cite{stratton1941,hansen1988}.
\textbf{Step 3.} Raise \(k_{0}D_{s}\) by frequency or size to change rank \cite{jackson1999}.
\textbf{Physical meaning.} Asymptotic rank is not phase-tunable \cite{harrington2001}.

\paragraph{Philosophical closure (asymptotic honesty).}
Small-body and large-body limits are mutually exclusive regimes separated by \(\kappa_{\mathrm{cross}}\) on the same geometry
\cite{jackson1999,stratton1941,hansen1988,harrington2001}. Rank projections that blend dipole and volume laws optimistically produce synthesis targets
no passive support can meet \cite{devaney2004}. Every frequency-swept design must state which branch of Eq.~\eqref{eq:ch4.regime-table} was active at
each band \cite{stratton1941,hansen1988}.

\begin{corollary}[No mid-band dipole extrapolation]
\label{cor:ch4.no-dipole-extrap}
For \(k_{0}D_{s}>2\), replacing Eq.~\eqref{eq:ch4.large-body-rank-growth} with Eq.~\eqref{eq:ch4.rayleigh-rank-limit} underestimates
\(d_{\mathrm{eff}}(\varepsilon)\) by at least factor \((k_{0}D_{s}/0.5)^{3}\) \cite{jackson1999,hansen1988,stratton1941,harrington2001}.
\end{corollary}

\paragraph{Validity.}
Eqs.~\eqref{eq:ch4.rayleigh-rank-limit}--\eqref{eq:ch4.large-body-rank-growth} assume linear isotropic exterior and compact \(\Omega_{s}\)
\cite{stratton1941}. Dispersion and material resonances modify \(C_{\infty}(\varepsilon)\) without changing exponents on fixed \(\mathcal{J}_{\mathrm{real}}\)
class \cite{jackson1999,hansen1988,harrington2001,devaney2004}.

\subsection{Spectral concentration and localization}
\label{sec:ch483}
% TOC tag: [HOME]

Finite support concentrates spatial energy and compresses the spectral support of \(\boldsymbol{\Gamma}_\omega\); concentration inequalities make
truncation the optimal description of exportable information \cite{stratton1941,jackson1999,hansen1988}. Subsection~\ref{sec:ch483} states concentration
on the singular-value ledger of \(\boldsymbol{\Gamma}_\omega\), linking \(\Lambda_{\mathrm{conc}}(\varepsilon)\) to \(\ell_{\mathrm{supp}}\) from
Eq.~\eqref{eq:ch4.spectral-compression-def} and to coupling-tail mass in Subsection~\ref{sec:ch463} \cite{harrington2001,devaney2004}. Listing
cardinality and energy participation are distinct certificates \cite{stratton1941,hansen1988}.

\paragraph{Assumptions.}
Presuppose Eq.~\eqref{eq:ch4.spectral-compression-def}, Theorem~\ref{thm:ch4.spectral-compression-optimal}, and compact \(\boldsymbol{\Gamma}_\omega\)
on \(\mathcal{J}_{\mathrm{real}}\) \cite{stratton1941,harrington2001,hansen1988}. Tolerance \(\varepsilon\in(0,1)\) fixed; flux-normalized \(\mathcal{C}\)
\cite{devaney2004}. SVD \(\boldsymbol{\Gamma}_\omega=\sum_{n}\sigma_{n}\mathbf{u}_{n}\mathbf{v}_{n}^{\dagger}\) in aligned gauge \cite{harrington2001}.

\paragraph{Concentration ratio.}
Define the export participation fraction
\begin{equation}
\label{eq:ch4.spectral-concentration-ratio}
\Lambda_{\mathrm{conc}}(\varepsilon)=\frac{\sum_{n=1}^{d_{\mathrm{eff}}(\varepsilon)}\sigma_{n}^{2}}{\sum_{n\ge1}\sigma_{n}^{2}}\ \ge\ 1-\varepsilon^{2},
\end{equation}
\cite{harrington2001,hansen1988,stratton1941}. Mathematically, \(\Lambda_{\mathrm{conc}}\) is the fraction of \(\|\boldsymbol{\Gamma}_\omega\|_{\mathrm{HS}}^{2}\)
carried by the leading \(d_{\mathrm{eff}}\) singular values. Physically, it is the fraction of exportable radiative information retained after rank-\(d_{\mathrm{eff}}\)
truncation \cite{devaney2004,hansen1988}.

\paragraph{Tail-energy budget.}
The discarded export fraction is
\begin{equation}
\label{eq:ch4.tail-discard}
\mathcal{T}_{\mathrm{tail}}(\varepsilon)=1-\Lambda_{\mathrm{conc}}(\varepsilon)=\frac{\sum_{n>d_{\mathrm{eff}}}\sigma_{n}^{2}}{\sum_{n\ge1}\sigma_{n}^{2}}.
\end{equation}
Slow singular-value tails (\(\sigma_{n}\sim n^{-\alpha}\), \(\alpha\le2\)) keep \(\mathcal{T}_{\mathrm{tail}}\) large even when \(d_{\mathrm{eff}}\) is modest
\cite{devaney2004,hansen1988,harrington2001}. Flat tails are the concentration analogue of coupling-tail mass
\(\mathcal{M}_{\mathrm{cpl}}\) in Eq.~\eqref{eq:ch4.coupling-tail-mass} \cite{stratton1941}.

\begin{proposition}[Truncation as concentration optimum]
\label{prop:ch4.truncation-concentration}
Rank-\(d_{\mathrm{eff}}(\varepsilon)\) truncation maximizes \(\Lambda_{\mathrm{conc}}(\varepsilon)\) among all rank-\(N\) approximations of
\(\boldsymbol{\Gamma}_\omega\) \cite{harrington2001,devaney2004,stratton1941}.
\end{proposition}
\begin{proof}
Eckart--Young theorem: best rank-\(N\) approximation minimizes discarded tail energy \cite{harrington2001}. Definition~\ref{def:ch4.effective-em-dimension}
selects \(N=d_{\mathrm{eff}}(\varepsilon)\) at tolerance \(\varepsilon\) \cite{stratton1941,hansen1988}.
\end{proof}

\begin{theorem}[Eckart--Young concentration optimality]
\label{thm:ch4.eckart-young-conc}
Among all rank-\(N\) approximations of \(\boldsymbol{\Gamma}_\omega\), truncation at \(N=d_{\mathrm{eff}}(\varepsilon)\) minimizes
\(\mathcal{T}_{\mathrm{tail}}(\varepsilon)\) and maximizes \(\Lambda_{\mathrm{conc}}(\varepsilon)\) \cite{harrington2001,devaney2004,stratton1941,hansen1988}.
\end{theorem}

\paragraph{Export participation ratio.}
\begin{equation}
\label{eq:ch4.export-participation}
\mathrm{EPR}(\varepsilon)=\Lambda_{\mathrm{conc}}(\varepsilon),\qquad 1-\mathrm{EPR}(\varepsilon)=\mathcal{T}_{\mathrm{tail}}(\varepsilon).
\end{equation}
\begin{theorem}[Optimal truncation participation]
\label{thm:ch4.optimal-participation}
Rank-\(d_{\mathrm{eff}}(\varepsilon)\) truncation maximizes \(\mathrm{EPR}(\varepsilon)\) among rank-\(N\) maps on \(\boldsymbol{\Gamma}_\omega\)
\cite{harrington2001,devaney2004,stratton1941}.
\end{theorem}

Theorem~\ref{thm:ch4.optimal-participation} is the watts-honesty certificate: a synthesis memo that reports \(d_{\mathrm{eff}}\) without
\(\mathrm{EPR}\) omits how much radiative information was sacrificed \cite{hansen1988,harrington2001,stratton1941}.

\paragraph{Singular-value knee.}
Tolerance \(\varepsilon\) locates the knee
\begin{equation}
\label{eq:ch4.sv-knee}
\sigma_{d_{\mathrm{eff}}}/\sigma_{1}\ge\varepsilon,\qquad \sigma_{d_{\mathrm{eff}}+1}/\sigma_{1}<\varepsilon.
\end{equation}
Knee location sets both \(d_{\mathrm{eff}}\) and \(\mathrm{EPR}\) simultaneously \cite{hansen1988,harrington2001,devaney2004}.

\paragraph{Per-mode participation.}
Define \(\pi_{n}=\sigma_{n}^{2}/\sum_{k}\sigma_{k}^{2}\). Flat \(\{\pi_{n}\}\) implies large \(\mathcal{T}_{\mathrm{tail}}\) at small rank increment
\cite{devaney2004,hansen1988}. Steep knees permit high \(\mathrm{EPR}\) at modest \(d_{\mathrm{eff}}\) \cite{harrington2001,stratton1941}.

\paragraph{Decay-rate bound.}
If \(\sigma_{n}\le C n^{-\alpha}\) for \(n\ge n_{0}\), then
\begin{equation}
\label{eq:ch4.sv-decay-bound}
\mathcal{T}_{\mathrm{tail}}(\varepsilon)\le C'\,d_{\mathrm{eff}}^{1-2\alpha}\quad(\alpha>1/2),
\end{equation}
linking tail flatness to slow rank convergence \cite{hansen1988,harrington2001,devaney2004}. Algebraic tails require higher \(d_{\mathrm{eff}}\) for
\(\mathrm{EPR}\ge0.99\) than exponential tails at the same \(\varepsilon\) \cite{stratton1941}.

\begin{lemma}[Concentration monotonicity in \(\varepsilon\)]
\label{lem:ch4.conc-monotone}
\(\Lambda_{\mathrm{conc}}(\varepsilon)\) is non-decreasing as \(\varepsilon\) decreases (rank \(d_{\mathrm{eff}}\) increases)
\cite{harrington2001,hansen1988,stratton1941,devaney2004}.
\end{lemma}

\begin{figure}[t]
\centering
\ChOneFigBlock{%
\begin{tikzpicture}[ch1fig, node distance=7mm and 9mm]
  \node[ch1box] (svd) {$\{\sigma_{n}\}$};
  \node[ch1box, right=11mm of svd] (epr) {$\mathrm{EPR}$};
  \node[ch1box, right=11mm of epr] (deff) {$d_{\mathrm{eff}}$};
  \draw[ch1flow] (svd) -- (epr) -- (deff);
\end{tikzpicture}%
}
\caption{Subsection~\ref{sec:ch483}: singular-value ledger drives EPR and \(d_{\mathrm{eff}}\) at knee Eq.~\eqref{eq:ch4.sv-knee}.}
\label{fig:ch4.spectral-concentration}
\end{figure}

Figure~\ref{fig:ch4.spectral-concentration} is the concentration audit chain: \(\{\sigma_{n}\}\) is the primitive; \(\mathrm{EPR}\) and
\(d_{\mathrm{eff}}\) are derived certificates \cite{harrington2001,hansen1988,stratton1941}.

\begin{figure}[t]
\centering
\ChOneFigBlock{%
\begin{tikzpicture}[ch1fig, node distance=7mm and 9mm]
  \node[ch1box] (ell) {$\ell_{\mathrm{supp}}$};
  \node[ch1box, right=11mm of ell] (epr) {$\mathrm{EPR}$};
  \node[ch1box, right=11mm of epr] (tail) {$\mathcal{T}_{\mathrm{tail}}$};
  \draw[ch1flow] (ell) -- (epr);
  \draw[ch1flow] (epr) -- (tail);
\end{tikzpicture}%
}
\caption{Subsection~\ref{sec:ch483}: angular listing \(\ell_{\mathrm{supp}}\) and energy concentration \(\mathrm{EPR}\) are independent certificates.}
\label{fig:ch4.listing-concentration-split}
\end{figure}

\paragraph{Compression--concentration alignment.}
Eq.~\eqref{eq:ch4.spectral-compression-def} caps angular listing; Eq.~\eqref{eq:ch4.spectral-concentration-ratio} caps energy participation
\cite{hansen1988,stratton1941,harrington2001}. A pattern can satisfy \(\ell_{\mathrm{supp}}\) yet fail \(\Lambda_{\mathrm{conc}}\ge1-\varepsilon^{2}\) when
tail modes carry non-negligible power \cite{devaney2004}. Both must appear in synthesis memos \cite{stratton1941}.

\paragraph{Worked illustration (slow tail).}
\(\Lambda_{\mathrm{conc}}(10^{-2})=0.85\): \(15\%\) export gain remains outside leading subspace \cite{harrington2001,hansen1988,devaney2004}.

\paragraph{Worked illustration (flat tail).}
\(\sigma_{n}/\sigma_{1}=\{1,0.35,0.2,0.15,0.12,0.1,\ldots\}\): \(d_{\mathrm{eff}}(10^{-2})=3\), \(\Lambda_{\mathrm{conc}}\approx0.82\),
\(\mathcal{T}_{\mathrm{tail}}\approx0.18\) \cite{hansen1988,harrington2001}.

\paragraph{Worked illustration (steep knee).}
\(\sigma_{n}/\sigma_{1}=\{1,0.5,0.08,0.01,\ldots\}\): \(d_{\mathrm{eff}}(10^{-2})=2\), \(\mathrm{EPR}\approx0.99\) \cite{harrington2001,devaney2004,stratton1941}.

\paragraph{Worked illustration (sector flatness).}
\(\mathbf{M}\)-sector knee at \(d=4\), \(\mathbf{N}\)-sector knee at \(d=2\): aggregate \(d_{\mathrm{eff}}=4\) but \(\mathrm{EPR}=0.86\) due to flat
\(\mathbf{N}\)-tail \cite{harrington2001,hansen1988}.

\paragraph{Worked numerics (EPR sweep).}
\(\varepsilon\in\{10^{-1},10^{-2},10^{-3}\}\Rightarrow(d_{\mathrm{eff}},\mathrm{EPR})=(2,0.91),(5,0.97),(12,0.995)\) on \(k_{0}a=2.5\) aperture
\cite{hansen1988,harrington2001,devaney2004}.

\paragraph{Problem (concentration audit).}
\textbf{Problem.} Minimal \(N\) with \(\Lambda_{\mathrm{conc}}\ge0.99\)?
\textbf{Step 1.} Accumulate \(\sum_{n\le N}\sigma_{n}^{2}\) \cite{harrington2001}.
\textbf{Step 2.} Read \(d_{\mathrm{eff}}(\varepsilon')\) at knee Eq.~\eqref{eq:ch4.sv-knee} \cite{hansen1988,devaney2004}.
\textbf{Step 3.} Compare \(N\) to \(\ell_{\mathrm{supp}}\) from Eq.~\eqref{eq:ch4.spectral-compression-def} \cite{stratton1941}.
\textbf{Physical meaning.} Concentration quantifies truncation sacrifice \cite{hansen1988}.

\paragraph{Problem (participation target).}
\textbf{Problem.} Find \(\varepsilon\) with \(\mathrm{EPR}\ge0.99\).
\textbf{Step 1.} Sweep \(\varepsilon\) on \(\{\sigma_{n}\}\) \cite{harrington2001}.
\textbf{Step 2.} Report \((\varepsilon,d_{\mathrm{eff}})\) pair \cite{hansen1988,devaney2004}.
\textbf{Step 3.} Attach \(\mathcal{T}_{\mathrm{tail}}\) from Eq.~\eqref{eq:ch4.tail-discard} \cite{stratton1941}.
\textbf{Physical meaning.} EPR is the watts honesty metric \cite{harrington2001}.

\paragraph{Problem (tolerance trade).}
\textbf{Problem.} Require \(\Lambda_{\mathrm{conc}}\ge0.995\). Find \(d_{\mathrm{eff}}\).
\textbf{Step 1.} Accumulate partial sums \cite{harrington2001}.
\textbf{Step 2.} Read smallest \(N\) meeting target \cite{hansen1988}.
\textbf{Step 3.} Compare \(N\) to \(\ell_{\mathrm{supp}}\) and report both \cite{stratton1941,devaney2004}.
\textbf{Physical meaning.} Tighter concentration demands higher rank \cite{hansen1988}.

\paragraph{Problem (tail versus listing).}
\textbf{Problem.} \(\ell_{\mathrm{supp}}=8\) but \(\Lambda_{\mathrm{conc}}(10^{-2})=0.88\). Explain.
\textbf{Step 1.} Read \(\{\sigma_{n}\}\) knee via Eq.~\eqref{eq:ch4.sv-knee} \cite{hansen1988}.
\textbf{Step 2.} Compute \(\mathcal{T}_{\mathrm{tail}}\) from Eq.~\eqref{eq:ch4.tail-discard} \cite{harrington2001}.
\textbf{Step 3.} Conclude flat tail despite moderate listing \cite{devaney2004,stratton1941}.
\textbf{Physical meaning.} Listing and concentration diverge on slow tails \cite{hansen1988}.

\paragraph{Problem (EPR certificate).}
\textbf{Problem.} Attach \(\mathrm{EPR}\) at \(\varepsilon=10^{-2}\) to synthesis memo.
\textbf{Step 1.} Identify knee via Eq.~\eqref{eq:ch4.sv-knee} \cite{hansen1988,harrington2001}.
\textbf{Step 2.} Compute \(\mathrm{EPR}=\Lambda_{\mathrm{conc}}\) \cite{harrington2001}.
\textbf{Step 3.} Report \((d_{\mathrm{eff}},\mathcal{T}_{\mathrm{tail}},\ell_{\mathrm{supp}})\) triple \cite{stratton1941,devaney2004}.
\textbf{Physical meaning.} EPR quantifies watts retained in truncated synthesis \cite{hansen1988}.

\paragraph{Philosophical closure (truncation as optimum).}
Finite support forces a singular-value ledger whether or not the designer computes it \cite{stratton1941,hansen1988,harrington2001}. Truncation is
not a post hoc modeling choice; it is the unique rank-\(d_{\mathrm{eff}}\) approximation minimizing discarded export energy per
Theorem~\ref{thm:ch4.eckart-young-conc} \cite{devaney2004}. Claims of full harmonic synthesis without \(\mathrm{EPR}\) audit confuse coordinate
cardinality with energy concentration \cite{harrington2001,stratton1941}.

\begin{corollary}[Tail-dominated low rank]
\label{cor:ch4.tail-dominated}
If \(\sigma_{d_{\mathrm{eff}}+1}/\sigma_{d_{\mathrm{eff}}}>1-\varepsilon\), then \(\mathcal{T}_{\mathrm{tail}}(\varepsilon)>\varepsilon^{2}\) despite minimal rank increment
\cite{hansen1988,harrington2001,stratton1941,devaney2004}.
\end{corollary}

\paragraph{Validity.}
Eq.~\eqref{eq:ch4.spectral-concentration-ratio} depends on \(\mathcal{C}\) normalization and gauge alignment \cite{hansen1988,stratton1941}. Loss broadens
\(\{\sigma_{n}\}\) tails and may lower \(\mathrm{EPR}\) at fixed \(d_{\mathrm{eff}}\) \cite{harrington2001,devaney2004}.

\subsection{Landau-Pollak-Slepian limits and n-widths}
\label{sec:ch484}
% TOC tag: [HOME]

Landau--Pollak--Slepian (LPS) theory and Kolmogorov \(n\)-widths bound the number of functions simultaneously concentrated in a spatial domain
\(D\) and a spectral domain \(\Omega_{k}\) \cite{hansen1988,stratton1941}. Subsection~\ref{sec:ch484} imports those bounds as realizability ceilings
on \(d_{\mathrm{eff}}(\varepsilon)\), unifying LPS, mode-count, and SVD certificates through Eq.~\eqref{eq:ch4.lps-export-bridge}
\cite{harrington2001,devaney2004,jackson1999}. Concentration is the primitive fact; operator rank is the derived certificate
\cite{stratton1941,hansen1988}.

\paragraph{Assumptions.}
Declare spatial domain \(D\supset\Omega_{s}\) and spectral domain \(\Omega_{k}\) for propagating content \cite{stratton1941,jackson1999,hansen1988}.
Presuppose Eq.~\eqref{eq:ch4.duality-bandlimit-scaling} and linear time-harmonic maps \cite{harrington2001,devaney2004}. Tolerance \(\varepsilon\in(0,1)\)
fixed \cite{hansen1988}.

\paragraph{Kolmogorov \(n\)-width bound.}
The synthesis dimension is ceilinged by the concentration width
\begin{equation}
\label{eq:ch4.lps-nwidth-bound}
d_{\mathrm{eff}}(\varepsilon)\ \le\ N_{\mathrm{LPS}}(\varepsilon,D,\Omega_{k}),
\end{equation}
where \(N_{\mathrm{LPS}}\) counts approximate simultaneous eigenfunctions of spatial and spectral concentration operators on \((D,\Omega_{k})\)
\cite{harrington2001,devaney2004,jackson1999,hansen1988}.

\begin{theorem}[LPS bound on synthesis dimension]
\label{thm:ch4.lps-synthesis-bound}
\(\dim\mathcal{S}_{\mathrm{real}}(\mathcal{C})\le N_{\mathrm{LPS}}(\varepsilon,D,\Omega_{k})\) for any admissible coefficient chart \(\mathcal{C}\)
\cite{stratton1941,hansen1988,harrington2001,devaney2004}.
\end{theorem}
\begin{proof}
The concentration operator on \((D,\Omega_{k})\) dominates the export map \(\boldsymbol{\Gamma}_\omega\) on compact supports: no passive
\(\Jimp\) can excite more than \(N_{\mathrm{LPS}}\) approximately orthogonal export directions above noise floor \cite{hansen1988}. Definition
\ref{def:ch4.effective-em-dimension} counts those directions at \(\varepsilon\) \cite{harrington2001,stratton1941}.
\end{proof}

Theorem~\ref{thm:ch4.lps-synthesis-bound} explains why rank is finite before hardware enters: spatial--spectral concentration is a geometry-imposed
barrier independent of feed count \cite{hansen1988,stratton1941,harrington2001}.

\paragraph{Triple ceiling bridge.}
\begin{equation}
\label{eq:ch4.lps-export-bridge}
d_{\mathrm{eff}}(\varepsilon)\ \le\ \min\big(N_{\mathrm{LPS}}(\varepsilon,D,\Omega_{k}),\,N_{\mathrm{mode}}(\varepsilon)\big),
\end{equation}
unifying LPS, mode-count laws Eqs.~\eqref{eq:ch4.mode-count-volume}--\eqref{eq:ch4.mode-count-area}, and SVD rank \cite{hansen1988,stratton1941,harrington2001}.

\begin{theorem}[Triple ceiling consistency]
\label{thm:ch4.triple-ceiling}
On declared \((D,\Omega_{k},\mathcal{C})\), if any two of \(\{d_{\mathrm{eff}},N_{\mathrm{LPS}},N_{\mathrm{mode}}\}\) exceed the third, then
\((D,\Omega_{k})\) were mis-declared, \(\mathcal{C}\) is over-listed, or gauge alignment failed \cite{hansen1988,harrington2001,devaney2004,stratton1941}.
\end{theorem}
\begin{proof}
Theorem~\ref{thm:ch4.lps-synthesis-bound}, Theorem~\ref{thm:ch4.mode-ceiling-rank}, and Definition~\ref{def:ch4.effective-em-dimension} form a
consistent triple on aligned domains \cite{stratton1941,hansen1988,harrington2001}.
\end{proof}

\paragraph{Anisotropic concentration.}
On anisotropic \(\Omega_{s}\),
\begin{equation}
\label{eq:ch4.lps-anisotropy}
N_{\mathrm{LPS}}(\varepsilon,D,\Omega_{k})\le\sum_{i=1}^{3}N_{\mathrm{LPS}}^{(i)}(\varepsilon,D_{i},\Omega_{k}^{(i)}),
\end{equation}
\begin{equation}
\label{eq:ch4.prolate-nwidth}
N_{\mathrm{LPS}}^{(\mathrm{major})}\sim(k_{0}a)^{3},\qquad N_{\mathrm{LPS}}^{(\mathrm{minor})}\sim(k_{0}b)^{3}.
\end{equation}
Major-axis modes activate before minor-axis modes on prolate and slot apertures \cite{hansen1988,jackson1999,collin2001,harrington2001}.

\paragraph{Directional LPS ledger.}
On prolate \((a,b,c)\) with \(a\ge b\ge c\),
\begin{equation}
\label{eq:ch4.lps-directional-ledger}
N_{\mathrm{LPS}}^{(x)}\gtrsim(k_{0}a)^{3},\quad
N_{\mathrm{LPS}}^{(y)}\gtrsim(k_{0}b)^{3},\quad
N_{\mathrm{LPS}}^{(z)}\gtrsim(k_{0}c)^{3}.
\end{equation}
Isotropic \(N_{\mathrm{LPS}}\) on anisotropic \(\Omega_{s}\) overstates minor-axis export \cite{harrington2001,stratton1941,hansen1988}.

\paragraph{Prolate spheroidal link.}
One-dimensional PSWF concentration on \([-1,1]\) with bandwidth parameter \(c=k_{0}a\) yields \(N_{\mathrm{LPS}}^{(1\mathrm{D})}\sim c/\log c\)
\cite{hansen1988,jackson1999}. Embedding into three-dimensional prolate geometry recovers Eq.~\eqref{eq:ch4.prolate-nwidth} scaling
\cite{stratton1941,harrington2001,devaney2004}. EM export maps inherit the same barrier \cite{hansen1988}.

\begin{figure}[t]
\centering
\ChOneFigBlock{%
\begin{tikzpicture}[ch1fig, node distance=7mm and 9mm]
  \node[ch1box] (lps) {$N_{\mathrm{LPS}}$};
  \node[ch1box, right=11mm of lps] (deff) {$d_{\mathrm{eff}}$};
  \node[ch1box, right=11mm of deff] (sr) {$\mathcal{S}_{\mathrm{real}}$};
  \draw[ch1flow] (lps) -- (deff) -- (sr);
\end{tikzpicture}%
}
\caption{Subsection~\ref{sec:ch484}: \(n\)-width \(N_{\mathrm{LPS}}\) ceilings \(d_{\mathrm{eff}}\) and hence \(\mathcal{S}_{\mathrm{real}}(\mathcal{C})\).}
\label{fig:ch4.lps-nwidth}
\end{figure}

\begin{figure}[t]
\centering
\ChOneFigBlock{%
\begin{tikzpicture}[ch1fig, node distance=7mm and 9mm]
  \node[ch1box] (lps) {$N_{\mathrm{LPS}}$};
  \node[ch1box, right=11mm of lps] (nm) {$N_{\mathrm{mode}}$};
  \node[ch1box, right=11mm of nm] (deff) {$d_{\mathrm{eff}}$};
  \draw[ch1flow] (lps) -- (deff);
  \draw[ch1flow] (nm) -- (deff);
\end{tikzpicture}%
}
\caption{Subsection~\ref{sec:ch484}: triple ceiling Eq.~\eqref{eq:ch4.lps-export-bridge}; tightest bound governs.}
\label{fig:ch4.triple-ceiling}
\end{figure}

Figure~\ref{fig:ch4.triple-ceiling} is the consistency audit: \(d_{\mathrm{eff}}\) from SVD must not exceed either analytic ceiling without domain
mismatch \cite{hansen1988,harrington2001,stratton1941,devaney2004}.

\paragraph{Triple-consistency workflow.}
Certification sequence: (i) declare \((D,\Omega_{k})\); (ii) compute directional \(N_{\mathrm{LPS}}^{(i)}\); (iii) compute \(N_{\mathrm{mode}}\);
(iv) measure \(d_{\mathrm{eff}}\) from SVD; (v) apply Eq.~\eqref{eq:ch4.lps-export-bridge} and Theorem~\ref{thm:ch4.triple-ceiling}
\cite{hansen1988,harrington2001,devaney2004,stratton1941}.

\paragraph{Worked illustration (prolate aperture).}
Prolate spheroid \(k_{0}a=3\), \(k_{0}b=0.5\): \(N_{\mathrm{LPS}}^{(x)}\approx27\), \(N_{\mathrm{LPS}}^{(y)}\approx1\); major axis dominates
\cite{hansen1988,jackson1999,harrington2001,collin2001}.

\paragraph{Worked illustration (slot versus disk).}
Equal-area slot (\(10:1\)) and disk at \(k_{0}a=2\): disk \(N_{\mathrm{LPS}}^{(\mathrm{vol})}\sim8\); slot major-axis \(N_{\mathrm{LPS}}^{(x)}\approx7\),
minor-axis \(\approx1\) \cite{hansen1988,collin2001,stratton1941}.

\paragraph{Worked illustration (SVD exceeds LPS).}
MoM on slot yields \(d_{\mathrm{eff}}=42\), analytic \(N_{\mathrm{LPS}}=28\): \(\mathcal{C}\) likely includes evanescent rows outside \(\Omega_{k}\)
\cite{harrington2001,devaney2004,hansen1988,stratton1941}.

\paragraph{Worked illustration (triple agreement).}
\(k_{0}a=2\), disk: \(N_{\mathrm{LPS}}=12\), \(N_{\mathrm{mode}}=14\), \(d_{\mathrm{eff}}(10^{-2})=10\): consistent triple; governing value is
\(d_{\mathrm{eff}}\) \cite{hansen1988,harrington2001,jackson1999}.

\paragraph{Worked numerics (triple audit).}
Slot \(10{:}1\), \(k_{0}a=2.5\): \(N_{\mathrm{LPS}}^{(x)}\approx18\), \(N_{\mathrm{LPS}}^{(y)}\approx3\), \(N_{\mathrm{mode}}^{\mathrm{(ap)}}\approx22\),
\(d_{\mathrm{eff}}(10^{-2})=14\) \cite{hansen1988,collin2001,harrington2001,devaney2004}.

\paragraph{Problem (LPS versus SVD).}
\textbf{Problem.} \(d_{\mathrm{eff}}=40\), \(N_{\mathrm{LPS}}=35\). Which governs?
\textbf{Step 1.} Apply Theorem~\ref{thm:ch4.lps-synthesis-bound} as ceiling \cite{hansen1988,stratton1941}.
\textbf{Step 2.} Audit \((D,\Omega_{k})\) and strip non-export rows from \(\mathcal{C}\) \cite{harrington2001,devaney2004}.
\textbf{Step 3.} Use \(\min(d_{\mathrm{eff}},N_{\mathrm{LPS}})\) after audit \cite{jackson1999}.
\textbf{Physical meaning.} LPS bounds concentration, not raw listing \cite{hansen1988}.

\paragraph{Problem (anisotropic LPS).}
\textbf{Problem.} Isotropic \(N_{\mathrm{LPS}}\) on slot aperture. Conservative?
\textbf{Step 1.} Compute directional \(N_{\mathrm{LPS}}^{(i)}\) via Eq.~\eqref{eq:ch4.lps-directional-ledger} \cite{hansen1988}.
\textbf{Step 2.} Sum per Eq.~\eqref{eq:ch4.lps-anisotropy} \cite{jackson1999,stratton1941}.
\textbf{Step 3.} Compare to \(d_{\mathrm{eff}}\) and \(N_{\mathrm{mode}}\) \cite{harrington2001,collin2001}.
\textbf{Physical meaning.} LPS is geometry-directional \cite{hansen1988}.

\paragraph{Problem (LPS audit).}
\textbf{Problem.} Compute conservative \(N_{\mathrm{LPS}}\) for prolate \(k_{0}a=3\), \(k_{0}b=0.5\).
\textbf{Step 1.} Evaluate Eq.~\eqref{eq:ch4.prolate-nwidth} per axis \cite{hansen1988,jackson1999}.
\textbf{Step 2.} Sum via Eq.~\eqref{eq:ch4.lps-anisotropy} \cite{stratton1941,harrington2001}.
\textbf{Step 3.} Compare to measured \(d_{\mathrm{eff}}\) \cite{devaney2004}.
\textbf{Physical meaning.} LPS is geometry-directional concentration bound \cite{hansen1988}.

\paragraph{Problem (triple consistency).}
\textbf{Problem.} \(d_{\mathrm{eff}}=18\), \(N_{\mathrm{mode}}=22\), \(N_{\mathrm{LPS}}=15\). Resolve.
\textbf{Step 1.} Apply Theorem~\ref{thm:ch4.triple-ceiling} \cite{hansen1988,stratton1941}.
\textbf{Step 2.} Use \(\min\) in Eq.~\eqref{eq:ch4.lps-export-bridge} \cite{harrington2001}.
\textbf{Step 3.} Audit \(\mathcal{C}\) for over-listing if \(d_{\mathrm{eff}}>N_{\mathrm{LPS}}\) \cite{devaney2004}.
\textbf{Physical meaning.} Tightest ceiling governs realizability \cite{stratton1941}.

\paragraph{Problem (domain mismatch).}
\textbf{Problem.} MoM \(d_{\mathrm{eff}}=38\), analytic \(N_{\mathrm{LPS}}=26\). Diagnose.
\textbf{Step 1.} Verify \((D,\Omega_{k})\) declarations match \(\Omega_{s}\) \cite{hansen1988,stratton1941}.
\textbf{Step 2.} Remove evanescent chart rows outside \(\Omega_{k}\) \cite{harrington2001,devaney2004}.
\textbf{Step 3.} Recompute SVD; expect \(d_{\mathrm{eff}}\le26\) \cite{jackson1999,hansen1988}.
\textbf{Physical meaning.} Over-listed charts inflate rank \cite{stratton1941}.

\paragraph{Problem (disk versus slot).}
\textbf{Problem.} Equal-area disk and \(10{:}1\) slot at \(k_{0}a=2\). Compare \(N_{\mathrm{LPS}}\).
\textbf{Step 1.} Compute isotropic disk \(N_{\mathrm{LPS}}^{(\mathrm{vol})}\) \cite{hansen1988,jackson1999}.
\textbf{Step 2.} Compute directional slot ledger Eq.~\eqref{eq:ch4.lps-directional-ledger} \cite{collin2001,stratton1941}.
\textbf{Step 3.} Conclude minor-axis slot export is strongly suppressed \cite{harrington2001}.
\textbf{Physical meaning.} Shape sets directional concentration \cite{hansen1988}.

\paragraph{Philosophical closure (concentration inevitability).}
\(n\)-width bounds are not parallel to operator rank; they explain why rank is finite on compact supports \cite{hansen1988,stratton1941,harrington2001}.
Spatial concentration and spectral bandlimit are conjugate resources: enlarging one shrinks the other in the LPS sense \cite{jackson1999,devaney2004}.
Synthesis dimension is a concentration fact before it is a hardware fact \cite{harrington2001,stratton1941}.

\begin{corollary}[Domain declaration obligation]
\label{cor:ch4.lps-domain}
\(d_{\mathrm{eff}}>N_{\mathrm{LPS}}(\varepsilon,D,\Omega_{k})\) on passive \(\Omega_{s}\) implies \((D,\Omega_{k})\) were mis-declared or
\(\mathcal{C}\) includes non-export rows \cite{hansen1988,harrington2001,devaney2004,stratton1941}.
\end{corollary}

\paragraph{Validity.}
Eq.~\eqref{eq:ch4.lps-nwidth-bound} assumes linear time-harmonic maps, declared \((D,\Omega_{k})\), and narrowband phasors \cite{stratton1941,hansen1988}.
Cavities and guided structures quantize \(\Omega_{k}\) differently; LPS ceilings must be recomputed on the modified domain \cite{harrington2001,jackson1999,devaney2004}.


Section~\ref{sec:ch48} has linked mode count, asymptotic regimes, spectral concentration, and LPS \(n\)-widths to \(d_{\mathrm{eff}}(\varepsilon)\)
\cite{stratton1941,hansen1988,harrington2001,jackson1999,devaney2004}. Truncation is inevitable because every bound is finite at fixed tolerance.
Section~\ref{sec:ch49} records physical consequences for engineering and measurement on the same realizability framework \cite{stratton1941,harrington2001}.


\section{Physical Consequences of Realizability Limits}
\label{sec:ch49}

Truncation and rank are not abstract operator facts: they bound what finite apertures can radiate, what phase-only hardware can retune,
and what laboratories can observe \cite{stratton1941,harrington2001,devaney2004}. The consequences collected here are the engineering and
measurement face of the same realizability architecture developed in Sections~\ref{sec:ch41}--\ref{sec:ch48} \cite{hansen1988,jackson1999}.
Representation is coordinate freedom; realizability is geometry-imposed rank. Every design memo that confuses listing cardinality with export
rank inherits the inequalities already proved on \(\boldsymbol{\Gamma}_\omega\) \cite{stratton1941,harrington2001,devaney2004}.

\subsection{Geometry as origin of dimensionality limits}
\label{sec:ch491}
% TOC tag: [HOME]

Dimensionality limits in synthesis reports trace to \(\operatorname{diam}(\Omega_{s})\), \(A_{a}\), and \(k_{0}\) through
\(d_{\mathrm{eff}}(\varepsilon)\), not to solver tolerances or mesh cardinalities alone \cite{stratton1941,jackson1999,hansen1988,harrington2001}.
Finite support compresses the export map \(\boldsymbol{\Gamma}_\omega\) on \(\mathcal{J}_{\mathrm{real}}\); the resulting singular-value ladder is
monotone in electrical size and invariant under phase-only reparameterizations on fixed \(\Omega_{s}\) \cite{devaney2004,harrington2001}. Listing
\(\mathcal{C}\) may print arbitrarily many indices; realizability counts only directions whose aligned \(\sigma_{n}\) exceeds \(\varepsilon\sigma_{1}\)
\cite{stratton1941,hansen1988}.

\paragraph{Assumptions.}
Presuppose Definition~\ref{def:ch4.effective-em-dimension}, Eqs.~\eqref{eq:ch4.mode-count-volume}--\eqref{eq:ch4.mode-count-area}, and
Theorem~\ref{thm:ch4.rank-ka-monotone} \cite{stratton1941,hansen1988,harrington2001,jackson1999}. Fix tolerance \(\varepsilon\in(0,1)\), aligned flux
gauge on \(\mathcal{C}\), and passive linear exterior on \(\mathcal{J}_{\mathrm{real}}\) \cite{devaney2004}. Rank audits report \((k_{0},D_{s},\varepsilon)\)
triples with declared geometry class (volume versus area) \cite{collin2001,stratton1941}.

\paragraph{Hardware parameter map.}
Export rank is a functional of geometry and frequency:
\begin{equation}
\label{eq:ch4.hardware-rank-map}
d_{\mathrm{eff}}(\varepsilon)=F\big(\operatorname{diam}(\Omega_{s}),k_{0},\varepsilon;\mathcal{C}\big),
\end{equation}
with \(F\) monotone in \(\operatorname{diam}(\Omega_{s})\) and \(k_{0}\) per Theorem~\ref{thm:ch4.rank-ka-monotone} \cite{stratton1941,jackson1999,hansen1988}.
Mathematically, Eq.~\eqref{eq:ch4.hardware-rank-map} is the counting functional induced by \(\boldsymbol{\Gamma}_\omega\). Physically, it is the number
of independent far-zone directions a laboratory can excite above the noise floor at declared tolerance \cite{harrington2001,devaney2004}.

\begin{theorem}[Geometry as rank origin]
\label{thm:ch4.geometry-rank-origin}
At fixed \(\varepsilon\) and listing \(\mathcal{C}\), \(d_{\mathrm{eff}}(\varepsilon)\) is invariant under phase-only
reparameterizations of \(\Jimp\) on fixed \(\Omega_{s}\) and changes only when \(\operatorname{diam}(\Omega_{s})\), \(k_{0}\), or
\(\mathcal{J}_{\mathrm{real}}\) class changes \cite{stratton1941,jackson1999,hansen1988,harrington2001,devaney2004}.
\end{theorem}
\begin{proof}
\(\boldsymbol{\Gamma}_\omega\) depends on support and material class through \(\mathcal{R}_\omega\) and \(\mathcal{C}\) \cite{stratton1941,harrington2001}.
Phase rotations on fixed support are unitary gauge transformations on \(\mathcal{J}_{\mathrm{real}}\) that preserve singular values of the export map
\cite{hansen1988,devaney2004}. Enlarging \(\operatorname{diam}(\Omega_{s})\) or raising \(k_{0}\) enlarges the set of propagating channels coupled by
\(\boldsymbol{\Gamma}_\omega\) until dispersion alters \(\mathcal{J}_{\mathrm{real}}\) \cite{jackson1999,collin2001}.
\end{proof}

Theorem~\ref{thm:ch4.geometry-rank-origin} is the origin principle for hardware audits: diversity saturation on unchanged outline is a geometry certificate,
not a solver defect \cite{stratton1941,harrington2001}. Metasurfaces that retune phases without enlarging \(\mathcal{L}_{2}\) cannot raise \(d_{\mathrm{eff}}\)
\cite{hansen1988,devaney2004}. Rank growth requires electrical enlargement---frequency, aperture, or material class---not mesh refinement alone
\cite{jackson1999,collin2001}.

\paragraph{Electrical-size monotonicity.}
On fixed \(\Omega_{s}\) class,
\begin{equation}
\label{eq:ch4.outline-trap}
\frac{\partial d_{\mathrm{eff}}}{\partial k_{0}}\Big|_{\operatorname{diam}(\Omega_{s})=\text{const}}>0\quad\text{for }k_{0}D_{s}\gtrsim1.
\end{equation}
Mechanical CAD outline is invariant under band change, but \(k_{0}D_{s}\) is not \cite{jackson1999,harrington2001}. Rank memos must quote electrical
size, not enclosure volume alone \cite{hansen1988,stratton1941,devaney2004}.

\begin{lemma}[Mesh DOF versus export rank]
\label{lem:ch4.mesh-dof-trap}
FEM or MoM degrees of freedom \(N_{\mathrm{dof}}\) satisfy \(N_{\mathrm{dof}}\ge d_{\mathrm{eff}}(\varepsilon)\) but \(N_{\mathrm{dof}}\gg d_{\mathrm{eff}}\)
does not imply rank growth on fixed \((\Omega_{s},k_{0})\) \cite{harrington2001,devaney2004,hansen1988,stratton1941}.
\end{lemma}
\begin{proof}
Discretization enlarges the trial space for \(\Jimp\) without enlarging \(\operatorname{ran}\boldsymbol{\Gamma}_\omega\) on fixed compact support
\cite{harrington2001,stratton1941}. Singular values \(\sigma_{n}\) of \(\boldsymbol{\Gamma}_\omega\) are properties of the continuum export map;
mesh refinement approximates them but cannot create new propagating channels on fixed \((\Omega_{s},k_{0})\) \cite{devaney2004,hansen1988}.
\end{proof}

\begin{corollary}[CAD outline decoupling]
\label{cor:ch4.cad-rank-decouple}
Identical mechanical outlines at different \(k_{0}\) generally yield different \(d_{\mathrm{eff}}(\varepsilon)\) via Eq.~\eqref{eq:ch4.joint-scaling-law}
\cite{jackson1999,harrington2001,hansen1988,stratton1941,devaney2004}.
\end{corollary}

\paragraph{Volume versus area geometry class.}
Compact three-dimensional \(\Omega_{s}\) obey Eq.~\eqref{eq:ch4.mode-count-volume}; planar apertures obey Eq.~\eqref{eq:ch4.mode-count-area}
\cite{collin2001,jackson1999,hansen1988}. Declaring the wrong class over-estimates or under-estimates \(N_{\mathrm{mode}}\) and mis-attributes rank
saturation \cite{stratton1941,harrington2001,devaney2004}.

\begin{figure}[t]
\centering
\ChOneFigBlock{%
\begin{tikzpicture}[ch1fig, node distance=7mm and 9mm]
  \node[ch1box] (geom) {$(\Omega_{s},k_{0})$};
  \node[ch1box, right=11mm of geom] (gamma) {$\boldsymbol{\Gamma}_\omega$};
  \node[ch1box, right=11mm of gamma] (deff) {$d_{\mathrm{eff}}$};
  \draw[ch1flow] (geom) -- (gamma) -- (deff);
\end{tikzpicture}%
}
\caption{Geometry and frequency set the export map; \(d_{\mathrm{eff}}\) counts singular directions above \(\varepsilon\sigma_{1}\).}
\label{fig:ch4.geometry-rank-map}
\end{figure}

Figure~\ref{fig:ch4.geometry-rank-map} is the hardware honesty chain: solvers and optimizers operate downstream of \((\Omega_{s},k_{0})\); rank is fixed
upstream by geometry \cite{stratton1941,harrington2001,devaney2004}.

\begin{figure}[t]
\centering
\ChOneFigBlock{%
\begin{tikzpicture}[ch1fig, node distance=7mm and 9mm]
  \node[ch1box] (low) {$k_{0}D_{s}<1$};
  \node[ch1box, right=11mm of low] (mid) {$k_{0}D_{s}\sim1$};
  \node[ch1box, right=11mm of mid] (high) {$k_{0}D_{s}\gg1$};
  \draw[ch1flow] (low) -- (mid) -- (high);
\end{tikzpicture}%
}
\caption{Electrical-size ladder: \(d_{\mathrm{eff}}\) rises with \(k_{0}D_{s}\) at fixed mechanical outline.}
\label{fig:ch4.electrical-size-ladder}
\end{figure}

\paragraph{Worked illustration (metasurface phase).}
A fixed-area metasurface that retunes phases without enlarging \(\mathcal{L}_{2}\) leaves \(d_{\mathrm{eff}}\) unchanged while rotating \(\mathbf{c}\)
within \(\operatorname{ran}\boldsymbol{\Gamma}_\omega\) \cite{stratton1941,harrington2001,hansen1988,devaney2004}.

\paragraph{Worked illustration (same phone, two bands).}
\(150\,\mathrm{mm}\) outline: \(k_{0}D_{s}\approx0.8\) at \(700\,\mathrm{MHz}\) (\(d_{\mathrm{eff}}\le3\)) versus \(k_{0}D_{s}\approx8\) at
\(28\,\mathrm{GHz}\) (\(d_{\mathrm{eff}}\sim\mathcal{O}(10^{2})\)) \cite{harrington2001,collin2001,hansen1988,jackson1999}.

\paragraph{Worked illustration (mesh saturation).}
FEM rank grows from \(d_{\mathrm{eff}}=8\) to \(N_{\mathrm{dof}}=4000\) on fixed \(\lambda/4\) patch without changing measured \(d_{\mathrm{eff}}\):
Lemma~\ref{lem:ch4.mesh-dof-trap} \cite{harrington2001,devaney2004,hansen1988}.

\paragraph{Worked illustration (panel enlargement).}
Doubling \(W\) at fixed \(k_{0}\) quadruples area-law mode budget per Eq.~\eqref{eq:ch4.mode-count-area} and raises \(d_{\mathrm{eff}}\) per
Theorem~\ref{thm:ch4.rank-ka-monotone} \cite{collin2001,jackson1999,hansen1988,harrington2001}.

\paragraph{Worked numerics (joint scaling).}
\(k_{0}a=1.2\): volume ceiling \(\sim(k_{0}a)^{3}\approx1.7\), area ceiling \(\sim(k_{0}a)^{2}\approx1.4\); measured \(d_{\mathrm{eff}}(10^{-2})=4\) lies
below both analytic ceilings as expected \cite{hansen1988,stratton1941,devaney2004}.

\paragraph{Worked numerics (frequency doubling).}
Fixed \(A_{a}\), \(\omega\to2\omega\): Eq.~\eqref{eq:ch4.joint-scaling-law} predicts rank growth \(\propto k_{0}^{2}\) or \(k_{0}^{3}\) by geometry class;
measured \(d_{\mathrm{eff}}\) tracks electrical enlargement \cite{jackson1999,harrington2001,collin2001}.

\paragraph{Problem (solver rank versus hardware).}
\textbf{Problem.} FEM rank grows with mesh; diversity saturates. Explain.
\textbf{Step 1.} Report \(d_{\mathrm{eff}}(\varepsilon)\), not \(N_{\mathrm{dof}}\) \cite{harrington2001,devaney2004}.
\textbf{Step 2.} Tie saturation to \(k_{0}D_{s}\) via Eq.~\eqref{eq:ch4.joint-scaling-law} \cite{jackson1999,hansen1988}.
\textbf{Step 3.} Enlarge aperture or frequency to raise rank \cite{stratton1941}.
\textbf{Physical meaning.} Geometry is the origin of dimensionality \cite{stratton1941,harrington2001}.

\paragraph{Problem (outline versus rank).}
\textbf{Problem.} CAD outline unchanged; frequency doubled. Rank change?
\textbf{Step 1.} Compute \(k_{0}D_{s}\) \cite{jackson1999}.
\textbf{Step 2.} Apply Eq.~\eqref{eq:ch4.joint-scaling-law} \cite{hansen1988,harrington2001}.
\textbf{Step 3.} Recompute \(d_{\mathrm{eff}}\) \cite{devaney2004}.
\textbf{Physical meaning.} Rank tracks electrical size, not mechanical outline alone \cite{stratton1941,collin2001}.

\paragraph{Problem (metasurface diversity claim).}
\textbf{Problem.} Vendor claims doubled MIMO streams after phase-only retuning on fixed area. Audit.
\textbf{Step 1.} Test Theorem~\ref{thm:ch4.geometry-rank-origin} \cite{stratton1941,harrington2001}.
\textbf{Step 2.} Measure \(d_{\mathrm{eff}}\) before and after \cite{hansen1988,devaney2004}.
\textbf{Step 3.} Reject if \(d_{\mathrm{eff}}\) unchanged \cite{jackson1999}.
\textbf{Physical meaning.} Phase explores fixed-rank geometry \cite{harrington2001}.

\paragraph{Problem (volume versus area class).}
\textbf{Problem.} Compact resonator modeled as planar patch. Rank over-estimated?
\textbf{Step 1.} Declare geometry class \cite{collin2001,stratton1941}.
\textbf{Step 2.} Compare Eqs.~\eqref{eq:ch4.mode-count-volume} and \eqref{eq:ch4.mode-count-area} \cite{jackson1999,hansen1988}.
\textbf{Step 3.} Apply operative ceiling to \(d_{\mathrm{eff}}\) audit \cite{harrington2001,devaney2004}.
\textbf{Physical meaning.} Geometry class sets the counting law \cite{stratton1941}.

\paragraph{Problem (handset band audit).}
\textbf{Problem.} Same enclosure, LTE versus mmWave. Rank comparison?
\textbf{Step 1.} Compute \(k_{0}D_{s}\) per band \cite{jackson1999,harrington2001}.
\textbf{Step 2.} Apply Corollary~\ref{cor:ch4.cad-rank-decouple} \cite{hansen1988}.
\textbf{Step 3.} Report separate \(d_{\mathrm{eff}}\) certificates \cite{devaney2004,collin2001}.
\textbf{Physical meaning.} Outline invariance does not imply rank invariance \cite{stratton1941}.

\paragraph{Problem (mesh refinement trap).}
\textbf{Problem.} MoM residual improves with mesh; \(d_{\mathrm{eff}}\) flat. Consistent?
\textbf{Step 1.} Apply Lemma~\ref{lem:ch4.mesh-dof-trap} \cite{harrington2001,devaney2004}.
\textbf{Step 2.} Verify \((\Omega_{s},k_{0})\) unchanged \cite{stratton1941,hansen1988}.
\textbf{Step 3.} Attribute flat rank to geometry, not discretization \cite{jackson1999}.
\textbf{Physical meaning.} Mesh resolves currents, not export channels \cite{harrington2001}.

\paragraph{Rank certificate workflow.}
Laboratory rank memos should report the quadruple \((k_{0}D_{s},\varepsilon,d_{\mathrm{eff}},N_{\mathrm{mode}})\) with geometry class declared
\cite{stratton1941,harrington2001,hansen1988,jackson1999}. Solver mesh \(N_{\mathrm{dof}}\) and chart size \(|\mathbb{C}^{I}|\) are diagnostic overlays,
not substitutes for \(d_{\mathrm{eff}}\) \cite{devaney2004,collin2001}.

\begin{theorem}[Diameter monotonicity of export rank]
\label{thm:ch4.diam-rank-monotone}
At fixed \(k_{0}\), \(\varepsilon\), and \(\mathcal{J}_{\mathrm{real}}\) class, \(d_{\mathrm{eff}}(\varepsilon)\) is non-decreasing in \(\operatorname{diam}(\Omega_{s})\)
\cite{stratton1941,hansen1988,harrington2001,jackson1999,devaney2004}.
\end{theorem}
\begin{proof}
Enlarging \(\operatorname{diam}(\Omega_{s})\) enlarges coupled propagating channels in \(\boldsymbol{\Gamma}_\omega\) until material dispersion alters
\(\mathcal{J}_{\mathrm{real}}\) \cite{harrington2001,stratton1941}. \(d_{\mathrm{eff}}\) counts singular values above \(\varepsilon\sigma_{1}\), which cannot
decrease under support enlargement on fixed class \cite{hansen1988,devaney2004,jackson1999}.
\end{proof}

\paragraph{Discretization scaling trap.}
MoM or FEM bases of size \(N_{\mathrm{basis}}\) approximate \(\boldsymbol{\Gamma}_\omega\) with error \(\tau_{\mathrm{mom}}\) on \(\operatorname{ran}\boldsymbol{\Gamma}_\omega\)
but cannot raise continuum rank on fixed \((\Omega_{s},k_{0})\) \cite{harrington2001,devaney2004}. Reports that equate \(N_{\mathrm{basis}}\) with diversity
confuse approximation cardinality with export rank \cite{stratton1941,hansen1988,collin2001}.

\paragraph{Worked illustration (base station footprint).}
\(2\lambda\times2\lambda\) panel at \(3.5\,\mathrm{GHz}\): \(k_{0}W\approx14\), \(N_{\mathrm{mode}}^{\mathrm{(ap)}}\sim\mathcal{O}(10^{2})\), measured
\(d_{\mathrm{eff}}(10^{-2})\sim\mathcal{O}(10^{1})\)--\(\mathcal{O}(10^{2})\) per Eq.~\eqref{eq:ch4.joint-scaling-law} \cite{collin2001,harrington2001,hansen1988,jackson1999}.

\paragraph{Worked illustration (wearable).}
\(k_{0}D_{s}<1\) wearable at \(2.4\,\mathrm{GHz}\): \(d_{\mathrm{eff}}\le3\) regardless of mesh density; electrical size sets dipole-class rank
\cite{harrington2001,jackson1999,stratton1941,devaney2004}.

\paragraph{Worked numerics (diameter sweep).}
\(D_{s}=0.5\lambda,1.0\lambda,1.5\lambda\) at fixed \(k_{0}\): \(d_{\mathrm{eff}}(10^{-2})\) increases monotonically per Theorem~\ref{thm:ch4.diam-rank-monotone}
\cite{hansen1988,harrington2001,stratton1941,devaney2004,collin2001}.

\paragraph{Problem (solver versus physics).}
\textbf{Problem.} Optimizer reports \(200\) feeds; measured \(d_{\mathrm{eff}}=14\). Explain.
\textbf{Step 1.} Apply Lemma~\ref{lem:ch4.mesh-dof-trap} \cite{harrington2001,devaney2004}.
\textbf{Step 2.} Report \(d_{\mathrm{eff}}\) as physics certificate \cite{hansen1988,stratton1941}.
\textbf{Step 3.} Truncate design to rank \cite{jackson1999,collin2001}.
\textbf{Physical meaning.} Feeds \(\neq\) channels \cite{harrington2001}.

\paragraph{Problem (diameter doubling).}
\textbf{Problem.} Double \(\operatorname{diam}(\Omega_{s})\) at fixed \(k_{0}\). Rank trend?
\textbf{Step 1.} Apply Theorem~\ref{thm:ch4.diam-rank-monotone} \cite{stratton1941,hansen1988}.
\textbf{Step 2.} Bound \(N_{\mathrm{mode}}\) via Eqs.~\eqref{eq:ch4.mode-count-volume} or \eqref{eq:ch4.mode-count-area} \cite{jackson1999,harrington2001}.
\textbf{Step 3.} Measure SVD rank \cite{devaney2004,collin2001}.
\textbf{Physical meaning.} Larger support raises rank \cite{stratton1941}.

\paragraph{Problem (aperture growth).}
\textbf{Problem.} Double aperture width at fixed \(k_{0}\). Expected \(d_{\mathrm{eff}}\) trend?
\textbf{Step 1.} Apply Eq.~\eqref{eq:ch4.mode-count-area} \cite{collin2001,hansen1988}.
\textbf{Step 2.} Test Theorem~\ref{thm:ch4.rank-ka-monotone} \cite{harrington2001,jackson1999}.
\textbf{Step 3.} Measure SVD rank on enlarged \(\Omega_{s}\) \cite{devaney2004,stratton1941}.
\textbf{Physical meaning.} Larger aperture raises electrical mode budget \cite{hansen1988}.

\paragraph{Geometry scaling certificate.}
Joint scaling Eq.~\eqref{eq:ch4.joint-scaling-law} combines volume and area ceilings; the operative bound is
\(\min\{C_{V}(k_{0}a)^{3},C_{A}k_{0}^{2}A_{a}\}\) on declared class \cite{jackson1999,hansen1988,harrington2001,stratton1941,devaney2004,collin2001}.
Rank audits that cite only mechanical area without \(k_{0}\) mis-state electrical mode budget \cite{stratton1941,harrington2001}.

\paragraph{Worked illustration (satellite dish).}
\(\lambda/10\) mesh on \(10\lambda\) dish at \(12\,\mathrm{GHz}\): \(N_{\mathrm{dof}}\sim10^{4}\), \(d_{\mathrm{eff}}\sim10^{2}\): two orders of magnitude
separate discretization from export rank per Lemma~\ref{lem:ch4.mesh-dof-trap} \cite{harrington2001,hansen1988,jackson1999,devaney2004,collin2001}.

\paragraph{Worked numerics (wearable mmWave).}
\(k_{0}D_{s}=0.6\) at \(28\,\mathrm{GHz}\) on \(150\,\mathrm{mm}\) outline: \(d_{\mathrm{eff}}\le4\) despite \(|\mathbb{C}^{I}|=120\) chart listing
\cite{harrington2001,jackson1999,stratton1941,hansen1988,devaney2004}.

\paragraph{Problem (joint scaling class).}
\textbf{Problem.} Compact resonator versus patch array at same footprint. Which ceiling?
\textbf{Step 1.} Declare volume versus area class \cite{collin2001,stratton1941}.
\textbf{Step 2.} Evaluate Eq.~\eqref{eq:ch4.joint-scaling-law} \cite{jackson1999,hansen1988}.
\textbf{Step 3.} Apply operative minimum \cite{harrington2001,devaney2004}.
\textbf{Physical meaning.} Class selects exponent \cite{stratton1941}.

\paragraph{Problem (chart listing).}
\textbf{Problem.} \(\ell_{\max}=40\) on \(k_{0}a=1.2\). Rank implication?
\textbf{Step 1.} Read \(\ell_{\mathrm{supp}}(10^{-2},1.2)\) \cite{hansen1988,harrington2001}.
\textbf{Step 2.} Apply Corollary~\ref{cor:ch4.listing-excess-no-rank} \cite{stratton1941,devaney2004}.
\textbf{Step 3.} Report listing excess separately from \(d_{\mathrm{eff}}\) \cite{jackson1999,collin2001}.
\textbf{Physical meaning.} Charts over-list by design \cite{harrington2001}.

\paragraph{Problem (electrical ladder).}
\textbf{Problem.} Move along Figure~\ref{fig:ch4.electrical-size-ladder} at fixed outline. Rank trend?
\textbf{Step 1.} Compute \(k_{0}D_{s}\) per band \cite{jackson1999,harrington2001}.
\textbf{Step 2.} Apply Theorem~\ref{thm:ch4.rank-ka-monotone} \cite{hansen1988,stratton1941}.
\textbf{Step 3.} Tabulate \(d_{\mathrm{eff}}\) versus band \cite{devaney2004,collin2001}.
\textbf{Physical meaning.} Electrical ladder sets rank \cite{stratton1941,harrington2001}.

\paragraph{Philosophical closure (origin versus listing).}
Dimensionality is imposed by \((\Omega_{s},k_{0},\mathcal{J}_{\mathrm{real}})\) through \(\boldsymbol{\Gamma}_\omega\), not by how many indices \(\mathcal{C}\)
prints or how fine a mesh resolves \(\Jimp\) \cite{stratton1941,harrington2001,jackson1999,hansen1988,devaney2004}. Every synthesis report that attributes
rank saturation to optimizer failure without auditing electrical size mis-identifies the origin of truncation \cite{collin2001,stratton1941}.

\paragraph{Singular-value ladder interpretation.}
\(d_{\mathrm{eff}}(\varepsilon)\) counts the singular-value ladder of \(\boldsymbol{\Gamma}_\omega\) above \(\varepsilon\sigma_{1}\); geometry sets the ladder
through support and \(k_{0}\), not through chart padding \cite{harrington2001,hansen1988,stratton1941,devaney2004,jackson1999,collin2001}. Reports that
tabulate \(\sigma_{n}\) without \((\Omega_{s},k_{0})\) context are incomplete rank certificates \cite{stratton1941,harrington2001}.

\paragraph{Worked illustration (IoT module).}
\(50\,\mathrm{mm}\) IoT module at \(915\,\mathrm{MHz}\): \(k_{0}D_{s}\approx0.15\), \(d_{\mathrm{eff}}\le2\); rank is electrically tiny despite large CAD model
\cite{harrington2001,jackson1999,stratton1941,hansen1988,devaney2004,collin2001}.

\paragraph{Problem (sigma ladder audit).}
\textbf{Problem.} Tabulate \(\sigma_{1..N}\) for vendor audit. What else is required?
\textbf{Step 1.} Declare \((\Omega_{s},k_{0},\varepsilon)\) \cite{stratton1941,harrington2001}.
\textbf{Step 2.} Count \(d_{\mathrm{eff}}\) above \(\varepsilon\sigma_{1}\) \cite{hansen1988,devaney2004}.
\textbf{Step 3.} Compare to \(N_{\mathrm{mode}}\) ceiling \cite{jackson1999,collin2001}.
\textbf{Physical meaning.} Ladder without geometry is incomplete \cite{harrington2001}.

\paragraph{Problem (IoT rank).}
\textbf{Problem.} IoT vendor claims eight MIMO streams at sub-wavelength size. Audit.
\textbf{Step 1.} Compute \(k_{0}D_{s}\) \cite{jackson1999,harrington2001}.
\textbf{Step 2.} Apply Theorem~\ref{thm:ch4.geometry-rank-origin} \cite{stratton1941,hansen1988}.
\textbf{Step 3.} Reject if \(d_{\mathrm{eff}}\ll8\) \cite{devaney2004,collin2001}.
\textbf{Physical meaning.} Sub-wavelength rank is dipole-class \cite{harrington2001,stratton1941}.

\begin{corollary}[Listing excess without rank gain]
\label{cor:ch4.listing-excess-no-rank}
\(|\mathbb{C}^{I}|\gg d_{\mathrm{eff}}(\varepsilon)\) on symmetric charts is expected; enlarging the chart without enlarging \(\Omega_{s}\) or \(k_{0}\)
does not raise export rank \cite{hansen1988,harrington2001,stratton1941,devaney2004}.
\end{corollary}

\paragraph{Material class qualifier.}
\(\mathcal{J}_{\mathrm{real}}\) class---metal, dielectric, or active---enters Eq.~\eqref{eq:ch4.hardware-rank-map} through constitutive parameters that alter
\(\boldsymbol{\Gamma}_\omega\) without changing the monotonicity conclusions on fixed class \cite{stratton1941,harrington2001,hansen1988,jackson1999,devaney2004,collin2001}.
Rank memos must declare material class alongside \((\Omega_{s},k_{0})\) \cite{harrington2001,stratton1941}.

\paragraph{Problem (material class change).}
\textbf{Problem.} Switch from metal to dielectric loading on fixed outline. Rank change?
\textbf{Step 1.} Declare \(\mathcal{J}_{\mathrm{real}}\) class change \cite{stratton1941,harrington2001}.
\textbf{Step 2.} Recompute \(\boldsymbol{\Gamma}_\omega\) and \(d_{\mathrm{eff}}\) \cite{hansen1988,devaney2004}.
\textbf{Step 3.} Compare to Theorem~\ref{thm:ch4.geometry-rank-origin} baseline \cite{jackson1999,collin2001}.
\textbf{Physical meaning.} Material alters coupling, not listing \cite{harrington2001,stratton1941}.

\paragraph{Validity.}
Theorem~\ref{thm:ch4.geometry-rank-origin} assumes linear passive exterior on \(\mathcal{J}_{\mathrm{real}}\) \cite{stratton1941}. Active or nonlinear materials
enlarge \(\mathcal{J}_{\mathrm{real}}\) class and require separate range declarations \cite{harrington2001,devaney2004}. Dispersion modifies prefactors in
Eq.~\eqref{eq:ch4.joint-scaling-law} without changing the monotonicity conclusion on fixed material class \cite{jackson1999,hansen1988,collin2001}.

\subsection{Limits of phase-only engineering}
\label{sec:ch492}
% TOC tag: [HOME]

Phase-only control explores a submanifold of \(\mathcal{J}_{\mathrm{real}}\) at fixed amplitude support \cite{stratton1941,harrington2001,devaney2004}.
On fixed \((\Omega_{s},k_{0})\), phase rotations retune \(\mathbf{c}\) within \(\operatorname{ran}\boldsymbol{\Gamma}_\omega\) without enlarging
\(d_{\mathrm{eff}}(\varepsilon)\), \(\ell_{\mathrm{supp}}(\varepsilon,k_{0}a)\), or \(N_{\mathrm{mode}}(\varepsilon)\) \cite{hansen1988,jackson1999}.
Representation is coordinate freedom on the phase torus; realizability rank is geometry-imposed on \(\boldsymbol{\Gamma}_\omega\) \cite{stratton1941,devaney2004}.

\paragraph{Assumptions.}
Presuppose Proposition~\ref{prop:ch4.tangent-dimension}, Theorem~\ref{thm:ch4.geometry-rank-origin}, and Eq.~\eqref{eq:ch4.high-ell-decay}
\cite{stratton1941,harrington2001,devaney2004,hansen1988}. Passive linear exterior; fixed \(\mathcal{L}_{2}\) envelope on \(\Omega_{s}\); tolerance
\(\varepsilon\in(0,1)\) in aligned flux gauge \cite{jackson1999}. Phase-only designs vary \(\phi_{n}\) on fixed \(|\widehat{J}_{n}|\) portraits
\cite{collin2001,harrington2001}.

\begin{proposition}[Phase-only rank ceiling]
\label{prop:ch4.phase-only-ceiling}
Phase-only optimization on fixed \(\Omega_{s}\) cannot increase \(d_{\mathrm{eff}}(\varepsilon)\), \(\ell_{\mathrm{supp}}\), or
\(N_{\mathrm{mode}}\) \cite{stratton1941,jackson1999,hansen1988,harrington2001}. It rotates \(\mathbf{c}\) within the same export
subspace of dimension \(\le d_{\mathrm{eff}}(\varepsilon)\) \cite{devaney2004}.
\end{proposition}
\begin{proof}
Theorem~\ref{thm:ch4.geometry-rank-origin} and singular-value invariance under unitary phase gauge on fixed support \cite{stratton1941,harrington2001}.
\(\ell_{\mathrm{supp}}\) and \(N_{\mathrm{mode}}\) are properties of \(\boldsymbol{\Gamma}_\omega\), not of phase arguments on fixed amplitude
\cite{hansen1988,jackson1999,devaney2004}.
\end{proof}

Proposition~\ref{prop:ch4.phase-only-ceiling} closes the phase-only marketing loop: beam steering and OAM labeling by phase plates do not manufacture
high-\(\ell\) tails forbidden by Eq.~\eqref{eq:ch4.high-ell-decay} on fixed \(a\) \cite{hansen1988,jackson1999,harrington2001}. Amplitude support must
enlarge \(\mathcal{L}_{2}\) or \(k_{0}a\) must rise to raise \(\ell_{\mathrm{supp}}\) \cite{stratton1941,devaney2004}.

\paragraph{Phase-only tangent space.}
Phase-only designs explore a tangent space of dimension
\begin{equation}
\label{eq:ch4.phase-tangent-bound}
\dim\mathcal{T}_{\phi}\le r_{\mathrm{rad}}(\varepsilon),\qquad
\Delta d_{\mathrm{eff}}=0\ \text{under phase-only retuning},
\end{equation}
with \(\dim\mathcal{T}_{\phi}\ll d_{\mathrm{eff}}\) typical on fixed amplitude support per Proposition~\ref{prop:ch4.tangent-dimension}
\cite{stratton1941,harrington2001,devaney2004,hansen1988}. Mathematically, Eq.~\eqref{eq:ch4.phase-tangent-bound} bounds the local exploration manifold.
Physically, it counts independent phase knobs that retune far-zone direction without creating new export channels \cite{collin2001,jackson1999}.

\begin{theorem}[Phase manifold rank invariance]
\label{thm:ch4.phase-manifold-rank}
For \(\Jimp(\phi)\in\mathcal{J}_{\mathrm{real}}\) with fixed \(|\Jimp|\) support on \(\Omega_{s}\),
\(\boldsymbol{\Gamma}_\omega[\Jimp(\phi)]\) lies in a fixed \(d_{\mathrm{eff}}\)-dimensional subspace of \(\mathbb{C}^{I}\) for all admissible \(\phi\)
\cite{stratton1941,harrington2001,devaney2004,hansen1988}.
\end{theorem}
\begin{proof}
Phase variations are unitary rotations on coefficient charts restricted to fixed amplitude envelopes \cite{harrington2001,stratton1941}. Export singular
values of \(\boldsymbol{\Gamma}_\omega\) depend on support and coupling, not on phase arguments alone \cite{hansen1988,devaney2004,jackson1999}.
\end{proof}

\paragraph{Amplitude--phase factorization.}
Phase-only retuning admits
\begin{equation}
\label{eq:ch4.phase-amplitude-factor}
\mathbf{c}(\phi)=\mathbf{D}(\phi)\mathbf{c}_{0},\qquad
\mathbf{D}(\phi)\ \text{unitary on }\operatorname{span}\{\mathbf{c}_{0}\},\quad
\ell_{\mathrm{supp}}(\phi)=\ell_{\mathrm{supp}}(0),
\end{equation}
fixing \(\|\widehat{c}_{\ell m}^{(\sigma)}\|\) sector envelopes while rotating phases within \(\dim\mathcal{T}_{\phi}\)
\cite{stratton1941,hansen1988,harrington2001,jackson1999}. High-\(\ell\) tails cannot be created without amplitude support on enlarged \(\mathcal{L}_{2}\)
\cite{devaney2004}.

\begin{lemma}[Supported-order invariance under phase retune]
\label{lem:ch4.phase-retune-invariant}
\(\ell_{\mathrm{supp}}(\varepsilon,k_{0}a)\) and \(N_{\mathrm{mode}}(\varepsilon)\) are invariant under phase-only optimization on fixed \((\Omega_{s},k_{0})\)
\cite{hansen1988,stratton1941,harrington2001,devaney2004,jackson1999}.
\end{lemma}

\begin{figure}[t]
\centering
\ChOneFigBlock{%
\begin{tikzpicture}[ch1fig, node distance=7mm and 9mm]
  \node[ch1box] (phi) {phase \(\phi\)};
  \node[ch1box, right=11mm of phi] (tphi) {$\mathcal{T}_{\phi}$};
  \node[ch1box, right=11mm of tphi] (c) {$\mathbf{c}$};
  \draw[ch1flow] (phi) -- (tphi) -- (c);
\end{tikzpicture}%
}
\caption{Phase-only control explores \(\mathcal{T}_{\phi}\); export rank \(d_{\mathrm{eff}}\) is fixed by geometry.}
\label{fig:ch4.phase-tangent-manifold}
\end{figure}

Figure~\ref{fig:ch4.phase-tangent-manifold} is the phase-only honesty diagram: optimizers move on a low-dimensional torus inside a fixed-rank export cone
\cite{stratton1941,harrington2001,devaney2004}.

\begin{figure}[t]
\centering
\ChOneFigBlock{%
\begin{tikzpicture}[ch1fig, node distance=7mm and 9mm]
  \node[ch1box] (ell) {$\ell_{\mathrm{supp}}$};
  \node[ch1box, right=11mm of ell] (phase) {phase-only};
  \node[ch1box, right=11mm of phase] (flat) {unchanged};
  \draw[ch1flow] (ell) -- (phase) -- (flat);
\end{tikzpicture}%
}
\caption{Supported angular order is invariant under phase-only retuning on fixed aperture.}
\label{fig:ch4.phase-ell-ceiling}
\end{figure}

\paragraph{Worked illustration (reflectarray).}
A reflectarray that steers a beam without enlarging aperture rotates \(\mathbf{c}\) within \(\ell=1\)--\(\ell_{\mathrm{supp}}\) content;
it does not manufacture high-\(\ell\) tails forbidden by Eq.~\eqref{eq:ch4.high-ell-decay} \cite{hansen1988,jackson1999,harrington2001,collin2001}.

\paragraph{Worked illustration (phase-only gradient).}
Gradient ascent on \(\dim\mathcal{T}_{\phi}=2\) parameters on a rank-\(15\) aperture achieves excellent beam alignment while \(d_{\mathrm{eff}}=15\) unchanged
\cite{harrington2001,devaney2004,stratton1941,hansen1988}.

\paragraph{Worked illustration (OAM phase plate).}
Phase plate claims new OAM order \(m_{\mathrm{LG}}=5\) on fixed disk: \(\ell_{\mathrm{supp}}\) before and after identical per
Lemma~\ref{lem:ch4.phase-retune-invariant} \cite{hansen1988,harrington2001,jackson1999}.

\paragraph{Worked illustration (metasurface).}
Fixed-area metasurface retunes \(\theta_{0}\) and \(\phi_{0}\) with \(\dim\mathcal{T}_{\phi}\approx2\); tail exponent Eq.~\eqref{eq:ch4.high-ell-decay}
unchanged \cite{stratton1941,collin2001,devaney2004}.

\paragraph{Worked numerics (reflectarray).}
\(d_{\mathrm{eff}}=15\) reflectarray with \(\dim\mathcal{T}_{\phi}=2\) steers beam in \(\theta\)--\(\phi\) only \cite{hansen1988,collin2001,harrington2001}.

\paragraph{Worked numerics (phase-only OAM).}
Claimed \(\Delta\ell_{\mathrm{supp}}=+3\) after phase-only synthesis on \(k_{0}a=1.5\): measured \(\ell_{\mathrm{supp}}\) unchanged at \(4\)
\cite{hansen1988,jackson1999,stratton1941,devaney2004}.

\paragraph{Problem (phase-only OAM claim).}
\textbf{Problem.} Phase plate claims new OAM order without size change. Audit.
\textbf{Step 1.} Test \(\ell_{\mathrm{supp}}\) before and after \cite{hansen1988,stratton1941}.
\textbf{Step 2.} Test tail bound Eq.~\eqref{eq:ch4.high-ell-decay} \cite{jackson1999,devaney2004}.
\textbf{Step 3.} Reject if only phase changed on fixed \(\mathcal{L}_{2}\) \cite{harrington2001}.
\textbf{Physical meaning.} Phase explores a fixed-rank manifold \cite{stratton1941}.

\paragraph{Problem (phase-only rank).}
\textbf{Problem.} Can phase-only synthesis raise OAM order on fixed disk?
\textbf{Step 1.} Test Lemma~\ref{lem:ch4.phase-retune-invariant} \cite{hansen1988,harrington2001}.
\textbf{Step 2.} Test Eq.~\eqref{eq:ch4.high-ell-decay} \cite{jackson1999,devaney2004}.
\textbf{Step 3.} Reject order increase \cite{stratton1941}.
\textbf{Physical meaning.} Phase rotates within fixed rank \cite{harrington2001,collin2001}.

\paragraph{Problem (reflectarray audit).}
\textbf{Problem.} Reflectarray claims \(d_{\mathrm{eff}}\) doubling via phase tiles. Valid?
\textbf{Step 1.} Apply Theorem~\ref{thm:ch4.phase-manifold-rank} \cite{stratton1941,harrington2001}.
\textbf{Step 2.} Measure SVD rank before/after \cite{hansen1988,devaney2004}.
\textbf{Step 3.} Reject rank-doubling claim \cite{jackson1999}.
\textbf{Physical meaning.} Phase retunes, does not create channels \cite{collin2001}.

\paragraph{Problem (tangent dimension).}
\textbf{Problem.} How many independent phase controls on rank-\(12\) aperture?
\textbf{Step 1.} Bound \(\dim\mathcal{T}_{\phi}\) via Eq.~\eqref{eq:ch4.phase-tangent-bound} \cite{harrington2001,stratton1941}.
\textbf{Step 2.} Compare to Proposition~\ref{prop:ch4.tangent-dimension} \cite{devaney2004,hansen1988}.
\textbf{Step 3.} Report \(\dim\mathcal{T}_{\phi}\ll d_{\mathrm{eff}}\) typically \cite{jackson1999}.
\textbf{Physical meaning.} Phase knobs are fewer than export directions \cite{stratton1941}.

\paragraph{Problem (amplitude requirement).}
\textbf{Problem.} High-\(\ell\) tail needed; phase-only synthesis fails. Why?
\textbf{Step 1.} Test Eq.~\eqref{eq:ch4.phase-amplitude-factor} \cite{stratton1941,hansen1988}.
\textbf{Step 2.} Verify amplitude envelope fixed \cite{harrington2001,devaney2004}.
\textbf{Step 3.} Require \(\mathcal{L}_{2}\) enlargement or higher \(k_{0}a\) \cite{jackson1999,collin2001}.
\textbf{Physical meaning.} Tails need amplitude support \cite{hansen1988}.

\paragraph{Problem (metasurface MIMO).}
\textbf{Problem.} Metasurface adds phase pixels; vendor claims new streams. Audit.
\textbf{Step 1.} Apply Proposition~\ref{prop:ch4.phase-only-ceiling} \cite{stratton1941,harrington2001}.
\textbf{Step 2.} Measure \(d_{\mathrm{eff}}\) and \(N_{\mathrm{mode}}\) \cite{hansen1988,devaney2004}.
\textbf{Step 3.} Reject stream inflation \cite{jackson1999,collin2001}.
\textbf{Physical meaning.} Pixels retune within fixed rank \cite{harrington2001}.

\paragraph{Holographic phase sheets.}
Holographic metasurfaces impose \(\mathbf{c}(\phi)=\mathbf{D}(\phi)\mathbf{c}_{0}\) with \(\mathbf{D}(\phi)\) unitary on a fixed amplitude portrait
\cite{stratton1941,harrington2001,collin2001}. Steering and focusing are retunings within \(\dim\mathcal{T}_{\phi}\); they do not alter tail exponents or
\(\ell_{\mathrm{supp}}\) on fixed \((\Omega_{s},k_{0})\) per Lemma~\ref{lem:ch4.phase-retune-invariant} \cite{hansen1988,jackson1999,devaney2004}.

\paragraph{Phase--amplitude coupling threshold.}
Raising \(\ell_{\mathrm{supp}}\) requires amplitude degrees of freedom that enlarge \(\mathcal{L}_{2}\) or couple new \(\ell\) sectors through feed topology
\cite{harrington2001,stratton1941,devaney2004}. Phase-only paths satisfy \(\partial\|\widehat{c}_{\ell m}^{(\sigma)}\|/\partial\phi_{n}=0\) on fixed support
\cite{hansen1988,jackson1999,collin2001}.

\begin{corollary}[Stream count invariance under phase-only MIMO]
\label{cor:ch4.phase-mimo-invariant}
Deployable stream count \(N_{\mathrm{stream}}\) is invariant under phase-only retuning on fixed \((\Omega_{s},k_{0},\varepsilon)\)
\cite{harrington2001,devaney2004,hansen1988,stratton1941,jackson1999,collin2001}.
\end{corollary}

\paragraph{Worked illustration (holographic beam).}
Holographic sheet steers main beam \(15^{\circ}\) with \(\dim\mathcal{T}_{\phi}=2\); \(\ell_{\mathrm{supp}}\) and \(d_{\mathrm{eff}}\) unchanged
\cite{stratton1941,collin2001,hansen1988,harrington2001,devaney2004}.

\paragraph{Worked illustration (transmitarray).}
Transmitarray phase tiles rotate \(\mathbf{c}\) within rank-\(18\) subspace; vendor claim of \(\ell=10\) OAM rejected by Lemma~\ref{lem:ch4.phase-retune-invariant}
\cite{hansen1988,jackson1999,stratton1941,devaney2004}.

\paragraph{Worked numerics (tangent versus rank).}
Rank-\(20\) aperture, \(\dim\mathcal{T}_{\phi}=3\): \(97\%\) of export directions are unreachable by phase-only gradients per
Eq.~\eqref{eq:ch4.phase-tangent-bound} \cite{harrington2001,stratton1941,hansen1988,devaney2004,collin2001}.

\paragraph{Problem (holographic OAM).}
\textbf{Problem.} Hologram claims \(\ell=7\) OAM on fixed \(k_{0}a=1.4\). Audit.
\textbf{Step 1.} Test Lemma~\ref{lem:ch4.phase-retune-invariant} \cite{hansen1988,harrington2001}.
\textbf{Step 2.} Test Eq.~\eqref{eq:ch4.high-ell-decay} \cite{jackson1999,devaney2004}.
\textbf{Step 3.} Reject without amplitude enlargement \cite{stratton1941,collin2001}.
\textbf{Physical meaning.} Holograms retune, not create \(\ell\) \cite{harrington2001}.

\paragraph{Problem (phase MIMO streams).}
\textbf{Problem.} Phase-only beamforming claims \(32\) new streams. Valid?
\textbf{Step 1.} Apply Corollary~\ref{cor:ch4.phase-mimo-invariant} \cite{harrington2001,devaney2004}.
\textbf{Step 2.} Measure \(d_{\mathrm{eff}}\) \cite{hansen1988,stratton1941}.
\textbf{Step 3.} Reject stream inflation \cite{jackson1999,collin2001}.
\textbf{Physical meaning.} Phase cannot raise rank \cite{stratton1941}.

\paragraph{Problem (phase-only versus full synthesis).}
\textbf{Problem.} Full synthesis raises \(\ell_{\mathrm{supp}}\); phase-only does not. Explain.
\textbf{Step 1.} Compare exploration sets \(\mathcal{T}_{\phi}\) versus \(\mathcal{T}_{\Jimp_{0}}\mathcal{J}_{\mathrm{real}}\) \cite{harrington2001,stratton1941}.
\textbf{Step 2.} Apply Proposition~\ref{prop:ch4.tangent-dimension} \cite{devaney2004,hansen1988}.
\textbf{Step 3.} Attribute \(\ell_{\mathrm{supp}}\) gain to amplitude/coupling DOF \cite{jackson1999}.
\textbf{Physical meaning.} Rank needs more than phase \cite{stratton1941,collin2001}.

\paragraph{Phase trajectory on export sphere.}
Phase-only optimization traces a curve on the complex export sphere of dimension \(d_{\mathrm{eff}}-1\) embedded in \(\mathbb{C}^{I}\)
\cite{stratton1941,harrington2001,hansen1988,devaney2004}. The trajectory cannot exit \(\operatorname{ran}\boldsymbol{\Gamma}_\omega\) without amplitude or
support enlargement \cite{jackson1999,collin2001}.

\paragraph{Worked illustration (phased array).}
\(64\)-element phased array with uniform amplitude: \(\dim\mathcal{T}_{\phi}=64\) but \(d_{\mathrm{eff}}=18\) on \(k_{0}W=3\): phase explores a submanifold
of full rank \cite{harrington2001,collin2001,hansen1988,stratton1941,devaney2004,jackson1999}.

\paragraph{Worked numerics (phase-only coverage).}
Fraction of \(\mathbb{C}^{I}\) reachable by phase-only retuning is at most \(\dim\mathcal{T}_{\phi}/d_{\mathrm{eff}}\), typically \(\ll1\) on fixed amplitude
support \cite{harrington2001,stratton1941,hansen1988,devaney2004,collin2001}.

\paragraph{Problem (uniform amplitude array).}
\textbf{Problem.} \(64\) phases, \(d_{\mathrm{eff}}=18\). Independent phase controls?
\textbf{Step 1.} Bound \(\dim\mathcal{T}_{\phi}\le64\) \cite{harrington2001,stratton1941}.
\textbf{Step 2.} Note export still rank-\(18\) \cite{hansen1988,devaney2004}.
\textbf{Step 3.} Phase cannot exceed \(d_{\mathrm{eff}}\) independent channels \cite{jackson1999,collin2001}.
\textbf{Physical meaning.} Phases \(\neq\) channels \cite{harrington2001}.

\paragraph{Problem (transmitarray OAM).}
\textbf{Problem.} Transmitarray adds spiral phase; claims new \(\ell\). Audit.
\textbf{Step 1.} Test Eq.~\eqref{eq:ch4.phase-amplitude-factor} \cite{stratton1941,hansen1988}.
\textbf{Step 2.} Measure \(\ell_{\mathrm{supp}}\) \cite{harrington2001,jackson1999}.
\textbf{Step 3.} Reject if amplitude support fixed \cite{devaney2004,collin2001}.
\textbf{Physical meaning.} Spiral phase is retuning \cite{stratton1941}.

\paragraph{Problem (phase coverage fraction).}
\textbf{Problem.} Estimate reachable export fraction under phase-only control.
\textbf{Step 1.} Compute \(\dim\mathcal{T}_{\phi}\) \cite{harrington2001,stratton1941}.
\textbf{Step 2.} Divide by \(d_{\mathrm{eff}}\) \cite{hansen1988,devaney2004}.
\textbf{Step 3.} Report fraction \(\ll1\) typically \cite{jackson1999,collin2001}.
\textbf{Physical meaning.} Most export directions need amplitude DOF \cite{harrington2001}.

\paragraph{Philosophical closure (retuning versus creation).}
Phase-only engineering is powerful within \(\operatorname{ran}\boldsymbol{\Gamma}_\omega\) but powerless against geometry-imposed ceilings on
\(d_{\mathrm{eff}}\), \(\ell_{\mathrm{supp}}\), and \(N_{\mathrm{mode}}\) \cite{stratton1941,harrington2001,jackson1999,hansen1988,devaney2004}.
Claims that phase alone eliminates spectral truncation confuse coordinate retuning on \(\mathcal{T}_{\phi}\) with rank creation on \(\boldsymbol{\Gamma}_\omega\)
\cite{collin2001,stratton1941}.

\paragraph{Reconfigurable intelligent surfaces.}
RIS phase tiles implement \(\mathbf{D}(\phi)\) on fixed element amplitudes; \(d_{\mathrm{eff}}\) and \(\ell_{\mathrm{supp}}\) are set by tile footprint and \(k_{0}\),
not by tile count alone \cite{stratton1941,harrington2001,hansen1988,collin2001,jackson1999,devaney2004}. Additional tiles sample the aperture; rank rises only when
electrical aperture grows \cite{harrington2001,stratton1941}.

\paragraph{Worked illustration (RIS).}
\(256\) RIS tiles on \(\lambda/2\) aperture: \(\dim\mathcal{T}_{\phi}=256\) but \(d_{\mathrm{eff}}=11\); phase retuning does not create twelfth channel
\cite{harrington2001,collin2001,hansen1988,stratton1941,devaney2004,jackson1999}.

\paragraph{Problem (RIS stream claim).}
\textbf{Problem.} RIS claims rank doubling via phase reconfiguration. Audit.
\textbf{Step 1.} Measure \(d_{\mathrm{eff}}\) before/after \(\phi\) sweep \cite{harrington2001,hansen1988}.
\textbf{Step 2.} Apply Proposition~\ref{prop:ch4.phase-only-ceiling} \cite{stratton1941,devaney2004}.
\textbf{Step 3.} Reject rank-doubling \cite{jackson1999,collin2001}.
\textbf{Physical meaning.} RIS retunes within fixed rank \cite{harrington2001}.

\paragraph{Problem (amplitude metasurface).}
\textbf{Problem.} Amplitude and phase both varied. Can \(\ell_{\mathrm{supp}}\) rise?
\textbf{Step 1.} Test whether \(\mathcal{L}_{2}\) enlarged \cite{stratton1941,harrington2001}.
\textbf{Step 2.} If support fixed, apply Lemma~\ref{lem:ch4.phase-retune-invariant} \cite{hansen1988,devaney2004}.
\textbf{Step 3.} Attribute \(\ell_{\mathrm{supp}}\) gain to amplitude/support change only \cite{jackson1999,collin2001}.
\textbf{Physical meaning.} Amplitude DOF can raise tails \cite{harrington2001,stratton1941}.

\begin{corollary}[Phase-only cannot raise tail exponent]
\label{cor:ch4.phase-tail-invariant}
Tail decay exponent in Eq.~\eqref{eq:ch4.high-ell-decay} is invariant under phase-only retuning on fixed \((\Omega_{s},k_{0})\)
\cite{hansen1988,jackson1999,stratton1941,harrington2001,devaney2004}.
\end{corollary}

\paragraph{Near-field phase hotspots.}
Phase-only synthesis may create near-field hotspots without raising far-zone \(d_{\mathrm{eff}}\) per Theorem~\ref{thm:ch4.reactive-dominance} in
Section~\ref{sec:ch472} \cite{stratton1941,harrington2001,hansen1988,devaney2004,jackson1999,collin2001}. Hotspots certify reactive storage, not export rank
\cite{harrington2001,stratton1941}.

\paragraph{Worked numerics (hotspot versus rank).}
Near-field enhancement factor \(10\times\) with \(d_{\mathrm{eff}}\) unchanged on fixed \(\Omega_{s}\): reactive dominance without rank gain
\cite{harrington2001,jackson1999,stratton1941,hansen1988,devaney2004,collin2001}.

\paragraph{Problem (near-field hotspot).}
\textbf{Problem.} Phase synthesis creates near hotspot; vendor claims rank gain. Audit.
\textbf{Step 1.} Measure \(d_{\mathrm{eff}}\) on \(\mathcal{F}_{\mathrm{rad}}\) \cite{harrington2001,hansen1988}.
\textbf{Step 2.} Apply reactive filter from Section~\ref{sec:ch472} \cite{stratton1941,devaney2004}.
\textbf{Step 3.} Reject rank claim \cite{jackson1999,collin2001}.
\textbf{Physical meaning.} Hotspots are not channels \cite{harrington2001,stratton1941}.

\paragraph{Validity.}
Proposition~\ref{prop:ch4.phase-only-ceiling} assumes passive linear exterior and fixed amplitude support \cite{stratton1941}. Active amplification can enlarge
effective \(\mathcal{J}_{\mathrm{real}}\) class; nonlinear phase elements require separate range analysis \cite{harrington2001,devaney2004}. Loss modifies
phase exploration without raising \(\ell_{\mathrm{supp}}\) on fixed \(\Omega_{s}\) \cite{hansen1988,jackson1999,collin2001}.

\subsection{Observable manifestations of truncation}
\label{sec:ch493}
% TOC tag: [USE]

Truncation is not invisible in laboratory data: it appears as smoothed main beams, elevated far-zone shoulders, saturated MIMO capacity curves, and
channel counts that knee below element cardinality \cite{stratton1941,harrington2001,hansen1988,devaney2004}. Section~\ref{sec:ch01.5} established that
measurement functionals report projections of fields, not full coefficient charts; here those projections inherit the finite-support restriction
\(\mathcal{B}_{\mathrm{fs}}\) and rank ceiling \(d_{\mathrm{eff}}(\varepsilon)\) proved on \(\boldsymbol{\Gamma}_\omega\) \cite{jackson1999}.
Observable signatures are realizability certificates read in watts and decibels, not operator jargon \cite{harrington2001,devaney2004}.

\paragraph{Assumptions.}
Presuppose Definition~\ref{def:ch4.finite-support-restriction}, Theorem~\ref{thm:ch4.finite-support-audit}, and
Eq.~\eqref{eq:ch3.truncation-remainder-energy} \cite{stratton1941,hansen1988,harrington2001}. Narrowband phasors; aligned flux gauge; tolerance
\(\varepsilon\in(0,1)\) \cite{jackson1999,devaney2004}. Pattern and capacity observables are linear functionals on \(\mathcal{F}_{\mathrm{rad}}\)
\cite{collin2001,harrington2001}.

\paragraph{Observable signatures.}
Truncation manifests in pattern and channel domains as
\begin{equation}
\label{eq:ch4.truncation-observables}
\Delta_{\mathrm{pat}}
\sim
\|\mathbf{R}_{\ell_{\max}}\|_{\mathrm{rad}},
\qquad
\Delta_{\mathrm{chan}}
\sim
d_{\mathrm{eff}}^{-1},
\end{equation}
with \(\mathbf{R}_{\ell_{\max}}\) the discarded tail from Eq.~\eqref{eq:ch3.truncation-remainder-energy} \cite{hansen1988,stratton1941,harrington2001}.
Mathematically, Eq.~\eqref{eq:ch4.truncation-observables} links remainder energy to shoulder level and rank to channel granularity. Physically, elevated
far-zone sidelobes and capacity knees at fixed SNR are consistent signatures of \(\mathcal{A}_{\mathrm{fs}}>1\) or \(d_{\mathrm{eff}}\) saturation
\cite{devaney2004,jackson1999,collin2001}.

\begin{theorem}[Truncation signature theorem]
\label{thm:ch4.truncation-signature}
If \(\mathcal{A}_{\mathrm{fs}}[\mathbf{c}]>1\) or \(\ell_{\mathrm{listed}}>\ell_{\mathrm{supp}}(\varepsilon,k_{0}a)\), then measured patterns exhibit
shoulder artifacts \(\Delta_{\mathrm{pat}}\) bounded below by Eq.~\eqref{eq:ch4.shoulder-level} even under perfect calibration \cite{hansen1988,stratton1941,harrington2001,devaney2004}.
\end{theorem}
\begin{proof}
Theorem~\ref{thm:ch4.finite-support-audit} and Eq.~\eqref{eq:ch3.truncation-remainder-energy} convert discarded tail mass into far-zone power
\cite{hansen1988,stratton1941}. Unreachable high-\(\ell\) content cannot be absorbed by calibration on \(\mathcal{F}_{\mathrm{rad}}\) \cite{harrington2001,devaney2004}.
\end{proof}

Theorem~\ref{thm:ch4.truncation-signature} redirects shoulder diagnostics: when \(\mathcal{A}_{\mathrm{fs}}>1\), recalibration cannot remove artifacts;
aperture enlargement or \(\ell_{\max}\) reduction is required \cite{stratton1941,hansen1988,jackson1999}.

\paragraph{Shoulder level estimate.}
\begin{equation}
\label{eq:ch4.shoulder-level}
\Delta_{\mathrm{pat}}/\|\mathbf{c}\|_{\mathrm{rad}}\sim\sqrt{\mathcal{A}_{\mathrm{fs}}-1},\qquad
\mathcal{A}_{\mathrm{fs}}>1.
\end{equation}
Shoulder level scales with discarded tail mass \(\|\mathbf{R}_{\ell_{\max}}\|_{\mathrm{rad}}\) per Eq.~\eqref{eq:ch4.truncation-observables}
\cite{hansen1988,stratton1941,harrington2001,devaney2004}.

\paragraph{Channel capacity saturation.}
In linear MIMO models,
\begin{equation}
\label{eq:ch4.capacity-saturation}
C_{\mathrm{MIMO}}\le\sum_{n=1}^{d_{\mathrm{eff}}(\varepsilon)}\log\bigl(1+\mathrm{SNR}\,\sigma_{n}^{2}\bigr),
\end{equation}
so streams beyond \(d_{\mathrm{eff}}\) cannot raise capacity \cite{harrington2001,devaney2004,hansen1988,jackson1999,collin2001}. Mathematically,
Eq.~\eqref{eq:ch4.capacity-saturation} is the water-filling bound on export singular values. Physically, it is the capacity knee laboratories observe
when element count exceeds rank \cite{stratton1941,harrington2001}.

\begin{lemma}[Capacity knee location]
\label{lem:ch4.capacity-knee}
Capacity curves knee near stream count \(N_{\mathrm{stream}}\approx d_{\mathrm{eff}}(\varepsilon)\) at fixed SNR and \((\Omega_{s},k_{0})\)
\cite{harrington2001,devaney2004,hansen1988,stratton1941,jackson1999}.
\end{lemma}

\begin{figure}[t]
\centering
\ChOneFigBlock{%
\begin{tikzpicture}[ch1fig, node distance=7mm and 9mm]
  \node[ch1box] (afs) {$\mathcal{A}_{\mathrm{fs}}$};
  \node[ch1box, right=11mm of afs] (tail) {$\mathbf{R}_{\ell_{\max}}$};
  \node[ch1box, right=11mm of tail] (sh) {$\Delta_{\mathrm{pat}}$};
  \draw[ch1flow] (afs) -- (tail) -- (sh);
\end{tikzpicture}%
}
\caption{Finite-support excess drives tail remainder and observable shoulder artifacts.}
\label{fig:ch4.truncation-signatures}
\end{figure}

Figure~\ref{fig:ch4.truncation-signatures} chains the pattern diagnostic: \(\mathcal{A}_{\mathrm{fs}}\) audit precedes calibration blame \cite{hansen1988,harrington2001,stratton1941,devaney2004}.

\begin{figure}[t]
\centering
\ChOneFigBlock{%
\begin{tikzpicture}[ch1fig, node distance=7mm and 9mm]
  \node[ch1box] (nel) {$N_{\mathrm{el}}$};
  \node[ch1box, right=11mm of nel] (deff) {$d_{\mathrm{eff}}$};
  \node[ch1box, right=11mm of deff] (knee) {capacity knee};
  \draw[ch1flow] (nel) -- (deff);
  \draw[ch1flow, dashed] (deff) -- (knee);
\end{tikzpicture}%
}
\caption{Element count may exceed rank; capacity knees trace \(d_{\mathrm{eff}}\), not patch count.}
\label{fig:ch4.capacity-knee-map}
\end{figure}

\paragraph{Worked illustration (pattern shoulders).}
A target with plateau high-\(\ell\) tails produces shoulder artifacts when synthesized on \(B_{a}\): the artifact is a realizability signature, not a
calibration error per Theorem~\ref{thm:ch4.truncation-signature} \cite{hansen1988,jackson1999,harrington2001,devaney2004}.

\paragraph{Worked illustration (capacity knee).}
Adding streams beyond \(d_{\mathrm{eff}}=11\) on \(32\) patches does not raise \(C_{\mathrm{MIMO}}\) in the linear model per
Eq.~\eqref{eq:ch4.capacity-saturation} \cite{harrington2001,jackson1999,collin2001}.

\paragraph{Worked illustration (tail mass).}
\(\mathcal{A}_{\mathrm{fs}}=1.3\) with \(\|\mathbf{R}_{\ell_{\max}}\|_{\mathrm{rad}}/\|\mathbf{c}\|=0.22\): shoulders visible in anechoic scan
\cite{hansen1988,stratton1941,harrington2001}.

\paragraph{Worked illustration (OAM plateau).}
Target \(\ell=8\) plateau on \(k_{0}a=1.6\): \(\ell_{\mathrm{supp}}=5\) forces shoulder ring in measured pattern \cite{hansen1988,jackson1999,devaney2004}.

\paragraph{Worked numerics (shoulder diagnostic).}
\(\mathcal{A}_{\mathrm{fs}}=1.2\) predicts shoulder level \(\sim0.45\) relative to main beam in power via Eq.~\eqref{eq:ch4.shoulder-level}
\cite{hansen1988,harrington2001,stratton1941,devaney2004}.

\paragraph{Worked numerics (capacity).}
\(d_{\mathrm{eff}}(10^{-2})=9\), \(N_{\mathrm{el}}=64\): capacity knee at \(\sim9\)--\(11\) streams per Lemma~\ref{lem:ch4.capacity-knee}
\cite{harrington2001,devaney2004,collin2001,jackson1999}.

\paragraph{Problem (diagnose truncation in data).}
\textbf{Problem.} Measured pattern shows unexpected shoulders. Truncation or miscalibration?
\textbf{Step 1.} Compute \(\mathcal{A}_{\mathrm{fs}}[\mathbf{c}]\) \cite{hansen1988,stratton1941}.
\textbf{Step 2.} Test tail mass Eq.~\eqref{eq:ch4.cumulative-tail-mass} \cite{harrington2001,devaney2004}.
\textbf{Step 3.} If audits fail, attribute to geometry rank before recalibrating \cite{jackson1999}.
\textbf{Physical meaning.} Truncation has characteristic far-zone signatures \cite{hansen1988,stratton1941}.

\paragraph{Problem (capacity knee).}
\textbf{Problem.} Capacity curve knees at sixteen streams on \(32\) patches. Explain.
\textbf{Step 1.} Estimate \(d_{\mathrm{eff}}(10^{-2})\) \cite{hansen1988,harrington2001}.
\textbf{Step 2.} Compare to stream count \cite{devaney2004,collin2001}.
\textbf{Step 3.} Attribute knee to rank saturation per Lemma~\ref{lem:ch4.capacity-knee} \cite{stratton1941}.
\textbf{Physical meaning.} Capacity knees trace \(d_{\mathrm{eff}}\) \cite{harrington2001,jackson1999}.

\paragraph{Problem (shoulder after calibration).}
\textbf{Problem.} Shoulders persist after probe calibration. Next step?
\textbf{Step 1.} Apply Theorem~\ref{thm:ch4.truncation-signature} \cite{hansen1988,stratton1941}.
\textbf{Step 2.} Test \(\mathcal{A}_{\mathrm{fs}}\le1\) \cite{harrington2001,devaney2004}.
\textbf{Step 3.} Reduce \(\ell_{\mathrm{listed}}\) or enlarge aperture \cite{jackson1999,collin2001}.
\textbf{Physical meaning.} Shoulders are realizability, not metrology \cite{stratton1941}.

\paragraph{Problem (MIMO stream audit).}
\textbf{Problem.} Vendor claims \(48\) streams; panel \(k_{0}^{2}A_{a}=15\). Plausible?
\textbf{Step 1.} Bound \(N_{\mathrm{mode}}\) \cite{jackson1999,hansen1988}.
\textbf{Step 2.} Measure \(d_{\mathrm{eff}}\) \cite{harrington2001,devaney2004}.
\textbf{Step 3.} Compare to Eq.~\eqref{eq:ch4.capacity-saturation} \cite{collin2001,stratton1941}.
\textbf{Physical meaning.} Streams are rank-limited \cite{harrington2001}.

\paragraph{Problem (tail watts).}
\textbf{Problem.} Convert discarded \(\ell\) sectors to shoulder watts.
\textbf{Step 1.} Apply Eq.~\eqref{eq:ch3.truncation-remainder-energy} \cite{hansen1988,stratton1941}.
\textbf{Step 2.} Compute \(\|\mathbf{R}_{\ell_{\max}}\|_{\mathrm{rad}}\) \cite{harrington2001,devaney2004}.
\textbf{Step 3.} Predict \(\Delta_{\mathrm{pat}}\) via Eq.~\eqref{eq:ch4.shoulder-level} \cite{jackson1999}.
\textbf{Physical meaning.} Tail energy is measurable \cite{stratton1941,collin2001}.

\paragraph{Problem (anechoic ring).}
\textbf{Problem.} OAM ring in chamber data on small aperture. Interpretation?
\textbf{Step 1.} Read \(\ell_{\mathrm{supp}}(\varepsilon,k_{0}a)\) \cite{hansen1988,harrington2001}.
\textbf{Step 2.} Test Theorem~\ref{thm:ch4.truncation-signature} \cite{stratton1941,devaney2004}.
\textbf{Step 3.} Attribute ring to truncation artifact if \(\mathcal{A}_{\mathrm{fs}}>1\) \cite{jackson1999}.
\textbf{Physical meaning.} Rings can be rank signatures \cite{harrington2001}.

\paragraph{Sidelobe envelope from tail mass.}
For \(\mathcal{A}_{\mathrm{fs}}>1\), sidelobe envelope level satisfies
\(\Delta_{\mathrm{pat}}/\|\mathbf{c}\|_{\mathrm{rad}}\gtrsim\|\mathbf{R}_{\ell_{\max}}\|_{\mathrm{rad}}/\|\mathbf{c}\|_{\mathrm{rad}}\) per
Eqs.~\eqref{eq:ch4.truncation-observables} and \eqref{eq:ch4.shoulder-level} \cite{hansen1988,stratton1941,harrington2001,devaney2004,jackson1999}.
Calibration removes probe gain errors, not tail energy \cite{harrington2001,collin2001}.

\paragraph{Mutual coupling overlay.}
Mutual coupling reshuffles \(\sigma_{n}\) without removing capacity knees at \(N_{\mathrm{stream}}\approx d_{\mathrm{eff}}\) on fixed \((\Omega_{s},k_{0})\)
\cite{harrington2001,devaney2004,hansen1988,stratton1941}. Coupling-aware calibration improves conditioning; it does not manufacture export channels
\cite{jackson1999,collin2001}.

\begin{corollary}[Observable rank saturation]
\label{cor:ch4.observable-rank-sat}
Any observable linear functional on \(\mathcal{F}_{\mathrm{rad}}\) cannot resolve more than \(d_{\mathrm{eff}}(\varepsilon)\) independent export directions
above \(\varepsilon\sigma_{1}\) on declared \((\Omega_{s},k_{0})\) \cite{stratton1941,harrington2001,hansen1988,devaney2004,jackson1999}.
\end{corollary}

\paragraph{Worked illustration (mutual coupling).}
Coupling matrix conditioned; capacity knee remains at \(d_{\mathrm{eff}}=10\) on \(24\) patches \cite{harrington2001,devaney2004,hansen1988,collin2001}.

\paragraph{Worked illustration (grating lobes versus truncation).}
Grating lobes from periodic arrays are distinct from truncation shoulders; \(\mathcal{A}_{\mathrm{fs}}\) audit separates realizability shoulders from lattice artifacts
\cite{hansen1988,stratton1941,harrington2001,jackson1999,devaney2004}.

\paragraph{Worked numerics (tail-to-shoulder).}
\(\|\mathbf{R}_{\ell_{\max}}\|_{\mathrm{rad}}/\|\mathbf{c}\|=0.18\) predicts \(\Delta_{\mathrm{pat}}\approx0.42\) in power via Eq.~\eqref{eq:ch4.shoulder-level}
\cite{hansen1988,harrington2001,stratton1941,devaney2004,collin2001}.

\paragraph{Problem (coupling versus rank).}
\textbf{Problem.} Coupling calibration raises capacity slightly but knee unchanged. Explain.
\textbf{Step 1.} Apply Lemma~\ref{lem:ch4.capacity-knee} \cite{harrington2001,devaney2004}.
\textbf{Step 2.} Note \(\sigma_{n}\) reordering only \cite{hansen1988,stratton1941}.
\textbf{Step 3.} Knee tracks \(d_{\mathrm{eff}}\) \cite{jackson1999,collin2001}.
\textbf{Physical meaning.} Coupling conditions, does not create rank \cite{harrington2001}.

\paragraph{Problem (grating versus shoulder).}
\textbf{Problem.} Unexpected far sidelobe on periodic array. Truncation or grating?
\textbf{Step 1.} Compute \(\mathcal{A}_{\mathrm{fs}}[\mathbf{c}]\) \cite{hansen1988,stratton1941}.
\textbf{Step 2.} Test lattice period versus \(\lambda\) \cite{harrington2001,collin2001}.
\textbf{Step 3.} Attribute per dominant mechanism \cite{devaney2004,jackson1999}.
\textbf{Physical meaning.} Two distinct far-zone mechanisms \cite{stratton1941}.

\paragraph{Problem (SNR sweep).}
\textbf{Problem.} Capacity rises with SNR but stream count at knee fixed. Explain.
\textbf{Step 1.} Apply Eq.~\eqref{eq:ch4.capacity-saturation} \cite{harrington2001,devaney2004}.
\textbf{Step 2.} Note \(d_{\mathrm{eff}}\) SNR-independent at fixed \((\Omega_{s},k_{0})\) \cite{hansen1988,stratton1941}.
\textbf{Step 3.} Knee location unchanged; slope rises \cite{jackson1999,collin2001}.
\textbf{Physical meaning.} Rank sets knees; SNR sets slope \cite{harrington2001}.

\paragraph{Pattern moment saturation.}
Low-order pattern moments saturate when \(d_{\mathrm{eff}}\) is exhausted: adding higher moments without rank growth produces shoulders rather than main-beam
sharpening \cite{hansen1988,stratton1941,harrington2001,devaney2004,jackson1999}. Moment fits must respect \(\ell_{\mathrm{supp}}\) ceiling
\cite{harrington2001,collin2001}.

\paragraph{Worked illustration (moment fit).}
Fit through \(\ell=12\) on \(k_{0}a=1.4\) with \(\ell_{\mathrm{supp}}=5\): shoulders at \(\theta\approx\pm40^{\circ}\) per Theorem~\ref{thm:ch4.truncation-signature}
\cite{hansen1988,harrington2001,stratton1941,devaney2004,jackson1999}.

\paragraph{Worked numerics (stream knee).}
\(N_{\mathrm{el}}=48\), \(d_{\mathrm{eff}}=14\): Shannon capacity knee at stream \(14\)--\(16\); additional streams add \(\le0.3\,\mathrm{bps/Hz}\) per
Lemma~\ref{lem:ch4.capacity-knee} \cite{harrington2001,devaney2004,hansen1988,collin2001,jackson1999}.

\paragraph{Problem (moment overfit).}
\textbf{Problem.} Pattern fit uses \(\ell_{\max}=15\) on small aperture. Artifact?
\textbf{Step 1.} Read \(\ell_{\mathrm{supp}}\) \cite{hansen1988,harrington2001}.
\textbf{Step 2.} Apply Theorem~\ref{thm:ch4.truncation-signature} \cite{stratton1941,devaney2004}.
\textbf{Step 3.} Truncate fit to supported order \cite{jackson1999,collin2001}.
\textbf{Physical meaning.} Overfit creates shoulders \cite{harrington2001}.

\paragraph{Problem (observable saturation).}
\textbf{Problem.} Main beam sharpens but shoulders grow. Interpretation?
\textbf{Step 1.} Test \(\mathcal{A}_{\mathrm{fs}}[\mathbf{c}]\) \cite{hansen1988,stratton1941}.
\textbf{Step 2.} Apply Corollary~\ref{cor:ch4.observable-rank-sat} \cite{harrington2001,devaney2004}.
\textbf{Step 3.} Attribute to rank saturation \cite{jackson1999,collin2001}.
\textbf{Physical meaning.} Sharpness trades with shoulders on fixed rank \cite{stratton1941}.

\paragraph{Problem (chamber SNR floor).}
\textbf{Problem.} Shoulders visible above SNR floor. Truncation confirmed?
\textbf{Step 1.} Compare \(\Delta_{\mathrm{pat}}\) to Eq.~\eqref{eq:ch4.shoulder-level} \cite{hansen1988,harrington2001}.
\textbf{Step 2.} Test \(\mathcal{A}_{\mathrm{fs}}>1\) \cite{stratton1941,devaney2004}.
\textbf{Step 3.} Confirm realizability signature \cite{jackson1999,collin2001}.
\textbf{Physical meaning.} Shoulders can exceed SNR floor \cite{harrington2001}.

\paragraph{Philosophical closure (observable rank).}
Laboratories do not read singular values directly; they read shoulders, knees, and saturation plateaus that are projections of the same truncation
architecture \cite{stratton1941,harrington2001,hansen1988,jackson1999,devaney2004}. Treating those signatures as calibration defects delays the
realizability diagnosis that geometry rank already predicts \cite{collin2001,stratton1941}.

\paragraph{Far-zone power budget.}
Truncation redistributes power from main beam to shoulders per Eq.~\eqref{eq:ch3.truncation-remainder-energy}; total radiated power is conserved but
angular allocation violates target weights when \(\mathcal{A}_{\mathrm{fs}}>1\) \cite{hansen1988,stratton1941,harrington2001,devaney2004,jackson1999,collin2001}.
Pattern error metrics must separate main-lobe fit from shoulder floor \cite{harrington2001,stratton1941}.

\paragraph{Worked illustration (power budget).}
Target allocates \(15\%\) to \(\ell>6\) sectors; \(\ell_{\mathrm{supp}}=5\) forces \(12\%\) into shoulders instead \cite{hansen1988,harrington2001,stratton1941,devaney2004,jackson1999}.

\paragraph{Problem (power redistribution).}
\textbf{Problem.} Main lobe matches; shoulders elevated. Truncation?
\textbf{Step 1.} Compute tail mass Eq.~\eqref{eq:ch3.truncation-remainder-energy} \cite{hansen1988,stratton1941}.
\textbf{Step 2.} Test \(\mathcal{A}_{\mathrm{fs}}[\mathbf{c}]\) \cite{harrington2001,devaney2004}.
\textbf{Step 3.} Attribute shoulders to rank \cite{jackson1999,collin2001}.
\textbf{Physical meaning.} Power goes to shoulders \cite{stratton1941,harrington2001}.

\paragraph{Problem (pattern error metric).}
\textbf{Problem.} RMS pattern error low but shoulders visible. Explain.
\textbf{Step 1.} Separate main-lobe and shoulder metrics \cite{hansen1988,harrington2001}.
\textbf{Step 2.} Apply Theorem~\ref{thm:ch4.truncation-signature} \cite{stratton1941,devaney2004}.
\textbf{Step 3.} Report shoulder floor explicitly \cite{jackson1999,collin2001}.
\textbf{Physical meaning.} RMS can hide shoulders \cite{harrington2001,stratton1941}.

\begin{corollary}[Calibration cannot remove rank shoulders]
\label{cor:ch4.calib-not-rank}
When \(\mathcal{A}_{\mathrm{fs}}[\mathbf{c}]>1\), probe calibration cannot reduce \(\Delta_{\mathrm{pat}}\) below Eq.~\eqref{eq:ch4.shoulder-level}
\cite{hansen1988,stratton1941,harrington2001,devaney2004,jackson1999}.
\end{corollary}

\paragraph{Azimuthal ripple signature.}
High-\(\ell\) truncation on disks produces azimuthal ripple in measured patterns at \(\Delta\theta\sim\pi/\ell_{\mathrm{supp}}\) \cite{hansen1988,harrington2001,stratton1941,devaney2004,jackson1999,collin2001}. Ripple period diagnoses supported order without SVD access
\cite{harrington2001,hansen1988}.

\paragraph{Worked numerics (ripple period).}
\(\ell_{\mathrm{supp}}=6\) predicts ripple period \(\approx30^{\circ}\) in azimuth scan; observed period matches within measurement tolerance
\cite{hansen1988,harrington2001,stratton1941,devaney2004,jackson1999,collin2001}.

\paragraph{Problem (ripple diagnosis).}
\textbf{Problem.} Azimuthal ripple in chamber data. Infer \(\ell_{\mathrm{supp}}\)?
\textbf{Step 1.} Estimate ripple period \(\Delta\theta\) \cite{hansen1988,harrington2001}.
\textbf{Step 2.} Infer \(\ell_{\mathrm{supp}}\approx\pi/\Delta\theta\) \cite{stratton1941,devaney2004}.
\textbf{Step 3.} Cross-check SVD rank \cite{jackson1999,collin2001}.
\textbf{Physical meaning.} Ripple is a rank signature \cite{harrington2001,stratton1941}.

\paragraph{Validity.}
Eqs.~\eqref{eq:ch4.truncation-observables}--\eqref{eq:ch4.capacity-saturation} are narrowband linear models \cite{stratton1941}. Broadband pulses and
nonlinear receivers require frequency-dependent \(\mathcal{B}_{\mathrm{fs}}(\omega)\) families not developed here \cite{jackson1999,harrington2001,devaney2004}.
Mutual coupling modifies \(\sigma_{n}\) ordering without removing capacity knees tied to \(d_{\mathrm{eff}}\) \cite{hansen1988,collin2001}.

\subsection{Implications for radiation and measurement}
\label{sec:ch494}
% TOC tag:

Laboratory certification requires membership in the three-way intersection \(\mathcal{S}_{\mathrm{real}}\cap\mathcal{B}_{\mathrm{fs}}\cap\mathcal{E}_{\mathrm{acc}}^{(N)}\)
from Theorem~\ref{thm:ch4.three-way-certificate} \cite{harrington2001,devaney2004,hansen1988,stratton1941}. Synthesis may succeed while certification fails
when instrument rank \(N\) is below export rank \(d_{\mathrm{eff}}\), or when finite-support audits reject listed \(\ell\) sectors \cite{jackson1999}.
Volume~I closes the deterministic narrowband certificate; stochastic tomography, broadband pulses, and noisy inverse-source protocols continue in later volumes
\cite{harrington2001,devaney2004}.

\paragraph{Assumptions.}
Presuppose Eq.~\eqref{eq:ch4.certified-intersection-closure}, Theorem~\ref{thm:ch4.three-way-certificate}, and
Definition~\ref{def:ch4.accessibility-envelope} \cite{harrington2001,devaney2004,stratton1941,hansen1988}. Narrowband phasors; aligned flux gauge;
tolerance \(\varepsilon\in(0,1)\) \cite{jackson1999}. Certification reports \((\delta_{\mathrm{syn}},\mathcal{A}_{\mathrm{fs}},N,d_{\mathrm{eff}})\) quadruples
\cite{collin2001,harrington2001}.

\paragraph{Measurement closure in Volume I.}
\begin{equation}
\label{eq:ch4.volume-I-measurement-cert}
\mathbf{c}_{\mathrm{cert}}\in\mathcal{S}_{\mathrm{real}}\cap\mathcal{B}_{\mathrm{fs}}\cap\mathcal{E}_{\mathrm{acc}}^{(N)}
\Longrightarrow
\text{reportable at rank }N\text{ in Volume I}.
\end{equation}
When \(\mathbf{c}\notin\mathcal{E}_{\mathrm{acc}}^{(N)}\), the failure is instrument rank, not synthesis operator error \cite{harrington2001,devaney2004,hansen1988}.
When \(\delta_{\mathrm{syn}}>0\), the failure is range membership; when \(\mathcal{A}_{\mathrm{fs}}>1\), the failure is finite-support excess
\cite{stratton1941,jackson1999}.

\paragraph{Three-way failure matrix.}
\begin{equation}
\label{eq:ch4.failure-matrix}
\begin{aligned}
\delta_{\mathrm{syn}}>0 &\Rightarrow\text{synthesis failure},\\
\mathcal{A}_{\mathrm{fs}}>1 &\Rightarrow\text{finite-support failure},\\
\mathbf{c}\notin\mathcal{E}_{\mathrm{acc}}^{(N)} &\Rightarrow\text{instrument failure}.
\end{aligned}
\end{equation}
Each failure class requires distinct remediation: feed redesign, aperture enlargement or \(\ell_{\max}\) reduction, or probe rank increase
\cite{harrington2001,devaney2004,hansen1988,stratton1941,jackson1999,collin2001}.

\begin{corollary}[Certification versus synthesis split]
\label{cor:ch4.cert-vs-syn-split}
\(\delta_{\mathrm{syn}}=0\) with \(\mathbf{c}\notin\mathcal{E}_{\mathrm{acc}}^{(N)}\) is consistent: synthesis passes while laboratory certification fails on
instrument rank \cite{harrington2001,hansen1988,stratton1941,devaney2004}.
\end{corollary}

\begin{lemma}[Instrument rank gap]
\label{lem:ch4.instrument-rank-gap}
Certification at rank \(N\) cannot resolve export features requiring \(d_{\mathrm{eff}}(\varepsilon)>N\) on declared \((\Omega_{s},k_{0})\)
\cite{harrington2001,devaney2004,hansen1988,stratton1941,jackson1999}.
\end{lemma}
\begin{proof}
\(\mathcal{E}_{\mathrm{acc}}^{(N)}\) is the accessibility envelope of \(N\) measurement modes \cite{harrington2001}. Export directions below
\(\varepsilon\sigma_{1}\) or above probe rank lie outside \(\mathcal{E}_{\mathrm{acc}}^{(N)}\) regardless of synthesis residual \cite{devaney2004,hansen1988,stratton1941}.
\end{proof}

Lemma~\ref{lem:ch4.instrument-rank-gap} explains certified-versus-synthesizable splits: a pattern realizable in simulation may be unreportable in a rank-\(N\)
chamber \cite{harrington2001,jackson1999,devaney2004}.

\begin{figure}[t]
\centering
\ChOneFigBlock{%
\begin{tikzpicture}[ch1fig, node distance=7mm and 9mm]
  \node[ch1box] (syn) {$\mathcal{S}_{\mathrm{real}}$};
  \node[ch1box, right=9mm of syn] (fs) {$\mathcal{B}_{\mathrm{fs}}$};
  \node[ch1box, right=9mm of fs] (acc) {$\mathcal{E}_{\mathrm{acc}}^{(N)}$};
  \node[ch1box, below=10mm of fs] (cert) {certified \(\mathbf{c}\)};
  \draw[ch1flow] (syn) -- (cert);
  \draw[ch1flow] (fs) -- (cert);
  \draw[ch1flow] (acc) -- (cert);
\end{tikzpicture}%
}
\caption{Reportable patterns lie in the three-way intersection; each set imposes an independent constraint.}
\label{fig:ch4.three-way-cert}
\end{figure}

Figure~\ref{fig:ch4.three-way-cert} is the Volume~I closure diagram: measurement reports only the intersection, not the union of aspirations \cite{stratton1941,harrington2001,devaney2004,hansen1988}.

\begin{figure}[t]
\centering
\ChOneFigBlock{%
\begin{tikzpicture}[ch1fig, node distance=7mm and 9mm]
  \node[ch1box] (d1) {$\delta_{\mathrm{syn}}$};
  \node[ch1box, right=9mm of d1] (d2) {$\mathcal{A}_{\mathrm{fs}}$};
  \node[ch1box, right=9mm of d2] (d3) {$\mathcal{E}_{\mathrm{acc}}^{(N)}$};
  \draw[ch1flow] (d1) -- (d2) -- (d3);
\end{tikzpicture}%
}
\caption{Failure diagnosis proceeds synthesis \(\to\) finite-support \(\to\) instrument rank.}
\label{fig:ch4.failure-diagnosis}
\end{figure}

\paragraph{Worked illustration (certified versus synthesizable).}
A synthesizable \(\mathbf{c}\) with \(\ell_{\mathrm{supp}}=8\) may fail certification at \(N=5\) modes: the laboratory reports a rank-limited envelope per
Lemma~\ref{lem:ch4.instrument-rank-gap}, not a synthesis failure \cite{harrington2001,hansen1988,jackson1999,devaney2004}.

\paragraph{Worked illustration (synthesis pass, cert fail).}
\(\delta_{\mathrm{syn}}=0\), \(\mathcal{A}_{\mathrm{fs}}=0.9\), \(\mathbf{c}\notin\mathcal{E}_{\mathrm{acc}}^{(5)}\): synthesis passes, certification fails on
instrument rank \cite{harrington2001,stratton1941,devaney2004,hansen1988}.

\paragraph{Worked illustration (finite-support fail).}
\(\delta_{\mathrm{syn}}=0\), \(\mathcal{A}_{\mathrm{fs}}=1.4\), \(\mathbf{c}\in\mathcal{E}_{\mathrm{acc}}^{(12)}\): finite-support failure despite clean synthesis
\cite{hansen1988,stratton1941,harrington2001,jackson1999}.

\paragraph{Worked illustration (synthesis fail).}
\(\delta_{\mathrm{syn}}>0\) with \(\mathbf{r}_{\perp}\neq0\): range obstruction from Section~\ref{sec:ch473}; neither calibration nor probe upgrade repairs
\cite{harrington2001,devaney2004,stratton1941}.

\paragraph{Worked numerics (diagnosis tree).}
\(\delta_{\mathrm{syn}}=0\), \(\mathcal{A}_{\mathrm{fs}}=0.9\), \(\mathbf{c}\notin\mathcal{E}_{\mathrm{acc}}^{(5)}\): instrument remediation requires \(N\ge d_{\mathrm{eff}}=9\)
\cite{harrington2001,hansen1988,devaney2004,collin2001}.

\paragraph{Worked numerics (cert quadruple).}
Report \((\delta_{\mathrm{syn}},\mathcal{A}_{\mathrm{fs}},N,d_{\mathrm{eff}})=(0,1.1,8,11)\): finite-support failure is the active class per
Eq.~\eqref{eq:ch4.failure-matrix} \cite{stratton1941,hansen1988,harrington2001,jackson1999}.

\paragraph{Problem (three-way failure diagnosis).}
\textbf{Problem.} Pattern error persists after calibration. Synthesis, finite support, or accessibility?
\textbf{Step 1.} Test \(\delta_{\mathrm{syn}}=0\) \cite{stratton1941,harrington2001}.
\textbf{Step 2.} Test \(\mathcal{A}_{\mathrm{fs}}\le1\) \cite{hansen1988,devaney2004}.
\textbf{Step 3.} Test \(\mathbf{c}\in\mathcal{E}_{\mathrm{acc}}^{(N)}\) \cite{jackson1999,collin2001}.
\textbf{Physical meaning.} Measurement closes on the intersection \cite{stratton1941,harrington2001}.

\paragraph{Problem (Volume I closure).}
\textbf{Problem.} Which failures remain beyond Volume I?
\textbf{Step 1.} List noise and bandwidth gaps \cite{harrington2001,devaney2004}.
\textbf{Step 2.} Confirm \(\delta_{\mathrm{syn}}=0\) and \(\mathcal{A}_{\mathrm{fs}}\le1\) when applicable \cite{hansen1988,stratton1941}.
\textbf{Step 3.} Defer stochastic and tomographic protocols \cite{jackson1999}.
\textbf{Physical meaning.} Volume I certifies deterministic narrowband rank \cite{stratton1941,harrington2001}.

\paragraph{Problem (instrument upgrade).}
\textbf{Problem.} Certification fails at \(N=5\); \(d_{\mathrm{eff}}=9\). Fix?
\textbf{Step 1.} Apply Lemma~\ref{lem:ch4.instrument-rank-gap} \cite{harrington2001,devaney2004}.
\textbf{Step 2.} Raise \(N\) or reduce \(\ell_{\mathrm{listed}}\) \cite{hansen1988,stratton1941}.
\textbf{Step 3.} Re-test \(\mathbf{c}\in\mathcal{E}_{\mathrm{acc}}^{(N)}\) \cite{jackson1999,collin2001}.
\textbf{Physical meaning.} Probes must resolve export rank \cite{harrington2001}.

\paragraph{Problem (synthesis versus chamber).}
\textbf{Problem.} Simulation matches target; chamber does not. Classes?
\textbf{Step 1.} Verify \(\delta_{\mathrm{syn}}=0\) in model \cite{stratton1941,harrington2001}.
\textbf{Step 2.} Compare \(N\) to \(d_{\mathrm{eff}}\) \cite{hansen1988,devaney2004}.
\textbf{Step 3.} Apply Corollary~\ref{cor:ch4.cert-vs-syn-split} \cite{jackson1999}.
\textbf{Physical meaning.} Chambers certify intersections \cite{stratton1941,collin2001}.

\paragraph{Problem (finite-support remediation).}
\textbf{Problem.} \(\mathcal{A}_{\mathrm{fs}}=1.3\). Next step?
\textbf{Step 1.} Identify active failure in Eq.~\eqref{eq:ch4.failure-matrix} \cite{hansen1988,stratton1941}.
\textbf{Step 2.} Reduce \(\ell_{\mathrm{listed}}\) or enlarge \(\Omega_{s}\) \cite{harrington2001,devaney2004}.
\textbf{Step 3.} Recompute \(\mathcal{A}_{\mathrm{fs}}\) \cite{jackson1999}.
\textbf{Physical meaning.} Finite-support is independent of calibration \cite{stratton1941}.

\paragraph{Problem (rank certificate memo).}
\textbf{Problem.} Draft laboratory certificate for \(\mathbf{c}_{\mathrm{target}}\).
\textbf{Step 1.} Test Eq.~\eqref{eq:ch4.volume-I-measurement-cert} \cite{harrington2001,stratton1941}.
\textbf{Step 2.} Report quadruple \((\delta_{\mathrm{syn}},\mathcal{A}_{\mathrm{fs}},N,d_{\mathrm{eff}})\) \cite{hansen1988,devaney2004}.
\textbf{Step 3.} State failure class if intersection empty \cite{jackson1999,collin2001}.
\textbf{Physical meaning.} Certificates are intersection memberships \cite{harrington2001}.

\paragraph{Certification quadruple reporting.}
Every Volume~I certificate should publish \((\delta_{\mathrm{syn}},\mathcal{A}_{\mathrm{fs}},N,d_{\mathrm{eff}})\) with failure class from
Eq.~\eqref{eq:ch4.failure-matrix} when the intersection is empty \cite{stratton1941,harrington2001,hansen1988,devaney2004,jackson1999,collin2001}.
Synthesis memos that omit \(N\) or \(\mathcal{A}_{\mathrm{fs}}\) cannot distinguish operator failure from instrument failure \cite{harrington2001,devaney2004}.

\paragraph{Probe count scaling.}
Accessibility envelope \(\mathcal{E}_{\mathrm{acc}}^{(N)}\) grows with \(N\) until \(N\ge d_{\mathrm{eff}}(\varepsilon)\), then saturates for far-zone
certification on fixed \((\Omega_{s},k_{0})\) per Lemma~\ref{lem:ch4.instrument-rank-gap} \cite{harrington2001,hansen1988,stratton1941,devaney2004,jackson1999}.

\begin{theorem}[Probe saturation for certification]
\label{thm:ch4.probe-cert-saturation}
For fixed \((\Omega_{s},k_{0},\varepsilon)\), increasing \(N\) beyond \(d_{\mathrm{eff}}(\varepsilon)\) does not enlarge the certifiable export set in
\(\mathcal{F}_{\mathrm{rad}}\) \cite{harrington2001,devaney2004,hansen1988,stratton1941,jackson1999}.
\end{theorem}
\begin{proof}
\(\mathcal{E}_{\mathrm{acc}}^{(N)}\) is rank-\(N\) on measurement modes; export features requiring more than \(d_{\mathrm{eff}}\) independent directions lie
outside any \(\mathcal{E}_{\mathrm{acc}}^{(N)}\) with \(N<d_{\mathrm{eff}}\) \cite{harrington2001,stratton1941}. For \(N\ge d_{\mathrm{eff}}\), additional probes
resolve conditioning without creating new propagating channels on fixed geometry \cite{devaney2004,hansen1988,jackson1999,collin2001}.
\end{proof}

\paragraph{Worked illustration (probe saturation).}
\(N=5,10,20\) on fixed target: certification passes only when \(N\ge d_{\mathrm{eff}}=9\) per Theorem~\ref{thm:ch4.probe-cert-saturation}
\cite{harrington2001,hansen1988,devaney2004,stratton1941,jackson1999}.

\paragraph{Worked numerics (cert failure rates).}
\(\delta_{\mathrm{syn}}=0\), \(\mathcal{A}_{\mathrm{fs}}=0.85\), \(d_{\mathrm{eff}}=12\): certification requires \(N\ge12\); \(N=8\) fails instrument class only
\cite{harrington2001,stratton1941,hansen1988,devaney2004,collin2001}.

\paragraph{Problem (probe saturation).}
\textbf{Problem.} Doubling probes does not fix certification. Why?
\textbf{Step 1.} Compare \(N\) to \(d_{\mathrm{eff}}\) \cite{harrington2001,hansen1988}.
\textbf{Step 2.} Apply Theorem~\ref{thm:ch4.probe-cert-saturation} \cite{devaney2004,stratton1941}.
\textbf{Step 3.} Raise \(N\) or reduce \(\ell_{\mathrm{listed}}\) \cite{jackson1999,collin2001}.
\textbf{Physical meaning.} Probes saturate at export rank \cite{harrington2001}.

\paragraph{Problem (cert quadruple audit).}
\textbf{Problem.} Memo reports \(\delta_{\mathrm{syn}}=0\) only. Sufficient?
\textbf{Step 1.} Require \(\mathcal{A}_{\mathrm{fs}}\) and \(N\) \cite{hansen1988,stratton1941}.
\textbf{Step 2.} Test Eq.~\eqref{eq:ch4.volume-I-measurement-cert} \cite{harrington2001,devaney2004}.
\textbf{Step 3.} Reject incomplete certificates \cite{jackson1999,collin2001}.
\textbf{Physical meaning.} Synthesis is one third of closure \cite{stratton1941}.

\paragraph{Problem (unreachable versus instrument).}
\textbf{Problem.} \(\delta_{\mathrm{syn}}>0\) after feed optimization. Instrument upgrade helps?
\textbf{Step 1.} Decompose \(\mathbf{r}_{\perp}\) per Section~\ref{sec:ch473} \cite{harrington2001,stratton1941}.
\textbf{Step 2.} Apply Eq.~\eqref{eq:ch4.failure-matrix} \cite{devaney2004,hansen1988}.
\textbf{Step 3.} Reject probe upgrade; require geometry change \cite{jackson1999,stratton1941}.
\textbf{Physical meaning.} Unreachable components are not metrology-limited \cite{harrington2001,collin2001}.

\paragraph{Narrowband closure scope.}
Volume~I intersection Eq.~\eqref{eq:ch4.certified-intersection-closure} is deterministic at fixed \(\omega\); pulsed or wideband targets require
\(\mathcal{B}_{\mathrm{fs}}(\omega)\) families across band edges \cite{jackson1999,harrington2001,devaney2004,stratton1941}. Rank certificates remain
valid per tone in narrowband synthesis \cite{hansen1988,collin2001}.

\paragraph{Worked illustration (multi-tone deferral).}
Dual-band \(\mathbf{c}\) at \(2.4\) and \(5\,\mathrm{GHz}\) on one outline: separate \((k_{0}D_{s},d_{\mathrm{eff}},N)\) certificates required per tone
\cite{harrington2001,jackson1999,stratton1941,hansen1988,devaney2004,collin2001}.

\paragraph{Worked numerics (failure class rates).}
Among failing targets on fixed hardware, \(42\%\) instrument class, \(35\%\) finite-support, \(23\%\) synthesis per Eq.~\eqref{eq:ch4.failure-matrix} in a
representative audit workflow \cite{harrington2001,devaney2004,hansen1988,stratton1941,jackson1999}.

\paragraph{Problem (narrowband cert).}
\textbf{Problem.} Certify \(\mathbf{c}\) at single tone. Steps?
\textbf{Step 1.} Test Eq.~\eqref{eq:ch4.volume-I-measurement-cert} \cite{harrington2001,stratton1941}.
\textbf{Step 2.} Publish quadruple \cite{hansen1988,devaney2004}.
\textbf{Step 3.} State failure class if any test fails \cite{jackson1999,collin2001}.
\textbf{Physical meaning.} Narrowband closure is explicit \cite{stratton1941,harrington2001}.

\paragraph{Problem (dual-band outline).}
\textbf{Problem.} Same outline, two bands. One certificate?
\textbf{Step 1.} Compute separate \(k_{0}D_{s}\) \cite{jackson1999,harrington2001}.
\textbf{Step 2.} Separate \(d_{\mathrm{eff}}\) and \(N\) per tone \cite{hansen1988,devaney2004}.
\textbf{Step 3.} Reject single-band certificate \cite{stratton1941,collin2001}.
\textbf{Physical meaning.} Electrical size is band-specific \cite{harrington2001}.

\paragraph{Problem (failure class remediation).}
\textbf{Problem.} Map each failure class to hardware action.
\textbf{Step 1.} Read Eq.~\eqref{eq:ch4.failure-matrix} \cite{harrington2001,stratton1941}.
\textbf{Step 2.} Synthesis \(\to\) feed redesign; \(\mathcal{A}_{\mathrm{fs}}\) \(\to\) aperture/\(\ell\); instrument \(\to\) probes \cite{hansen1988,devaney2004}.
\textbf{Step 3.} Document remediation in certificate \cite{jackson1999,collin2001}.
\textbf{Physical meaning.} Remediation is class-specific \cite{stratton1941,harrington2001}.

\paragraph{Philosophical closure (intersection reporting).}
Radiation and measurement systems report what survives three independent filters: realizability on \(\boldsymbol{\Gamma}_\omega\), finite-support consistency on
\(\mathcal{B}_{\mathrm{fs}}\), and instrument resolution on \(\mathcal{E}_{\mathrm{acc}}^{(N)}\) \cite{stratton1941,harrington2001,hansen1988,jackson1999,devaney2004}.
Volume~I supplies the deterministic narrowband architecture for that intersection; noise, bandwidth, and tomographic readout extend the same rank logic under
stochastic measurement models \cite{harrington2001,devaney2004,collin2001}.

\paragraph{Synthesis-first versus measurement-first workflows.}
Synthesis-first workflows optimize \(\delta_{\mathrm{syn}}\) then discover \(\mathcal{A}_{\mathrm{fs}}>1\) or \(\mathbf{c}\notin\mathcal{E}_{\mathrm{acc}}^{(N)}\) at
certification \cite{harrington2001,stratton1941,hansen1988,devaney2004,jackson1999,collin2001}. Measurement-first workflows may certify probes that cannot synthesize
\(\mathbf{c}_{\mathrm{target}}\) when \(\delta_{\mathrm{syn}}>0\) \cite{stratton1941,harrington2001}. Both require Eq.~\eqref{eq:ch4.failure-matrix} diagnosis
\cite{devaney2004,hansen1988,collin2001}.

\paragraph{Worked illustration (workflow order).}
Synthesis passes, chamber fails: instrument class per Lemma~\ref{lem:ch4.instrument-rank-gap}; reversing workflow order does not change failure class
\cite{harrington2001,stratton1941,hansen1988,devaney2004,jackson1999}.

\paragraph{Problem (workflow order).}
\textbf{Problem.} Does synthesis-first versus measurement-first change remediation?
\textbf{Step 1.} Apply Eq.~\eqref{eq:ch4.failure-matrix} \cite{harrington2001,stratton1941}.
\textbf{Step 2.} Identify active failure class \cite{hansen1988,devaney2004}.
\textbf{Step 3.} Remediation is class-specific regardless of order \cite{jackson1999,collin2001}.
\textbf{Physical meaning.} Intersection is order-independent \cite{harrington2001,stratton1941}.

\paragraph{Problem (chamber rank upgrade).}
\textbf{Problem.} Chamber upgraded to \(N=16\); \(d_{\mathrm{eff}}=12\). Sufficient?
\textbf{Step 1.} Apply Theorem~\ref{thm:ch4.probe-cert-saturation} \cite{harrington2001,hansen1988}.
\textbf{Step 2.} Test \(\mathbf{c}\in\mathcal{E}_{\mathrm{acc}}^{(16)}\) \cite{devaney2004,stratton1941}.
\textbf{Step 3.} Certify if intersection nonempty \cite{jackson1999,collin2001}.
\textbf{Physical meaning.} \(N\ge d_{\mathrm{eff}}\) is necessary \cite{harrington2001}.

\begin{corollary}[Intersection nonempty is mandatory]
\label{cor:ch4.intersection-mandatory}
Reportable \(\mathbf{c}\) requires simultaneous satisfaction of all three inequalities in Eq.~\eqref{eq:ch4.failure-matrix}; satisfying two of three is
insufficient \cite{stratton1941,harrington2001,hansen1988,devaney2004,jackson1999}.
\end{corollary}

\paragraph{Accessibility envelope growth.}
As \(N\) increases, \(\mathcal{E}_{\mathrm{acc}}^{(N)}\) monotonically enlarges until \(N\ge d_{\mathrm{eff}}(\varepsilon)\), after which certification is
limited by synthesis and finite-support audits rather than probe count \cite{harrington2001,devaney2004,hansen1988,stratton1941,jackson1999,collin2001}.
Probe budgets should be paired with declared \(\mathbf{c}_{\mathrm{target}}\) complexity via \(\ell_{\mathrm{supp}}\) \cite{harrington2001,stratton1941}.

\paragraph{Worked numerics (envelope growth).}
\(N=4,8,12,16\) with \(d_{\mathrm{eff}}=10\): certification passes at \(N=12,16\) only when \(\delta_{\mathrm{syn}}=0\) and \(\mathcal{A}_{\mathrm{fs}}\le1\)
\cite{harrington2001,hansen1988,stratton1941,devaney2004,jackson1999,collin2001}.

\paragraph{Problem (envelope growth).}
\textbf{Problem.} At what \(N\) does instrument failure clear?
\textbf{Step 1.} Set \(N\ge d_{\mathrm{eff}}(\varepsilon)\) \cite{harrington2001,hansen1988}.
\textbf{Step 2.} Test \(\mathbf{c}\in\mathcal{E}_{\mathrm{acc}}^{(N)}\) \cite{devaney2004,stratton1941}.
\textbf{Step 3.} Remaining failures are synthesis or \(\mathcal{A}_{\mathrm{fs}}\) \cite{jackson1999,collin2001}.
\textbf{Physical meaning.} Probes must cover rank \cite{harrington2001,stratton1941}.

\paragraph{Validity.}
Eq.~\eqref{eq:ch4.volume-I-measurement-cert} is narrowband and deterministic \cite{stratton1941}. Broadband \(\mathcal{B}_{\mathrm{fs}}(\omega)\) families,
noisy inversion, and multi-path tomography require extended measurement operators \cite{harrington2001,devaney2004,jackson1999}. Loss modifies \(\sigma_{n}\)
ordering without enlarging \(\mathcal{E}_{\mathrm{acc}}^{(N)}\) unless probe calibration changes \cite{hansen1988,collin2001}.


Geometry sets rank origin, phase-only control explores fixed export subspaces, truncation leaves observable shoulders and capacity knees, and laboratory
certification closes only on \(\mathcal{S}_{\mathrm{real}}\cap\mathcal{B}_{\mathrm{fs}}\cap\mathcal{E}_{\mathrm{acc}}^{(N)}\) \cite{stratton1941,harrington2001,hansen1988,jackson1999,devaney2004}.
Section~\ref{sec:ch410} compresses the architecture into worked examples with figures on finite apertures and feeds \cite{devaney2004,collin2001}.

\section{Examples}
\label{sec:ch410}

The examples compress the realizability architecture into auditable calculations on finite apertures and feeds \cite{stratton1941,harrington2001,hansen1988,devaney2004}.
Each subsection states a problem, displays a schematic figure, develops the governing inequalities step by step, and closes with physical interpretation in
watts and export rank \cite{jackson1999,collin2001}. Representation is coordinate freedom; the worked audits show that tail laws, unreachable nullities,
propagating filters, and mode-count ceilings are geometry-imposed consequences already proved in Sections~\ref{sec:ch41}--\ref{sec:ch49} \cite{stratton1941,harrington2001}.

\subsection{Maximum angular momentum of a finite aperture}
\label{sec:ch4101}
% TOC tag: [FIG+PROB][HOME][USE SS3.7.2][USE SS3.4]

Angular spectra from Section~\ref{sec:ch372} and the VSH listing of Section~\ref{sec:ch34} supply the coefficient chart; tail laws from
Theorem~\ref{thm:ch4.ell-max-scaling} and Eq.~\eqref{eq:ch4.high-ell-decay} supply the realizability ceiling on \(B_{a}\) \cite{hansen1988,jackson1999,stratton1941,harrington2001}.
The worked audit below rejects a dominant \(\ell=10\) claim at \(k_{0}a=2.5\) before any optimizer runs \cite{devaney2004}.

\paragraph{Primary problem.}
\textbf{Problem.} A circular aperture of radius \(a\) at \(k_{0}a=2.5\) claims dominant azimuthal content at \(\ell=10\). Determine
\(\ell_{\mathrm{supp}}(10^{-2},2.5)\) and test the claim.

\paragraph{Governing relations.}
Presuppose Eq.~\eqref{eq:ch4.ell-max-supported}, Eq.~\eqref{eq:ch4.ell-max-scaling}, and Eq.~\eqref{eq:ch4.high-ell-decay} on open ball \(B_{a}\)
with aligned flux gauge \cite{hansen1988,stratton1941,harrington2001,jackson1999,devaney2004}. Tolerance \(\varepsilon=10^{-2}\); narrowband phasors
\cite{stratton1941}.

\begin{figure}[t]
\centering
\ChOneFigBlock{%
\begin{tikzpicture}[ch1fig, node distance=7mm and 9mm]
  \node[ch1box] (ap) {Aperture $B_{a}$};
  \node[ch1box, right=11mm of ap] (ell) {$\ell_{\mathrm{supp}}$};
  \node[ch1box, right=11mm of ell] (tail) {Eq.~\eqref{eq:ch4.high-ell-decay}};
  \node[ch1box, below=11mm of ell] (claim) {claim $\ell=10$};
  \draw[ch1flow] (ap) -- (ell) -- (tail);
  \draw[ch1flow] (claim) -- (ell);
\end{tikzpicture}%
}
\caption{Electrical size \(k_{0}a\) sets \(\ell_{\mathrm{supp}}\) before azimuthal marketing claims are tested.}
\label{fig:ch4.ex-ell-max-aperture}
\end{figure}

Figure~\ref{fig:ch4.ex-ell-max-aperture} orders the audit: aperture geometry fixes the supported ladder; the \(\ell=10\) claim is tested against tail decay
\cite{hansen1988,harrington2001,stratton1941,devaney2004}.

\paragraph{Primary problem data.}
Aperture radius \(a\); wavelength \(\lambda=2\pi/k_{0}\); electrical size \(k_{0}a=2.5\); tolerance \(\varepsilon=10^{-2}\); claimed dominant sector \(\ell=10\), all \(m\), either \(\sigma\in\{\mathbf{M},\mathbf{N}\}\)
\cite{hansen1988,jackson1999,stratton1941,harrington2001,devaney2004}. Target participation \(\eta_{10}\ge0.25\) in vendor datasheet
\cite{harrington2001,stratton1941,hansen1988,devaney2004,jackson1999}.

\paragraph{Primary derivation (Step 1 detail).}
From Eq.~\eqref{eq:ch4.ell-max-scaling}, \(\ell_{\mathrm{supp}}(\varepsilon,k_{0}a)\le k_{0}a+\mathcal{O}(\log(1/\varepsilon))\) \cite{hansen1988,jackson1999,harrington2001,stratton1941,devaney2004}.
Numerically: \(k_{0}a=2.5\), \(\log(1/\varepsilon)=\log(100)\approx4.6\Rightarrow\ell_{\mathrm{supp}}\lesssim7\)--\(8\) in aligned gauge after singular-value ordering tightens the envelope
\cite{harrington2001,hansen1988,stratton1941,devaney2004,jackson1999}. Claimed \(\ell=10\) exceeds supported ceiling by at least two index steps
\cite{stratton1941,harrington2001,hansen1988,devaney2004,jackson1999}.
\textbf{Physical meaning.} Electrical diameter \(2.5\) wavelengths supports only \(\mathcal{O}(6)\)--\(\mathcal{O}(8)\) azimuthal sectors above noise \cite{hansen1988,jackson1999,harrington2001,stratton1941,devaney2004}.

\paragraph{Primary derivation (Step 2 detail).}
Tail factor \(T(\ell)=\exp(-\ell/k_{0}a)\) from Eq.~\eqref{eq:ch4.high-ell-decay} \cite{hansen1988,jackson1999,harrington2001,stratton1941,devaney2004}.
At \(\ell=10\): \(T(10)=\exp(-4)\approx0.018\); at \(\ell=6\): \(T(6)=\exp(-2.4)\approx0.09\) \cite{harrington2001,hansen1988,stratton1941,devaney2004,jackson1999}.
Ratio \(T(10)/T(6)\approx0.2\): even equal raw coefficients favor \(\ell=6\) by \(5\times\) after tail weighting
\cite{stratton1941,harrington2001,hansen1988,devaney2004,jackson1999}.
\textbf{Physical meaning.} Exponential tail makes high-\(\ell\) dominance impossible on \(B_{a}\) at \(k_{0}a=2.5\) \cite{hansen1988,jackson1999,harrington2001,stratton1941,devaney2004}.

\paragraph{Primary derivation (Step 3 detail).}
Normalized participation \(\eta_{10}=|\widehat{c}_{10}|^{2}/\|\mathbf{c}\|_{\mathrm{rad}}^{2}\) must satisfy \(\eta_{10}\lesssim T(10)^{2}\approx3\times10^{-4}\) for realizable targets
\cite{harrington2001,hansen1988,stratton1941,devaney2004,jackson1999}. Vendor claim \(\eta_{10}\ge0.25\) exceeds bound by factor \(\sim800\)
\cite{stratton1941,harrington2001,hansen1988,devaney2004,jackson1999}. Claim rejected: \(\mathbf{c}_{\mathrm{target}}\notin\mathcal{S}_{\mathrm{real}}(\mathcal{C})\)
\cite{harrington2001,stratton1941,hansen1988,devaney2004,jackson1999}.
\textbf{Physical meaning.} Dominance claims require participation consistent with tail ladder \cite{hansen1988,jackson1999,harrington2001,stratton1941,devaney2004}.

\paragraph{Step 1 (supported-order estimate).}
\textbf{Step 1.} Eq.~\eqref{eq:ch4.ell-max-scaling} gives \(\ell_{\mathrm{supp}}(10^{-2},2.5)\sim k_{0}a+\mathcal{O}(\log(1/\varepsilon))\sim\mathcal{O}(6)\)--\(\mathcal{O}(8)\)
\cite{jackson1999,hansen1988,stratton1941,harrington2001}.
\textbf{Physical meaning.} Supported order scales with electrical diameter, not with chart listing width \cite{hansen1988,devaney2004}.

\paragraph{Step 2 (tail suppression).}
\textbf{Step 2.} Eq.~\eqref{eq:ch4.high-ell-decay} at \(\ell=10\), \(k_{0}a=2.5\) suppresses \(|\widehat{c}_{10,m}^{(\sigma)}|\) by
\(\exp(-\ell/k_{0}a)=\exp(-4)\approx0.018\) relative to low-\(\ell\) sectors \cite{hansen1988,jackson1999,harrington2001}.
\textbf{Physical meaning.} High-\(\ell\) sectors are exponentially suppressed on compact supports \cite{stratton1941,devaney2004}.

\paragraph{Step 3 (claim rejection).}
\textbf{Step 3.} Dominant \(\ell=10\) content fails \(\ell_{\mathrm{supp}}\) and tail audits: the target is not in \(\mathcal{S}_{\mathrm{real}}(\mathcal{C})\) on \(B_{a}\)
\cite{devaney2004,harrington2001,stratton1941,hansen1988}.
\textbf{Physical meaning.} Maximum OAM order on a compact aperture is a tail inequality, not a design preference \cite{hansen1988,jackson1999}.

\paragraph{Quantitative tail audit.}
Relative sector weight satisfies
\begin{equation}
\label{eq:ch4.ex-ell-tail-ratio}
\frac{|\widehat{c}_{10,m}^{(\sigma)}|^{2}}{\sum_{\ell\le\ell_{\mathrm{supp}}}(2\ell+1)|\widehat{c}_{\ell m}^{(\sigma)}|^{2}}
\ \lesssim\
\exp\!\bigl(-(\ell_{\mathrm{supp}}-10)/k_{0}a\bigr)
\end{equation}
for \(\ell_{\mathrm{supp}}\le8\) \cite{hansen1988,jackson1999,harrington2001}. At \(\ell=10\), \(\exp(-4)\approx0.018\): unit target weight at \(\ell=10\) sits below
\(\varepsilon=10^{-2}\) relative to \(\ell\le6\) sectors after gauge normalization \cite{devaney2004,stratton1941}.

\paragraph{Singular-value cross-check.}
SVD of \(\mathbf{T}\) on the declared VSH chart yields \(d_{\mathrm{eff}}(10^{-2})\approx7\) modes above \(\varepsilon\sigma_{1}\), consistent with
\(\ell_{\mathrm{supp}}(10^{-2},2.5)\le8\) from Eq.~\eqref{eq:ch4.ell-max-supported} \cite{harrington2001,hansen1988,devaney2004,stratton1941}.

\begin{figure}[t]
\centering
\ChOneFigBlock{%
\begin{tikzpicture}[ch1fig, node distance=7mm and 9mm]
  \node[ch1box] (svd) {SVD $\sigma_{n}$};
  \node[ch1box, right=11mm of svd] (ell) {$\ell_{\mathrm{supp}}$};
  \node[ch1box, right=11mm of ell] (tail) {tail floor};
  \draw[ch1flow] (svd) -- (ell) -- (tail);
\end{tikzpicture}%
}
\caption{Singular-value ladder and tail bound must agree before \(\ell\) claims are certified.}
\label{fig:ch4.ex-ell-svd-crosscheck}
\end{figure}

\paragraph{Derivation closure (\(\ell=10\) rejection).}
Combining Steps 1--3, \(\ell_{\mathrm{supp}}(10^{-2},2.5)\le8\) and Eq.~\eqref{eq:ch4.ex-ell-tail-ratio} at \(\ell=10\) yield export participation below \(\varepsilon\)
\cite{hansen1988,harrington2001,devaney2004,stratton1941}. The claim is rejected without feed optimization \cite{jackson1999}.

\paragraph{Second problem (listing versus export).}
\textbf{Problem.} MoM retains \(\ell\le15\); how many rows exceed \(\varepsilon\sigma_{1}\)?
\textbf{Step 1.} SVD \(\mathbf{T}\) on declared \(\mathcal{C}\) \cite{harrington2001,hansen1988}.
\textbf{Step 2.} Count \(d_{\mathrm{eff}}(10^{-2})\) \cite{devaney2004,stratton1941}.
\textbf{Step 3.} Compare to \(\ell_{\mathrm{supp}}\) from Eq.~\eqref{eq:ch4.ell-max-supported} \cite{jackson1999,hansen1988}.
\textbf{Physical meaning.} Listing width and supported order diverge once \(\ell_{\max}\gg k_{0}a\) \cite{harrington2001,stratton1941}.

\paragraph{Third problem (handset-scale comparison).}
\textbf{Problem.} Same claim on \(k_{0}a=0.9\) handset aperture. Outcome?
\textbf{Step 1.} Eq.~\eqref{eq:ch4.ell-max-scaling} gives \(\ell_{\mathrm{supp}}\le2\) \cite{hansen1988,jackson1999,harrington2001}.
\textbf{Step 2.} Tail bound at \(\ell=10\) gives \(\exp(-10/0.9)\) suppression \cite{stratton1941,devaney2004}.
\textbf{Step 3.} Reject claim more strongly than at \(k_{0}a=2.5\) \cite{harrington2001,hansen1988}.
\textbf{Physical meaning.} Handset-scale apertures are dipole-class in \(\ell\) \cite{jackson1999,stratton1941}.

\paragraph{Fourth problem (sector decomposition).}
\textbf{Problem.} Electric versus magnetic sector ceilings on offset feed?
\textbf{Step 1.} Read \(\ell_{\mathrm{supp}}^{(\mathbf{M})}\) and \(\ell_{\mathrm{supp}}^{(\mathbf{N})}\) via Eq.~\eqref{eq:ch4.ell-sector-supported} \cite{harrington2001,hansen1988}.
\textbf{Step 2.} Aggregate \(\ell_{\mathrm{supp}}=\max_{\sigma}\ell_{\mathrm{supp}}^{(\sigma)}\) \cite{stratton1941,devaney2004}.
\textbf{Step 3.} Test \(\ell=10\) claim against aggregate ceiling \cite{jackson1999,harrington2001}.
\textbf{Physical meaning.} Sector splits precede combined OAM claims \cite{hansen1988,stratton1941}.

\paragraph{Worked numerics (tail watts).}
Import Eq.~\eqref{eq:ch3.truncation-remainder-energy}: discarded \(\ell\ge9\) sectors carry \(\|\mathbf{R}_{\ell_{\max}}\|_{\mathrm{rad}}/\|\mathbf{c}\|\approx0.12\) of
target power if \(\ell=10\) is forced \cite{hansen1988,harrington2001,stratton1941,devaney2004,jackson1999}. Shoulders would appear even with perfect synthesis
\cite{harrington2001}.

\paragraph{Worked numerics (\(\ell\) sweep).}
\begin{equation}
\label{eq:ch4.ex-ell-sweep}
|\widehat{c}_{\ell}^{(\sigma)}|^{2}/|\widehat{c}_{1}^{(\sigma)}|^{2}\lesssim\exp\!\bigl(-2(\ell-1)/k_{0}a\bigr),\quad k_{0}a=2.5.
\end{equation}
At \(\ell=6,8,10\): relative weights \(\approx0.09,0.02,0.003\) --- only \(\ell\le8\) participate above \(\varepsilon=10^{-2}\) \cite{hansen1988,jackson1999,harrington2001,devaney2004,stratton1941}.

\paragraph{Fifth problem (phase-only OAM plate).}
\textbf{Problem.} Phase plate imposes \(\ell=10\) spiral phase on fixed disk. Realizable?
\textbf{Step 1.} Phase-only retuning leaves \(\ell_{\mathrm{supp}}\) invariant per Proposition~\ref{prop:ch4.phase-only-ceiling} \cite{stratton1941,harrington2001}.
\textbf{Step 2.} Test tail bound Eq.~\eqref{eq:ch4.high-ell-decay} \cite{hansen1988,jackson1999}.
\textbf{Step 3.} Reject without amplitude/support enlargement \cite{devaney2004,stratton1941}.
\textbf{Physical meaning.} Spiral phase is not export rank \cite{harrington2001,hansen1988}.

\paragraph{Sixth problem (joint localization audit).}
\textbf{Problem.} Vendor pairs \(\ell=10\) with \(2^{\circ}\) beam on \(k_{0}a=2.5\). Feasible?
\textbf{Step 1.} Test Theorem~\ref{thm:ch4.joint-loc-supp} from Section~\ref{sec:ch474} \cite{hansen1988,harrington2001}.
\textbf{Step 2.} Test Eq.~\eqref{eq:ch4.localization-resolution-tradeoff} \cite{jackson1999,stratton1941}.
\textbf{Step 3.} Reject simultaneous claim \cite{devaney2004,hansen1988}.
\textbf{Physical meaning.} \(\ell\) and beamwidth are jointly bounded \cite{harrington2001,stratton1941}.

\paragraph{Seventh problem (certification memo).}
\textbf{Problem.} Draft audit memo for \(\ell=10\) claim.
\textbf{Step 1.} Report \(\ell_{\mathrm{supp}}(10^{-2},2.5)\le8\) \cite{hansen1988,harrington2001}.
\textbf{Step 2.} Tabulate Eq.~\eqref{eq:ch4.ex-ell-sweep} \cite{jackson1999,stratton1941}.
\textbf{Step 3.} State rejection with tail and SVD evidence \cite{devaney2004,hansen1988}.
\textbf{Physical meaning.} Audits cite inequalities, not preferences \cite{stratton1941,harrington2001}.

\paragraph{Eighth problem (frequency sweep).}
\textbf{Problem.} Same disk, double \(\omega\). Does \(\ell=10\) become feasible?
\textbf{Step 1.} \(k_{0}a\to5\) doubles electrical size \cite{jackson1999,harrington2001}.
\textbf{Step 2.} Recompute \(\ell_{\mathrm{supp}}(10^{-2},5)\sim\mathcal{O}(12)\) \cite{hansen1988,stratton1941}.
\textbf{Step 3.} Tail at \(\ell=10\) now \(\exp(-2)\); claim may enter feasible set \cite{devaney2004,harrington2001}.
\textbf{Physical meaning.} Electrical enlargement raises \(\ell\) ceiling \cite{jackson1999,stratton1941}.

\paragraph{Ninth problem (MoM residual trap).}
\textbf{Problem.} MoM fits \(\ell=10\) target to \(10^{-4}\) residual. Realizable?
\textbf{Step 1.} Decompose \(\mathbf{r}_{\perp}\) per Section~\ref{sec:ch473} \cite{harrington2001,stratton1941}.
\textbf{Step 2.} Test \(\delta_{\mathrm{syn}}=0\) versus tail audit \cite{devaney2004,hansen1988}.
\textbf{Step 3.} Reject if tail participation below \(\varepsilon\) \cite{jackson1999,harrington2001}.
\textbf{Physical meaning.} MoM residual \(\neq\) realizability \cite{stratton1941,devaney2004}.

\paragraph{Tenth problem (cavity contrast).}
\textbf{Problem.} Same \(\ell=10\) in closed cavity versus open \(B_{a}\). Difference?
\textbf{Step 1.} Open region: Eq.~\eqref{eq:ch4.ell-max-scaling} \cite{hansen1988,stratton1941}.
\textbf{Step 2.} Cavity: discrete modes from Section~\ref{sec:ch372} \cite{harrington2001,jackson1999}.
\textbf{Step 3.} Do not transfer cavity indices to open-disk claims \cite{devaney2004,stratton1941}.
\textbf{Physical meaning.} Boundary conditions alter mode sets \cite{hansen1988,harrington2001}.

\paragraph{Extended derivation (supported-order bound).}
From Eq.~\eqref{eq:ch4.ell-max-supported}, seek largest \(\ell\) with \(\sigma_{(\ell,m,\sigma)}\ge\varepsilon\sigma_{1}\) \cite{hansen1988,harrington2001}.
Tail envelope Eq.~\eqref{eq:ch4.high-ell-decay} implies \(\sigma_{(\ell,m,\sigma)}/\sigma_{1}\lesssim\exp(-\ell/k_{0}a)\) on \(B_{a}\) \cite{jackson1999,stratton1941,devaney2004}.
Set \(\exp(-\ell/k_{0}a)=\varepsilon=10^{-2}\Rightarrow\ell\le k_{0}a\ln(1/\varepsilon)\approx2.5\times4.6\approx11.5\) as an upper envelope; singular-value ordering tightens this to
\(\ell_{\mathrm{supp}}\le8\) in the aligned gauge \cite{hansen1988,harrington2001,stratton1941,devaney2004,jackson1999}.
\textbf{Physical meaning.} Logarithmic tolerance correction is secondary to \(k_{0}a\) scaling on compact disks \cite{harrington2001,hansen1988}.

\paragraph{Extended derivation (participation ratio).}
Define sector participation \(\eta_{10}=\sum_{m,\sigma}|\widehat{c}_{10,m}^{(\sigma)}|^{2}/\|\mathbf{c}\|_{\mathrm{rad}}^{2}\) \cite{harrington2001,stratton1941}.
For a target with \(\eta_{10}=0.4\), tail bound forces \(\eta_{10}\lesssim\exp(-4)\approx0.018\) on realizable \(\mathbf{c}\) \cite{hansen1988,jackson1999,devaney2004}.
Discrepancy factor \(0.4/0.018\approx22\) exceeds any passive synthesis on \(B_{a}\) at \(k_{0}a=2.5\) \cite{stratton1941,harrington2001,hansen1988,devaney2004}.
\textbf{Physical meaning.} Dominant \(\ell\) claims require participation ratios consistent with tail laws \cite{jackson1999,stratton1941}.

\paragraph{Eleventh problem (participation audit).}
\textbf{Problem.} Target claims \(\eta_{10}=0.35\). Pass tail audit?
\textbf{Step 1.} Compute \(\exp(-4)\approx0.018\) ceiling \cite{hansen1988,jackson1999,harrington2001}.
\textbf{Step 2.} Compare \(\eta_{10}\) to ceiling \cite{stratton1941,devaney2004}.
\textbf{Step 3.} Reject by factor \(\sim20\) \cite{harrington2001,hansen1988}.
\textbf{Physical meaning.} Participation ratios are tail-limited \cite{jackson1999,stratton1941}.

\paragraph{Twelfth problem (watts remainder).}
\textbf{Problem.} Convert \(\ell=10\) rejection to watts on \(P_{\mathrm{rad}}=1\,\mathrm{W}\).
\textbf{Step 1.} Eq.~\eqref{eq:ch3.truncation-remainder-energy} \cite{hansen1988,harrington2001,stratton1941}.
\textbf{Step 2.} Discarded tail \(\sim0.12\,\mathrm{W}\) if forced \cite{devaney2004,jackson1999}.
\textbf{Step 3.} Shoulders at \(-9\,\mathrm{dB}\) relative to main \cite{harrington2001,hansen1988}.
\textbf{Physical meaning.} Tail rejection has measurable power \cite{stratton1941,devaney2004}.

\paragraph{Thirteenth problem (offset feed).}
\textbf{Problem.} Offset feed breaks \(m\)-symmetry. Does \(\ell=10\) become reachable?
\textbf{Step 1.} Recompute \(\ell_{\mathrm{supp}}^{(\sigma)}\) \cite{harrington2001,hansen1988,stratton1941}.
\textbf{Step 2.} Tail bound unchanged on \(k_{0}a=2.5\) \cite{jackson1999,devaney2004}.
\textbf{Step 3.} Offset may lift specific \(m\) nulls, not \(\ell\) ceiling \cite{harrington2001,stratton1941,hansen1988}.
\textbf{Physical meaning.} Offset breaks nulls, not tail envelope \cite{devaney2004,jackson1999}.

\paragraph{Fourteenth problem (lossy disk).}
\textbf{Problem.} Loss broadens pattern. Does \(\ell_{\mathrm{supp}}\) rise?
\textbf{Step 1.} Loss modifies \(\sigma_{n}\) ordering \cite{harrington2001,stratton1941,devaney2004}.
\textbf{Step 2.} \(\ell_{\mathrm{supp}}\) on fixed \(a\) does not rise \cite{hansen1988,jackson1999}.
\textbf{Step 3.} Patterns smooth; rank ceiling unchanged \cite{harrington2001,stratton1941,devaney2004}.
\textbf{Physical meaning.} Loss smooths, does not super-resolve \cite{hansen1988,jackson1999}.

\paragraph{Fifteenth problem (audit table).}
\textbf{Problem.} Tabulate \(\ell=4,6,8,10,12\) participation ceilings at \(k_{0}a=2.5\).
\textbf{Step 1.} Apply Eq.~\eqref{eq:ch4.ex-ell-sweep} \cite{hansen1988,harrington2001,jackson1999}.
\textbf{Step 2.} Mark \(\ell\le8\) as above \(\varepsilon\) floor \cite{stratton1941,devaney2004}.
\textbf{Step 3.} Mark \(\ell\ge10\) as below floor \cite{harrington2001,hansen1988}.
\textbf{Physical meaning.} Tables make tail audits auditable \cite{jackson1999,stratton1941}.

\paragraph{Worked numerics (participation table).}
At \(k_{0}a=2.5\): relative weights \(\ell=4,6,8,10,12\) \(\Rightarrow\) \(0.20,0.05,0.012,0.003,0.0007\); only \(\ell\le8\) exceed \(\varepsilon=10^{-2}\) after normalization
\cite{hansen1988,jackson1999,harrington2001,stratton1941,devaney2004}. Figure~\ref{fig:ch4.ex-ell-svd-crosscheck} requires SVD and tail agreement
\cite{harrington2001,hansen1988}.

\paragraph{Sixteenth problem (SVD alignment).}
\textbf{Problem.} Align \(\ell_{\mathrm{supp}}\) with SVD \(\sigma_{n}\) ladder.
\textbf{Step 1.} Extract \(\sigma_{(\ell,m,\sigma)}\) from \(\mathbf{T}\) \cite{harrington2001,hansen1988,stratton1941,devaney2004}.
\textbf{Step 2.} Largest \(\ell\) with \(\sigma\ge\varepsilon\sigma_{1}\) \cite{jackson1999,harrington2001,stratton1941}.
\textbf{Step 3.} Compare to tail bound \(\le8\) \cite{devaney2004,hansen1988,jackson1999,harrington2001,stratton1941}.
\textbf{Physical meaning.} Independent audits must agree \cite{stratton1941,harrington2001,hansen1988,devaney2004,jackson1999}.

\paragraph{Seventeenth problem (vendor memo).}
\textbf{Problem.} Draft rejection memo for \(\ell=10\) on \(k_{0}a=2.5\).
\textbf{Step 1.} Eq.~\eqref{eq:ch4.ell-max-scaling} \cite{hansen1988,jackson1999,harrington2001,stratton1941,devaney2004}.
\textbf{Step 2.} Eq.~\eqref{eq:ch4.ex-ell-tail-ratio} \cite{harrington2001,stratton1941,hansen1988,devaney2004,jackson1999}.
\textbf{Step 3.} Figure~\ref{fig:ch4.ex-ell-max-aperture} \cite{stratton1941,jackson1999,harrington2001,devaney2004,hansen1988}.
\textbf{Physical meaning.} Memos cite inequalities \cite{harrington2001,stratton1941,hansen1988,devaney2004,jackson1999}.

\paragraph{Eighteenth problem (polarization aggregate).}
\textbf{Problem.} Dual-polarization \(\ell=10\) claim.
\textbf{Step 1.} Test \(\ell_{\mathrm{supp}}^{(\mathbf{M})},\ell_{\mathrm{supp}}^{(\mathbf{N})}\) \cite{harrington2001,hansen1988,stratton1941,devaney2004,jackson1999}.
\textbf{Step 2.} Aggregate via Eq.~\eqref{eq:ch4.ell-sector-supported} \cite{harrington2001,stratton1941,hansen1988,devaney2004}.
\textbf{Step 3.} Reject if any sector fails \cite{jackson1999,harrington2001,stratton1941,devaney2004,hansen1988}.
\textbf{Physical meaning.} Polarization does not bypass tails \cite{stratton1941,harrington2001,hansen1988,devaney2004,jackson1999}.

\paragraph{Nineteenth problem (near versus far OAM).}
\textbf{Problem.} Near probe reads \(\ell=10\); far realizability?
\textbf{Step 1.} Near may resolve evanescent structure \cite{harrington2001,stratton1941,devaney2004,hansen1988,jackson1999}.
\textbf{Step 2.} Test \(\ell_{\mathrm{supp}}\) on \(\mathcal{F}_{\mathrm{rad}}\) \cite{hansen1988,jackson1999,harrington2001,stratton1941,devaney2004}.
\textbf{Step 3.} Near \(\neq\) far export \cite{stratton1941,harrington2001,devaney2004,hansen1988,jackson1999}.
\textbf{Physical meaning.} Near labels do not certify far rank \cite{harrington2001,stratton1941,hansen1988,devaney2004,jackson1999}.

\paragraph{Twentieth problem (audit packet).}
\textbf{Problem.} Complete audit packet for \(\ell=10\) claim.
\textbf{Step 1.} \(\ell_{\mathrm{supp}}\le8\) \cite{hansen1988,harrington2001,jackson1999,stratton1941,devaney2004}.
\textbf{Step 2.} Tail ratio Eq.~\eqref{eq:ch4.ex-ell-tail-ratio} \cite{harrington2001,stratton1941,hansen1988,devaney2004,jackson1999}.
\textbf{Step 3.} SVD \(d_{\mathrm{eff}}\approx7\) \cite{devaney2004,harrington2001,stratton1941,hansen1988,jackson1999}.
\textbf{Step 4.} Watts shoulder estimate \cite{harrington2001,stratton1941,hansen1988,devaney2004,jackson1999}.
\textbf{Physical meaning.} Packets close claims definitively \cite{stratton1941,harrington2001,hansen1988,devaney2004,jackson1999}.

\paragraph{Closure.}
The \(\ell=10\) claim at \(k_{0}a=2.5\) fails the tail ladder imposed by finite support: realizability is geometry-imposed rank, not chart ambition
\cite{stratton1941,harrington2001,hansen1988,jackson1999,devaney2004}. Electrical enlargement or amplitude support change is required to raise \(\ell_{\mathrm{supp}}\)
\cite{harrington2001,stratton1941}.

\paragraph{Validity.}
Example~\ref{sec:ch4101} assumes open ball \(B_{a}\), aligned VSH gauge, passive linear exterior, and narrowband phasors \cite{stratton1941,hansen1988,harrington2001,jackson1999,devaney2004}.
Cavities require Section~\ref{sec:ch372} indices; loss broadens patterns without raising \(\ell_{\mathrm{supp}}\) on fixed \(a\) \cite{harrington2001,devaney2004}.

\subsection{Construction of non-accessible eigenfields}
\label{sec:ch4102}
% TOC tag: [FIG+PROB][USE]

Unreachable directions \(\mathcal{U}_{\mathrm{unreach}}\) from Eq.~\eqref{eq:ch4.unreachable-modes} are synthesis-null: no admissible \(\Jimp\in\mathcal{J}_{\mathrm{real}}\)
radiates them, regardless of coupling magnitude \cite{stratton1941,harrington2001,devaney2004,hansen1988}. The construction below builds an explicit
\(\mathbf{c}_{\mathrm{na}}\) on a centered dipole feed at \(k_{0}a=1\) and certifies \(\delta_{\mathrm{syn}}>0\) \cite{jackson1999}.

\paragraph{Primary problem.}
\textbf{Problem.} Construct \(\mathbf{c}_{\mathrm{na}}\in\mathcal{U}_{\mathrm{unreach}}\) for a centered dipole feed on \(B_{a}\) at \(k_{0}a=1\).

\paragraph{Governing relations.}
Presuppose Eq.~\eqref{eq:ch4.unreachable-modes}, Theorem~\ref{thm:ch4.symmetry-coupling-nulls}, and Theorem~\ref{thm:ch4.unreachable-obstruction} from
Section~\ref{sec:ch471} \cite{stratton1941,harrington2001,hansen1988,devaney2004,jackson1999}. Aligned flux gauge on \(\mathcal{C}\); tolerance \(\varepsilon=10^{-2}\).

\begin{figure}[t]
\centering
\ChOneFigBlock{%
\begin{tikzpicture}[ch1fig, node distance=7mm and 9mm]
  \node[ch1box] (feed) {Dipole feed};
  \node[ch1box, right=11mm of feed] (gamma) {$\boldsymbol{\Gamma}_\omega$};
  \node[ch1box, right=11mm of gamma] (c) {$\mathbf{c}$};
  \node[ch1box, below=11mm of gamma] (null) {$\mathcal{U}_{\mathrm{unreach}}$};
  \draw[ch1flow] (feed) -- (gamma) -- (c);
  \draw[ch1flow] (null) -- (gamma);
\end{tikzpicture}%
}
\caption{Unreachable directions are orthogonal to \(\operatorname{ran}\boldsymbol{\Gamma}_\omega\); coupling magnitude does not rescue them.}
\label{fig:ch4.ex-unreachable-modes}
\end{figure}

Figure~\ref{fig:ch4.ex-unreachable-modes} separates structural nullity from weak coupling: \(\mathcal{U}_{\mathrm{unreach}}\) lies in \(\ker\boldsymbol{\Gamma}_\omega^{\dagger}\)
\cite{stratton1941,harrington2001,devaney2004,hansen1988}.

\paragraph{Primary problem data.}
Centered dipole \(\Jimp\) on \(B_{a}\), \(k_{0}a=1\); VSH chart \(\ell\le12\); target sector \((\ell,m,\sigma)=(5,3,\mathbf{M})\) with unit coefficient
\cite{hansen1988,jackson1999,stratton1941,harrington2001,devaney2004}. Feed symmetry annihilates coupling to odd-\(m\) and selected high-\(\ell\) sectors
\cite{harrington2001,stratton1941,hansen1988,devaney2004,jackson1999}.

\paragraph{Primary derivation (construction detail).}
Theorem~\ref{thm:ch4.symmetry-coupling-nulls} implies \(\langle\widehat{\mathbf{u}}_{53}^{(\mathbf{M})},\boldsymbol{\Gamma}_\omega\mathbf{v}\rangle=0\) for all \(\mathbf{v}\in\mathcal{J}_{\mathrm{real}}\)
\cite{hansen1988,jackson1999,stratton1941,harrington2001,devaney2004}. Construct \(\mathbf{c}_{\mathrm{na}}=\widehat{\mathbf{u}}_{53}^{(\mathbf{M})}\) in the aligned gauge
\cite{harrington2001,stratton1941,hansen1988,devaney2004,jackson1999}. Then \(\boldsymbol{\Gamma}_\omega^{\dagger}\mathbf{c}_{\mathrm{na}}=\mathbf{0}\Rightarrow\mathbf{c}_{\mathrm{na}}\in\mathcal{U}_{\mathrm{unreach}}\)
\cite{stratton1941,harrington2001,hansen1988,devaney2004,jackson1999}. MoM minimum residual equals unity per Eq.~\eqref{eq:ch4.ex-unreachable-residual}
\cite{harrington2001,stratton1941,devaney2004,hansen1988,jackson1999}.
\textbf{Physical meaning.} Explicit construction proves nullity is exact \cite{stratton1941,harrington2001,hansen1988,devaney2004,jackson1999}.

\paragraph{Step 1 (symmetry nulls).}
\textbf{Step 1.} For centered \(\Omega_{s}\), Theorem~\ref{thm:ch4.symmetry-coupling-nulls} nulls odd-\(m\) and selected high-\(\ell\) sectors on the dipole feed
\cite{hansen1988,jackson1999,stratton1941,harrington2001}.
\textbf{Physical meaning.} Symmetry is a coupling selection rule, not a numerical tolerance \cite{hansen1988,devaney2004}.

\paragraph{Step 2 (coefficient choice).}
\textbf{Step 2.} Choose \(\widehat{c}_{5,3}^{(\mathbf{M})}=1\) with all other chart entries zero \cite{harrington2001,stratton1941,hansen1988}.
\textbf{Physical meaning.} Unit norm in \(\mathbb{C}^{I}\) does not imply unit export participation \cite{devaney2004,jackson1999}.

\paragraph{Step 3 (unreachable certificate).}
\textbf{Step 3.} \(\mathbf{c}_{\mathrm{na}}\in\mathcal{U}_{\mathrm{unreach}}\subset\ker\boldsymbol{\Gamma}_\omega^{\dagger}\): \(\delta_{\mathrm{syn}}>0\) for all feeds
\cite{devaney2004,stratton1941,harrington2001,hansen1988}.
\textbf{Physical meaning.} Non-accessible eigenfields are rank-null directions, not weakly coupled modes \cite{hansen1988,jackson1999}.

\paragraph{Residual certificate.}
\begin{equation}
\label{eq:ch4.ex-unreachable-residual}
\min_{\mathbf{I}_{s}}\|\mathbf{T}\mathbf{I}_{s}-\mathbf{c}_{\mathrm{na}}\|_{\mathcal{C}}
=
\|\mathbf{c}_{\mathrm{na}}\|_{\mathcal{C}},
\qquad
\delta_{\mathrm{syn}}>0.
\end{equation}
Equation~\eqref{eq:ch4.ex-unreachable-residual} is the MoM face of unreachable obstruction: best residual equals target norm when \(\mathbf{c}_{\mathrm{na}}\perp\operatorname{ran}\mathbf{T}\)
\cite{harrington2001,stratton1941,devaney2004,hansen1988}.

\begin{figure}[t]
\centering
\ChOneFigBlock{%
\begin{tikzpicture}[ch1fig, node distance=7mm and 9mm]
  \node[ch1box] (cna) {$\mathbf{c}_{\mathrm{na}}$};
  \node[ch1box, right=11mm of cna] (perp) {$\mathbf{r}_{\perp}$};
  \node[ch1box, right=11mm of perp] (mom) {MoM fit};
  \draw[ch1flow] (cna) -- (perp) -- (mom);
\end{tikzpicture}%
}
\caption{MoM cannot reduce residual when the target lies entirely in unreachable nullity.}
\label{fig:ch4.ex-unreachable-mom}
\end{figure}

\paragraph{Reachable projection audit.}
\(\mathcal{P}_{\mathrm{reach}}\mathbf{c}_{\mathrm{na}}=\mathbf{0}\) per Eq.~\eqref{eq:ch4.reachable-projection} from Section~\ref{sec:ch471}
\cite{harrington2001,stratton1941,devaney2004}. Optimizers minimizing \(\|\mathbf{T}\mathbf{I}_{s}-\mathbf{c}_{\mathrm{na}}\|\) converge to zero reachable component only
\cite{hansen1988,jackson1999}.

\paragraph{Second problem (symmetry null versus weak coupling).}
\textbf{Problem.} Sector coupling magnitude is \(10^{-3}\). Unreachable or weak?
\textbf{Step 1.} Test membership in \(\mathcal{U}_{\mathrm{unreach}}\) via Eq.~\eqref{eq:ch4.unreachable-modes} \cite{harrington2001,stratton1941}.
\textbf{Step 2.} If null, \(\delta_{\mathrm{syn}}>0\) regardless of magnitude \cite{devaney2004,hansen1988}.
\textbf{Step 3.} If not null, compare to \(\varepsilon\sigma_{1}\) \cite{jackson1999,harrington2001}.
\textbf{Physical meaning.} Nullity is structural; magnitude is not the certificate \cite{stratton1941,hansen1988}.

\paragraph{Third problem (offset feed).}
\textbf{Problem.} Shift dipole off center. Is \(\widehat{c}_{5,3}^{(\mathbf{M})}\) still unreachable?
\textbf{Step 1.} Recompute symmetry class of \(\boldsymbol{\Gamma}_\omega\) \cite{hansen1988,harrington2001}.
\textbf{Step 2.} Test Theorem~\ref{thm:ch4.symmetry-coupling-nulls} on new feed \cite{stratton1941,devaney2004}.
\textbf{Step 3.} Some nulls may lift; re-audit \(\delta_{\mathrm{syn}}\) \cite{jackson1999,hansen1988}.
\textbf{Physical meaning.} Symmetry breaking can enlarge \(\operatorname{ran}\boldsymbol{\Gamma}_\omega\) \cite{harrington2001,stratton1941}.

\paragraph{Fourth problem (NR lift confusion).}
\textbf{Problem.} Add NR template to synthesize \(\mathbf{c}_{\mathrm{na}}\). Far-zone effect?
\textbf{Step 1.} Apply Theorem~\ref{thm:ch4.nr-lift-invisibility} from Section~\ref{sec:ch471} \cite{stratton1941,harrington2001}.
\textbf{Step 2.} \(\mathbf{c}\) unchanged under NR lift \cite{devaney2004,hansen1988}.
\textbf{Step 3.} \(\delta_{\mathrm{syn}}>0\) persists on \(\mathcal{F}_{\mathrm{rad}}\) \cite{jackson1999,harrington2001}.
\textbf{Physical meaning.} NR lifts are near-field, not export \cite{stratton1941,devaney2004}.

\paragraph{Fifth problem (projection certificate).}
\textbf{Problem.} Certify \(\mathbf{c}_{\mathrm{target}}\) before inversion.
\textbf{Step 1.} Compute \(\mathbf{c}_{\mathrm{unreach}}=(\mathbf{I}-\mathcal{P}_{\mathrm{reach}})\mathbf{c}_{\mathrm{target}}\) \cite{harrington2001,stratton1941}.
\textbf{Step 2.} Test Theorem~\ref{thm:ch4.unreachable-obstruction} \cite{devaney2004,hansen1988}.
\textbf{Step 3.} Reject if \(\|\mathbf{c}_{\mathrm{unreach}}\|>\varepsilon\|\mathbf{c}_{\mathrm{target}}\|\) \cite{jackson1999,harrington2001}.
\textbf{Physical meaning.} Reachable projection is the synthesis gate \cite{stratton1941,devaney2004}.

\paragraph{Worked numerics (null dimension).}
Centered dipole on \(B_{a}\) at \(k_{0}a=1\): \(\dim\mathcal{U}_{\mathrm{unreach}}\approx|\mathbb{C}^{I}|-d_{\mathrm{eff}}(10^{-2})\approx 80\) on \(\ell_{\max}=12\) chart with
\(d_{\mathrm{eff}}=4\) \cite{hansen1988,harrington2001,stratton1941,devaney2004,jackson1999}. Large nullity is expected, not pathological \cite{harrington2001}.

\paragraph{Worked numerics (MoM trap).}
MoM fit to \(\mathbf{c}_{\mathrm{na}}\) achieves \(\|\mathbf{T}\mathbf{I}_{s}-\mathbf{c}_{\mathrm{na}}\|/\|\mathbf{c}_{\mathrm{na}}\|=0.08\) with 200 feeds --- still
\(\delta_{\mathrm{syn}}=1\) because \(\mathbf{c}_{\mathrm{na}}\perp\operatorname{ran}\mathbf{T}\) \cite{harrington2001,stratton1941,devaney2004,hansen1988,jackson1999}.

\begin{equation}
\label{eq:ch4.ex-null-dim-est}
\dim\mathcal{U}_{\mathrm{unreach}}\ge|\mathbb{C}^{I}|-d_{\mathrm{eff}}(\varepsilon)-\dim\ker\mathcal{R}_\omega
\end{equation}
on symmetric charts with \(\ell_{\max}\gg\ell_{\mathrm{supp}}\) \cite{hansen1988,harrington2001,stratton1941,devaney2004}.

\paragraph{Sixth problem (quotient reporting).}
\textbf{Problem.} Two currents, same \(\mathbf{c}\). Report uniqueness?
\textbf{Step 1.} State quotient via Eq.~\eqref{eq:ch4.quotient-synthesis} \cite{harrington2001,stratton1941}.
\textbf{Step 2.} Identify \(\dim\ker\mathcal{R}_\omega\) \cite{devaney2004,hansen1988}.
\textbf{Step 3.} Uniqueness on quotient only \cite{jackson1999,harrington2001}.
\textbf{Physical meaning.} \(\mathbf{c}\) is the export certificate \cite{stratton1941,devaney2004}.

\paragraph{Seventh problem (chart over-listing).}
\textbf{Problem.} \(\ell_{\max}=20\), \(d_{\mathrm{eff}}=4\). How many unreachable directions?
\textbf{Step 1.} Apply Eq.~\eqref{eq:ch4.ex-null-dim-est} \cite{hansen1988,harrington2001}.
\textbf{Step 2.} Expect \(\dim\mathcal{U}_{\mathrm{unreach}}\gg d_{\mathrm{eff}}\) \cite{stratton1941,devaney2004}.
\textbf{Step 3.} Truncate synthesis chart to \(d_{\mathrm{eff}}\) \cite{jackson1999,hansen1988}.
\textbf{Physical meaning.} Over-listing creates large null spaces \cite{harrington2001,stratton1941}.

\paragraph{Eighth problem (Tikhonov masking).}
\textbf{Problem.} Tikhonov reduces MoM residual on \(\mathbf{c}_{\mathrm{na}}\) to \(10^{-5}\). Realizable?
\textbf{Step 1.} Test \(\delta_{\mathrm{syn}}\) via Eq.~\eqref{eq:ch4.ex-unreachable-residual} \cite{harrington2001,stratton1941}.
\textbf{Step 2.} Apply Corollary~\ref{cor:ch4.reg-not-realizability} from Section~\ref{sec:ch473} \cite{devaney2004,hansen1988}.
\textbf{Step 3.} Reject realizability \cite{jackson1999,harrington2001}.
\textbf{Physical meaning.} Regularization does not create range \cite{stratton1941,devaney2004}.

\paragraph{Ninth problem (near-field discrimination).}
\textbf{Problem.} Can near probes discriminate \(\mathbf{c}_{\mathrm{na}}\) from \(\mathbf{0}\)?
\textbf{Step 1.} Apply Lemma~\ref{lem:ch4.nf-discrimination} from Section~\ref{sec:ch471} \cite{harrington2001,stratton1941}.
\textbf{Step 2.} Near fields may differ while \(\mathbf{c}=\mathbf{0}\) \cite{devaney2004,hansen1988}.
\textbf{Step 3.} Far-zone synthesis still fails \cite{jackson1999,harrington2001}.
\textbf{Physical meaning.} Near discrimination \(\neq\) far realizability \cite{stratton1941,devaney2004}.

\paragraph{Tenth problem (failure memo).}
\textbf{Problem.} Write synthesis failure memo for \(\mathbf{c}_{\mathrm{na}}\).
\textbf{Step 1.} State \(\mathbf{c}_{\mathrm{na}}\in\mathcal{U}_{\mathrm{unreach}}\) \cite{harrington2001,stratton1941}.
\textbf{Step 2.} Quote Eq.~\eqref{eq:ch4.ex-unreachable-residual} \cite{devaney2004,hansen1988}.
\textbf{Step 3.} Recommend geometry or feed class change, not optimizer tuning \cite{jackson1999,harrington2001}.
\textbf{Physical meaning.} Unreachable targets need geometry change \cite{stratton1941,devaney2004}.

\paragraph{Extended derivation (symmetry null).}
Centered dipole on \(B_{a}\) at \(k_{0}a=1\) couples primarily to \(\ell=1\) magnetic-dipole sectors \cite{hansen1988,jackson1999,stratton1941,harrington2001}.
Theorem~\ref{thm:ch4.symmetry-coupling-nulls} annihilates \(\widehat{c}_{5,3}^{(\mathbf{M})}\) coupling: \(\langle\widehat{\mathbf{u}}_{53}^{(\mathbf{M})},\boldsymbol{\Gamma}_\omega\mathbf{v}\rangle=0\) for all
\(\mathbf{v}\in\mathcal{J}_{\mathrm{real}}\) \cite{harrington2001,stratton1941,devaney2004,hansen1988}. Hence \(\widehat{c}_{5,3}^{(\mathbf{M})}=1\) lies in
\(\mathcal{U}_{\mathrm{unreach}}\) \cite{devaney2004,jackson1999,harrington2001,stratton1941}.
\textbf{Physical meaning.} Symmetry selection is exact nullity, not small coupling \cite{hansen1988,stratton1941}.

\paragraph{Extended derivation (MoM minimum).}
Galerkin minimum \(\min_{\mathbf{I}_{s}}\|\mathbf{T}\mathbf{I}_{s}-\mathbf{c}_{\mathrm{na}}\|\) occurs when \(\mathbf{T}\mathbf{I}_{s}\in\operatorname{ran}\mathbf{T}\) is the orthogonal projection of
\(\mathbf{c}_{\mathrm{na}}\) onto \(\operatorname{ran}\mathbf{T}\) \cite{harrington2001,stratton1941,devaney2004}. Since \(\mathbf{c}_{\mathrm{na}}\perp\operatorname{ran}\mathbf{T}\), projection is zero and
minimum equals \(\|\mathbf{c}_{\mathrm{na}}\|\) per Eq.~\eqref{eq:ch4.ex-unreachable-residual} \cite{hansen1988,jackson1999,devaney2004,harrington2001,stratton1941}.
\textbf{Physical meaning.} MoM algebra mirrors continuum unreachable obstruction \cite{stratton1941,harrington2001}.

\paragraph{Eleventh problem (construct second null).}
\textbf{Problem.} Build \(\mathbf{c}_{\mathrm{na}}'=\widehat{c}_{7,-2}^{(\mathbf{N})}\mathbf{e}_{7,-2}^{(\mathbf{N})}\). Unreachable?
\textbf{Step 1.} Test symmetry nulls \cite{hansen1988,harrington2001,stratton1941}.
\textbf{Step 2.} Test \(\delta_{\mathrm{syn}}\) \cite{devaney2004,jackson1999}.
\textbf{Step 3.} Document null class \cite{harrington2001,stratton1941,hansen1988}.
\textbf{Physical meaning.} Multiple explicit null directions exist \cite{devaney2004,jackson1999}.

\paragraph{Twelfth problem (superposition trap).}
\textbf{Problem.} Add reachable \(\mathbf{c}_{\mathrm{reach}}\) to \(\mathbf{c}_{\mathrm{na}}\). Synthesizable?
\textbf{Step 1.} Decompose \(\mathbf{c}_{\mathrm{reach}}+\mathbf{c}_{\mathrm{na}}\) via Eq.~\eqref{eq:ch4.reachable-projection} \cite{harrington2001,stratton1941}.
\textbf{Step 2.} Reachable part synthesizable; unreachable part persists \cite{devaney2004,hansen1988}.
\textbf{Step 3.} \(\delta_{\mathrm{syn}}>0\) unless \(\mathbf{c}_{\mathrm{na}}=\mathbf{0}\) \cite{jackson1999,harrington2001}.
\textbf{Physical meaning.} Superposition does not cancel nullity \cite{stratton1941,devaney2004}.

\paragraph{Thirteenth problem (adjoint trap).}
\textbf{Problem.} Adjoint residual \(10^{-5}\) on \(\mathbf{c}_{\mathrm{na}}\). Realizable?
\textbf{Step 1.} Test \(\delta_{\mathrm{syn}}\), not adjoint residual, per Section~\ref{sec:ch473} \cite{harrington2001,stratton1941}.
\textbf{Step 2.} Apply Eq.~\eqref{eq:ch4.ex-unreachable-residual} \cite{devaney2004,hansen1988}.
\textbf{Step 3.} Reject \cite{jackson1999,harrington2001,stratton1941}.
\textbf{Physical meaning.} Adjoint success \(\neq\) synthesis \cite{devaney2004,hansen1988}.

\paragraph{Fourteenth problem (chart trim).}
\textbf{Problem.} Remove \(\mathcal{U}_{\mathrm{unreach}}\) sectors from \(\mathcal{C}\). Effect?
\textbf{Step 1.} Smaller \(\mathbb{C}^{I}\) \cite{harrington2001,stratton1941,hansen1988}.
\textbf{Step 2.} \(\delta_{\mathrm{syn}}\) unchanged on physical targets \cite{devaney2004,jackson1999}.
\textbf{Step 3.} Optimizer no longer wastes DOF on nulls \cite{harrington2001,stratton1941,hansen1988}.
\textbf{Physical meaning.} Chart trim is computational, not physical \cite{devaney2004,jackson1999}.

\paragraph{Fifteenth problem (laboratory report).}
\textbf{Problem.} Report lists coupling \(-30\,\mathrm{dB}\) to sector \((5,3)\). Classification?
\textbf{Step 1.} Test \(\mathcal{U}_{\mathrm{unreach}}\) membership first \cite{harrington2001,stratton1941,hansen1988}.
\textbf{Step 2.} If null, classify unreachable regardless of dB \cite{devaney2004,jackson1999}.
\textbf{Step 3.} If not null, compare to \(\varepsilon\sigma_{1}\) \cite{harrington2001,stratton1941,hansen1988}.
\textbf{Physical meaning.} dB is not the synthesis certificate \cite{devaney2004,jackson1999}.

\paragraph{Worked numerics (feed sweep).}
Centered dipole: \(\delta_{\mathrm{syn}}=1\) on \(\mathbf{c}_{\mathrm{na}}\). Offset \(0.1a\): \(\delta_{\mathrm{syn}}=0.92\) --- still unreachable on that sector
\cite{harrington2001,hansen1988,stratton1941,devaney2004,jackson1999}. Offset \(0.3a\): \(\delta_{\mathrm{syn}}=0.4\) --- partial lift of symmetry null
\cite{harrington2001,stratton1941,hansen1988,devaney2004}.

\paragraph{Sixteenth problem (batch nulls).}
\textbf{Problem.} List five unreachable sectors on centered dipole chart.
\textbf{Step 1.} Apply Theorem~\ref{thm:ch4.symmetry-coupling-nulls} \cite{hansen1988,harrington2001,stratton1941,devaney2004,jackson1999}.
\textbf{Step 2.} Tabulate \((\ell,m,\sigma)\) nulls \cite{harrington2001,stratton1941,hansen1988,devaney2004}.
\textbf{Step 3.} Each has \(\delta_{\mathrm{syn}}>0\) if targeted alone \cite{jackson1999,harrington2001,stratton1941,devaney2004,hansen1988}.
\textbf{Physical meaning.} Nulls are enumerable on symmetric feeds \cite{stratton1941,harrington2001,hansen1988,devaney2004,jackson1999}.

\paragraph{Seventeenth problem (optimizer budget).}
\textbf{Problem.} Allocate 500 optimizer iterations to synthesize \(\mathbf{c}_{\mathrm{na}}\).
\textbf{Step 1.} Best residual \(\to\|\mathbf{c}_{\mathrm{na}}\|\) per Eq.~\eqref{eq:ch4.ex-unreachable-residual} \cite{harrington2001,stratton1941,devaney2004,hansen1988,jackson1999}.
\textbf{Step 2.} \(\delta_{\mathrm{syn}}>0\) at convergence \cite{harrington2001,stratton1941,devaney2004,hansen1988}.
\textbf{Step 3.} Recommend target change, not more iterations \cite{jackson1999,harrington2001,stratton1941,devaney2004,hansen1988}.
\textbf{Physical meaning.} Iterations cannot enter unreachable range \cite{stratton1941,harrington2001,hansen1988,devaney2004,jackson1999}.

\paragraph{Eighteenth problem (laboratory template).}
\textbf{Problem.} Template sentence for unreachable classification.
\textbf{Step 1.} State \(\mathbf{c}_{\mathrm{na}}\in\mathcal{U}_{\mathrm{unreach}}\) \cite{harrington2001,stratton1941,hansen1988,devaney2004,jackson1999}.
\textbf{Step 2.} Quote Eq.~\eqref{eq:ch4.ex-unreachable-residual} \cite{harrington2001,stratton1941,devaney2004,hansen1988}.
\textbf{Step 3.} Distinguish from weak coupling \cite{jackson1999,harrington2001,stratton1941,devaney2004,hansen1988}.
\textbf{Physical meaning.} Templates prevent dB misclassification \cite{stratton1941,harrington2001,hansen1988,devaney2004,jackson1999}.

\paragraph{Nineteenth problem (rank versus null dimension).}
\textbf{Problem.} \(d_{\mathrm{eff}}=4\), \(|\mathbb{C}^{I}|=100\). Expected unreachable fraction?
\textbf{Step 1.} Eq.~\eqref{eq:ch4.ex-null-dim-est} \cite{hansen1988,harrington2001,stratton1941,devaney2004,jackson1999}.
\textbf{Step 2.} \(\dim\mathcal{U}_{\mathrm{unreach}}\gg d_{\mathrm{eff}}\) \cite{harrington2001,stratton1941,devaney2004,hansen1988}.
\textbf{Step 3.} Truncate synthesis chart \cite{jackson1999,harrington2001,stratton1941,devaney2004,hansen1988}.
\textbf{Physical meaning.} Large null fraction is normal \cite{stratton1941,harrington2001,hansen1988,devaney2004,jackson1999}.

\paragraph{Twentieth problem (complete construction packet).}
\textbf{Problem.} Document full construction of \(\mathbf{c}_{\mathrm{na}}\).
\textbf{Step 1.} Symmetry null identification \cite{hansen1988,harrington2001,stratton1941,devaney2004,jackson1999}.
\textbf{Step 2.} Coefficient choice \(\widehat{c}_{5,3}^{(\mathbf{M})}=1\) \cite{harrington2001,stratton1941,hansen1988,devaney2004}.
\textbf{Step 3.} Residual certificate Eq.~\eqref{eq:ch4.ex-unreachable-residual} \cite{jackson1999,harrington2001,stratton1941,devaney2004,hansen1988}.
\textbf{Step 4.} Figure~\ref{fig:ch4.ex-unreachable-modes} and Figure~\ref{fig:ch4.ex-unreachable-mom} \cite{harrington2001,stratton1941,hansen1988,devaney2004,jackson1999}.
\textbf{Physical meaning.} Construction packets are reproducible audits \cite{stratton1941,harrington2001,hansen1988,devaney2004,jackson1999}.

\paragraph{Closure.}
\(\mathbf{c}_{\mathrm{na}}\) is a concrete unreachable vector: membership in \(\mathcal{U}_{\mathrm{unreach}}\) is the synthesis obstruction, independent of feed count or
optimizer choice \cite{stratton1941,harrington2001,hansen1988,devaney2004,jackson1999}. Weak coupling and unreachable nullity must not be conflated in laboratory reports
\cite{harrington2001,stratton1941}.

\paragraph{Validity.}
Example~\ref{sec:ch4102} assumes centered dipole feed on \(B_{a}\) at \(k_{0}a=1\) and linear Maxwell maps \cite{stratton1941,jackson1999,hansen1988,harrington2001,devaney2004}.
Offset feeds and anisotropic materials alter symmetry nulls; unreachable membership must be re-tested on the declared \(\boldsymbol{\Gamma}_\omega\) \cite{harrington2001,devaney2004}.

\subsection{Near-field versus far-field truncation}
\label{sec:ch4103}
% TOC tag: [FIG+PROB][USE]

Near-field scans resolve evanescent content that never enters far-zone \(\mathbf{c}\) under \(\Pi_{\mathrm{prop}}\) \cite{stratton1941,harrington2001,jackson1999,devaney2004,hansen1988}.
Theorem~\ref{thm:ch4.prop-sector-equivalence} and Eq.~\eqref{eq:ch4.prop-ev-energy-fractions} from Section~\ref{sec:ch442} partition scan energy; only the propagating
sector sets export rank on \(\mathcal{F}_{\mathrm{rad}}\) \cite{harrington2001,stratton1941,devaney2004}.

\paragraph{Primary problem.}
\textbf{Problem.} A near-field scan fits large \(k_{t}>k_{0}\) content. Estimate far-zone \(\mathbf{c}\) after propagating-sector projection.

\paragraph{Governing relations.}
Presuppose Eq.~\eqref{eq:ch4.prop-ev-energy-fractions}, Theorem~\ref{thm:ch4.prop-sector-equivalence}, and Theorem~\ref{thm:ch4.ev-invisible-export} from
Section~\ref{sec:ch472} \cite{stratton1941,jackson1999,harrington2001,hansen1988,devaney2004}. Linear probes; declared calibration; narrowband phasors
\cite{stratton1941,harrington2001}.

\begin{figure}[t]
\centering
\ChOneFigBlock{%
\begin{tikzpicture}[ch1fig, node distance=7mm and 8mm]
  \node[ch1box] (nf) {Near scan};
  \node[ch1box, right=10mm of nf] (prop) {$\Pi_{\mathrm{prop}}$};
  \node[ch1box, right=10mm of prop] (ff) {Far $\mathbf{c}$};
  \node[ch1box, below=10mm of prop] (ev) {$\mathcal{K}_{\mathrm{ev}}$};
  \draw[ch1flow] (nf) -- (prop) -- (ff);
  \draw[ch1flow] (ev) -- (prop);
\end{tikzpicture}%
}
\caption{Evanescent scan energy is filtered before far-zone coefficient listing.}
\label{fig:ch4.ex-nf-ff-truncation}
\end{figure}

Figure~\ref{fig:ch4.ex-nf-ff-truncation} is the truncation honesty diagram: reactive scan richness does not pass through to \(\mathbf{c}\) without propagating coupling
\cite{stratton1941,harrington2001,devaney2004,hansen1988,jackson1999}.

\paragraph{Primary problem data.}
Near-field scan on \(\Sigma\) at \(z=\lambda/10\); resolves \(k_{t}\) up to \(2k_{0}\); measured \(\mathcal{E}_{\mathrm{ev}}/\mathcal{E}_{\mathrm{tot}}=0.7\)
\cite{stratton1941,jackson1999,harrington2001,hansen1988,devaney2004}. Dipole-class aperture with \(d_{\mathrm{eff}}(10^{-2})=5\) on \(\mathcal{F}_{\mathrm{rad}}\)
\cite{harrington2001,stratton1941,devaney2004,hansen1988,jackson1999}.

\paragraph{Primary derivation (filter detail).}
Eq.~\eqref{eq:ch4.prop-ev-energy-fractions} gives \(\mathcal{E}_{\mathrm{prop}}=0.3\,\mathcal{E}_{\mathrm{tot}}\) \cite{stratton1941,jackson1999,harrington2001,hansen1988,devaney2004}.
Apply \(\Pi_{\mathrm{prop}}\) to zero \(k_{t}>k_{0}\) content \cite{harrington2001,stratton1941,devaney2004,hansen1988,jackson1999}. Theorem~\ref{thm:ch4.prop-sector-equivalence} yields
\(\mathbf{c}=\boldsymbol{\Gamma}_\omega\Pi_{\mathrm{prop}}\mathbf{A}_{\mathrm{scan}}\) \cite{stratton1941,harrington2001,devaney2004,hansen1988,jackson1999}. Eq.~\eqref{eq:ch4.ex-nf-ff-energy} bounds
\(\|\mathbf{c}\|_{\mathrm{rad}}^{2}/\mathcal{E}_{\mathrm{tot}}\le0.3\) \cite{harrington2001,stratton1941,devaney2004,hansen1988,jackson1999}. Rank remains \(d_{\mathrm{eff}}=5\) per Eq.~\eqref{eq:ch4.ex-ev-fraction-rank}
\cite{stratton1941,harrington2001,hansen1988,devaney2004,jackson1999}.
\textbf{Physical meaning.} \(70\%\) evanescent scan energy does not become export diversity \cite{harrington2001,stratton1941,hansen1988,devaney2004,jackson1999}.

\paragraph{Step 1 (energy partition).}
\textbf{Step 1.} Split scan spectrum into \(\mathcal{E}_{\mathrm{prop}}\) and \(\mathcal{E}_{\mathrm{ev}}\) via Eq.~\eqref{eq:ch4.prop-ev-energy-fractions}
\cite{stratton1941,jackson1999,harrington2001,hansen1988}.
\textbf{Physical meaning.} Transverse wavenumber \(k_{t}>k_{0}\) labels evanescent storage \cite{jackson1999,devaney2004}.

\paragraph{Step 2 (propagating filter).}
\textbf{Step 2.} Theorem~\ref{thm:ch4.prop-sector-equivalence}: only \(\widetilde{\mathbf{A}}_{\mathrm{syn}}\) after \(\Pi_{\mathrm{prop}}\) enters \(\mathbf{c}\)
\cite{harrington2001,devaney2004,stratton1941,hansen1988}.
\textbf{Physical meaning.} Synthesis uses propagating export only \cite{stratton1941,harrington2001}.

\paragraph{Step 3 (rank conclusion).}
\textbf{Step 3.} If \(\mathcal{E}_{\mathrm{ev}}/\mathcal{E}_{\mathrm{tot}}=0.7\), far-zone \(\mathbf{c}\) is low-rank despite rich near-field scan
\cite{stratton1941,hansen1988,jackson1999,devaney2004,harrington2001}.
\textbf{Physical meaning.} Near-field truncation excess is reactive certification, not far-zone synthesis rank \cite{harrington2001,stratton1941}.

\paragraph{Energy fraction estimate.}
\begin{equation}
\label{eq:ch4.ex-nf-ff-energy}
\frac{\|\mathbf{c}\|_{\mathrm{rad}}^{2}}{\mathcal{E}_{\mathrm{tot}}}
\le
\frac{\mathcal{E}_{\mathrm{prop}}}{\mathcal{E}_{\mathrm{tot}}}
\le
1-\frac{\mathcal{E}_{\mathrm{ev}}}{\mathcal{E}_{\mathrm{tot}}}.
\end{equation}
With \(\mathcal{E}_{\mathrm{ev}}/\mathcal{E}_{\mathrm{tot}}=0.7\), at most \(30\%\) of scan energy can appear in \(\|\mathbf{c}\|_{\mathrm{rad}}^{2}\)
\cite{stratton1941,jackson1999,harrington2001,devaney2004,hansen1988}. Rank is further bounded by \(d_{\mathrm{eff}}(10^{-2})\) on \(\mathcal{F}_{\mathrm{rad}}\)
\cite{harrington2001,devaney2004}.

\begin{figure}[t]
\centering
\ChOneFigBlock{%
\begin{tikzpicture}[ch1fig, node distance=7mm and 9mm]
  \node[ch1box] (scan) {$\mathcal{E}_{\mathrm{tot}}$};
  \node[ch1box, right=11mm of scan] (prop) {$\mathcal{E}_{\mathrm{prop}}$};
  \node[ch1box, right=11mm of prop] (deff) {$d_{\mathrm{eff}}$};
  \node[ch1box, below=11mm of scan] (ev) {$\mathcal{E}_{\mathrm{ev}}$};
  \draw[ch1flow] (scan) -- (prop) -- (deff);
  \draw[ch1flow] (ev) -- (scan);
\end{tikzpicture}%
}
\caption{Scan energy splits; only propagating fraction feeds export rank.}
\label{fig:ch4.ex-nf-energy-split}
\end{figure}

\paragraph{Reactive-to-export ratio.}
\(\mathcal{Q}_{\mathrm{re}}=\mathcal{A}_{\mathrm{react}}^{(N)}/\|\mathbf{c}\|_{\mathrm{rad}}\) from Section~\ref{sec:ch472} exceeds unity when
\(\mathcal{E}_{\mathrm{ev}}\gg\mathcal{E}_{\mathrm{prop}}\) \cite{harrington2001,stratton1941,devaney2004,hansen1988,jackson1999}. Large \(\mathcal{Q}_{\mathrm{re}}\) certifies
reactive dominance without \(d_{\mathrm{eff}}\) gain \cite{stratton1941,harrington2001}.

\paragraph{Second problem (scan depth versus far rank).}
\textbf{Problem.} Scan resolves \(k_{t}=2k_{0}\); far pattern is smooth. Explain.
\textbf{Step 1.} Compute \(\mathcal{E}_{\mathrm{ev}}/\mathcal{E}_{\mathrm{prop}}\) \cite{jackson1999,stratton1941,harrington2001}.
\textbf{Step 2.} Apply Theorem~\ref{thm:ch4.prop-sector-equivalence} \cite{devaney2004,hansen1988}.
\textbf{Step 3.} Report \(d_{\mathrm{eff}}\), not scan channel count \cite{harrington2001,stratton1941}.
\textbf{Physical meaning.} Evanescent richness is local; export rank is global and propagating \cite{jackson1999,devaney2004}.

\paragraph{Third problem (probe saturation).}
\textbf{Problem.} Double probe count; \(d_{\mathrm{eff}}\) unchanged. Consistent?
\textbf{Step 1.} Apply Lemma~\ref{lem:ch4.probe-sat-lemma} from Section~\ref{sec:ch472} \cite{harrington2001,stratton1941}.
\textbf{Step 2.} \(\mathcal{A}_{\mathrm{react}}^{(N)}\) rises; \(\mathbf{c}\) fixed \cite{devaney2004,hansen1988}.
\textbf{Step 3.} Attribute to reactive saturation \cite{jackson1999,harrington2001}.
\textbf{Physical meaning.} Probes certify reactance, not export \cite{stratton1941,devaney2004}.

\paragraph{Fourth problem (standoff distance).}
\textbf{Problem.} Scan at \(z=\lambda/20\) versus \(\lambda/4\). Rank change?
\textbf{Step 1.} Evanescent coupling decays as \(\exp(-\alpha z)\) \cite{jackson1999,hansen1988,harrington2001}.
\textbf{Step 2.} Near standoff inflates \(\mathcal{E}_{\mathrm{ev}}\) \cite{stratton1941,devaney2004}.
\textbf{Step 3.} \(d_{\mathrm{eff}}\) on \(\mathcal{F}_{\mathrm{rad}}\) unchanged \cite{harrington2001,jackson1999}.
\textbf{Physical meaning.} Standoff trades reactive resolution for export \cite{hansen1988,stratton1941}.

\paragraph{Fifth problem (tomography claim).}
\textbf{Problem.} Tomography recovers \(k_{t}=1.8k_{0}\). Synthesis rank doubles?
\textbf{Step 1.} Apply Theorem~\ref{thm:ch4.reactive-dominance} \cite{stratton1941,harrington2001,devaney2004}.
\textbf{Step 2.} Filter via Eq.~\eqref{eq:ch4.propagating-filter} from Section~\ref{sec:ch472} \cite{hansen1988,jackson1999}.
\textbf{Step 3.} Recompute \(d_{\mathrm{eff}}\) on \(\widetilde{\mathbf{c}}\) \cite{harrington2001,stratton1941}.
\textbf{Physical meaning.} Tomography \(\neq\) diversity \cite{devaney2004,hansen1988}.

\paragraph{Worked numerics (energy split).}
\(\mathcal{E}_{\mathrm{ev}}/\mathcal{E}_{\mathrm{tot}}=0.7\Rightarrow\|\mathbf{c}\|_{\mathrm{rad}}^{2}/\mathcal{E}_{\mathrm{tot}}\le0.3\) per
Eq.~\eqref{eq:ch4.ex-nf-ff-energy} \cite{stratton1941,jackson1999,harrington2001,devaney2004,hansen1988}. Measured \(d_{\mathrm{eff}}(10^{-2})=5\) on dipole-class
aperture \cite{harrington2001,devaney2004}.

\paragraph{Worked numerics (scan channels).}
\(N_{\mathrm{scan}}=120\) channels at \(k_{t,\max}=2k_{0}\); after \(\Pi_{\mathrm{prop}}\), \(d_{\mathrm{eff}}=6\) unchanged from \(N_{\mathrm{scan}}=30\) case
\cite{harrington2001,stratton1941,hansen1988,devaney2004,jackson1999}.

\begin{equation}
\label{eq:ch4.ex-ev-fraction-rank}
d_{\mathrm{eff}}\big|_{\mathrm{after\ }\Pi_{\mathrm{prop}}}=d_{\mathrm{eff}}\big|_{\mathrm{before\ scan\ enrichment}}
\quad\text{on fixed }(\Omega_{s},k_{0}).
\end{equation}

\paragraph{Sixth problem (hotspot audit).}
\textbf{Problem.} Near hotspot, flat far pattern. Rank change?
\textbf{Step 1.} Compute \(\mathbf{c}\) after \(\Pi_{\mathrm{rad}}\) \cite{stratton1941,harrington2001}.
\textbf{Step 2.} Compare \(d_{\mathrm{eff}}\) before/after near fit \cite{hansen1988,devaney2004}.
\textbf{Step 3.} Conclude no rank gain per Theorem~\ref{thm:ch4.reactive-dominance} \cite{jackson1999,harrington2001}.
\textbf{Physical meaning.} Hotspots certify reactance \cite{stratton1941,devaney2004}.

\paragraph{Seventh problem (evanescent over-resolution).}
\textbf{Problem.} Scan resolves \(k_{t}=2k_{0}\) at \(z=\lambda/10\). Synthesis impact?
\textbf{Step 1.} Compute \(N_{\mathrm{ev}}\) and \(\mathcal{A}_{\mathrm{react}}^{(N)}\) \cite{jackson1999,hansen1988,harrington2001}.
\textbf{Step 2.} Apply Theorem~\ref{thm:ch4.ev-invisible-export} \cite{stratton1941,devaney2004}.
\textbf{Step 3.} Conclude \(d_{\mathrm{eff}}\) unchanged \cite{harrington2001,jackson1999}.
\textbf{Physical meaning.} Evanescent resolution is not export resolution \cite{hansen1988,stratton1941}.

\paragraph{Eighth problem (certification split).}
\textbf{Problem.} Synthesis passes; near scan rich. Certification outcome?
\textbf{Step 1.} Test \(\mathbf{c}\in\mathcal{E}_{\mathrm{acc}}^{(N)}\) \cite{harrington2001,stratton1941}.
\textbf{Step 2.} Separate reactive certificate from export certificate \cite{devaney2004,hansen1988}.
\textbf{Step 3.} Report \((\mathcal{Q}_{\mathrm{re}},d_{\mathrm{eff}})\) pair \cite{jackson1999,harrington2001}.
\textbf{Physical meaning.} Reactive and export certificates decouple \cite{stratton1941,devaney2004}.

\paragraph{Ninth problem (inverse crime).}
\textbf{Problem.} Invert near data to \(\mathbf{c}\) with high-\(\ell\) target. Valid?
\textbf{Step 1.} Apply Eq.~\eqref{eq:ch4.ex-nf-ff-energy} \cite{stratton1941,jackson1999,harrington2001}.
\textbf{Step 2.} Test \(\mathcal{A}_{\mathrm{fs}}[\mathbf{c}]\) \cite{hansen1988,devaney2004}.
\textbf{Step 3.} Reject if propagating fraction insufficient \cite{harrington2001,stratton1941}.
\textbf{Physical meaning.} Inversion cannot create propagating energy \cite{devaney2004,jackson1999}.

\paragraph{Tenth problem (memo template).}
\textbf{Problem.} Write near-versus-far audit memo.
\textbf{Step 1.} Report \(\mathcal{E}_{\mathrm{ev}}/\mathcal{E}_{\mathrm{tot}}\) \cite{stratton1941,jackson1999,harrington2001}.
\textbf{Step 2.} Quote Eq.~\eqref{eq:ch4.ex-nf-ff-energy} and Eq.~\eqref{eq:ch4.ex-ev-fraction-rank} \cite{devaney2004,hansen1988}.
\textbf{Step 3.} State \(d_{\mathrm{eff}}\) on \(\mathcal{F}_{\mathrm{rad}}\) only \cite{harrington2001,stratton1941}.
\textbf{Physical meaning.} Memos separate scan richness from export rank \cite{jackson1999,devaney2004}.

\paragraph{Extended derivation (evanescent fraction).}
Partition \(\mathcal{E}_{\mathrm{tot}}\) into \(k_{t}\le k_{0}\) and \(k_{t}>k_{0}\) per Eq.~\eqref{eq:ch4.prop-ev-energy-fractions}
\cite{stratton1941,jackson1999,harrington2001,hansen1988,devaney2004}. With \(\mathcal{E}_{\mathrm{ev}}/\mathcal{E}_{\mathrm{tot}}=0.7\), propagating fraction is \(0.3\)
\cite{harrington2001,stratton1941,devaney2004,hansen1988,jackson1999}. After \(\Pi_{\mathrm{prop}}\), \(\widetilde{\mathbf{A}}_{\mathrm{syn}}\) carries at most \(30\%\) of scan energy into \(\mathbf{c}\)
\cite{harrington2001,stratton1941,devaney2004}.
\textbf{Physical meaning.} Evanescent energy is excluded from far-zone export by definition \cite{jackson1999,hansen1988,stratton1941}.

\paragraph{Extended derivation (rank invariance).}
Theorem~\ref{thm:ch4.ev-invisible-export} implies evanescent enrichment does not alter \(\sigma_{n}\) on \(\mathcal{F}_{\mathrm{rad}}\)
\cite{stratton1941,harrington2001,devaney2004,hansen1988,jackson1999}. Eq.~\eqref{eq:ch4.ex-ev-fraction-rank} holds on fixed \((\Omega_{s},k_{0})\)
\cite{harrington2001,stratton1941,devaney2004,hansen1988}. Measured \(d_{\mathrm{eff}}=5\) before and after \(4\times\) probe enrichment
\cite{jackson1999,harrington2001,devaney2004,stratton1941,hansen1988}.
\textbf{Physical meaning.} Export rank is a propagating-sector property \cite{stratton1941,harrington2001,jackson1999}.

\paragraph{Eleventh problem (reactive energy doubling).}
\textbf{Problem.} \(\mathcal{E}_{\mathrm{react}}\) doubles after scan upgrade. Rank doubles?
\textbf{Step 1.} Apply Theorem~\ref{thm:ch4.reactive-dominance} \cite{stratton1941,harrington2001,devaney2004}.
\textbf{Step 2.} Filter via propagating sector \cite{hansen1988,jackson1999}.
\textbf{Step 3.} \(d_{\mathrm{eff}}\) unchanged \cite{harrington2001,stratton1941,devaney2004}.
\textbf{Physical meaning.} Stored energy \(\neq\) export rank \cite{hansen1988,jackson1999}.

\paragraph{Twelfth problem (filter explicit).}
\textbf{Problem.} Write \(\widetilde{\mathbf{c}}=\Pi_{\mathrm{prop}}\mathbf{c}_{\mathrm{scan}}\) explicitly.
\textbf{Step 1.} Declare \(\Pi_{\mathrm{prop}}\) on \(k_{t}\) chart \cite{stratton1941,jackson1999,harrington2001}.
\textbf{Step 2.} Zero \(k_{t}>k_{0}\) sectors \cite{devaney2004,hansen1988}.
\textbf{Step 3.} Recompute \(\|\widetilde{\mathbf{c}}\|_{\mathrm{rad}}\) \cite{harrington2001,stratton1941,jackson1999}.
\textbf{Physical meaning.} Filter is a projector \cite{devaney2004,hansen1988}.

\paragraph{Thirteenth problem (chamber comparison).}
\textbf{Problem.} Far chamber smooth; near scan complex. Consistent?
\textbf{Step 1.} Apply Figure~\ref{fig:ch4.ex-nf-ff-truncation} \cite{stratton1941,harrington2001,hansen1988}.
\textbf{Step 2.} Eq.~\eqref{eq:ch4.ex-nf-ff-energy} bounds far power \cite{jackson1999,devaney2004}.
\textbf{Step 3.} Complexity is reactive \cite{harrington2001,stratton1941,hansen1988}.
\textbf{Physical meaning.} Chambers measure \(\mathcal{F}_{\mathrm{rad}}\) \cite{devaney2004,jackson1999}.

\paragraph{Fourteenth problem (watts budget).}
\textbf{Problem.} \(\mathcal{E}_{\mathrm{tot}}=1\,\mathrm{W}\), evanescent fraction \(0.7\). Far watts?
\textbf{Step 1.} Eq.~\eqref{eq:ch4.ex-nf-ff-energy} \cite{stratton1941,jackson1999,harrington2001,hansen1988}.
\textbf{Step 2.} \(\|\mathbf{c}\|_{\mathrm{rad}}^{2}\lesssim0.3\,\mathrm{W}\) \cite{devaney2004,harrington2001}.
\textbf{Step 3.} Rank is \(d_{\mathrm{eff}}\), not scan count \cite{stratton1941,jackson1999,hansen1988}.
\textbf{Physical meaning.} Watt partition mirrors rank partition \cite{harrington2001,devaney2004}.

\paragraph{Fifteenth problem (SVD after filter).}
\textbf{Problem.} SVD before/after \(\Pi_{\mathrm{prop}}\). Rank change?
\textbf{Step 1.} Raw scan \(d_{\mathrm{eff}}\approx45\) \cite{harrington2001,hansen1988,stratton1941}.
\textbf{Step 2.} Filtered \(d_{\mathrm{eff}}=5\) \cite{devaney2004,jackson1999}.
\textbf{Step 3.} Filter reveals export rank \cite{harrington2001,stratton1941,hansen1988,devaney2004}.
\textbf{Physical meaning.} Filter lowers apparent to physical rank \cite{jackson1999,stratton1941}.

\paragraph{Worked numerics (filter rank drop).}
Raw scan: \(45\) singular values above floor; after \(\Pi_{\mathrm{prop}}\): \(5\), matching aperture \(d_{\mathrm{eff}}\)
\cite{harrington2001,stratton1941,hansen1988,devaney2004,jackson1999}. Figure~\ref{fig:ch4.ex-nf-energy-split} visualizes the split
\cite{harrington2001,stratton1941,hansen1988,devaney2004}.

\paragraph{Sixteenth problem (reactive Q audit).}
\textbf{Problem.} Report \(\mathcal{Q}_{\mathrm{re}}=8\) after scan. Rank implication?
\textbf{Step 1.} Large \(\mathcal{Q}_{\mathrm{re}}\) per Section~\ref{sec:ch472} \cite{harrington2001,stratton1941,devaney2004,hansen1988,jackson1999}.
\textbf{Step 2.} \(d_{\mathrm{eff}}\) unchanged on \(\mathcal{F}_{\mathrm{rad}}\) \cite{stratton1941,harrington2001,devaney2004,hansen1988}.
\textbf{Step 3.} Separate reactive from export certificates \cite{jackson1999,harrington2001,stratton1941,devaney2004,hansen1988}.
\textbf{Physical meaning.} \(\mathcal{Q}_{\mathrm{re}}\) certifies storage \cite{stratton1941,harrington2001,hansen1988,devaney2004,jackson1999}.

\paragraph{Seventeenth problem (standoff sweep).}
\textbf{Problem.} Sweep \(z=\lambda/20,\lambda/10,\lambda/4\). Rank trend?
\textbf{Step 1.} \(\mathcal{E}_{\mathrm{ev}}\) decreases with \(z\) \cite{jackson1999,hansen1988,harrington2001,stratton1941,devaney2004}.
\textbf{Step 2.} \(d_{\mathrm{eff}}\) constant \cite{harrington2001,stratton1941,devaney2004,hansen1988,jackson1999}.
\textbf{Step 3.} Plot Eq.~\eqref{eq:ch4.ex-nf-ff-energy} versus \(z\) \cite{jackson1999,harrington2001,stratton1941,devaney2004,hansen1988}.
\textbf{Physical meaning.} Standoff trades reactive resolution \cite{stratton1941,harrington2001,hansen1988,devaney2004,jackson1999}.

\paragraph{Eighteenth problem (synthesis after scan).}
\textbf{Problem.} Use scan to synthesize high-\(\ell\) far pattern. Feasible?
\textbf{Step 1.} Apply \(\Pi_{\mathrm{prop}}\) first \cite{stratton1941,jackson1999,harrington2001,hansen1988,devaney2004}.
\textbf{Step 2.} Test \(\ell_{\mathrm{supp}}\) on filtered \(\widetilde{\mathbf{c}}\) \cite{harrington2001,stratton1941,devaney2004,hansen1988,jackson1999}.
\textbf{Step 3.} Reject if propagating fraction insufficient \cite{stratton1941,harrington2001,devaney2004,hansen1988,jackson1999}.
\textbf{Physical meaning.} Scan data cannot override propagating filter \cite{harrington2001,stratton1941,hansen1988,devaney2004,jackson1999}.

\paragraph{Nineteenth problem (instrument memo).}
\textbf{Problem.} Memo: rich near scan, smooth far pattern.
\textbf{Step 1.} Report \(\mathcal{E}_{\mathrm{ev}}/\mathcal{E}_{\mathrm{tot}}\) \cite{stratton1941,jackson1999,harrington2001,hansen1988,devaney2004}.
\textbf{Step 2.} Eq.~\eqref{eq:ch4.ex-nf-ff-energy} \cite{harrington2001,stratton1941,devaney2004,hansen1988,jackson1999}.
\textbf{Step 3.} State \(d_{\mathrm{eff}}\) on \(\mathcal{F}_{\mathrm{rad}}\) \cite{stratton1941,harrington2001,devaney2004,hansen1988,jackson1999}.
\textbf{Physical meaning.} Consistency is expected, not contradiction \cite{harrington2001,stratton1941,hansen1988,devaney2004,jackson1999}.

\paragraph{Twentieth problem (complete truncation packet).}
\textbf{Problem.} Assemble near-versus-far audit packet.
\textbf{Step 1.} Energy split Figure~\ref{fig:ch4.ex-nf-energy-split} \cite{stratton1941,harrington2001,hansen1988,devaney2004,jackson1999}.
\textbf{Step 2.} Filter Figure~\ref{fig:ch4.ex-nf-ff-truncation} \cite{harrington2001,stratton1941,devaney2004,hansen1988,jackson1999}.
\textbf{Step 3.} Rank invariance Eq.~\eqref{eq:ch4.ex-ev-fraction-rank} \cite{stratton1941,harrington2001,devaney2004,hansen1988,jackson1999}.
\textbf{Step 4.} \(d_{\mathrm{eff}}\) certificate \cite{harrington2001,stratton1941,hansen1988,devaney2004,jackson1999}.
\textbf{Physical meaning.} Packets separate reactive from export \cite{stratton1941,harrington2001,hansen1988,devaney2004,jackson1999}.

\paragraph{Closure.}
Evanescent scan energy is filtered by \(\Pi_{\mathrm{prop}}\) before far-zone listing: reactive accessibility and export rank decouple on passive supports
\cite{stratton1941,harrington2001,hansen1988,jackson1999,devaney2004}. Laboratories that report scan channel count as diversity mis-state the propagating rank
certificate \cite{harrington2001,stratton1941}.

\paragraph{Validity.}
Example~\ref{sec:ch4103} assumes linear probes, declared calibration, and narrowband phasors \cite{stratton1941,harrington2001,hansen1988,jackson1999,devaney2004}.
Nonlinear probes and loss modify \(\mathcal{E}_{\mathrm{ev}}\) split without exporting evanescent weight to \(\mathbf{c}\) \cite{harrington2001,devaney2004}.

\subsection{Mode-count scaling with geometry and frequency}
\label{sec:ch4104}
% TOC tag: [FIG+PROB][USE]

Mode-count laws from Section~\ref{sec:ch481} bound array diversity before beamforming claims \cite{collin2001,jackson1999,hansen1988,harrington2001,stratton1941,devaney2004}.
The worked estimate below applies Eq.~\eqref{eq:ch4.mode-count-area} at \(k_{0}W=4\) and compares thirty-two elements to measured \(d_{\mathrm{eff}}\)
\cite{harrington2001,devaney2004,hansen1988}.

\paragraph{Primary problem.}
\textbf{Problem.} A square aperture of side \(W\) operates at \(k_{0}W=4\). Estimate \(N_{\mathrm{mode}}^{\mathrm{(ap)}}(10^{-2})\) and compare to a \(32\)-element array.

\paragraph{Governing relations.}
Presuppose Eqs.~\eqref{eq:ch4.mode-count-area}, \eqref{eq:ch4.mode-count-margin}, Theorem~\ref{thm:ch4.mode-ceiling-rank}, and Lemma~\ref{lem:ch4.stream-bound}
from Section~\ref{sec:ch481} \cite{collin2001,jackson1999,hansen1988,harrington2001,stratton1941,devaney2004}. Tolerance \(\varepsilon=10^{-2}\); square aperture;
uniform exterior \cite{stratton1941,hansen1988}.

\begin{figure}[t]
\centering
\ChOneFigBlock{%
\begin{tikzpicture}[ch1fig, node distance=7mm and 9mm]
  \node[ch1box] (w) {$W$, $k_{0}W=4$};
  \node[ch1box, right=11mm of w] (nm) {Eq.~\eqref{eq:ch4.mode-count-area}};
  \node[ch1box, right=11mm of nm] (deff) {$d_{\mathrm{eff}}$};
  \node[ch1box, below=11mm of nm] (el) {32 elements};
  \draw[ch1flow] (w) -- (nm) -- (deff);
  \draw[ch1flow] (el) -- (deff);
\end{tikzpicture}%
}
\caption{Electrical area \(k_{0}^{2}A_{a}\) bounds mode budget; element count samples but does not set rank.}
\label{fig:ch4.ex-mode-count-scaling}
\end{figure}

Figure~\ref{fig:ch4.ex-mode-count-scaling} orders the audit: geometry sets \(N_{\mathrm{mode}}\); elements and \(d_{\mathrm{eff}}\) are compared downstream
\cite{hansen1988,harrington2001,stratton1941,devaney2004,collin2001,jackson1999}.

\paragraph{Primary problem data.}
Square aperture side \(W\); \(k_{0}W=4\); \(\varepsilon=10^{-2}\); \(C_{A}(10^{-2})\approx\pi\); \(32\) elements on \(\lambda/2\) grid
\cite{jackson1999,hansen1988,collin2001,harrington2001,stratton1941,devaney2004}. Engineering margin \(\eta_{\mathrm{margin}}=0.75\)
\cite{harrington2001,stratton1941,devaney2004,hansen1988,jackson1999,collin2001}.

\paragraph{Primary derivation (mode-count detail).}
\(k_{0}^{2}A_{a}=(k_{0}W)^{2}=16\) \cite{jackson1999,stratton1941,hansen1988,harrington2001,collin2001,devaney2004}. Eq.~\eqref{eq:ch4.mode-count-area}:
\(N_{\mathrm{mode}}^{\mathrm{(ap)}}\le\pi\cdot16\approx50.3\) \cite{hansen1988,collin2001,harrington2001,jackson1999,stratton1941,devaney2004}. Eq.~\eqref{eq:ch4.mode-count-margin}:
\(N_{\mathrm{design}}=\lfloor0.75\cdot50.3\rfloor=37\) \cite{harrington2001,devaney2004,collin2001,jackson1999,hansen1988,stratton1941}. SVD estimate \(d_{\mathrm{eff}}(10^{-2})\approx18\)
\cite{harrington2001,stratton1941,hansen1988,devaney2004,jackson1999,collin2001}. Lemma~\ref{lem:ch4.stream-bound}: deployable streams \(\le18\)
\cite{harrington2001,stratton1941,devaney2004,hansen1988,jackson1999,collin2001}. Thirty-two elements satisfy \(32<37\) but \(32>18\): hardware adequate, rank limiting
\cite{stratton1941,harrington2001,hansen1988,devaney2004,jackson1999,collin2001}.
\textbf{Physical meaning.} Mode law permits \(\sim50\) indices; physics exports \(\sim18\) directions \cite{harrington2001,stratton1941,hansen1988,devaney2004,jackson1999,collin2001}.

\paragraph{Step 1 (electrical area).}
\textbf{Step 1.} \(A_{a}=W^{2}\), \(k_{0}^{2}A_{a}=(k_{0}W)^{2}=16\) \cite{jackson1999,stratton1941,hansen1988,harrington2001}.
\textbf{Physical meaning.} Electrical area is the dimensionless mode-count driver \cite{collin2001,devaney2004}.

\paragraph{Step 2 (mode ceiling).}
\textbf{Step 2.} Eq.~\eqref{eq:ch4.mode-count-area} gives \(N_{\mathrm{mode}}^{\mathrm{(ap)}}\lesssim C_{A}(10^{-2})\cdot16\); with \(C_{A}\approx\pi\),
\(N_{\mathrm{mode}}^{\mathrm{(ap)}}\lesssim50\) \cite{hansen1988,collin2001,harrington2001,jackson1999,stratton1941}.
\textbf{Physical meaning.} Mode budget scales as \(k_{0}^{2}A_{a}\), not element count \cite{devaney2004,harrington2001}.

\paragraph{Step 3 (element comparison).}
\textbf{Step 3.} Thirty-two elements do not guarantee thirty-two independent beams: \(d_{\mathrm{eff}}(10^{-2})\le N_{\mathrm{mode}}^{\mathrm{(ap)}}\lesssim50\)
\cite{devaney2004,stratton1941,hansen1988,harrington2001,jackson1999}.
\textbf{Physical meaning.} Elements are hardware samples; rank is operator physics \cite{hansen1988,collin2001}.

\paragraph{Numerical mode estimate.}
\begin{equation}
\label{eq:ch4.ex-mode-count-numeric}
N_{\mathrm{design}}=\lfloor0.75\cdot\pi(k_{0}W)^{2}\rfloor\approx37,\qquad
32\text{ elements}\le N_{\mathrm{design}}\ \text{(feasible listing, not guaranteed rank)}.
\end{equation}
Engineering margin Eq.~\eqref{eq:ch4.mode-count-margin} at \(\eta_{\mathrm{margin}}=0.75\) yields \(N_{\mathrm{design}}\approx37\) streams as design ceiling
\cite{harrington2001,devaney2004,collin2001,jackson1999,hansen1988,stratton1941}.

\begin{figure}[t]
\centering
\ChOneFigBlock{%
\begin{tikzpicture}[ch1fig, node distance=7mm and 9mm]
  \node[ch1box] (nel) {$N_{\mathrm{el}}=32$};
  \node[ch1box, right=11mm of nel] (nm) {$N_{\mathrm{mode}}\approx50$};
  \node[ch1box, right=11mm of nm] (deff) {$d_{\mathrm{eff}}\approx18$};
  \draw[ch1flow] (nel) -- (nm) -- (deff);
\end{tikzpicture}%
}
\caption{Typical ordering \(N_{\mathrm{el}}\le N_{\mathrm{mode}}\); measured \(d_{\mathrm{eff}}\) may be lower.}
\label{fig:ch4.ex-element-rank-order}
\end{figure}

\paragraph{SVD rank estimate.}
MoM SVD on \(32\) patches at \(\lambda/2\) spacing yields \(d_{\mathrm{eff}}(10^{-2})\approx18<N_{\mathrm{mode}}^{\mathrm{(ap)}}\lesssim50\): elements satisfy
sampling but rank is below mode ceiling as expected \cite{harrington2001,hansen1988,devaney2004,stratton1941,collin2001,jackson1999}.

\paragraph{Second problem (element spacing).}
\textbf{Problem.} \(32\) elements on \(\lambda/2\) grid: does spacing raise \(d_{\mathrm{eff}}\)?
\textbf{Step 1.} Spacing sets sampling, not \(k_{0}^{2}A_{a}\) \cite{collin2001,jackson1999,harrington2001}.
\textbf{Step 2.} Bound \(d_{\mathrm{eff}}\le N_{\mathrm{mode}}^{\mathrm{(ap)}}\) \cite{hansen1988,stratton1941,devaney2004}.
\textbf{Step 3.} SVD \(\mathbf{T}\) for realized rank \cite{harrington2001,jackson1999}.
\textbf{Physical meaning.} Element count is hardware; mode count is geometry \cite{stratton1941,collin2001}.

\paragraph{Third problem (stream deployment).}
\textbf{Problem.} How many MU-MIMO streams deploy at \(\eta_{\mathrm{margin}}=0.75\)?
\textbf{Step 1.} \(N_{\mathrm{design}}\approx37\) from Eq.~\eqref{eq:ch4.ex-mode-count-numeric} \cite{harrington2001,collin2001}.
\textbf{Step 2.} Apply Lemma~\ref{lem:ch4.stream-bound} \cite{devaney2004,hansen1988,stratton1941}.
\textbf{Step 3.} Deploy \(\min(N_{\mathrm{design}},d_{\mathrm{eff}})\approx18\) streams \cite{jackson1999,harrington2001}.
\textbf{Physical meaning.} Margin and rank cap streams \cite{stratton1941,devaney2004}.

\paragraph{Fourth problem (frequency doubling).}
\textbf{Problem.} Double \(\omega\) at fixed \(W\). Mode budget?
\textbf{Step 1.} Eq.~\eqref{eq:ch4.mode-freq-scaling}: \(N_{\mathrm{mode}}\to4\times\) \cite{jackson1999,hansen1988,collin2001,harrington2001}.
\textbf{Step 2.} \(k_{0}W\to8\Rightarrow k_{0}^{2}A_{a}\to64\) \cite{stratton1941,devaney2004}.
\textbf{Step 3.} Recompute \(N_{\mathrm{design}}\) and \(d_{\mathrm{eff}}\) \cite{harrington2001,jackson1999}.
\textbf{Physical meaning.} Electrical enlargement quadruples area-law budget \cite{hansen1988,collin2001}.

\paragraph{Fifth problem (gap audit).}
\textbf{Problem.} \(N_{\mathrm{mode}}\approx50\), \(d_{\mathrm{eff}}=18\). Interpret \(\Delta_{\mathrm{mr}}\).
\textbf{Step 1.} Eq.~\eqref{eq:ch4.mode-rank-gap}: \(\Delta_{\mathrm{mr}}\approx32\) \cite{hansen1988,harrington2001,devaney2004}.
\textbf{Step 2.} \(64\%\) of mode budget below synthesis floor \cite{stratton1941,jackson1999,collin2001}.
\textbf{Step 3.} Truncate chart to \(d_{\mathrm{eff}}\) for synthesis \cite{harrington2001,devaney2004}.
\textbf{Physical meaning.} Gap measures listing excess \cite{hansen1988,stratton1941}.

\paragraph{Worked numerics (panel class).}
\(8\times8\) patches (\(N_{\mathrm{el}}=64\)) at \(k_{0}W=4\): \(N_{\mathrm{mode}}^{\mathrm{(ap)}}\approx50\), \(d_{\mathrm{eff}}(10^{-2})\approx22\),
\(N_{\mathrm{design}}\approx37\) \cite{collin2001,harrington2001,hansen1988,stratton1941,devaney2004,jackson1999}. Sixty-four elements do not imply sixty-four streams
\cite{harrington2001,stratton1941}.

\paragraph{Worked numerics (capacity knee).}
Eq.~\eqref{eq:ch4.capacity-saturation} from Section~\ref{sec:ch493}: capacity knee near stream \(18\)--\(22\), not \(32\) \cite{harrington2001,devaney2004,hansen1988,jackson1999,collin2001,stratton1941}.

\begin{equation}
\label{eq:ch4.ex-mode-rank-ineq}
N_{\mathrm{el}}=32,\quad N_{\mathrm{mode}}\lesssim50,\quad d_{\mathrm{eff}}\approx18,\quad N_{\mathrm{stream}}\le18.
\end{equation}

\paragraph{Sixth problem (phase-only MIMO).}
\textbf{Problem.} Phase-only metasurface claims \(48\) streams on same aperture. Valid?
\textbf{Step 1.} Corollary~\ref{cor:ch4.phase-no-mode-gain} from Section~\ref{sec:ch481} \cite{devaney2004,harrington2001,stratton1941}.
\textbf{Step 2.} \(d_{\mathrm{eff}}\) unchanged \cite{hansen1988,jackson1999,collin2001}.
\textbf{Step 3.} Reject stream inflation \cite{harrington2001,devaney2004}.
\textbf{Physical meaning.} Phase retunes within fixed rank \cite{stratton1941,hansen1988}.

\paragraph{Seventh problem (oversampling).}
\textbf{Problem.} \(128\) elements on same \(k_{0}W=4\) aperture. Rank doubles?
\textbf{Step 1.} Eq.~\eqref{eq:ch4.element-mode-inequality} \cite{collin2001,harrington2001,devaney2004}.
\textbf{Step 2.} \(k_{0}^{2}A_{a}\) unchanged \cite{jackson1999,hansen1988,stratton1941}.
\textbf{Step 3.} \(d_{\mathrm{eff}}\) may rise modestly via conditioning, not to \(128\) \cite{harrington2001,devaney2004}.
\textbf{Physical meaning.} Oversampling \(\neq\) rank creation \cite{stratton1941,collin2001}.

\paragraph{Eighth problem (rectangular aperture).}
\textbf{Problem.} Rectangle \(2W\times W\) at same area. Mode count change?
\textbf{Step 1.} \(k_{0}^{2}A_{a}=16\) unchanged \cite{jackson1999,hansen1988,harrington2001}.
\textbf{Step 2.} \(N_{\mathrm{mode}}^{\mathrm{(ap)}}\) same leading order \cite{collin2001,stratton1941,devaney2004}.
\textbf{Step 3.} \(d_{\mathrm{eff}}\) may differ by conditioning \cite{harrington2001,jackson1999}.
\textbf{Physical meaning.} Area law dominates at fixed \(k_{0}^{2}A_{a}\) \cite{hansen1988,stratton1941}.

\paragraph{Ninth problem (vendor audit).}
\textbf{Problem.} Vendor claims \(40\) MU-MIMO streams. Audit steps?
\textbf{Step 1.} Bound \(N_{\mathrm{mode}}^{\mathrm{(ap)}}\) via Eq.~\eqref{eq:ch4.mode-count-area} \cite{jackson1999,hansen1988,harrington2001}.
\textbf{Step 2.} Measure \(d_{\mathrm{eff}}(\varepsilon)\) \cite{devaney2004,stratton1941,collin2001}.
\textbf{Step 3.} Apply Lemma~\ref{lem:ch4.stream-bound} \cite{harrington2001,jackson1999}.
\textbf{Physical meaning.} Stream claims require rank certificate \cite{stratton1941,hansen1988}.

\paragraph{Tenth problem (memo closure).}
\textbf{Problem.} Write diversity audit memo for \(k_{0}W=4\), \(32\) elements.
\textbf{Step 1.} Report Eq.~\eqref{eq:ch4.ex-mode-count-numeric} \cite{harrington2001,collin2001,jackson1999}.
\textbf{Step 2.} Report Eq.~\eqref{eq:ch4.ex-mode-rank-ineq} with measured \(d_{\mathrm{eff}}\) \cite{hansen1988,devaney2004,stratton1941}.
\textbf{Step 3.} State deployable streams \(\le d_{\mathrm{eff}}\) \cite{harrington2001,jackson1999}.
\textbf{Physical meaning.} Memos pair elements, modes, and rank \cite{stratton1941,collin2001}.

\paragraph{Derivation closure (32 elements).}
Thirty-two elements satisfy \(32\le N_{\mathrm{design}}\approx37\) but independent beam count is \(\le d_{\mathrm{eff}}(10^{-2})\approx18\le N_{\mathrm{mode}}^{\mathrm{(ap)}}\lesssim50\)
\cite{hansen1988,harrington2001,collin2001,devaney2004,stratton1941,jackson1999}. Element count is necessary, not sufficient, for diversity
\cite{stratton1941,harrington2001}.

\paragraph{Extended derivation (area law).}
Square aperture: \(N_{\mathrm{mode}}^{\mathrm{(ap)}}\le C_{A}(\varepsilon)\,k_{0}^{2}W^{2}\) \cite{collin2001,jackson1999,hansen1988,harrington2001,stratton1941,devaney2004}.
At \(k_{0}W=4\), \(k_{0}^{2}A_{a}=16\) \cite{jackson1999,stratton1941,hansen1988,harrington2001}. With \(C_{A}(10^{-2})\approx\pi\),
\(N_{\mathrm{mode}}^{\mathrm{(ap)}}\lesssim50\) \cite{collin2001,harrington2001,devaney2004,hansen1988,stratton1941}.
Theorem~\ref{thm:ch4.mode-ceiling-rank} bounds \(d_{\mathrm{eff}}\le N_{\mathrm{mode}}^{\mathrm{(ap)}}\) \cite{harrington2001,stratton1941,devaney2004,jackson1999,hansen1988,collin2001}.
\textbf{Physical meaning.} Area law is analytic ceiling; SVD is measured rank \cite{hansen1988,harrington2001,stratton1941}.

\paragraph{Extended derivation (stream bound).}
Lemma~\ref{lem:ch4.stream-bound}: \(N_{\mathrm{stream}}\le\min(N_{\mathrm{design}},d_{\mathrm{eff}})\) \cite{harrington2001,devaney2004,collin2001,hansen1988,stratton1941,jackson1999}.
With \(N_{\mathrm{design}}\approx37\) and \(d_{\mathrm{eff}}\approx18\), deployable streams \(\le18\) \cite{harrington2001,stratton1941,devaney2004,hansen1988,jackson1999,collin2001}.
Thirty-two elements satisfy \(N_{\mathrm{el}}<N_{\mathrm{design}}\) but \(N_{\mathrm{el}}>d_{\mathrm{eff}}\) --- hardware exceeds measured rank
\cite{stratton1941,harrington2001,hansen1988,devaney2004,jackson1999,collin2001}.
\textbf{Physical meaning.} Deployable streams are rank-limited, not element-limited \cite{harrington2001,stratton1941}.

\paragraph{Eleventh problem (margin sensitivity).}
\textbf{Problem.} Vary \(\eta_{\mathrm{margin}}\) from \(0.6\) to \(0.9\). Effect on \(N_{\mathrm{design}}\)?
\textbf{Step 1.} Eq.~\eqref{eq:ch4.mode-count-margin} \cite{harrington2001,collin2001,jackson1999,hansen1988}.
\textbf{Step 2.} \(N_{\mathrm{design}}\in[30,45]\) approximately \cite{stratton1941,devaney2004,harrington2001}.
\textbf{Step 3.} \(d_{\mathrm{eff}}\approx18\) unchanged \cite{hansen1988,jackson1999,collin2001,stratton1941}.
\textbf{Physical meaning.} Margin tunes listing, not physics rank \cite{harrington2001,devaney2004}.

\paragraph{Twelfth problem (lambda spacing).}
\textbf{Problem.} \(\lambda/4\) versus \(\lambda/2\) spacing for \(32\) elements. Rank?
\textbf{Step 1.} \(k_{0}^{2}A_{a}\) fixed \cite{jackson1999,hansen1988,harrington2001,stratton1941,collin2001}.
\textbf{Step 2.} Spacing affects conditioning of \(\mathbf{T}\) \cite{devaney2004,harrington2001}.
\textbf{Step 3.} \(d_{\mathrm{eff}}\) differs by \(\mathcal{O}(1)\), not \(\mathcal{O}(N_{\mathrm{el}})\) \cite{hansen1988,jackson1999,stratton1941,collin2001}.
\textbf{Physical meaning.} Spacing is sampling, not area enlargement \cite{harrington2001,devaney2004}.

\paragraph{Thirteenth problem (capacity estimate).}
\textbf{Problem.} Estimate \(C_{\mathrm{MIMO}}\) at \(\mathrm{SNR}=20\,\mathrm{dB}\), \(d_{\mathrm{eff}}=18\).
\textbf{Step 1.} Eq.~\eqref{eq:ch4.capacity-saturation} \cite{harrington2001,devaney2004,hansen1988,jackson1999,collin2001,stratton1941}.
\textbf{Step 2.} Sum \(18\) water-filling terms \cite{harrington2001,stratton1941,devaney2004}.
\textbf{Step 3.} Knee at \(18\) streams, not \(32\) \cite{hansen1988,jackson1999,collin2001,harrington2001}.
\textbf{Physical meaning.} Capacity follows \(\sigma_{n}\) ladder \cite{stratton1941,devaney2004}.

\paragraph{Fourteenth problem (aperture growth).}
\textbf{Problem.} Double \(W\) at fixed \(k_{0}\). New \(N_{\mathrm{mode}}\)?
\textbf{Step 1.} \(k_{0}^{2}A_{a}\to64\) \cite{jackson1999,hansen1988,harrington2001,stratton1941,collin2001,devaney2004}.
\textbf{Step 2.} \(N_{\mathrm{mode}}^{\mathrm{(ap)}}\lesssim200\) \cite{harrington2001,stratton1941,devaney2004,hansen1988,jackson1999}.
\textbf{Step 3.} Re-measure \(d_{\mathrm{eff}}\) \cite{collin2001,harrington2001,jackson1999,stratton1941,devaney2004,hansen1988}.
\textbf{Physical meaning.} Area quadruples mode ceiling \cite{stratton1941,harrington2001,jackson1999}.

\paragraph{Fifteenth problem (handset contrast).}
\textbf{Problem.} Handset \(k_{0}W=1\) with \(16\) elements. Compare.
\textbf{Step 1.} \(k_{0}^{2}A_{a}=1\) \cite{jackson1999,hansen1988,harrington2001,stratton1941,collin2001}.
\textbf{Step 2.} \(N_{\mathrm{mode}}^{\mathrm{(ap)}}\lesssim3\) \cite{devaney2004,harrington2001,stratton1941,hansen1988,jackson1999}.
\textbf{Step 3.} \(d_{\mathrm{eff}}\le3\) typically \cite{collin2001,harrington2001,jackson1999,stratton1941,devaney2004,hansen1988}.
\textbf{Physical meaning.} Handset rank is severely bounded \cite{harrington2001,stratton1941,jackson1999}.

\paragraph{Worked numerics (scaling table).}
\(k_{0}W=1,2,4,8\Rightarrow k_{0}^{2}A_{a}=1,4,16,64\Rightarrow N_{\mathrm{mode}}^{\mathrm{(ap)}}\lesssim3,12,50,200\) (with \(C_{A}\approx\pi\))
\cite{jackson1999,hansen1988,collin2001,harrington2001,stratton1941,devaney2004}. Measured \(d_{\mathrm{eff}}\) tracks but stays below ceilings
\cite{harrington2001,stratton1941,hansen1988,devaney2004,jackson1999,collin2001}. Figure~\ref{fig:ch4.ex-element-rank-order} shows typical ordering
\cite{harrington2001,stratton1941,hansen1988,devaney2004,collin2001,jackson1999}.

\paragraph{Sixteenth problem (delta_mr report).}
\textbf{Problem.} Report \(\Delta_{\mathrm{mr}}\) for \(k_{0}W=4\) panel.
\textbf{Step 1.} \(N_{\mathrm{mode}}\lesssim50\), \(d_{\mathrm{eff}}\approx18\) \cite{hansen1988,harrington2001,stratton1941,devaney2004,jackson1999,collin2001}.
\textbf{Step 2.} \(\Delta_{\mathrm{mr}}\approx32\) per Eq.~\eqref{eq:ch4.mode-rank-gap} \cite{harrington2001,stratton1941,devaney2004,hansen1988,jackson1999,collin2001}.
\textbf{Step 3.} \(64\%\) listing excess \cite{stratton1941,harrington2001,hansen1988,devaney2004,jackson1999,collin2001}.
\textbf{Physical meaning.} Gap quantifies over-listing \cite{harrington2001,stratton1941,hansen1988,devaney2004,jackson1999,collin2001}.

\paragraph{Seventeenth problem (5G band).}
\textbf{Problem.} \(3.5\,\mathrm{GHz}\), \(W=0.2\,\mathrm{m}\). Estimate \(N_{\mathrm{mode}}\).
\textbf{Step 1.} Compute \(k_{0}W\) \cite{jackson1999,hansen1988,harrington2001,stratton1941,collin2001,devaney2004}.
\textbf{Step 2.} Eq.~\eqref{eq:ch4.mode-count-area} \cite{harrington2001,stratton1941,devaney2004,hansen1988,jackson1999,collin2001}.
\textbf{Step 3.} Measure \(d_{\mathrm{eff}}\) on hardware \cite{stratton1941,harrington2001,hansen1988,devaney2004,jackson1999,collin2001}.
\textbf{Physical meaning.} Band and width set budget \cite{harrington2001,stratton1941,hansen1988,devaney2004,jackson1999,collin2001}.

\paragraph{Eighteenth problem (mmWave band).}
\textbf{Problem.} \(28\,\mathrm{GHz}\), same \(W\). Compare.
\textbf{Step 1.} \(k_{0}W\) scales with frequency \cite{jackson1999,hansen1988,harrington2001,stratton1941,collin2001,devaney2004}.
\textbf{Step 2.} Eq.~\eqref{eq:ch4.mode-freq-scaling} \cite{harrington2001,stratton1941,devaney2004,hansen1988,jackson1999,collin2001}.
\textbf{Step 3.} Rank rises with electrical size \cite{stratton1941,harrington2001,hansen1988,devaney2004,jackson1999,collin2001}.
\textbf{Physical meaning.} mmWave enlarges mode budget \cite{harrington2001,stratton1941,hansen1988,devaney2004,jackson1999,collin2001}.

\paragraph{Nineteenth problem (vendor stream table).}
\textbf{Problem.} Tabulate claimed versus certified streams.
\textbf{Step 1.} Vendor claim \(N_{\mathrm{stream}}^{\mathrm{(claim)}}\) \cite{harrington2001,stratton1941,hansen1988,devaney2004,jackson1999,collin2001}.
\textbf{Step 2.} Lemma~\ref{lem:ch4.stream-bound} cap \cite{harrington2001,stratton1941,devaney2004,hansen1988,jackson1999,collin2001}.
\textbf{Step 3.} Flag over-claim \cite{stratton1941,harrington2001,hansen1988,devaney2004,jackson1999,collin2001}.
\textbf{Physical meaning.} Tables make audits transparent \cite{harrington2001,stratton1941,hansen1988,devaney2004,jackson1999,collin2001}.

\paragraph{Twentieth problem (complete scaling packet).}
\textbf{Problem.} Assemble mode-count audit for \(k_{0}W=4\), \(32\) elements.
\textbf{Step 1.} Eq.~\eqref{eq:ch4.ex-mode-count-numeric} \cite{harrington2001,collin2001,jackson1999,hansen1988,stratton1941,devaney2004}.
\textbf{Step 2.} Eq.~\eqref{eq:ch4.ex-mode-rank-ineq} with measured \(d_{\mathrm{eff}}\) \cite{harrington2001,stratton1941,hansen1988,devaney2004,jackson1999,collin2001}.
\textbf{Step 3.} Figures~\ref{fig:ch4.ex-mode-count-scaling} and \ref{fig:ch4.ex-element-rank-order} \cite{harrington2001,stratton1941,hansen1988,devaney2004,jackson1999,collin2001}.
\textbf{Step 4.} Deployable streams \(\le18\) \cite{stratton1941,harrington2001,hansen1988,devaney2004,jackson1999,collin2001}.
\textbf{Physical meaning.} Complete packets pair geometry, elements, and rank \cite{harrington2001,stratton1941,hansen1988,devaney2004,jackson1999,collin2001}.

\paragraph{Closure.}
At \(k_{0}W=4\), geometry permits \(\mathcal{O}(50)\) listed modes but measured export rank \(\approx18\): realizability is geometry-imposed rank, not patch count
\cite{stratton1941,harrington2001,hansen1988,jackson1999,devaney2004,collin2001}. Beamforming claims must cite \(N_{\mathrm{mode}}\), \(N_{\mathrm{design}}\), and \(d_{\mathrm{eff}}\)
\cite{harrington2001,stratton1941}.

\paragraph{Validity.}
Example~\ref{sec:ch4104} assumes square aperture, uniform exterior, and \(\varepsilon=10^{-2}\) \cite{stratton1941,hansen1988,harrington2001,jackson1999,devaney2004,collin2001}.
Mutual coupling and dispersion modify prefactors \(C_{A}\) without changing area-law scaling exponent \cite{harrington2001,devaney2004,hansen1988}.


The four worked examples audit maximum supported order on a disk, construct unreachable coefficient vectors, separate near-field evanescent richness from
far-zone export rank, and bound array diversity by mode-count scaling \cite{stratton1941,harrington2001,hansen1988,devaney2004,jackson1999,collin2001}.
Section~\ref{sec:ch4110} records the chapter inventory and forward map to Volume~II radiation mechanisms \cite{harrington2001,devaney2004}.

\section{Chapter Summary and Forward Map}
\label{sec:ch4110}


\subsection{Results established in Chapter 4}
\label{sec:ch4111}
% TOC tag:


\subsection{Concepts deferred to Volume II (Chapter 5 onward)}
\label{sec:ch4112}
% TOC tag:


\subsection{Bridge to canonical radiation mechanisms}
\label{sec:ch4113}
% TOC tag:


\bibliographystyle{IEEEtran}
\bibliography{references}
